\documentclass[11pt]{amsart}
\usepackage[margin=1.05in]{geometry}
\usepackage{amsmath,amssymb,mathrsfs}
\usepackage{xcolor}
\usepackage{hyperref}
\hypersetup{hidelinks,pdftitle={The field-strength Osterwalder--Schrader axioms for QED in three dimensions},pdfauthor={}}
\allowdisplaybreaks
\numberwithin{equation}{section}
\newtheorem{theorem}{Theorem}[section]
\newtheorem{proposition}[theorem]{Proposition}
\newtheorem{lemma}[theorem]{Lemma}
\newtheorem{corollary}[theorem]{Corollary}
\theoremstyle{definition}
\newtheorem{definition}[theorem]{Definition}
\theoremstyle{remark}
\newtheorem{remark}[theorem]{Remark}
\newcommand{\R}{\mathbb R}
\newcommand{\C}{\mathbb C}
\newcommand{\Z}{\mathbb Z}
\newcommand{\N}{\mathbb N}
\newcommand{\E}{\mathbb E}
\newcommand{\Sspace}{\mathcal S}
\newcommand{\cS}{\mathcal S}
\newcommand{\dd}{\mathrm d}
\DeclareMathOperator{\supp}{supp}
\DeclareMathOperator{\tr}{tr}
\DeclareMathOperator{\diag}{diag}
\title[Field-strength OS axioms for QED in three dimensions]{The field-strength Osterwalder--Schrader axioms\newline for QED in three dimensions}
\author{}
\date{}
\begin{document}
\begin{abstract}
We construct non-Gaussian laws of closed two-form fields from noncompact
lattice QED with chamber fermions. Their Schwinger functions satisfy the
Osterwalder--Schrader axioms, including clustering.
\end{abstract}
\maketitle
\setcounter{tocdepth}{1}
\tableofcontents

\section{Introduction}

The observable in this paper is the field strength. We regard it as a real
tempered two-form on $\R^3$, and write $F(f)$ for its pairing with a test
two-form $f$. The regularization uses a noncompact gauge field and
four-component chamber fermions, with all flux sectors included.
The order of limits is part of the construction: the volume tends to
infinity at a fixed lattice spacing, and the spacing then tends to zero.

\begin{theorem}[Field-strength OS theorem]\label{thm:main}
Fix $0<\rho<1$ and $0<\bar m\le1$, the chamber Wilson parameter and
reference mass parameter in physical units. Use the standard chamber weights
$(6,6\sqrt2,6\sqrt2,6)$. There are an odd scale factor $L\ge5$,
numbers $e_{\mathrm{phys}},\sigma>0$, and real volume-independent bare mass shifts
$m_{0,N}(e)$ for $N\in\Z_{\ge0}$ and $0<e\le e_{\mathrm{phys}}$, with
$|m_{0,N}(e)|<1/8$, with the following property.

Fix a real $\beta\in C_c^\infty(\R^3;\Lambda^2)$ such that
$f=dd^*\beta$ has
\begin{equation}\label{eq:positive-coefficient}
 b(f):=-6\tr\bigl[(S_{\bar m}\,i\gamma\cdot d^*\beta)^4\bigr]>0,
 \qquad S_{\bar m}=(\bar m+\gamma\cdot\partial)^{-1}.
\end{equation}
Such probes exist. There is $e_\beta\in(0,e_{\mathrm{phys}}]$ such that, for $0<e\le e_\beta$,
the all-flux chamber laws on the coordinate cubic tori of unit
side $n=R/\eta$, with periodic gauge fields and fermions
antiperiodic in all three coordinate directions, and with
\begin{equation}\label{eq:cutoff-scales}
 \eta=L^{-N},\quad R=8hL^T,\quad
 h=L^{\lfloor\log_L\bar m^{-1}\rfloor},\quad
 q=e\eta^{1/2},\quad \boldsymbol\mu_N=\bar m\eta+m_{0,N}(e)
\end{equation}
with $N,T\in\Z_{\ge0}$ have volume-first limit points.
Let $\mu$ be any such limit point and
$\mu=\int\mu_\xi\,\dd\nu(\xi)$ its decomposition into
translation-ergodic laws. For $\nu$-almost every $\xi$, $\mu_\xi$ is
carried by closed tempered two-forms and satisfies
\begin{equation}\label{eq:main-domination}
 \int e^{F(g)}\,\dd\mu_\xi(F)\le e^{\|g\|_{L^2}^2/2}
 \qquad(g\in\Sspace(\R^3;\Lambda^2)\text{ real}).
\end{equation}
Its Schwinger functions
$S_n^\xi(g_1\otimes\cdots\otimes g_n)
=\int\prod_{a=1}^nF(g_a)\,\dd\mu_\xi(F)$
extend to the complex tensor Schwartz spaces
$\Sspace_n=\Sspace(\R^{3n};(\Lambda^2\C^3)^{\otimes n})$
and obey the following conditions, with $S_0^\xi=1$.
\begin{enumerate}
\item[\textup{(E0$'$)}] For $n\ge1$ and $\Psi\in\Sspace_n$,
\begin{equation}\label{eq:main-growth}
 |S_n^\xi(\Psi)|\le2\,66^n\sqrt{n!}\,q_{6n}(\Psi),\qquad
 q_m(\Psi)=\max_{|\alpha|\le m}\sup_y
 (1+|y|^2)^{m/2}|\partial^\alpha\Psi(y)|,
\end{equation}
where the maximum is also over tensor components.
\item[\textup{(E1)}] For every Euclidean motion $g(x)=Qx+a$,
$Q\in O(3)$, $S_n^\xi(g\cdot\Psi)=S_n^\xi(\Psi)$, where
\[
 (g\cdot\Psi)(x_1,\ldots,x_n)
 =(\Lambda^2Q)^{\otimes n}\Psi(g^{-1}x_1,\ldots,g^{-1}x_n).
\]
\item[\textup{(E2)}] Let $r$ be reflection in a plane through the origin
and $H$ either closed half-space bounded by that plane.
If $\Psi_0\in\C$ and $\Psi_k\in\Sspace_k$, $1\le k\le M$,
have $\supp\Psi_k\subset H^k$, then
\begin{equation}\label{eq:main-rp}
 \sum_{k,l=0}^M S_{k+l}^\xi
 \bigl(\Theta_k\Psi_k\otimes\Psi_l\bigr)\ge0,
 \qquad \Theta_k\Psi_k=\overline{r\cdot\Psi_k},\quad
 \Theta_0c=\overline c.
\end{equation}
\item[\textup{(E3)}] $S_n^\xi$ is unchanged by simultaneous permutation
of the arguments and their tensor indices.
\item[\textup{(E4)}] For $\Psi\in\Sspace_k$ and $\Phi\in\Sspace_l$,
\begin{equation}\label{eq:main-cluster}
 S_{k+l}^\xi(\Psi\otimes\tau_a^{\otimes l}\Phi)
 \longrightarrow S_k^\xi(\Psi)S_l^\xi(\Phi)
 \quad\text{as }|a|\longrightarrow\infty,
\end{equation}
where
\[
 (\tau_a^{\otimes l}\Phi)(x_1,\ldots,x_l)
 =\Phi(x_1-a,\ldots,x_l-a).
\]
\end{enumerate}
Moreover, for a constant $C_f$ independent of the limit point and of
$\xi$,
\begin{equation}\label{eq:main-cumulant}
 \bigl|\kappa_4^{\mu_\xi}(f)-e^4b(f)\bigr|
 \le C_f e^{4+\sigma},\qquad
 \kappa_4^{\mu_\xi}(f)\ge\tfrac12e^4b(f)>0.
\end{equation}
In particular, $\mu_\xi$ is non-Gaussian.
\end{theorem}

Here $\kappa_4^\lambda(f)$ is the fourth cumulant of $F(f)$ under a
law $\lambda$. For the Gaussian-dominated laws considered here,
$F(f)$ has mean zero, so
$\kappa_4^\lambda(f)=S_4^\lambda(f^{\otimes4})
-3S_2^\lambda(f^{\otimes2})^2$.
The charge threshold may depend on the fixed probe.
The parent law $\mu$ satisfies
\textup{(E0$'$)}--\textup{(E3)}. The factorization in
\textup{(E4)} is a conclusion for its ergodic components.

The tensor action in \textup{(E1)} is essential to
\textup{(E2)}. For reflection in $x_1=0$, the components $F_{12}$ and
$F_{13}$ change sign, while $F_{23}$ does not. Replacing this action
by scalar reflection changes the positivity assertion.
The growth condition is stated on all tensor tests, including tests
meeting the diagonals; bounds on products of separately smeared fields
alone do not state \textup{(E0$'$)}.

\subsection{The two parts of the proof}

There are two principal tasks. The first is to obtain a Euclidean law
from the chamber regularization. The second is to show that the law
is non-Gaussian. Gaussian domination and reflection positivity supply
the compactness and analytic bounds used in the first task. The
renormalization analysis controls the fourth cumulant in the second.

The chamber sites lie away from the lattice mirrors. Every bond
crossing a mirror is normal to it, and its fermionic crossing
coefficients are $w_i(1\pm\rho)/2$. This gives reflection positivity
on tori adapted to the mirror. Coordinate tori are adapted to the
axis mirrors, but not to the diagonal mirrors. The comparison of
these geometries is therefore a substantive step. The fold inequality
bounds a cubic-torus expectation by an adapted-torus expectation plus
a vanishing pressure difference divided by the perturbation strength.
The volume limit must therefore precede removal of the perturbation.
Amplification controls local insertions; a two-layer reduction and
elimination then prove reflection and translation invariance of the
comparison states. Here the two-layer reduction is obtained by
eliminating the outer layers, using the fermion equations of motion
and an imaginary shift of the gauge integration contour; the remaining
error tends to zero with the perturbation strength.
Translation removes the restriction to an
observable inside one cell, and convex separation gives diagonal
positivity of the original limit. Taking volume before spacing
preserves this input in the continuum.

The nine mirror directions do not generate all rotations. Their
reflection-positive reconstructions give instead the sharp bound
$H_j\ge|\mathsf P_i|$, $i\ne j$. This spectral bound and Gaussian domination
control correlations in a complex tube after smearing in a transverse
coordinate. At two points their continuation in the rotation angle
$t$ is entire, with growth at most $C\exp(2|\operatorname{Im}t|)$.
Only angular frequencies of modulus at most two can occur, and the
Bianchi identity eliminates the noncovariant terms.

At higher orders let $\mathfrak a(n)$ measure the largest deviation
from the required tensor rotation weight. Reflection positivity gives
convexity of this non-negative integer index; positivity on regular
polygons gives a divisibility condition. With the linear angular bound,
these yield
\[
 \mathfrak a(2k)=k\mathfrak a(2),\qquad
 \mathfrak a(n)\le\tfrac n2\mathfrak a(2).
\]
The two-point result gives $\mathfrak a(2)=0$, so every order has the
required rotation weight. Covariance is first obtained for separated planar projections;
an analytic extension across coincidences of projected coordinates
gives it for disjoint supports. Locality and vanishing on gradients
then exclude contact terms, allowing passage to general tests.
This proves full covariance of the law, and hence reflection
positivity at every plane.

For clustering, the sharp cone identifies invariance under one
coordinate translation with invariance under all translations.
The local correlation estimate then yields factorization in an
ergodic law. The required bounds and mirror positivity are first
proved for almost every component; the rotation argument is applied
to each such component. Rotation invariance of the parent mixture
alone would not justify this step.

The non-Gaussianity argument uses only two reference volumes for its
renormalization estimates. The chamber operator has a four-component
light sector and a gapped heavy sector. Exponential decay and H\"older
bounds for its local inverses control the summability and locality of
the polymer activities. Extraction of the vacuum and mass terms leaves
an irrelevant activity that contracts under block averaging; the tuned
flow gives a mass assignment independent of volume.
The reference analysis controls fourth-order
response to an external connection. Its transfer to cutoff volumes
gives, beyond charge-dependent cutoff thresholds,
\[
 \bigl|\kappa_4(X_{\eta,R}(f))-e^4b(f)\bigr|
 \le C_f\bigl(e^{4+\sigma}+\eta^2+e^4\eta^{1/4}\bigr),
\]
uniformly in the volume; $X_{\eta,R}(f)$ is the lattice field smeared
with $f$. The last term includes the chamber-to-continuum error in the
free fourth coefficient. A separate continuum calculation produces a
positive probe: writing $a=d^*\beta$, its large-mass quartic term is a positive multiple of
$\int|\operatorname{curl}a|^4$, and dilation gives positivity at the
fixed mass $\bar m$.
Chessboard estimates also transfer products
of translated Wick powers. These product
bounds control invariant conditional expectations and are what makes
the cumulant estimate hold in almost every component. A fourth-cumulant
bound on the parent law by itself would be insufficient.

\subsection{Relation to earlier work}\label{subsec:earlier-work}

Theorem~\ref{thm:main} states the tensor form of the axioms of
Osterwalder and Schrader \cite{OS1973,OS1975}.
For reflection-positive lattice gauge theories and their transfer
matrices, see Osterwalder and Seiler \cite{OsterwalderSeiler1978};
the general chessboard framework is developed by
Fr\"ohlich, Israel, Lieb, and Simon \cite{FILS1978}.
The construction-specific steps here are the chamber crossings,
the comparison that gives diagonal positivity for the original
limiting law, and the angular argument for covariance at every order.

Dimock \cite{DimockUVS2022} proves ultraviolet stability at weak
coupling in a fixed finite volume using Ba{\l}aban's renormalization
group. Our block-integration analysis adapts this framework to
chamber fermions and the two reference tori used in the field-strength
estimates. The volume is then taken to infinity before the spacing
is removed; the conclusion concerns field-strength laws and their
fourth cumulants.

\subsection{Conventions}

The components of a two-form are indexed by $1\le i<j\le3$; their
$L^2$ inner product is the sum of the three component inner products.
We complexify test spaces when needed, while the random field is real.
The adjoint $d^*$ is taken in this Euclidean inner product.
The matrices $\gamma_i=\diag(\sigma_i,-\sigma_i)$ use the Pauli
matrices, and satisfy $\gamma_i\gamma_j+\gamma_j\gamma_i=2\delta_{ij}$.

For a Euclidean motion $g(x)=Qx+a$, its action on a one-argument
test is $(g\cdot f)(x)=(\Lambda^2Q)f(g^{-1}x)$ and on a
distribution is $(gF)(f)=F(g^{-1}\cdot f)$.
Translations follow this convention:
$(\tau_a f)(x)=f(x-a)$. Observable reflection includes complex
conjugation. The Hodge dual of a two-form transforms as an axial
vector, by $\det(Q)Q$; this convention will also be used in the
rotation argument.

\subsection{Scales and notation}

The physical and level coordinates have different unit lengths.
The following dictionary collects the conventions used in the
renormalization estimates.
\begin{center}
\begin{tabular}{@{}p{0.25\linewidth}p{0.69\linewidth}@{}}
Physical cutoff &
$\eta=L^{-N}$, $R=8hL^T$ and unit-lattice side $n=R/\eta$.\\[4pt]
Level $j$ &
$a_j=L^{j-N}$, $\eta_j=\eta/a_j$ and $\ell_j=R/a_j$;
one level unit has physical length $a_j$.\\[4pt]
Charges &
$e$ is the physical charge, $q=e\sqrt\eta$ the fine unit
charge, and $e_j=e\sqrt{a_j}$ the running charge.\\[4pt]
Masses &
$\mu_j=\bar m a_j$ is the positive reference mass and $m_j$
the mass interaction. The total fine mass is
$\boldsymbol\mu_N=\mu_0+m_{0,N}$, which is real.\\[4pt]
Dirac operators &
$\mathcal D_A$ is the finite matrix including the mass;
$D^{\mathrm{ch}}[B]$ is the unit-mesh massless matrix with
bond phase $B$; $D_A$ is the massless operator at the level
mesh; $D_j(A)$ is the blocked coarse operator.\\[4pt]
Curvatures &
$F^{\mathrm u}$ is the unit lattice curvature, $F^\sharp$
its rescaled step field, and $F$ the limiting tempered
two-form; see Section~\ref{sec:continuum-limits}.\\[4pt]
Reference cells &
$R^*$ is the reference side, $\ell_c=R^*/(2L)$ its cell
side, and $N_*=8L^3$ the number of cells.\\
\end{tabular}
\end{center}

\section{The chamber regularization}\label{sec:regularization}

The photon field is a real one-cochain on the cubic lattice. The fermion
has one site in each chamber of the lattice mirrors. This choice has a
specific consequence: a fermion bond crossing a mirror joins a site to
its reflected site and is perpendicular to the mirror. The hopping term
then admits a positive reflected expansion. We prove this expansion and
the resulting finite-volume reflection positivity below. The construction
also requires a nonzero partition function when the bare mass has
undetermined sign. For this purpose we identify a nonzero coefficient
of the fermion determinant, independent of the mass.

Throughout this section, $0<\rho<1$, $q>0$, and the mass parameter $M$
is real. The hopping weights $w_0,\ldots,w_3$ are positive. The estimates
and identities below apply, in particular, to the standard weights
used in the continuum construction.

\subsection{The photon field and the flux sum}\label{subsec:reg-photon}

We orient each link $(x,\mu)$ from $x$ to $x+e_\mu$ and each plaquette
$(x;\mu\nu)$, $\mu<\nu$, by the ordered pair $(e_\mu,e_\nu)$.
The real $k$-cochains are denoted by $C^k$, and the integer cochains by
$C^k(\Z)$. The coboundary $d$ and its adjoint $d^*$ use these
orientations. Thus
\begin{equation}\label{eq:reg-coboundary}
 (dA)_{\mu\nu}(x)=A_\mu(x)+A_\nu(x+e_\mu)
                     -A_\mu(x+e_\nu)-A_\nu(x).
\end{equation}
An oriented link may also be written as a pair of its endpoints;
$A(y,x)=-A(x,y)$ and $A(x,x)=0$. Pullback by a lattice isometry includes
the orientation of the image cell. In particular, pullback commutes
with $d$. For a coordinate reflection, this gives a minus sign on
the curvature components containing the normal direction and a
plus sign on the tangential component. These are the two-form
reflection signs used throughout.

Let $\Gamma\subset\Z^3$ be a rank-three subgroup with
\[
 \min\{|g|_\infty:g\in\Gamma\setminus\{0\}\}\ge4.
\]
We write $C^k(\Gamma)$ for the cochains on $\Z^3/\Gamma$ and identify
them with their periodic extensions. The torus has
$V=[\Z^3:\Gamma]$ photon sites, $3V$ links and $3V$ plaquettes. The
inner product on real cochains is the sum over positively oriented
cells, and $\|\cdot\|$ denotes its norm.

A photon configuration is a class in
\begin{equation}\label{eq:reg-configuration}
 \operatorname{Conf}(\Gamma)
 =\{(A,z):A\in C^1(\Gamma),\ z\in C^2(\Gamma;\Z),\ dz=0\}/\!\sim,
 \qquad
 (A,z)\sim\left(A+\frac{2\pi}{q}k,z-dk\right),
\end{equation}
where $k\in C^1(\Gamma;\Z)$. Its curvature and link phases are
\begin{equation}\label{eq:reg-curvature}
 F=dA+\frac{2\pi}{q}z,
 \qquad u_b=\exp(iqA(b)).
\end{equation}
Both are independent of the representative, and $dF=0$. The integer
cochain is required to be closed; there is no restriction on its flux.
The curvature therefore remains a real, unbounded field.

Put $r_q=2\pi/q$. Each class has a unique representative with
$A(b)\in[0,r_q)$ at every link. The measure on the quotient is defined
by
\begin{equation}\label{eq:reg-configuration-integral}
 \int_{\operatorname{Conf}(\Gamma)}\Phi(A,z)\,\dd(A,z)
 =\sum_{\substack{z\in C^2(\Gamma;\Z)\\dz=0}}
       \int_{[0,r_q)^{3V}}\Phi(A,z)\,\dd A.
\end{equation}
This formula holds for nonnegative measurable functions and for
absolutely integrable complex functions. Any measurable fundamental
domain may be used instead. Indeed, the changes of representative
preserve Lebesgue measure in $A$ and counting measure in $z$.

We will repeatedly use the finite number
\begin{equation}\label{eq:reg-gaussian-sum}
 C_q(\varepsilon)=\sup_{y\in\R}\sum_{j\in\Z}
                  e^{-\varepsilon(y+r_qj)^2},\qquad \varepsilon>0.
\end{equation}
It satisfies
\[
 C_q(\varepsilon)\le2+\frac{q}{2\pi}\sqrt{\frac{\pi}{\varepsilon}}.
\]
To see this, retain the two terms nearest the origin and bound each
remaining term by $r_q^{-1}$ times the Gaussian integral over the
adjacent interval of length $r_q$ closer to the origin. Those intervals
have disjoint interiors. Write $C_q=C_q(1/2)$. Dropping the constraint
$dz=0$ in the flux sum gives
\begin{equation}\label{eq:reg-maxwell-integral}
 0<\int_{\operatorname{Conf}(\Gamma)}e^{-\|F\|^2/2}\,\dd(A,z)
 \le(r_qC_q)^{3V}.
\end{equation}
The lower bound follows already from $z=0$; for the upper bound, sum
each $z_\sigma$ separately at fixed $A$.

The periodic extensions of the curvatures lie in
\[
 \Omega=\{F\in C^2:dF=0\},
\]
with the topology of pointwise convergence on plaquettes. This is a
closed subspace of a countable product of copies of $\R$. Local
observables depend on finitely many plaquette values. For a finitely
supported real two-cochain $\beta$, the character
$\exp(i\langle\beta,F\rangle)$ is bounded and continuous on $\Omega$.

\subsection{Chambers and the fermion graph}\label{subsec:reg-chambers}

The lattice mirrors are
\begin{equation}\label{eq:reg-mirrors}
 \{x_\mu=c\},\quad c\in\tfrac12\Z,
 \qquad
 \{x_\mu\pm x_\nu=c\},\quad c\in\Z,\quad\mu<\nu.
\end{equation}
They form three coordinate families and six diagonal families. Let
$\mathcal L$ be the corresponding linear forms $\ell$ and write
$\ell(x)=n_\ell\cdot x$. The admissible positions are
$\operatorname{Adm}(x_\mu)=\frac12\Z$ and
$\operatorname{Adm}(x_\mu\pm x_\nu)=\Z$. The reflection in a mirror
$\Pi=(\ell,c)$ is
\begin{equation}\label{eq:reg-reflection}
 \theta_\Pi x=x-\frac{2(\ell(x)-c)}{|n_\ell|^2}n_\ell.
\end{equation}
Its upward normal is $n_\ell/|n_\ell|$. Coordinate mirrors at integer
positions are called site planes, and those at half-integer positions
are called link planes.

The group generated by these reflections is
\[
 W=\Z^3\rtimes O_{\mathrm{cub}},
 \qquad O_{\mathrm{cub}}=O(3)\cap\operatorname{GL}(3,\Z).
\]
Here $O_{\mathrm{cub}}$ consists of the $48$ signed permutation matrices.
A fundamental chamber is
\[
 C_0=\{0<x_3<x_2<x_1<\tfrac12\}.
\]
Its wall reflections are
\begin{align*}
 s_0(x)&=(x_1,x_2,-x_3),&
 s_1(x)&=(x_1,x_3,x_2),\\
 s_2(x)&=(x_2,x_1,x_3),&
 s_3(x)&=(1-x_1,x_2,x_3).
\end{align*}
Denote the walls by $\Pi_i$ and their inward unit normals by
\begin{equation}\label{eq:reg-wall-normals}
 n_0=e_3,\quad n_1=(e_2-e_3)/\sqrt2,
 \quad n_2=(e_1-e_2)/\sqrt2,\quad n_3=-e_1.
\end{equation}

Fix $v_0=(v_1,v_2,v_3)\in C_0$. The fermion sites and their neighbours
are
\begin{equation}\label{eq:reg-fermion-sites}
 \mathcal X=Wv_0,\qquad (gv_0)^{(i)}=gs_iv_0.
\end{equation}
The unordered bond $\{x,x^{(i)}\}$ has type $i$. A site has the unique
form $x=k+rv_0$, with $k\in\Z^3$ and $r\in O_{\mathrm{cub}}$; put
$\pi(x)=k$ and $y(x)=rv_0$. Thus $\pi(x)$ is its nearest photon site.

\begin{proposition}\label{prop:reg-geometry}
The chambers are precisely $gC_0$, $g\in W$, and $W$ acts on them
simply transitively. No fermion site lies on a mirror. Every bond
crosses exactly one mirror, its chamber wall, and crosses it
orthogonally at its midpoint. Consequently a bond crossing $\Pi$ is
$\{x,\theta_\Pi x\}$.

The projection $\pi$ is $W$-equivariant. Bonds of types $0,1,2$ have
the same photon site at both ends; a type-$3$ bond joins two nearest
photon sites. Exactly eight fermion bonds lie over each photon link.
If a mirror does not cut the link and its open segment lies on one
side, all eight bonds lie on that side.
The graph is connected and bipartite, and each site has one neighbour
of each type.
\end{proposition}

\begin{proof}
Coordinate reflections and coordinate transpositions generate
$O_{\mathrm{cub}}$, and
$\theta_{x_\mu,1/2}\theta_{x_\mu,0}=\tau_{e_\mu}$ gives the
translations. Conversely, \eqref{eq:reg-reflection} is an affine signed
permutation with integer translation, so the mirror group is $W$.

For a point on no mirror, subtract the unique nearest lattice point.
The resulting three coordinates are nonzero, have distinct absolute
values, and lie in $(-1/2,1/2)$. There is exactly one signed
permutation making them positive and decreasing. This places the
point in exactly one translate of $C_0$. These translates are disjoint
open convex sets, meet no mirror, and cover its complement. They are
therefore the chambers, with the asserted free transitive action.

For the four bonds from $v_0$, the midpoint is, respectively,
\[
 (v_1,v_2,0),\quad
 (v_1,(v_2+v_3)/2,(v_2+v_3)/2),\quad
 ((v_1+v_2)/2,(v_1+v_2)/2,v_3),\quad
 (1/2,v_2,v_3).
\]
Each belongs to the relative interior of its wall and to no other
mirror. Each half of the bond, except its midpoint, lies in one of
the two adjacent chambers. Transport by $g\in W$ proves the bond
assertion at every site.

In the coordinates $x=k+rv_0$, the first three neighbours are
$k+rs_iv_0$, whereas the fourth is
$k+re_1+rr_{e_1}v_0$, where $r_{e_1}$ changes the sign of the first
coordinate. This proves the statements about $\pi$ and its
equivariance. Over $(k,\mu)$ the lower ends are exactly the eight
sites for which $y_\mu=v_1$: the remaining two absolute values can
be ordered in two ways, with four choices of their signs. Their
type-$3$ neighbours are the upper ends of these eight bonds.

The bond over $b=(k,\mu)$ has the segment
$k+te_\mu+y_\perp$, $v_1\le t\le1-v_1$, where the other two
coordinates of $y_\perp$ have absolute values $v_2,v_3$. For a
diagonal height involving $x_\mu$, its varying part is
$t\pm y_\nu\in(0,1)$, the same interval between integers as on the
open photon link. For heights not involving $x_\mu$, the changes
are $y_\nu$ or $y_\nu\pm y_{\nu'}$, whose absolute values are
smaller than $1/2$ or $1$, respectively. Away from a mirror
containing $b$ these preserve its side. The coordinate height
$x_\mu$ gives the same side except at the unique link plane through
the midpoint, which cuts $b$. This proves the assertion about
photon links.

The generators $s_0,s_1,s_2$ generate the signed permutations, and
$s_3r_{e_1}=\tau_{e_1}$; conjugation gives all coordinate translations.
Thus the four wall reflections generate $W$, proving connectivity.
The two parts are the sites $gv_0$ with $\det r=1$ and with
$\det r=-1$, since each neighbour changes this determinant.
\end{proof}

The geometric location is useful for reflection, while the coordinates
$(k,r)$ are useful for the renormalisation analysis. If
$d_i=\operatorname{dist}(v_0,\Pi_i)$, then
\[
 x^{(i)}-x=-2d_i\,rn_i,
 \qquad \frac{x^{(i)}-x}{|x^{(i)}-x|}=-rn_i.
\]
The directions, bond types and projected endpoints in these
coordinates do not depend on $v_0$. Hence neither the operator below
nor its determinant depends on $v_0$ after this identification.
The standard weights are
\begin{equation}\label{eq:reg-standard-weights}
 w_i=48|\Pi_i\cap\overline C_0|,
 \qquad (w_0,w_1,w_2,w_3)=(6,6\sqrt2,6\sqrt2,6).
\end{equation}
Their inward normals satisfy $\sum_iw_in_i=0$.

\subsection{The chamber operator and the finite-volume law}
\label{subsec:reg-law}

Let $\sigma_\mu$ be the Pauli matrices and
\[
 \gamma_\mu=\diag(\sigma_\mu,-\sigma_\mu),\qquad
 \gamma_v=\sum_\mu v_\mu\gamma_\mu.
\]
They form a Hermitian Clifford triple. For a unit vector $n$, put
$P_n^\pm=(1\pm\gamma_n)/2$. Both projections have rank two, and
\begin{equation}\label{eq:reg-clifford-reflection}
 \gamma_n\gamma_v\gamma_n=-\gamma_{r_nv},\qquad
 r_nv=v-2(v\cdot n)n.
\end{equation}
This follows by multiplying
$\gamma_n\gamma_v+\gamma_v\gamma_n=2n\cdot v$ by $\gamma_n$.

For an oriented bond $(x,x')$ of type $i$ define
\begin{equation}\label{eq:reg-hop}
 U_A(x,x')=\exp(iqA(\pi(x),\pi(x'))),\qquad
 w(x,x')=\frac{w_i}{2}\left(\rho-\gamma_{(x'-x)/|x'-x|}\right).
\end{equation}
The chamber operator is
\begin{equation}\label{eq:reg-dirac}
 (\mathcal D_A\psi)(x)=M\psi(x)-\sum_{i=0}^3
                   w(x,x^{(i)})U_A(x,x^{(i)})\psi(x^{(i)}).
\end{equation}
Its lattice mass is
\begin{equation}\label{eq:reg-lattice-mass}
 m=M-\frac\rho2\sum_iw_i.
\end{equation}
For the standard weights, the weighted sum of the Clifford terms in
the differences $\psi(x)-U_A(x,x^{(i)})\psi(x^{(i)})$ vanishes on the
diagonal. Thus \eqref{eq:reg-dirac}, in coordinates $(k,r)$, is exactly
the matrix $D^{\mathrm{ch}}[qA]+m$, where
$D^{\mathrm{ch}}[B]$ is the unit-mesh massless chamber
matrix with external hopping phase $\exp(iB(b))$.
At a continuum cutoff this
means $m=\bar m\eta+m_{0,N}(e)$ and
$M=\bar m\eta+m_{0,N}(e)+\frac\rho2\sum_iw_i$.

A fermion boundary condition is a character
$\chi:\Gamma\to\{1,-1\}$. Write $\mathbb T=(\Gamma,\chi)$.
Its spinors satisfy $\psi(x+g)=\chi(g)\psi(x)$ for $g\in\Gamma$.
{\color{red}The restriction $\mathcal D_A^{\mathbb T}$ acts on a space of dimension $192V$,
with the inner product obtained by summing over representatives of
$\mathcal X/\Gamma$. For each site introduce generators
$\psi_\alpha(x),\bar\psi_\alpha(x)$, $1\le\alpha\le4$, obeying the same
boundary conditions.} Their Grassmann algebra is denoted by
$\operatorname{Gr}_{\mathbb T}$. Its Berezin integral is normalized by
\begin{equation}\label{eq:reg-berezin}
 \int_{\mathbb T}\prod_{x\in\mathcal X/\Gamma}
                  \prod_{\alpha=1}^4\psi_\alpha(x)\bar\psi_\alpha(x)=1.
\end{equation}
Products of pairs are even, so their ordering is immaterial. With
$S_{\mathbb T}(A)=\sum_x\bar\psi(x)(\mathcal D_A\psi)(x)$, this normalization
gives $\int_{\mathbb T}e^{-S_{\mathbb T}(A)}=\det \mathcal D_A^{\mathbb T}$.

The unnormalized functional is
\begin{equation}\label{eq:reg-functional}
 Z_{\mathbb T}[G]=\int_{\operatorname{Conf}(\Gamma)}
      e^{-\|F\|^2/2}\int_{\mathbb T}e^{-S_{\mathbb T}(A)}G(A,z)
      \,\dd(A,z).
\end{equation}
For a scalar $G$ the inner integral is
$G(A,z)\det \mathcal D_A^{\mathbb T}$. For a Grassmann-valued integrand,
absolute convergence is understood coefficient by coefficient.
Bounded measurable coefficients suffice: $e^{-S_{\mathbb T}(A)}$
has finitely many coefficients, each a polynomial in the link phases,
and \eqref{eq:reg-maxwell-integral} supplies an integrable majorant.

\begin{proposition}\label{prop:reg-determinant-sign}
For every real $A$ and $M$,
\[
 \det \mathcal D_A^{\mathbb T}\ge0,
 \qquad
 \|\mathcal D_A^{\mathbb T}\|\le |M|+\frac{1+\rho}{2}\sum_iw_i.
\]
Moreover $\Re \mathcal D_A^{\mathbb T}=m+\frac\rho2\nabla_A^*\nabla_A$, where
\[
 (\nabla_A\psi)(x,i)=\sqrt{w_i/2}\,
           \bigl(\psi(x)-U_A(x,x^{(i)})\psi(x^{(i)})\bigr).
\]
The determinant is not identically zero as a function of the link
phases, for every real mass parameter $M$.
\end{proposition}

\begin{proof}
Write
\[
 H_A\psi(x)=\sum_iw_iU_A(x,x^{(i)})\psi(x^{(i)}),\qquad
 J_A\psi(x)=\sum_iw_i\gamma_{(x^{(i)}-x)/|x^{(i)}-x|}
                         U_A(x,x^{(i)})\psi(x^{(i)}).
\]
Reversing a bond conjugates its phase and reverses its direction.
Reindexing by the involution $x\mapsto x^{(i)}$ therefore gives
$H_A^*=H_A$ and $J_A^*=-J_A$. Hence
\[
 \mathcal D_A=M-\frac\rho2H_A+\frac12J_A,\qquad
 \nabla_A^*\nabla_A=\sum_iw_i-H_A.
\]
This proves the real-part identity. In the splitting
$\C^4=\C^2\oplus\C^2$, the two blocks of $H_A$ coincide, and the two
blocks of $J_A$ have opposite signs. It follows that
\begin{equation}\label{eq:reg-two-blocks}
 \mathcal D_A^{\mathbb T}=\diag(\check{\mathcal D}_A^{\mathbb T},
                                      (\check{\mathcal D}_A^{\mathbb T})^*),
 \qquad \det \mathcal D_A^{\mathbb T}=|\det\check{\mathcal D}_A^{\mathbb T}|^2.
\end{equation}
Each hopping matrix has norm $w_i(1+\rho)/2$, while the associated
translation with phase is unitary. Their sum gives the norm bound.

For the last assertion we view the determinant as a Laurent
polynomial $P_{\mathbb T}$ in the independent variables $u_b$.
Only type-$3$ bonds contain a variable, with $u_b$ at the lower end
and $u_b^{-1}$ at the upper end. Fix a coordinate direction $\mu$.
The coefficient of
\begin{equation}\label{eq:reg-extreme-monomial}
 \prod_{b\parallel e_\mu}u_b^{32}
\end{equation}
will be computed without using $M$.

Partition the sites modulo $\Gamma$ into six layers according to
the increasing values
\[
 y_\mu=-v_1,-v_2,-v_3,v_3,v_2,v_1.
\]
Each layer has $8V$ sites. Each site has one upward neighbour,
one downward neighbour, and two neighbours in its own layer. The
upward map is a permutation of the sites and visits successive layers
modulo six. Its bond types are
\begin{equation}\label{eq:reg-upward-types}
 (i_1,i_2,i_3,i_4,i_5,i_6)=(2,1,0,1,2,3).
\end{equation}
For layers $1,\ldots,5$ the upward phase is one; in layer six it is
$u_b$ for the link $(\pi(x),\mu)$. These facts follow by changing the
smallest sign or exchanging successive absolute values in the
neighbour formulas above. The sixth step changes $v_1$ to $-v_1$
and increases the photon coordinate by one.

There are eight lower and eight upper fermion ends over each photon
link, each carrying four matrix rows. An exponent $32$ at a link
therefore requires every lower row to take the upward type-$3$
block and every upper row to avoid the reverse block. Imposing this
at all links in direction $\mu$ forces every row in every layer to
take its upward block. Indeed, layer-six rows fill all columns in
layer one. A layer-one row cannot then use a column in its own layer,
and the reverse block is excluded by the exponent condition; it must
go upward to layer two. Inductively the columns in the current and
preceding layers are already occupied, so layers two through five
must also go upward. The selected terms use no other link variables.

The remaining freedom is a permutation of the four spin indices at
each site. The site permutation, repeated four times, has sign one.
Summing over the spin permutations therefore gives the product of
the $4\times4$ determinants of the selected hopping blocks. The
boundary signs have fourth power one, and
\[
 \det w(x,x^\uparrow)
   =(w_{i_k}/2)^4(1-\rho^2)^2
\]
at a site in layer $k$. Consequently the coefficient of
\eqref{eq:reg-extreme-monomial} is $C_*^V$, where
\begin{equation}\label{eq:reg-extreme-coefficient}
 C_*=(w_0/2)^{32}(w_1/2)^{64}(w_2/2)^{64}(w_3/2)^{32}
                                           (1-\rho^2)^{96}>0.
\end{equation}
This coefficient is independent of $M$. In particular the determinant
polynomial, and its restriction to the phase torus, are nonzero.
\end{proof}

The real-part estimate alone would require control of the sign of
$m$ to prove invertibility. The coefficient computation instead
proves that the determinant has positive integral at every real $M$,
which is all that is needed to define the law.

\begin{proposition}\label{prop:reg-partition}
The partition function satisfies $0<Z_{\mathbb T}[1]<\infty$.
More precisely,
\begin{equation}\label{eq:reg-partition-bounds}
 r_q^{3V}e^{-6r_q^2V}C_*^V
 \le Z_{\mathbb T}[1]
 \le(r_qC_q)^{3V}\left(|M|+\frac{1+\rho}{2}\sum_iw_i\right)^{192V}.
\end{equation}
Thus $\E_{\mathbb T}[G]=Z_{\mathbb T}[G]/Z_{\mathbb T}[1]$ defines a
probability law on scalar observables. Its curvature law
$\mu_{\mathbb T}$ is the law of the periodic extension $\widetilde F$
on $\Omega$.
\end{proposition}

\begin{proof}
For a nonnegative Laurent polynomial $P$ on the phase torus, every
Fourier coefficient satisfies
\[
 |\widehat P(\lambda)|
 =\left|\int P(u)\overline{u^\lambda}\,\dd u\right|
 \le\int P(u)\,\dd u,
\]
where Haar measure is normalized. By the preceding proposition and
\eqref{eq:reg-extreme-coefficient},
\[
 \int_{[0,r_q)^{3V}}\det \mathcal D_A^{\mathbb T}\,\dd A
 \ge r_q^{3V}C_*^V.
\]
On this box $|(dA)_\sigma|\le2r_q$ and
$\|dA\|^2\le12Vr_q^2$. The $z=0$ term in
\eqref{eq:reg-functional} gives the lower bound in
\eqref{eq:reg-partition-bounds}. For the upper bound use
$|\det \mathcal D_A^{\mathbb T}|\le\|\mathcal D_A^{\mathbb T}\|^{192V}$ and
\eqref{eq:reg-maxwell-integral}. The density
$e^{-\|F\|^2/2}\det \mathcal D_A^{\mathbb T}$ is nonnegative. After division
by its strictly positive integral it is a probability density, whose
image under $(A,z)\mapsto\widetilde F$ is the asserted law.
\end{proof}

\subsection{Symmetries and reflected Grassmann products}
\label{subsec:reg-symmetries}

The field strength transforms as a two-cochain, with the orientation
of each plaquette included. For an antiperiodic cubic torus the law is invariant
under $W$ and under $F\mapsto-F$. We record the corresponding
fermionic identities, since reflection positivity uses their signs.

Gauge covariance follows directly from
\[
 U_{A-d\lambda}(x,x')
 =e^{iq\lambda(\pi(x))}U_A(x,x')e^{-iq\lambda(\pi(x'))}.
\]
For lattice covariance set
$\gamma_5=\left(\begin{smallmatrix}0&1_2\\1_2&0\end{smallmatrix}\right)$.
It anticommutes with the Clifford triple, and
\[
 (\gamma_n\gamma_5)\gamma_v(\gamma_n\gamma_5)^{-1}
 =\gamma_{r_nv}.
\]
If the linear part of $g\in W$ is $r_{n_1}\cdots r_{n_k}$, the
corresponding product of these spin matrices, followed by the
permutation $x\mapsto g^{-1}x$, intertwines $\mathcal D_A^{\mathbb T}$ and
$\mathcal D_{(g^{-1})^*A}^{g\mathbb T}$, where
$g\mathbb T=(r\Gamma,\chi\circ r^{-1})$ and $r$ is the linear part
of $g$. Taking determinants and changing
variables in \eqref{eq:reg-functional} proves
\[
 \mu_{g\mathbb T}=((g^{-1})^*)_*\mu_{\mathbb T}.
\]
The configuration measure is preserved because pullback permutes
cochain coordinates up to signs and maps one fundamental domain
to another. All translations preserve $\mathbb T$, and all elements
of $W$ preserve an antiperiodic cubic torus, proving the stated
invariance.

For charge conjugation let $C=\sigma_2\oplus\sigma_2$. The Pauli
identity $C\gamma_v^\top C^{-1}=-\gamma_v$, together with reversal
of the oriented bond, gives
\[
 C(\mathcal D_A^{\mathbb T})^\top C^{-1}=\mathcal D_{-A}^{\mathbb T}.
\]
Thus the determinant is unchanged by $A\mapsto-A$. The substitution
$(A,z)\mapsto(-A,-z)$ preserves the configuration measure and sends
$F$ to $-F$, proving charge symmetry of the curvature law.

At a mirror $\Pi$ with upward normal $n$, the Grassmann reflection
$\Theta_\Pi$ is antilinear and reverses products. In components it is
\begin{equation}\label{eq:reg-grassmann-reflection}
 \Theta_\Pi\psi_\alpha(x)
  =-\sum_\beta\bar\psi_\beta(\theta_\Pi x)(\gamma_n)_{\beta\alpha},
 \qquad
 \Theta_\Pi\bar\psi_\alpha(x)
  =-\sum_\beta(\gamma_n)_{\alpha\beta}\psi_\beta(\theta_\Pi x).
\end{equation}
On scalar coefficient functions it acts as
$\Phi(A,z)\mapsto\overline{\Phi(\theta_\Pi^*A,\theta_\Pi^*z)}$.
The Clifford identities give $\Theta_\Pi^2=1$ and
\begin{equation}\label{eq:reg-reflected-bilinear}
 \Theta_\Pi(\bar\psi(x)Y\psi(x'))
 =\bar\psi(\theta_\Pi x')\gamma_nY^*\gamma_n\psi(\theta_\Pi x).
\end{equation}
For a hopping term,
$\gamma_nw(x,x')\gamma_n=w(\theta_\Pi x',\theta_\Pi x)$, by
\eqref{eq:reg-clifford-reflection}. Reflection therefore permutes the
oriented terms of the action, including the real mass terms.

If $Y$ is a finite set of sites, then
\begin{equation}\label{eq:reg-reflected-berezin}
 \int_{\theta_\Pi Y}\Theta_\Pi\xi=\overline{\int_Y\xi}.
\end{equation}
Only the top-degree coefficient needs to be checked. At one site
reflection maps $\prod_\alpha\psi_\alpha\bar\psi_\alpha$ to the same
product at the reflected site, multiplied by $(\det\gamma_n)^2=1$.
Products from different sites are even. This proves
\eqref{eq:reg-reflected-berezin}, and a measure-preserving change of
photon variables gives
\begin{equation}\label{eq:reg-reflected-functional}
 Z_{\theta_\Pi\mathbb T}[\Theta_\Pi G]
 =\overline{Z_{\mathbb T}[G]}
\end{equation}
for bounded coefficient integrands, and for the integrable
Gaussian-growth classes considered below.

On even integrands, compositions of these reflections define an
action of $W$. The ambiguity of a reflection word acts only by the
common sign on all fermion generators. To verify this, compose two
words representing the same lattice isometry. The resulting even
word has trivial geometric action, and its spin matrix is
$\diag(u,u)$ with $u$ a product of an even number of Pauli matrices
$\sigma_n$. It commutes with all Pauli matrices, hence is scalar;
its determinant is one, so $u=\pm1_2$. Its action on an even monomial
is consequently the identity. We write $G^g$ for this action when
folding even integrands. For a real scalar observable it is
$G^g(F)=G(g^*F)$.

\subsection{The fermion crossing at a mirror}\label{subsec:reg-crossing}

Let a bond of type $j$ cross $\Pi$, with lower endpoint $x$ and upper
endpoint $x'=\theta_\Pi x$. Its direction is the upward normal $n$.
The two hopping matrices are
\begin{equation}\label{eq:reg-crossing-projectors}
 \begin{split}
 w(x,x')&=\tfrac{w_j}{2}\bigl((1+\rho)P_n^--(1-\rho)P_n^+\bigr),\\
 w(x',x)&=\tfrac{w_j}{2}\bigl((1+\rho)P_n^+-(1-\rho)P_n^-\bigr).
 \end{split}
\end{equation}
Choose orthonormal bases $e_{n,1}^\pm,e_{n,2}^\pm$ of the two
projected spaces. For each $k=1,2$ introduce four crossing modes.
Their coefficients and phases are displayed together:
\begin{equation}\label{eq:reg-crossing-modes}
\begin{array}{c|c|c|c}
 &\zeta&\omega&u\\ \hline
 1&(e_{n,k}^-)^*\psi(x')&\tfrac{w_j}{2}(1+\rho)&U_A(x,x')\\
 2&(e_{n,k}^+)^*\psi(x')&\tfrac{w_j}{2}(1-\rho)&U_A(x,x')\\
 3&\bar\psi(x')e_{n,k}^+&\tfrac{w_j}{2}(1+\rho)&\overline{U_A(x,x')}\\
 4&\bar\psi(x')e_{n,k}^-&\tfrac{w_j}{2}(1-\rho)&\overline{U_A(x,x')}
\end{array}
\end{equation}
All coefficients are positive in the range $0<\rho<1$.

\begin{proposition}\label{prop:reg-crossing-expansion}
Let $J$ be a finite collection of bonds crossing $\Pi$, and let
$S_\Pi^\times(J)$ be their part of the action. Index all their crossing
modes by a finite set $\mathcal I$. Then
\begin{equation}\label{eq:reg-crossing-positive}
 -S_\Pi^\times(J)=\sum_{i\in\mathcal I}\omega_i u_i\Theta_\Pi(\zeta_i)\zeta_i.
\end{equation}
Choose an order on $\mathcal I$ and, for $I\subset\mathcal I$, put
$\zeta_I=\prod_{i\in I}\zeta_i$, $\omega_I=\prod_{i\in I}\omega_i$
and $u_I=\prod_{i\in I}u_i$. Then
\begin{equation}\label{eq:reg-crossing-exponential}
 e^{-S_\Pi^\times(J)}=\sum_{I\subset\mathcal I}\omega_Iu_I\Theta_\Pi(\zeta_I)\zeta_I.
\end{equation}
At site and diagonal mirrors all $u_i$ are one. At a link plane,
groups one and two have charge $+1$, groups three and four charge
$-1$, and their phases are $e^{iq\mathfrak q(i)A(b)}$ at the cut
photon link $b$.
\end{proposition}

\begin{proof}
The reflections of the four modes in \eqref{eq:reg-crossing-modes}
are, respectively,
\[
 \bar\psi(x)e_{n,k}^-,\quad-\bar\psi(x)e_{n,k}^+,
 \quad-(e_{n,k}^+)^*\psi(x),\quad(e_{n,k}^-)^*\psi(x).
\]
Multiplying each reflected mode by the original mode and using
anticommutation gives the four projected terms in
\eqref{eq:reg-crossing-projectors}, with the indicated phases.
Summation over $k$ uses
$P_n^\pm=\sum_ke_{n,k}^\pm(e_{n,k}^\pm)^*$ and gives the negative of
the two bond terms. Summation over bonds proves
\eqref{eq:reg-crossing-positive}.

Each $\Theta_\Pi(\zeta_i)\zeta_i$ is even and has square zero.
These elements commute. Reversal of the order under $\Theta_\Pi$
gives
\[
 \prod_{i\in I}\Theta_\Pi(\zeta_i)\zeta_i
 =\Theta_\Pi(\zeta_I)\zeta_I,
\]
with no remaining sign. Expanding the product of the finite
exponentials proves \eqref{eq:reg-crossing-exponential}.
Only type-$3$ bonds carry a photon phase, and they cross link
planes. Their projected endpoints are the oriented endpoints of
the cut link, giving the last assertion.
\end{proof}

\subsection{Adapted tori and the second mirror}\label{subsec:reg-adapted}

An adapted torus for $\ell\in\mathcal L$ has the form
\begin{equation}\label{eq:reg-adapted-torus}
 \Gamma=\Gamma_\parallel\oplus\Z v,
 \quad \Gamma_\parallel=\Gamma\cap\{\ell=0\},
 \quad v=\frac{p}{|n_\ell|^2}n_\ell\in\Z^3,
 \quad \chi(v)=-1,
\end{equation}
where $p=\ell(v)>0$ is the height period. For an admissible $c$ set
\[
 c'=c+p/2,\quad \Pi=(\ell,c),\quad\Pi'=(\ell,c'),\quad
 U_+=\{c<\ell<c'\},\quad U_-=\{c-p/2<\ell<c\}.
\]
The two mirrors are admissible: for coordinate families $p\ge4$,
while for diagonal families $p$ is even and $p\ge8$. Their
reflections satisfy $\theta_{\Pi'}=\tau_v\theta_\Pi$ and induce the
same geometric reflection on the torus, exchanging its two halves.
An antiperiodic coordinate torus is adapted to each coordinate
family. The torus with periods $ae_1-ae_2$, $be_1+be_2$ and $ce_3$,
antiperiodic on these periods, is adapted to the corresponding two
diagonal families and the third coordinate family. Cubic images
give adapted tori for all nine families.

Choose representatives $\mathcal X_+$ in $U_+$ and
$\mathcal X_-=\theta_\Pi\mathcal X_+$ in $U_-$. They partition the
fermion sites of the torus. Let $S_+$ contain the mass terms and
oriented bonds entirely in $U_+$, and set
$\theta=\theta_\Pi$, $\Theta=\Theta_\Pi$; the second boundary mirror
induces the same geometric reflection on the torus.
The reflected part is $\Theta S_+$. Its coefficient functions
involve only links in the closed upper slab.

The crossing modes at $\Pi$ lie in
$\operatorname{Gr}(\mathcal X_+)$. At $\Pi'$ their upper endpoints
are outside $U_+$; write
$\check\zeta_{i'}=\Theta_{\Pi'}\zeta_{i'}$, so the checked modes
lie in $\operatorname{Gr}(\mathcal X_+)$. The boundary condition gives
\begin{equation}\label{eq:reg-second-reflection}
 \Theta_{\Pi'}=\chi(v)^{\deg}\Theta_\Pi
                 \quad\hbox{on }\operatorname{Gr}_{\mathbb T}.
\end{equation}
In the crossing term at $\Pi'$ we therefore have
\[
 \Theta_{\Pi'}(\zeta_{i'})\zeta_{i'}
 =-\chi(v)\Theta(\check\zeta_{i'})\check\zeta_{i'}.
\]
The minus sign is from exchanging the two odd factors. Since
$\chi(v)=-1$, it is canceled by antiperiodicity. Both boundary
crossings thus have the same positive coefficients, and
\begin{equation}\label{eq:reg-two-crossings}
 e^{-S_{\mathbb T}}=e^{-S_+}\Theta(e^{-S_+})
 \sum_{I,I'}\omega_I\omega_{I'}u_Iu_{I'}
       \Theta(\zeta_I)\zeta_I
       \Theta(\check\zeta_{I'})\check\zeta_{I'}.
\end{equation}
Here $I$ and $I'$ range over the finite crossing-mode sets at the
two mirrors.

The reflected Berezin integral separates the two halves:
\begin{equation}\label{eq:reg-half-berezin}
 \int_{\mathbb T}\Theta(\xi')\xi
 =\overline{\int_{\mathcal X_+}\xi'}
                  \int_{\mathcal X_+}\xi,
 \qquad \xi,\xi'\in\operatorname{Gr}(\mathcal X_+).
\end{equation}
Indeed, the top monomial is the product of the two half-top
monomials, and reflection maps one to the other by
\eqref{eq:reg-reflected-berezin}. We must still separate the photon
flux constraints; these couple the halves at the cut cubes.

\subsection{The photon crossing and Gauss' law}\label{subsec:reg-gauss}

Let $\mathcal M$ be the two boundary mirrors, understood periodically.
Classify a cell by its relative interior: it is upper or lower if
that interior is in the respective half, lies on the mirrors if the
whole cell does, and is cut if its relative interior meets a mirror
without being contained in it. Write $B_+,B_-,B_0,B_\times$ for the
four sets of links, $P_+,P_-,P_0,P_\times$ for the plaquettes, and
$Q_+,Q_-,Q_\times$ for the cubes.

Reflection exchanges the upper and lower cells. It fixes a cell
lying on the mirrors with orientation sign $+1$, and a cut cell
with orientation sign $-1$. Each cell meets at most one boundary
mirror: its height variation is at most $|n_\ell|^2$, smaller than
$p/2$. To justify the orientation assertion, a point in the
intersection of the relative interior with the mirror is fixed
and lies in the relative interior of the reflected cell. Relative
interiors of different cells are disjoint, so the cell is fixed.
Reflection is the identity on its affine hull in the first case
and a reflection of that affine hull in the cut case.

An upper cell has only upper faces or faces lying on the mirrors;
a cut cell has only upper, lower or cut faces. Moreover a plaquette
containing a cut link is cut, and a cube containing a cut plaquette
is cut. An upper face can touch the boundary at a relative interior
point only when its affine height is constant. A face of a cut cell
cannot lie in the mirror: reflection in the supporting plane of a
face would not preserve that cell. Finally, the height takes both
signs on a cut face, hence on the relative interior of its coface.
In particular, the incidences at a cut plaquette and a cut cube give
\begin{equation}\label{eq:reg-cut-differentials}
 (dA)_\sigma=(d_+A)_\sigma-(d_+\theta^*A)_\sigma
                                      +(d_\times A)_\sigma,
 \qquad
 (dz)(Q)=(d_+z)(Q)-(d_+\theta^*z)(Q)+(d_\times z)(Q),
\end{equation}
where $d_+$ and $d_\times$ retain the incidences from upper and cut
faces, respectively. For example,
$(d_+A)_\sigma=\sum_{b\in B_+}[\sigma:b]A(b)$.
The minus sign in the lower term follows from orientation reversal
of the cut cell and the identity $d\theta^*=\theta^*d$.

Use representatives
\[
 A=(A_+,A_-,A_0,a),\qquad z=(z_+,z_-,z_0,z_\times).
\]
Choose the interval $[0,r_q)$ for upper and mirror links, its signed
reflection for lower links, and $[-r_q/2,r_q/2)$ for cut links.
This is a fundamental domain. Reflection maps the lower variables
measure-preservingly to an independent copy of the upper variables,
fixes $(A_0,z_0)$, and changes the sign of $a$.

The upper and mirror photon weights are
\begin{equation}\label{eq:reg-half-photon-weights}
 W_+=\prod_{\sigma\in P_+}e^{-F_\sigma^2/2}
               \prod_{Q\in Q_+}1_{\{(dz)(Q)=0\}},
 \qquad W_0=\prod_{\sigma\in P_0}e^{-F_\sigma^2/2}.
\end{equation}
The full Maxwell weight with the closedness constraint is
$W_0W_+\Theta(W_+)$ times the Gaussian and constraint factors at
the cut plaquettes and cubes. The half weights contain only upper
and mirror variables. Their integrals are bounded uniformly in the
mirror variables by
\[
 \int W_+\,\dd(A_+,z_+)
       \le r_q^{|B_+|}C_q^{|P_+|},\qquad
 \int W_0\,\dd(A_0,z_0)
       \le r_q^{|B_0|}C_q^{|P_0|}.
\]
The proof drops the cube constraints and uses
\eqref{eq:reg-gaussian-sum} on each plaquette.

At a cut cube impose integer closedness by the Fourier multiplier
\[
 1_{\{(dz)(Q)=0\}}
   =\int_0^{2\pi}e^{i\alpha_Q(dz)(Q)}\frac{\dd\alpha_Q}{2\pi}.
\]
For $\alpha\in[0,2\pi)^{Q_\times}$ and
$k\in\Z^{P_\times}$ define
\begin{equation}\label{eq:reg-dual-crossing-variables}
 \begin{aligned}
 \alpha'_\sigma=\sum_{Q\in Q_\times}[Q:\sigma]\alpha_Q,
 \quad \beta_\sigma=\frac{q}{2\pi}(\alpha'_\sigma-2\pi k_\sigma),
 \\
 u_{\alpha,k}
 =\exp\left(i\sum_Q\alpha_Q(d_+z)(Q)
                  -i\sum_\sigma\beta_\sigma(d_+A)_\sigma\right).
 \end{aligned}
\end{equation}
The function $u_{\alpha,k}$ depends only on the upper variables and
has modulus one.

Poisson summation, in the normalization needed here, reads
\begin{equation}\label{eq:reg-poisson-crossing}
 \sum_{j\in\Z}e^{-(y+r_qj)^2/2+ij\alpha}
 =\frac{q}{\sqrt{2\pi}}\sum_{k\in\Z}
          e^{-\beta_k^2/2-i\beta_ky},
 \qquad\beta_k=\frac{q}{2\pi}(\alpha-2\pi k).
\end{equation}
It follows by Fourier transforming the Gaussian under
$t=y+r_qs$. Insert the cube multipliers and apply this formula at
each cut plaquette. Splitting $dA$ and $dz$ by
\eqref{eq:reg-cut-differentials} gives the exact identity
\begin{align}\label{eq:reg-photon-crossing}
 &\sum_{z_\times}
       \prod_{\sigma\in P_\times}e^{-F_\sigma^2/2}
       \prod_{Q\in Q_\times}1_{\{(dz)(Q)=0\}}
 \nonumber\\
 &\quad=\int\prod_{Q\in Q_\times}\frac{\dd\alpha_Q}{2\pi}
       \sum_{k\in\Z^{P_\times}}
       \prod_{\sigma\in P_\times}
               \frac{q}{\sqrt{2\pi}}e^{-\beta_\sigma^2/2}\,
       u_{\alpha,k}\Theta(u_{\alpha,k})
       \exp\left(iq\sum_{b\in B_\times}a_b(d^*k)(b)\right).
\end{align}
Here $k$ is extended by zero outside $P_\times$. To obtain the final
factor, observe that at a cut link $b$,
\[
 \sum_{\sigma\in P_\times}[\sigma:b]\alpha'_\sigma
 =\sum_{Q\in Q_\times}\alpha_Q
                    \sum_\sigma[Q:\sigma][\sigma:b]=0.
\]
The restriction to cut plaquettes may be dropped because any
plaquette containing a cut link is cut; the last sum is zero by
$d^2=0$. Thus the part of $\beta$ involving $\alpha'$ cancels from
the cut-link factor.

All the weights before the two reflected factors on the right of
\eqref{eq:reg-photon-crossing} are positive. Their total integral is
at most $C_q^{|P_\times|}$. Indeed, with $y=0$ in
\eqref{eq:reg-poisson-crossing},
\[
 \frac{q}{\sqrt{2\pi}}\sum_k e^{-\beta_k^2/2}
 \le\sum_j e^{-(r_qj)^2/2}\le C_q
\]
at each cut plaquette. The original Gaussian sum also has an
absolute bound $C_q^{|P_\times|}$, justifying insertion of the
Fourier multipliers and the exchanges of their sums and integrals.

Cut links occur only at coordinate link planes. A link fixed by a
mirror reflection but not contained in its plane has its endpoints
exchanged; hence it is perpendicular to the mirror, whose normal
is a coordinate vector. A diagonal mirror has cut plaquettes and
cubes, but no cut links.

For crossing-mode subsets $I,I'$ let
$\mathfrak q(I\cup I')(b)$ be the sum of their charges at modes lying
over $b$. The fermionic phases in \eqref{eq:reg-two-crossings} are
\[
 u_Iu_{I'}=\exp\left(iq\sum_{b\in B_\times}
                              a_b\mathfrak q(I\cup I')(b)\right).
\]
Integration of the cut links over one period gives
\begin{equation}\label{eq:reg-gauss-indicator}
 \int_{[-r_q/2,r_q/2)^{B_\times}}
 e^{iq\sum_ba_b((d^*k)(b)+\mathfrak q(I\cup I')(b))}\,\dd a
 =r_q^{|B_\times|}
   1_{\{d^*k+\mathfrak q(I\cup I')=0\text{ on }B_\times\}}.
\end{equation}
The quantities in the exponent are integers. This is precisely
Gauss' law at the cut links, with a nonnegative multiplier. When
there are no cut links the condition is empty.

\subsection{The reflected form on an adapted torus}
\label{subsec:reg-reflected-form}

Let $\mathcal B^+$ be the Grassmann-valued integrands whose
generators are in $\mathcal X_+$ and whose bounded measurable
coefficient functions depend only on the upper and mirror
curvatures and link phases. For the later operator estimates we
also use $\mathcal G^+$: its coefficients $\Phi$ satisfy, for some
$C<\infty$ and $\varepsilon>0$,
\begin{equation}\label{eq:reg-upper-growth}
 W_+|\Phi|\le C\prod_{\sigma\in P_+}e^{-\varepsilon F_\sigma^2}.
\end{equation}
This condition is imposed on the upper and mirror variables and
places no growth factor on mirror plaquettes. The bound is relevant
only where the upper cube constraints hold. Such half-configurations
can be extended to full closed configurations by reflecting their
upper $z$, keeping $z_0$, and taking $z_\times=0$; the cut-cube
identity \eqref{eq:reg-cut-differentials} proves closedness. Thus
these coefficient functions have an unambiguous reading wherever
$W_+\ne0$. In particular $\mathcal B^+\subset\mathcal G^+$.

\begin{theorem}\label{thm:reg-adapted-positivity}
For a torus adapted to $\ell$ and $G,G'\in\mathcal G^+$, the integral
$Z_{\mathbb T}[\Theta(G')G]$ converges absolutely, coefficient by
coefficient, and admits a representation
\begin{equation}\label{eq:reg-positive-representation}
 Z_{\mathbb T}[\Theta(G')G]
   =\int\overline{\widehat{G'}(s)}\widehat G(s)\,\dd\mathfrak m(s),
\end{equation}
where $\mathfrak m$ is a finite positive measure and
$G\mapsto\widehat G$ is linear with bounded measurable values.
This holds for even, odd and mixed-parity integrands. In particular
$Z_{\mathbb T}[\Theta(G)G]\ge0$.

For every bounded scalar observable $\Phi$ depending only on curvatures
of plaquettes wholly contained in $U_+$, the curvature law satisfies
\begin{equation}\label{eq:reg-law-positivity}
 \int_\Omega\overline{\Phi(\theta^*F)}\Phi(F)\,\dd\mu_{\mathbb T}(F)
 \ge0.
\end{equation}
It is invariant under this reflection.
\end{theorem}

\begin{proof}
The measure space has coordinates
\[
 s=(A_0,z_0,\alpha,k,I,I'),
\]
with mirror variables $(A_0,z_0)$, cut-cube variables
$\alpha\in[0,2\pi)^{Q_\times}$, integers $k\in\Z^{P_\times}$, and
the two finite mode subsets. Define
\begin{align}\label{eq:reg-positive-measure}
 \dd\mathfrak m(s)={}&W_0\,\dd(A_0,z_0)
     \prod_{Q\in Q_\times}\frac{\dd\alpha_Q}{2\pi}
     \prod_{\sigma\in P_\times}
                         \frac{q}{\sqrt{2\pi}}e^{-\beta_\sigma^2/2}
 \nonumber\\
 &\quad\cdot\omega_I\omega_{I'}r_q^{|B_\times|}
        1_{\{d^*k+\mathfrak q(I\cup I')=0\text{ on }B_\times\}},
\end{align}
where $z_0,k,I,I'$ carry counting measure. Put
\begin{equation}\label{eq:reg-half-transform}
 \widehat G(s)=\int W_+u_{\alpha,k}
         \int_{\mathcal X_+}G e^{-S_+}\zeta_I\check\zeta_{I'}
                                      \,\dd(A_+,z_+).
\end{equation}
The action $S_+$ has bounded coefficients depending only on upper
and mirror links. Thus the integrand in this formula is a
half-integrand of the stated kind.

We first prove the identity for bounded coefficients and definite
parities. Inserting \eqref{eq:reg-two-crossings} into the reflected
functional combines each term with the half factors
\[
 \xi=G e^{-S_+}\zeta_I\check\zeta_{I'},\qquad
 \xi'=G'e^{-S_+}\zeta_I\check\zeta_{I'}.
\]
The sign from this rearrangement is
$(-1)^{(|G|+|G'|)(|I|+|I'|)}$. It is one when the parities agree.
When they disagree, one of $\xi,\xi'$ is odd, so its half Berezin
integral vanishes; the corresponding full integral also vanishes
by \eqref{eq:reg-half-berezin}. The sign consequently disappears
from the integrated formula in all cases.

Now insert \eqref{eq:reg-photon-crossing} and integrate the cut
links. Formula \eqref{eq:reg-gauss-indicator} supplies the last
factor in \eqref{eq:reg-positive-measure}. At fixed $s$, the
remaining factors are
\[
 \Theta\bigl(W_+u_{\alpha,k}\xi'\bigr)
                    W_+u_{\alpha,k}\xi.
\]
The lower photon variables map measure-preservingly onto an
independent set of upper variables. Combining this change of
variables with \eqref{eq:reg-half-berezin} gives
$\overline{\widehat{G'}(s)}\widehat G(s)$, proving
\eqref{eq:reg-positive-representation} once the exchanges of
integration have been justified.

For that justification let $C$ bound all the finitely many half
Berezin coefficients of $G e^{-S_+}\zeta_I\check\zeta_{I'}$ and of
the corresponding expression with $G'$. The scalar phase factors
have modulus one. Before the cut-link integration, an absolute
majorant therefore has total integral at most
\begin{equation}\label{eq:reg-representation-majorant}
 C^2(r_qC_q)^{3V}
       \prod_{i\text{ at }\Pi}(1+\omega_i)
       \prod_{i'\text{ at }\Pi'}(1+\omega_{i'}).
\end{equation}
The upper, lower and mirror variables are bounded by their
Gaussian sums; the cut-link intervals contribute their lengths;
the dual photon weights are bounded by $C_q^{|P_\times|}$; and the
finite mode sums give the two products. This proves absolute
integrability and all required uses of Fubini's theorem. The
same bounds without the half factors show that $\mathfrak m$ is
finite. It is positive by construction. Formula
\eqref{eq:reg-half-transform} is measurable and bounded uniformly
in $s$ by $Cr_q^{|B_+|}C_q^{|P_+|}$. Mixed parities follow by
sesquilinearity.

For $G,G'\in\mathcal G^+$, use
\eqref{eq:reg-upper-growth} with exponents $\varepsilon$ and
$\varepsilon'$. The algebraic identities above are unchanged.
The absolute majorant now replaces the upper and lower Gaussian
sums by $C_q(\varepsilon)^{|P_+|}$ and
$C_q(\varepsilon')^{|P_-|}$, respectively, while the mirror and
cut sums retain exponent $1/2$. These quantities are finite by
\eqref{eq:reg-gaussian-sum}. This proves coefficientwise
integrability, the same representation, and the boundedness of
$\widehat G$. Taking $G'=G$ proves positivity.

For a scalar $\Phi$ with this plaquette dependence, divide this positivity
by $Z_{\mathbb T}[1]>0$ and use the definition of the curvature
law. Reflection invariance follows from
\eqref{eq:reg-reflected-functional}, since the adapted torus is
preserved by reflection. This proves
\eqref{eq:reg-law-positivity} and completes the theorem.
\end{proof}

The same form remains positive after multiplication of the weight
by $\Phi_+\Theta(\Phi_+)$, with $\Phi_+$ bounded and scalar on the
upper half: apply the theorem to $G\Phi_+$. This observation will
be used for the reflected arrays of observables. It uses positivity
of the joint reflected form before integrating out the fermion.
The determinant's pointwise nonnegativity by itself does not imply
reflection positivity.


\section{Slab operators and finite-volume comparison}\label{sec:slab-comparison}

A cubic torus admits the coordinate reflection-positive representations
of the preceding section. A diagonal mirror requires a different period
lattice. To compare the two geometries, we express the functional as a
trace of operators associated with slabs. Reflection gives the adjoint
of a slab operator; repeated reflected copies give its Schatten norm.
The resulting comparison changes the closing vector, and, when needed,
the fermion character. Three such changes connect the cubic torus with
the diagonal torus and with a second coordinate torus of the same volume.
Their pressures differ by a quantity tending to zero. This is the
finite-volume input for diagonal reflection positivity.

We also prove Gaussian domination on coordinate tori. Its trace formula
uses real slice connections and a sum over magnetic fluxes. A real
external two-cochain enters through unitary electric translations and
magnetic multiplication operators. Schatten's inequality first removes
the electric translations. A second choice of time direction removes
the remaining magnetic perturbation.

Throughout, $q>0$, $0<\rho<1$, $w_i>0$, and $M\in\R$. The mass may
depend on the mesh, but it is fixed when the volume tends to infinity.
We use $r_q=2\pi/q$ and the counting Berezin convention
\eqref{eq:reg-berezin}. The determinant, and hence every curvature law
below, agrees with the convention used in the renormalization argument.
More precisely, there are $192V$ fermion modes on a torus of photon
volume $V$. With the pairing
\[
 \langle\bar\Psi,O\Psi\rangle_{\mathrm{ch}}
       =\frac1{48}\sum_x\bar\Psi(x)(O\Psi)(x),
 \qquad \mathrm D\Psi=48^{192V}\int_{\mathbb T},
\]
one has
\[
 \int\mathrm D\Psi\,
       e^{-\langle\bar\Psi,O\Psi\rangle_{\mathrm{ch}}}=
       \det O.
\]
The factor $48^{192V}$ cancels the factor from the determinant of the
weighted quadratic form. Thus no additional normalization of the
curvature law is needed at this interface. The slab modes in this
section use the counting convention.

\subsection{Local observables and their lifts}\label{subsec:slab-local}

We need local observables that may contain fermion variables and whose
coefficients grow with the curvature. Their definition must make sense
in every flux sector. Let $B=y+[-r,r]^3$, and write $B^{(j)}$ for the
oriented lattice $j$-cells contained in $B$. The fermion sites assigned
to $B$ are
\[
 \mathcal X(B)=\{x\in\mathcal X:
          \pi(x)+[-\tfrac12,\tfrac12]^3\subset B\}.
\]
This support condition includes the photon site at which the fermion
transforms under gauge transformations.

Let $\mathcal G(B)$ consist of the Borel maps
$G:\R^{B^{(1)}}\to\operatorname{Gr}(\mathcal X(B))$ with gauge-invariant
values and with the following growth bound. If
$G(\widehat A)=\sum_\Psi a_\Psi(\widehat A)\Psi$, then, for some
$C<\infty$ and $\varepsilon>0$,
\begin{equation}\label{eq:slab-box-growth}
 |a_\Psi(\widehat A)|\le C
   \exp\left((\tfrac12-\varepsilon)
                    \sum_{\sigma\in B^{(2)}}(d\widehat A)_\sigma^2\right).
\end{equation}
Here gauge invariance means
\begin{equation}\label{eq:slab-box-gauge}
 a_\Psi(\widehat A-d\lambda)
        e^{iq\langle\lambda,\mathfrak q(\Psi)\rangle}
       =a_\Psi(\widehat A)
\end{equation}
for every real function $\lambda$ on the photon sites of $B$;
$\mathfrak q(\Psi)$ counts the $\psi$ generators minus the
$\bar\psi$ generators at fermion sites projecting to each photon site.
We denote the even subspace by $\mathcal G^{\mathrm{ev}}(B)$.

A local lift of a configuration $(A,z)$ on $B$ is a real one-cochain
$\widehat A$ on $B^{(1)}$ satisfying
\begin{equation}\label{eq:slab-local-lift}
 d\widehat A=F\quad\text{on }B^{(2)},\qquad
 e^{iq\widehat A(b)}=e^{iqA(b)}\quad(b\in B^{(1)}).
\end{equation}
We evaluate $G$ at this lift and then impose the fermion identifications
of the torus or cylinder. We now verify that this procedure is Borel
and independent of the lift.

Replace $B$ by its integer box
$\prod_\mu[a_\mu,b_\mu]$, where $a_\mu=\lceil y_\mu-r\rceil$ and
$b_\mu=\lfloor y_\mu+r\rfloor$. If this box has no sites, use the
unique empty lift. Otherwise its axial spanning tree consists of all
links in direction $3$, the links in direction $2$ on $x_3=a_3$, and the
links in direction $1$ on $x_2=a_2$, $x_3=a_3$. Prescribe a one-cochain
on this tree. Its values on the remaining links are determined by
\begin{align}\label{eq:slab-lift-recursion}
 \widehat A_1(x+e_2)&=\widehat A_1(x)+\widehat A_2(x+e_1)
                                  -\widehat A_2(x)-F_{12}(x),
                                  &&x_3=a_3,\notag\\
 \widehat A_1(x+e_3)&=\widehat A_1(x)+\widehat A_3(x+e_1)
                                  -\widehat A_3(x)-F_{13}(x),\\
 \widehat A_2(x+e_3)&=\widehat A_2(x)+\widehat A_3(x+e_2)
                                  -\widehat A_3(x)-F_{23}(x).\notag
\end{align}
First use the first line in increasing $x_2$, then the last two lines
in increasing $x_3$. Each link outside the tree occurs as an unknown
once. The construction is linear in $F$ and the tree values. It gives
the prescribed $13$- and $23$-curvatures and the $12$-curvature on the
bottom face. If $dF=0$, induction over the cubes gives the remaining
$12$-curvatures: on each cube the two closed two-cochains $F$ and
$d\widehat A$ agree on five faces, hence on the sixth.

On the tree choose $\widehat A(b)=\arg_q(e^{iqA(b)})$, where
$\arg_q$ takes values in $[0,r_q)$. This is a Borel choice. Subtract
the recursion for $A$ from \eqref{eq:slab-lift-recursion}. The tree
differences lie in $r_q\Z$, and every subsequent difference is a signed
sum of earlier ones and of a value of $r_qz$. Thus
$\widehat A-A$ is integer-valued in units of $r_q$ on every link. This
proves \eqref{eq:slab-local-lift}, with Borel dependence on $F$ and the
tree phases.

Two local lifts differ by $r_qk$, where $k$ is an integer closed
one-cochain. Sum $k$ along the tree paths from its bottom corner to
define an integer function $\lambda$. Then $d\lambda=k$ on the tree.
The zero-curvature recursion shows that $k-d\lambda=0$ on every link.
Equation \eqref{eq:slab-box-gauge}, applied with $-r_q\lambda$, gives
$a_\Psi(\widehat A+r_qd\lambda)=a_\Psi(\widehat A)$, since
$\langle\lambda,\mathfrak q(\Psi)\rangle$ is integer. The value of $G$
is therefore independent of the lift.

For an open set $U$, let $\mathcal A(U)$ be the linear span of the maps
\begin{equation}\label{eq:slab-admissible-product}
 \widehat A\longmapsto\Phi(d\widehat A)
                   \prod_{j=1}^rG_j(\widehat A|_{B_j^{(1)}}),
\end{equation}
where the boxes $B_j\subset U$ are pairwise disjoint,
$G_j\in\mathcal G^{\mathrm{ev}}(B_j)$, and $\Phi$ is a bounded Borel
function of finitely many curvature values at plaquettes contained in
$U$. A box $B$ embeds in a period group $\Gamma$ when the sets
$B+[-1,1]^3+\gamma$, $\gamma\in\Gamma$, are pairwise disjoint.
Boxes $B_1,\ldots,B_r$ embed disjointly in $\Gamma$ when each
embeds and all sets $B_j+\gamma$, indexed by distinct pairs
$(j,\gamma)$, are pairwise disjoint. A product is read on the
quotient only when its boxes embed disjointly.

It remains to check that equal maps \eqref{eq:slab-admissible-product}
have equal values on configurations. This requires one lift for every
box in a representation. On $\Z^3$, a closed integer two-cochain $z$
has the integer primitive
\begin{align*}
 k_1(x)&=0,\\
 k_2(x)&=\sum_{s=0}^{x_1-1}z_{12}(s,x_2,x_3),\\
 k_3(x)&=\sum_{s=0}^{x_1-1}z_{13}(s,x_2,x_3)
                    +\sum_{s=0}^{x_2-1}z_{23}(0,s,x_3).
\end{align*}
The sums are oriented for negative endpoints. The first two components
of $dk=z$ follow by taking differences. In the third, $dz=0$ changes
the summand into a difference of consecutive $z_{23}$ values, so the
sum telescopes. Therefore $\widehat A=A+r_qk$ is a lift on the whole
lattice. It need not be periodic, even when $(A,z)$ is periodic.

The same conclusion holds on a closed slab. Choose a coordinate
$\mu_0$ absent from the normal $n_\ell$. The slab cell complex is a
planar strip times $\Z$ in direction $\mu_0$.
Let $\mu<\nu$ be the other coordinates, denoted by $(a,r)$.
Write $Q(a,r)$ for the unit square with lower corner $(a,r)$ and
$I_r$ for the interval of $a$ such that $[(a,r),(a,r+1)]$ lies
in the strip. When $I_r$ is nonempty choose $a_r\in I_r$.
On the base sheet set the $\mu$-link values of $k$ to zero and put
\[
 k([(a,r),(a,r+1)]\times\{0\})
       =\sum_{a_r\le s<a}z(Q(s,r)\times\{0\}).
\]
All squares in this oriented sum lie in the strip because $I_r$
is an interval. Set $k=0$ on the $\mu_0$-links. For a horizontal
link $b$ in direction $\nu'\in\{\mu,\nu\}$, extend by
\[
 k(b\times\{j\})=k(b\times\{0\})
       +s_{\nu'}\sum_{0\le j'<j}z(b\times[j',j'+1]),
 \qquad
 s_{\nu'}=\begin{cases}-1,&\nu'<\mu_0,\\+1,&\nu'>\mu_0.\end{cases}
\]
The sums with reversed endpoints have the opposite sign.
The resulting primitive has $dk=z$ on all vertical
plaquettes and on the base sheet. The closed two-cochain $z-dk$ then
has vanishing vertical components; closedness makes its horizontal
component constant between sheets. It vanishes on the base sheet and
hence everywhere. Again the primitive is taken on the lifted strip,
not on its periodic quotient.

Evaluation at this common lift, followed by the quotient homomorphism
on fermion generators, agrees with the separate local evaluations.
It is linear in the map $G$, so it is independent of any representation
of $G$ as a sum of products. Moreover, different box plaquettes have
different representatives on the quotient. Multiplying
\eqref{eq:slab-box-growth} by the Maxwell density therefore leaves
\begin{equation}\label{eq:slab-local-majorant}
 C\exp(-\varepsilon_G\|F\|^2),\qquad
 \varepsilon_G=\min(\tfrac12,\varepsilon_1,\ldots,\varepsilon_r)>0.
\end{equation}
The finite Grassmann expansion of the action has bounded coefficients
for real $A$. The Gaussian sum \eqref{eq:reg-gaussian-sum} now proves
absolute integrability, coefficient by coefficient. This also proves
the required Borel measurability on the quotient.

\subsection{Boundary spaces and slab kernels}\label{subsec:slab-kernels}

Fix a mirror family $\ell$ and a cross-section
$\Xi=(\Gamma_\parallel,\chi_\parallel)$, where
$\Gamma_\parallel\subset\Z^3\cap\{\ell=0\}$ has rank two and systole
at least four. Here and below, the systole of a period lattice
$\Gamma$ is $\min_{\gamma\in\Gamma\setminus\{0\}}|\gamma|_\infty$.
Work on the cylinder $\Z^3/\Gamma_\parallel$. The
admissible mirror positions are
$c\in\frac12|n_\ell|^2\Z$. Write $\mathcal M^j(c)$ for the $j$-cells
contained in $\{\ell=c\}$, $\mathcal X^j(c)$ for those cut through
their relative interiors by this plane, and $\mathcal B^j(c,c')$ for
the $j$-cells contained in $\{c\le\ell\le c'\}$ and in neither
bounding plane. In particular a slab cell may touch either boundary.
All these sets are finite on a bounded slab.

For a chain of admissible mirrors, the cylinder cells partition into
the cells in its mirrors, those cut by them, and those lying between
consecutive mirrors. A cell cut by $c$ has its height range
$[c-|n_\ell|^2/2,c+|n_\ell|^2/2]$. Its non-cut faces lie on one
side of the plane; the other faces are cut by the same plane. This
partition applies also when two mirrors are consecutive. A fermion
bond crosses at most one mirror of such a chain, and no fermion site
lies in a mirror.

A closed-slab configuration $\varpi=(A,z)$ assigns real link values
and integer plaquette values to all cells of the closed slab, with
$dz=0$ on its cubes. We integrate its interior values using
$A(b)\in[0,r_q)$ and counting measure for $z$. Its interior photon
weight is
\begin{equation}\label{eq:slab-interior-weight}
 W_{c,c'}(\varpi)=
    \prod_{\sigma\in\mathcal B^2(c,c')}e^{-F_\sigma^2/2}
    \prod_{Q\in\mathcal B^3(c,c')}\mathbf1[(dz)(Q)=0].
\end{equation}
We use the larger class $\mathcal I_\Xi[c,c']$ of Borel
Grassmann-valued functions of closed-slab configurations, with fermion
support $c<\ell<c'$, invariant under
$(A,z)\mapsto(A+r_qk,z-dk)$, whose coefficients obey
\begin{equation}\label{eq:slab-integrand-growth}
 |a_\Psi(\varpi)|\le C_G
   \exp\left((\tfrac12-\varepsilon_G)
          \sum_{\sigma\in\mathcal B^2(c,c')}F_\sigma^2\right),
 \qquad 0<\varepsilon_G\le\tfrac12.
\end{equation}
The local products supported in the open slab give elements of this
class by the lift construction above. Products in consecutive slabs
remain in the class, with the smaller of the two positive
$\varepsilon$'s. Reflection and translation also preserve the class.

At the mirror $c$, let $\mathcal J(c)$ be the ordered set of crossing
modes from \eqref{eq:reg-crossing-modes}, taken modulo
$\Gamma_\parallel$. Let $\mathcal F(c)$ be the exterior space with
orthonormal basis $e_I$, $I\subset\mathcal J(c)$. Put
$\omega_c e_I=\omega_Ie_I$ and $U_{\mathrm f}e_I=(-1)^{|I|}e_I$.
Translation of the modes includes both the character
$\chi_\parallel$ and the sign required to reorder their product:
\[
 U_ve_I=s_v(I)e_{r_v(I)},\qquad
 \tau_v\zeta_I^{(c)}=s_v(I)\zeta_{r_v(I)}^{(c+\ell(v))}.
\]
The $U_v$ are unitary, satisfy the translation law, and commute with
parity. For an integer translation parallel to the normal, choose the
orders compatibly so that $s_v(I)=1$.

For a Grassmann element $H$ supported in the open slab define
\begin{equation}\label{eq:slab-fermion-matrix}
 N_{c,c'}[H]_{I',I}
    =\int_{(c,c')}\zeta_I^{(c)}H
                        \Theta_{c'}(\zeta_{I'}^{(c')}).
\end{equation}
An operator $\Upsilon$ on $\mathcal F(c)$ is represented at the
mirror by
\[
 \operatorname{Ins}_c[\Upsilon]
       =\sum_{I,J}\Upsilon_{IJ}\Theta_c(\zeta_J^{(c)})\zeta_I^{(c)}.
\]
The particular insertion
$\operatorname{Ins}_c[\omega_c\diag(u_I(A))]$ is the crossing
exponential \eqref{eq:reg-crossing-exponential}. The ordering in these
definitions gives
\begin{equation}\label{eq:slab-matrix-composition}
 N_{c,c''}[H\operatorname{Ins}_{c'}[\Upsilon]H']
             =N_{c',c''}[H']\Upsilon N_{c,c'}[H].
\end{equation}
Indeed, an entry of the left side is a sum of products
$[\zeta_IH\Theta_{c'}\zeta_J]
 [\zeta_KH'\Theta_{c''}\zeta_{I''}]$ in disjoint sets of slab
generators. Berezin factorization gives the entries on the right
without moving any of these factors. Reflecting the integrand reverses
their order and conjugates the Berezin coefficient; with
$v=2(c'-c)n_\ell/|n_\ell|^2$ this gives
\begin{equation}\label{eq:slab-matrix-adjoint}
 N_{c',2c'-c}[\Theta_{c'}H]=U_vN_{c,c'}[H]^*.
\end{equation}

We next include the photon crossing in the boundary data. The cut
incidence maps $d_c$ and $d_c^+$ are the restrictions of the coboundary
to cut faces and to upper faces, as in
\eqref{eq:reg-cut-differentials}. Introduce one angle $\vartheta_Q$
for each cut cube and one integer $k_\sigma$ for each cut plaquette,
and set
\begin{align}\label{eq:slab-dual-measure}
 \varkappa_\sigma&=r_q^{-1}
       ((d_c^*\vartheta)_\sigma-2\pi k_\sigma),\notag\\
 \dd\mu_c^\times&=
      \prod_{Q\in\mathcal X^3(c)}\frac{\dd\vartheta_Q}{2\pi}
      \prod_{\sigma\in\mathcal X^2(c)}
                     \frac{q}{\sqrt{2\pi}}e^{-\varkappa_\sigma^2/2},\\
 e_c((\vartheta,k);\varpi)&=
       e^{i\langle\vartheta,d_c^+z\rangle
                             -i\langle\varkappa,d_c^+A\rangle},\notag\\
 D_c(k,I)&=r_q^{|\mathcal X^1(c)|}
       \mathbf1[d_c^*k+\mathfrak q(I)=0].\notag
\end{align}
Here $\mathfrak q(I)(b)$ is the number of modes of groups $1,2$ over
the cut link $b$, minus the number of groups $3,4$ over it. There are
no cut links at diagonal mirrors. At a coordinate site plane there
are no cut cells at all; its Fourier space is a point.

The one-mirror photon identity proved by the Fourier and Poisson
calculation of the preceding section is
\begin{multline}\label{eq:slab-photon-gluing}
 \sum_{z^\times}\int\dd A^\times\,
       \prod_{\sigma\in\mathcal X^2(c)}e^{-F_\sigma^2/2}
       \prod_{Q\in\mathcal X^3(c)}\mathbf1[(dz)(Q)=0]\,u_I(A^\times)
       \\
 =\int\dd\mu_c^\times\,D_c(k,I)
           e_c((\vartheta,k);\varpi_+)
            \overline{e_c((\vartheta,k);\theta_c^*\varpi_-)}.
\end{multline}
The upper and lower data on the right have the same values on the
mirror. Their upper and reflected lower faces provide the non-cut
faces on the left. Cube Fourier multipliers enforce closedness,
Poisson summation gives the factors $e^{-\varkappa^2/2}$, and compact
cut-link integration gives $D_c(k,I)$. This identity is used in the
direction shown when constructing a kernel and in the opposite
direction when gluing two kernels.

Mirror variables are
$\xi=(\mathfrak m,\vartheta,k,I)$, with
$\mathfrak m=(A_0,z_0)\in[0,r_q)^{\mathcal M^1(c)}
\times\Z^{\mathcal M^2(c)}$. On their space $Y_c$ put
\begin{equation}\label{eq:slab-boundary-measure}
 \dd\mu_c(\xi)=\dd A_0
      \prod_{\sigma\in\mathcal M^2(c)}e^{-F_\sigma(\mathfrak m)^2/2}
                  \dd\mu_c^\times(\vartheta,k)D_c(k,I),
 \qquad \mathcal H(c)=L^2(Y_c,\mu_c).
\end{equation}
The integer and mode variables carry counting measure. This is a
finite positive measure. Summing the mirror plaquette Gaussians gives
$C_q^{|\mathcal M^2(c)|}$, while Poisson summation at fixed angles
gives
\[
 \frac q{\sqrt{2\pi}}\sum_{k_\sigma}e^{-\varkappa_\sigma^2/2}
 \le\sum_{j\in\Z}e^{-r_q^2j^2/2}.
\]
The remaining factors are bounded by
$r_q^{|\mathcal M^1(c)|+|\mathcal X^1(c)|}2^{|\mathcal J(c)|}$.
The induced signed translations $V_v:\mathcal H(c)\to
\mathcal H(c+\ell(v))$ are unitary and satisfy the group law. The
parity operator multiplies the component $I$ by $(-1)^{|I|}$.

Let $S_{(c,c')}(A)$ be the fermion action containing precisely the mass
and bond terms with fermion endpoints in the open slab. For
$G\in\mathcal I_\Xi[c,c']$, define the kernel
\begin{multline}\label{eq:slab-kernel}
 \mathfrak S[G](\xi',\xi)=\sqrt{\omega_{I'}\omega_I}
      \int\dd\varpi^\circ\,W_{c,c'}(\varpi)
      e_c((\vartheta,k);\varpi)
       \overline{e_{c'}((\vartheta',k');\theta_{c'}^*\varpi)}
       \\
       \cdot N_{c,c'}[G(\varpi)e^{-S_{(c,c')}(A)}]_{I',I},
       \qquad \varpi=\mathfrak m\vee\varpi^\circ\vee\mathfrak m'.
\end{multline}
Set the whole integrand to zero where $W_{c,c'}(\varpi)=0$;
$G(\varpi)$ is evaluated only on the closed configurations in its
domain.
The associated integral operator maps $\mathcal H(c)$ to
$\mathcal H(c')$. The square roots of the mode weights belong in the
kernel; the measure \eqref{eq:slab-boundary-measure} contains none of
these weights. This places one full crossing weight at a glued mirror.

\begin{theorem}\label{thm:slab-operators}
Every $\mathfrak S[G]$ is Hilbert--Schmidt and depends linearly on $G$.
If $G$ has constants $C_G,\varepsilon_G$ in
\eqref{eq:slab-integrand-growth}, then
\begin{align}\label{eq:slab-kernel-bound}
 |\mathfrak S[G](\xi',\xi)|
    &\le C_{c,c'}C_GC_q(\varepsilon_G)^{|\mathcal B^2(c,c')|},\notag\\
 \|\mathfrak S[G]\|_2
    &\le\sqrt{\mu_c(Y_c)\mu_{c'}(Y_{c'})}
                C_{c,c'}C_GC_q(\varepsilon_G)^{|\mathcal B^2(c,c')|}.
\end{align}
The constant depends on the fixed cross-section, mirror positions,
$q,M,\rho,w_i$, and chamber base point. For consecutive slabs,
\begin{align}\label{eq:slab-operator-rules}
 \mathfrak S[G']\mathfrak S[G]&=\mathfrak S[GG'],\notag\\
 \mathfrak S[\tau_vG]&=V_v\mathfrak S[G]V_v^{-1},\\
 \mathfrak S[\Theta_{c'}G]&=V_{2(c'-c)n_\ell/|n_\ell|^2}
                                         \mathfrak S[G]^*.\notag
\end{align}
If $\mathbb T=(\Gamma,\chi)$ has cross-section $\Xi$, $g\in\Gamma$
is a closing vector with minimal positive height $p=\ell(g)$, and
$c_0<\cdots<c_r=c_0+p$ are admissible, $r\ge2$, and
$G_j\in\mathcal I_\Xi[c_{j-1},c_j]$ for every $j$, then
\begin{equation}\label{eq:slab-trace}
 Z_{\mathbb T}[G_1\cdots G_r]
    =\operatorname{Tr}_{\mathcal H(c_0)}
       \left(V_g^{-1}U_{\mathrm f}^{(1+\chi(g))/2}
               \mathfrak S[G_r]\cdots\mathfrak S[G_1]\right).
\end{equation}
These identities apply to arbitrary parities. For homogeneous $G$,
$U_{\mathrm f}\mathfrak S[G]=(-1)^{\deg G}\mathfrak S[G]U_{\mathrm f}$.
\end{theorem}

\begin{proof}
The coefficient of the finite Grassmann polynomial
$e^{-S_{(c,c')}(A)}$ is bounded uniformly in real $A$. Each coefficient
in the matrix in \eqref{eq:slab-kernel} is consequently bounded by a
fixed multiple of the majorant for $G$. The phases have modulus one.
After multiplication by the photon weight, the absolute integrand is
bounded by
\[
 C C_G\prod_{\sigma\in\mathcal B^2(c,c')}
                           e^{-\varepsilon_GF_\sigma^2}.
\]
Drop the cube constraints and sum the integers separately. Each sum
is at most $C_q(\varepsilon_G)$; each compact link integral contributes
$r_q$. The mode factors and finite matrix dimension enter $C_{c,c'}$.
This proves the pointwise bound in \eqref{eq:slab-kernel-bound}.
The boundary measures are finite, so the kernel is in
$L^2(Y_{c'}\times Y_c)$, with the stated norm. These bounds also
justify all the gluing integrals below before any cancellation.

For composition, multiply the two kernels and integrate over the common
boundary. Their two square roots give $\omega_{I'}$ for the middle
mode. The two phases and its boundary measure are precisely the right
side of \eqref{eq:slab-photon-gluing}. Applying that identity restores
the photon variables on the cut cells, their Maxwell weights, their
cube constraints, and the phase $u_{I'}$ of the crossing modes. Sum
the middle modes by \eqref{eq:slab-matrix-composition}, with
$\Upsilon=\omega_{c'}\diag(u_{I'})$. This restores the crossing
fermion action. The plaquettes in the middle mirror come from its
boundary measure, and the cell partition restores the full weight on
the combined slab. No cube lies in a mirror, so every closedness
constraint is present once. This is the kernel of $\mathfrak S[GG']$.

Translation acts by a signed relabeling of the modes and an
orientation-preserving relabeling of cells, proving covariance.
For reflection use \eqref{eq:slab-matrix-adjoint}; the two boundary
phases exchange their roles and the interior weight is preserved.
Reflection may reverse a compact integration interval, but integer
shift invariance makes any interval of length $r_q$ equivalent to
$[0,r_q)$. The reflected upper boundary is the translate of the
original lower one by $2(c'-c)n_\ell/|n_\ell|^2$. This gives the last
identity of \eqref{eq:slab-operator-rules}. Parity follows because a
Berezin integral of odd total degree vanishes.

For the trace, composition reduces the claim to two slabs. A product
of two Hilbert--Schmidt kernels is trace class, and its trace is the
absolutely convergent integral of their cyclic product. The closing
operator identifies the upper mirror data with the translate by $g$
of the lower data. Its contribution to this cyclic kernel is
\[
 (-\chi(g))^{|I_0|}s_g(I_0).
\]
Here $s_g$ is the mode relabeling sign; $(-\chi(g))^{|I_0|}$ is the
remaining Grassmann closing sign. It is $1$ for antiperiodic closure
and the fermion parity for periodic closure. To see this sign directly,
translate the lower crossing product to the upper boundary. Each of
its $|I_0|$ generators contributes $\chi(g)$, while moving it across
the adjacent upper product contributes $(-1)^{|I_0|}$ in every term
of even total degree. Their product is $(-\chi(g))^{|I_0|}$.

Apply \eqref{eq:slab-photon-gluing} at the two mirrors. For fixed
photon variables the remaining mode sum is
\[
 \tr_{\mathcal F(c_0)}
    (U_g^{-1}U_{\mathrm f}^{(1+\chi(g))/2}
                       N_2\Upsilon_1N_1\Upsilon_0),
 \qquad \Upsilon_j=\omega_{c_j}\diag(u_I(A)).
\]
Berezin composition and closure replace this trace by the full torus
Berezin integral. The two open-slab actions and the two crossing
actions include each mass and hopping term once. The photon variables
likewise include every torus link and plaquette once, and the cut and
interior cube constraints impose $dz=0$ everywhere. Their measure is
therefore the measure in \eqref{eq:reg-functional}. This proves
\eqref{eq:slab-trace} for even total degree. With odd total degree both
sides vanish, the operator under the trace anticommuting with parity.
Finally, expand nonhomogeneous entries into their even and odd parts.
\end{proof}

\subsection{Reflected folds and Schatten estimates}\label{subsec:slab-folds}

For an operator $Y$ between Hilbert spaces, write
$\|Y\|_s=(\operatorname{Tr}(Y^*Y)^{s/2})^{1/s}$. We use Schatten's
inequality in the following form. If $1\le s_j<\infty$,
$Y_j:H_{j-1}\to H_j$ belongs to the Schatten class of exponent $s_j$,
$H_r=H_0$, and $\sum_js_j^{-1}=1$, then the cyclic product is trace
class and
\begin{equation}\label{eq:slab-schatten-holder}
 |\operatorname{Tr}(Y_r\cdots Y_1)|
                \le\prod_{j=1}^r\|Y_j\|_{s_j}.
\end{equation}
The statement for different spaces follows from the usual inequality
on their orthogonal direct sum: extend each $Y_j$ by zero on all other
summands. The nonzero block of the product is the displayed cyclic
product, and the singular values of each extension are those of $Y_j$.
A unitary closing identification can be absorbed into $Y_r$ without
changing its Schatten norm.

Let $\mathbb T$ have height period $p$ in the family $\ell$. Its
straightened torus $\mathbb T^0$ has the same cross-section and the
closing vector
\begin{equation}\label{eq:slab-straight-vector}
 g^0=\frac p{|n_\ell|^2}n_\ell,
 \qquad \Gamma^0=\Gamma_\parallel\oplus\Z g^0,
 \qquad\chi^0(g^0)=-1.
\end{equation}
We use this torus only when $g^0\in\Z^3$ and its lattice has systole
at least four. Its volume is the volume of $\mathbb T$: both lattices
have the same covolume in the mirror and the same normal height period.

Divide one period into $2k$ slabs of equal height $h'$, with all
$c_j=c_0+jh'$ admissible and $p=2kh'$. For an integrand $G$ in one of
these slabs, form its fold by reflecting it successively in the
bounding planes, taking all $2k$ copies. This is an integrand on
$\mathbb T^0$. Its fermion boxes are disjoint: they lie strictly
between different consecutive mirrors; within each slab the
cross-section embedding separates their translates.

\begin{proposition}\label{prop:slab-fold-holder}
For integrands $G_j\in\mathcal I_\Xi[c_{j-1},c_j]$,
\begin{align}\label{eq:slab-fold-holder}
 Z_{\mathbb T^0}[\operatorname{fold}_\ell G_j]
              &=\|\mathfrak S[G_j]\|_{2k}^{2k},\notag\\
 |Z_{\mathbb T}[G_1\cdots G_{2k}]|
              &\le\prod_{j=1}^{2k}
               Z_{\mathbb T^0}[\operatorname{fold}_\ell G_j]^{1/(2k)}.
\end{align}
In particular every folded factor on the right is nonnegative.
\end{proposition}

\begin{proof}
Translate the starting slab, if necessary, to $[c_0,c_1]$, and put
$Y=\mathfrak S[G]$. Let
$v=2h'n_\ell/|n_\ell|^2$, which is integer because both mirror
positions are admissible. The first reflected pair has product
$V_vY^*Y$ by \eqref{eq:slab-operator-rules}. Its translate by $sv$
has product $V_{(s+1)v}Y^*YV_{sv}^{-1}$. Multiplying all $k$ pairs
telescopes to
\[
 V_{kv}(Y^*Y)^k.
\]
Since $kv=g^0$ and $\chi^0(g^0)=-1$, the closing translation in
\eqref{eq:slab-trace} cancels $V_{g^0}$ and introduces no parity.
The trace is $\operatorname{Tr}(Y^*Y)^k$, proving the first identity.
Translation of a different starting slab changes neither this trace
nor the singular values.

For the original torus, the operator
$V_g^{-1}U_{\mathrm f}^{(1+\chi(g))/2}$ is unitary. Apply
\eqref{eq:slab-schatten-holder} to its trace formula, with every
exponent equal to $2k$. The first identity identifies the resulting
norms with folded partition functions. The Hilbert--Schmidt estimate
ensures all operators belong to the required Schatten classes.
\end{proof}

We next apply this inequality to cells. A frame is a triple of the
mirror families with mutually orthogonal normals $n_1,n_2,n_3$.
The two frames needed below are $(x_1,x_2,x_3)$ and
$(x_1-x_2,x_1+x_2,x_3)$. Choose admissible offsets $o_a$ and an integer
height width $h>0$. The open cells are
\begin{equation}\label{eq:slab-cells}
 Q_\kappa=\{x:o_a+h\kappa_a<\ell^a(x)
                         <o_a+h(\kappa_a+1),\ a=1,2,3\},
 \qquad |Q|=\frac{h^3}{|n_1||n_2||n_3|}.
\end{equation}
The reflections in the planes $\ell^a=o_a+hj$ generate a product of
three infinite dihedral groups. Exactly one element of this group
carries a given open cell to any specified open cell. Its action on
even Grassmann-valued functions is generated by the reflections
$\Theta$ and translations used above. Denote the resulting copy of
$G\in\mathcal A(Q_\kappa)$ in $Q_{\kappa'}$ by $G^{[\kappa']}$.

A rectangular torus in this frame has periods
$\gamma_a=p_an_a/|n_a|^2\in\Z^3$, with $\chi(\gamma_a)=-1$.
If each $p_a\in2h\Z$, the copies fit its period lattice. Let
$N=|\mathbb T|/|Q|$ be the number of cells, and define
\begin{equation}\label{eq:slab-chessboard-norm}
 \operatorname{Fold}_{\mathbb T}G
          =\prod_{\kappa'\bmod\Gamma}G^{[\kappa']},
 \qquad \|G\|_{\mathbb T}
          =Z_{\mathbb T}[\operatorname{Fold}_{\mathbb T}G]^{1/N}.
\end{equation}
This is an unnormalized seminorm. Its defining partition function is
nonnegative, by the first identity of
\eqref{eq:slab-fold-holder} applied to one layer of the folded array.
The value is independent of the chosen representatives of the cells.

\begin{theorem}\label{thm:slab-chessboard}
If a rectangular torus $\mathbb T$ carries the cells
\eqref{eq:slab-cells}, then
\begin{equation}\label{eq:slab-chessboard-rectangular}
 \left|Z_{\mathbb T}\left[\prod_{\kappa}G_\kappa\right]\right|
                       \le\prod_\kappa\|G_\kappa\|_{\mathbb T},
 \qquad G_\kappa\in\mathcal A(Q_\kappa).
\end{equation}
Let $C(a,b,c)$ have coordinate periods $ae_1,be_2,ce_3$,
all with character $-1$.
Let $e_-=e_1-e_2$, $e_+=e_1+e_2$, and let $D(a,b,c)$ have periods
$ae_-,be_+,ce_3$, all with character $-1$. Let $D'(a,b,c)$ have the
same periods and characters, except for $\chi(be_+)=+1$.
For the diagonal frame and $n\in2h\Z$, $n\ge8$, the cells have
representatives
\[
 0\le\kappa_1<n/h,\qquad 0\le\kappa_2<2n/h,
                     \qquad0\le\kappa_3<n/h.
\]
They represent the cells of $C(n,n,n)$, $D'(n/2,n,n)$, and
$D(n/2,n,n)$, and
\begin{equation}\label{eq:slab-chessboard-diagonal}
 \left|Z_{\mathbb T}\left[\prod_\kappa G_\kappa\right]\right|
           \le\prod_\kappa\|G_\kappa\|_{D(n/2,n,n)},
 \quad \mathbb T=C(n,n,n)\ \text{or }D'(n/2,n,n).
\end{equation}
For coordinate cells of width $h$, $a,c\in2h\Z$, $a,c\ge4$, one has
\begin{equation}\label{eq:slab-chessboard-coordinate}
 \left|Z_{D(a,2a,c)}\left[\prod_\kappa G_\kappa\right]\right|
                \le\prod_\kappa\|G_\kappa\|_{C(a,4a,c)}.
\end{equation}
\end{theorem}

\begin{proof}
Group the cell integrands into layers perpendicular to the first normal,
and apply \eqref{eq:slab-fold-holder}. Its right side consists of the
arrays obtained by retaining one layer and reflecting it throughout
the straightened torus. Next apply the same inequality to each such
array in direction two, and then in direction three. The reflected
array after these three steps is the full fold of the integrand from
one original cell. The three exponents multiply to
\[
 \frac h{p_1}\frac h{p_2}\frac h{p_3}
       =\frac1N.
\]
Thus the last product consists of one factor
$\|G_\kappa\|_{\mathbb T_3}$ for each original cell, where
$\mathbb T_3$ is the final rectangular torus.

We verify that all three arrays have consistent periods and that the
same cells represent them. If $\Gamma_0$ is the initial lattice, then
straightening successively in the frame gives
\begin{align*}
 \Gamma_1&=(\Gamma_0\cap\{\ell^1=0\})\oplus\Z\gamma_1,\\
 \Gamma_2&=(\Gamma_0\cap\{\ell^1=\ell^2=0\})
                                  \oplus\Z\gamma_1\oplus\Z\gamma_2,\\
 \Gamma_3&=\Z\gamma_1\oplus\Z\gamma_2\oplus\Z\gamma_3,
 \qquad \gamma_a=p_an_a/|n_a|^2.
\end{align*}
The kernel in the second line is a rank-one group along $n_3$; its
positive generator is $\gamma_3$, by the definition of the third
height period. Characters already fixed at $\gamma_1,\gamma_2$ are
retained in subsequent cross-sections, so all three final characters
are $-1$. The brick $0\le\kappa_a<p_a/h$ represents $\Gamma_3$.
It represents $\Gamma_2$ as well; then its section $\kappa_2=0$
represents the second cross-section, and its section $\kappa_1=0$
represents the first one. Restoring the normal coordinate in each
case gives representatives for the preceding lattice. These groups
all have the same covolume.

We also check the stronger box embedding condition used in the local
classes. If $B=y+[-r,r]^3$ lies in an open cell, then
$2|n_a|^2r<h$ for every $a$, by taking the range of $\ell^a$ on $B$.
For a nonzero period $\gamma$, at least one height is nonzero and
$|\ell^a(\gamma)|\ge2h$. Consequently
$|\gamma|_\infty>4r$, and the systole also gives
$|\gamma|_\infty\ge4$. These inequalities imply
$|\gamma|_\infty>2r+2$: use the second if $r<1$, the first if
$r\ge1$. Thus translates of $B+[-1,1]^3$ are disjoint. Boxes in
different cell orbits cannot meet after period translations, since
the open cells themselves are disjoint. It follows that every layer
belongs to its slab class. This justifies the three applications.

For a rectangular torus, every straightening leaves the torus unchanged,
which proves \eqref{eq:slab-chessboard-rectangular}. For a cubic torus
with even side $n$, the diagonal operations are
\begin{equation}\label{eq:slab-three-untwistings}
 C(n,n,n)\xrightarrow{\ell_-}D'(n/2,n,n)
          \xrightarrow{\ell_+}D(n/2,n,n),
 \qquad D(n/2,n,n)\xrightarrow{x_1}C(n/2,2n,n).
\end{equation}
For the first operation the cross-section is generated by $ne_+,ne_3$,
with characters $+1,-1$; its height period is $n$, and straightening
replaces the closing vector $ne_1$ by $(n/2)e_-$. For the second, the
cross-section is generated by $(n/2)e_-,ne_3$, both with character
$-1$. Its closing vector $ne_+$ already points along the normal,
and the operation changes its character from $+1$ to $-1$. For the
third, the cross-section is generated by $2ne_2,ne_3$, both with
character $-1$; the height period is $n/2$ and the straight closing
vector is $(n/2)e_1$. The lattices have systole at least four.

The first two diagonal operations, followed by the unchanged $x_3$
operation, have height periods $(n,2n,n)$. Hence they give
\eqref{eq:slab-chessboard-diagonal}; starting at $D'$ instead, use
the order $(\ell_+,x_3,\ell_-)$ to obtain the same final fold.
Likewise, starting at $D(a,2a,c)$ and using coordinate directions
gives periods $(a,4a,c)$ and ends at $C(a,4a,c)$, proving
\eqref{eq:slab-chessboard-coordinate}.
\end{proof}

\subsection{Pressure and the comparison error}\label{subsec:slab-pressure}

Write $C(n)=C(n,n,n)$ and
\[
 f(n_1,n_2,n_3)=\frac1{n_1n_2n_3}\log Z_{C(n_1,n_2,n_3)}[1].
\]
The mass and charge in this function are held fixed. The following
error will compare local expectations in the cubic and diagonal
geometries.

\begin{proposition}\label{prop:slab-pressure}
The pressure $f$ is nonincreasing in each side, and
\[
 f(n_1,n_2,n_3)\longrightarrow f_\infty:=\inf f\in\R
                     \quad\text{as }\min_\mu n_\mu\longrightarrow\infty.
\]
For every even $n\ge8$,
\begin{equation}\label{eq:slab-partition-chain}
 Z_{C(n)}[1]\le Z_{D'(n/2,n,n)}[1]
       \le Z_{D(n/2,n,n)}[1]\le Z_{C(n/2,2n,n)}[1].
\end{equation}
With $\Delta(n)=f(n/2,2n,n)-f(n,n,n)$,
\begin{equation}\label{eq:slab-pressure-error}
 0\le\Delta(n)\le f(n/2,n/2,n/2)-f_\infty\longrightarrow0.
\end{equation}
The logarithms of any two partition functions in
\eqref{eq:slab-partition-chain} differ in absolute value by at most
$n^3\Delta(n)$.
\end{proposition}

\begin{proof}
Fix the cross-section with sides $n_2,n_3$ and let
$Y=\mathfrak S[1]:\mathcal H(0)\to\mathcal H(1/2)$ be the coordinate
half-step. The first identity of \eqref{eq:slab-fold-holder}, applied
to $2n_1$ half-slabs, gives
\begin{equation}\label{eq:slab-pressure-trace}
 Z_{C(n_1,n_2,n_3)}[1]=\operatorname{Tr}(Y^*Y)^{n_1}.
\end{equation}
The positive trace-class operator $Y^*Y$ is independent of $n_1$ and
is nonzero, since the partition function is positive. Consequently
\[
 f(n_1,n_2,n_3)=\frac1{n_2n_3}\log\|Y^*Y\|_{n_1}.
\]
Schatten norms decrease as their exponent increases, so $f$ decreases
in $n_1$. Permuting the time direction proves the other two claims.
The same trace identity shows separate convexity of $\log Z$ in the
side lengths: H\"older's inequality gives
$\operatorname{Tr}A^{(1-t)s+ts'}\le
 (\operatorname{Tr}A^s)^{1-t}(\operatorname{Tr}A^{s'})^t$ for $A=Y^*Y$.

The lower partition bound \eqref{eq:reg-partition-bounds} gives the
volume-independent floor
\[
 f\ge\log C_*+3\log r_q-6r_q^2.
\]
Together with the upper bound in the same equation, it makes
$f_\infty$ finite. Monotonicity now gives convergence in any manner
with all sides tending to infinity: choose a finite triple at which
$f<f_\infty+\delta$ and then exceed its three sides.

Apply \eqref{eq:slab-fold-holder} with constant integrands to each of
the three operations \eqref{eq:slab-three-untwistings}. A fold of
the constant is the constant, so $Z_{\mathbb T}[1]\le Z_{\mathbb T^0}[1]$.
This proves the chain. Each torus in it has volume $n^3$, so its first
and last logarithms differ by $n^3\Delta(n)$. Side monotonicity gives
$f(n/2,2n,n)\le f(n/2,n/2,n/2)$, while
$f(n,n,n)\ge f_\infty$, proving \eqref{eq:slab-pressure-error}.
All intermediate logarithms lie between the endpoints, which proves
the last assertion.
\end{proof}

Thus the estimate concerns the logarithm per unit volume. It does not
give a bounded ratio of partition functions: the retained bound is
$e^{n^3\Delta(n)}$, with $\Delta(n)\to0$. The local comparison
divides the logarithmic error bound $n^3\Delta(n)$ by the number of
copies of the observable and by the perturbation strength.

\subsection{Coordinate slices and the all-flux trace}\label{subsec:slab-coordinate-trace}

For Gaussian domination we use a slice representation with real
spatial connections. Take direction $1$ as time on
$C_{\mathbf n}=C(n_1,n_2,n_3)$, $\mathbf n=(n_1,n_2,n_3)$,
$n_\mu\ge4$, and let
$\Sigma=\Z/n_2\Z\times\Z/n_3\Z$. Its links and plaquettes are
$\Sigma^{(1)}$ and $\Sigma^{(2)}$, its coboundary is $d_s$, and
$v=|\Sigma|=n_2n_3$. A connection has spatial slices $A_j$ and
temporal links $a_j$, and its curvature has the parts
\begin{equation}\label{eq:slab-slice-curvatures}
 F_j^{\mathrm{el}}=A_{j+1}-A_j-d_sa_j+r_qz_j^{\mathrm{el}},
 \qquad F_j^{\mathrm{mag}}=d_sA_j+r_qz_j^{\mathrm{mag}}.
\end{equation}
Closedness reads
$z_{j+1}^{\mathrm{mag}}-z_j^{\mathrm{mag}}=d_sz_j^{\mathrm{el}}$.

We first unfold the integer variables. Put
$k_0=0$, $k_{j+1}=k_j+z_j^{\mathrm{el}}$. Closedness is equivalent to
\[
 z_j^{\mathrm{mag}}=z_0^{\mathrm{mag}}+d_sk_j,
                       \qquad d_sk_{n_1}=0.
\]
The integer shift with spatial values $k_j$ and zero temporal values
removes all electric integers except the last one, $k_{n_1}$, and
changes every spatial slice to $A_j+r_qk_j$. Summing $k_j$ while
integrating $A_j$ over its compact box unfolds the intermediate
slices to all of $\R^{\Sigma^{(1)}}$. The last integer remains as a
winding vector
\[
 g=r_qk_{n_1}\in\Gamma_s,
 \qquad\Gamma_s=r_q(\Z^{\Sigma^{(1)}}\cap\ker d_s).
\]

Fix a plaquette $\sigma_0$ of $\Sigma$ and write
$k=\sum_\sigma z_0^{\mathrm{mag}}(\sigma)$. Then
$z_0^{\mathrm{mag}}-k\mathbf1_{\sigma_0}=d_sk'$ for an integer
one-cochain $k'$. Integer exactness follows from the dual graph whose
vertices are the plaquettes. A link of
$\Sigma$ is a dual edge directed toward the plaquette where its
coefficient in $d_s$ is positive. The dual graph is connected. Fix
a root and a spanning tree. Unit signed flow along the path from the
root to a vertex $\sigma$ has divergence
$\delta_\sigma-\delta_{\sigma_0}$. Summing these paths with the
integer coefficients of a zero-sum two-cochain gives precisely that
cochain as its divergence, hence as $d_sk'$. Two choices of $k'$
differ by an integer element of $\ker d_s$.

Shift every slice by $r_qk'$. The magnetic integers now equal
$k\mathbf1_{\sigma_0}$ in every slice. Summation over $k'$ modulo
closed integer cochains replaces the compact box for $A_0$ by a
fundamental domain of $\Gamma_s$. We arrive at paths
\[
 A_0\in\R^{\Sigma^{(1)}}/\Gamma_s,\qquad
 A_j\in\R^{\Sigma^{(1)}}\ (0<j<n_1),\qquad
 A_{n_1}=A_0+g,\quad g\in\Gamma_s,
\]
with $a_j\in[0,r_q)^{\Sigma}$ and flux $k\in\Z$. Their curvatures are
\begin{equation}\label{eq:slab-unfolded-curvatures}
 F_j^{\mathrm{el}}=A_{j+1}-A_j-d_sa_j,
 \qquad F_j^{\mathrm{mag}}=d_sA_j+r_qk\mathbf1_{\sigma_0}.
\end{equation}
Thus integration over the original configuration quotient equals
summation over $k,g$ and integration over these paths. For a
nonnegative integrand the identity follows by Tonelli's theorem. For
an absolutely integrable complex integrand it follows by applying
that identity to its modulus and then using Fubini's theorem. Each
change above preserves link phases and curvature, so it applies to
the determinant and to all curvature observables used here.

There are two finite Fock spaces in this representation: $\mathcal F'$
at a coordinate site plane and $\mathcal F$ at a link plane. The
matrix of the open half-slab gives a map
$N(A):\mathcal F'\to\mathcal F$, and we include both crossing weights
by writing
\[
 \widehat N(A)=\omega^{1/2}N(A)(\omega')^{1/2}.
\]
Let $u(a)$ be the diagonal unitary on $\mathcal F$ with entries
$e^{iq\langle a,\mathfrak q(I)\rangle}$. The fermion trace is
\begin{equation}\label{eq:slab-slice-determinant}
 \det \mathcal D_A^{C_{\mathbf n}}=
    \tr_{\mathcal F'}\prod_{j=n_1-1,\ldots,0}
             \bigl(\widehat N(A_{j+1})^*u(a_j)\widehat N(A_j)\bigr).
\end{equation}
This formula also holds on an unfolded path, since
$A_{n_1}-A_0\in\Gamma_s$ leaves $\widehat N$ unchanged.

To derive \eqref{eq:slab-slice-determinant}, split the torus at each
site and link plane. The half-slab above a link plane is the reflected
half-slab below it, hence its matrix is $N(A_{j+1})^*$. The crossing
at a site plane has no link phase and gives $\omega'$. At a link
plane it gives $\omega^{1/2}u(a_j)\omega^{1/2}$. Compose the matrices
and close antiperiodically; the closing parity is the identity.
Cyclicity moves half of the final $\omega'$ to the front, giving
\eqref{eq:slab-slice-determinant}. Every entry of $\widehat N(A)$ is
a polynomial in the link phases: the half-slab action is quadratic
in finitely many generators, its exponential terminates, and
Berezin integration extracts a coefficient. Consequently
\begin{equation}\label{eq:slab-half-matrix-bound}
 C_N:=\sup_{A\in\R^{\Sigma^{(1)}}}\|\widehat N(A)\|<\infty,
 \qquad\widehat N(A+r_qk')=\widehat N(A)
       \quad(k'\in\Z^{\Sigma^{(1)}}).
\end{equation}

On real slice connections the two Hilbert spaces are
\[
 H^s=\bigoplus_{k\in\Z}
      L^2(\R^{\Sigma^{(1)}}/\Gamma_s;\mathcal F'),\qquad
 H^l=\bigoplus_{k\in\Z}
      L^2(\R^{\Sigma^{(1)}}/\Gamma_s;\mathcal F).
\]
All quotient measures are induced by Lebesgue measure. Define commuting
self-adjoint electric operators $E_b$ by
$e^{-i\langle x,E\rangle}\Phi(A)=\Phi(A-x)$. On $H^l$ let
\begin{equation}\label{eq:slab-electric-operators}
 U_G(a)\Phi(A)=u(a)\Phi(A-d_sa),\qquad
 P_G=r_q^{-v}\int_{[0,r_q)^{\Sigma}}U_G(a)\,\dd a,
 \qquad K=r_q^ve^{-|E|^2/2}P_G.
\end{equation}
The integer mode charges make $U_G$ a representation of the compact
gauge group $\R^{\Sigma}/r_q\Z^{\Sigma}$. Its average $P_G$ is the
orthogonal projection onto invariant vectors. It commutes with every
function of $E$, because $u(a)$ acts on the Fock index and all spatial
translations commute. Hence
\begin{equation}\label{eq:slab-electric-square-root}
 K^{1/2}=r_q^{v/2}e^{-|E|^2/4}P_G.
\end{equation}

For a real slice two-cochain $b$, define the multiplication operator
$M_b:H^s\to H^l$ by
\begin{equation}\label{eq:slab-magnetic-multiplier}
 (M_b\Phi)_k(A)=
     e^{-\|d_sA+r_qk\mathbf1_{\sigma_0}+b\|^2/4}
                                      \widehat N(A)\Phi_k(A),
 \qquad X_b=K^{1/2}M_b.
\end{equation}
The magnetic weight is bounded and $\widehat N$ is bounded, so $M_b$
is bounded. We must prove that $X_b$ is Hilbert--Schmidt; the real
quotient is not compact and there are infinitely many flux sectors.

The periodized Gaussian
\[
 \widetilde g_{1/2}([x])
   =\sum_{\gamma\in\Gamma_s}\prod_{b\in\Sigma^{(1)}}
                       \pi^{-1/2}e^{-(x_b+\gamma_b)^2}
\]
has integral one on the quotient and is bounded. Indeed
$\Gamma_s\subset r_q\Z^{\Sigma^{(1)}}$, so the series is dominated
by the product of one-dimensional Gaussian sums on that larger lattice.
It is therefore in $L^1\cap L^\infty\subset L^2$. The kernel of
$K^{1/2}$ is
\[
 K^{1/2}(A',A)=r_q^{-v/2}\int_{[0,r_q)^{\Sigma}}
       \widetilde g_{1/2}([A'-A-d_sa])u(a)\,\dd a.
\]
Its operator norm is bounded by the nonnegative function
\[
 \widehat g([x])=r_q^{-v/2}\int_{[0,r_q)^{\Sigma}}
                  \widetilde g_{1/2}([x-d_sa])\,\dd a,
\]
which has finite $L^1$ and $L^\infty$ norms and hence finite $L^2$ norm.
Integrating the square of the kernel of $X_b$ over $A'$ now gives
\begin{equation}\label{eq:slab-shifted-hs}
 \|X_b\|_2^2\le
    (\dim\mathcal F')C_N^2\|\widehat g\|_2^2
       \sum_{k\in\Z}\int_{\R^{\Sigma^{(1)}}/\Gamma_s}
             e^{-\|d_sA+r_qk\mathbf1_{\sigma_0}+b\|^2/2}\,\dd A.
\end{equation}

To bound this last sum, split the quotient as
\[
 (\ker d_s/\Gamma_s)\times(\ker d_s)^\perp.
\]
The first factor is compact: $\ker d_s$ is the kernel of an integer
matrix, hence has an integer basis, and $\Gamma_s$ has full rank in
it. On the second factor $d_s$ is injective, so
$\|d_sA\|^2$ is positive definite. Moreover
$\ker d_s^*$ consists of constant plaquette functions. This follows
from the two equations
\[
 (d_s^*b)(y,2)=b(y)-b(y-e_3),\qquad
 (d_s^*b)(y,3)=b(y-e_2)-b(y).
\]
Thus $\operatorname{ran}d_s$ is the orthogonal complement of the
constant two-cochain. The orthogonal constant component of
$r_qk\mathbf1_{\sigma_0}+b$ has squared norm
$v^{-1}(r_qk+\sum_\sigma b_\sigma)^2$. Translate the variable in
$(\ker d_s)^\perp$ to remove the remaining component. We obtain
\begin{equation}\label{eq:slab-flux-gaussian}
 \sum_k\int e^{-\|d_sA+r_qk\mathbf1_{\sigma_0}+b\|^2/2}\,\dd A
    =C_s\sum_{k\in\Z}
               e^{-(r_qk+\sum_\sigma b_\sigma)^2/(2v)}<\infty,
\end{equation}
where
$C_s=\operatorname{vol}(\ker d_s/\Gamma_s)
\int_{(\ker d_s)^\perp}e^{-\|d_sA\|^2/2}\,\dd A<\infty$.
This constant is independent of $k$ and $b$. Equations
\eqref{eq:slab-shifted-hs}--\eqref{eq:slab-flux-gaussian} prove the
Hilbert--Schmidt property for every real magnetic perturbation. The
bounds here depend on the fixed cross-section; their purpose is to
justify the trace argument.

For a real external two-cochain $\beta$ on $C_{\mathbf n}$, put
\begin{equation}\label{eq:slab-external-partition}
 Z_{C_{\mathbf n}}(\beta)=\int_{\operatorname{Conf}(C_{\mathbf n})}
          e^{-\|F+\beta\|^2/2}\det \mathcal D_A^{C_{\mathbf n}}\,\dd(A,z).
\end{equation}
Let $\beta_j^{\mathrm{el}}$ and $\beta_j^{\mathrm{mag}}$ be its
electric and magnetic slices. Then
\begin{equation}\label{eq:slab-external-trace}
 Z_{C_{\mathbf n}}(\beta)=(2\pi)^{n_1n_2n_3}
       \operatorname{Tr}_{H^s}\prod_{j=n_1-1,\ldots,0}
            (X_{j+1}^*U_jX_j),
 \quad X_j=X_{\beta_j^{\mathrm{mag}}},\quad
 U_j=e^{i\langle\beta_j^{\mathrm{el}},E\rangle},
 \quad X_{n_1}=X_0.
\end{equation}

Here is the conversion from the unfolded path integral. The electric
Maxwell factors are the kernels
\[
 e^{-|A_{j+1}-A_j-d_sa_j+\beta_j^{\mathrm{el}}|^2/2}
       =(2\pi)^v g_1^{\otimes 2v}
                       (A_{j+1}-A_j-d_sa_j+\beta_j^{\mathrm{el}}).
\]
After summing the winding vector they are periodized on $\Gamma_s$.
Gaussian convolution with a translated kernel is
$e^{i\langle\beta_j^{\mathrm{el}},E\rangle}e^{-|E|^2/2}$.
Integration over $a_j$ includes the matrix $u(a_j)$ from
\eqref{eq:slab-slice-determinant}; it is exactly the gauge average in
\eqref{eq:slab-electric-operators}. The magnetic Gaussian in slice
$j$ splits into the two quarter-Gaussians in $M_j$ and $M_j^*$.
The determinant half-slab matrices supply their other factors.
Thus the cyclic path integral is the trace on $H^l$ of
\[
 \prod_{j=n_1-1,\ldots,0}
       (K^{1/2}U_jK^{1/2}M_jM_j^*).
\]
Each time slice contributed $(2\pi)^v$, giving the constant in
\eqref{eq:slab-external-trace}. Move the final $M_0^*$ to the front
and regroup the products as $X_{j+1}^*U_jX_j$. The
Hilbert--Schmidt property above and the absolute Gaussian path
majorant justify this cyclicity and every interchange. In particular
the external-field integral is finite; alternatively this follows
directly by absorbing a linear factor into a smaller Gaussian.

We spell out the path and kernel identification used here. On a
fundamental domain write every intermediate real slice as
$A_j=A_j^\flat+g_j$, $g_j\in\Gamma_s$, and set $g_0=0$ and
$g_{n_1}=A_{n_1}-A_0$. Each unperiodized electric kernel then has
the argument
$A_{j+1}^\flat-A_j^\flat+g_{j+1}-g_j-d_sa_j$.
The map $(g_1,\ldots,g_{n_1})\mapsto
(g_{j+1}-g_j)_{0\le j<n_1}$ is a bijection of $\Gamma_s^{n_1}$.
Summing these differences independently is therefore precisely
periodization of every cyclic kernel. The magnetic factors and
$\widehat N$ are unchanged by the shifts $g_j$. The finite Fock trace
is changed from $\mathcal F'$ to $\mathcal F$ by moving its last
$\widehat N(A_0)^*$ to the front; periodicity identifies it with
$\widehat N(A_{n_1})^*$.

For absolute convergence replace every matrix factor by its norm,
and bound each shifted electric Gaussian by a fixed constant times
$e^{-|A'-A|^2/4}$ and each shifted magnetic factor by a constant
times $e^{-\|F^{\mathrm{mag}}\|^2/4}$. This defines nonnegative
scalar kernels. Their periodized convolutions are bounded and in
$L^2$, and multiplication by the magnetic Gaussian is square
integrable after summing fluxes by \eqref{eq:slab-flux-gaussian}.
Hence each scalar majorant operator is Hilbert--Schmidt. The
integral of their cyclic product is finite: group the first factor
against the product of the others and use the Hilbert--Schmidt
Cauchy--Schwarz inequality. Tonelli first gives this conclusion
for the nonnegative majorants, and Fubini then permits the kernel
expansions and periodization changes for the original integrand.
This proves the trace identity without a formal interchange of
the infinite flux sum and the operator trace.

\subsection{Gaussian domination}\label{subsec:slab-gaussian}

\begin{theorem}\label{thm:slab-gaussian}
For every coordinate torus with sides at least four and every real
two-cochain $\beta$,
\begin{equation}\label{eq:slab-gaussian-partition}
 0<Z_{C_{\mathbf n}}(\beta)\le Z_{C_{\mathbf n}}(0).
\end{equation}
Consequently its normalized curvature law obeys
\begin{equation}\label{eq:slab-gaussian-domination}
 \E_{C_{\mathbf n}}e^{\langle\beta,F\rangle}
                              \le e^{\|\beta\|^2/2}.
\end{equation}
For real two-cochains $\beta_i$ and complex numbers $z_i$,
\begin{equation}\label{eq:slab-complex-gaussian}
 \left|\E_{C_{\mathbf n}}e^{\sum_i z_i\langle\beta_i,F\rangle}\right|
       \le\exp\left(\tfrac12
                    \left\|\sum_i(\operatorname{Re}z_i)\beta_i\right\|^2\right).
\end{equation}
The constants are independent of the volume, the charge, and the mass.
\end{theorem}

\begin{proof}
The determinant is nonnegative and positive on a set of positive
configuration measure, since $Z_{C_{\mathbf n}}(0)>0$. The shifted photon
density is strictly positive everywhere, proving the strict lower
bound in \eqref{eq:slab-gaussian-partition}.

In \eqref{eq:slab-external-trace} move $X_0^*$ through the trace to
obtain a trace on $H^l$:
\[
 (2\pi)^{-n_1n_2n_3}Z_{C_{\mathbf n}}(\beta)
       =\operatorname{Tr}_{H^l}
          (U_{n_1-1}X_{n_1-1}X_{n_1-1}^*\cdots U_0X_0X_0^*).
\]
Every $U_j$ is unitary. Apply Schatten's inequality with all exponents
equal to $n_1$; the trace on the left is nonnegative. Since $X_jX_j^*$
and $X_j^*X_j$ have the same nonzero eigenvalues, this gives
\begin{equation}\label{eq:slab-magnetic-reduction}
 Z_{C_{\mathbf n}}(\beta)
       \le\prod_{j=0}^{n_1-1}
            \bigl((2\pi)^{n_1n_2n_3}
                    \operatorname{Tr}(X_j^*X_j)^{n_1}\bigr)^{1/n_1}
       =\prod_{j=0}^{n_1-1}Z_{C_{\mathbf n}}(\beta^{(j)})^{1/n_1}.
\end{equation}
Here $\beta^{(j)}$ has zero electric part and the same magnetic
part $\beta_j^{\mathrm{mag}}$ in every time slice. The equality uses
the trace formula again, with the same $X_j$ in each slice and all
unitary factors equal to the identity.

Each $\beta^{(j)}$ is supported only on $23$-plaquettes. Interchange
directions $1$ and $2$. The torus becomes $C(n_2,n_1,n_3)$ and the
perturbation is supported only on $13$-plaquettes, so its magnetic
part for the new time direction is zero. The reduction
\eqref{eq:slab-magnetic-reduction} on this torus bounds its partition
function by the unperturbed one. Lattice isometry covariance of the
determinant and photon measure identifies these partition functions
with $Z_{C_{\mathbf n}}(\beta^{(j)})$ and $Z_{C_{\mathbf n}}(0)$. Substitution in
\eqref{eq:slab-magnetic-reduction} proves the upper bound in
\eqref{eq:slab-gaussian-partition}.

Replacing $\beta$ by $-\beta$ and completing the square gives
\eqref{eq:slab-gaussian-domination}. For complex parameters, take
the absolute value inside the probability integral and apply that
bound to $\sum_i(\operatorname{Re}z_i)\beta_i$, proving
\eqref{eq:slab-complex-gaussian}.
\end{proof}

We record the consequences used below. Charge conjugation centers
$F(\beta)=\langle\beta,F\rangle$. The Gaussian bound and the
nonnegative Taylor coefficients of $\cosh$ give
\begin{align}\label{eq:slab-gaussian-moments}
 \E F(\beta)^2&\le\|\beta\|^2,\notag\\
 \E F(\beta)^{2r}&\le(2r\exp(1)\|\beta\|^2)^r,\qquad r\ge1,\\
 \E F(\beta)^4&\le\tfrac32\exp(2)\|\beta\|^4.\notag
\end{align}
For the second moment let the parameter in the quadratic term tend
to zero. For the $2r$-th moment use
$(2r)!s^{-2r}e^{s^2\|\beta\|^2/2}$ and choose
$s^2=2r/\|\beta\|^2$. The fourth-moment constant uses $4!=24$
before estimating the factorial. H\"older and Cauchy--Schwarz in the
fourth-cumulant identity then give
\begin{equation}\label{eq:slab-gaussian-cumulant}
 |\kappa_4(F(\beta_1),\ldots,F(\beta_4))|
                              \le15\prod_{i=1}^4\|\beta_i\|.
\end{equation}

For nonzero $\beta_i$ with disjoint supports and thresholds
$\xi_i\ge0$, Markov's inequality
applied to $\sum_i\pm\xi_i\beta_i/\|\beta_i\|$, for each choice of
signs, gives
\begin{equation}\label{eq:slab-gaussian-joint-tail}
 \Pr\{|F(\beta_i)|\ge\xi_i\|\beta_i\|\text{ for all }i\}
                        \le\prod_i2e^{-\xi_i^2/2}.
\end{equation}
Disjointness makes the squared norm of this test equal to
$\sum_i\xi_i^2$. Finally, if a finite-support lattice test embeds in
the torus, its periodization has the same norm and the same pairing
with the periodic curvature. Thus \eqref{eq:slab-gaussian-domination}
also applies to these local tests on the lifted field law.

\subsection{The finite-volume fold inequality}\label{subsec:slab-finite-fold}

We give the local form of the comparison that will be used to obtain
diagonal positivity. Fix diagonal cells of even height width $h$,
and let $O$ be a real trigonometric polynomial in finitely many
curvature values at plaquettes contained in the base cell $Q_0$.
Thus $O$ is bounded and depends only on this open cell. Let $n\ge8$
be divisible by $4h$, and abbreviate $D_n=D(n/2,n,n)$. Its cell
count, and that of the cubic torus, is
\[
 N_n=\frac{n^3}{|Q|}=\frac{2n^3}{h^3}.
\]
The translations between cells form the integer lattice
\[
 \Lambda_h=\frac h2\Z e_-+\frac h2\Z e_++h\Z e_3.
\]
On $C(n)$ form the translated copies, and on $D_n$ the reflected
copies:
\begin{equation}\label{eq:slab-two-arrays}
 \widetilde O_n(F)=\sum_{v\in\Lambda_h/n\Z^3}O(\tau_v^*F),
 \qquad O_n^{\mathrm{fold}}(F)=\sum_{\kappa\bmod\Gamma_{D_n}}
                                                O(g_\kappa^*F).
\end{equation}
Each sum has $N_n$ terms. The first sum is read on cubic-periodic
curvatures, the second on $D_n$-periodic curvatures. Here $g_\kappa$
is the reflected map from the base cell to $Q_\kappa$.

Let $\mu_{C(n)}$ and $\mu_{D_n}$ be the normalized curvature laws.
For $t>0$ define the tilted diagonal law
\begin{equation}\label{eq:slab-tilted-law}
 \mu_{D_n}^t(\Phi)=
 \frac{\mu_{D_n}(e^{tO_n^{\mathrm{fold}}}\Phi)}
      {\mu_{D_n}(e^{tO_n^{\mathrm{fold}}})}.
\end{equation}
The folded array is invariant under the group generated by the cell
walls: that group permutes the cells of $D_n$ and their reflected
copies. Thus the tilt retains these symmetries.

\begin{proposition}\label{prop:slab-finite-fold}
For the preceding data and every $t>0$,
\begin{equation}\label{eq:slab-finite-fold}
 \mu_{C(n)}(O)\le\mu_{D_n}^t(O)+\frac{|Q|\Delta(n)}{t}.
\end{equation}
The error is the fixed-parameter pressure error of
\eqref{eq:slab-pressure-error}.
\end{proposition}

\begin{proof}
Jensen's inequality and translation invariance on $C(n)$ give
\begin{equation}\label{eq:slab-fold-lower-tangent}
 \log\frac{Z_{C(n)}[e^{t\widetilde O_n}]}{Z_{C(n)}[1]}
                     \ge t\mu_{C(n)}(\widetilde O_n)
                     =tN_n\mu_{C(n)}(O).
\end{equation}
The diagonal chessboard inequality gives
\begin{equation}\label{eq:slab-fold-change}
 Z_{C(n)}[e^{t\widetilde O_n}]
               \le\|e^{tO}\|_{D_n}^{N_n}
               =Z_{D_n}[e^{tO_n^{\mathrm{fold}}}].
\end{equation}
To check that its $N_n$ factors have the same seminorm, translate
the base observable by $v\in\Lambda_h$. Translation permutes the cells
and conjugates their reflected maps, so the fold of the translated
observable is the translate of the original full fold. The law on
$D_n$ is translation invariant, hence its partition function and
the seminorm are unchanged. All observables here are bounded scalar
functions supported in distinct open cells, so the chessboard
inequality applies directly.

Set $L(t)=\log\mu_{D_n}(e^{tO_n^{\mathrm{fold}}})$. Boundedness
permits differentiation, giving
$L'(t)=\mu_{D_n}^t(O_n^{\mathrm{fold}})$. Convexity and $L(0)=0$
give $L(t)\le tL'(t)$. Wall-group invariance of the tilted law makes
the expectation of every reflected copy equal to that of $O$; hence
\begin{equation}\label{eq:slab-fold-upper-tangent}
 \log\frac{Z_{D_n}[e^{tO_n^{\mathrm{fold}}}]}{Z_{D_n}[1]}
                          \le tN_n\mu_{D_n}^t(O).
\end{equation}
Combining \eqref{eq:slab-fold-lower-tangent}--\eqref{eq:slab-fold-upper-tangent}
leaves, in addition, the normalization difference
\[
 \log Z_{D_n}[1]-\log Z_{C(n)}[1]\le n^3\Delta(n).
\]
Division by $tN_n$ proves \eqref{eq:slab-finite-fold}.
\end{proof}

For fixed $O$, its cell geometry, and $t>0$, the error in
\eqref{eq:slab-finite-fold} tends to zero as $n\to\infty$ along the
admissible sides. The estimate compares the cubic law with a tilted
diagonal law. The next step must control these tilted laws as
$t\downarrow0$; the pressure estimate alone does not identify their
local limits.


\section{Diagonal reflection positivity}\label{sec:diagonal-positivity}

The finite-volume fold inequality compares a cubic expectation with
an expectation in a tilted adapted law. At each wall of the chosen
cell structure, the tilted functional is reflection positive for
observables in the corresponding open positive half-period slab.
The walls are separated by the cell width, and this positivity does not yet place
the original volume limit among the diagonal-positive laws. We must
first show that removing the tilt also removes its residual dependence
on layers. The argument uses local operators inserted into a coordinate
crossing and then eliminates the fields outside its two adjacent
fermion layers. The resulting comparison states have enough translation
invariance to apply the fold inequality to an arbitrary local observable.
Convex separation will then give the required diagonal positivity.

Fix $q>0$, $0<\rho<1$, positive weights $w_i$ and a real mass $M$.
All volume limits in this section keep these parameters fixed. We
use cubic sides $n_i\to\infty$ such that eight divides all sufficiently
large sides and the set
\[
 \mathcal D=\{m\in\N:m\mid n_i\text{ for all sufficiently large }i\}
\]
is infinite. The set $\mathcal D$ is closed under divisors and least
common multiples. The cutoff sides $8L^{a+T+N}$ have this property
at fixed $N$, since $8L^j\in\mathcal D$ for each $j\ge0$.

\subsection{The tilted law at its walls}\label{subsec:diag-wall-law}

Fix diagonal cells with offsets $(c_0,c_+,c_3)$, even width $h$, and
$c_0,c_+\in\Z$, $c_3\in\Z+\tfrac12$, and $4h\in\mathcal D$.
Let $O$ be the real trigonometric polynomial used in the finite fold
comparison. Its support is contained in a box of radius $r_O$, and
we require $h>8r_O+8$. Such widths can be chosen because
$\mathcal D$ contains arbitrarily large multiples of eight. For
large divisible sides, put $D_i=D(n_i/2,n_i,n_i)$ and
$\mu_i^t=\mu_{D_i}^t$, as in \eqref{eq:slab-tilted-law}. Write
\[
 O(F)=\sum_\beta a_\beta e^{i\langle\beta,F\rangle},
 \qquad |O|_1=\sum_\beta|a_\beta|.
\]
The sum is over the distinct exponents in the chosen finite expansion.
In particular $|O|\le|O|_1$. Let $V_i=O_{n_i}^{\mathrm{fold}}$ and
$N_i=n_i^3/|Q|$. For an admissible Grassmann integrand we use
the same notation for the normalized functional
$\mu_i^t(G)=Z_{D_i}[e^{tV_i}G]/Z_{D_i}[e^{tV_i}]$.
On bosonic integrands it agrees with integration against the
tilted probability law. The normalizing denominator satisfies
\begin{equation}\label{eq:diag-tilt-denominator}
 Z_{D_i}[e^{tV_i}]\ge e^{-tN_i|O|_1}Z_{D_i}[1]>0,
 \qquad t\ge0.
\end{equation}
This uses only the positive bosonic measure; it makes no positivity
assertion about a Grassmann integrand.

At a wall $\Pi=(\ell^a,o_a+jh)$, the adapted torus has height period
$p_a$, with $(p_1,p_2,p_3)=(n_i,2n_i,n_i)$. The number $p_a/h$ is
even. Split the cells into the two half-slabs adjacent to this wall,
and define
\[
 \Phi_+(F)=\exp\left(t\sum_{Q_\kappa\subset
              \{o_a+jh<\ell^a<o_a+jh+p_a/2\}}O(g_\kappa^*F)\right),
\]
with representatives in one half-period. Reflection in the wall
maps these cells onto the other half-period modulo the normal closing
vector. Because the reflected copies of $O$ form one wall-invariant
array, on periodic curvatures
\begin{equation}\label{eq:diag-wall-factorization}
 e^{tV_i}=\Phi_+\Theta_\Pi\Phi_+.
\end{equation}
The function $\Phi_+$ is bounded, real, and supported strictly within
the upper half-slab. The support lies there because each plaquette
of a copy of $O$ is contained in its open cell.

Apply the adapted reflected representation to $\Phi_+G$ and
$\Phi_+G'$. For local even integrands $G,G'$ in the upper half-slab,
\begin{equation}\label{eq:diag-tilted-positive-form}
 \mu_i^t(\Theta_\Pi G'\,G)=
 \frac{\operatorname{Tr}
          (\mathfrak S[\Phi_+G']^*\mathfrak S[\Phi_+G])}
      {Z_{D_i}[e^{tV_i}]}.
\end{equation}
The numerator is the Hilbert--Schmidt inner product of the two slab
operators. Thus, on $D_i$, the tilted functional is reflection positive
at each cell wall for local even integrands supported in the corresponding
open positive half-period slab. Wall reflection also gives
$\mu_i^t(\Theta_\Pi G)=\overline{\mu_i^t(G)}$: the unperturbed
functional is reflection covariant, the scalar weight $e^{tV_i}$
is invariant, and its denominator is real and positive.

The bosonic law is invariant under the full group generated by the
cell walls. Indeed, that group preserves $D_i$ and permutes the
reflected copies in $V_i$, while the unperturbed bosonic law is
invariant under its lattice isometries. In particular the tilt is
invariant under translations by $he_-$, $he_+$ and $2he_3$.
These are the periods of the reflected cell array, rather than all
lattice translations.

To compare weak limits we use the Polish space
\[
 \Omega=\{F\in\R^{\{\text{oriented plaquettes of }\Z^3\}}:dF=0\},
\]
with its product topology. The pullback by a lattice reflection or
translation is continuous, so invariance passes to weak limits.
Positivity passes for each bounded continuous local positive-side test
once its support lies in the finite positive half-period slab for all
large volumes. Positivity against bounded Borel local tests follows once reflection invariance
has passed: on their finite-dimensional marginal approximate the
test in $L^2$ by compactly supported continuous functions. The
reflected-form difference is at most
\[
 \|\Phi-\Phi_k\|_2(\|\Phi\|_2+\|\Phi_k\|_2),
\]
because reflection preserves the $L^2$ norm. The Gaussian domination
bound also passes to weak limits, by truncating
$e^{\langle\beta,F\rangle}$ and then using monotone convergence.
These facts will be applied after obtaining tightness of the tilted
laws; their extensive weight alone gives no uniform moment estimate.

\subsection{Local operators at a coordinate link plane}\label{subsec:diag-local-operators}

An insertion must carry an operator norm that is uniform in the
cross-section and in electric characters. Fix a finite set $S$ of
links of $\Z^2$ and a nonempty finite set $\Lambda$ of sites
containing their endpoints. Let $\mathcal F_\Lambda$ be the Fock
space of crossing modes over $\Lambda$, and put
\[
 H_S=L^2(\R^S)\otimes\mathcal F_\Lambda,
 \qquad E_b=-i\partial_{A(b)}.
\]
The diagonal unitary $u(a)$ on $\mathcal F_\Lambda$ has integer mode
charges. An operator $\Upsilon$ has charge $z\in\Z^\Lambda$ when
$u(a)\Upsilon u(a)^{-1}=e^{iq\langle a,z\rangle}\Upsilon$.

Let $\mathfrak A(S)$ be the span of the functions
$e^{-Q(\xi_{S'})+\langle\lambda,\xi_{S'}\rangle}$, where
$S'\subset S$, $Q$ is real and positive definite on $\R^{S'}$,
and $\lambda\in\C^{S'}$. The empty subset gives the constant one.
These functions are bounded and uniformly continuous. Their span
is closed under products, conjugation and real translations: products
add the positive quadratic forms on the union of their supports,
while translation completes the same quadratic exponent.

For $\delta\in\R^S$, extended by zero to the full slice, require
$d_s^*\delta=qz(\delta)$ with $z(\delta)\in\Z^\Lambda$. The local
operator algebra consists of the finite sums
\begin{equation}\label{eq:diag-operator-normal-form}
 X=\sum_\delta e^{i\langle\delta,A\rangle}\Upsilon_\delta(E_S),
 \qquad\Upsilon_\delta\in\mathfrak A(S)\otimes
                      \operatorname{End}_{z(\delta)}(\mathcal F_\Lambda).
\end{equation}
We denote it by $\mathfrak W(S,\Lambda)$. The Weyl identity gives
\begin{align}\label{eq:diag-operator-rules}
 (e^{i\delta A}\Upsilon(E_S))(e^{i\delta'A}\Upsilon'(E_S))
       &=e^{i(\delta+\delta')A}\Upsilon(E_S+\delta')\Upsilon'(E_S),
          \notag\\
 (e^{i\delta A}\Upsilon(E_S))^*
       &=e^{-i\delta A}\Upsilon^*(E_S-\delta).
\end{align}
Here $\delta A$ abbreviates the real pairing. These formulas make
$\mathfrak W(S,\Lambda)$ a unital $*$-algebra. The normal form is
unique. In Fourier variables its terms send a function supported
near $\xi_0$ to functions supported near $\xi_0+\delta$; small
enough neighborhoods separate the finitely many distinct exponents.
If the sum vanishes as an operator, every continuous symbol vanishes
on each such neighborhood and hence everywhere.

Every $X$ commutes with the local gauge transformations
$\Phi(A)\mapsto u(a)\Phi(A-d_sa)$. The phase change
$e^{-i\langle\delta,d_sa\rangle}$ cancels the Fock charge phase.
The total charge of $z(\delta)$ is zero, since it is the divergence
of a finitely supported cochain. Each Fock term is therefore even.
For $\beta\in\R^S$, define
\begin{equation}\label{eq:diag-alpha}
 \alpha_\beta X=e^{i\langle\beta,A\rangle}Xe^{-i\langle\beta,A\rangle}
    =\sum_\delta e^{i\langle\delta,A\rangle}
                                      \Upsilon_\delta(E_S-\beta).
\end{equation}
This is an isometric $*$-automorphism on $H_S$.

For a rectangular cross-section $\Xi$ containing the smallest integer
rectangle spanned by $\Lambda$,
the same formulas define $X_\Xi$ on the real slice quotient space
$H^l$ of \S\ref{subsec:slab-coordinate-trace}. The Fock operators
are localized into the full crossing Fock space by $\iota$.
The localization is the tensor inclusion for the exterior space, with the reordering
sign and cross-section character attached to each representative;
it is a unital injective $*$-homomorphism. The phase
$e^{i\langle\delta,A\rangle}$ is well defined on the quotient.
Indeed, on the support hull an integer closed one-cochain is an
integer coboundary, by the box Poincare argument. Thus for
$\gamma\in\Gamma_s$,
$\langle\delta,\gamma\rangle\in2\pi\Z$, using
$d_s^*\delta=qz(\delta)$. The local Weyl formulas therefore remain
valid on the quotient. Uniqueness of the normal form makes the
representation independent of its presentation. It commutes with
$P_G$ and with functions of the electric variables outside $S$.

Define the norm needed for characters by
\begin{equation}\label{eq:diag-sharp-norm}
 \|X\|_\#=\sup_{\Xi,\beta}\| (\alpha_\beta X)_\Xi\|.
\end{equation}
The support hull must embed in the cross-section; the supremum is
over such cross-sections and all $\beta\in\R^S$. The electric spectrum
of a quotient can be smaller than $\R^S$, so its unshifted operator
norm need not equal the local norm.

\begin{proposition}\label{prop:diag-local-norm}
For $X\in\mathfrak W(S,\Lambda)$,
\[
 \|X\|_\#=\|X\|_{\mathcal L(H_S)}.
\]
The norm is unchanged by enlarging $S,\Lambda$. In particular it is
a $C^*$-norm, $\|1\|_\#=1$, and every quotient realization satisfies
$\|(\alpha_\beta X)_\Xi\|\le\|X\|_\#$.
\end{proposition}

\begin{proof}
First remove the Fock localization. Its tensor extension is an
injective $*$-homomorphism and hence isometric; we may work on
$L^2(\R^{\Sigma^{(1)}}/\Gamma_s)\otimes\mathcal F_\Lambda$ with the
original finite Fock operators. Call this realization $X_\Xi^0$.
Introduce an additional variable $z\in\R^S$, with electric operators
$E'=-i\partial_z$, on
\[
 L^2(\R^S_z)\otimes
 L^2(\R^{\Sigma^{(1)}}/\Gamma_s)\otimes\mathcal F_\Lambda.
\]
The operator $X(z,E')$ has norm $\|X\|$, since it is $X$ tensored
with the identity. Choose a Borel representative $A^\flat$ of the
quotient and use the two unitaries
\[
 (U_1\Phi)(z,[A])=\Phi(z+A^\flat|_S,[A]),\qquad
 (U_2\Phi)(z,[A])=\Phi(z,[A-z]),
\]
where $z$ is extended by zero outside $S$. The first conjugation
multiplies $e^{i\delta z}$ by $e^{i\delta A}$. The second replaces
$E'$ by $E'+E_S$ and cancels the factor $e^{i\delta z}$. Hence
\[
 U_2U_1X(z,E')(U_2U_1)^{-1}
       =\sum_\delta e^{i\delta A}\Upsilon_\delta(E_S+E').
\]
Fourier transform in $z$. This becomes a direct integral of
$T_p=(\alpha_{-p}X)_\Xi^0$, $p\in\R^S$. Uniform continuity of the
symbols makes $p\mapsto T_p$ norm continuous. Its direct-integral
norm is therefore $\sup_p\|T_p\|$: the reverse inequality follows
by using a unit vector nearly maximizing $T_{p_0}$ on a small ball
about $p_0$. It follows that
\[
 \|X\|=\sup_p\| (\alpha_{-p}X)_\Xi^0\|
        =\sup_p\| (\alpha_{-p}X)_\Xi\|.
\]
This holds for each admissible cross-section and proves the asserted
supremum identity. Enlarging $S$ tensors with an identity on the
additional $L^2$ factor; enlarging $\Lambda$ uses the same isometric
Fock inclusion. Both preserve the norm.
\end{proof}

\subsection{Insertions and a coordinate trace bound}\label{subsec:diag-insertions}

Let $c=j+\tfrac12$. Write $\operatorname{Ins}_c$ for the crossing
insertion map, and $\omega$ for its diagonal mode weights. For an
even matrix kernel $k(A',A)$ on $\R^S$, set
\begin{equation}\label{eq:diag-kernel-insertion}
 \begin{split}
 G[k](\widehat A)={}&\operatorname{Ins}_c\!\left[
   \omega^{1/2}k(\widehat A_{j+1}|_S,
           (\widehat A_j+d_s\widehat a_j)|_S)
             u(\widehat a_j)\omega^{1/2}\right]\\
 &\quad\cdot\prod_{b\in S}g_1(F_j^{\mathrm{el}}(b))^{-1}
                      \exp(S_c^\times(\Lambda)).
 \end{split}
\end{equation}
Here $S_c^\times(\Lambda)$ is the part of the fermion action crossing
the plane over $\Lambda$. The two factors on the second line remove
the original electric and fermion crossing weights on the support
of $k$. Thus $G[k]$ inserts $k$ in their place. Its definition uses
the infinite lattice and is independent of a torus.

We require
\begin{align}\label{eq:diag-kernel-conditions}
 u(a)^{-1}k(A'+d_sa,A+d_sa)u(a)&=k(A',A),\notag\\
 \|k(A',A)\|&\le C\exp(-\epsilon|A'-A|^2),\qquad\epsilon>0.
\end{align}
Choose a box containing the two copies of $S$, the temporal links
over $\Lambda$, and the unit cubes centered at $(j,y)$ and $(j+1,y)$,
$y\in\Lambda$. In what follows we use the cube $B^\#(c,S,\Lambda)$
obtained by taking twice the largest half-width of the smallest
rectangular box containing these sets, with the same center. This
choice is covariant under coordinate translations and reflections.
The insertion reads only these links and the crossing fields, so
it is supported in $B^\#$.

The insertion is an even, gauge invariant local integrand. Indeed,
an even kernel has entries only between occupation vectors of the
same parity, so every monomial
$\Theta_c\zeta_J\,\zeta_I$ in its insertion is even. If the gauge
parameter is $\lambda$, its upper and lower slice values change
the matrix in \eqref{eq:diag-kernel-insertion} to
\[
 u(\lambda_{j+1})^{-1}
   k(\widehat A_{j+1},\widehat A_j+d_s\widehat a_j)
       u(\widehat a_j)u(\lambda_j).
\]
The matrix entry acquires precisely the inverse of the phase of
$\Theta_c\zeta_J\,\zeta_I$. The electric curvature and the crossing
action are gauge invariant. Finally, the kernel estimate and
$g_1(x)^{-1}=\sqrt{2\pi}e^{x^2/2}$ bound every Grassmann coefficient
by
\begin{equation}\label{eq:diag-insertion-growth}
 C'\exp\left((\tfrac12-\epsilon')
                 \sum_{b\in S}F_j^{\mathrm{el}}(b)^2\right),
 \qquad\epsilon'=\min(\epsilon,\tfrac12).
\end{equation}
The other factors have bounded coefficients. This is the local
integrand bound of \S\ref{subsec:slab-local}, since its electric
plaquettes lie in the box.

For $X\in\mathfrak W(S,\Lambda)$ let $k_X^{(s)}$ be the continuous
kernel of $e^{-s|E_S|^2}Xe^{-s|E_S|^2}$, with $s=\tfrac14$ or
$\tfrac18$, and write
\begin{equation}\label{eq:diag-operator-insertions}
 G_X=G[k_X^{(1/4)}],\qquad G'_X=G[k_X^{(1/8)}].
\end{equation}
Both kernels satisfy \eqref{eq:diag-kernel-conditions}. To see their
decay, consider one normal-form term $e^{i\delta A}\Upsilon(E)$ with
$\Upsilon(\xi)=e^{-Q(\xi)+\lambda\xi}\Upsilon_0$. Identify $Q$
with its symmetric matrix, so $Q(\xi)=\langle\xi,Q\xi\rangle$,
extending $Q$ by
zero to the other links. With $Q_s=4sI+2Q$ and
$\lambda_s=\lambda-2s\delta$, its kernel is
\begin{equation}\label{eq:diag-regularized-kernel}
 \begin{split}
 k_X^{(s)}(A',A)={}&e^{i\delta A'}(2\pi)^{-|S|/2}
          (\det Q_s)^{-1/2}e^{-s|\delta|^2}\\
 &\cdot\exp\left(\tfrac12
  \langle\lambda_s+i(A'-A),Q_s^{-1}
                  (\lambda_s+i(A'-A))\rangle\right)\Upsilon_0.
 \end{split}
\end{equation}
The real part of this exponent has a strictly negative quadratic
form in $A'-A$, with a fixed linear term. It gives the required
Gaussian decay after decreasing its exponent. Gauge covariance
follows by conjugating the regularized operator by the gauge
unitary; equality of their continuous kernels gives the first
identity in \eqref{eq:diag-kernel-conditions}. Finite sums preserve
these properties. The identity operator has kernel
$\prod_{b\in S}g_1(A'(b)-A(b))$ when $s=\tfrac14$, so
$G_1=1$; unused links and modes cancel in the same way.

Reflection acts on these insertions by adjoints:
\begin{equation}\label{eq:diag-insertion-reflection}
 \Theta_cG_X=G_{X^*},\qquad\Theta_cG'_X=G'_{X^*}.
\end{equation}
The kernel is changed to $k(A,A')^*$. At the same time the two
slices exchange, $\widehat a_j$ changes sign, and covariance moves
$u(\widehat a_j)$ to the required side of the kernel. This gives
\eqref{eq:diag-insertion-reflection}, including its field signs.
We also need coordinate translations and spatial reflections to
preserve the operator norm, since each coordinate fold contains
such images.

For a spatial reflection $\theta$ in the second or third coordinate,
the spin matrix exchanges the two eigenspaces of $\gamma_1$.
Its unitary matrix $v$ on crossing modes preserves their weights
and reverses their charges. On exterior products the reordering
sign is $r_I=(-1)^{|I|(|I|-1)/2}$; hence the Fock unitary is
$U=(\bigwedge v)r$. The anticommutation of the two field reflection
maps contributes the sign $(-1)^{\deg}$, and moving the two mode
products past one another contributes $(-1)^{|I||J|}$. For the
parity-matched entries of an even kernel these signs leave exactly
$r_Ir_J$. Consequently
\[
 \Theta_\theta\operatorname{Ins}_c[\Upsilon(\widehat A)]
 =\operatorname{Ins}_c[
       U\overline{\Upsilon(\theta^*\widehat A)}U^*].
\]
The antiunitary
$(J\Phi)(A)=U\overline{\Phi(\theta_s^*A)}$ gives this kernel
transformation. It sends an oriented electric operator to the
negative of the reflected electric operator, and sends
$e^{i\delta A}$ to $e^{-i\langle\theta_s^*\delta,A\rangle}$.
Thus it preserves the charge constraint, the local algebra and
the heat factors $e^{-s|E|^2}$. Translations act by the corresponding
unitary relabeling. Together with adjoints at the link plane, these
facts show that each translate or signed coordinate reflection
of $G_X$ is an insertion $G_{X^g}$ with
$\|X^g\|=\|X\|$. The same holds for $G'_X$. The box $B^\#$
transforms with the insertion.

An even local integrand $G$ in a box $B$ satisfies the coordinate
bound $C\ge0$ if, on every rectangular coordinate torus
$T=(\Gamma,\chi)=C(n_1,n_2,n_3)$ with $n_\mu\ge4$,
for every nonempty finite family of images $g_aB$ under lattice
translations and signed coordinate reflections that embeds
disjointly in $\Gamma$, and every $\beta\in C^2(\Gamma)$,
\begin{equation}\label{eq:diag-coordinate-bound}
 \left|Z_T\left[\varphi_\beta\prod_{a=1}^rG^{g_a}\right]\right|
       \le C^rZ_T[1],\qquad
       \varphi_\beta(F)=e^{i\langle\beta,F\rangle},
\end{equation}
provided that, in each coordinate direction $\mu$, there is a
plaquette $\sigma$ orthogonal to $e_\mu$ such that none of the
boxes or their period translates contains any plaquette
$\sigma+k e_\mu$, $k\in\Z$. This condition places the flux
representative used in the slice trace on plaquettes read by
none of the insertions.
When the coordinate cell structure cuts $T$ and its walls are at
half-integers, boxes contained in its open cells satisfy the
free-column condition. Choose a plaquette orthogonal to the chosen
direction that meets a wall normal to another direction.
Period translations permute the open cells, while every plaquette
in this column meets that wall. Thus none of these plaquettes
is contained in a box or any of its period translates.

\begin{proposition}\label{prop:diag-coordinate-insertions}
The insertion $G_X$ satisfies \eqref{eq:diag-coordinate-bound} with
$C=\|X\|$. The insertion $G'_X$ satisfies it with
$C=2^{|S|/2}\|X\|$. A bounded bosonic integrand satisfies it with
$C=\|G\|_\infty$. These bounds survive box enlargement,
multiplication by a curvature character, and transport by a
signed permutation of the coordinates.
\end{proposition}

\begin{proof}
We give the trace calculation, including the support condition.
Choose the time direction to be the first coordinate. By the
covariance just proved, the images of the insertions have the same
operator norms and are at link planes $j+\tfrac12$. The free column
lets us choose the magnetic flux plaquette $\sigma_0$ so that no
plaquette in its column is contained in an insertion box.
At a fixed step $j$ their spatial supports and modes
are disjoint. Let $S_j$ be the union of the supports of the primed
insertions at that step.

The unfolded path variables are genuine local lifts on each
$[j,j+1]\times\operatorname{hull}(\Lambda_a)$. On electric
plaquettes their coboundary is $A_{j+1}-A_j-d_sa_j$. On magnetic
plaquettes it is $d_sA_j$; the flux correction is supported at
$\sigma_0$, which is not a plaquette of the local hull.
Two local lifts with
these curvatures and phases differ by an integer exact cochain
$\frac{2\pi}{q}d\lambda$. Gauge invariance therefore allows the
insertion to be read directly at the path variables. This is why
the free-column hypothesis is needed.

At each crossing, the factors $e^{S_c^\times(\Lambda_a)}$ remove
the crossing weights on the local modes. The remaining crossing
factors, together with their insertions, form
\[
 \operatorname{Ins}_c\left[
   \omega^{1/2}\left(\prod_a\iota(k_a(A_{j+1},A_j+d_sa_j))\right)
                   u(a_j)\omega^{1/2}\right].
\]
The even localized matrices commute because their mode sets are
disjoint. The factors $g_1^{-1}$ similarly remove the electric
Gaussian densities on the local links. The slice determinant
factorization of \S\ref{subsec:slab-coordinate-trace} now gives
the path product with these modified crossing kernels. Its
Gaussian majorant is integrable: every modified kernel has
strict decay in its electric difference, all remaining links
retain $g_1$, and the magnetic multiplier has its Gaussian
weight. Hence unfolding, periodization and the path trace are
valid for this product.

Separate the curvature character into electric factors
$e^{i\langle\beta_j,F_j^{\mathrm{el}}\rangle}$ and magnetic
unitaries $u_j$. Multiplying a regularized kernel by
$e^{i\langle\beta_j,A'-A\rangle}$ changes the corresponding
operator to
\[
 e^{-s|E-\beta_j|^2}(\alpha_{\beta_j}X)e^{-s|E-\beta_j|^2}.
\]
Fourier transformation of the disjoint link kernels therefore
gives
\begin{equation}\label{eq:diag-mixed-trace}
 Z_T\left[\varphi_\beta\prod_aG_a\right]
 =(2\pi)^{|T|}\operatorname{Tr}_{H^l}
  \prod_{j=n_1-1,\ldots,0}
  \left(K_{\beta_j,S_j}^{1/2}B_jK_{\beta_j,S_j}^{1/2}
                         M_0u_jM_0^*\right),
\end{equation}
where
\[
 B_j=\prod_{a:j_a=j}(\alpha_{\beta_j}X_a)_\Xi,\qquad
 K_{\beta,S}=r_q^v
     e^{-|E-\beta|^2/2+|E-\beta|_S^2/4}P_G.
\]
Here $r_q^v=(2\pi/q)^{|\Sigma|}$ is the gauge-volume factor in
\eqref{eq:slab-electric-operators}; the primed kernels change
the electric Gaussian variance only on $S$. Both the kernels and the local operators
commute with $P_G$. In particular
$\|B_j\|\le\prod_{a:j_a=j}\|X_a\|$ by the norm identity above.

Set $T_j=K_{\beta_j,S_j}^{1/2}M_0$. These are
Hilbert--Schmidt by the same quotient Gaussian calculation as
\eqref{eq:slab-shifted-hs}. Move the first square root cyclically
to the end of the trace. The product then has alternating factors
$B_jT_ju_jT_{j-1}^*$. Schatten H\"older, with $2n_1$ equal
exponents, bounds its absolute trace by
\[
 \prod_j\|B_j\|\prod_j\|T_j\|_{2n_1}^2.
\]
The slice trace applied to $(T_j^*T_j)^{n_1}$ gives exactly
\begin{equation}\label{eq:diag-variance-trace}
 \operatorname{Tr}(T_j^*T_j)^{n_1}
 =(2\pi)^{-|T|}Z_T\left[
    \prod_{l=0}^{n_1-1}e^{i\langle\beta_j,F_l^{\mathrm{el}}\rangle}
    \prod_{l=0}^{n_1-1}\prod_{b\in S_j}
                     \sqrt2e^{-F_l^{\mathrm{el}}(b)^2/2}\right].
\end{equation}
The second product changes $g_1$ into $g_{1/2}$ on those links.
The positive bosonic law bounds the absolute value by
$2^{|S_j|n_1/2}(2\pi)^{-|T|}Z_T[1]$. Consequently
\[
 \prod_j\|T_j\|_{2n_1}^2
 \le 2^{\frac12\sum_j|S_j|}(2\pi)^{-|T|}Z_T[1].
\]
Together with \eqref{eq:diag-mixed-trace} this proves the two
operator bounds. For insertions in another time direction use
an orientation-preserving coordinate permutation. The scalar
bound follows directly from the positive law. Multiplication by
a character merely changes $\beta$, and enlargement restricts
the allowed placements. A signed permutation transports the
torus, the free columns and the placements; conjugation of a
character still has modulus one. These observations prove the
remaining assertions.
\end{proof}

\subsection{Amplification in two cell structures}\label{subsec:diag-amplification}

A diagonal fold exchanges the first and second time directions.
Its copies of a coordinate insertion therefore cannot all be put
in one coordinate trace. We will place these copies in different
coordinate cells and apply a second chessboard estimate. Each
resulting coordinate fold contains copies of only one insertion,
all with the same time direction, so
\eqref{eq:diag-coordinate-bound} applies.

The placement requires two scales depending on the support radius
$r$, but not its center or the tilt. Choose $h''\in h\Z$ and
$k\in\N$ such that
\begin{equation}\label{eq:diag-coarse-widths}
 h''\ge2r+6h,\qquad kh\ge16h'',\qquad
 kh\in h''\Z,\qquad4h'',2kh\in\mathcal D.
\end{equation}
These choices exist: take $4h''$ to be the least common multiple
of $4h$ and an element of $\mathcal D$ exceeding $4(2r+6h)$,
then take $2kh$ to be the least common multiple of $4h''$ and an
element exceeding $32h''$. Fix choices depending only on $h$, $r$
and the side sequence. For all sufficiently large $i$, the needed
divisibility conditions then hold simultaneously.

Here is the geometry of the placement. Let $B=y+[-r,r]^3$. Choose
coarse diagonal offsets $c'_\pm\in c_\pm+h\Z$ (with
$c_-=c_0$) and $c'_3\in c_3+h\Z$ so that
\begin{equation}\label{eq:diag-coarse-center}
 \ell_-(y)-c'_-\in[2h'',2h''+h),\quad
 \ell_+(y)-c'_+\in[4h'',4h''+h),\quad
 y_3-c'_3\in[h'',h''+h).
\end{equation}
The box lies strictly in the coarse cell $Q'_0$ of width $kh$.
Every element of its wall group has the form
\[
 g=\tau_{kh(z_-e_-+z_+e_++2z_3e_3)}
            \theta_-^{a_-}\theta_+^{a_+}\theta_3^{a_3},
 \qquad z\in\Z^3,\quad a\in\{0,1\}^3,
\]
where the reflections are at the lower walls of $Q'_0$. From
\eqref{eq:diag-coarse-center} one obtains
\begin{equation}\label{eq:diag-coordinate-cell-index}
 gy=y+h''\kappa(g)+s(a),\qquad |s(a)|_\infty<2h,
\end{equation}
with
\[
 \kappa(g)=(-2a_--4a_+,\,2a_--4a_+,\,-2a_3)
       +\frac{kh}{h''}(z_-+z_+,\,z_+-z_-,\,2z_3).
\]
The eight first vectors are distinct, and their pairwise
differences have size at most six. Since $kh/h''\ge16$, equality
of two indices first forces equality of their $a$'s and then of
their $z$'s. Thus $g\mapsto\kappa(g)$ is injective.

Choose half-integer coordinate offsets $o_\mu$ in
\[
 I_\mu=(y_\mu+r+2h-h'',\ y_\mu-r-2h).
\]
These intervals have length at least $2h$. Equation
\eqref{eq:diag-coordinate-cell-index} places every $gB$ strictly
inside coordinate cell $\kappa(g)$ of width $h''$. We can also
choose the offsets so that all these coordinate walls avoid
every reflected copy of $\supp O$. Under the original diagonal
wall group, the centers of the copies of a box containing $O$
have at most four residue classes modulo $h$ in either of the
first two coordinates, and at most two in the third. Its radius
$r_O$ therefore forbids at most $4(2r_O+1)=8r_O+4$ half-integer
classes. Since $h>8r_O+8$, an allowed class exists. Each $I_\mu$
contains a representative, and $h''\in h\Z$ keeps every wall
in the allowed class.

For large $i$, the adapted torus $D_i$ carries the original and
coarse diagonal structures; its period lattice lies in the
translation subgroup of the coarse wall group. The coordinate
structure of width $h''$ cuts $D_i$ and is carried by
$T_i=C(n_i/2,2n_i,n_i)$. Formula
\eqref{eq:diag-coordinate-cell-index} also shows that coarse
diagonal cells distinct modulo the periods have their boxes in
coordinate cells distinct modulo those periods. Indeed a period
translation changes $\kappa$ by that translation divided by
$h''$; injectivity then gives the converse. Thus the counts are
\begin{equation}\label{eq:diag-coarse-cell-counts}
 N'_i=\frac{n_i^3}{|Q|k^3},\qquad
 N''_i=\frac{n_i^3}{(h'')^3}.
\end{equation}
The first is the number of coarse diagonal cells in $D_i$; the
second is the number of coordinate cells in either torus. The
shortest periods are at least $n_i/2\ge kh\ge16h''$, so the
boxes, with the margins used for local lifts, embed. Folding
one such box in the coordinate structure gives a disjoint
placement with free columns, since every image stays inside
its open cell.

\begin{theorem}\label{thm:diag-amplification}
Put $\Delta_i=\Delta(n_i)$, with $\Delta$ defined in
\eqref{eq:slab-pressure-error}. It tends to zero as $i\to\infty$,
with the charge, mass and weights fixed.
If $G$ satisfies the coordinate bound $C$ in a box of radius $r$,
then, for $k$ as above and all sufficiently large $i$,
\begin{equation}\label{eq:diag-amplification}
 |\mu_i^t(G)|\le C\gamma_{t,i}(k),\qquad
 \gamma_{t,i}(k)=
      \exp(2t|O|_1k^3+|Q|k^3\Delta_i),\qquad t\ge0.
\end{equation}
The threshold and $k$ depend only on the side sequence, $h$ and
$r$. In particular they are uniform in the location of $G$.
\end{theorem}

\begin{proof}
Each coarse diagonal cell contains $k^3$ original cells. Put
$P=\exp(t\sum_{Q_\kappa\subset Q'_0}O^{[\kappa]})$. The tilt
$e^{tV_i}$ is exactly the product of the reflected copies of $P$
over the coarse cells: their walls are original walls, so each
reflection carries the original array to the same array. Apply
the adapted chessboard estimate to the integrands $PG$ in one
coarse cell and $P$ in the others. The fold of $P$ is $e^{tV_i}$.
After dividing its powers out, this gives
\begin{equation}\label{eq:diag-amplified-power}
 |\mu_i^t(G)|^{N'_i}
 \le\frac{Z_{D_i}[e^{tV_i}\prod_{a=1}^{N'_i}G^{g_a}]}
              {Z_{D_i}[e^{tV_i}]}.
\end{equation}
The right side is nonnegative by the reflected fold form. This
step uses reflection positivity of even integrands and does not
require $G$ itself to be a positive scalar function.

Expand $e^{tV_i}$ in curvature characters. The coefficient sum
is at most $\exp(tN_i|O|_1)$: the character algebra is closed
under multiplication, and the exponential series is absolutely
summable in its coefficient norm. Every plaquette in a character
of this expansion lies strictly within a coordinate cell,
because the chosen walls avoid the support of every copy of $O$.
For each character term, collect its factors by coordinate cells.
Each cell then contains a character alone or a character times
one $G^{g_a}$; precisely $N'_i$ cells have the latter form. Period
translations can be used to place the integrands in a chosen
coordinate brick without changing their value on $D_i$.

Use the chessboard estimate from the adapted to the coordinate
frame, \eqref{eq:slab-chessboard-coordinate}. A cell with a
character alone has fold bounded by $Z_{T_i}[1]$. A cell with
a copy of $G$ has a fold consisting of a character and its
coordinate images; the geometry above supplies the free columns
and disjoint embedding, so its fold is bounded by
$C^{N''_i}Z_{T_i}[1]$. Taking the chessboard roots and multiplying
the cell contributions gives, for every character term,
\[
 \left|Z_{D_i}\left[\varphi_\beta
                    \prod_{a=1}^{N'_i}G^{g_a}\right]\right|
                \le C^{N'_i}Z_{T_i}[1].
\]
The coefficient sum therefore bounds the numerator of
\eqref{eq:diag-amplified-power} by
$e^{tN_i|O|_1}C^{N'_i}Z_{T_i}[1]$. The denominator bound
\eqref{eq:diag-tilt-denominator} contributes a second
$e^{tN_i|O|_1}$. Finally, the pressure comparison gives
$Z_{T_i}[1]/Z_{D_i}[1]\le e^{n_i^3\Delta_i}$.
The root of order $N'_i$ gives
\eqref{eq:diag-amplification}.
\end{proof}

We also need exponential moments, for which a coordinate operator
bound is unnecessary. Repeat \eqref{eq:diag-amplified-power} with
$G=e^{\langle\beta,F\rangle}$ supported in $B$. It is an admissible
local integrand, since
$\langle\beta,F\rangle\le\|\beta\|^2+
\frac14\sum_{\supp\beta}F_\sigma^2$. Its diagonal copies have
disjoint supports. Remove the tilt from the numerator by its
pointwise bound, and apply the second chessboard estimate to
their coordinate cell integrands. The fold of one exponential
has $N''_i$ disjoint copies of $\beta$, and hence squared norm
$N''_i\|\beta\|^2$. Gaussian domination on $T_i$ bounds its
chessboard root by
$e^{\|\beta\|^2/2}Z_{T_i}[1]^{1/N''_i}$. There are $N'_i$
nonempty cells. The same denominator and pressure bounds
therefore yield
\begin{equation}\label{eq:diag-tilted-domination}
 \mu_i^t(e^{\langle\beta,F\rangle})
      \le e^{\|\beta\|^2/2}\gamma_{t,i}(k).
\end{equation}
For each fixed $t$ this proves tightness on $\Omega$. Indeed,
Markov's inequality applied to $\beta=\pm s1_\sigma$ gives a
uniform coordinate tail $2e^{-sR+s^2/2}\gamma_{t,i}(k(h,1))$.
Choose radii $R_\sigma$ with these tails summable over plaquettes;
the intersection of $\Omega$ with their product of compact
intervals is compact. The finitely many initial laws are tight
individually. Any weak limit $\nu_t$ consequently satisfies
\begin{equation}\label{eq:diag-limit-domination}
 \nu_t(e^{\langle\beta,F\rangle})
      \le e^{\|\beta\|^2/2+2t|O|_1k^3}.
\end{equation}
This uses lower semicontinuity of the exponential integral,
as explained after \eqref{eq:diag-tilted-positive-form}.

\subsection{The two adjacent fermion layers}\label{subsec:diag-two-layers}

The remaining comparison argument uses gauge invariant polynomials
\[
 P(\widehat A)=\sum_a b_a e^{i\langle\delta_a,\widehat A\rangle}
                                     \Psi_a,\qquad
 d^*\delta_a=q\mathfrak q(\Psi_a),
\]
where the cochains have finite support and the $\Psi_a$ are ordered
Grassmann monomials. Denote their class by $\mathcal P$. The charge
condition is exactly gauge invariance. It also makes each term
even: the total divergence is zero, so its numbers of $\psi$ and
$\bar\psi$ factors agree. The expansion is unique after equal
pairs $(\delta_a,\Psi_a)$ have been collected, because the monomials
and the distinct characters are linearly independent. Each such
polynomial is an admissible local integrand with bounded coefficients.

At a coordinate link plane $c=j+\tfrac12$, a two-layer polynomial
reads photon links in the planes $j$ and $j+1$ and between them,
and fermion fields only at the two ends of the bonds crossing
$c$. These are the thin layers $\mathcal X_{j,6}$ and
$\mathcal X_{j+1,1}$; the other thin layers of the two chambers do
not occur. Call this subclass $\mathcal P_1(c)$. For a monomial
$Y=e^{i\delta\widehat A}\Psi$, write its exponent as
$(\delta_0,\delta_1,\delta_\perp)$ on the lower spatial, upper
spatial and temporal links. If $\mathfrak q_\pm$ are the lower
and upper charge densities, gauge invariance gives
\begin{equation}\label{eq:diag-two-layer-gauss}
 d_s^*\delta_0-\delta_\perp=q\mathfrak q_-,\qquad
 d_s^*\delta_1+\delta_\perp=q\mathfrak q_+.
\end{equation}
Let $S=\supp\delta_0\cup\supp\delta_1$ and choose a nonempty $\Lambda$
containing their endpoints, the support of $\delta_\perp$ and
all sites over which the fields of $\Psi$ occur.

Every such $Y$ is a kernel insertion, but the operator between
the usual quarter-heat factors need not be bounded. We give the
different regularization that permits approximation by bounded
operators. The crossing modes and their reflected modes form a
basis of the two-layer generators. Since all their weights are
positive when $0<\rho<1$, there is a unique finite matrix
$\Upsilon_Y$ with
\begin{equation}\label{eq:diag-two-layer-matrix}
 \left.\Psi e^{-S_c^\times(\Lambda)}\right|_{\widehat a_j=0}
       =\operatorname{Ins}_c[
                      \omega^{1/2}\Upsilon_Y\omega^{1/2}].
\end{equation}
To obtain it, expand the left side in the basis
$\Theta_c\zeta_J\zeta_I$ and multiply the resulting matrix by
$\omega^{-1/2}$ on both sides. Its nonzero entries have
$\mathfrak q(I)-\mathfrak q(J)=\mathfrak q_-+\mathfrak q_+$.
Thus $\Upsilon_Y$ has charge
$q^{-1}d_s^*(\delta_0+\delta_1)$ and is even.

Set
\begin{equation}\label{eq:diag-two-layer-kernel}
 k_Y^0(A',A)=e^{i\delta_1A'}
             \prod_{b\in S}g_1(A'(b)-A(b))e^{i\delta_0A}
                         \Upsilon_Y.
\end{equation}
It is the continuous kernel of
$e^{i(\delta_0+\delta_1)A}e^{-|E_S+\delta_0|^2/2}\Upsilon_Y$.
We also write $k_Y^0$ for this integral operator.
The charge identity and unchanged Gaussian difference prove
its covariance. Moreover
\begin{equation}\label{eq:diag-two-layer-is-insertion}
 Y=G[k_Y^0].
\end{equation}
Indeed, the Gaussian factors in the insertion cancel. Applying
the automorphism that multiplies the lower fields by their
charge phases to \eqref{eq:diag-two-layer-matrix} replaces its
matrix by $\omega^{1/2}\Upsilon_Yu(\widehat a_j)\omega^{1/2}$
and its left side by
$e^{-iq\langle\mathfrak q_-,\widehat a_j\rangle}
 \Psi e^{-S_c^\times(\Lambda)}$. The remaining temporal exponent
in $G[k_Y^0]$ is
$d_s^*\delta_0-q\mathfrak q_-=\delta_\perp$ by
\eqref{eq:diag-two-layer-gauss}, proving the identity.

Put $\widetilde E=E_S+\delta_0$, $D=\delta_1-\delta_0$ and
$a=(|\delta_0|^2+|\delta_1|^2)/8$. Define the bounded operator
\begin{equation}\label{eq:diag-two-layer-prime}
 X'_Y=e^{i(\delta_0+\delta_1)A}
       e^{-\frac14|\widetilde E|^2+
                         \frac14\langle\widetilde E,D\rangle+a}
                                                      \Upsilon_Y.
\end{equation}
It belongs to $\mathfrak W(S,\Lambda)$, and completion of the
quadratic exponent gives
\[
 e^{-|E_S|^2/8}X'_Ye^{-|E_S|^2/8}=k_Y^0,
 \qquad
 \|X'_Y\|=\|\Upsilon_Y\|e^{|D|^2/16+a}.
\]
For the first identity move the left heat factor through the
phase using the Weyl rule. Its resulting squares are
$\frac18|\widetilde E+\delta_1|^2+
 \frac18|\widetilde E-\delta_0|^2$; their linear and constant
terms cancel those in \eqref{eq:diag-two-layer-prime}. Therefore
$Y=G'_{X'_Y}$ and it satisfies the coordinate bound
$2^{|S|/2}\|X'_Y\|$. This bound will also control the approximation
error uniformly in the torus.

For $\epsilon>0$ let $P_\epsilon=e^{-\epsilon|E_S|^2}$ and
\begin{equation}\label{eq:diag-two-layer-epsilon}
 X_{Y,\epsilon}=e^{i(\delta_0+\delta_1)A}
  \exp\left(-\epsilon|\widetilde E+\delta_1|^2
       -\epsilon|\widetilde E-\delta_0|^2
       +\tfrac12\langle\widetilde E,D\rangle+2a\right)\Upsilon_Y.
\end{equation}
The quadratic part is strictly negative, so this too is in
$\mathfrak W(S,\Lambda)$. The same Weyl calculation gives
\begin{align}\label{eq:diag-two-layer-approximation}
 e^{-|E_S|^2/4}X_{Y,\epsilon}e^{-|E_S|^2/4}
                      &=P_\epsilon k_Y^0P_\epsilon,\notag\\
 Y-G_{X_{Y,\epsilon}}&=
               G'_{X'_Y-P_\epsilon X'_YP_\epsilon}.
\end{align}
The second equality follows from the first and the primed
representation, since their continuous kernels agree. Its
coordinate bound is $2^{|S|/2}\xi_Y(\epsilon)$, where
\begin{equation}\label{eq:diag-two-layer-error}
 \begin{split}
 \xi_Y(\epsilon)={}&\|\Upsilon_Y\|
       \sup_{x\in\R^S}e^{-|x|^2/4+\langle x,D\rangle/4+a}\\
 &\quad\cdot(1-e^{-\epsilon|x+\delta_1|^2
                         -\epsilon|x-\delta_0|^2})
       \le C_Y\epsilon.
 \end{split}
\end{equation}
The last inequality uses $1-e^{-u}\le u$ and the boundedness
of a Gaussian times a quadratic polynomial. By contrast,
\begin{equation}\label{eq:diag-two-layer-operator-size}
 \|X_{Y,\epsilon}\|
 =\|\Upsilon_Y\|
   \exp\left(\frac{(1/2-2\epsilon)^2|D|^2}{8\epsilon}
                  +(1/4-\epsilon)(|\delta_0|^2+|\delta_1|^2)\right).
\end{equation}
It can grow exponentially in $1/\epsilon$. This forces the
order of the limits in the next result.

Reflection changes $(\delta_0,\delta_1,\delta_\perp)$ to
$(-\delta_1,-\delta_0,\delta_\perp)$ and
$\Upsilon_Y$ to $\Upsilon_Y^*$. The latter follows by reflecting
\eqref{eq:diag-two-layer-matrix}; $\operatorname{Ins}_c[\omega]$
is even and invariant. Taking adjoints in the two heat identities
then gives
$X'_{\Theta_cY}=(X'_Y)^*$ and
$X_{\Theta_cY,\epsilon}=X_{Y,\epsilon}^*$; the heat operator
is injective with dense range, so these identities determine
the bounded middle operators uniquely. All constructions extend
linearly to two-layer polynomials. Coordinate permutations
give the same assertions at every coordinate link plane.

\begin{proposition}\label{prop:diag-two-layer-principle}
For every $Y\in\mathcal P_1(c)$ at a coordinate link plane,
\begin{equation}\label{eq:diag-two-layer-principle}
 \lim_{t\downarrow0}\limsup_{i\to\infty}
    |\mu_i^t(\Theta_cY)-\overline{\mu_i^t(Y)}|=0.
\end{equation}
The assertion concerns the finite-volume functionals on fermion
polynomials; it does not presume that they converge to fermion
laws.
\end{proposition}

\begin{proof}
For a fixed support, $o_i^t(X)=\mu_i^t(G_X)$ is a unital linear
functional on the local operator algebra with
$|o_i^t(X)|\le\gamma\|X\|$, where
$\gamma=\gamma_{t,i}(k)$. A unital functional this close in norm
to one is nearly Hermitian. If $X_0=X_0^*$ and $\|X_0\|\le1$,
write $o_i^t(X_0)=u+iv$. For each real $s$,
\[
 (v+s)^2\le |o_i^t(X_0+is)|^2
            \le\gamma^2\|X_0+is\|^2
            \le\gamma^2(1+s^2).
\]
For $\gamma>1$ take $s=v/(\gamma^2-1)$; this gives
$|v|\le\sqrt{\gamma^2-1}$. For $\gamma=1$ the same conclusion
follows by letting $s$ tend to either infinity. Decomposing a
general $X$ into its real and imaginary self-adjoint parts yields
\begin{equation}\label{eq:diag-near-hermitian}
 |o_i^t(X^*)-\overline{o_i^t(X)}|
           \le4\sqrt{\gamma^2-1}\|X\|.
\end{equation}
Use \eqref{eq:diag-insertion-reflection} to interpret this as
reflection covariance of the insertion expectation.

Apply \eqref{eq:diag-near-hermitian} to $X_{Y,\epsilon}$ and
use \eqref{eq:diag-two-layer-approximation} for $Y$ and its
reflection. The amplification theorem controls the two errors
and gives
\begin{equation}\label{eq:diag-two-layer-finite-error}
 \begin{split}
 |\mu_i^t(\Theta_cY)-\overline{\mu_i^t(Y)}|
 \le{}&4\sqrt{\gamma^2-1}\|X_{Y,\epsilon}\|\\
      &+2\cdot2^{|S|/2}\gamma\xi_Y(\epsilon).
 \end{split}
\end{equation}
The support box and its coarse scales do not depend on
$\epsilon$. First take $i\to\infty$, and then $t\downarrow0$
with $\epsilon$ fixed. Since $\Delta_i\to0$, the first term
vanishes and the remaining upper bound is
$2\cdot2^{|S|/2}\xi_Y(\epsilon)$. Let $\epsilon\downarrow0$
using \eqref{eq:diag-two-layer-error}. Finite sums and coordinate
transport give the stated result.
\end{proof}

\subsection{Removing fields from a thin layer}\label{subsec:diag-fermion-recursion}

The fermion graph has six thin layers in each thick layer $j$,
at heights $j-v_1,j-v_2,j-v_3,j+v_3,j+v_2,j+v_1$.
Each site has one neighbor in the next thin layer, one in the
preceding one and two in its own layer. The forward matching
bonds have types $2,1,0,1,2,3$ in succession. This follows
directly from the chamber reflections: type zero reverses
the coordinate of modulus $v_3$, types one and two interchange
the adjacent moduli, and type three reverses the coordinate of
modulus $v_1$ and shifts its nearest photon site by the corresponding
unit vector. The bonds of types zero, one and two carry no photon
phase. Type three bonds within thin layers $2,\ldots,5$ carry
spatial phases in plane $j$, while those between layer six and
the next layer one carry temporal phases between $j$ and $j+1$.
These facts control which planes the recursion can read.

For a photon site $y$, let $\mathcal P^\pm(y)$ consist of the
same exponential monomials as $\mathcal P$, with
$d^*\delta=q(\mathfrak q(\Psi)\mp1_y)$. They transform with
gauge phase $e^{\pm iq\lambda(y)}$ and are odd. Removing a
$\bar\psi(x)$ from an invariant monomial leaves an element of
$\mathcal P^+(\pi(x))$; removing a $\psi(x)$ leaves one in
$\mathcal P^-(\pi(x))$. Multiplication by a hopping phase transports
the specified charge to the other endpoint. These assertions
are the identities
\[
 d^*1_{(y,y')}=1_{y'}-1_y,\qquad
 U_A(x,x')=e^{iq\langle1_{(\pi(x),\pi(x'))},A\rangle}.
\]
They show that every term of the following replacements stays
in $\mathcal P$.

Berezin integration by parts, with left derivatives, gives
\begin{align}\label{eq:diag-equations-of-motion}
 Z_T[\Phi(\bar\psi \mathcal D_A)_\alpha(x')G]
       &=-Z_T[\Phi\partial_{\psi_\alpha(x')}G],\notag\\
 Z_T[\Phi(\mathcal D_A\psi)_\alpha(x')G]
       &= Z_T[\Phi\partial_{\bar\psi_\alpha(x')}G].
\end{align}
The appropriate charge of $G$ makes each integrand invariant.
These identities first hold at fixed complex $A$: the
derivatives of $e^{-S_T(A)}$ are respectively
$(\bar\psi \mathcal D_A)_\alpha e^{-S_T(A)}$ and
$-(\mathcal D_A\psi)_\alpha e^{-S_T(A)}$, and the Berezin integral
annihilates a derivative. Integrating against any bounded
bosonic factor $\Phi(F)$ then proves
\eqref{eq:diag-equations-of-motion}. The enlarged local box
must embed so that all generators involved are distinct torus
generators and the phases can be read on the local lift.

Let $x'$ be the neighbor of $x$ in the next thin layer. Its
matching hopping matrix $w(x,x')$ has eigenvalues
$\frac{w_i}{2}(\rho\mp1)$ and is invertible for our fixed
$0<\rho<1$. Solve the equation at $x'$ for the field at $x$.
With the replaced field moved to the left of the monomial,
the two rules are
\begin{equation}\label{eq:diag-field-replacement}
 \begin{split}
 \bar\psi(x)\longmapsto{}&
 \left(M\bar\psi(x')-
   \sum_{x''\sim x',\,x''\ne x}\bar\psi(x'')w(x'',x')U_A(x'',x')
                  +\partial_{\psi(x')}\right)
                     w(x,x')^{-1}U_A(x',x),\\
 \psi(x)\longmapsto{}&U_A(x,x')w(x',x)^{-1}
 \left(M\psi(x')-
   \sum_{x''\sim x',\,x''\ne x}w(x',x'')U_A(x',x'')\psi(x'')
                  -\partial_{\bar\psi(x')}\right).
 \end{split}
\end{equation}
The derivatives act on the monomial to their right. Multiplying
by the inverse matching matrix and its inverse phase in
\eqref{eq:diag-equations-of-motion} proves equality of the
expectations before and after either replacement, for every
bosonic factor $\Phi$. This equality is independent of a choice
of time direction or a symmetry of the torus.

For a thin layer $\mathbb L$ define a linear map $R^+_{\mathbb L}$
by recursion on the number of its fields in an ordered monomial.
Select the first such field, move it to the front with its
Grassmann sign, and apply \eqref{eq:diag-field-replacement}
through its matching neighbor in $\mathbb L^+$. That neighbor
has only one neighbor in $\mathbb L$, the field just removed.
All other fields created by the rule lie in $\mathbb L^+$ or
$\mathbb L^{++}$; its derivatives remove fields in
$\mathbb L^+$. Thus the number in $\mathbb L$ falls by exactly
one at every recursive call. The map is well-defined, preserves
$\mathcal P$, removes all fields in $\mathbb L$, and creates
no new fields outside its next two layers. A downward map
$R^-_{\mathbb L}$ is obtained using the preceding matching
neighbor.

If the initial polynomial is read in a box $B$, the entire
one-layer recursion is read in $B+[-1,1]^3$. The fields being
removed stay in the original box, and their next neighbors
and the other neighbors of those sites have photon projections
equal or nearest to the original projection: among two successive
different bonds at most one has type three. Hence these sites
and every added phase are in the same enlarged box, irrespective
of the number of recursive calls. Induction therefore gives
\begin{equation}\label{eq:diag-thin-layer-identity}
 Z_T[\Phi P]=Z_T[\Phi R^\pm_{\mathbb L}P]
\end{equation}
whenever this enlarged box embeds. The new photon links are
only those of type three bonds with an endpoint in
$\mathbb L^\pm$. The recursion commutes with multiplication
by a curvature character, since it differentiates only fields
and its choices depend only on the monomial.

We first use this recursion to bound an arbitrary polynomial;
here its photon links need not yet be reduced. Move its fields
upwards through all thin layers below $\mathcal X_{j,6}$ and
downwards through those above $\mathcal X_{j+1,1}$. The last
downward replacement can reach $\mathcal X_{j,6}$, and cannot
reach a lower layer. We obtain, on every sufficiently large
torus and with every bounded bosonic factor, an equality
\begin{equation}\label{eq:diag-polynomial-normal-form}
 Z_T[\Phi P]=\sum_a b_a Z_T[\Phi\varphi_{\beta_a}Y_a],
 \qquad Y_a\in\mathcal P_1(c).
\end{equation}
Here the $Y_a$ can be chosen with integer exponents $qj_a$.
Indeed each remaining monomial has integer charge of total
zero on the two adjacent photon planes. In their connected
rectangular support graph choose paths from a fixed vertex
to each charged vertex; the charge-weighted sum of these
paths is an integer cochain $j_a$ with the required divergence.
Then $u_a=\delta_a-qj_a$ is divergence free in the same
integer box and can be written $d^*\beta_a$ there.

The filling must stay inside the support box for the polynomial
bound below. In
$\prod_\mu[a_\mu,a'_\mu]$, set electric components
$\beta_j^{\mathrm{el}}=-\sum_{l=a_1}^j(u_a)_l$ for
$a_1\le j<a'_1$. Subtracting $d^*\beta$ removes all spatial
links below the top plane; divergence zero then removes
the temporal links. On that top plane a second partial-sum
sweep in the second coordinate removes the third-direction
links below its top edge. Divergence zero removes the
second-direction links and leaves only a terminal segment
in direction three. Its divergence equation and vanishing
boundary value make it zero. The two sweeps use only
plaquettes of the box. This proves the asserted filling
and thus the splitting $\delta_a=d^*\beta_a+qj_a$ in
\eqref{eq:diag-polynomial-normal-form}.

Apply the coordinate bound to each $Y_a$, absorb the extra
character and use amplification. For any fixed box $B_0$
there are constants $C_P$, $k_P$ and a volume threshold with
\begin{equation}\label{eq:diag-polynomial-apriori}
 \sup_{\supp\beta\subset B_0}
       |\mu_i^t(\varphi_\beta P)|
             \le C_P\gamma_{t,i}(k_P),\qquad t\ge0.
\end{equation}
One may take
$C_P=\sum_a|b_a|2^{|S_a|/2}\|X'_{Y_a}\|$, with the summand
one for $Y_a=1$. All support boxes in the finite normal form
and $B_0$ fit in one fixed box; its radius determines $k_P$.
The size of $\beta$ does not enter the bound. Reflected
polynomials have the same kind of bound, since the normal
form identity can be reflected on the reflected torus and
the coordinate bounds are preserved by reflection.

\subsection{The contour shift and an outer plane}\label{subsec:diag-contour}

A fermion recursion removes fields from a plane but leaves its
photon exponent. An imaginary shift of links in the next plane
removes that exponent together with the outgoing temporal one.
The all-flux measure is essential for this shift: summation over
exact integer fluxes unfolds the compact link intervals to
real lines on which the contour can move.

Let $u$ be a real cochain supported on finitely many links $S$
of one coordinate plane, and let $P\in\mathcal P$. Choose a
box $B$ reading $P$ and containing $b+[-1,1]^3$ for every
$b\in S$. Put $P^{(u)}(\widehat A)=P(\widehat A-iu)$ by
its entire exponential expansion. Define
\begin{equation}\label{eq:diag-shift-fermion-polynomial}
 \begin{split}
 Q_u=\sum_{b\in S}\sum_{(x,x')\text{ over }b}\big[&
 (e^{qu_b}-1)\bar\psi(x)w(x,x')U_A(x,x')\psi(x')\\
 &+(e^{-qu_b}-1)\bar\psi(x')w(x',x)U_A(x',x)\psi(x)\big],
 \qquad V_u=e^{Q_u}.
 \end{split}
\end{equation}
The bonds are oriented with $b$. Each term is invariant and
even, and the exponential is a finite polynomial by nilpotence.
It reads only those links of $S$ with $u_b\ne0$ and their
overlying bonds. For a plane normal to $e_1$, those fields
are in thin layers $2,\ldots,5$ of its thick layer. The change
in the action is $S_T(A-iu)-S_T(A)=-Q_u$, so
$e^{-S_T(A-iu)}=e^{-S_T(A)}V_u$.

If $\Phi(F)$ is entire and bounded on each tube of bounded
imaginary curvature, then on every torus in which $B$ embeds
we have
\begin{equation}\label{eq:diag-contour-identity}
 Z_T[\Phi P]=e^{\|du\|^2/2}
       Z_T[\Phi(F-idu)\varphi_{du}(F)V_uP^{(u)}].
\end{equation}
The norm here is the lattice norm, equal to the torus norm
because its support embeds. We supply the contour justification.
For each $b$ choose its electric plaquette leaving the plane
in the normal direction. Its only side in $S$ is $b$, with
coefficient $\pm1$. Consequently $d$ is injective on cochains
supported in $S$, and the action
$z\mapsto z+dk$, $k\in\Z^S$, on closed integer fluxes is free.
Choose orbit representatives $z_0$. The compact $A_S$ intervals
and the sum over $k$ become the integral over $\R^S$, since
$(A,z_0+dk)$ represents the same configuration as
$(A+\frac{2\pi}{q}k,z_0)$.

For fixed $z_0$ and the other real links, choose an integer
primitive of $z_0$ in the local box. Adding it to $A$ gives
a local lift affine in the complex variables $\zeta=A_S$.
The Berezin coefficients of $e^{-S_T(A)}$ are polynomials
in the link phases, so the full integrand extends to an
entire function $f(\zeta)$. On the strip
$|\Im\zeta_b|\le\|u\|_\infty$ the selected plaquettes give
\begin{equation}\label{eq:diag-contour-strip}
 |f(\zeta)|\le C
          \prod_{b\in S}e^{-(\Re\zeta_b-a_b)^2/2}.
\end{equation}
Indeed each of their real curvatures is $\pm\Re\zeta_b+a'_b$,
all imaginary curvatures are bounded in this finite strip,
the factor $\Phi$ is bounded there, and each remaining
exponential phase is bounded there. The constant can be
chosen independent of the remaining real links and the
orbit representative; only the real Gaussian centers vary.

Move one variable at a time along the rectangle with real
endpoints $\pm R$ and imaginary displacement $-iu_b$. The
vertical sides tend to zero by \eqref{eq:diag-contour-strip}.
Fubini then gives $\int f(\xi)\,\dd\xi=
\int f(\xi-iu)\,\dd\xi$. At the shifted arguments the
Maxwell exponent becomes
\[
 -\tfrac12\langle F-idu,F-idu\rangle
       =-\tfrac12\|F\|^2+i\langle du,F\rangle+\tfrac12\|du\|^2.
\]
The change in the fermion action is the polynomial $V_u$,
and the observable becomes $P^{(u)}$. Both the initial
integrand and the resulting real integrand are absolutely
integrable by their local bounds. Hence we may sum the
identities over $z_0$ and integrate the remaining links,
which proves \eqref{eq:diag-contour-identity}.

We now eliminate the lowest plane $j$ read by $P$. This means
that its fields have photon projection height at least $j$
and its links have no endpoint below plane $j$. Apply the
six upward thin-layer maps of this thick layer to obtain
$P^\uparrow=\sum_a b_a e^{i\delta'_a\widehat A}\Psi_a$.
It reads no fields of thick layer $j$ and no links below it;
the only new links are spatial links in plane $j$ and temporal
links from $j$ to $j+1$. At each term, gauge invariance in
plane $j$ therefore says
\begin{equation}\label{eq:diag-lowest-plane-charge}
 (\delta'_a)^\perp_j=d_s^*(\delta'_a)_j.
\end{equation}
For this term let $u_a$ be supported on plane $j+1$ with
$(u_a)_{j+1}=(\delta'_a)_j$. Define
\begin{equation}\label{eq:diag-outer-polynomial}
 P_a=b_a e^{\|du_a\|^2/2+\langle\delta'_a,u_a\rangle}
       e^{i\langle\delta'_a+d^*du_a,\widehat A\rangle}
                                    V_{u_a}\Psi_a.
\end{equation}
Then \eqref{eq:diag-thin-layer-identity} and
\eqref{eq:diag-contour-identity} give the exact relation
\begin{equation}\label{eq:diag-outer-identity}
 Z_T[\Phi P]=\sum_a Z_T[\Phi(F-idu_a)P_a].
\end{equation}
The construction is on the infinite lattice; the same relation
holds for every $\Phi$ entire in the plaquette variables and
bounded on each tube $|\Im F_\sigma|\le r$ for all $\sigma$,
$r>0$, and every torus containing the fixed enlarged support box.

Here is the cancellation of the eliminated exponent. Write
$v=(\delta'_a)_j$. From the slice formulas
$(du)^{\mathrm{el}}_l=u_{l+1}-u_l-d_su_l^\perp$ and
$(d^*\beta)_l=\beta^{\mathrm{el}}_{l-1}
-\beta^{\mathrm{el}}_l+d_s^*\beta^{\mathrm{mag}}_l$, the only
nonzero components of $d^*du_a$ are
\begin{align}\label{eq:diag-outer-cancellation}
 (d^*du_a)_j&=-v,&(d^*du_a)^\perp_j&=-d_s^*v,\notag\\
 (d^*du_a)_{j+1}&=(2+d_s^*d_s)v,
       &(d^*du_a)^\perp_{j+1}&=d_s^*v,\notag\\
 (d^*du_a)_{j+2}&=-v.
\end{align}
Thus the spatial exponent in plane $j$ vanishes, and
\eqref{eq:diag-lowest-plane-charge} cancels its outgoing
temporal exponent as well. The polynomial $V_{u_a}$ reads
only thick layer $j+1$. Hence $P_a$ has lowest plane at
least $j+1$ and highest plane at most the maximum of $j+2$
and the previous highest plane. The six field recursions
enlarge the box by six, and the shifted links with their
plaquettes require two more. The whole step is read in
$B+[-8,8]^3$. Elimination from the highest plane is obtained
by the downward maps and a shift in plane $j-1$. The corresponding
incoming temporal exponent is
$-d_s^*(\delta'_a)_j$, and the same two components cancel.
The highest plane falls by one, while the lowest can fall
only to $j-2$.

The source of error in the tilted law is precisely the shifted
bosonic factor. For a finite cochain $\eta$, only finitely
many copies in the array $V_i$ meet its support. Their change
is the fixed trigonometric polynomial
\begin{equation}\label{eq:diag-array-shift}
 H_\eta(F)=\sum_{\kappa,b}a_b
       (e^{\langle(g_\kappa^{-1})^*\beta_b,\eta\rangle}-1)
                        \varphi_{(g_\kappa^{-1})^*\beta_b}(F),
\end{equation}
where $O=\sum_ba_b\varphi_{\beta_b}$ and only contributing
cells occur. At most $|\supp\eta|$ cells contribute; each
lies in the support box of $\eta$ enlarged by $h$. Its
coefficient norm has the bound
\begin{equation}\label{eq:diag-array-shift-size}
 |H_\eta|_1\le C(\eta):=|\supp\eta|\,|O|_1
        (e^{\max_b\|\beta_b\|_1\|\eta\|_\infty}-1).
\end{equation}
If that enlarged box embeds, no two contributing cells are
period equivalent and the same polynomial is the torus
difference $V_i(F-i\eta)-V_i(F)$. Expand its exponential
in characters:
\begin{equation}\label{eq:diag-shifted-tilt}
 e^{tV_i(F-i\eta)}=e^{tV_i(F)}
        \left(1+\sum_\alpha b_\alpha(t)\varphi_{\beta_\alpha}(F)\right),
 \qquad\sum_\alpha|b_\alpha(t)|\le e^{tC(\eta)}-1.
\end{equation}
The coefficient norm is submultiplicative, so the exponential
series proves absolute uniform convergence. All character
supports lie in one fixed enlarged box, independent of $i$
and $t$; sums of character exponents do not enlarge this support.

Set $P^\flat=\sum_aP_a$ in \eqref{eq:diag-outer-identity}.
Insert \eqref{eq:diag-shifted-tilt} with $\eta=du_a$ and divide
by the positive denominator. Termwise integration is valid:
at a fixed torus the functional on bounded bosonic factors
multiplying $P_a$ is continuous for their supremum norm, by
the absolute integrability of its Berezin coefficients.
We get
\begin{equation}\label{eq:diag-outer-error}
 |\mu_i^t(P)-\mu_i^t(P^\flat)|
 \le(e^{Ct}-1)\sum_a
       \sup_{\supp\beta\subset B'}|\mu_i^t(\varphi_\beta P_a)|,
\end{equation}
where $C=\max_aC(du_a)$ and $B'$ is a fixed box enlarged
by $h$. The a priori bound
\eqref{eq:diag-polynomial-apriori} makes this $O(t)$ uniformly
for all large $i$ and $0\le t\le1$: each of the finitely many
amplification factors is bounded once $\Delta_i\le1$, and
$e^{Ct}-1\le(e^C-1)t$ on this interval.

The same estimate applies after transporting the identity.
This needs care because $D_i$ need not be invariant under
a coordinate permutation. Reflect the equality on the
reflected torus and then reflect it back, using functional
covariance. The conjugations change a shift $\eta$ to
$-\theta^*\eta$. Indeed, if
$\Phi^\theta(F)=\overline{\Phi(\theta^*\overline F)}$, then
$(\Phi^\theta(\,\cdot-i\eta))^\theta=
\Phi(\,\cdot+i\theta^*\eta)$. Thus
\eqref{eq:diag-outer-identity} becomes an exact identity
for $\Theta P$ with shifts $-\theta^*du_a$ and polynomials
$\Theta P_a$ on the original torus. Iterating reflections
transports it by any $g$, with shifts
$\det(\vec g)(g^{-1})^*du_a$, where $\vec g$ is the linear
part of $g$. Its support box is transported at the same time. Apply \eqref{eq:diag-shifted-tilt} and
\eqref{eq:diag-polynomial-apriori} on $D_i$ itself to obtain
the transported $O(t)$ estimate. No invariance of $D_i$
or of its tilted array under $g$ is used.

\begin{theorem}\label{thm:diag-polynomial-reduction}
For each $P\in\mathcal P$ and coordinate link plane $c$ there
is $Y_P\in\mathcal P_1(c)$, independent of the torus and tilt,
such that, for every fixed lattice isometry $g$,
\begin{equation}\label{eq:diag-polynomial-reduction}
 |\mu_i^t(P^g)-\mu_i^t((Y_P)^g)|\le C_{P,g}t
 \quad(0\le t\le1, i\ge i_0(P,g)).
\end{equation}
At a plane normal to either other coordinate the corresponding
reduction and its reflected estimate hold as well.
\end{theorem}

\begin{proof}
For $c=j_0+\tfrac12$, eliminate lowest planes until the
lowest is at least $j_0$. Every step raises it by at least
one, so this takes finitely many steps. Then eliminate
highest planes until the highest is at most $j_0+1$.
During this phase an eliminated plane $j\ge j_0+2$ can
create links only as low as $j-2\ge j_0$, preserving the
first reduction. The resulting polynomial reads just
the two thick layers and their connecting links. Apply
the upward recursions in thin layers $1,\ldots,5$ of
the lower thick layer, followed by the downward recursions
in thin layers $6,\ldots,2$ of the upper one. New fields
remain at or above the lower layer six and at or below
the upper layer one. The new type three phases are
spatial phases in these two planes or the temporal
phases between them. Thus the result $Y_P$ is in
$\mathcal P_1(c)$.

There are finitely many outer-plane steps, each with
\eqref{eq:diag-outer-error} and its transported version;
the final ten thin-layer recursions are exact with the
unchanged bosonic factor and contribute no error. Take
one volume threshold for all their fixed support boxes
and a priori bounds, and sum the finite $O(t)$ errors.
This proves \eqref{eq:diag-polynomial-reduction}.
For another direction first apply the construction to
$P^{g_\mu^{-1}}$ and then transport its exact identities
by the cyclic coordinate permutation $g_\mu$. The general
transported estimate, evaluated on $D_i$, proves the
claim there and its reflected version.
\end{proof}

\subsection{Comparison states and translation invariance}\label{subsec:diag-comparison-states}

A comparison state for the chosen cell structure and $O$ is
a law $\nu_0$ obtained by the ordered limits
\begin{equation}\label{eq:diag-comparison-state}
 \mu_{i_m}^{t_k}\longrightarrow\nu_{t_k}
       \quad(m\to\infty)\text{ for each }k,
 \qquad
 \nu_{t_k}\longrightarrow\nu_0\quad(k\to\infty),
 \qquad t_k>0,\quad t_k\to0.
\end{equation}
All convergences here are weak convergences of bosonic laws on
$\Omega$, and the volume subsequence is common to all strengths.
Given a prescribed subsequence of volumes, a further subsequence
produces such a state. To see this, take strengths $t_k=1/k$,
extract nested convergent volume subsequences by the tightness
proved above, and choose their diagonal subsequence. The limiting
laws $\nu_{t_k}$ are themselves uniformly tight: applying
\eqref{eq:diag-limit-domination} to $\beta=\pm1_\sigma$ and
using $t_k\le1$ bounds their exponential coordinate moments
by a common constant. The same product-of-intervals argument
gives a weakly convergent strength subsequence. This proves
existence of \eqref{eq:diag-comparison-state} subordinate to
any prescribed volume subsequence.

Every comparison state has Gaussian domination
\begin{equation}\label{eq:diag-comparison-domination}
 \nu_0(e^{\langle\beta,F\rangle})\le e^{\|\beta\|^2/2},
                         \qquad\beta\text{ of finite support}.
\end{equation}
This follows from \eqref{eq:diag-limit-domination} by letting
$t_k\to0$ and truncating the exponential in the second weak
limit. Wall invariance and wall reflection positivity also
pass through both limits. For each local positive-side test,
its support is contained in the positive half-period slab
once $i_m$ is large, so \eqref{eq:diag-tilted-positive-form}
applies. The continuity and $L^2$ extension described in
\S\ref{subsec:diag-wall-law} give positivity for bounded Borel
local tests. In particular $\nu_0$ is invariant and reflection
positive at $(x_1-x_2,c_0)$.

Suppose the unperturbed cubic laws $\mu_{C(n_i)}$ converge to
$\nu$ on the prescribed subsequence. For a comparison state
subordinate to it, the finite fold inequality gives
\begin{equation}\label{eq:diag-comparison-order}
 \nu(O)\le\nu_0(O).
\end{equation}
Indeed fix $t_k>0$, take the volume limit in
\[
 \mu_{C(n_{i_m})}(O)\le\mu_{i_m}^{t_k}(O)
                              +|Q|\Delta_{i_m}/t_k,
\]
and use bounded continuity of $O$ and $\Delta_{i_m}\to0$.
Then let $t_k\to0$. The denominator $t_k$ in the error
precludes reversing this order.

We now obtain the extra invariance required to use
\eqref{eq:diag-comparison-order} for observables at arbitrary
positions. Fix a coordinate link reflection $\theta_c$
and a curvature character $P=\varphi_\beta$. Its polynomial
representation is $e^{i\langle d^*\beta,\widehat A\rangle}$;
on every sufficiently large torus this equals
$e^{i\langle\beta,\widetilde F\rangle}$ on the local lift.
Apply Theorem~\ref{thm:diag-polynomial-reduction} to $P$
and its reflection, and then
Proposition~\ref{prop:diag-two-layer-principle} to its two-layer
reduction. The finite errors are at most two $O(t)$ terms
plus a finite sum of the terms in
\eqref{eq:diag-two-layer-finite-error}. Along the common
volume subsequence in \eqref{eq:diag-comparison-state},
take volume to infinity, then $t_k\to0$ with the regularization
fixed, and finally remove the regularization. Only the
character expectations are passed to weak limits. The
conclusion is
\begin{equation}\label{eq:diag-character-reflection}
 \nu_0(\Theta_c\varphi_\beta)
                           =\overline{\nu_0(\varphi_\beta)}.
\end{equation}
There is no need to extract limiting fermion expectations
in this argument.

For a bosonic character
$\Theta_c\varphi_\beta=\varphi_{-\theta_c^*\beta}$, while
positivity of the probability law gives
$\overline{\nu_0(\varphi_\beta)}=\nu_0(\varphi_{-\beta})$.
Thus \eqref{eq:diag-character-reflection}, with $-\beta$
in place of $\beta$, is invariance under the tensor pullback
$F\mapsto\theta_c^*F$. Characters determine all finite
plaquette marginals by uniqueness of their Fourier
transforms; these marginals determine the law because
cylinder sets generate the Borel sigma algebra of $\Omega$.
Consequently every comparison state is invariant under
all coordinate link reflections, in all three directions.
The product of reflections at $c$ and $c+1$ is translation
by $2e_\mu$, so
\begin{equation}\label{eq:diag-even-translation}
 (\tau_v^*)_*\nu_0=\nu_0,\qquad v\in2\Z^3.
\end{equation}
We need only these even translations. The original wall
periods alone would not have permitted the next placement.

\subsection{Convex separation and the nine mirror families}\label{subsec:diag-separation}

Fix an integer $c_0$. Let $\mathcal R(c_0)$ be the set of
laws on $\Omega$ satisfying Gaussian domination and both
invariance and reflection positivity at
$(x_1-x_2,c_0)$. This is convex and weakly compact. Gaussian
domination gives a common tail
$\lambda(|F_\sigma|\ge R)\le2e^{-R^2/2}$ and hence uniform
tightness. It is weakly closed by truncation of its exponential
integrals. Mirror invariance and positivity against continuous
local tests are also closed, and the $L^2$ approximation
extends the latter to Borel tests. Each condition is affine
or an affine inequality in the law, proving convexity.

We claim that, for every real local trigonometric polynomial
$O$ and every cubic limit $\nu$ along our divisible side
sequence, some $\nu_0\in\mathcal R(c_0)$ satisfies
\eqref{eq:diag-comparison-order}. The polynomial need not
initially fit in a cell having its lower wall at $c_0$.
Let $\supp O\subset y_0+[-r,r]^3$, with $r\ge1$, and choose
an admissible even width $h>8r+8$. Choose an even integer
$a$ in
\[
 (c_0+2r-\ell_-(y_0),\ c_0+h-2r-\ell_-(y_0)).
\]
The interval has length greater than two. For $v=ae_1$ and
$y=y_0+v$ its defining inequalities place the translated
box between the two $\ell_-$ walls. Choose integer $c_+$
and half-integer $c_3$ so that
\[
 c_+<\ell_+(y)-2r<\ell_+(y)+2r<c_++h,
 \qquad c_3<y_3-r<y_3+r<c_3+h.
\]
The available intervals have length greater than one.
The box therefore lies in the open diagonal cell, since
$|\ell_\pm(x)-\ell_\pm(y)|\le2r$ on it. Form the tilted
law with the translated polynomial $O^{\tau_v}$.

Choose a comparison state subordinate to the volume
subsequence defining $\nu$. It belongs to $\mathcal R(c_0)$,
and the fold inequality gives
$\nu(O^{\tau_v})\le\nu_0(O^{\tau_v})$. The original cubic
laws are invariant under every lattice translation, and
this invariance passes to $\nu$. Equation
\eqref{eq:diag-even-translation} gives the same invariance
under this particular $v$ for $\nu_0$. Hence
\begin{equation}\label{eq:diag-dominating-state}
 \nu(O)=\nu(O^{\tau_v})
       \le\nu_0(O^{\tau_v})=\nu_0(O).
\end{equation}
The comparison state may depend on $O$. Convex separation
is what turns these individual inequalities into reflection
positivity of the original law.

\begin{lemma}\label{lem:diag-separation}
Let $\mathcal R$ be a convex weakly compact set of laws on
$\Omega$. If for every real local trigonometric polynomial
$O$ some $\lambda\in\mathcal R$ satisfies
$\nu(O)\le\lambda(O)$, then $\nu\in\mathcal R$.
\end{lemma}

\begin{proof}
Give the real vector space of finite signed measures on
$\Omega$ the weak topology defined by integration against
these polynomials. It is locally convex and Hausdorff:
if every such integral of a signed measure vanishes, its
finite-coordinate marginals have zero Fourier transform,
so its positive and negative parts agree on cylinder sets
and hence on every Borel set. Its continuous linear
functionals are precisely integration against the test
space; a continuous linear functional for a topology
defined by these seminorms depends on finitely many of
them, and their finite real span is still that space.
Since the tests are bounded and continuous, the inclusion
from the usual weak topology of laws is continuous.
Thus $\mathcal R$ is a compact convex subset of this
locally convex Hausdorff vector space. It is nonempty
by the hypothesis at $O=0$.

If $\nu$ were outside $\mathcal R$, strict separation of
a compact convex set and a disjoint closed point would
give a continuous linear functional with
$\max_{\lambda\in\mathcal R}\lambda(O)<\nu(O)$ for some
test $O$. Its maximum exists by compactness. This contradicts
the assumed inequality for that $O$.
\end{proof}

\begin{theorem}\label{thm:diag-lattice-limits}
At fixed $q>0$, $0<\rho<1$, positive chamber weights and
real mass, the cubic curvature laws along the divisible
side sequence are tight on $\Omega$. Every weak limit
satisfies Gaussian domination, is invariant under lattice
translations, signed coordinate permutations and $F\mapsto-F$,
and is reflection positive at all mirrors
\[
 x_\mu=c,\quad c\in\tfrac12\Z;
 \qquad x_\mu\pm x_\nu=c,\quad c\in\Z,\quad\mu<\nu.
\]
Thus it has all nine lattice mirror families.
\end{theorem}

\begin{proof}
The coordinate Gaussian domination theorem gives the
uniform one-plaquette tails and hence tightness. Any
finitely supported $\beta$ embeds without overlap once
the side is large, so its periodization has unchanged
norm. Truncation passes Gaussian domination to every
limit. The laws are supported on the closed subspace
$dF=0$. Finite-volume translation, signed permutation
and charge-conjugation invariances pass to the limit
under their continuous actions on $\Omega$.

For a coordinate mirror $x_\mu=c$, a local positive-side
continuous test lies in its positive half-period slab
for all large sides. The cubic torus is adapted to this
mirror, and its reflected slab trace gives nonnegativity
of the reflected form. Weak convergence and the $L^2$
extension give coordinate reflection positivity of the
limit. For the diagonal mirror $(x_1-x_2,c_0)$,
\eqref{eq:diag-dominating-state} holds for every real
trigonometric polynomial. Lemma~\ref{lem:diag-separation}
therefore gives $\nu\in\mathcal R(c_0)$, proving its
diagonal positivity. This is valid for every integer
$c_0$.

A signed permutation transports this diagonal mirror
and its positive half-space to any plane
$x_\mu\pm x_\nu=c$. The corresponding reflection is the
conjugate of the original one. Pull a positive-side test
back by this isometry; invariance of $\nu$ identifies
its reflected form with the already nonnegative one.
This gives the six diagonal families and completes
the proof.
\end{proof}

The fixed-parameter qualification is used throughout:
the pressure error and every support-dependent amplification
scale are allowed to depend on $q$, $M$ and the weights.
At spacing $\eta=L^{-N}$ the construction has
$q=e\eta^{1/2}$ and
$M=\bar m\eta+m_{0,N}(e)+\frac\rho2\sum_iw_i$, by
\eqref{eq:reg-lattice-mass}. The side sequence in the volume
parameter has the divisibility assumed here. Applying this
theorem at fixed $N$ therefore proves positivity for the
volume limits used in the construction. The subsequent mesh limit must
preserve these volume limits; the argument above supplies
no estimate permitting simultaneous variation of spacing
and volume.


\section{Gaussian domination and tensor moments}\label{sec:probability}

The Gaussian bound is used at three different points: to obtain
compactness of the curvature laws, to control products of field
insertions, and to pass moments to limiting and conditional laws.
We record its consequences in a form independent of the regulator.
The tensor extension below supplies the stated growth axiom on general
tests, rather than only on tensor products.

\subsection{Random linear forms}

Let $\lambda$ be a probability law on the real tempered two-forms,
with its cylinder $\sigma$-algebra, and suppose
\begin{equation}\label{eq:probability-G}
 \E_\lambda e^{F(f)}\le e^{\|f\|_2^2/2}
 \qquad(f\in\Sspace(\R^3;\Lambda^2)\text{ real}).
\end{equation}
Set $X(f)=F(f)$ and extend $X$ complex linearly. For $p>0$ put
$c_p=2\,2^{p/2}\Gamma(1+p/2)$.

\begin{proposition}\label{prop:subgaussian}
The map $X$ extends uniquely and continuously to
$L^2(\R^3;\Lambda^2)$ with values in $L^p(\lambda)$ for every
$p\ge1$. For real $f$ and $u>0$,
\begin{align}
 \E_\lambda X(f)&=0,&
 \E_\lambda X(f)^2&\le\|f\|_2^2,\label{eq:field-variance}\\
 \lambda(|X(f)|\ge u)&\le2e^{-u^2/(2\|f\|_2^2)},&
 \E_\lambda|X(f)|^p&\le c_p\|f\|_2^p\quad(p>0).
 \label{eq:field-moments}
\end{align}
If $\|f\|_2=0$, the tail bound is interpreted as zero for $u>0$.
For real $f_1,\ldots,f_n$,
\begin{equation}\label{eq:joint-moment-bound}
 \E_\lambda\prod_{j=1}^n|X(f_j)|
 \le c_n\prod_j\|f_j\|_2,
 \qquad c_n\le2\sqrt{2^n n!}.
\end{equation}
For complex tests the same bound holds with an additional factor
$2^{n/2}$. The generating function of any finite family of tests is
entire, and
\begin{equation}\label{eq:entire-generating-bound}
 \left|\E_\lambda e^{X(f)}\right|
 \le e^{\|\operatorname{Re}f\|_2^2/2}.
\end{equation}
\end{proposition}

\begin{proof}
For a real test, \eqref{eq:probability-G} applied to $tf$ gives
$\E_\lambda e^{tX(f)}\le e^{t^2\|f\|_2^2/2}$. Also
\[
 \E_\lambda e^{t|X(f)|}\le
 \E_\lambda e^{tX(f)}+\E_\lambda e^{-tX(f)}
 \le2e^{t^2\|f\|_2^2/2}.
\]
Thus the exponential series can be differentiated at zero. Comparing
its linear and quadratic terms with the Gaussian upper bound, for
both signs of $t$, gives \eqref{eq:field-variance}.
Markov's inequality with $t=u/\|f\|_2^2$ gives the upper tail;
applying it to $-f$ gives the lower tail. Integrating this bound yields
\[
 \E_\lambda|X(f)|^p
 =p\int_0^\infty u^{p-1}\lambda(|X(f)|>u)\,\dd u
 \le2\,2^{p/2}\Gamma(1+p/2)\|f\|_2^p.
\]
H\"older's inequality gives \eqref{eq:joint-moment-bound};
log-convexity of $\Gamma$ gives
$\Gamma(1+n/2)\le\sqrt{n!}$.
For $f=g+ih$, the triangle inequality in $L^n(\lambda)$ and
$\|g\|_2+\|h\|_2\le\sqrt2\|f\|_2$ give the complex bound.
These estimates applied to differences extend $X$ from real Schwartz
tests to $L^2$, consistently for all $p\ge1$.

On a compact polydisc in the variables of a finite generating
function, its integrand is bounded by a product of exponentials of
$|X(\operatorname{Re}f_j)|$ and $|X(\operatorname{Im}f_j)|$.
H\"older's inequality and the preceding exponential bound make this
majorant integrable. Differentiation under the integral proves the
entire extension.
Taking the modulus of the integrand and applying
\eqref{eq:probability-G} proves \eqref{eq:entire-generating-bound}.
\end{proof}

For fixed $n$, the moment bound controls differences of products by
expanding one factor at a time. In particular,
$S_n(f_1,\ldots,f_n)=\E_\lambda\prod_jX(f_j)$ is a bounded
multilinear form on complex $L^2$, symmetric in its arguments and
real on real tests.

\subsection{Generating functions and cumulants}

Near the origin the logarithm of a generating function is well defined
uniformly in the law. If $f=g+ih$, the elementary remainder bound
$|e^z-1-z|\le e^{|z|}-1-|z|$, convexity, and
$\E_\lambda X(f)=0$ give
\begin{equation}\label{eq:zero-free-bound}
 \bigl|\E_\lambda e^{X(f)}-1\bigr|
 \le e^{2\|f\|_2^2}-1.
\end{equation}
Indeed, writing $\varphi(t)=e^t-1-t$, one has
$\varphi(t)\le2(\cosh t-1)$ for $t\ge0$, and
\begin{align*}
 \E_\lambda\varphi(|X(f)|)
 &\le\tfrac12\E_\lambda\varphi(2|X(g)|)
       +\tfrac12\E_\lambda\varphi(2|X(h)|)\\
 &\le e^{2\|g\|_2^2}-1+e^{2\|h\|_2^2}-1
 \le e^{2\|f\|_2^2}-1.
\end{align*}
Consequently
\[
 W(z)=\log\E_\lambda\exp\Bigl(\sum_{j=1}^nz_jX(f_j)\Bigr)
\]
is analytic and satisfies $|W(z)|\le\log2$ on
$\sum_j|z_j|\|f_j\|_2\le2/5$, using the branch with $W(0)=0$.
Cauchy's formula on the polydisc of radii
$2/(5n\|f_j\|_2)$ gives
\begin{equation}\label{eq:cumulant-Cauchy}
 |\partial_1\cdots\partial_nW(0)|
 \le(\log2)(5n/2)^n\prod_j\|f_j\|_2.
\end{equation}
If a test has zero norm, its random variable vanishes and the
corresponding derivative is zero. These derivatives are the joint
cumulants used in the fourth-order response calculation.

The entire extension also gives moment determinacy. Two laws obeying
\eqref{eq:probability-G} and having the same moments have the same
entire generating functions for every finite family of real tests.
Their characteristic functions, and hence their finite-dimensional
distributions, coincide. Cylinder sets generate the chosen
$\sigma$-algebra, so the two laws coincide. We will use this after
proving rotation covariance of all moments.

\subsection{Extension to tensor tests}

Let $A=1+|x|^2-\Delta$ and let
$h_{\kappa}=h_k e_{ij}$, $k\in\Z_{\ge0}^3$, $i<j$, be the
real Hermite basis of $L^2(\R^3;\Lambda^2)$. Its eigenvalues are
$\lambda_\kappa=2|k|_1+4$. For a tensor test write
\[
 \Psi=\sum_{\boldsymbol\kappa}c_{\boldsymbol\kappa}(\Psi)
 h_{\kappa_1}\otimes\cdots\otimes h_{\kappa_n}.
\]
The Hermite expansion converges in Schwartz space. The estimate we
need is
\begin{equation}\label{eq:Hermite-l1}
 \sum_{\boldsymbol\kappa}|c_{\boldsymbol\kappa}(\Psi)|
 \le4^{-n}(34\pi\sqrt3)^n q_{6n}(\Psi).
\end{equation}
To see the derivative order and the dependence on $n$, set
$Z_4=\sum_\kappa\lambda_\kappa^{-4}<1/16$. Weighted
Cauchy--Schwarz and Parseval give
\[
 \sum_{\boldsymbol\kappa}|c_{\boldsymbol\kappa}(\Psi)|
 \le Z_4^{n/2}\|(A^2)^{\otimes n}\Psi\|_2.
\]
With $M(x)=1+|x|^2$,
\[
 A^2=M^2-2M\Delta-4x\cdot\nabla-6+\Delta^2.
\]
The sum of the moduli of the differential coefficients of $A^2$ is
bounded by $34M^2$. Applying this expression in each of the $n$ arguments gives
the componentwise bound
\[
 |(A^2)^{\otimes n}\Psi(x)|
 \le34^n q_{6n}(\Psi)\prod_{a=1}^n(1+|x_a|^2)^{-1}.
\]
Since $\int_{\R^3}(1+|x|^2)^{-2}\,\dd x=\pi^2$ and there
are $3^n$ components, \eqref{eq:Hermite-l1} follows.
The bound on $Z_4$ follows by counting Hermite eigenvalues:
\[
 Z_4=\frac3{32}\sum_{d\ge0}\frac{d+1}{(d+2)^3}
 \le\frac3{32}\Bigl(\frac{\pi^2}{6}-1\Bigr)<\frac1{16}.
\]

Define the random tensor evaluation by the $L^2(\lambda)$ sum
\begin{equation}\label{eq:tensor-random-evaluation}
 Q_\Psi=\sum_{\boldsymbol\kappa}c_{\boldsymbol\kappa}(\Psi)
 \prod_{a=1}^n X(h_{\kappa_a}).
\end{equation}
The product moment bound, now applied to $2n$ real tests of norm one,
shows that this sum converges absolutely in $L^2(\lambda)$ and
\[
 \|Q_\Psi\|_{L^2(\lambda)}
 \le\sqrt{c_{2n}}\sum_{\boldsymbol\kappa}|c_{\boldsymbol\kappa}(\Psi)|.
\]
For a decomposable test it equals $\prod_aX(f_a)$, by the Hermite
expansion of each $f_a$ and the same product bounds.
Thus $S_n(\Psi):=\E_\lambda Q_\Psi$ is the unique continuous
extension of the moments to $\Sspace_n$. Moreover,
\begin{align*}
 |S_n(\Psi)|
 &\le c_n\sum_{\boldsymbol\kappa}|c_{\boldsymbol\kappa}(\Psi)|\\
 &\le2\sqrt{n!}\left(\frac{34\pi\sqrt6}{4}\right)^nq_{6n}(\Psi)
 \le2\,66^n\sqrt{n!}\,q_{6n}(\Psi).
\end{align*}
This is \textup{(E0$'$)}. The construction also provides the
$L^2$ tensor evaluations needed when passing the cluster estimate
from local observables to Schwinger functions.

\subsection{The topology of the limits}

The Hilbert scale used for the limit laws is
\[
 \|F\|_{-s}^2=\sum_\kappa\lambda_\kappa^{-2s}|F(h_\kappa)|^2,
 \qquad H_{-s}=\{F\in\Sspace':\|F\|_{-s}<\infty\}.
\]
A closed bounded ball of $H_{-3}$ is compact in $H_{-4}$. In Hermite
coordinates, the tail of the squared $H_{-4}$ norm on the indices
$\lambda_\kappa>M$ is bounded by $M^{-2}$ times the squared
$H_{-3}$ norm; the remaining coordinates form a finite-dimensional
space. A diagonal subsequence argument proves compactness.

The lattice application will give the uniform bound
$\E\|F^\sharp\|_{-3}^2\le61$. Its consequence is
\[
 P_{N,T}(\|F\|_{-3}>r)\le61r^{-2},
\]
so the laws are tight on the separable Hilbert space $H_{-4}$.
Prokhorov's theorem gives subsequential weak limits, and the inclusion
$H_{-4}\hookrightarrow\Sspace'$ is continuous for the strong
distribution topology.

For fixed $N$, let $\mathcal V_N$ be the set of weak limits in
$H_{-4}$ obtained as $T\to\infty$. A volume-first limit is a weak
limit of laws $\lambda_i\in\mathcal V_{N_i}$ with $N_i\to\infty$.
The uniform compactness bound makes each $\mathcal V_N$ nonempty.
The $H_{-3}$ balls are closed in $H_{-4}$, so the portmanteau theorem
passes their probability lower bound to every fixed-spacing volume
limit. The family $\bigcup_{N\ge0}\mathcal V_N$ therefore obeys the
same tightness bound on $H_{-4}$, giving volume-first limit points.
Their additional diagonal positivity depends on taking the volume
limit at each fixed spacing; compactness alone would not supply it.


\section{Volume and continuum limits}\label{sec:continuum-limits}

Gaussian domination gives compactness of the curvature laws at all
cutoffs. Diagonal reflection positivity uses a further identification:
at a fixed spacing, a volume limit is the step-field image of an
infinite lattice law. This permits the nine-mirror theorem to be
applied before the spacing is removed.

\subsection{The test map and its bounds}

Write $\eta=L^{-N}$, $R=8hL^T$, and $n=R/\eta$, as in
\eqref{eq:cutoff-scales}. A plaquette of the unit covering lattice is
$\sigma=(x;\mu\nu)$, with $\mu<\nu$ and centre
$y_\sigma=x+(e_\mu+e_\nu)/2$. Its averaging cube is
\[
 Q_\eta(\sigma)=\eta\bigl(y_\sigma+[-1/2,1/2)^3\bigr).
\]
For each fixed orientation these cubes partition $\R^3$.
If $u$ is a real lattice curvature, its physical step field is
$\eta^{-3/2}u(\sigma)$ on $Q_\eta(\sigma)$ in component $\mu\nu$.
For a Schwartz two-form $f$, define
\begin{align}
 f_{[\eta]}(\sigma)
   &=\eta^{-3/2}\int_{Q_\eta(\sigma)}f_{\mu\nu}(y)\,\dd y,
                                                   \label{eq:limit-test-cover}\\
 f^R(y)&=\sum_{k\in\Z^3}f(y+Rk),\notag\\
 f_\eta([\sigma])
   &=\eta^{-3}\int_{Q_\eta(\sigma)}f^R_{\mu\nu}(y)\,\dd y.
                                                   \label{eq:limit-test-torus}
\end{align}
Here $[\sigma]$ is the plaquette class modulo $n\Z^3$.
The torus cochain norm is $|v|_\eta^2=\eta^3\sum_{[\sigma]}|v([\sigma])|^2$.
The covering cochain norm $\|\cdot\|$ is the unweighted $\ell^2$ norm.
For the periodic extension $\widetilde u$ of $u$,
\begin{equation}\label{eq:limit-step-pairing}
 F^\sharp(f)=\eta^{3/2}\sum_{[\sigma]}f_\eta([\sigma])u(\sigma)
            =\sum_\sigma f_{[\eta]}(\sigma)\widetilde u(\sigma).
\end{equation}
The normalization in the first equality combines the cell volume
$\eta^3$ with the physical field factor $\eta^{-3/2}$.

The estimates needed below are
\begin{align}
 \|f_{[\eta]}\|&\le\|f\|_2, \label{eq:limit-cover-norm}\\
 \sum_\sigma|f_{[\eta]}(\sigma)|(1+|x(\sigma)|_\infty)
 &\le\eta^{-3/2}\int_{\R^3}\sum_{\mu<\nu}|f_{\mu\nu}(y)|
                         (2+|y|_\infty/\eta)\,\dd y, \label{eq:limit-cover-sum}\\
 |f_\eta|_\eta&\le\|f^R\|_{L^2(\R^3/R\Z^3)}\le p_R(f),\label{eq:limit-torus-norm}
\end{align}
where, with $Q_R=[-R/2,R/2)^3$,
\begin{equation}\label{eq:limit-periodic-seminorm}
 p_R(f)=\|f\|_2+
        26R^{-2}\||y|^2f\|_{L^2(\R^3\setminus Q_R)}.
\end{equation}
One also has $|f_\eta|_\eta\le\|f\|_2$ if the support of $f$ lies
in a translate of an open cube of side $R$. For $R\ge1$,
$p_R(f)\le p_1(f)$, and $p_R(f)\to\|f\|_2$ as $R\to\infty$.
The seminorm $p_1$ is continuous on Schwartz space.

The first bound follows from Cauchy--Schwarz on each averaging
cube. On that cube $1+|x(\sigma)|_\infty\le2+|y|_\infty/\eta$,
which gives the second. Absolute convergence permits unfolding:
\[
 \eta^{3/2}f_\eta([\sigma])
       =\sum_{k\in\Z^3}f_{[\eta]}(\sigma+nk).
\]
This proves \eqref{eq:limit-step-pairing}. Applying the cube
Cauchy--Schwarz inequality to $f^R$ gives the first inequality in
\eqref{eq:limit-torus-norm}. If the support lies in one open
$R$-cube, its distinct $R\Z^3$ translates are disjoint, so
$\|f^R\|_{L^2(\R^3/R\Z^3)}=\|f\|_2$.

For the remaining estimate put $a_k=\|f\|_{L^2(Q_R+Rk)}$.
Minkowski's inequality bounds the periodized norm by $\sum_k a_k$.
On $Q_R+Rk$, $k\ne0$, one has $|y|\ge R|k|_\infty/2$; hence
\[
 \sum_{k\ne0}a_k
 \le4R^{-2}\left(\sum_{k\ne0}|k|_\infty^{-4}\right)^{1/2}
                   \||y|^2f\|_{L^2(\R^3\setminus Q_R)}.
\]
The number of points with $|k|_\infty=j$ is $24j^2+2$, so
the sum in parentheses is $4\pi^2+\pi^4/45<41.65$ and its
square root times four is less than $26$.
This proves \eqref{eq:limit-torus-norm} and
\eqref{eq:limit-periodic-seminorm}. The comparison with $p_1$
and the limit of $p_R$ follow by inclusion of the cubes and
integrability of $|y|^4|f(y)|^2$.

For the Hermite two-forms of Section~\ref{sec:probability},
\begin{equation}\label{eq:limit-Hermite-norm}
 p_1(h_\kappa)\le31\lambda_\kappa.
\end{equation}
Here is the weighted estimate behind this bound. Let
$\mathsf M=1+|y|^2$, $P_j=-i\partial_j$, and
$\mathsf A=\mathsf M-\Delta$. Since
$[P_j,[P_j,\mathsf M]]=-2$,
\[
 \|\mathsf A v\|_2^2
 =\|\mathsf Mv\|_2^2+\|\Delta v\|_2^2
      +2\sum_j\langle P_jv,\mathsf MP_jv\rangle-6\|v\|_2^2
 \ge\|\mathsf Mv\|_2^2-6\|v\|_2^2.
\]
The Hermite eigenvalues satisfy $\lambda_\kappa\ge4$, and therefore
$\|\mathsf Mv\|_2^2\le(11/8)\|\mathsf Av\|_2^2$.
Taking $v=h_\kappa$ gives
$p_1(h_\kappa)\le1+26\sqrt{11/8}\lambda_\kappa
\le31\lambda_\kappa$.

The test map respects the affine cubic symmetries. For a signed
permutation and lattice translation $g$ of the unit lattice, let
$\mathfrak g(y)=\eta g(y/\eta)$. With the exterior-square action
on two-forms and the cochain pull-back,
\begin{equation}\label{eq:limit-test-equivariance}
 (\mathfrak g\cdot f)_{[\eta]}=(g^{-1})^*f_{[\eta]},
 \qquad
 (\mathfrak g\cdot f)_\eta=(g^{-1})^*f_\eta.
\end{equation}
Indeed, $g$ sends centres to centres and the averaging cubes to
averaging cubes up to null boundaries. Its orientation sign is
the sign of $\wedge^2\vec g$ on the corresponding plaquette.
Changing variables gives the first formula; periodization and
the bijection $k\mapsto\vec gk$ give the second.

\subsection{Compactness and convergence of moments}

Let $P_{N,T}$ be the law of $F^\sharp$ for the all-flux chamber
measure. The mass assignment may be any real sequence independent of
$T$. Theorem~\ref{thm:slab-gaussian} and
\eqref{eq:limit-step-pairing} give
\begin{equation}\label{eq:limit-cutoff-G}
 \E_{N,T}e^{F^\sharp(f)}
       \le e^{|f_\eta|_\eta^2/2}\le e^{p_R(f)^2/2}
       \qquad(f\text{ real}).
\end{equation}
The random-linear-form proof of Proposition~\ref{prop:subgaussian}
applies here with the Hilbert seminorm $f\mapsto|f_\eta|_\eta$.
In particular its variance bound and \eqref{eq:limit-Hermite-norm}
give
\begin{equation}\label{eq:limit-uniform-tightness}
 \E_{N,T}\|F^\sharp\|_{-3}^2
 \le\sum_\kappa\lambda_\kappa^{-6}p_1(h_\kappa)^2
 \le961\sum_\kappa\lambda_\kappa^{-4}<61.
\end{equation}
Tonelli's theorem justifies the first sum. Every finite lattice
cochain has a step field in $H_{-3}$: Cauchy--Schwarz in
\eqref{eq:limit-step-pairing} bounds its $\kappa$-th Hermite
coordinate by $31\lambda_\kappa\|u\|_{\ell^2(C(n)^{(2)})}$,
with the unweighted norm on the unit-lattice torus plaquettes.
The step map from this finite-dimensional cochain space to $H_{-3}$
is consequently continuous.

The compactness argument at the end of Section~\ref{sec:probability}
now applies. The laws $P_{N,T}$ are uniformly tight on $H_{-4}$,
and their fixed-spacing volume limits are uniformly tight there.
Thus the sets $\mathcal V_N$ and the set of volume-first limits
are nonempty.

For a joint sequence $N_k,T_k\to\infty$ suppose
$P_{N_k,T_k}\to\mu$ weakly on $H_{-4}$. Truncating $e^{F(f)}$
above by a constant, passing to the weak limit, and then removing
the truncation gives
\begin{equation}\label{eq:limit-continuum-G}
 \E_\mu e^{F(f)}\le e^{\|f\|_2^2/2}.
\end{equation}
The limiting constant follows from $p_{R_k}(f)\to\|f\|_2$.
For a fixed finite family of complex tests $f_1,\ldots,f_s$, define
\[
 Z_k(z)=\E_{N_k,T_k}\exp\left(\sum_a z_aF^\sharp(f_a)\right),
 \qquad
 Z(z)=\E_\mu\exp\left(\sum_a z_aF(f_a)\right).
\]
The evaluations $F\mapsto(F(f_a))_a$ are continuous on $H_{-4}$.
Moreover,
\[
 \E_{N_k,T_k}\left|\exp\left(\sum_a z_aF^\sharp(f_a)\right)\right|^2
 \le\exp\left(2p_1\left(\operatorname{Re}\sum_a z_af_a\right)^2\right).
\]
This bound is uniform in $k$ and on compact sets of $z$.
Uniform integrability therefore gives pointwise convergence
$Z_k\to Z$. The same exponential bounds give local uniform
bounds on the entire functions $Z_k$. Cauchy's estimates make them
equi-Lipschitz on compact sets, so their convergence is locally
uniform. Cauchy's formula then passes every derivative to the limit.
In particular all joint moments converge.

On each compact subset of
$\{\sum_a|z_a|\|f_a\|_2<2/5\}$, the test-map bounds give
$\sum_a|z_a||(f_a)_{\eta_k}|_{\eta_k}<2/5$ for all large $k$.
The zero-free estimate \eqref{eq:zero-free-bound}, with this
seminorm, bounds $|Z_k-1|$ by $1/2$. The principal logarithms
therefore converge locally uniformly near zero, and their
derivatives give convergence of every joint cumulant.

\subsection{Closedness and the affine cubic action}

The step fields are closed on the lattice. To pass this identity
to distributions, let $\xi\in C_c^\infty(\R^3;\R)$ represent the
three-form $\xi\,\dd y_1\wedge\dd y_2\wedge\dd y_3$. Put
\[
 \partial_j^\eta\xi(y)
  =\eta^{-1}\bigl(\xi(y+\eta e_j/2)-\xi(y-\eta e_j/2)\bigr)
\]
and define the two-form $v^\eta$ in the component order $23,13,12$
by
\[
 (v^\eta_{23},v^\eta_{13},v^\eta_{12})
       =(-\partial_1^\eta\xi,\partial_2^\eta\xi,-\partial_3^\eta\xi).
\]
It satisfies
\begin{equation}\label{eq:limit-Bianchi-error}
 F^\sharp(v^\eta)=0,\qquad
 \|v^\eta-d^*\xi\|_2
        \le\frac{\eta^2}{24}\sum_{j=1}^3\|\partial_j^3\xi\|_2.
\end{equation}

For the exact identity, assign to each unit cube the number
$\xi_{[\eta]}(x)=\eta^{-5/2}\int_{\eta(x+[0,1)^3)}\xi(y)\,\dd y$.
The three plaquette coefficients of $v^\eta_{[\eta]}$ are
\[
 -(\xi_{[\eta]}(x)-\xi_{[\eta]}(x-e_1)),\quad
  \xi_{[\eta]}(x)-\xi_{[\eta]}(x-e_2),\quad
 -(\xi_{[\eta]}(x)-\xi_{[\eta]}(x-e_3)).
\]
In the pairing with $\widetilde u$ these regroup to
$\sum_x\xi_{[\eta]}(x)(d\widetilde u)(x)=0$.
The weighted summability in \eqref{eq:limit-cover-sum} justifies
regrouping; at finite volume $\widetilde u$ is bounded.
For the error bound, the Fourier multiplier of
$\partial_j^\eta-\partial_j$ is
$i(2\eta^{-1}\sin(\eta p_j/2)-p_j)$, whose modulus is at most
$\eta^2|p_j|^3/24$. Plancherel's theorem gives
\eqref{eq:limit-Bianchi-error}.

For all large $k$, $d^*\xi-v^{\eta_k}$ is supported in one open
$R_k$-cube. The variance bound and \eqref{eq:limit-Bianchi-error}
then give
\[
 \E_{N_k,T_k}F^\sharp(d^*\xi)^2
 \le\|d^*\xi-v^{\eta_k}\|_2^2
 \le\frac{\eta_k^4}{24^2}
                \left(\sum_j\|\partial_j^3\xi\|_2\right)^2.
\]
Passing bounded truncations of the square to the weak limit shows
$F(d^*\xi)=0$ almost surely. Choose a countable dense family of
compactly supported smooth three-form tests. Continuity of each
tempered distribution extends the identity from this family to all
tests, so $\mu(dF=0)=1$.

The finite laws are invariant under lattice translations, signed
coordinate permutations and $u\mapsto-u$. By
\eqref{eq:limit-test-equivariance} their step-field characteristic
functionals have the corresponding invariances. Signed permutations
and change of sign pass directly to the weak limit. For a real
translation $b$, choose $b_k\in\eta_k\Z^3$ with
$|b_k-b|_\infty\le\eta_k/2$. The variance bound gives
\[
 \left|\E_{N_k,T_k}e^{iF^\sharp(\tau_bf)}
             -\E_{N_k,T_k}e^{iF^\sharp(\tau_{b_k}f)}\right|
       \le p_1(\tau_bf-\tau_{b_k}f)\longrightarrow0.
\]
Thus the limiting characteristic functional is invariant under
every translation. Characteristic functionals determine the law,
so $\mu$ has these invariances.

\subsection{Identifying a fixed-spacing volume limit}

Fix $N$, hence $\eta$, and take any $\mu^{(N)}\in\mathcal V_N$.
Choose a defining sequence $T_j\to\infty$ with
$P_{N,T_j}\to\mu^{(N)}$ weakly on $H_{-4}$.
Let $\nu_j$ be the law on $\Omega$ of the periodically extended
unit-lattice curvature. We shall identify a product-topology limit
of these laws and prove
\begin{equation}\label{eq:limit-volume-identification}
 \nu_j\longrightarrow\nu^{(N)},\qquad
 \mu^{(N)}=(\sharp_\eta)_*\nu^{(N)}.
\end{equation}
The step map in this identity is only Borel; the identification
requires more than the continuous mapping theorem.

For a closed cochain $u$ put
\[
 \Omega'=\left\{u\in\Omega:
       \sup_\sigma\frac{|u(\sigma)|}{1+|x(\sigma)|_\infty}<\infty\right\}.
\]
On $\Omega'$ the componentwise step field $\sharp_\eta u$ is
tempered, and
$(\sharp_\eta u)(f)=\sum_\sigma f_{[\eta]}(\sigma)u(\sigma)$
is absolutely convergent by \eqref{eq:limit-cover-sum}.
Let $\Omega^{(\eta)}$ be the subset on which this distribution
has finite $H_{-3}$ norm. Regard $\sharp_\eta$ as a map
$\Omega\to H_{-4}$ by setting it to zero outside $\Omega^{(\eta)}$.
Both subsets and the map are Borel. Indeed, the growth bound is
a countable supremum; each test evaluation is a pointwise limit
of finite coordinate sums; and the squared $H_{-3}$ norm is a
countable sum of nonnegative evaluations. The Hermite evaluations
generate the Borel sigma algebra of $H_{-4}$.

For each plaquette choose a real smooth bump $b^\sigma$, with
only its corresponding two-form component nonzero, supported in
the interior of $Q_\eta(\sigma)$ and of integral $\eta^{3/2}$.
Then $(b^\sigma)_{[\eta]}=\mathbf1_\sigma$. The map
\begin{equation}\label{eq:limit-step-inverse}
 \mathcal J_\eta:H_{-4}\longrightarrow\R^{\{\text{plaquettes}\}},
 \qquad \mathcal J_\eta(F)(\sigma)=F(b^\sigma)
\end{equation}
is continuous for the product topology, and
$\mathcal J_\eta\sharp_\eta u=u$ on $\Omega^{(\eta)}$.
The finite-volume identities are
$P_{N,T_j}=(\sharp_\eta)_*\nu_j$ and
$\nu_j=(\mathcal J_\eta)_*P_{N,T_j}$. Thus the continuous mapping theorem
gives $\nu_j\to\nu^{(N)}:=(\mathcal J_\eta)_*\mu^{(N)}$ in the product
space. Closedness of $\Omega$ and portmanteau show that
$\nu^{(N)}$ is supported on $\Omega$.

This lattice law satisfies Gaussian domination for every
finitely supported cochain $\beta$. Once $n_j$ is large, the
support of $\beta$ embeds without overlap in the cubic torus,
and its periodization has unchanged $\ell^2$ norm.
Theorem~\ref{thm:slab-gaussian}, weak convergence and bounded
truncations of $e^{\langle\beta,u\rangle}$ give
\begin{equation}\label{eq:limit-lattice-G}
 \int_\Omega e^{\langle\beta,u\rangle}\,\dd\nu^{(N)}(u)
                          \le e^{\|\beta\|^2/2}.
\end{equation}

For any law $\nu$ on $\Omega$ with this bound, the one-coordinate
tail estimate is $\nu(|u(\sigma)|\ge t)\le2e^{-t^2/2}$.
There are at most $3(2k+1)^3$ plaquettes with base point of
sup norm $k$, and hence
\[
 \sum_\sigma\nu\bigl(|u(\sigma)|\ge1+|x(\sigma)|_\infty\bigr)
 \le6\sum_{k\ge0}(2k+1)^3e^{-(k+1)^2/2}<\infty.
\]
Borel--Cantelli gives $\nu(\Omega')=1$.
Truncate $f_{[\eta]}$ to a finite set of plaquettes.
Absolute convergence on $\Omega'$ and Fatou's lemma pass
\eqref{eq:limit-lattice-G} to this cochain, giving
\[
 \int e^{(\sharp_\eta u)(f)}\,\dd\nu(u)
       \le e^{\|f_{[\eta]}\|^2/2}\le e^{\|f\|_2^2/2}.
\]
Its variance bound and Tonelli's theorem yield
\[
 \int\|\sharp_\eta u\|_{-3}^2\,\dd\nu(u)
      \le\sum_\kappa\lambda_\kappa^{-6}.
\]
In particular $\nu(\Omega^{(\eta)})=1$.

It remains to recover $\mu^{(N)}$ from $\nu^{(N)}$.
Choose a countable Schwartz-dense family $\mathcal D_0$ of real
compactly supported smooth two-forms. For $f\in\mathcal D_0$ set
\[
 f'=f-\sum_\sigma f_{[\eta]}(\sigma)b^\sigma.
\]
The sum is finite and $f'_{[\eta]}=0$. Consequently
$F^\sharp(f')=0$ at every finite volume. The hyperplane
$\{F:F(f')=0\}$ is closed in $H_{-4}$, so it has full
$\mu^{(N)}$ measure. Intersecting these sets for
$f\in\mathcal D_0$ with the full-measure set
$\mathcal J_\eta(F)\in\Omega^{(\eta)}$ gives
\[
 F(f)=\sum_\sigma f_{[\eta]}(\sigma)F(b^\sigma)
              =(\sharp_\eta \mathcal J_\eta(F))(f)
                  \qquad(f\in\mathcal D_0)
\]
on one set of full measure. Both sides are tempered distributions,
so density extends this equality to every test. This proves
\eqref{eq:limit-volume-identification}.

\subsection{Nine mirror families in the continuum}

At fixed $N$ the charge $q=e\eta^{1/2}$ and the chamber mass
parameter $M=\bar m\eta+m_{0,N}+\rho\sum_iw_i/2$ are independent of volume.
The sides $n_j=8hL^{T_j+N}$ satisfy the divisibility condition
of Theorem~\ref{thm:diag-lattice-limits}: every required power
of $L$ eventually divides the side, and $8$ divides it.
That theorem applies to $\nu^{(N)}$. It is invariant under
the affine cubic group and change of sign, and reflection
positive at the three coordinate and six diagonal mirror
families. The identification
\eqref{eq:limit-volume-identification} transports its
characteristic functional to $\mu^{(N)}$.

Write $\ell(y)=n_\ell\cdot y$ for one of
$y_j,y_j-y_k,y_j+y_k$, $j<k$. The admissible unit-lattice
heights are $\tfrac12\Z$ for a coordinate family and $\Z$
for a diagonal family. If $c$ is an admissible physical
height, the cochain test map intertwines the lattice
reflection with $r_{\ell,c}$ by
\eqref{eq:limit-test-equivariance}. Suppose that real
$f_1,\ldots,f_s\in C_c^\infty$ have support in
$\{\ell\ge c+\delta\}$ with $\delta>2\eta$.
If $f_{a,[\eta]}(\sigma)\ne0$, its averaging cube meets
this support. Every corner $\eta x$ of the plaquette is
within sup distance $\eta$ of a point $y$ in that intersection.
Since $|n_\ell|_1\le2$,
\[
 \ell(\eta x)\ge\ell(y)-2\eta>c.
\]
Thus $f_{a,[\eta]}$ has finite support strictly on the
positive lattice side. Lattice reflection positivity gives
\begin{equation}\label{eq:limit-mirror-characters}
 \sum_{a,b}\overline{z_a}z_b
   \E_{\mu^{(N)}}e^{iF(f_b-r_{\ell,c}\cdot f_a)}\ge0.
\end{equation}
This is the reflected form of the character sum
$\sum_bz_be^{i\langle f_{b,[\eta]},u\rangle}$.

\begin{proposition}\label{prop:limit-nine-mirrors}
Every volume-first limit of the chamber cutoff laws is a
translation-invariant law of closed tempered two-forms with
Gaussian domination. It is invariant and reflection positive
at every plane in the nine families
$y_j=c$ and $y_j\pm y_k=c$, $j<k$, $c\in\R$.
All finite-family moments and cumulants converge along a
joint cutoff sequence realizing that limit.
\end{proposition}

\begin{proof}
Suppose $\mu^{(N_i)}\to\mu$ weakly on $H_{-4}$, with
$N_i\to\infty$ and $\mu^{(N_i)}\in\mathcal V_{N_i}$.
Fix a metric for weak convergence of probability laws on
$H_{-4}$. For each $i$, choose a volume index $T_i\ge i$
along the defining volume subsequence such that the distance
between $P_{N_i,T_i}$ and $\mu^{(N_i)}$ is less than $1/i$.
Then $P_{N_i,T_i}\to\mu$ with both indices tending to infinity.
The preceding compactness, moment convergence, closedness and
invariance arguments apply to this joint sequence.

Fix $\ell$ and first take $c\in\Z[L^{-1}]$.
Then $c\in\eta_i\Z$ for all large $i$, and hence it is an
admissible height for each lattice mirror family.
Any finite family of compact tests in $\{\ell>c\}$ has
a common support margin $\delta>0$.
For all large $i$, $2\eta_i<\delta$, so
\eqref{eq:limit-mirror-characters} holds for $\mu^{(N_i)}$.
The functions in this finite sum are bounded and continuous
on $H_{-4}$; weak convergence passes its nonnegativity to $\mu$.

For an arbitrary height $c$, choose $c'\in\Z[L^{-1}]$ and
$v=(c-c')|n_\ell|^{-2}n_\ell$. Then
$r_{\ell,c}=\tau_vr_{\ell,c'}\tau_{-v}$.
Translate the test forms by $-v$; they are now supported on
the positive side of $\ell=c'$. Translation invariance
identifies the reflected form at $c$ with that at $c'$.
Thus \eqref{eq:limit-mirror-characters} holds at every
height. Character sums are dense in the corresponding
positive-side $L^2$ space, and the reflected form is bounded
there by Cauchy--Schwarz. Its nonnegativity therefore extends
to that space. Finally $r_{\ell,0}$ is a signed permutation,
and $r_{\ell,c}$ is its composition with a translation.
The invariances already proved give invariance at each mirror.
\end{proof}


\section{Reflection-positive reconstruction, the spectral cone, and clustering}
\label{sec:continuum-cone}

The diagonal mirrors supply an energy--momentum bound with constant
one. We prove this directly by comparing two reflection-positive
reconstructions. This avoids using rotational covariance in the step
that will later provide its spectral input.

\subsection{The reflection-positive Hilbert space}

For an open set $O$, let $\Sigma_O$ be generated by $F(f)$ with
real $f\in C_c^\infty(O;\Lambda^2)$. Let $\lambda$ be a
translation-invariant law and let $r_j$ reflect in $x_j=0$.
Suppose that $\lambda$ is invariant and reflection positive at this
plane. For a Euclidean motion $g$ write
$\widehat g A(F)=A(g^{-1}F)$ and $U(a)=\widehat{\tau_a}$.
For a reflection $r$ put
$\mathcal R_r=\widehat r$. Write $\mathcal R_j=\mathcal R_{r_j}$ and
$\Theta_j A=\overline{\mathcal R_j A}$.
On $\mathcal E_j^+=L^2(\Sigma_{\{x_j>0\}},\lambda)$ the form
\[
 b_j(A,B)=\E_\lambda[(\Theta_j A)B]
\]
is positive semidefinite and $b_j(A,A)\le\|A\|_{L^2(\lambda)}^2$.
Let $\mathcal H_j$ be its quotient completion, $[A]_j$ the quotient
map, and $\Omega_j=[1]_j$.
Reflection positivity stated on the characters $e^{iF(f)}$, with
real $f\in C_c^\infty(\{x_j>0\};\Lambda^2)$, extends to this form:
their linear span is dense in the positive-side $L^2$ space, and the
form is bounded by the $L^2$ norm.

\begin{proposition}\label{prop:axis-reconstruction}
There are a non-negative self-adjoint $H_j$ and a strongly continuous
unitary representation $V_j$ of $e_j^\perp$ on $\mathcal H_j$ such
that
\begin{equation}\label{eq:axis-semigroup}
 e^{-tH_j}V_j(y)[A]_j=[U(te_j+y)A]_j,
 \qquad t\ge0,\quad y\perp e_j.
\end{equation}
The two representations commute, and $H_j\Omega_j=0$.
\end{proposition}

\begin{proof}
Translations are strongly continuous on $L^2(\lambda)$. On
characters this follows from continuity of translation of a Schwartz
test and bounded convergence; their density gives it on all of
$L^2(\lambda)$. Positive time translations preserve the positive-side
algebra. With $|A|_b=b_j(A,A)^{1/2}$, reflection and translation
invariance give
\[
 |U(te_j)A|_b^2=b_j(A,U(2te_j)A)
 \le |A|_b\,|U(2te_j)A|_b.
\]
Iteration yields
\[
 |U(te_j)A|_b
 \le |A|_b^{1-2^{-k}}\|A\|_{L^2(\lambda)}^{2^{-k}}.
\]
Letting $k\to\infty$ proves that $[A]_j\mapsto[U(te_j)A]_j$
is a well-defined contraction. It is symmetric, obeys the semigroup
law, and is strongly continuous by the contraction of the quotient
map. Writing it as $T_j(t)$, one has
$T_j(t)=T_j(t/2)^2\ge0$. If $T_j(t)\psi=0$, repeated square roots
and strong continuity give $\psi=0$. Spectral calculus therefore
defines $H_j=-\log T_j(1)\ge0$; uniqueness of positive roots gives
$T_j(t)=e^{-tH_j}$ first at dyadic times and then at all times.

For $y\perp e_j$, $U(y)$ preserves the half-space and commutes with
reflection, so it is an isometry for $b_j$, with inverse $U(-y)$.
Its quotient extension is the unitary $V_j(y)$. Strong continuity
and commutation with $T_j(t)$ follow from the corresponding Euclidean
translation identities. The constant observable is fixed throughout.
\end{proof}

Write $V_j(se_i)=e^{is\mathsf P_i}$ for $i\ne j$. The joint spectral
measure $\mathsf E_j$ of $H_j$ and these commuting momenta lies on
$[0,\infty)\times e_j^\perp$. For $\psi\in\mathcal H_j$ let
$\varrho_\psi(A)=\langle\psi,\mathsf E_j(A)\psi\rangle$.
Then
\begin{equation}\label{eq:joint-spectral-integral}
 \langle\psi,e^{-tH_j}V_j(y)\psi\rangle
 =\int e^{-tE+ip\cdot y}\,\dd\varrho_\psi(E,p).
\end{equation}

\subsection{A diagonal mirror}

\begin{proposition}\label{prop:diagonal-cone}
Let $i\ne j$, and suppose in addition that $\lambda$ is invariant
and reflection positive at $x_j+x_i=0$. Then
\begin{equation}\label{eq:sharp-coordinate-cone}
 \mathsf E_j\{(E,p):E<|p_i|\}=0.
\end{equation}
Thus $H_j\ge|\mathsf P_i|$ in the joint spectral calculus.
\end{proposition}

\begin{proof}
Fix $D>0$ and $A\in L^2(\Sigma_B,\lambda)$, where
$B=\{0<x_j<D,\ |x_i|<D\}$, and set $\psi=[A]_j$.
Let $\varrho$ be the image of $\varrho_\psi$ under
$(E,p)\mapsto(E,p_i)$. This is a finite positive measure on the
right half-plane, identified by $z=E-ip_i$.
For $a,c\ge0$ consider
\[
 K(a,c)=\E_\lambda\bigl[(\Theta_j A)
 U(2De_j+a(e_j+e_i)+c(e_j-e_i))A\bigr].
\]
The axis reconstruction gives
\begin{equation}\label{eq:cone-first-picture}
 K(a,c)=\int e^{-2DE}e^{-az-c\overline z}\,\dd\varrho(E,p_i).
\end{equation}

There is a second representation with the diagonal as time direction.
Put $n_\pm=(e_j\pm e_i)/\sqrt2$, $y=D(e_j+e_i)/2$, and let
$r_+$ reflect in $x_j+x_i=0$. The observables
\[
 A_1=\mathcal R_{r_+}U(-y)\mathcal R_j A,
 \qquad A_2=U(2De_j-y)A
\]
are supported on the positive side of this diagonal. Indeed,
$r_jB-y$ lies in its negative side and $B+2De_j-y$ in its positive
side. Their $L^2$ norms equal $\|A\|_2$.
Writing $\Theta_+A=\overline{\mathcal R_{r_+}A}$, translation of
the original expectation by $-y$ gives
\[
 K(a,c)=\E_\lambda\bigl[(\Theta_+A_1)
 U(a(e_j+e_i)+c(e_j-e_i))A_2\bigr].
\]
Choose $Q\in O(3)$ carrying $e_j,e_i$ to $n_+,n_-$ and apply
Proposition~\ref{prop:axis-reconstruction} to the transformed law
$(Q^{-1})_*\lambda$. If $\mathsf E'$ is its joint spectral measure,
set $\psi_a=[\widehat{Q^{-1}}A_a]'$, $a=1,2$, and let
$\varrho'$ be the image of
$\langle\psi_1,\mathsf E'(\cdot)\psi_2\rangle$
under $(E,p)\mapsto(\sqrt2E,\sqrt2p_i)$. This gives
\begin{equation}\label{eq:cone-second-picture}
 K(a,c)=\int e^{-aE'+icp'}\,\dd\varrho'(E',p'),
 \qquad \|\varrho'\|_{\mathrm{TV}}\le\|A\|_2^2,
\end{equation}
where $E'\ge0$. The variation bound is spectral
Cauchy--Schwarz applied to the two quotient vectors of $A_1,A_2$.
The factors $\sqrt2$ from the diagonal coordinates have been absorbed
into $E',p'$.

For $h\in L^1(0,\infty)$ write
$\widetilde h(z)=\int_0^\infty h(a)e^{-az}\,\dd a$.
Integrating the two representations against
$h(a)\overline{h(c)}$ yields the decisive bound
\begin{equation}\label{eq:two-Laplace-bound}
 \int e^{-2DE}|\widetilde h(E-ip_i)|^2\,\dd\varrho
 \le\|A\|_2^2
 \sup_{u\ge0}|\widetilde h(u)|
 \sup_{v\in\R}|\widetilde h(iv)|.
\end{equation}
All integrations are justified by the finite variation of the
spectral measures and the $L^1$ assumption on $h$.

We choose transforms whose two axis suprema have bounded product,
but whose modulus grows exponentially outside $|\arg z|\le\pi/4$.
For $s=\pm1$ and $n\ge1$ use the principal-power functions
\begin{equation}\label{eq:cone-Laplace-tests}
 \widetilde h_n(z)
 =e^{-n}z^{1-4isn/\pi}(1+z)^{-3},\qquad
 \widetilde h_n(0)=0.
\end{equation}
Their suprema on the positive and imaginary axes are at most
$e^{-n}$ and $e^n$, respectively. Each is the Laplace transform of
an $L^1$ function. Here is the factorization that verifies this
point, including the boundary of the half-plane. For
$\nu=1-4isn/\pi$,
\[
 z^\nu(1+z)^{-3}
 =\bigl(1-(1+z)^{-1}\bigr)^\nu(1+z)^{-(3-\nu)}.
\]
The second factor is the transform of
$a^{2-\nu}e^{-a}/\Gamma(3-\nu)$, an integrable function.
The binomial coefficients of the first factor are absolutely
summable, since $|\binom\nu k|=O(k^{-2})$. It is therefore the
transform of the finite complex measure
\[
 \delta_0+
 \sum_{k\ge1}(-1)^k\binom\nu k
 \frac{a^{k-1}e^{-a}}{(k-1)!}\,\dd a.
\]
Convolution gives $h_n\in L^1$. Absolute convergence and
$\operatorname{Re}\nu=1$ also give the stated boundary value at
$z=0$.

Fix $\delta>0$ and the sector
$S_\delta=\{z\ne0:\operatorname{Re}z\ge0,
\ s\arg z\ge\pi/4+\delta\}$. On this sector,
\[
 |\widetilde h_n(z)|^2
 \ge e^{8\delta n/\pi}|z|^2|1+z|^{-6}.
\]
Substitution in \eqref{eq:two-Laplace-bound} gives
\[
 e^{8\delta n/\pi}
 \int_{S_\delta}e^{-2DE}|z|^2|1+z|^{-6}\,\dd\varrho
 \le\|A\|_2^2.
\]
The integrand is positive on $S_\delta$, so letting $n\to\infty$
shows $\varrho(S_\delta)=0$. Taking both signs and rational
$\delta>0$ gives $|\arg z|\le\pi/4$ on its support, or
$E\ge|p_i|$.

Finite linear combinations of characters with compact positive-side
tests lie in some such $B$ and have dense quotient classes in
$\mathcal H_j$. The spectral projection of $\{E<|p_i|\}$
vanishes on this dense set and hence vanishes identically.
\end{proof}

Applying this result to the nine-mirror laws gives
$H_j\ge|\mathsf P_i|$ for every ordered pair $i\ne j$. This coordinate
cone is the input for the analytic rotation argument and for the
identification of invariant projections used in clustering.

\subsection{Invariant observables and reflection}
\label{subsec:invariant-reflection}

Let $P_j$ be the orthogonal projection in $L^2(\lambda)$ onto the
observables fixed by all $U(te_j)$, and let $P$ project onto those
fixed by all translations. These are the corresponding conditional
expectations. The averages
$T^{-1}\int_0^TU(te_j)\,\dd t$ converge strongly to $P_j$.
The projections commute with one another and with every translation.

An observable fixed by translations along $e_j$ can be represented
in either half-space bounded by any plane $x_j=c$. Indeed, approximate
it in $L^2$ by finite sums of characters with compact tests, and translate
each approximant far enough into the chosen half-space. Translation
fixes the limiting observable and preserves the approximation error.
The half-space $L^2$ subspace is closed.

The linear geometric reflection $\mathcal R_j$ restricts to a positive
self-adjoint involution on this invariant subspace. Positivity follows
from its inclusion in the positive-side space and reflection
positivity. Its spectrum is therefore $\{1\}$, so
\begin{equation}\label{eq:invariant-reflection-fixed}
 \mathcal R_jA=A\qquad(A\in\operatorname{ran}P_j).
\end{equation}
Thus the quotient map is an isometry there. Its range is precisely
$\ker H_j$. To identify it, average
\[
 \langle[A]_j,e^{-tH_j}[B]_j\rangle
     =\langle\mathcal R_jA,U(te_j)B\rangle_{L^2(\lambda)}
\]
over $0\le t\le T$. Spectral dominated convergence on the left and
the mean ergodic theorem on the right give
\begin{equation}\label{eq:zero-energy-class}
 P_0[B]_j=[P_jB]_j,\qquad P_0=\mathbf1_{\{0\}}(H_j).
\end{equation}
The image of $\operatorname{ran}P_j$ is closed because the map is
isometric, and contains the dense subspace $P_0[\mathcal E_j^+]_j$
of $\ker H_j$. This proves the asserted identification.

For a law invariant and reflection positive at the nine mirrors,
the cone forces the joint spectral measure at zero energy to be
concentrated at zero momentum. Hence $V_j(y)$ fixes $\ker H_j$.
If $A\in\operatorname{ran}P_j$ and $y\perp e_j$, then
\[
 [U(y)A-A]_j=0.
\]
The difference is again in $\operatorname{ran}P_j$, so the isometry
gives $U(y)A=A$. Consequently
\begin{equation}\label{eq:invariant-projections-equal}
 P_1=P_2=P_3=P.
\end{equation}

\subsection{Correlations in all directions}

\begin{proposition}\label{prop:all-direction-mixing}
Suppose $\lambda$ is translation invariant and invariant and
reflection positive at the nine mirrors. For all $A,B\in L^2(\lambda)$,
\begin{equation}\label{eq:all-direction-mixing}
 \langle A,U(a)B\rangle\longrightarrow\langle A,PB\rangle
                      \quad\text{as }|a|\longrightarrow\infty.
\end{equation}
If $\lambda$ is translation ergodic, the limit is
$\overline{\E_\lambda A}\,\E_\lambda B$ and
$\ker H_j=\C\Omega_j$ for every $j$.
\end{proposition}

\begin{proof}
First fix $j$. For $A=\mathcal R_jA'$ and $A',B\in\mathcal E_j^+$,
\eqref{eq:zero-energy-class} gives, for $t\ge0$ and $y\perp e_j$,
\begin{align*}
 |\langle A,U(te_j+y)(1-P_j)B\rangle|
 &=|\langle[A']_j,V_j(y)e^{-tH_j}(1-P_0)[B]_j\rangle|\\
 &\le\|[A']_j\|\,\|e^{-tH_j}(1-P_0)[B]_j\|.
\end{align*}
The last norm tends to zero by spectral dominated convergence;
the bound is uniform in $y$. Approximate general $A,B$ by local
character sums. Translate their supports into opposite half-spaces,
at a cost of replacing $t$ by $t-2r$ for one fixed $r$.
The correlation error is bounded uniformly in $t,y$ by
$\|A-A_1\|_2\|B\|_2+\|A_1\|_2\|B-B_1\|_2$.
The adjoint identity gives the same conclusion as $t\to-\infty$.
Thus
\[
 \lim_{|t|\to\infty}\sup_{y\perp e_j}
 |\langle A,U(te_j+y)(1-P_j)B\rangle|=0.
\]
By \eqref{eq:invariant-projections-equal}, $P_j=P$ for all $j$.
For any sequence $|a|\to\infty$, a subsequence has one fixed
coordinate tending in modulus to infinity. The displayed uniform
bound applies with that coordinate, while $U(a)PB=PB$.
Every subsequence has such a further subsequence, proving
\eqref{eq:all-direction-mixing}. In an ergodic law $P$ is the
projection onto constants. The isometry onto $\ker H_j$ then gives
its one-dimensionality.
\end{proof}

\subsection{Passing the hypotheses to ergodic components}
\label{subsec:component-hypotheses}

The conditional laws in an ergodic decomposition require their own
reflection and Gaussian bounds. We prove these from invariant
weights; covariance of a mixture alone does not give covariance
of its individual components.

Let $\mathcal I$ be the translation-invariant event sigma algebra.
The space of tempered distributions, with the sigma algebra generated
by its evaluations, is standard Borel. It therefore admits regular
conditional laws $\lambda_F$ given $\mathcal I$, with
\begin{equation}\label{eq:component-kernel}
 \lambda_F(\Phi)=\E_\lambda[\Phi\mid\mathcal I](F),\qquad
 \lambda=\int\lambda_F\,\dd\lambda(F).
\end{equation}
The first equality is understood almost surely for each integrable
$\Phi$ (and for nonnegative $\Phi$ by monotone convergence).
The image measure on laws is the translation-ergodic decomposition.
We recall why a single full-measure set of ergodic laws can be used.

For an open set $O$, $\Sspace(\overline O;\Lambda^2)$ denotes
the subspace of $\Sspace(\R^3;\Lambda^2)$ consisting of forms
supported in $\overline O$, with the induced Schwartz topology.
Fix a countable real test space $\mathcal D_0\subset C_c^\infty(\R^3;\Lambda^2)$
over $\mathbb Q$, stable under rational translations and cubic
isometries, with $\mathcal D_0\cap C_c^\infty(O;\Lambda^2)$ dense
in $\Sspace(\overline O;\Lambda^2)$ for $\R^3$ and the six coordinate
half-spaces $O$. Take the rational span of translated cubic images
of countable Schwartz-dense compact-test families in these seven
open sets to obtain it.
Let $\mathcal A_0$ consist of finite sums of $e^{iF(f)}$,
$f\in\mathcal D_0$, with coefficients in $\mathbb Q+i\mathbb Q$.
This same countable family is dense in $L^2$ for every law on the
evaluation sigma algebra and determines such laws. For density,
condition an observable orthogonal to its span on finitely many
evaluations. Rational combinations of the tests show that the
Fourier transform of the resulting finite complex measure vanishes
on rational frequencies. Continuity and Fourier uniqueness make it
zero. Conditional expectations on the increasing finite evaluation
sigma algebras converge to the observable, proving density.
The same finite-dimensional Fourier argument proves determination
of laws. For the elements
$A_p$, choose cube-average times $T_k\to\infty$ so that
\[
 \|M_{T_k}A_p-PA_p\|_2^2\le2^{-k}\quad(p\le k),\qquad
 M_T=T^{-3}\int_{[0,T]^3}U(a)\,\dd a.
\]
Equation \eqref{eq:component-kernel} and summability give convergence
in $L^2(\lambda_F)$ for every $p$ on one full-measure set of $F$.
Moreover $PA_p$ has zero conditional variance, and thus equals
the constant $\lambda_F(A_p)$ in that space. The conditional laws
are translation invariant: test the parent invariance against
$\mathcal I$-measurable bounded weights, first for rational
translations and the countable characters, then use continuity and
Fourier uniqueness. The mean ergodic theorem for each such law now
shows that its invariant projection is the constant projection,
first on the dense algebra and then on all $L^2(\lambda_F)$.
Thus these conditional laws are ergodic. This also gives uniqueness
of their image measure: the same cube averages converge pointwise
along a subsequence and in $L^2$ under almost every law in any
ergodic decomposition, so they determine the same character
expectations and hence the same component law.

Fix an axis mirror. A bounded invariant weight $h$ satisfies
$\mathcal R_jh=h$ by \eqref{eq:invariant-reflection-fixed} and has a
representative in either half-space. If $h\ge0$ and
$A\in\mathcal E_j^+$, reflection positivity applied to $\sqrt hA$
gives
\begin{equation}\label{eq:weighted-component-rp}
 \E_\lambda[h(\Theta_jA)A]\ge0.
\end{equation}
Reflection invariance similarly gives
$\E_\lambda[h\mathcal R_jA]=\E_\lambda[hA]$.
Taking all bounded $\mathcal I$-measurable weights in these identities
proves that the corresponding conditional reflection form is
nonnegative and the conditional character expectations are reflection
invariant. Countably many positive-side character sums suffice;
density extends the statements to the full positive-side algebra
in each conditional law. Translation invariance moves the mirror
to every real height.

This argument also transfers any further fixed mirror normal $n$.
Choose $Q\in O(3)$ with $Qe_1=n$ and apply the axis argument to
the transported law $(Q^{-1})_*\lambda$. The transported law is
reflection positive at $x_1=0$; it need not equal $\lambda$.
The map $Q_*$ sends translation-ergodic laws to translation-ergodic
laws and transports their reflection forms. By uniqueness of the
ergodic decomposition, its image decomposition is that of $\lambda$.
Applying this to the six diagonal normals, and intersecting the
resulting nine full-measure sets, transfers all required mirrors.

Suppose now that the parent law satisfies
\eqref{eq:probability-G}. To transfer this Gaussian domination with
the same constant, fix real
$f\in\Sspace$ and put $G=e^{F(f)}$, $c=e^{\|f\|_2^2/2}$,
$\phi=PG$. All moments of $G$ are finite. For each integer $n\ge2$,
Tonelli's theorem and the parent bound give
\begin{equation}\label{eq:component-exponential-average}
 \E_\lambda[(M_TG)^n]\le T^{-3n}
 \int_{([0,T]^3)^n}
 \exp\left(\frac n2\|f\|_2^2+
       \frac12\sum_{p\ne q}\langle f,\tau_{a_p-a_q}f\rangle\right)\,\dd a.
\end{equation}
The overlaps tend to zero as $|a_p-a_q|\to\infty$ and are bounded
by $\|f\|_2^2$. Fix $\epsilon>0$. Outside the set where one pair is within distance
$\delta$, the right side is at most
$c^n e^{n(n-1)\epsilon/2}$ once $\delta$ is sufficiently large.
The exceptional set has relative volume at most
$\binom n2(2\delta)^3T^{-3}$ and its integrand is at most
$e^{n^2\|f\|_2^2/2}$.
Since $M_TG\to\phi$ in $L^2$, and the averages and $\phi$ are
uniformly bounded in $L^{2n}$, interpolation gives convergence in
$L^n$. Sending $T\to\infty$ and then $\epsilon\to0$ in
\eqref{eq:component-exponential-average} yields
$\|\phi\|_n\le c$ for every $n$, hence $\phi\le c$ almost surely.
By \eqref{eq:component-kernel},
\begin{equation}\label{eq:component-G}
 \lambda_F(e^{F(f)})\le e^{\|f\|_2^2/2}\quad\text{almost surely}.
\end{equation}
Choose a countable Schwartz-dense real test family and intersect
these full-measure sets. For any other test, pointwise convergence
of the evaluations under each conditional law and Fatou's lemma
extend \eqref{eq:component-G} to it. Thus all real tests obey the
bound on one full-measure set of component laws.

Finally a measurable set of full parent measure has full measure in
almost every component by \eqref{eq:component-kernel}. In particular
closedness passes to the components; its event is measurable by
testing $dF=0$ on a countable dense family of three-form tests.
The nine reflections generate the cubic group, so each of these
components is also cubically invariant. Proposition
\ref{prop:all-direction-mixing} applies to each component and gives
clustering and $\ker H_j=\C\Omega_j$. The rotation theorem below
will be applied to these same component laws.


\section{Rotational covariance}\label{sec:rotations}

The rotation argument starts with a field insertion separated from
the two observables in a reflection form. Gaussian domination bounds
the insertion in the reconstructed Hilbert space, independently of
the number of factors in either observable. The spectral cone then
allows complex translations. Together these bounds give bounded
operators for insertions concentrated at a point of the
$(x_1,x_2)$-plane. Closedness removes the extra angular frequencies
at two points. At higher orders, reflection positivity gives a positive
kernel for the highest angular coefficient. Polygon divisibility and a
convexity estimate then reduce every order to the two-point result.

Let $\mu$ be a translation-invariant, Gaussian-dominated law,
invariant and reflection positive at $x_1=0$.
We first work with this reconstruction, writing
$\mathcal H=\mathcal H_1$, $H=H_1$, and
$U_t=U(te_1)$.
Let $\mathcal A_+$ denote the positive-side measurable functions
belonging to every $L^p(\mu)$, $p<\infty$. It is an algebra.
Their classes are dense in $\mathcal H$, since $\mathcal A_+$
contains the bounded positive-side functions.

Use the Hodge isometry
$*f=(f_{23},-f_{13},f_{12})$ to write
$\mathbb B(b)=F(*^{-1}b)$, complex linearly. Proposition
\ref{prop:subgaussian} extends this map to vector-valued $L^2$ tests.
The action on these tests is
\[
 (g\cdot b)(x)=\det(g)\,g\,b(g^{-1}x)
 \quad(g\in O(3)),\qquad
 (\tau_a b)(x)=b(x-a).
\]
In particular
$\Theta_1\mathbb B(b)=\mathbb B(\overline{r_1\cdot b})$
and $\Theta_1U_t=U_{-t}\Theta_1$.
Write $L^2_{[a,b]}$ for the vector tests vanishing almost everywhere
outside $a\le x_1\le b$.

\subsection{The insertion estimate}\label{subsec:rotation-insertion}

For $\ell>0$ and $h\in L^2_{[0,\ell]}$, consider
$M_\ell(h)\Phi=e^{\mathbb B(h)}U_\ell\Phi$ for bounded
positive-side $\Phi$. This function lies in $\mathcal A_+$.
We first prove
\begin{equation}\label{eq:rotation-source-bound}
 \|[M_\ell(h)\Phi]_1\|
       \le e^{\|\Re h\|_2^2/2}\|[\Phi]_1\|.
\end{equation}

Set $h^\#=\tau_{\ell e_1}(\overline{r_1\cdot h}+h)$.
Reflection and translation invariance give
\begin{equation}\label{eq:rotation-source-doubling}
 \|[M_\ell(h)\Phi]_1\|^2
   =\langle[\Phi]_1,[M_{2\ell}(h^\#)\Phi]_1\rangle.
\end{equation}
The two summands in $h^\#$ have disjoint slab interiors, and
$\tau_{2\ell e_1}\overline{r_1\cdot h^\#}=h^\#$.

For a source $h'$ on $[0,\ell']$ with
$\tau_{\ell'e_1}\overline{r_1\cdot h'}=h'$, let
$h'^{(n)}=\sum_{j=0}^{n-1}\tau_{j\ell'e_1}h'$.
Doubling $h'^{(n)}$ gives $h'^{(2n)}$. Thus, with
$m_a=\|[M_{2^a\ell'}(h'^{(2^a)})\Phi]_1\|$,
\[
 m_a^2\le\|[\Phi]_1\|m_{a+1},\qquad
 m_0\le\|[\Phi]_1\|^{1-2^{-a}}m_a^{2^{-a}}.
\]
On the other hand, the integrand in the doubling identity is bounded
in modulus by
\[
 \|\Phi\|_\infty^2e^{\mathbb B(\Re h'^{(2^{a+1})})}.
\]
The real translates have disjoint interiors, so Gaussian domination
gives
\[
 m_a^2\le\|\Phi\|_\infty^2
       e^{\,2^a\|\Re h'\|_2^2}.
\]
Taking the indicated powers and sending $a\to\infty$ proves
\eqref{eq:rotation-source-bound} for symmetric sources, including
$\|[\Phi]_1\|=0$. Apply this to $h^\#$, whose real part has squared
norm $2\|\Re h\|_2^2$, in
\eqref{eq:rotation-source-doubling}. Cauchy--Schwarz proves the
bound for arbitrary $h$.

\begin{proposition}\label{prop:rotation-insertion}
For $\ell>0$ and $f\in L^2_{[0,\ell]}$ there is a unique bounded
operator $\mathcal B_\ell(f)$ on $\mathcal H$ satisfying
\[
 \mathcal B_\ell(f)[\Phi]_1
       =[\mathbb B(f)U_\ell\Phi]_1\quad(\Phi\in\mathcal A_+).
\]
It is complex linear in $f$, and
\begin{equation}\label{eq:rotation-insertion-bound}
 \|\mathcal B_\ell(f)\|\le k_1\|f\|_2,\qquad
 k_1=\sqrt{2\exp(1)}.
\end{equation}
For real $f$, the constant is $\exp(1/2)$.
If $f\in L^2_{[c,c']}$ with $0\le c<c'\le\ell$, then
\begin{equation}\label{eq:rotation-insertion-local}
 \mathcal B_\ell(f)=e^{-cH}
       \mathcal B_{c'-c}(\tau_{-ce_1}f)e^{-(\ell-c')H}.
\end{equation}
For $f\in L^2_{[-d,d]}$, set
$\mathcal B_d^{\mathrm c}(f)=\mathcal B_{2d}(\tau_{de_1}f)$.
Its matrix elements and adjoint are
\begin{align}
 \langle[\Phi']_1,\mathcal B_d^{\mathrm c}(f)[\Phi]_1\rangle
   &=\E_\mu[(\Theta_1U_d\Phi')\mathbb B(f)U_d\Phi],
                         \label{eq:rotation-insertion-centered}\\
 \mathcal B_d^{\mathrm c}(f)^*
   &=\mathcal B_d^{\mathrm c}(\overline{r_1\cdot f}).\notag
\end{align}
For $r\ge|t|$ and $a\in\R$, one also has
\begin{multline}\label{eq:rotation-insertion-shift}
 e^{-(r+t)H+ia\mathsf P_2}\mathcal B_d^{\mathrm c}(f)
           e^{-(r-t)H-ia\mathsf P_2}\\
       =\mathcal B_{d+r}^{\mathrm c}(\tau_{te_1+ae_2}f).
\end{multline}
\end{proposition}

\begin{proof}
For real $f$ and bounded $\Phi$, the map
$z\mapsto[e^{z\mathbb B(f)}U_\ell\Phi]_1$ is entire.
Indeed its exponential series converges locally uniformly in
$L^2(\mu)$, dominated on $|z|\le R$ by
$\|\Phi\|_\infty e^{R|\mathbb B(f)|}$.
The quotient map is a contraction. On $|z|=\|f\|_2^{-1}$,
\eqref{eq:rotation-source-bound} bounds its norm by
$\exp(1/2)\|[\Phi]_1\|$. Cauchy's coefficient estimate gives the
real insertion bound. Applying it to $\Re f$ and $\Im f$, and using
$\|\Re f\|_2+\|\Im f\|_2\le\sqrt2\|f\|_2$, gives
\eqref{eq:rotation-insertion-bound}.
The estimate sends null classes to null classes, so it defines the
operator on the dense classes of bounded functions and extends it
uniquely to $\mathcal H$.

For $\Phi\in\mathcal A_+$, its bounded truncations converge in $L^4$.
Since $\mathbb B(f)\in L^4$, their products with the inserted field
converge in $L^2$ by H\"older's inequality. This proves the asserted
formula on $\mathcal A_+$ and complex linearity in $f$.

Equation \eqref{eq:rotation-insertion-local} follows by applying
$U_c$ to
$\mathbb B(\tau_{-ce_1}f)U_{\ell-c}\Phi$ and using
\eqref{eq:axis-semigroup}. Translation by $-de_1$ in the reflection
form gives \eqref{eq:rotation-insertion-centered}. Applying reflection
and conjugation to that integrand gives the adjoint formula.
For \eqref{eq:rotation-insertion-shift}, apply the two exterior
semigroups to a class $[\Phi]_1$. The resulting function is
\[
 \mathbb B(\tau_{(d+r)e_1+te_1+ae_2}f)\,
       U_{2(d+r)}\Phi.
\]
The test $\tau_{te_1+ae_2}f$ is supported in $[-d-r,d+r]$,
since $|t|\le r$. The displayed function therefore defines its centered
insertion at $d+r$, which is the right side. Equality on these dense
classes proves the operator identity.
\end{proof}

In particular, for $\Phi_1,\Phi_2\in\mathcal A_+$ measurable on
$x_1>d$,
\begin{equation}\label{eq:rotation-separated-insertion}
 |\E_\mu[(\Theta_1\Phi_1)\mathbb B(f)\Phi_2]|
 \le k_1\|f\|_2\,\|[U_{-d}\Phi_1]_1\|
                         \|[U_{-d}\Phi_2]_1\|.
\end{equation}
This is \eqref{eq:rotation-insertion-centered} applied to the two
translated observables. The constant does not depend on their degrees.

\subsection{Complex translations and point plates}
\label{subsec:rotation-point-plates}

From now on assume also the sharp cone
$\mathsf E_1\{(E,p_\perp):E<|p_2|\}=0$.
For the laws under construction, Proposition
\ref{prop:diagonal-cone} supplies this from the diagonal mirror.

For $\xi=(\xi^0,\xi^1)\in\C^2$, put
\[
 d_{\mathcal W}(\xi)=\Re\xi^0-|\Im\xi^1|,\qquad
 \mathcal W=\{d_{\mathcal W}>0\},\qquad
 \xi^*=(\overline{\xi^0},-\overline{\xi^1}),\quad n_0=(1,0).
\]
The cone \eqref{eq:sharp-coordinate-cone} gives contractions
\[
 T(\xi)=\int e^{-\xi^0E+i\xi^1p_2}\,\dd\mathsf E_1(E,p_\perp)
                  \quad(\xi\in\overline{\mathcal W}).
\]
They obey $T(\xi)T(\xi')=T(\xi+\xi')$,
$T(\xi)^*=T(\xi^*)$, and $T(\xi)\Omega_1=\Omega_1$.
They are strongly continuous on the closed tube by spectral
dominated convergence. In its interior they are holomorphic in
operator norm: on $d_{\mathcal W}\ge\epsilon>0$ the first and
second derivatives of the multiplier are bounded by constants
times $E e^{-\epsilon E}$ and $E^2e^{-\epsilon E}$.
Taylor's formula therefore gives a uniform quadratic remainder.
For real $y=(y^0,y^1)$ with $y^0\ge0$,
$T(y)[\Phi]_1=[U(y^0e_1+y^1e_2)\Phi]_1$.

Write $x=(y,x_3)$ with $y=(x_1,x_2)$. Set
\begin{gather*}
 e^+=e_1+ie_2,\quad e^-=e_1-ie_2,\quad e^3=e_3,\\
 w(+)=1,\quad w(-)=-1,\quad w(3)=0,\\
 s(+)=s(-)=1,\quad s(3)=-1.
\end{gather*}
A label $\lambda=(\sigma,\chi)$ has
$\sigma\in\{+,-,3\}$ and real $\chi\in L^2(\R)$.
For $\alpha\in L^2(\R^2;\C)$ its plate is
$(\alpha\otimes\lambda)(x)=\alpha(y)\chi(x_3)e^\sigma$.
Thus $\|\alpha\otimes\lambda\|_2\le\sqrt2\|\alpha\|_2\|\chi\|_2$.
For rotations $R_t$ in the $(x_1,x_2)$-plane and the coordinate
reflection $r_1$, direct multiplication gives
\begin{align}
 R_t\cdot(\alpha\otimes\lambda)
   &=e^{-iw(\sigma)t}(\alpha\circ R_{-t})\otimes\lambda,
                              \label{eq:rotation-plate-actions}\\
 \overline{r_1\cdot(\alpha\otimes\lambda)}
   &=s(\sigma)(\overline\alpha\circ r_1)\otimes\lambda.\notag
\end{align}

\begin{proposition}\label{prop:rotation-point-plates}
For each label $\lambda$ there are bounded operators
$\mathcal B_\lambda(\xi,\xi')$, $\xi,\xi'\in\mathcal W$,
holomorphic in operator norm and real linear in $\chi$, such that
\begin{equation}\label{eq:rotation-point-bound}
 \|\mathcal B_\lambda(\xi,\xi')\|
 \le\frac{k_2\|\chi\|_2}
          {\min(d_{\mathcal W}(\xi),d_{\mathcal W}(\xi'))},
 \qquad k_2=\frac8\pi\exp(1/2).
\end{equation}
For $\kappa,\kappa'\in\overline{\mathcal W}$,
\begin{align}
 T(\kappa)\mathcal B_\lambda(\xi,\xi')T(\kappa')
   &=\mathcal B_\lambda(\kappa+\xi,\xi'+\kappa'),
                                   \label{eq:rotation-point-shift}\\
 \mathcal B_\lambda(\xi,\xi')^*
   &=s(\sigma)\mathcal B_\lambda(\xi'^*,\xi^*).\notag
\end{align}
If $\alpha\in L^2(\R^2)$ is compactly supported in $|y^0|<d$, then
\begin{equation}\label{eq:rotation-point-integral}
 \mathcal B_d^{\mathrm c}(\alpha\otimes\lambda)
   =\int\alpha(y)\mathcal B_\lambda(dn_0+y,dn_0-y)\,\dd y.
\end{equation}
The integral is a Bochner integral in operator norm.
When $\supp\alpha\subset\{|y^0|\le d'\}$, $d'<d$, its norm is at
most $k_2\|\chi\|_2\|\alpha\|_1/(d-d')$.
\end{proposition}

\begin{proof}
Let $j_\delta=(\pi\delta^2)^{-1}\mathbf1_{\{|y|<\delta\}}$.
Set $A_\delta=\mathcal B_\delta^{\mathrm c}(j_\delta\otimes\lambda)$
and, when both tube distances exceed $\delta$, set
\[
 \mathcal B_\lambda^\delta(\xi,\xi')
   =T(\xi-\delta n_0)A_\delta T(\xi'-\delta n_0).
\]
These operators obey the shift and adjoint identities in
\eqref{eq:rotation-point-shift}. Here $j_\delta$ is real and
reflection invariant, so $A_\delta^*=s(\sigma)A_\delta$.
The real translation identity \eqref{eq:rotation-insertion-shift}
also gives, for $r\ge|y^0|$,
\begin{equation}\label{eq:rotation-point-regularized-translation}
 \mathcal B_\lambda^\delta((r+\delta)n_0+y,(r+\delta)n_0-y)
    =\mathcal B_{\delta+r}^{\mathrm c}((\tau_yj_\delta)\otimes\lambda).
\end{equation}

We prove a bound uniform as $\delta\downarrow0$.
Fix $\xi,\xi'$ and set
$d_0=\min(d_{\mathcal W}(\xi),d_{\mathcal W}(\xi'))$,
$r=(d_0-\delta)/2>0$. For unit vectors $\psi,\psi'$,
\[
 h(z)=\langle\psi,\mathcal B_\lambda^\delta(\xi+z,\xi'-z)\psi'\rangle
\]
is holomorphic on $|\Re z^0|<r$, $|\Im z^1|<r$.
Write $z=y+i(t,t')$. Factoring the exterior translations gives
\[
 h(y+i(t,t'))=
 \langle\psi_1,\mathcal B_{\delta+r}^{\mathrm c}
             ((\tau_yj_\delta)\otimes\lambda)\psi'_1\rangle,
\]
where
$\psi_1=T(\xi-(\delta+r)n_0+i(t,t'))^*\psi$ and
$\psi'_1=T(\xi'-(\delta+r)n_0-i(t,t'))\psi'$.
For $|t'|<r$ these are contractions of the original vectors;
they are independent of $y$.
For $\eta\in L^2((-r,r)^2)$, extended by zero, the integral of
$\eta(y)\tau_yj_\delta$ is $\eta*j_\delta$ in $L^2$.
Its support lies in $|y^0|\le r+\delta$.
The insertion bound and Young's inequality therefore give
\[
 \left|\int\eta(y)h(y+i(t,t'))\,\dd y\right|
   \le2\exp(1/2)\|\eta\|_2\|\chi\|_2.
\]
Hence
$\int_{(-r,r)^2}|h(y+i(t,t'))|^2\,\dd y
\le4\exp(1)\|\chi\|_2^2$.
Apply the mean-value inequality on a polydisc of radius $r'<r$
and bound its integral by that over the four-dimensional box.
This yields
\[
 |h(0)|^2\le(\pi r'^2)^{-2}(2r)^2\,4\exp(1)\|\chi\|_2^2.
\]
Letting $r'\uparrow r$ and taking arbitrary unit vectors proves
\begin{equation}\label{eq:rotation-regularized-point-bound}
 \|\mathcal B_\lambda^\delta(\xi,\xi')\|
       \le\frac{k_2\|\chi\|_2}{d_0-\delta}.
\end{equation}

To identify the limit, put $r_0=d_0/4$.
For $\delta<r_0$ the functions
$h_\delta(z)=\langle\psi,\mathcal B_\lambda^\delta
(\xi+z,\xi'-z)\psi'\rangle$ form a bounded normal family on
$|z^0|,|z^1|<r_0$ by
\eqref{eq:rotation-regularized-point-bound}.
For $\eta\in L^2((-r_0,r_0)^2)$, the real translation identity
and an additional equal time translation on both sides give
\begin{multline}\label{eq:rotation-point-limit-test}
 \int\eta(y)h_\delta(y)\,\dd y\\
 =\langle\psi,T(\xi-2r_0n_0)
       \mathcal B_{2r_0}^{\mathrm c}((\eta*j_\delta)\otimes\lambda)
       T(\xi'-2r_0n_0)\psi'\rangle.
\end{multline}
Since $\eta*j_\delta\to\eta$ in $L^2$, this has a unique limit.
Any two holomorphic subsequential limits thus have the same
integrals on the real square; continuity identifies them there,
and the identity theorem identifies them on the polydisc.
Consequently every matrix element at $z=0$ converges.
The bound \eqref{eq:rotation-regularized-point-bound} defines a
bounded limiting sesquilinear form, represented by
$\mathcal B_\lambda(\xi,\xi')$ with
\eqref{eq:rotation-point-bound}. This is weak operator convergence.
The shift, adjoint, and real linearity identities pass to these
limits.

For holomorphy, fix $(\xi_0,\xi'_0)$ and a positive $\epsilon$
smaller than both tube distances. Near this pair the shift identity
factors the operator as
\[
 T(\xi-\xi_0+\epsilon n_0)\,
 \mathcal B_\lambda(\xi_0-\epsilon n_0,\xi'_0-\epsilon n_0)\,
 T(\xi'-\xi'_0+\epsilon n_0).
\]
The two exterior factors are holomorphic in operator norm.

Finally choose $d'<d$ containing the time support of $\alpha$.
Integrating
\eqref{eq:rotation-point-regularized-translation} with $r=d-\delta$
gives
\[
 \int\alpha(y)\mathcal B_\lambda^\delta(dn_0+y,dn_0-y)\,\dd y
       =\mathcal B_d^{\mathrm c}((\alpha*j_\delta)\otimes\lambda).
\]
The right side converges in norm to the left side of
\eqref{eq:rotation-point-integral}. For matrix elements on the
left, \eqref{eq:rotation-regularized-point-bound} gives the
integrable majorant
$2k_2\|\chi\|_2|\alpha(y)|/(d-d')$ for sufficiently small $\delta$.
Dominated convergence identifies the limiting matrix elements.
Holomorphy makes the limiting integrand norm continuous on the
support, and \eqref{eq:rotation-point-bound} bounds its norm.
It is therefore Bochner integrable, proving the formula and its
stated norm bound.
\end{proof}

\subsection{Ordered field kernels}\label{subsec:rotation-ordered-kernels}

The intermediate arguments of a product of point plates matter
only through their sum. Indeed, for $\kappa_1,\kappa_2\in\mathcal W$
and sufficiently small $\epsilon>0$, the shift identity gives
\[
 \mathcal B_\lambda(\xi,\kappa_1)
       \mathcal B_{\lambda'}(\kappa_2,\xi')
 =\mathcal B_\lambda(\xi,\epsilon n_0)
       T(\kappa_1+\kappa_2-2\epsilon n_0)
       \mathcal B_{\lambda'}(\epsilon n_0,\xi').
\]
The same expression results from any other decomposition of
$\kappa_1+\kappa_2$ into two points of $\mathcal W$.
An exterior argument adjacent to the vacuum can likewise be
changed arbitrarily in $\mathcal W$, since $T(\kappa)\Omega_1=
T(\kappa)^*\Omega_1=\Omega_1$.

For $N\ge2$ and $\xi_1,\ldots,\xi_{N-1}\in\mathcal W$, define
\begin{multline}\label{eq:rotation-ordered-kernel}
 G_N(\xi;\lambda)=
 \langle\Omega_1,\mathcal B_{\lambda_1}(\kappa,\xi_1/2)
       \left(\prod_{a=2}^{N-1}
          \mathcal B_{\lambda_a}(\xi_{a-1}/2,\xi_a/2)\right)\\
       \mathcal B_{\lambda_N}(\xi_{N-1}/2,\kappa)\Omega_1\rangle,
       \qquad\kappa\in\mathcal W.
\end{multline}
Set $G_1(\lambda)=\langle\Omega_1,
\mathcal B_\lambda(\kappa,\kappa)\Omega_1\rangle$.
Both definitions are independent of $\kappa$.
For $x_1,x_{a+1}-x_a\in\mathcal W$, define the vector
\begin{multline}\label{eq:rotation-ordered-vector}
 \Psi(x;\lambda)=
 \mathcal B_{\lambda_1}(x_1,(x_2-x_1)/2)
       \left(\prod_{a=2}^{n-1}
          \mathcal B_{\lambda_a}((x_a-x_{a-1})/2,(x_{a+1}-x_a)/2)\right)\\
       \mathcal B_{\lambda_n}((x_n-x_{n-1})/2,\kappa)\Omega_1.
\end{multline}
Products here and below are ordered with increasing index from
left to right; an empty product is the identity.
For $n=1$ this means
$\mathcal B_{\lambda_1}(x_1,\kappa)\Omega_1$; for $n=0$ it means
$\Omega_1$. It too is independent of $\kappa$.
Put $s(\lambda)=\prod_a s(\sigma_a)$, $w(\lambda)=\sum_a w(\sigma_a)$,
and let $\bar\lambda$ exchange the $+$ and $-$ labels.

\begin{proposition}\label{prop:rotation-ordered-kernels}
The kernels and vectors are holomorphic on their stated domains.
For $N\ge2$,
\begin{equation}\label{eq:rotation-ordered-kernel-bound}
 |G_N(\xi;\lambda)|
 \le\prod_{a=1}^N(2k_2\|\chi_a\|_2)
       \left(\min_a d_{\mathcal W}(\xi_a)\right)^{-N}.
\end{equation}
Let $\alpha_1,\ldots,\alpha_N\in L^2(\R^2)$ have compact
time-ordered supports:
$\sup_{\supp\alpha_a}y^0<\inf_{\supp\alpha_{a+1}}y^0$.
Then
\begin{equation}\label{eq:rotation-ordered-density}
 \E_\mu\prod_a\mathbb B(\alpha_a\otimes\lambda_a)
 =\int\prod_a\alpha_a(y_a)\,
       G_N(y_2-y_1,\ldots,y_N-y_{N-1};\lambda)\,\dd Y.
\end{equation}
If the supports are also in $y^0>0$, their product class is
the Bochner integral of $\prod_a\alpha_a(y_a)\Psi(Y;\lambda)$.
One has $G_1=0$,
$\overline{G_N(\xi;\lambda)}=G_N(\bar\xi;\bar\lambda)$, and
\begin{multline}\label{eq:rotation-ordered-Gram}
 \langle\Psi(x;\lambda),\Psi(z;\lambda')\rangle
 =s(\lambda)G_{n+n'}((x_n-x_{n-1})^*,\ldots,(x_2-x_1)^*,\\
       x_1^*+z_1,z_2-z_1,\ldots,z_{n'}-z_{n'-1};
       \lambda_n,\ldots,\lambda_1,\lambda'_1,\ldots,\lambda'_{n'})
\end{multline}
when $n,n'\ge1$. With a vacuum in the first slot the inner
product is $G_{n'}$ on the successive differences of $z$;
with a vacuum in the second slot it is $s(\lambda)G_n$
on the reversed starred differences and reversed labels.
Furthermore
$T(\kappa')\Psi(x;\lambda)=\Psi(x+\kappa';\lambda)$ for
$\kappa'\in\overline{\mathcal W}$, where every point is translated.
\end{proposition}

\begin{proof}
Holomorphy follows from that of the factors. In the definition of
$G_N$, choose $\kappa=tn_0$ with
$t\ge\frac12\min_a d_{\mathcal W}(\xi_a)$ and apply
\eqref{eq:rotation-point-bound} to every factor. This proves
\eqref{eq:rotation-ordered-kernel-bound}.

To prove the density formula, first put the supports in $y^0>0$.
Choose disjoint time slabs $[a_j,b_j]$, in order, with each support
strictly inside its slab, $a_1\ge0$, and $b_j\le a_{j+1}$.
Let $f_j=\alpha_j\otimes\lambda_j$,
$c_j=(a_j+b_j)/2$, $d_j=(b_j-a_j)/2$, and
$I_j=\mathcal B_{d_j}^{\mathrm c}(\tau_{-c_je_1}f_j)$.
The insertion identity and its localization give
\begin{equation}\label{eq:rotation-ordered-product-class}
 \left[\prod_j\mathbb B(f_j)\right]_1
 =T(a_1n_0)I_1T((a_2-b_1)n_0)I_2\cdots I_N\Omega_1.
\end{equation}
For precision, the backward induction uses
$\Phi'=\prod_{j>j_0}\mathbb B(\tau_{-b_{j_0}e_1}f_j)$,
which belongs to $\mathcal A_+$.
The tail product is $U_{b_{j_0}}\Phi'$, so insertion of $f_{j_0}$
is $\mathcal B_{b_{j_0}}(f_{j_0})[\Phi']_1$.
Localization supplies $T(a_{j_0}n_0)I_{j_0}$; the induction
on the translated tail supplies the next nonnegative gap
$a_{j_0+1}-b_{j_0}$ and the remaining factors.

Use \eqref{eq:rotation-point-integral} to write
\[
 I_j=\int\alpha_j(y_j)
       \mathcal B_{\lambda_j}(y_j-a_jn_0,b_jn_0-y_j)\,\dd y_j.
\]
On the compact supports all these operators are norm bounded;
the $\alpha_j$ are integrable. The product in
\eqref{eq:rotation-ordered-product-class} is therefore one Bochner
integral. Absorb each gap semigroup into a neighbouring point plate.
The inner arguments between plates $j$ and $j+1$ then add to
$y_{j+1}-y_j$. Replace them by equal halves using the product
identity above. The initial argument becomes $y_1$, and the
terminal argument is removed against the vacuum. The integral is
exactly the stated vector formula. Pairing with $\Omega_1$ proves
\eqref{eq:rotation-ordered-density}.
A common real time translation puts any time-ordered supports in
the positive half-plane, without changing their differences or
their moment, proving the general formula.
For $N=1$, this formula and the vanishing mean field give $G_1=0$.

For the Gram identity, take the adjoint of the product defining
$\Psi(x;\lambda)$. It reverses the order, stars the arguments,
and multiplies by $s(\lambda)$. In the middle of the two products
the arguments $x_1^*$ and $z_1$ add to $x_1^*+z_1$; replace them
by equal halves. This gives
\eqref{eq:rotation-ordered-Gram}, including the stated vacuum cases.
The translation identity follows by shifting the first argument.

Conjugating \eqref{eq:rotation-ordered-density} replaces the plates
by $\bar\alpha_a\otimes\bar\lambda_a$. Testing in small ordered
product neighbourhoods identifies
$\overline{G_N(\xi;\lambda)}=G_N(\xi;\bar\lambda)$ on real
$\xi\in\mathcal W^{N-1}$. Both sides, with $\xi$ conjugated on
the first, are holomorphic, and this domain is connected.
The identity theorem proves the claimed conjugation formula.
\end{proof}

\subsection{The two-point angular estimate}
\label{subsec:rotation-two-point-estimate}

The operator bound for $N$ point plates gives angular growth of
order $N$. At two points a positive spectral measure gives the
precise growth order two. Closedness will remove the unwanted
endpoint frequencies.

Fix a coordinate reconstruction $j$. For scalar
$u\in L^1\cap L^2(\R^3)$ supported in $x_j\ge0$, put
\[
 \widetilde u(E,p_\perp)=
       \int u(x)e^{-x_jE+ip_\perp\cdot x_\perp}\,\dd x.
\]
There are complex Radon measures $\varrho^{(j)}_{ab}$,
$a,b\in\{1,2,3\}$, with
\begin{equation}\label{eq:rotation-one-field-spectrum}
 \langle[\mathbb B_a(u)]_j,[\mathbb B_b(u')]_j\rangle
       =\int\overline{\widetilde u}\,\widetilde u'\,
                                    \dd\varrho^{(j)}_{ab}.
\end{equation}
They have a common positive Radon measure $\varrho_0$ such that
$\varrho^{(j)}_{ab}=r_{ab}\varrho_0$ and the matrix $r$ is
positive semidefinite and Hermitian almost everywhere.
In particular, for vector coefficients $c,c'$,
the measures $\varrho_{cc'}=\sum_{a,b}\bar c_ac'_b\varrho_{ab}$
satisfy
\begin{equation}\label{eq:rotation-spectral-CS}
 \int|G_1G_2|\,\dd|\varrho_{cc'}|
 \le\left(\int|G_1|^2\,\dd\varrho_{cc}\right)^{1/2}
       \left(\int|G_2|^2\,\dd\varrho_{c'c'}\right)^{1/2}.
\end{equation}
All these measures are carried by the joint spectral cone.

Here is the construction, including the fact that the measures
need only be locally finite. Convolution and
\eqref{eq:axis-semigroup} give
\[
 [\mathbb B_a(u*\alpha)]_j
       =\widetilde u(H_j,\mathsf P)[\mathbb B_a(\alpha)]_j
\]
for $\alpha\in L^2$ supported in $x_j\ge0$.
This follows by applying the bounded field and quotient maps to
the $L^2$ Bochner integral
$u*\alpha=\int u(x)\tau_x\alpha\,\dd x$;
Fubini applies to its scalar spectral matrix elements since
$u\in L^1$.
Let $P_1(v)=(\pi(1+v^2))^{-1}$, whose Fourier transform with
either sign is $e^{-|p|}$, and choose
\[
 \alpha_0(x)=\mathbf1_{(0,1)}(x_j)\prod_{i\ne j}P_1(x_i),
 \qquad
 a_0(E,p_\perp)=\widetilde\alpha_0
       =\frac{1-e^{-E}}{E}e^{-|p_\perp|_1}>0.
\]
The quotient at $E=0$ is one. With
$\psi_a=[\mathbb B_a(\alpha_0)]_j$ and
$m_{ab}=\langle\psi_a,\mathsf E_j(\cdot)\psi_b\rangle$,
commutativity of convolution gives
\begin{equation}\label{eq:rotation-spectral-reference}
 a_0(H_j,\mathsf P)[\mathbb B_a(u)]_j
       =\widetilde u(H_j,\mathsf P)\psi_a.
\end{equation}
The finite matrix measure $m$ is positive. If $m_0=\sum_a m_{aa}$,
spectral Cauchy--Schwarz shows $m_{ab}\ll m_0$.
Its density matrix is positive Hermitian on one common full-measure
set: test its quadratic form on the countable rational complex
vectors and then use continuity. Define
$\varrho_{ab}=a_0^{-2}m_{ab}$ and $\varrho_0=a_0^{-2}m_0$.
The strictly positive continuous $a_0$ has a positive lower bound
on bounded sets, so these measures are Radon.
Pointwise matrix Cauchy--Schwarz followed by its integral version
gives \eqref{eq:rotation-spectral-CS}.

On $A_n=\{a_0>1/n\}$, which is bounded, one may divide
\eqref{eq:rotation-spectral-reference} by $a_0$ using bounded
spectral calculus. This proves \eqref{eq:rotation-one-field-spectrum}
with the measures restricted to $A_n$.
As $A_n$ increase to the whole spectral space,
the diagonal identity follows by monotone convergence; its integral
is $\|[\mathbb B_a(u)]_j\|^2\le\|u\|_2^2$.
Equation \eqref{eq:rotation-spectral-CS} then makes the mixed
integral absolutely convergent and gives the unrestricted identity
by dominated convergence. Cone support follows from that of
$\mathsf E_j$.

Now use $j=1$, and take real profiles
$\chi,\chi'\in C_c^\infty(\R)$. Their transverse transforms in this
spectral convention are
$\widehat\chi_+(p_3)=\int\chi(x_3)e^{ip_3x_3}\,\dd x_3$.
For labels $\lambda=(\sigma,\chi)$ and
$\lambda'=(\sigma',\chi')$, set
\[
 \varrho_{\lambda\lambda'}=
 \sum_{a,b}\overline{e^\sigma_a}e^{\sigma'}_b
       \overline{\widehat\chi_+}\widehat{\chi'}_+\,\varrho^{(1)}_{ab}.
\]
In this formula $\widehat{\chi'}_+$ means the same positive-sign
transform of $\chi'$. These are Radon measures on $(E,p_2,p_3)$,
and $\varrho_{\lambda\lambda}\ge0$.
For $\alpha,\alpha'\in L^1\cap L^2(\R^2;\C)$ vanishing almost
everywhere on $y^0<0$, put
$\alpha^L(E,p_2)=\int\alpha(y)e^{-y^0E+iy^1p_2}\,\dd y$.
Equation \eqref{eq:rotation-one-field-spectrum} gives
\begin{equation}\label{eq:rotation-plate-spectrum}
 \langle[\mathbb B(\alpha\otimes\lambda)]_1,
          [\mathbb B(\alpha'\otimes\lambda')]_1\rangle
 =\int\overline{\alpha^L}(\alpha')^L\,
                                     \dd\varrho_{\lambda\lambda'}.
\end{equation}

The crucial estimate is
\begin{equation}\label{eq:rotation-plate-Laplace-bound}
 I_\lambda(t):=\int e^{-tE}\,\dd\varrho_{\lambda\lambda}
       \le\frac{36\|\chi\|_2^2}{\pi t^2}\quad(t>0).
\end{equation}
Indeed, for $c>0$ take
$\alpha_c(y)=c^{-2}\mathbf1_{(c,2c)}(y^0)P_1(y^1/c)$.
On $E\ge|p_2|$ its transform is bounded below by
$e^{-2cE-c|p_2|}\ge e^{-3cE}$.
Since $\|\alpha_c\|_2^2=(2\pi c^2)^{-1}$, the variance bound and
the quotient contraction in \eqref{eq:rotation-plate-spectrum} give
$I_\lambda(6c)\le\|\chi\|_2^2/(\pi c^2)$.
Choosing $c=t/6$ proves the estimate. Its mixed version is
\begin{equation}\label{eq:rotation-plate-mixed-Laplace}
 \int e^{-tE}\,\dd|\varrho_{\lambda\lambda'}|
       \le\sqrt{I_\lambda(t)I_{\lambda'}(t)}.
\end{equation}

Define, for $y^0>0$,
\[
 G^{[1]}_{\lambda\lambda'}(y)
       =s(\sigma)\int e^{-y^0E+iy^1p_2}\,
                                   \dd\varrho_{\lambda\lambda'}.
\]
It is smooth, since on $y^0\ge t>0$ every derivative of order $q$
is bounded in modulus by
$E^qe^{-tE}\le q!(2/t)^qe^{-tE/2}$ on the cone.
It computes the two-point function on time-ordered compact plane
supports. To see this, translate the supports to opposite sides
of $y^0=0$ and put $\alpha^\theta=\bar\alpha\circ r_1$.
Then
\[
 \E_\mu[\mathbb B(\alpha\otimes\lambda)
                 \mathbb B(\alpha'\otimes\lambda')]
 =s(\sigma)\int\overline{(\alpha^\theta)^L}(\alpha')^L\,
                                    \dd\varrho_{\lambda\lambda'}.
\]
Expanding the transforms gives the integral against
$G^{[1]}_{\lambda\lambda'}(y'-y)$. Fubini is justified by
\[
 \|\alpha\|_1\|\alpha'\|_1e^{-\delta E},
\]
where $\delta>0$ is the time gap, and
\eqref{eq:rotation-plate-mixed-Laplace}.
Thus this kernel agrees with $G_2$ on real positive-time differences.

For $r>0$ and $n_\zeta=(\cos\zeta,\sin\zeta)$, it has the
holomorphic continuation
\[
 h^{[1]}_{\lambda\lambda'}(\zeta)
   =s(\sigma)\int
       e^{-rE\cos\zeta+irp_2\sin\zeta}\,
                                \dd\varrho_{\lambda\lambda'}
       \quad(|\Re\zeta|<\pi/2).
\]
If $\zeta=t+i\beta$, the real part of the exponent is
$-r\cos t(E\cosh\beta+p_2\sinh\beta)$.
The cone bounds this above by
$-r\cos t\,e^{-|\beta|}E$.
On every compact substrip this gives an integrable exponential
majorant; dominated convergence and Morera's theorem give
holomorphy. Equations
\eqref{eq:rotation-plate-Laplace-bound}--\eqref{eq:rotation-plate-mixed-Laplace}
give
\begin{equation}\label{eq:rotation-two-point-strip}
 |h^{[1]}_{\lambda\lambda'}(t+i\beta)|
 \le\frac{36\|\chi\|_2\|\chi'\|_2}
          {\pi(r\cos t)^2}\,e^{2|\beta|}.
\end{equation}

\subsection{Closedness removes the unwanted two-point frequencies}
\label{subsec:rotation-two-point-frequencies}

Assume in addition that $\mu$ is closed and cubically invariant.
In particular, its two-point form is invariant under
$R_{\pi/2}$. We use only this quarter-turn consequence here.

For $\upsilon\in(\pi/2)\Z$ and $n_\upsilon\cdot y>0$, define
\[
 G^{(\upsilon)}_{\lambda\lambda'}(y)
   =e^{iw\upsilon}
       G^{[1]}_{\lambda\lambda'}(R_{-\upsilon}y),
       \qquad w=w(\sigma)+w(\sigma').
\]
By quarter-turn invariance and
\eqref{eq:rotation-plate-actions}, this is the two-point density
for supports ordered along $n_\upsilon$.
Two such densities agree on overlaps: test them on small product
neighbourhoods ordered in both directions and use continuity.
The four half-planes cover $\R^2\setminus\{0\}$, so they define
one smooth kernel $G_{\lambda\lambda'}$ there.
It computes the two-point function for any two disjoint compact
plane supports. Indeed, partition them into sufficiently small
half-open squares. Each pair of squares is ordered along one of
the four coordinate directions, and there are only finitely many
pairs. Sum the ordered density identities.

For fixed $r>0$, the functions
\[
 h^{(\upsilon)}(\zeta)=e^{iw\upsilon}
       h^{[1]}_{\lambda\lambda'}(\zeta-\upsilon)
       \quad(|\Re\zeta-\upsilon|<\pi/2)
\]
agree on adjacent strips, by the identity theorem applied on
their real overlap. They give an entire function $h$ with
\begin{equation}\label{eq:rotation-two-point-entire-angle}
 h(\vartheta)=G_{\lambda\lambda'}(rn_\vartheta),\qquad
 h(\zeta+\pi/2)=i^w h(\zeta),\qquad
 |h(t+i\beta)|\le C_{\chi,\chi'}r^{-2}e^{2|\beta|},
\end{equation}
where $C_{\chi,\chi'}=72\|\chi\|_2\|\chi'\|_2/\pi$.
For the bound, choose a strip centre within $\pi/4$ of $t$ and
apply \eqref{eq:rotation-two-point-strip}.

The entire function $e^{-iw\zeta}h(\zeta)$ has period $\pi/2$,
so it is $H(e^{4i\zeta})$ for a holomorphic $H$ on
$\C\setminus\{0\}$. The Laurent series of $H$ gives
frequencies $k=w+4m$, $m\in\Z$. Cauchy's formula on the circle
parametrized by $e^{4i(t+i\beta)}$, $0\le t\le\pi/2$,
bounds the coefficient $a^k(r)$ by
$C_{\chi,\chi'}r^{-2}e^{2|\beta|+k\beta}$.
Sending $\beta$ to the appropriate infinity kills $|k|>2$.
Consequently
\begin{equation}\label{eq:rotation-two-point-window}
 G_{\lambda\lambda'}(rn_\vartheta)
   =\sum_{\substack{|k|\le2\\k-w\in4\Z}}a^k(r)e^{ik\vartheta},
 \qquad |a^k(r)|\le C_{\chi,\chi'}r^{-2}.
\end{equation}
These coefficients are smooth in $r$ and real bilinear in the
profiles, by their Fourier integral formula.

Every pair of labels has only its expected frequency $k=w$,
except for the pairs $(+,+)$ and $(-,-)$, whose windows contain
both $2$ and $-2$. To eliminate the opposite frequency, use
closedness in its vector form $\mathbb B(\nabla u)=0$.
Take $\eta\in C_c^\infty(\R^2;\C)$ and
$\chi_0\in C_c^\infty(\R;\R)$. With $z=y^0+iy^1$ and
$\partial_z=(\partial_{y^0}-i\partial_{y^1})/2$,
\[
 \nabla(\eta(y)\chi_0(x_3))
   =(\partial_z\eta)\otimes(+,\chi_0)
       +(\partial_{\bar z}\eta)\otimes(-,\chi_0)
       +\eta\otimes(3,\chi'_0).
\]
Pair this gradient with a second plate whose plane support is
disjoint. Apply the density identities and integrate by parts.
Testing on small disjoint product neighbourhoods gives
\begin{equation}\label{eq:rotation-two-point-Bianchi}
 \partial_zG_{(+,\chi_0)\lambda'}
       +\partial_{\bar z}G_{(-,\chi_0)\lambda'}
       +G_{(3,\chi'_0)\lambda'}=0
                       \quad\text{on }\R^2\setminus\{0\}.
\end{equation}
There is no diagonal contribution in this equation.
In polar coordinates,
\begin{align*}
 \partial_z(f(r)e^{ik\vartheta})
       &=\tfrac12e^{i(k-1)\vartheta}(f'(r)+kf(r)/r),\\
 \partial_{\bar z}(f(r)e^{ik\vartheta})
       &=\tfrac12e^{i(k+1)\vartheta}(f'(r)-kf(r)/r).
\end{align*}
For $\lambda'=(+,\chi')$, the frequency $-3$ in
\eqref{eq:rotation-two-point-Bianchi} occurs only in the derivative
of $a^{-2}_{(+,\chi_0)(+,\chi')}$.
It obeys $f'=2f/r$.
For $\lambda'=(-,\chi')$, the frequency $3$ likewise gives
$f'=2f/r$ for $a^2_{(-,\chi_0)(-,\chi')}$.
Each solution is $f(r)=f(1)r^2$. The bound
$|f(r)|\le C_{\chi_0,\chi'}r^{-2}$ for all $r>0$ forces
$f(1)=0$ by sending $r\to\infty$. Thus
\begin{equation}\label{eq:rotation-two-point-phase}
 G_{\lambda\lambda'}(R_ty)
       =e^{i(w(\sigma)+w(\sigma'))t}G_{\lambda\lambda'}(y).
\end{equation}
The plate action cancels this phase in the two-point density
formula, proving rotational covariance for disjoint compact
plane supports.

\subsection{Extending identities across coincident plane supports}
\label{subsec:rotation-contact}

The preceding density argument does not yet address coincident
plane supports. We use the following $L^2$ argument,
which will also be applied at every higher order.

\begin{proposition}\label{prop:rotation-no-contact}
Let $U\subset\R^2$ be open, and write $L^2_U$ for the vector
fields supported in $U\times\R$. Let $D:(L^2_U)^N\to\C$ be a
bounded multilinear form. Suppose it vanishes when the entries
are supported over pairwise disjoint compact subsets of $U$,
and whenever one entry is $\nabla u$ with
$u\in C_c^\infty(U\times\R)$. Then $D=0$.
\end{proposition}

\begin{proof}
We argue simultaneously for all open $U$ by induction on $N$.
At $N=1$, fields with compact plane support are dense, so the
first assumption suffices.

At $N=2$, represent the bilinear form as
$D(f_1,f_2)=\int f_1(x)\cdot(Tf_2)(x)\,\dd x$ for a bounded
operator $T$ on $L^2_U$. The pairing here is bilinear.
For disjoint compact $K,K'\subset U$ the first assumption gives
$\mathbf1_KT\mathbf1_{K'}=0$.
Inner regularity and strong convergence of the corresponding
multiplication projections extend this to disjoint Borel sets.
Using $A$ and $U\setminus A$ on both sides gives
$T\mathbf1_A=\mathbf1_AT$ for every Borel $A$.
Uniform approximation by simple functions then gives $Th=hT$
for every bounded plane multiplier $h$. Hence
\[
 D(h\nabla u,f_2)=D(\nabla u,hf_2)=0.
\]

Fix a compact $A\subset U$, choose $u_0\in C_c^\infty(U)$ equal
to one near $A$, and let $\chi_0\in C_c^\infty(\R)$.
On $A\times\R$, the gradients of
$u_0\chi_0$, $u_0y^0\chi_0$, and $u_0y^1\chi_0$ are respectively
\[
 \chi_0'e_3,\qquad
 \chi_0e_1+y^0\chi_0'e_3,\qquad
 \chi_0e_2+y^1\chi_0'e_3.
\]
The preceding identity with multiplier $h\mathbf1_A$ first
annihilates $h\mathbf1_A\chi'_0e_3$.
Applying it with $hy^0\mathbf1_A$ and $hy^1\mathbf1_A$ then
annihilates $h\mathbf1_A\chi_0e_1$ and
$h\mathbf1_A\chi_0e_2$. These multipliers are bounded since $A$
is compact. Such plane-transverse products span a dense subspace.
For the transverse derivative family, its orthogonal complement
in $L^2(\R)$ has distributional derivative zero, hence is a
constant and therefore zero. Thus these three vector families
are total in $L^2_U$, and $D=0$.

Suppose $N\ge3$ and the assertion is proved at every smaller order.
Two consequences of the induction hypothesis separate unequal
plane cells. First, if one entry is supported over a compact
$K\subset U$ and the others over $U\setminus K$, fix that entry.
The form on the remaining $N-1$ slots satisfies both assumptions
on the open set $U\setminus K$ and is zero.
Second, if the slots divide into nonempty groups $I,J$, with
the $J$ entries supported over a compact $K_J$ and the $I$
entries over its complement, fix the $J$ entries.
The form on the $I$ slots inherits the gradient assumption.
When those slots have pairwise disjoint compact supports, the
first consequence, applied to any one of them, gives zero.
The induction hypothesis on $U\setminus K_J$ therefore makes
this form zero as well.

By density put all entries over one compact $K_0\subset U$.
Partition the plane into half-open squares $Q$ of side $d$ and
expand $D$ by multilinearity. A tuple of squares which are not
all equal has zero contribution. To justify compact separation,
replace each square by its concentric closed square of side
$d-2\epsilon$. Separate the slots whose square equals the first
from the remaining slots, use the second consequence with
$K_J=K_0$ intersected with the union of their closed squares,
and send $\epsilon\downarrow0$ in $L^2$.
Only the equal-square terms remain, and Cauchy--Schwarz gives
\[
 |D(f_1,\ldots,f_N)|
 \le\|D\|\max_Q\|\mathbf1_Qf_3\|_2
       \|f_1\|_2\|f_2\|_2\prod_{i=4}^N\|f_i\|_2.
\]
Indeed, sum the first two square norms by Cauchy--Schwarz and
bound the others by their whole norms. The maximum tends to
zero as $d\downarrow0$: its square is the integral over an
area-$d^2$ cell of the fixed integrable function
$y\mapsto\int|f_3(y,x_3)|^2\,\dd x_3$.
Absolute continuity of the integral proves the assertion.
\end{proof}

Apply this proposition with $N=2$ to
\[
 D(b,b')=\E_\mu[\mathbb B(R_tb)\mathbb B(R_tb')]
              -\E_\mu[\mathbb B(b)\mathbb B(b')],
\]
where $R_tb$ means $R_t\cdot b$. Gaussian domination makes it a
bounded bilinear form. Finite sums of plates with smooth real
transverse profiles are dense in the fields over any fixed
compact plane set: simple plane-transverse tensors are dense
in $L^2$, and $e^+,e^-,e^3$ span $\C^3$.
Thus \eqref{eq:rotation-two-point-phase} supplies the disjoint-support
assumption for general fields. Closedness supplies the gradient
assumption, since rotations carry gradients to gradients.
Consequently the two-point function is invariant under every
$R_t$, on all $L^2$ tests.

\subsection{The angular expansion at every order}
\label{subsec:rotation-angular-expansion}

Fix $N\ge2$ and a tuple of labels $\lambda$, with total weight
$w=w(\lambda)$. Put $\mathcal W_\upsilon=R_\upsilon\mathcal W$
for $\upsilon\in(\pi/2)\Z$. For a permutation $\mathfrak s$,
let
\[
 \mathcal Z(\mathfrak s,\upsilon)=
 \{Z=(z_1,\ldots,z_N):
       z_{\mathfrak s(a+1)}-z_{\mathfrak s(a)}
                    \in\mathcal W_\upsilon,\ 1\le a<N\}.
\]
On this open convex set define
\[
 G_\upsilon^{\mathfrak s}(Z;\lambda)=e^{iw\upsilon}
 G_N\left(
       (R_{-\upsilon}(z_{\mathfrak s(a+1)}-z_{\mathfrak s(a)}))_{a=1}^{N-1};
       \mathfrak s\lambda\right).
\]
Here $(\mathfrak s\lambda)_a=\lambda_{\mathfrak s(a)}$.
Moment symmetry, quarter-turn invariance, and
\eqref{eq:rotation-ordered-density} identify this as the moment
density on compact product supports in its real part.
Two such holomorphic functions agree on their intersection.
Indeed, the intersection is convex and is preserved by taking
real parts, so a nonempty intersection contains an open set of
real configurations. Testing both densities there identifies
them by continuity, and the holomorphic identity theorem
identifies them on the full intersection.
They therefore define one holomorphic function
$\mathcal G_N(Z;\lambda)$ on
$\mathcal Z_N=\bigcup_{\mathfrak s,\upsilon}\mathcal Z(\mathfrak s,\upsilon)$.

The real set $\mathcal O_N=\mathcal Z_N\cap(\R^2)^N$ consists
of configurations whose first coordinates are all distinct
or whose second coordinates are all distinct.
Its complement is the finite union of the codimension-two
subspaces
$\{y_a^0=y_b^0,\ y_k^1=y_l^1\}$.
Thus $\mathcal O_N$ is open, dense, and connected.
The functions $\mathcal G_N$ are symmetric under simultaneous
permutation of points and labels, invariant under common
translations, and satisfy
\[
 \mathcal G_N(R_{\pi/2}Z;\lambda)=i^w\mathcal G_N(Z;\lambda),
 \qquad
 \overline{\mathcal G_N(Z;\lambda)}
       =\mathcal G_N(\bar Z;\bar\lambda).
\]
These identities follow from the defining densities and
the corresponding identities for $G_N$.

For $\beta\in\R$ write $\Lambda_\beta=R_{i\beta}$.
In the coordinates $\xi^\pm=\xi^0\pm i\xi^1$,
$\mathcal W=\{\Re\xi^+>0,\Re\xi^->0\}$ and
$\Lambda_\beta\xi^\pm=e^{\mp\beta}\xi^\pm$.
Thus every tube $\mathcal Z(\mathfrak s,\upsilon)$ is preserved
by $\Lambda_\beta$. Moreover
$d_{\mathcal W}(\Lambda_\beta\xi)\ge
e^{-|\beta|}d_{\mathcal W}(\xi)$ on the closed tube.

\begin{proposition}\label{prop:rotation-finite-angular-expansion}
There are holomorphic coefficients $c_m(Z;\lambda)$ on
$\mathcal Z_N$, vanishing unless $|w+4m|\le N$, such that
\begin{equation}\label{eq:rotation-finite-angular-expansion}
 \mathcal G_N(\Lambda_\beta Z;\lambda)
       =\sum_m c_m(Z;\lambda)e^{-(w+4m)\beta}
                         \quad(\beta\in\R).
\end{equation}
They are unique, symmetric under simultaneous permutation,
and invariant under common translations. They satisfy
\begin{align}
 c_m(\Lambda_\beta Z;\lambda)
       &=e^{-(w+4m)\beta}c_m(Z;\lambda),
                           \label{eq:rotation-angular-coefficient-identities}\\
 c_m(R_{\pi/2}Z;\lambda)&=i^wc_m(Z;\lambda),\qquad
 \sum_m c_m(Z;\lambda)=\mathcal G_N(Z;\lambda),\notag\\
 \overline{c_m(X;\lambda)}&=c_{-m}(X;\bar\lambda)
                           \quad(X\in\mathcal O_N).\notag
\end{align}
\end{proposition}

\begin{proof}
Fix $0<\vartheta_0<\pi/4$ and first consider real configurations
whose consecutive differences have positive lengths and angles
$|\vartheta_a|<\vartheta_0$ from $n_0$.
For $\zeta=t+i\beta$,
\[
 \Re(R_{\zeta-\upsilon}\xi_a)^\pm
       =|\xi_a|e^{\mp\beta}\cos(\vartheta_a+t-\upsilon).
\]
The strips $|t-\upsilon|<\pi/2-\vartheta_0$ cover the complex
$\zeta$-plane. Thus $h_X(\zeta)=\mathcal G_N(R_\zeta X;\lambda)$
is entire. Choosing $\upsilon$ within $\pi/4$ of $t$ bounds every
gap distance below by
$r_Xe^{-|\beta|}\cos(\vartheta_0+\pi/4)$, where
$r_X=\min_a|\xi_a|$.
Equation \eqref{eq:rotation-ordered-kernel-bound} gives
$|h_X(t+i\beta)|\le C_{X,\lambda}e^{N|\beta|}$.
Its quarter-period is $h_X(\zeta+\pi/2)=i^wh_X(\zeta)$.
The Laurent coefficient argument used at two points, with
growth exponent $N$, gives the finite frequencies
$k=w+4m$, $|k|\le N$.

To extend this statement to a whole tube, let
$\mathcal M=\{m:|w+4m|\le N\}$ and
$P(u)=\prod_{m\in\mathcal M}(u+w+4m)$.
The functions
$Z\mapsto P(\partial_\beta)G_\upsilon^{\mathfrak s}
(\Lambda_\beta Z;\lambda)$ are holomorphic for each real $\beta$.
They vanish on the nonempty open real subset obtained by rotating
and permuting the configurations just used.
The identity theorem gives zero on the entire convex tube.
The resulting constant-coefficient ordinary differential equation
has distinct characteristic roots $-(w+4m)$.
Its coefficients in the exponential expansion are fixed linear
combinations of its derivatives at $\beta=0$, and hence are
holomorphic in $Z$.
On overlapping tubes the exponentials are linearly independent,
so the coefficients agree and define $c_m$ on $\mathcal Z_N$.
This also proves uniqueness. Comparing coefficients after a
permutation, common translation, or addition of the parameter
$\beta$ gives the stated identities. For the conjugation identity
use $\overline{\Lambda_\beta X}=\Lambda_{-\beta}X$ and
$w(\bar\lambda)=-w$.
\end{proof}

The coefficients have exact real rotation phases whenever both
configurations are in $\mathcal O_N$:
\begin{equation}\label{eq:rotation-angular-coefficient-phase}
 c_m(R_tX;\lambda)=e^{i(w+4m)t}c_m(X;\lambda).
\end{equation}
Locally in one tube, the allowed angles form an interval:
each ordered difference has a positive projection precisely
when $\cos(\vartheta_a+t)>0$.
The identity theorem in $t+i\beta$ extends
\eqref{eq:rotation-finite-angular-expansion} from the imaginary
axis to that interval. Comparing the exponential coefficients
gives \eqref{eq:rotation-angular-coefficient-phase} locally.

To pass through changes of ordering, first take distinct points
with no two distinct chords perpendicular. Such configurations
are a dense open set, the complement of finitely many nonzero
polynomial zero sets. Every rotation of one of them belongs to
$\mathcal O_N$: simultaneous coincidences of first and second
coordinates would give perpendicular chords, or identical points.
The local phase identity therefore continues along its whole
rotation orbit. Approximation proves it whenever $X,R_tX$
are in $\mathcal O_N$.

For an arbitrary configuration $X$ of distinct plane points,
choose $t$ so that all first coordinates of $R_tX$ are distinct.
Only finitely many angles modulo $\pi$ are excluded.
Define
\[
 \widetilde c_m(X;\lambda)
       =e^{-i(w+4m)t}c_m(R_tX;\lambda).
\]
The phase identity makes this independent of the choice.
Locally one may keep $t$ fixed, so it is real analytic on the
connected space $(\R^2)^N_{\ne}$ of distinct configurations.
It inherits permutation symmetry, translation invariance, and
\[
 \widetilde c_m(R_tX;\lambda)
       =e^{i(w+4m)t}\widetilde c_m(X;\lambda).
\]
Put $\widetilde G_N=\sum_m\widetilde c_m$.
It is real analytic and agrees with $\mathcal G_N$ on $\mathcal O_N$.

Define $\mathfrak a(0)=\mathfrak a(1)=0$, and, for $N\ge2$,
let $\mathfrak a(N)$ be the largest $|m|$ for which $c_m$ is
not identically zero for some tuple of $N$ labels; set it to zero
if all coefficients vanish. Since $|w|\le N$ and $|w+4m|\le N$,
\begin{equation}\label{eq:rotation-angular-linear-bound}
 \mathfrak a(N)\le N/2.
\end{equation}
A nonzero coefficient can be detected on a strictly time-ordered
real configuration after a quarter turn and permutation:
the identity theorem on each tube and the coefficient identities
above give this reduction. Conjugation then shows that if
$\mathfrak a(N)>0$, both indices
$\mathfrak a(N)$ and $-\mathfrak a(N)$ occur.

\subsection{Reflection bounds for the angular coefficients}
\label{subsec:rotation-angular-reflections}

We next prove the convexity estimate
\begin{equation}\label{eq:rotation-angular-convexity}
 2\mathfrak a(a+b)\le\mathfrak a(2a)+\mathfrak a(2b)
                          \quad(a,b\ge0).
\end{equation}
Let $N=a+b$ and $M=\mathfrak a(N)>0$.
Choose a strictly time-ordered $X=(y_1,\ldots,y_N)$ and labels
with $c_{-M}(X;\lambda)\ne0$. Translate its time cut so that the
first $a$ points have negative times and the remaining $b$
have positive times. Reflect the first group and reverse its
order to obtain $x$, with reversed labels $\lambda^x$;
call the second group $z$, with labels $\lambda^z$.
Their weights satisfy $w=w_x+w_z$.
The Gram identity gives
\[
 \langle\Psi(\Lambda_\beta x;\lambda^x),
             \Psi(\Lambda_\beta z;\lambda^z)\rangle
       =s(\lambda^x)\mathcal G_N(\Lambda_\beta X;\lambda).
\]
An empty group gives the vacuum.
The squared norm of the first vector is, up to the sign
$s(\lambda^x)$, the kernel of its reflected double:
the configuration $(r_1x_a,\ldots,r_1x_1,x_1,\ldots,x_a)$
with the same labels in reverse then forward order.
Its weight is $2w_x$. The other squared norm has weight $2w_z$.
These are direct instances of
\eqref{eq:rotation-ordered-Gram}; the reflection-conjugation
keeps the labels.

Apply Cauchy--Schwarz to this matrix element and expand the
three kernels by \eqref{eq:rotation-finite-angular-expansion}.
Multiplication by $e^{2w\beta}$ cancels their weights and gives
\[
 \left|\sum_m c_m(X;\lambda)e^{-4m\beta}\right|^2
 \le
 \left|\sum_{|m|\le\mathfrak a(2a)}d_me^{-4m\beta}\right|
 \left|\sum_{|m|\le\mathfrak a(2b)}d'_me^{-4m\beta}\right|.
\]
The coefficients $d_m,d'_m$ belong to the two specified doubles.
A vacuum factor is one. As $\beta\to+\infty$, the left side
has a nonzero leading term of rate $8M$, and the right side is
bounded by a constant times
$e^{4(\mathfrak a(2a)+\mathfrak a(2b))\beta}$.
This proves \eqref{eq:rotation-angular-convexity}.
The cases $M=0$ or $N\le1$ already satisfy it.

There is also a positive kernel for the highest coefficient.
For $a\ge1$, set $M=\mathfrak a(2a)$. On labelled configurations
of $a$ distinct points in $y^0>0$, define
\begin{equation}\label{eq:rotation-top-coefficient-kernel}
 K_a((x;\lambda),(z;\lambda'))=
       s(\lambda)\widetilde c_{-M}^{(2a)}
                        (r_1x,z;\lambda,\lambda').
\end{equation}
The superscript indicates the number of points.
This kernel is Hermitian and positive semidefinite.
First suppose the times in each configuration are distinct and
order them increasingly. By the Gram formula,
\[
 \langle\Psi(\Lambda_\beta x;\lambda),
                 \Psi(\Lambda_\beta z;\lambda')\rangle
 =s(\lambda)\sum_{|m|\le M}
       \widetilde c_m^{(2a)}(r_1x,z;\lambda,\lambda')
       e^{-(w(\lambda)+w(\lambda')+4m)\beta}.
\]
Reversing the reflected points is harmless by labelled symmetry.
Multiply this Gram kernel by the positive factors
$e^{w(\lambda)\beta}e^{w(\lambda')\beta}e^{-4M\beta}$.
Its pointwise limit as $\beta\to+\infty$ is $K_a$.
Positive semidefiniteness passes to this limit on every finite
matrix. Approximating configurations with tied times by ones
with distinct times proves the same assertion for all the
stated configurations.

The exact phase of $\widetilde c_{-M}$ transports this consequence
to any line $L$ in the plane. If $x,z$ lie on one open side of
$L$ and $r_L$ is its reflection, translation to a point of $L$
and rotation of its normal to $n_0$ express
$\widetilde c_{-M}(r_Lx,z)$ as $K_a$ times a unit-modulus factor.
Taking absolute values of the positive-kernel
Cauchy--Schwarz inequality gives
\begin{multline}\label{eq:rotation-top-coefficient-CS}
 |\widetilde c_{-M}^{(2a)}(r_Lx,z;\lambda,\lambda')|^2\\
 \le
 |\widetilde c_{-M}^{(2a)}(r_Lx,x;\lambda,\lambda)|
 |\widetilde c_{-M}^{(2a)}(r_Lz,z;\lambda',\lambda')|.
\end{multline}
In particular a zero reflected diagonal forces every such
off-diagonal coefficient to vanish.

\subsection{Regular polygons and the reduction to two points}
\label{subsec:rotation-polygons}

We prove the arithmetic restriction
\begin{equation}\label{eq:rotation-polygon-divisibility}
 a\mid 2\mathfrak a(2a)\quad(a\ge1).
\end{equation}
Fix $a$ and put $M=\mathfrak a(2a)$,
$c_*=\widetilde c_{-M}^{(2a)}$.
The coefficients depend only on the labelled set of points.
If a rotation through $t$ permutes such a set without changing
its labels, its exact phase gives
$(1-e^{i(w-4M)t})c_*=0$.
For a regular $2a$-gon with the same label $(\sigma,\chi)$ at
each vertex, use $t=\pi/a$ and $w=2aw(\sigma)$.
The label weight cancels, leaving $e^{-4\pi iM/a}$.
Thus $c_*$ vanishes on this polygon whenever $a$ fails to divide $2M$.

That selection rule concerns a special configuration. The
positive kernel propagates its zero to arbitrary configurations
by the following halving argument. For a labelled set $S$ of
at most $2a$ distinct points, let $\mathcal N(S)$ mean that
$c_*$ vanishes on every $2a$-point labelled configuration
containing $S$.
Let $S'$ have $j'\le a$ points on one open side $L_+$ of a line
$L$, and suppose $S\subset r_LS'\cup S'$.
Then
\begin{equation}\label{eq:rotation-halving}
 \mathcal N(S)\ \Longrightarrow\
       \mathcal N(S')\ \text{and}\ \mathcal N(r_LS').
\end{equation}
To see this, add $a-j'$ arbitrary points $R$ in $L_+$ to $S'$,
forming $x=S'\cup R$. The reflected double $r_Lx\cup x$
contains $S$, so its diagonal coefficient is zero.
Equation \eqref{eq:rotation-top-coefficient-CS} implies
$c_*(r_Lx\cup z)=0$ for every $a$-point labelled set $z$ in $L_+$.
Thus, with $r_LS'$ fixed, the remaining $a-j'$ points on
$L_-$ and $a$ points on $L_+$ may vary in an open set.

For each choice of their labels, the space of their distinct
positions, avoiding the fixed points, is the complement of
finitely many codimension-two affine subspaces and is connected.
The coefficient is real analytic there. Its zero on the
nonempty open set just obtained extends to all positions.
Every $2a$-point configuration containing $r_LS'$ is obtained
this way. This proves $\mathcal N(r_LS')$.
Interchange the sides of $L$ and repeat the argument to obtain
$\mathcal N(S')$.

Suppose now $a\nmid2M$. Then $M>0$ and $a\ge3$.
Fix an arbitrary label and any regular $2a$-gon with vertices
$P_k$, $k\in\Z/(2a)\Z$, in cyclic order.
Let $\mathcal N_j$ assert $\mathcal N(S)$ for every arc of $j$
consecutive vertices, with this label, in every such polygon.
Selection gives $\mathcal N_{2a}$.
For $j\ge2$, put $j'=\lceil j/2\rceil\le a$ and take an arc
$S'=\{P_{k_0+1},\ldots,P_{k_0+j'}\}$.
Reflect in the perpendicular bisector of
$[P_{k_0},P_{k_0+1}]$.
The $a$ vertices $P_{k_0+1},\ldots,P_{k_0+a}$ lie on one
open side: their angles from the edge-midpoint radius are
$(2k-1)\pi/(2a)\in(0,\pi)$.
Their images satisfy $r_LP_{k_0+k}=P_{k_0+1-k}$.
Consequently $r_LS'\cup S'$ is an arc of $2j'$ consecutive
vertices and contains a $j$-vertex arc.
Equation \eqref{eq:rotation-halving} gives
$\mathcal N_j\Rightarrow\mathcal N_{j'}$.

Iteration reaches $j=1$. Any plane point is a vertex of some
regular $2a$-gon, so the coefficient vanishes whenever a
configuration has even one point with the fixed label.
Since the label was arbitrary, $c_*$ vanishes identically for
all label tuples. This contradicts the occurrence of the
index $-M$, proving \eqref{eq:rotation-polygon-divisibility}.

Now set $\nu(a)=\mathfrak a(2a)$. It is a nonnegative integer
sequence, with $\nu(0)=0$, $\nu(a)\le a$, and
$2\nu(a)\le\nu(a-1)+\nu(a+1)$ by
\eqref{eq:rotation-angular-convexity}.
Its increments $d_a=\nu(a)-\nu(a-1)$ are nondecreasing integers.
The linear bound keeps them bounded above, so they are eventually
constant, say $d$. Thus $\nu(a)=da-c$ for large $a$, where
$c=\sum_k(d-d_k)\ge0$ is an integer.
Divisibility \eqref{eq:rotation-polygon-divisibility} makes every
large $a$ divide $2c$, forcing $c=0$.
All increments equal $d=\nu(1)$, and
\[
 \mathfrak a(2a)=a\mathfrak a(2),\qquad
 2\mathfrak a(N)\le N\mathfrak a(2).
\]
For the second inequality split $N=\lfloor N/2\rfloor+
\lceil N/2\rceil$ in \eqref{eq:rotation-angular-convexity}.

The two-point result gives $\mathfrak a(2)=0$.
For smooth transverse profiles this follows from
\eqref{eq:rotation-two-point-phase} and the ordered density formula:
near any strictly ordered two-point configuration, only the
frequency $w$ occurs. Linear independence of the finite angular
exponentials makes $c_m=0$ there for $m\ne0$, and the identity
theorem on each tube gives zero throughout.
For arbitrary $L^2$ profiles, the all-$L^2$ two-point covariance
proved above gives the same local density identity.
Hence
\begin{equation}\label{eq:rotation-no-extra-angular-coefficients}
 \mathfrak a(N)=0\quad\text{for every }N,\qquad
 \widetilde G_N(R_tX;\lambda)=e^{iw(\lambda)t}
                                  \widetilde G_N(X;\lambda).
\end{equation}

\subsection{Moment densities on all distinct configurations}
\label{subsec:rotation-all-distinct-density}

The analytic continuation $\widetilde G_N$ must still be
identified as the moment density at distinct configurations
outside $\mathcal O_N$. This is a separate step; continuation
of a density does not by itself exclude an additional singular
part on that smaller set.

We first record the analytic displacement used for this step.
If polynomial field observables $\Phi'$ and $\Phi''$ have test
supports in $x_1<c-\delta$ and $x_1>c+\delta$, respectively,
then
\[
 s\longmapsto\E_\mu[\Phi'U(se_2)\Phi'']
\]
extends holomorphically to $|\Im s|<2\delta$.
Translate $c$ to zero and put
$\psi'=[U_{-\delta}\Theta_1\Phi']_1$,
$\psi''=[U_{-\delta}\Phi'']_1$.
The quotient and semigroup identities express the real
correlation as
$\langle\psi',T((2\delta,s))\psi''\rangle$.
This is holomorphic on the stated strip by the cone.

For compact, pairwise disjoint plane supports we claim
\begin{equation}\label{eq:rotation-distinct-density}
 \E_\mu\prod_a\mathbb B(\alpha_a\otimes\lambda_a)
 =\int\prod_a\alpha_a(y_a)\widetilde G_N(Y;\lambda)\,\dd Y.
\end{equation}
The integral is absolutely convergent: the product support
is a compact subset of $(\R^2)^N_{\ne}$, the kernel is bounded
there, and each $\alpha_a\in L^2$ is integrable on its compact
support.
Call the difference of the two sides $D(\alpha_1,\ldots,\alpha_N)$
on these tuples.

We show local vanishing near each distinct configuration $X$
by induction on the number of equal-second-coordinate pairs.
When there are none, $X\in\mathcal O_N$ and the ordered tube
density gives the identity in small product discs.
Otherwise take a pair with equal second coordinates and
different times. Choose a time cut between them, different
from all point times, and split the points into groups $I,J$
below and above the cut.
Move the $J$ group by $se_2$. The points remain distinct for
every real $s$, since the groups are separated in time.
A sufficiently small positive $s_0$ removes every
second-coordinate coincidence across the cut and creates none:
take it smaller than every nonzero initial difference of
second coordinates.
There are fewer coincident pairs at $X(s_0)$.

The induction hypothesis gives a radius $2r$ near $X(s_0)$
on which $D$ vanishes. Reduce $r$ so that the original discs
of radius $r$ are disjoint and preserve the time gap.
Fix tests in these original discs, and translate those in
group $J$. The resulting difference $F(s)$ is real analytic
on $\R$.
For its moment term this follows from the separated-cluster
strip argument above. For its density term, change variables
to integrate
$\widetilde G_N(Y+sv_J;\lambda)$ against the fixed product
of tests, where $(v_J)_a=e_2$ for $a\in J$ and zero otherwise.
Over any bounded real $s$-interval these configurations form
a compact set of distinct points. Finite local holomorphic
extensions of the real analytic kernel give a common radius
$\epsilon>0$ and a bound $C$ there. Cauchy's estimates bound
the integrated derivative of order $q$ by
$Cq!\epsilon^{-q}\prod_a\|\alpha_a\|_1$.
Taylor's formula therefore gives real analyticity of this
integral.

For $|s-s_0|<r$ the translated tests remain in the discs of
radius $2r$ at $X(s_0)$, so $F(s)=0$.
The real analytic identity theorem gives $F(0)=0$.
This proves local vanishing near $X$ and completes the induction.
For general disjoint compact supports, finitely many of these
product-disc neighbourhoods cover their product.
Partition each support into sufficiently small half-open
squares. Every product of nonzero pieces lies in one of the
larger vanishing neighbourhoods: choose a point of the
product support in it and make each square diameter smaller
than the available half-radius.
Multilinearity and finite summation prove
\eqref{eq:rotation-distinct-density}.

\newpage
\subsection{Euclidean covariance and all reflection planes}
\label{subsec:rotation-conclusion}

\begin{theorem}\label{thm:rotation-covariance}
Let $\mu$ be a Gaussian-dominated law of closed tempered two-forms.
Suppose it is invariant under translations and the cubic group,
and reflection positive at the nine planes
$x_j=0$ and $x_j\pm x_i=0$, $i<j$.
Then it is invariant under all Euclidean motions and reflection
positive at every plane, on the full $L^2$ space of either open side.
\end{theorem}

\begin{proof}
Proposition \ref{prop:diagonal-cone} supplies the sharp coordinate
cones. For the plane $(1,2)$, fix $t$ and define the bounded
$N$-linear form
\[
 D(b_1,\ldots,b_N)=
 \E_\mu\prod_a\mathbb B(R_t\cdot b_a)
       -\E_\mu\prod_a\mathbb B(b_a).
\]
Its boundedness follows from Proposition \ref{prop:subgaussian}.
For plates over pairwise disjoint compact plane supports,
\eqref{eq:rotation-distinct-density} and
\eqref{eq:rotation-no-extra-angular-coefficients} prove $D=0$:
the factor $e^{-iw(\lambda)t}$ from the plate action cancels
the factor $e^{iw(\lambda)t}$ from the kernel.
Density of finite sums of plates gives the same assertion for
arbitrary $L^2$ fields over those compact sets.
The form also vanishes on gradients, by closedness and the
fact that rotations carry gradients to gradients.
Proposition \ref{prop:rotation-no-contact} therefore gives $D=0$
on all $L^2$ entries, for every $N$.

The entire generating function of each finite real Schwartz
test family is determined by these moments. Thus its
characteristic function is unchanged by $R_t$, and Fourier
uniqueness on the evaluation sigma algebra gives
$(R_t)_*\mu=\mu$.
Permuting coordinates applies the same argument in the other
coordinate planes: the nine mirrors give their cones, and
cubic invariance gives their quarter turns.
These plane rotations generate $SO(3)$.
A coordinate reflection already preserves $\mu$, so the full
$O(3)$ acts invariantly. Translation invariance gives the
Euclidean group.

Let $g$ carry $x_1=0$ to any chosen plane and its positive side
to the chosen open side. The composition operator of $g$ is
unitary on $L^2(\mu)$, maps these side spaces onto each other,
and conjugates the axis reflection to the chosen reflection.
Transporting its nonnegative reflection form proves the final
assertion.
\end{proof}

Theorem \ref{thm:rotation-covariance} applies separately to the
component laws of Section \ref{subsec:component-hypotheses}.
They retain Gaussian domination, closedness, the cubic group,
and the nine mirror forms. Together with their all-direction
clustering, it supplies the hypotheses for the tensor OS
verification.


\section{Chamber propagators}\label{sec:propagators}

The block integration requires bounds for the chamber Dirac operator which
are uniform in the mesh, the volume, and the running mass. There are two
sources of coercivity. The internal chamber graph controls the part of a
spinor whose mean over a cluster vanishes. The averaging penalty controls
the remaining four components. The first of these bounds also identifies
the singular part of the free propagator. This permits a local regularity
estimate with constants independent of the mesh, which is needed when
the inverses on small tori are combined into an inverse on the full torus.

\subsection{The cluster decomposition}\label{subsec:cluster-decomposition}

Let $a>0$ be a mesh. A fermion site is written $x=(c,g)$, where
$c\in a\Z^3$ and $g\in O_{\mathrm{cub}}$. In these coordinates the four
neighbours are
\[
 x^{(i)}=(c,gs_i)\quad(0\le i\le2),\qquad
 x^{(3)}=(c+age_1,gr_{e_1}).
\]
Here $s_0$ changes the sign of the third coordinate, $s_1$ interchanges
the second and third coordinates, $s_2$ interchanges the first and second
coordinates. The matrix $r_{e_1}$ changes the sign of the first
coordinate; it is the linear part of the affine wall reflection
$s_3$ used in the regulator. The
inward normals of the four walls are
\[
 n_0=e_3,\qquad n_1=\frac{e_2-e_3}{\sqrt2},\qquad
 n_2=\frac{e_1-e_2}{\sqrt2},\qquad n_3=-e_1.
\]
Thus $\widehat d_i(c,g)=-gn_i$ and
\[
 M_i(c,g)=\frac{w_i}{2}(\rho-\gamma\cdot\widehat d_i(c,g)),
 \qquad
 (w_0,w_1,w_2,w_3)=(6,6\sqrt2,6\sqrt2,6).
\]
Throughout this section $0<\rho<1$ is fixed.

The discussion applies on the covering lattice and on a torus of side
$R\in a\N$, with any unitary character
$f(c+Re_\mu,g)=-\exp(i\theta_\mu)f(c,g)$. Pairings on the torus are
formed over a period cell. The notation $d(c,c')$ denotes the sup
distance modulo $R\Z^3$, or the covering sup distance when $R=\infty$.
They use the measure
\[
 \int\dd x\,v(x)=\frac{a^3}{48}\sum_c\sum_gv(c,g),
 \qquad
 \|f\|_2^2=\int\dd x\,|f(x)|_2^2.
\]
On four-component spinors depending only on $c$, the measure is
$a^3\sum_c$. The cluster injection and its adjoint are
\[
 (J\xi)(c,g)=\xi(c),\qquad
 (J^*f)(c)=\frac1{48}\sum_gf(c,g).
\]
Set $P=JJ^*$ and $Q=1-P$. The decomposition $f=J\xi+h$, with
$\xi=J^*f$ and $h=Qf$, is orthogonal. In particular,
\begin{equation}\label{eq:cluster-orthogonal}
 \|f\|_2^2=\|\xi\|_2^2+\|h\|_2^2.
\end{equation}
The normalization by $48$ makes $J$ an isometry. Kernels on the chamber
lattice are taken with respect to $\dd x$, so the kernel of the identity
is $48a^{-3}\mathbf1[x=x']1_4$.

Let $A$ be a connection and write
$u_A(c,c')=\exp(ie_a aA(c,c'))$ on a photon bond. The charge $e_a$ in
this notation is the charge at the level under consideration. The
phases on the internal bonds are one. Put
\[
 R_i(A)f(x)=U_A(x,x^{(i)})f(x^{(i)}),\qquad
 \mathcal L_A=\sum_{i=0}^3w_i(1-R_i(A)).
\]
The massless chamber operator is
\begin{equation}\label{eq:propagator-dirac}
 D_Af(x)=-a^{-1}\sum_{i=0}^3M_i(x)
       \big(U_A(x,x^{(i)})f(x^{(i)})-f(x)\big).
\end{equation}
At level $k$, $a=\eta_k$ and dilation to the unit lattice gives
\begin{equation}\label{eq:propagator-phase-dictionary}
 D_A=\eta_k^{-1}D^{\mathrm{ch}}[B_k],\qquad
 B_k(c/\eta_k,c'/\eta_k)=e_k\eta_k A(c,c').
\end{equation}
The reference mass is added separately as $D_A+\mu_k$.
For real $A$, the involutions $R_i(A)$ are unitary and self-adjoint,
and the bond expansion gives
\begin{equation}\label{eq:dirac-realpart}
 \operatorname{Re}(f,D_Af)=\frac{\rho}{2a}(f,\mathcal L_Af),
 \qquad D_A^*=\gamma_5D_A\gamma_5,
 \qquad \gamma_5=\begin{pmatrix}0&1_2\\1_2&0\end{pmatrix}.
\end{equation}
Here $\gamma_5$ is Hermitian, $\gamma_5^2=1$, and it anticommutes with
each $\gamma_\mu$. Each unoriented bond contributes
$a^3w_i|U_A(x,x^{(i)})f(x^{(i)})-f(x)|_2^2/48$ to the energy.

On one cluster let
\[
 \mathcal L_{\mathrm{int}}=\sum_{i=0}^2w_i(1-R_i),
 \qquad R_iv(g)=v(gs_i).
\]
Its kernel consists of the constants: the three reflections
$s_0,s_1,s_2$ generate $O_{\mathrm{cub}}$. Denote its least positive
eigenvalue by $\lambda_{\mathrm{int}}$ and set
\[
 g_{\mathrm h}=\frac{\rho\lambda_{\mathrm{int}}}{2},\qquad
 c_* =\frac{\min(w_3,\lambda_{\mathrm{int}})}{18},\qquad
 c_{\mathrm{lh}}=\sqrt3(1+\rho).
\]
No numerical evaluation of $\lambda_{\mathrm{int}}$ is needed.

\begin{proposition}\label{prop:cluster-energy}
For real $A$ and $\xi=J^*f$,
\begin{equation}\label{eq:cluster-energy}
 (f,\mathcal L_Af)\ge\lambda_{\mathrm{int}}\|Qf\|_2^2,
 \qquad
 a^{-2}(f,\mathcal L_Af)\ge c_*\|\nabla_A\xi\|_2^2,
\end{equation}
where
$\nabla_{A,\mu}\xi(c)=a^{-1}(u_A(c,c+ae_\mu)\xi(c+ae_\mu)-\xi(c))$.
The same inequalities hold on a finite set $E$ of photon sites, with the
energy restricted to bonds whose endpoints lie over $E$ and the photon
differences restricted to bonds in $E$. Moreover,
\begin{equation}\label{eq:offdiagonal-bound}
 \|QD_AJ\xi\|_2\le c_{\mathrm{lh}}\|\nabla_A\xi\|_2,
 \qquad
 |(\xi,J^*D_AQf)|\le
 c_{\mathrm{lh}}\|\nabla_A\xi\|_2\|f\|_2.
\end{equation}
On $\operatorname{ran}Q$, for $t>-g_{\mathrm h}/a$,
\begin{equation}\label{eq:heavy-inverse}
 \|(QD_AQ+t)^{-1}\|_{2\to2}\le(t+g_{\mathrm h}/a)^{-1}.
\end{equation}
All constants are independent of $a,R$, the character, and the real
connection.
\end{proposition}

\begin{proof}
The internal bonds over $c$ give
$a^3(f(c,\cdot),\mathcal L_{\mathrm{int}}f(c,\cdot))$ and hence at
least $\lambda_{\mathrm{int}}a^3\|Qf(c,\cdot)\|^2$. Summing over
$c\in E$ proves the first bound in \eqref{eq:cluster-energy}.

For the second bound, consider a photon bond $c'=c+av$. The eight
labels with $ge_1=v$ correspond to the eight chamber bonds crossing it.
Write $\mathcal O_v(c)$ for this octet. Since their endpoints form
$\mathcal O_{-v}(c')$, the difference of the cluster means is the sum
of the mean of the eight bond differences and two means of $Qf$:
\begin{align*}
 u_A(c,c')\xi(c')-\xi(c)
 ={}&\frac18\sum_{x\in\mathcal O_v(c)}
       \big(U_A(x,x^{(3)})f(x^{(3)})-f(x)\big)\\
 &-\frac{u_A(c,c')}8\sum_{x'\in\mathcal O_{-v}(c')}Qf(x')
   +\frac18\sum_{x\in\mathcal O_v(c)}Qf(x).
\end{align*}
The square is bounded by $3/8$ times the sum of the squared terms in
these three octets. Distinct photon bonds use distinct chamber bonds;
each octet of cluster remainders occurs at most once. Consequently,
with $\mathcal E_3$ and $\mathcal E_{\mathrm{int}}$ denoting the two
parts of the energy,
\[
 a^2\|\nabla_A\xi\|_{2,E}^2
 \le18\big(w_3^{-1}\mathcal E_3+\|Qf\|_{2,E}^2\big)
 \le18\big(w_3^{-1}\mathcal E_3+
                \lambda_{\mathrm{int}}^{-1}\mathcal E_{\mathrm{int}}\big).
\]
This proves the local bound and therefore the global one.

The internal hops vanish on $J\xi$. Thus
\[
 (D_AJ\xi)(c,g)=-3(\rho-\gamma\cdot ge_1)
                   \nabla_{A,ge_1}\xi(c).
\]
Each of the six directions occurs eight times. The norm of the matrix
in this formula is $3(1+\rho)$, and the two directed differences on a
photon bond have equal norms. It follows that
$\|D_AJ\xi\|_2^2\le3(1+\rho)^2\|\nabla_A\xi\|_2^2$.
Projection gives the first bound in \eqref{eq:offdiagonal-bound}; the
second follows from \eqref{eq:dirac-realpart}, since $J$ and $Q$ commute
with $\gamma_5$.

For $Qf=f$, \eqref{eq:dirac-realpart} gives
$\operatorname{Re}(f,(QD_AQ+t)f)\ge(t+g_{\mathrm h}/a)\|f\|_2^2$.
The same lower bound holds for the adjoint. The operator has closed
range and its adjoint has zero kernel, so it is onto, also on the
infinite lattice. This proves \eqref{eq:heavy-inverse}.
\end{proof}

\subsection{The free symbol and its Schur complement}
\label{subsec:schur-symbol}

Use lattice units for the moment. Let $\mathcal D_m(p)$ be the symbol
of $D_0+m$, $0\le m\le1$, on
$V=\ell^2(O_{\mathrm{cub}};\C^4)$ with measure $48^{-1}\sum_g$:
\[
 (\mathcal D_m(p)v)(g)=m v(g)
 +\sum_{i=0}^2M_i(g)(v(g)-v(gs_i))
 +M_3(g)(v(g)-\exp(ip\cdot ge_1)v(gr_{e_1})).
\]
Let $D_{\mathrm{ll}},D_{\mathrm{lh}},D_{\mathrm{hl}},D_{\mathrm{hh}}$
be its blocks in $V=\operatorname{ran}J\oplus\operatorname{ran}Q$.
The mass is included in the diagonal blocks. Set
\[
 H_m(p)=D_{\mathrm{hh}}(p)^{-1},\qquad
 B_m(p)=D_{\mathrm{ll}}(p)-D_{\mathrm{lh}}(p)H_m(p)D_{\mathrm{hl}}(p),
 \qquad t(p)=\sum_\mu(1-\cos p_\mu).
\]
The compression of the full operator is the Wilson operator:
\[
 D_{\mathrm{ll}}(p)=m+i\gamma\cdot\sin p+\rho t(p).
\]
Its Schur complement also contains the contribution from the heavy
part. In particular, adding a mass to the chamber operator changes
$H_m$ as well as $D_{\mathrm{ll}}$.

\begin{proposition}\label{prop:schur-symbol}
There are positive constants $\theta_0,c_{\mathrm e},\kappa_{\mathrm{strip}}$, depending
on $\rho$ only, such that $H_m$ and $B_m$ are holomorphic in a
neighbourhood of $\{|\operatorname{Im}p_\mu|\le\theta_0\}$ and
$2\pi\Z^3$-periodic. On the real cube $[-\pi,\pi]^3$,
\begin{align}
 \|H_m(p)\|&\le(g_{\mathrm h}+m)^{-1},\label{eq:symbol-heavy}\\
 \|D_{\mathrm{lh}}(p)\|+\|D_{\mathrm{hl}}(p)\|
   &\le2c_{\mathrm{lh}}(2t(p))^{1/2},\label{eq:symbol-cross}\\
 B_m(p)&=m+i\gamma\cdot\sin p+E_m(p),\qquad
 \|E_m(p)\|\le K_\rho t(p),\label{eq:schur-expansion}\\
 \operatorname{Re}(v,B_m(p)v)&\ge(m+\rho c_*t(p))|v|_2^2,
                 \label{eq:schur-realpart}\\
 |B_m(p)v|_2&\ge c_{\mathrm e}(|p|_2^2+m^2)^{1/2}|v|_2,
                 \label{eq:schur-ellipticity}
\end{align}
where $K_\rho=\rho+2c_{\mathrm{lh}}^2/g_{\mathrm h}$. For every
multi-index $\nu$,
\begin{equation}\label{eq:schur-inverse-derivatives}
 \|\partial_p^\nu B_m(p)^{-1}\|
       \le C_\nu(\rho)(|p|_2^2+m^2)^{-(1+|\nu|)/2}.
\end{equation}
If $m>0$, the inverse exists also on
$|\operatorname{Im}p_\mu|\le\kappa_{\mathrm{strip}}m$, where the same bound holds
with $\operatorname{Re}p$ on the right, for
$\operatorname{Re}p\in[-\pi,\pi]^3$. At $m=0$ the real derivative
bound is asserted away from $p=0$.
\end{proposition}

\begin{proof}
Only the fourth hop depends on $p$. For $p=p'+ip''$,
\[
 \|\mathcal D_m(p)-\mathcal D_m(p')\|
       \le3(1+\rho)(\exp(|p''|_\infty)-1).
\]
Choose $\theta_0$ so that the right side is $g_{\mathrm h}/2$ when
$|p''|_\infty=\theta_0$. The heavy gap then gives
$\operatorname{Re}D_{\mathrm{hh}}(p)\ge m+g_{\mathrm h}/2$ on this
strip. Invertibility holds on a slightly larger strip, which proves
holomorphy up to the closed strip. For real $p$, the off-diagonal bounds
in Proposition~\ref{prop:cluster-energy} give
$\|D_{\mathrm{lh}}\|,\|D_{\mathrm{hl}}\|
\le c_{\mathrm{lh}}(2t(p))^{1/2}$. This proves
\eqref{eq:schur-expansion}.

To retain the real-part bound after eliminating the heavy part, put
$u=Jv-H_mD_{\mathrm{hl}}v$. Then
$J^*u=v$, $Q\mathcal D_m u=0$, and
$\mathcal D_m u=JB_mv$. Therefore
\[
 \operatorname{Re}(v,B_mv)
   =\operatorname{Re}(u,\mathcal D_m u)
   =m\|u\|^2+\frac\rho2(u,\mathcal L_0(p)u)
   \ge(m+\rho c_*t(p))|v|_2^2.
\]
This proves \eqref{eq:schur-realpart}.

The ellipticity estimate uses the Dirac part near zero and the real
part away from zero. Put $t_0=(4K_\rho^2)^{-1}$. One has
$K_\rho\ge48/\lambda_{\mathrm{int}}>1$, since
$\lambda_{\mathrm{int}}\le12+24\sqrt2<48$.
Hence $t_0\le1/4$. If $t(p)\le t_0$, then
$|\sin p|_2^2\ge t(p)$ and
$\|E_m(p)\|\le\tfrac12t(p)^{1/2}$. Since
\[
 |(m+i\gamma\cdot\sin p)v|_2^2=(m^2+|\sin p|_2^2)|v|_2^2,
\]
it follows that
$|B_m(p)v|_2\ge\tfrac12(m^2+t(p))^{1/2}|v|_2$.
On $[-\pi,\pi]^3$, $t(p)\ge2|p|_2^2/\pi^2$. This proves
\eqref{eq:schur-ellipticity} in this region. If $t(p)\ge t_0$, then
\eqref{eq:schur-realpart} gives
$|B_m(p)v|_2\ge(m+\rho c_*t_0)|v|_2$.
The inequality $(|p|_2^2+m^2)^{1/2}\le m+\pi\sqrt3$ closes the
estimate with, for example,
\[
 c_{\mathrm e}=\min\left(0.22,\frac{\rho c_*t_0}{\pi\sqrt3}\right).
\]

The strip bounds give uniform bounds on the derivatives of $B_m$ by
Cauchy's formula. Differentiating $B_mB_m^{-1}=1$ and inducting on
$|\nu|$ proves \eqref{eq:schur-inverse-derivatives}. For the complex
strip, let $K_s$ bound $B_m$ on the fixed strip. On half that strip,
$\|\partial_\mu B_m\|\le2K_s/\theta_0$. Choose
\[
 \kappa_{\mathrm{strip}}=\min\left(\frac{\theta_0}{2},
                \frac{c_{\mathrm e}\theta_0}{12K_s}\right).
\]
Integration along the segment from $p'$ to $p'+ip''$ shows that the
change in $B_m$ is at most $c_{\mathrm e}m/2$. The lower bound is thus
at least $c_{\mathrm e}(|p'|_2^2+m^2)^{1/2}/2$. Differentiating the
inverse identity once more gives the asserted complex derivative bounds.
\end{proof}

\subsection{Free kernels and coercivity}\label{subsec:free-kernels}

Solving the heavy equation gives the inverse symbol
\begin{equation}\label{eq:inverse-symbol-blocks}
 \mathcal D_m^{-1}=
 \begin{pmatrix}
 B_m^{-1}&-B_m^{-1}D_{\mathrm{lh}}H_m\\
 -H_mD_{\mathrm{hl}}B_m^{-1}&
 H_m+H_mD_{\mathrm{hl}}B_m^{-1}D_{\mathrm{lh}}H_m
 \end{pmatrix}.
\end{equation}
The first term is independent of the cluster labels after insertion by
$J$ and $J^*$. Every other term has a weaker singularity at $p=0$:
$D_{\mathrm{lh}}$ and $D_{\mathrm{hl}}$ vanish there, whereas $H_m$
and all its derivatives are uniformly bounded.

\begin{proposition}\label{prop:free-chamber-kernel}
On the unit covering lattice let
\[
 G_m(c,g;c',g')=\int_{[-\pi,\pi]^3}\frac{\dd p}{(2\pi)^3}
       \exp(ip\cdot(c-c'))[\mathcal D_m(p)^{-1}](g,g').
\]
The bracket is the kernel on $V$ relative to the label measure
$48^{-1}\sum_g$, consistently with the chamber kernel convention.
It is well-defined also at $m=0$. It decomposes as
$G_m(x,x')=G_m^{\mathrm l}(c-c')+G_m^{\mathrm h}(x,x')$, where
\begin{align}
 |G_m^{\mathrm l}(n)|_2&\le C_G(1+|n|_\infty)^{-2},\label{eq:free-light}\\
 |G_m^{\mathrm l}(n+v)-G_m^{\mathrm l}(n)|_2
       &\le C_G(1+|n|_\infty)^{-3},\quad v\in\{\pm e_\mu\},\\
 |G_m^{\mathrm h}(x,x')|_2&\le C_G(1+|c-c'|_\infty)^{-3}.
\end{align}
Consequently $|G_m(x,x')|_2\le2C_G(1+|c-c'|_\infty)^{-2}$ and
\begin{equation}\label{eq:free-bond-difference}
 |G_m(x^{(i)},x')-G_m(x,x')|_2
       \le20C_G(1+|c-c'|_\infty)^{-3}.
\end{equation}
For $m>0$ these bounds have the additional factor
$\exp(-\kappa_{\mathrm{strip}}m|c-c'|_\infty)$. The constants are uniform for
$0\le m\le1$.
\end{proposition}

\begin{proof}
The Fourier estimate used here can be proved directly. If a periodic
symbol $T$ satisfies
$|\partial^\nu T(p)|\le A|p|_2^{-q-|\nu|}$, where $q<3$ and
$|\nu|\le N>3-q$, then its Fourier coefficients are bounded by
$C A(1+|n|_\infty)^{-(3-q)}$. For $n\ne0$, put
$r=|n|_\infty^{-1}$ and split the integral using a cutoff equal to
one for $|p|_2\le r$ and zero for $|p|_2\ge2r$. The inner integral
is bounded by $CA r^{3-q}$. On the complement integrate $N$ times in
a coordinate with $|n_\mu|=|n|_\infty$. The differentiated integrand
has integral at most $CA r^{3-q-N}$; the prefactor is $r^N$.
Derivatives of the cutoff are supported in $r\le|p|_2\le2r$ and have
the same bound. Boundary terms cancel by periodicity.

Apply this estimate with $(q,N)=(1,3)$ to $B_m^{-1}$. Multiplication
by $\exp(ip\cdot v)-1$ lowers its singularity to order zero, for which
$(q,N)=(0,4)$ applies. In \eqref{eq:inverse-symbol-blocks}, the two
off-diagonal terms have order zero and the term with two off-diagonal
factors has order $-1$. The latter can also be bounded by the
order-zero estimate. The remaining term $H_m$ is holomorphic on the
fixed strip, and its Fourier coefficients decay exponentially.
The kernel normalization costs a factor $48$ for operator-valued
symbols on $V$; this is absorbed in $C_G$.

For $m>0$, shift one coordinate of the contour by
$i\kappa_{\mathrm{strip}}m\operatorname{sgn}(n_\mu)$, where
$|n_\mu|=|n|_\infty$. The complex derivative bounds of
Proposition~\ref{prop:schur-symbol} give the same Fourier estimates
on the shifted contour. This supplies the exponential factor.
The fixed-strip bound for $H_m$ is stronger than the required one.

An internal hop leaves $c$ fixed, so its difference cancels
$G_m^{\mathrm l}$. An external hop replaces $c$ by $c+ge_1$; use the
light difference bound and the two heavy bounds at these points.
The inequalities $1+|n+v|_\infty\ge(1+|n|_\infty)/2$ and
$\exp(\kappa_{\mathrm{strip}}m)<\sqrt3$ give the factor $20$ in
\eqref{eq:free-bond-difference}.
\end{proof}

The Green kernel at mesh $a$, for $0\le\mu a\le1$, is
$a^{-2}G_{\mu a}(x/a,x'/a)$. Thus
\begin{equation}\label{eq:mesh-free-kernel}
 |G(x,x')|_2\le2C_G(a+|c-c'|_\infty)^{-2},\qquad
 |G(x^{(i)},x')-G(x,x')|_2
       \le20C_Ga(a+|c-c'|_\infty)^{-3}.
\end{equation}
Fourier multiplication gives both inverse identities on finitely
supported functions. The factor $48$ in the identity kernel is
consistent with the measure $\dd x$ above.

For distinct sites put
$d_X(x,\widetilde x)=\max(|c-\widetilde c|_\infty,a)$.
The chamber graph has paths from $x$ to $\widetilde x$ of at most
$39d_X(x,\widetilde x)/a$ bonds, with photon sites in the coordinate
box spanned by $c,\widetilde c$. Indeed, the generators
$s_0,s_1,s_2$ are the simple reflections of the signed permutation
group, with nine positive roots. Their Coxeter length is at most nine
\cite[Theorem~22.14]{EtingofLieNotes}, so any two labels in one
cluster are joined by at most nine internal bonds. A monotone photon
path has at most $3|c-\widetilde c|_\infty/a$ steps. Before each step
move to a label with the required external direction, and after the
last step move to the final label. This uses at most
$10\cdot3|c-\widetilde c|_\infty/a+9\le39d_X/a$ bonds.
If $d_X$ is less than half the distance to $c'$, sum
\eqref{eq:mesh-free-kernel} along such a path. In the other case use
the two size bounds. The resulting interpolation gives, for
$0\le\alpha\le1$,
\begin{equation}\label{eq:free-holder-kernel}
 |G(x,x')-G(\widetilde x,x')|_2
 \le6240C_G d_X(x,\widetilde x)^\alpha
 \left((a+|c-c'|_\infty)^{-2-\alpha}
       +(a+|\widetilde c-c'|_\infty)^{-2-\alpha}\right).
\end{equation}

The following coercivity estimate is separate from the real-part bound.
It is needed to recover mesh-independent control of the block remainder.

\begin{proposition}\label{prop:free-coercivity}
There is $c_b=c_b(\rho)>0$ such that
\begin{equation}\label{eq:free-coercivity}
 \|D_0f\|_2^2\ge c_ba^{-2}(f,\mathcal L_0f).
\end{equation}
For a cube $E$ of side $\ell_E$, let $P_E$ be the mean over
$E\times O_{\mathrm{cub}}$. Then
\begin{equation}\label{eq:chamber-poincare}
 \|f-P_Ef\|_{2,E}^2\le C_P\ell_E^2a^{-2}(f,\mathcal L_0f)_E,
 \qquad C_P=\frac1{8c_*}+\frac1{\lambda_{\mathrm{int}}}.
\end{equation}
\end{proposition}

\begin{proof}
The Schur estimate at zero mass implies
$|B_0(p)v|\ge c_{\mathrm l}t(p)^{1/2}|v|$, with
$c_{\mathrm l}=\rho c_*/(2K_\rho)$: use the Dirac part for
$t(p)\le t_0$ and the real part for $t(p)\ge t_0$.
Write $u=Jv+h$ in a fibre and $f=\mathcal D_0(p)u$. The heavy equation
and then the light equation give
\[
 h=H_0(f_{\mathrm h}-D_{\mathrm{hl}}v),\qquad
 B_0v=f_{\mathrm l}-D_{\mathrm{lh}}H_0f_{\mathrm h}.
\]
Since $t(p)\le6$, these imply
$(2t(p))^{1/2}|v|\le A_1\|f\|$ and $\|h\|\le A_2\|f\|$, where
\[
 A_1=\frac{\sqrt2}{c_{\mathrm l}}
         \left(1+\frac{6(1+\rho)}{g_{\mathrm h}}\right),\qquad
 A_2=g_{\mathrm h}^{-1}(1+c_{\mathrm{lh}}A_1).
\]
The bond decomposition also gives
\[
 (u,\mathcal L_0(p)u)\le6t(p)|v|^2+(30+24\sqrt2)\|h\|^2.
\]
For the internal bonds the light part cancels, leaving at most
$2(w_0+w_1+w_2)\|h\|^2$. In an external bond difference, split off
the photon difference of $v$ and the two heavy endpoint values, and
bound the square by three times their squared norms. The octet count
then gives $6t(p)|v|^2+3w_3\|h\|^2$. This proves the displayed
upper bound.
This yields \eqref{eq:free-coercivity} in each fibre with
$c_b=(3A_1^2+(30+24\sqrt2)A_2^2)^{-1}$.
Fourier decomposition gives the assertion for all characters on the
torus and on the covering lattice, restoring the factor $a^{-2}$.

For \eqref{eq:chamber-poincare}, first apply the scalar lattice
Poincar\'e bound
$\|\xi-P_E\xi\|_{2,E}^2\le\ell_E^2\|\nabla_0\xi\|_{2,E}^2/8$.
In one dimension this follows by writing each difference as a sum of
nearest-neighbour differences; the coefficient of a bond at position
$j$ is $j(n-j)/2\le n^2/8$. In three dimensions use the three commuting
line averages, whose successive remainders are orthogonal.
Now $f-P_Ef=J(\xi-P_E\xi)+Qf$. Its two parts are orthogonal, and the
local bounds \eqref{eq:cluster-energy} give
\eqref{eq:chamber-poincare}. The heavy term uses $\ell_E\ge a$.
\end{proof}

\subsection{Averaging penalties and the inverse at zero field}
\label{subsec:blocked-zero-inverse}

Fix an odd $L\ge5$ and put $a=\eta_k=L^{-k}$. Write $e_k$ for the
charge $e_a$ at this mesh. The level data $k,\mu,R,e_k$ are independent,
with $0<e_k\le e_0$; the flow will specify their relation later.
Unit blocks have side
one, and the next blocks have side $L$. In this and the following
subsections a finite period $R$ is a positive integer multiple of $L$.
For $z\in\Z^3$ the unit block is
$\Delta_z=\{c\in a\Z^3:|c-z|_\infty<1/2\}$.
For $z'\in L\Z^3$ the $L$-block is
$E_{z'}=\{c\in a\Z^3:|c-z'|_\infty<L/2\}$.
These partitions are taken modulo $R\Z^3$ on the torus. Since $L$
is odd, no fine site lies on a block boundary. Their means, including
the cluster mean, are denoted $\Pi_1$ and $\Pi_L$. At nonzero connection
the means use the gauge phases specified below. Let
\[
 b=1,\qquad b_L=b/L,\qquad b_-=b(1-L^{-1}),\qquad
 b_k=\frac{b(1-L^{-1})}{1-L^{-k}}\quad(k\ge1).
\]
For $k\ge1$, $y\in\R$, define
\[
 \alpha_y=b_k i\gamma_5y(b_k+i\gamma_5y)^{-1},\qquad
 \beta_y=b_k^2b_L(b_k+b_L+i\gamma_5y)^{-1}(b_k+i\gamma_5y)^{-1},
\]
and set $\alpha_\infty=b_k$, $\beta_\infty=0$. The blocked operators
at zero connection are
\begin{equation}\label{eq:blocked-zero-operator}
 F_y(0)=D_0+\mu+\alpha_y\Pi_1+\beta_y\Pi_L\quad(k\ge1),
 \qquad
 F_y(0)=D_0+\mu+b_L\Pi_L+i\gamma_5y\quad(k=0).
\end{equation}
Here $0\le\mu\le1$. In the first formula $y\in\R\cup\{\infty\}$;
in the second $y\in\R$. We use $-\infty=\infty$ when an adjoint
identity reverses this parameter. The notation $F_y$ denotes an operator in this
section, not the curvature.

The penalty bounds follow by diagonalizing $\gamma_5$. On its eigenspace
of sign $\epsilon$, the two eigenvalues and their sum are
\[
 \frac{b_ki\epsilon y}{b_k+i\epsilon y},\quad
 \frac{b_k^2b_L}{(b_k+b_L+i\epsilon y)(b_k+i\epsilon y)},\quad
 \frac{b_k(b_L+i\epsilon y)}{b_k+b_L+i\epsilon y}.
\]
For the first eigenvalue $v_1$ and the sum $v_2$, direct calculation gives
\[
 |v_1|^2=b_k\operatorname{Re}v_1,\qquad
 |v_2|^2\le b_k\operatorname{Re}v_2,\qquad
 \operatorname{Re}v_2\ge\frac{b_kb_L}{b_k+b_L}\ge\frac{b_-}{L}.
\]
Also $|\alpha_y|\le b_k\le b$ and $|\beta_y|\le b_L$.
On an $L$-block the decomposition
$\operatorname{Pen}=\alpha_y(\Pi_1-\Pi_L)+(\alpha_y+\beta_y)\Pi_L$
has orthogonal ranges. The scalar inequalities on the two spin
eigenspaces therefore give
\begin{equation}\label{eq:penalty-estimates}
 \operatorname{Re}(f,\operatorname{Pen}f)_E
      \ge\frac{b_-}{L}\|P_Ef\|_{2,E}^2,
 \qquad
 \|\operatorname{Pen}f\|_{2,E}^2
      \le b\operatorname{Re}(f,\operatorname{Pen}f)_E.
\end{equation}
The kernel on $E$ is
$\alpha_y\mathbf1[c,c'\text{ in the same unit block}]+\beta_yL^{-3}$,
and its norm is at most $2b$.

\begin{proposition}\label{prop:zero-blocked-inverse}
There are $C_I,C_5,\gamma_5'>0$, depending on $L,\rho$ and a fixed
$\alpha\in(0,1)$ only, such that $F_y(0)$ is invertible and
\[
 \|F_y(0)^{-1}\|_{2\to2}\le LC_I\qquad(k\ge1).
\]
If $f$ is supported over one unit block $\Delta_{z'}$ and
$v=F_y(0)^{-1}f$, then for $x$ over $\Delta_z$ and
$0<d_X(x,\widetilde x)\le1$,
\begin{equation}\label{eq:zero-blocked-holder}
 |v(x)|_2+d_X(x,\widetilde x)^{-\alpha}|v(x)-v(\widetilde x)|_2
 \le C_5\exp(-\gamma_5'd(z,z'))\|f\|_\infty.
\end{equation}
This includes pairs in one cluster. The constants are uniform in
$k,R$, the character, $\mu$, and $y$.
\end{proposition}

\begin{proof}
First suppose $k\ge1$. Put $f=F_y(0)u$ and
$X=\operatorname{Re}(u,f)$. The real-part identity
and \eqref{eq:penalty-estimates} imply
\[
 \sum_E\|P_Eu\|_{2,E}^2\le\frac L{b_-}X,\qquad
 \|\operatorname{Pen}u\|_2^2\le bX,\qquad
 \|\mu u\|_2^2\le X.
\]
Therefore $\|D_0u\|_2\le\|f\|_2+(1+\sqrt b)X^{1/2}$.
Use \eqref{eq:free-coercivity} for $D_0u$ and
\eqref{eq:chamber-poincare} on each $L$-block. With
$K_1=\sqrt L+\sqrt{Lb}$ this gives
\[
 \|u\|_2^2\le L\left(\frac1{b_-}+\frac{2C_PK_1^2}{c_b}\right)
           \|u\|_2\|f\|_2
       +\frac{2C_P}{c_b}L^2\|f\|_2^2.
\]
Hence $\|u\|_2\le LC_I\|f\|_2$, where
$C_I=b_-^{-1}+2C_PK_1^2/c_b+(2C_P/c_b)^{1/2}$.
The identity $F_y(0)^*=\gamma_5F_{-y}(0)\gamma_5$ gives the same
lower bound for the adjoint. Thus $F_y(0)$ has closed, dense range
and is onto, including on the infinite lattice.

To obtain decay, let $\psi$ be a bounded real periodic function with
$|\psi(c)-\psi(c')|\le a_0|c-c'|_1/L$. The mass commutes with this
weight. The external hopping term gives
$\|\exp(\psi)[D_0,\exp(-\psi)]\|_{2\to2}
\le3(1+\rho)\exp(1)a_0/L$.
The penalty kernel is supported in one $L$-block, has norm at most
$2b$, and $|\psi(c)-\psi(c')|\le3a_0$ there. Its contribution is at
most $6L^4b\exp(1)a_0/L$. Put
\[
 c_w=\exp(1)(3(1+\rho)+6L^4b),\qquad
 a_0=\min(1/3,(2C_Ic_w)^{-1}).
\]
Applying the $L^2$ inverse bound and absorbing the commutator gives
\begin{equation}\label{eq:weighted-block-inverse}
 \|\exp(\psi)F_y(0)^{-1}f\|_2
       \le2LC_I\|\exp(\psi)f\|_2.
\end{equation}
Take truncated multiples of the $\ell^1$ distance, divided by $L$,
to the source $L$-block. Letting the truncation increase proves an
exponentially weighted $L^2$ bound.

The passage from $L^2$ to the sup norm uses the free kernel. On the
covering lattice choose $c^\circ\in a\Z^3$ and $r\ge8a$, and put
\[
 B=\{|c-c^\circ|_\infty<r\},\quad
 B^+=\{|c-c^\circ|_\infty<r+a\},\quad
 B'=\{|c-c^\circ|_\infty\le r/4\}.
\]
Suppose $u$ is defined over $B^+$ and $(D_0+\mu)u=h$ over $B$,
with $0\le\mu a\le1$. Both evaluation points below lie over $B'$,
and $C$ depends on the fixed $\alpha\in(0,1)$ and $\rho$.
Let $\chi$ equal one
on the concentric half cube and zero outside the cube, with Lipschitz
constant $2/r$. Then
\[
 \chi u=G(\chi h)+G([D_0,\chi]u).
\]
The commutator contains only the external hop. Its rows are supported
where $|c-c^\circ|_\infty>r/2-a$. It satisfies
$\|[D_0,\chi]u\|_1\le45r^{1/2}\|u\|_{2,B^+}$.
For a point of the quarter cube, this support is at distance at least
$r/4-a\ge r/8$. The bounds \eqref{eq:mesh-free-kernel} and
\eqref{eq:free-holder-kernel} consequently give
\begin{equation}\label{eq:interior-dirac-estimate}
 \begin{split}
 |u(x)|_2&\le C\big(r\|h\|_{\infty,B}+r^{-3/2}\|u\|_{2,B^+}\big),\\
 |u(x)-u(\widetilde x)|_2&\le
 C\left(\frac{d_X(x,\widetilde x)}r\right)^\alpha
       \big(r\|h\|_{\infty,B}+r^{-3/2}\|u\|_{2,B^+}\big).
 \end{split}
\end{equation}
For the first term, sum $(a+|c-c'|_\infty)^{-2}$ over the cube;
the result is $O(r)$. For the difference, the analogous sum with
exponent $2+\alpha$ is $O(r^{1-\alpha}/(1-\alpha))$.
The cluster average in $\dd x'$ contributes one to both sums.
For the commutator term, its support separation gives the size bound;
summing bond differences along the chamber path gives the difference
bound. These estimates explain why the constants do not grow when
$a\downarrow0$.

For $u=F_y(0)^{-1}f$, use
$h=f-\operatorname{Pen}u$ in \eqref{eq:interior-dirac-estimate}, with
$r=L/4$. The enlarged cube meets at most eight $L$-blocks. On each,
the penalty kernel and Cauchy--Schwarz give
$\|\operatorname{Pen}u\|_\infty\le2bL^{3/2}\|u\|_{2,E}$.
The weighted bound \eqref{eq:weighted-block-inverse} therefore bounds
both terms in \eqref{eq:interior-dirac-estimate} by
$CL\exp(-a_0d_1(x,E'))\|f\|_\infty$, where
$d_1(x,E')$ is the $\ell^1$ distance to the source $L$-block divided
by $L$. This applies for $k\ge2$.
For $k=1$, the pointwise $L^2$ bound costs only
$\sqrt{48}L^{3/2}$; $L$ is fixed. Larger Hölder separations are
controlled by the two sup bounds. Finally,
$d_1(x,E')\ge d(z,z')/L-3$, and the source unit block lies in $E'$.
This proves \eqref{eq:zero-blocked-holder} with
$\gamma_5'=a_0/L$. The argument uses only periodic moduli and so
applies to every unitary character.
\end{proof}

The first level has a different resolvent parameter. On an $L$-block
the real-part bound and Poincar\'e inequality imply
\[
 \operatorname{Re}(u,F_y(0)u)\ge\frac{c_r}{L}\|u\|_2^2,
 \qquad c_r=\min(b_-,\rho/(2C_PL)).
\]
The term $i\gamma_5y$ has zero real part. Conjugation by a weight with
Lipschitz constant $a_z$ changes the hopping and penalty terms by at
most $\exp(1)(6+3b)a_z$. Choose
\[
 a_z\le\min\bigl((3L)^{-1},c_r/(2\exp(1)L(6+3b))\bigr).
\]
The weighted inverse then has $L^2$ norm at most $2L/c_r$.
On the other hand its part other than $i\gamma_5y$ has norm at most
$60+b$. For $|y|\ge2(60+b)$, the inverse norm is at most $2/|y|$.
Combining the two ranges and converting to a kernel, with its factor
$48$, gives
\begin{equation}\label{eq:level-zero-kernel}
 |F_y(0)^{-1}(x,x')|_2\le C_{00}(1+|y|)^{-1}
                         \exp(-a_zd(c,c')).
\end{equation}
This factor in $y$ is needed when the determinant is represented by an
integral of resolvents.

\subsection{Small connections and local inverses}\label{subsec:local-inverses}

The covariant block means are defined by the same nested blocks.
Let $\tau A(c,c')$ be the arithmetic average of the oriented bond sums
over the six coordinate-order paths
from $c$ to $c'$, and let $c=c_0,c_1,\ldots,c_k=z$ be the centres
of the nested blocks containing $c$. Put
\[
 \varphi_A(z,c)=\exp\left(ie_k\eta_k
                    \sum_{j<k}\tau A(c_{j+1},c_j)\right).
\]
The unit-block average on four-component spinors and its inverse
injection are
\[
 (Q_k(A)\xi)(z)=\eta_k^3\sum_{c\in\Delta_z}\varphi_A(z,c)\xi(c),
 \qquad
 (Q_k^T(-A)v)(c)=\varphi_A(z,c)^{-1}v(z).
\]
For the next average set, on the unit lattice,
\[
 (Q(A)v)(z')=L^{-3}\sum_{|z-z'|_\infty<L/2}U_A(z',z)v(z),
 \qquad (Q^T(-A)v')(z)=U_A(z',z)^{-1}v'(z').
\]
On the chamber space, for $k\ge1$, put
\[
 \Pi_1(A)=JQ_k^T(-A)Q_k(A)J^*,\qquad
 \Pi_L(A)=JQ_k^T(-A)Q^T(-A)Q(A)Q_k(A)J^*.
\]
At $k=0$ only $\Pi_L(A)=JQ^T(-A)Q(A)J^*$ is used.
This defines $F_y(A)$ by \eqref{eq:blocked-zero-operator}.
The formulas permit complex connections; for real $A$ the
idempotents are orthogonal projections.

For a pair of sites define
\[
 U_A(x,x')=\exp(ie_k\eta_k\tau A(c,c')),\qquad
 (\delta_{\alpha,A}v)(x,x')=
 \frac{U_A(x,x')v(x')-v(x)}{d_X(x,x')^\alpha}.
\]
Only pairs with $0<d_X<1$ enter the local norms. If $c=c'$ then
$U_A=1$, so these quotients include variations between the labels in
one cluster. There are no such short pairs when $k=0$.

For an operator on the torus define
\[
 N_\gamma(T)=\sup_z\sum_{z'}\exp(\gamma d(z,z'))
                     \|\mathbf1_{\Delta_z}T\mathbf1_{\Delta_{z'}}\|_{\infty\to\infty}.
\]
For operators with values on pairs, the left indicator restricts the
first point. Inserting the block partition between two factors and
using the triangle inequality for $d$ proves
$N_\gamma(TT')\le N_\gamma(T)N_\gamma(T')$.
Write $w_k(y)=1$ for $k\ge1$ and $w_0(y)=(1+|y|)^{-1}$.
Enlarge $C_5,C_{00}$ to at least one, and choose
\[
 \gamma=\tfrac12\min(\gamma_5',a_z),\qquad
 a'=\max(C_5,C_{00})\sum_{n\in\Z^3}\exp(-\gamma|n|_\infty).
\]
The preceding estimates imply
\begin{equation}\label{eq:zero-row-bounds}
 N_\gamma(F_y(0)^{-1})\le a'w_k(y),\qquad
 N_\gamma(\delta_{\alpha,0}F_y(0)^{-1})\le a'.
\end{equation}
At $k=0$ the analogous bound holds in the kernel row norm
$N_\gamma^{\mathrm{ker}}$, obtained by replacing the localized
operator norm with the supremum of $|T(x,x')|_2$ over the two clusters.

\begin{proposition}\label{prop:small-connection-inverse}
Let $B$ be a complex connection with $e_k\|B\|_\infty\le\delta$, where
$3L\delta\le1$ and $2C_Va'\delta\le1$, with
$C_V=2\exp(1)(42+6b)$. Then $F_y(B)$ is invertible and
\begin{equation}\label{eq:small-row-bounds}
 N_\gamma(F_y(B)^{-1})\le2a'w_k(y),\qquad
 N_\gamma(\delta_{\alpha,B}F_y(B)^{-1})\le18a'.
\end{equation}
At $k=0$ the first bound holds also in $N_\gamma^{\mathrm{ker}}$.
\end{proposition}

\begin{proof}
Write $V=F_y(B)-F_y(0)$. The internal hops, the mass, and the
resolvent term at level zero cancel. For the external hop,
$|u_B-1|\le\eta_k\delta\exp(\eta_k\delta)$; its weighted row sum
is at most $42\exp(\gamma)\exp(\delta)\delta$. The two averaging
penalties contribute respectively
$3b\delta\exp(3\delta)$ and
$3b\delta\exp(\gamma(L-1))\exp(3L\delta)$. Indeed, the nested paths
have total length at most $3/2$, and the path to an $L$-block centre
has length at most $3L/2$. Summing their normalized block kernels
removes the numbers of sites. The bounds $\gamma L\le\log2$ and
$3L\delta\le1$ give $N_\gamma(V)\le C_V\delta$.

Thus $N_\gamma(VF_y(0)^{-1})\le1/2$. The Neumann series gives
\[
 F_y(B)^{-1}=F_y(0)^{-1}
                  \sum_{n\ge0}(-VF_y(0)^{-1})^n
             =\sum_{n\ge0}(-F_y(0)^{-1}V)^nF_y(0)^{-1}.
\]
It proves the first bound in \eqref{eq:small-row-bounds}. The first
form gives $N_\gamma(\delta_{\alpha,0}F_y(B)^{-1})\le2a'$.
For distinct photon sites in a short pair,
$d_X^{-\alpha}|U_B-1|\le3\delta\exp(3\delta)\le1/4$;
for intracluster pairs the expression is zero. Evaluation at the
second point of a short pair uses at most $27$ neighbouring blocks,
and costs at most $27\exp(\gamma)$. Its contribution to the covariant
quotient is at most $16a'$, proving the second bound. At level zero,
the second form of the Neumann series keeps the kernel norm in the
last factor, so the kernel estimate follows as well.
\end{proof}

Fix a positive exponent $\varepsilon$; its choice in the flow will
be made with the other renormalization parameters. The connection
domain used by the RG is the following. Write $\ell=R$ for the side,
and assume $\ell/(LM)$ is odd and at least $13$.
Here $M$ is a power of $L$ with $M\ge L^2$. The $M$-cubes have centres
$z_\square\in M\Z^3/\ell\Z^3$ and are the periodic images of
$\{c:|c-z_\square|_\infty<M/2\}$. The differential of
a scalar in a connection is normalized by
$(d\lambda)(c,c')=\eta_k^{-1}(\lambda(c')-\lambda(c))$.
A connection belongs to $\widetilde{\mathcal G}_k$ if $A=A_0+A_1$,
where $A_0$ is real, $\|A_1\|_\infty<e_k^{-1+\varepsilon}$, and
on every $11M$ enlargement of an $M$-cube,
$A_0=B_0+d\lambda$, with $B_0,\lambda$ real and
$\|B_0\|_\infty<e_k^{-1+\varepsilon}$.
Consequently, on the $3M$ enlargement $W_\square$,
\begin{equation}\label{eq:local-small-gauge}
 A=B+d\lambda,\qquad e_k\|B\|_\infty<2e_k^\varepsilon,
 \qquad e_k\|\operatorname{Im}A\|_\infty<e_k^\varepsilon.
\end{equation}
This follows by restricting $B_0+A_1$ and $\lambda$ to that cube.
The domain is open: a sufficiently small complex perturbation can be
added to $A_1$. It is invariant under real gauge transformations.

For a lift $z_\square$ of the cube centre put
$W_\square=\{|c-z_\square|_\infty<3M/2\}$ and
$W_\square^-=\{|c-z_\square|_\infty<3M/2-\eta_k\}$.
Embed $W_\square$ in the positive chamber of the periodic torus of
side $6M$, whose lattice is translated by $\eta_k(1,1,1)/2$.
Its level-$j$ block centres are
$L^j\eta_k(1,1,1)/2+L^j\eta_k\Z^3$ for $0\le j\le k+1$.
Translation by $L(1,1,1)/2$ identifies these blocks with the
unshifted nested blocks, because $L^{k+1-j}$ is odd. The zero-field
and small-connection estimates therefore apply on this torus with
the same constants.
Fold its photon coordinates across the planes at multiples of $3M$.
Adjacent sites fold either to adjacent sites in $W_\square$ or to
one site; in the latter case assign phase one. Pull back $A$ by this
fold to obtain $A^W$. The fold preserves the supremum norm and
commutes with $d$ on scalar functions. It also preserves the nested
blocks, since $3M/(L^j\eta_k)$ is odd for $j\le k+1$.

Let $\mathrm{ext}$ extend a function over $W_\square$ by zero in the
seven other chambers, and let $\mathrm{res}$ restrict to its embedded
copy. Labels are unchanged in these two maps. Define
\[
 S_\square(A)=\mathrm{res}\,F_y(A^W)^{-1}\mathrm{ext}.
\]
For $A=B+d\lambda$ as in \eqref{eq:local-small-gauge},
$A^W=B^W+d\lambda^W$, with $\lambda^W$ real. Gauge conjugation by
$\exp(ie_k\lambda^W)$ reduces the inverse to the small connection
$B^W$. If $6Le_0^\varepsilon\le1$ and
$4C_Va'e_0^\varepsilon\le1$, Proposition~\ref{prop:small-connection-inverse}
applies to every $0<e_k\le e_0$. Restriction, extension, and unitary
gauge factors do not increase the localized norms. Thus
\begin{equation}\label{eq:local-inverse-bounds}
 N_\gamma(S_\square(A))\le2a'w_k(y),\qquad
 N_\gamma(\delta_{\alpha,A}S_\square(A))\le18a'.
\end{equation}
The second operator here is formed on the local cube before extension
to the original torus. It does not differentiate a discontinuous
extension by zero.

The local inverse has both identities needed for patching. If $f$ is
defined over $W_\square$, then
\begin{equation}\label{eq:local-right-identity}
 F_y(A)S_\square(A)f=f\quad\text{over }W_\square^-.
\end{equation}
If $f$ is supported over $W_\square^-$, then $F_y(A)f$ is supported
over $W_\square$ and
\begin{equation}\label{eq:local-left-identity}
 S_\square(A)F_y(A)f=f.
\end{equation}
For \eqref{eq:local-right-identity}, the rows at an interior site use
only its internal neighbours, its external neighbour, and its
$L$-block. These lie in $W_\square$, where the original and folded
operators agree. For \eqref{eq:local-left-identity}, the support
margin ensures that applying $F_y$ cannot leave $W_\square$.
After embedding, applying the folded operator to the extension of
$f$ agrees with extending $F_y(A)f$. Apply the inverse and restrict.
For an antiperiodic spinor on the original torus, use one lift and
then its signed translates; the enlarged cubes do not overlap their
translates. These arguments also show that $S_\square$ reads $A$
only in $W_\square$ and is analytic in those connection variables.

\subsection{The parametrix and the walk expansion}\label{subsec:propagator-walks}

Choose the product cutoff $h_\square$ with centre $z_\square$ and
one-dimensional factor
\[
 h_0(t)=\begin{cases}
 1,&|t|\le1/4,\\
 \cos(\pi(2|t|-1/2)/2),&1/4\le|t|\le3/4,\\
 0,&|t|\ge3/4.
 \end{cases}
\]
On a lift of the cube set
\[
 h_\square(c)=\prod_{\mu=1}^3
 h_0\bigl((c_\mu-z_{\square,\mu})/M\bigr),\qquad
 h_\square(c,g)=h_\square(c).
\]
Extend this with period $R$ and set it to zero away from the periodic
images of the lifted cube neighbourhood.
The cutoffs satisfy $\sum_\square h_\square^2=1$,
$\supp h_\square\subset\{|c-z_\square|_\infty<3M/4\}$, and
$|h_\square(c)-h_\square(c')|\le\pi|c-c'|_1/M$.
Put $C_\square=-[F_y(A),h_\square]$. Internal hops commute with
$h_\square$, and the penalty has range at most $L$. The phase bounds
and the same normalized sums used in the perturbation estimate give
\begin{equation}\label{eq:cutoff-commutator-bound}
 N_\gamma(C_\square)\le\pi C_V/M.
\end{equation}
The kernel is
$(h_\square(c)-h_\square(c'))F_y(A)(x,x')$. Its two photon sites lie
within $3M/4+L$ of $z_\square$. If $M>4L$, both lie in
$W_\square^-$; multiplication by $h_{\square'}$ can give a nonzero
operator only when $d(z_\square,z_{\square'})\le M$.

Using \eqref{eq:local-right-identity} and
\eqref{eq:local-left-identity} gives
\begin{align}
 F_yh_\square S_\square h_\square
       &=h_\square^2-C_\square S_\square h_\square,
                                   \label{eq:parametrix-right}\\
 h_\square S_\square h_\square F_y
       &=h_\square^2+h_\square S_\square C_\square.
                                   \label{eq:parametrix-left}
\end{align}
Let $S^{\mathrm p}=\sum_\square h_\square S_\square h_\square$ and
$R^{\mathrm p}=\sum_\square C_\square S_\square h_\square$.
At most $27$ cubes contribute to one row of $R^{\mathrm p}$.
Consequently,
\[
 N_\gamma(R^{\mathrm p})\le q_0:=\frac{54\pi C_Va'}M.
\]
Choose $M\ge108\pi C_Va'$. Equation~\eqref{eq:parametrix-right}
now gives a right inverse by a convergent series:
\begin{equation}\label{eq:global-parametrix-inverse}
 S_y(A)=F_y(A)^{-1}=S^{\mathrm p}\sum_{n\ge0}(R^{\mathrm p})^n.
\end{equation}
The finite-dimensional operator $F_y(A)$ is therefore invertible.
The left identity gives the second form with
$R^{\mathrm p\prime}=-\sum_\square h_\square S_\square C_\square$.

The terms in \eqref{eq:global-parametrix-inverse} have a local
description. A walk of length $n$ is a sequence
$\omega=(\square_0,\ldots,\square_n)$ with successive centres at
distance at most $M$. Its term is
\[
 S_\omega=(h_0S_0h_0)(C_1S_1h_1)\cdots(C_nS_nh_n),
\]
where the subscript $i$ denotes the cube $\square_i$.
Other sequences give zero because of the support of the commutator.
Its connection variables lie in
$\bigcup_{i=0}^nW_{\square_i}$.

\begin{proposition}\label{prop:walk-propagator-bounds}
For $p\ge1$ with $q(p)=54\pi C_Va'p/M\le1/2$,
\begin{align}
 \sup_z\sum_{z'}\exp(\gamma d(z,z'))
       \sum_{|\omega|=n}p^{n+1}
           \|\mathbf1_{\Delta_z}S_\omega\mathbf1_{\Delta_{z'}}\|
       &\le16a'pw_k(y)q(p)^n,\label{eq:walk-row-bound}\\
 \sup_z\sum_{z'}\exp(\gamma d(z,z'))
       \sum_{|\omega|=n}p^{n+1}
           \|\mathbf1_{\Delta_z}\delta_{\alpha,A}S_\omega
                            \mathbf1_{\Delta_{z'}}\|
       &\le152a'pq(p)^n.\label{eq:walk-holder-bound}
\end{align}
The norms are sup operator norms. At $k=0$ the first bound holds
also with the supremum of the normalized kernel over the two clusters.
In particular, $S_y(A)$ is analytic on $\widetilde{\mathcal G}_k$ and
\begin{equation}\label{eq:global-propagator-bounds}
 N_\gamma(S_y(A))\le32a'w_k(y),\qquad
 N_\gamma(\delta_{\alpha,A}S_y(A))\le304a'.
\end{equation}
These constants are independent of the level, the side, the running
mass in $[0,1]$, and the connection in the stated domain.
\end{proposition}

\begin{proof}
The first factor in a walk has norm at most $2a'w_k(y)$.
Every additional factor has norm at most $2\pi C_Va'/M$ by
\eqref{eq:local-inverse-bounds} and \eqref{eq:cutoff-commutator-bound}.
For a fixed output block, at most eight initial cutoffs contribute;
there are at most $27$ choices at each further step. Apply
submultiplicativity row by row and sum these $8\cdot27^n$ bounds.
This proves \eqref{eq:walk-row-bound}.

For the covariant quotient, only the factor at the left end of the
product is differentiated. The product identity is
\[
 \delta_{\alpha,A}(h_\square v)(x,x')
 =h_\square(c)\delta_{\alpha,A}v(x,x')
   +d_X^{-\alpha}(h_\square(c')-h_\square(c))U_A(x,x')v(x').
\]
The first contribution has localized norm at most $18a'$.
The second has coefficient at most $3\pi\exp(1)/M$, and evaluation
at the second point costs at most $27\exp(\gamma)$. With the stated
choice of $M$ this contribution is at most $a'$. Thus
$N_\gamma(\delta_{\alpha,A}h_\square S_\square)\le19a'$.
For intracluster pairs the second term vanishes, since $c=c'$.
There are again at most eight possible initial cubes, including when
the second point of a short pair enters the cutoff support. The count
$8\cdot27^n$ now gives \eqref{eq:walk-holder-bound}.

At level zero keep the kernel norm in the rightmost local inverse;
the other factors are controlled in operator norm. The same
submultiplicative row estimate proves the kernel assertion.
Summing the geometric bounds with $p=1$ gives
\eqref{eq:global-propagator-bounds}. For a fixed torus the sum of the
operator norms over all walks is finite; its tail of lengths at least
$j$ is bounded by $4a'(\ell/M)^3 2^{-j}$. Each walk term is analytic
in the finitely many connection variables it reads. The uniform tail
bound therefore proves analyticity of the sum. The local bounds in
\eqref{eq:walk-row-bound} contain no factor depending on $\ell$.
\end{proof}

The walk expansion provides both the inverse and the localization
needed by the block integration. Its geometric decay survives summing
over the intermediate cubes. The covariant Hölder estimate treats
the short pairs within one chamber cluster on the same footing as
pairs of distinct photon sites. These are the two bounds used to
control the chamber fields and the local fermion activities.


\section{Photon block integration}\label{sec:photon-primitives}

The photon block transformation uses a Landau representative for the
background and an axial representative for the fluctuation. The
background minimizer must have uniform local derivative bounds. The
fluctuation form must be uniformly positive on the axial slice, and its
inverse square root must decay between the local fluctuation coordinates.
These estimates also permit analytic weakening parameters to be inserted
without losing locality or the field domains of the renormalization step.

\subsection{Coarse averages and the Landau minimizer}
\label{subsec:photon-Landau-minimizer}

Let $L\ge5$ be odd, $\eta=\eta_k=L^{-k}$, and let $\ell$ be a
positive integer divisible by $L$. The photon fields are periodic
cochains on $T_k=\eta\Z^3/\ell\Z^3$, with the pairing
$\eta^3\sum_c$ on each component. Here the differences are normalized
by the mesh:
\[
 \partial_\mu f(x)=\eta^{-1}(f(x+\eta e_\mu)-f(x)),\qquad
 (d\lambda)_\mu=\partial_\mu\lambda,\qquad
 (d\mathscr A)_{\mu\nu}
       =\partial_\mu\mathscr A_\nu-\partial_\nu\mathscr A_\mu.
\]
We put $\delta=d^T$ and $\Delta_{\mathrm{sc}}=\delta d$ on scalars.
Fix $0<\alpha<1$. The H\"older quotient $\delta_\alpha\partial$
uses pairs of lattice sites at distance at most one, as in
\eqref{eq:photon-kernel-regularity}.
The unit lattice is $T_k^0=\Z^3/\ell\Z^3$ and the
unit blocks are $\Delta_z=\{c:|c-z|_\infty<1/2\}$.
Set $N=L^k$. The cochain averages are denoted by $\mathbf Q_k$;
on scalars, bonds, and plaquettes they are respectively
\begin{align}
 (\mathbf Q_k\lambda)(z)&=\eta^3\sum_{c\in\Delta_z}\lambda(c),\notag\\
 (\mathbf Q_k\mathscr A)_\mu(z)
    &=\eta^3\sum_{c\in\Delta_z}\sum_{j=0}^{N-1}
                                \eta\mathscr A_\mu(c+j\eta e_\mu),
                                     \label{eq:photon-block-averages}\\
 (\mathbf Q_k^{(2)}F)_{\mu\nu}(z)
    &=\eta^3\sum_{c\in\Delta_z}\sum_{i,j=0}^{N-1}
                    \eta^2F_{\mu\nu}(c+i\eta e_\mu+j\eta e_\nu).
\notag
\end{align}
All sums are on the periodic covering lattice. These averages preserve
constants and the mean of a bond field, are onto on bond fields, and
satisfy
\begin{equation}\label{eq:photon-averaging-identities}
 \mathbf Q_kd=d\mathbf Q_k,\qquad
 d\mathbf Q_k=\mathbf Q_k^{(2)}d,\qquad
 \|\mathbf Q_k^{(2)}F\|_2\le\|F\|_2.
\end{equation}
On scalars $\mathbf Q_k^T\lambda$ is constant with value $\lambda(z)$
on $\Delta_z$, and $\mathbf Q_k\mathbf Q_k^T=1$.
The differential identities follow by telescoping along the segments
in \eqref{eq:photon-block-averages}. The plaquette average has total
weight one. Jensen's inequality and summation over coarse sites count
each fine plaquette with total weight $\eta^3$, proving its contraction.
To see surjectivity on bonds, give the fine bond leaving a unit block
in direction $\mu$ the value $N A_\mu(z)$, and give the other bonds
value zero. Every segment of length one starting in $\Delta_z$ crosses
exactly one such boundary bond, and its average is $A_\mu(z)$.

Let $\mathcal R_k$ be the orthogonal projection on
$\Delta_{\mathrm{sc}}(\ker\mathbf Q_k)$ in the scalar space.
Denote by $H_k^{\mathrm{ph}}A$ the minimizer of
$\|d\mathscr A\|_2^2$ subject to
\begin{equation}\label{eq:photon-Landau-constraints}
 \mathbf Q_k\mathscr A=A,\qquad\mathcal R_k\delta\mathscr A=0.
\end{equation}
The superscript distinguishes it from the fermion field map $H_k(A)$.
Define the coarse photon form
\begin{equation}\label{eq:photon-coarse-form}
 \Delta_k=(dH_k^{\mathrm{ph}})^TdH_k^{\mathrm{ph}}.
\end{equation}

\begin{proposition}\label{prop:photon-Landau-minimizer}
The minimizer exists uniquely. It is real and linear, and
$\mathbf Q_kH_k^{\mathrm{ph}}=1$. It preserves constant fields;
if $dA=0$, then $dH_k^{\mathrm{ph}}A=0$. At $k=0$ it is the identity.
Its characterization is
\begin{equation}\label{eq:photon-minimizer-Euler}
 \mathbf Q_k\mathscr A=A,\qquad
 \delta d\mathscr A\in\operatorname{ran}\mathbf Q_k^T
     \text{ on bonds},\qquad
 \Delta_{\mathrm{sc}}\delta\mathscr A
      \in\operatorname{ran}\mathbf Q_k^T\text{ on scalars}.
\end{equation}
\end{proposition}

\begin{proof}
Surjectivity makes the averaging constraint nonempty. The possible
curvatures form an affine subspace of the finite lattice cochain space,
on which the squared norm has a minimum.
Differentiating along the kernel of $\mathbf Q_k$ shows that a
minimizer satisfies
\[
 \langle\delta d\mathscr A,B\rangle=0;
\]
this is the middle
condition in \eqref{eq:photon-minimizer-Euler}. The parallelogram
identity shows that all minimizers have the same curvature.

For a closed periodic bond field the discrete Hodge decomposition is
$B=c+d\lambda$, where $c$ is its mean. One can see this directly in
Fourier variables: at nonzero momentum, closedness says that the
components of $\widehat B$ are proportional to the forward-difference
symbol, and the proportionality factor is
$|\widehat d|^{-2}\widehat d^*\cdot\widehat B$. At zero momentum
only the constant remains. Exact fields have mean zero. Hence a closed
$B\in\ker\mathbf Q_k$ is $d\lambda$ with
$\lambda\in\ker\mathbf Q_k$: average the decomposition, use
\eqref{eq:photon-averaging-identities}, and subtract the constant value
of $\mathbf Q_k\lambda$. Thus the unconstrained minimizers are
$\mathscr A_1+d(\ker\mathbf Q_k)$.

The scalar Laplacian is injective on $\ker\mathbf Q_k$. Indeed,
$\Delta_{\mathrm{sc}}\lambda=0$ gives $d\lambda=0$, so $\lambda$
is constant and its average forces it to be zero. For
$\lambda\in\ker\mathbf Q_k$,
\[
 \mathcal R_k\delta(\mathscr A_1+d\lambda)
       =\mathcal R_k\delta\mathscr A_1+\Delta_{\mathrm{sc}}\lambda.
\]
Since the first term lies in the range
$\Delta_{\mathrm{sc}}(\ker\mathbf Q_k)$, exactly one such gauge
correction imposes the Landau constraint. Moreover
$\mathcal R_kf=0$ is equivalent to
$\Delta_{\mathrm{sc}}f\perp\ker\mathbf Q_k$; this proves the last
condition of \eqref{eq:photon-minimizer-Euler}.
All these equations are real and linear, so uniqueness proves the
reality and linearity of $H_k^{\mathrm{ph}}$.

Constants satisfy the equations. If $A=c+d\lambda$ is closed, then
$c+d\mathbf Q_k^T\lambda$ is a closed fine field with average $A$,
so the minimum curvature energy is zero. This proves preservation of
closed fields. Finally $\mathbf Q_0=1$ gives the level-zero assertion.
\end{proof}

\subsection{Axial coercivity on a block}\label{subsec:photon-axial-coercivity}

On the unit lattice, let
$B(y)=\{z\in\Z^3:|z-y|_\infty<L/2\}$, with centres in $L\Z^3$.
The one-step bond average and axial functional are
\[
 (\mathbf QA)(y,y+Le_\mu)
      =L^{-4}\sum_{z\in B(y)}A(\Gamma_{z,z+Le_\mu}),\qquad
 (\tau A)(y,z)=\tfrac16\sum_\pi A(\Gamma^\pi(y,z)).
\]
Here the first path is straight, and $\Gamma^\pi$ changes the
coordinates in the order $\pi$. The sums $A(\Gamma)$ are oriented
bond sums. For an interior bond $b=(z,z+e_\nu)$, subtract the two
paths from $y$ to its endpoints and append $b$. For each order this
loop bounds the strips belonging to the directions that come after
$\nu$. Their total number of plaquettes is
$\sum_{\mu\text{ after }\nu}|z_\mu-y_\mu|$. If $\tau A=0$,
averaging over the six orders gives
\begin{equation}\label{eq:axial-interior-bond}
 |A(b)|\le\tfrac12\sum_{\mu\ne\nu}|z_\mu-y_\mu|
           \sup_{p\subset B(y)}|dA(p)|
       \le\tfrac{L-1}{2}\sup_{p\subset B(y)}|dA(p)|.
\end{equation}

For a face between $B(y)$ and $B(y+Le_\nu)$, let $a(z)$ be its face
bond values and let $p_c$ be the maximum curvature on plaquettes in
these two blocks and across their common face. A transverse face step
is one plaquette plus the difference of two interior bond values.
Thus $|a(z)-a(z')|\le L|z-z'|_1p_c$.
The defining average $\mathbf QA$ counts each face bond $L$ times
and $L^3(L-1)$ interior bond occurrences. Therefore its face mean
$\bar a$ satisfies
\[
 |\bar a|\le L|(\mathbf QA)(c)|
                       +\tfrac12(L-1)^2p_c.
\]
The average transverse distance from a face site is at most $L-1$,
so $|a(z)-\bar a|\le L(L-1)p_c$. Combining these gives
\begin{equation}\label{eq:axial-local-bond-bound}
 |A(b)|\le\tfrac32L^2p_c+L|(\mathbf QA)(c)|.
\end{equation}
For interior bonds use \eqref{eq:axial-interior-bond} and omit the
coarse-average term. In particular, on the axial fluctuation space,
\begin{equation}\label{eq:axial-fluctuation-coercivity}
 \mathbf QZ=0,\quad\tau Z=0
     \quad\Longrightarrow\quad
 \|Z\|_\infty\le\tfrac32L^2\|dZ\|_\infty,\qquad
 \|Z\|_2^2\le9L^7\|dZ\|_2^2.
\end{equation}
For the second inequality, square the local bound and replace each
local curvature maximum by its sum of squares. An interior plaquette
belongs to the neighbourhoods of at most $3L^3+3L^2\le4L^3$ bonds;
a face plaquette to those of at most $L^2$ bonds. Summation gives
$(9/4)L^4\cdot4L^3=9L^7$.

\subsection{The minimizer symbol and uniform regularity}
\label{subsec:photon-symbol-regularity}

The Fourier calculation must separate the zero alias before the
continuum mesh limit is taken. Set $h=(N-1)/2$ and, for $s\in\C$,
\[
 u(s)=N^{-1}\sum_{j=-h}^he^{isj/N},\qquad
 s_N(s)=2N\sin(s/(2N)),\qquad c(s)=2\sin(s/2).
\]
The identity $u(s)s_N(s)=c(s)$ follows by telescoping the geometric
sum. For $q\in\C^3$ define
\[
 d_\mu(q)=N(e^{iq_\mu/N}-1),\quad
 d_\mu^*(q)=N(e^{-iq_\mu/N}-1),\quad
 D(q)=\sum_\mu s_N(q_\mu)^2,\quad U(q)=\prod_\mu u(q_\mu),
\]
and
\[
 V_\mu(q)=e^{i(1-N^{-1})q_\mu/2}u(q_\mu),\quad
 V_\mu^*(q)=V_\mu(-q),\quad
 \rho_\mu(p)=e^{ip_\mu}-1,\quad\rho_\mu^*(p)=e^{-ip_\mu}-1.
\]
The stars in these formulas denote analytic symbols. At real arguments
they are complex conjugates. The aliases are $q_m=p+2\pi m$,
$m\in\{-h,\ldots,h\}^3$, and
\begin{equation}\label{eq:photon-alias-identities}
 V_\mu(q_m)d_\mu(q_m)=\rho_\mu(p),\qquad
 V_\mu^*(q_m)d_\mu^*(q_m)=\rho_\mu^*(p),\qquad
 V_\mu V_\mu^*=u(q_{m,\mu})^2.
\end{equation}
The scalar averaging symbol is $U(q_m)$; on bonds it is
$U(q_m)V_\mu(q_m)$. Thus the Euler equations
\eqref{eq:photon-minimizer-Euler} can be solved separately in each
alias class.

Write $D_0=D(p)$, $U_0=U(p)$, $u_\mu=u(p_\mu)$, and
\begin{align*}
 X_\mu&=\sum_{m\ne0}\frac{U(q_m)^2u(q_{m,\mu})^2}{D(q_m)},&
 Y&=\sum_{m\ne0}\frac{U(q_m)^2}{D(q_m)^2},\\
 N_\mu&=U_0^2u_\mu^2+D_0X_\mu,&
 N_G&=U_0^2+D_0^2Y,\\
 \mathcal N&=\prod_\mu N_\mu,&
 \Pi&=\mathcal N-\sum_\mu s_N(p_\mu)^2X_\mu
                                      \prod_{\nu\ne\mu}N_\nu.
\end{align*}
The coarse form symbol and the scalar multiplier for the Landau
condition are
\begin{align}
 \Lambda_{\mu\nu}
     &=\delta_{\mu\nu}\frac{D_0}{N_\mu}
       -\frac{\rho_\mu\rho_\nu^*U_0^2\mathcal N}
                    {N_\mu N_\nu\Pi},\notag\\
 \omega_\nu&=\frac{\rho_\nu^*D_0U_0^2\mathcal N}{N_\nu N_G\Pi},
                                                    \label{eq:photon-Lambda-symbol}\\
 h_{\mu\nu}(p,m)
    &=\frac{U(q_m)V_\mu^*(q_m)\Lambda_{\mu\nu}}{D(q_m)}
          +\frac{d_\mu(q_m)U(q_m)\omega_\nu}{D(q_m)^2}
                                      \quad(m\ne0).\notag
\end{align}
The zero-alias formula has its apparent pole cancelled:
\begin{equation}\label{eq:photon-zero-alias}
 h_{\mu\nu}(p,0)=\delta_{\mu\nu}\frac{U_0V_\mu^*(p)}{N_\mu}
       +\frac{d_\mu(p)\rho_\nu^*U_0^3
                 (X_\mu-u_\mu^2D_0Y)\mathcal N}
                  {N_\mu N_\nu N_G\Pi}.
\end{equation}
Indeed, $N_\mu-u_\mu^2N_G=D_0(X_\mu-u_\mu^2D_0Y)$, so
\eqref{eq:photon-zero-alias} equals the preceding formula wherever
$D_0\ne0$.

For real $p\in[-\pi,\pi]^3$, the nonzero aliases have $D(q_m)>0$,
and $X_\mu,Y\ge0$. The positivity of $u_\mu$ gives $N_\mu,N_G>0$.
If $p\ne0$, the identity
\begin{equation}\label{eq:photon-Pi-positive}
 \Pi=\sum_\mu\frac{s_N(p_\mu)^2}{D_0}
                 U_0^2u_\mu^2\prod_{\nu\ne\mu}N_\nu
\end{equation}
shows that $\Pi>0$; at zero $N_\mu=N_G=\Pi=1$.
The algebraic identity
$U_0^2\sum_\mu\rho_\mu\rho_\mu^*/N_\mu=D_0\Pi/\mathcal N$
implies $\sum_\mu\rho_\mu^*\Lambda_{\mu\nu}=0$.
Substitution in the alias equations now gives
\[
 \widehat{(H_k^{\mathrm{ph}}A)_\mu}(q_m)
      =\sum_\nu h_{\mu\nu}(p,m)\widehat A_\nu(p),\qquad
 \widehat{(\Delta_kA)_\mu}(p)
      =\sum_\nu\Lambda_{\mu\nu}(p)\widehat A_\nu(p).
\]
To check the first formula, average the displayed alias expression.
The sums including $m=0$ are $N_\mu/D_0$ and $N_G/D_0^2$.
Its two terms give $\widehat A_\mu$ plus opposite multiples of
$\rho_\mu$, so the average is $A$. Its divergence is
$U(q_m)\widehat\omega(p)/D(q_m)$; applying the scalar Laplacian
gives the scalar range condition, and subtracting its gradient from
$D(q_m)\widehat{H_k^{\mathrm{ph}}A}$ gives
$U(q_m)V_\mu^*(q_m)\widehat\Lambda_\mu(p)$.
These are the other two Euler conditions. At $p=0$ only the zero
alias survives, giving the constant field. Uniqueness of the real
minimizer proves the formula. Pairing its bond Euler equation with
$H_k^{\mathrm{ph}}A'$ and using $\mathbf Q_kH_k^{\mathrm{ph}}=1$
identifies the symbol of $\Delta_k$ with $\Lambda$.

There is a strip of fixed width, independent of $N$, in which these
symbols are analytic. Here is the estimate that provides that width
and the derivative bounds. On
$|\operatorname{Re}p_\mu|<5\pi/4$,
$|\operatorname{Im}p_\mu|<1/4$, put
$a=(8/(5\pi))\sin(5\pi/8)$ and $c_h=\cosh(1/8)$.
The elementary complex sine identities give
\[
 \operatorname{Re}s_N(s)^2\ge
       a^2(\operatorname{Re}s)^2-c_h^2(\operatorname{Im}s)^2,
 \qquad |u(s)|\le c_h
\]
in the required alias range. Consequently, for $m\ne0$,
\begin{equation}\label{eq:photon-alias-estimates}
 \operatorname{Re}D(q_m)\ge|m|_2^2,\qquad
 |u(q_{m,\mu})|\le C\langle m_\mu\rangle^{-1},\qquad
 |d_\mu(q_m)|\le C\langle m_\mu\rangle,
\end{equation}
where $\langle t\rangle=\max(1,|t|)$.
For the first estimate use
$|\operatorname{Re}q_{m,\mu}|\ge(3\pi/4)|m_\mu|$ when
$m_\mu\ne0$; the lower bound before simplification is
$a^2(3\pi/4)^2|m|_2^2-3c_h^2/16\ge|m|_2^2$.
The bound on $u$ follows from $u=c/s_N$, and the bound on $d_\mu$
from its sine expression. Thus $X_\mu,Y,N_\mu,N_G,\Pi$ are
uniformly bounded on the larger rectangle, since their alias sums
are bounded by constant multiples of
$\sum_m\prod_\mu\langle m_\mu\rangle^{-2}$.

On its real cube $u_\mu\ge a$, so
$N_\mu\ge a^8$, $N_G\ge a^6$, and
\eqref{eq:photon-Pi-positive} gives $\Pi\ge a^{24}$.
If an analytic function is bounded by $B$ on that rectangle and is
at least $c>0$ on the slightly smaller real cube, apply Schwarz's
lemma to $t\mapsto f(p'+itp'')-f(p')$ on the disk
$|t|<4B/c$. For $|p''|_\infty<c/(16B)$ its argument stays in the
larger rectangle, and $|f(p)-f(p')|\le c/2$.
Apply this to the three denominators. It gives a fixed width $g_2>0$
on $|\operatorname{Re}p_\mu|<5\pi/4-1/4$ where all remain
bounded away from zero. In particular, the zero-alias formula is
analytic there.

In this smaller strip the alias bounds imply
\begin{equation}\label{eq:photon-h-alias-bound}
 |h_{\mu\nu}(p,m)|\le C\left(
       \langle m_\mu\rangle^{-1}|m|_2^{-2}
          +\langle m_\mu\rangle|m|_2^{-4}\right)
                     \prod_\sigma\langle m_\sigma\rangle^{-1}
                 \quad(m\ne0),
\end{equation}
and $|h(p,0)|\le C$. For every $0\le\alpha<1$ this proves
\begin{equation}\label{eq:photon-alias-Holder-sum}
 \sup_p\sum_m|h_{\mu\nu}(p,m)|
       \max\{1,|d_\lambda(q_m)|,
                      |d_\lambda(q_m)||q_m|_1^\alpha\}\le S_\alpha,
\end{equation}
with $S_\alpha$ independent of $N$. In fact, after multiplication
by $\langle m\rangle_\infty^{1+\alpha}$, the first term in
\eqref{eq:photon-h-alias-bound} is bounded by
\[
 C\langle m_\mu\rangle^{-2}
            \prod_{\sigma\ne\mu}\langle m_\sigma\rangle^{-(3-\alpha)/2},
\]
and the second by $C\prod_\sigma\langle m_\sigma\rangle^{-4/3}$.
Both product sums converge; the first is where $\alpha<1$ is used.

For $x\in\eta\Z^3$, set
$K_{\mu\nu}(p;x)=\sum_me^{iq_m\cdot x}h_{\mu\nu}(p,m)$.
It is $2\pi$-periodic in $p$. To check periodicity, away from $D_0=0$
write $E_\mu=N_\mu/D_0$ and $G=N_G/D_0^2$ as sums over all aliases,
and $r=\sum_\mu\rho_\mu\rho_\mu^*/E_\mu$. Then
\[
 \Lambda_{\mu\nu}=\delta_{\mu\nu}/E_\mu
                  -\rho_\mu\rho_\nu^*/(E_\mu E_\nu r),\qquad
 \omega_\nu=\rho_\nu^*/(E_\nu Gr).
\]
Translation of $p$ by $2\pi$ permutes the alias classes modulo
$2\pi N\Z^3$, and every summand, including $e^{iq_m\cdot x}$,
is periodic modulo this latter lattice. Thus $K$ and $\Lambda$
agree on real overlaps of the rectangles; analyticity extends the
identity and defines them on the full strip. Define the covering
kernel by
$H_{\mu\nu}^\infty(x)=(2\pi)^{-3}\int_{[-\pi,\pi]^3}K_{\mu\nu}(p;x)\dd p$.
Shift the contour to $p+i\gamma_H\operatorname{sgn}(x)$, with
$\gamma_H=g_2/2$. The vertical faces cancel by periodicity, and
\eqref{eq:photon-alias-Holder-sum} gives
\[
 |H^\infty(x)|, |\partial H^\infty(x)|,
 |z|^{-\alpha}|\partial H^\infty(x+z)-\partial H^\infty(x)|
          \le C S_\alpha e^{-\gamma_H|x|_1},
 \qquad x,z\in\eta\Z^3,\quad0<|z|\le1.
\]
The last estimate uses
$|e^{iq\cdot z}-1|\le2.5|z|^\alpha|q|_1^\alpha$ on the shifted
contour. The torus kernel is
$H_k^{\mathrm{ph}}(b,b')=\sum_{n\in\Z^3}H^\infty(c-z'+\ell n)$,
with the corresponding direction indices. Summing the exponential
bound over the lifts proves
\begin{equation}\label{eq:photon-kernel-regularity}
 |H_k^{\mathrm{ph}}(b,b')|, |\partial H_k^{\mathrm{ph}}(b,b')|,
 |\delta_\alpha\partial H_k^{\mathrm{ph}}((b,\tilde b);b')|
                 \le C_H(\alpha)e^{-\gamma_Hd(b,b')}.
\end{equation}
Here $b,\tilde b$ have the same direction and their initial points
are at distance at most one. The same contour argument for the bounded
symbol $\Lambda$ gives
\begin{equation}\label{eq:photon-form-kernel-decay}
 |\Delta_k(b,b')|\le C_\Delta e^{-\gamma_Hd(b,b')},\qquad
 \|\Delta_k\|_{2\to2}\le C_\Delta^{(2)}.
\end{equation}
All three constants are independent of $k$ and $\ell$.
For the operator norm, Parseval's formula reduces the estimate to
$\sup_p\|\Lambda(p)\|$. A uniform bound $A_1$ for its nine entries
gives $\|\Lambda(p)\|\le3A_1$, so one may take
$C_\Delta^{(2)}=3A_1$.

\subsection{Fluctuation coordinates and the positive square root}
\label{subsec:photon-fluctuation-root}

The fluctuations in one block integration satisfy
$\mathbf Q Z=0$ and $\tau Z=0$. We use coordinates in which the
Gaussian measure becomes a product measure, while each coordinate
still has bounded support.

For an $L$-block $B(y)$, let $\mathrm{Bd}(y)$ denote its interior
bond functions, with the counting inner product. The functional
$\tau_y$ has rank $L^3-1$: given values $f(x)$ for $x\ne y$, extend
$f(y)=0$ and apply $\tau_y$ to the interior restriction of $df$.
Every path sum telescopes to $f(x)$. Thus
\[
 \dim\ker\tau_y=3L^2(L-1)-(L^3-1)=2L^3-3L^2+1.
\]
The three coordinate reflections preserve $\ker\tau_0$ and act as
commuting orthogonal involutions. Choose an orthonormal basis
$e_i^0$ of simultaneous eigenvectors and translate it to $e_i^y$.
This choice depends only on $L$.

Each face between neighboring blocks has a central bond, on the
line joining their centers. Let
\[
 V_k=\bigoplus_y\ker\tau_y\ \oplus\!\!\!
       \bigoplus_{b\text{ noncentral face bond}}\R\delta_b,
 \qquad
 \{\Upsilon\}=\{e_i^y\}\cup\{\delta_b\}.
\]
Here $\delta_b$ is the point mass on the positively oriented bond
$b$. These are orthonormal coordinates, and
\begin{equation}\label{eq:photon-fluctuation-dimension}
 \dim V_k=d_V\ell^3,\qquad d_V=2-2L^{-3}.
\end{equation}
The support of a basis vector lies in one $L$-block or consists of
the two ends of a bond.

For $z\in V_k$ define $Cz$ to agree with $z$ off the central bonds,
and set, on the central bond $b_c$ of the face
$c=\langle y,y+Le_\mu\rangle$,
\begin{equation}\label{eq:photon-central-bond-elimination}
 (Cz)(b_c)=-\sum_{\substack{b\text{ on }c\\b\ne b_c}}z(b)
                -\sum_{b\text{ interior}}\frac{n_c(b)}Lz(b).
\end{equation}
The second sum is over the bonds of direction $\mu$ in the two
adjacent blocks; $n_c(b)$ counts the straight averaging paths that
contain $b$. Numbering the bonds on such a line by
$t=0,\ldots,2L-2$ gives
$n_c(b)=\min(t+1,2L-1-t)$. The face bond has multiplicity $L$;
every interior bond has multiplicity at most $L-1$. Consequently
\[
 L^4(\mathbf QZ)(c)=L\sum_{b\text{ on }c}Z(b)
                            +\sum_b n_c(b)Z(b).
\]
Only $Z(b_c)$ is unknown in the equation for this face. Central
bonds occur in no axial path and in no other face equation.
It follows that $C$ is a real linear bijection
\begin{equation}\label{eq:photon-coordinate-bijection}
 C:V_k\longrightarrow
       \{Z:\mathbf QZ=0,\ \tau Z=0\}.
\end{equation}

The following coordinate estimates will be used both for the
covariance and for the small-field domain:
\begin{equation}\label{eq:photon-coordinate-estimates}
 \|Cz\|_\infty\le c_C\sup_\Upsilon|z(\Upsilon)|,
 \qquad
 \|z\|_2\le\|Cz\|_2\le(5L^3)^{1/2}\|z\|_2,
 \qquad
 \sum_b|(C\Upsilon)(b)|\le c_1,
\end{equation}
where $c_C=3L^{9/2}$ and $c_1=14L^{3/2}$. To prove the first
bound, normalize the coordinate supremum to one. An interior value
is at most $(2L^3)^{1/2}$ by Cauchy--Schwarz, and a noncentral face
value is at most one. Formula
\eqref{eq:photon-central-bond-elimination} bounds a central value
by $L^2+2L^3(2L^3)^{1/2}\le3L^{9/2}$.
For the second bound, the new central values are supported away
from $z$. A central row has at most $2L^3$ entries, of modulus at
most one, and an entry contributes to at most two such rows.
Cauchy--Schwarz therefore adds at most $4L^3\|z\|_2^2$.
For a face coordinate, $C\delta_b$ has two entries of modulus one.
For an interior unit vector, the interior $\ell^1$ norm is at most
$(3L^3)^{1/2}$, and the total contribution of the six adjacent
central bonds is at most six times this norm. This proves the
last estimate with $c_1=14L^{3/2}$.
The same formula gives
\begin{equation}\label{eq:photon-coordinate-support}
 d(x,x')\le L\quad
 (x\in\supp C\Upsilon,\ x'\in\supp\Upsilon).
\end{equation}

Coordinate sign changes and translations by $M\Z^3$ permute
$\{\pm\Upsilon\}$. For the interior vectors this follows from the
simultaneous eigenvector choice; for face coordinates it follows
from the preservation of central bonds. Hence the standard Gaussian
$\mu_I$ on $V_k$ and its normalized restriction $\widehat\mu_k$ to
\[
 \mathrm{Box}_k=
 \{z:\ |z(\Upsilon)|\le p_{0,k}\text{ for every }\Upsilon\}
\]
are invariant under these transformations and under $z\mapsto-z$.
Independence of the coordinates gives
\begin{equation}\label{eq:photon-Gaussian-box-mass}
 \mu_I(\mathrm{Box}_k)=\exp(-\ell^3\delta\varepsilon^0_k),
 \qquad
 \delta\varepsilon^0_k=-d_V\log\Pr(|g|\le p_{0,k})
            \le12\exp(-p_{0,k}^2/2),
\end{equation}
when $p_{0,k}\ge1$, where $g$ is standard normal.
Indeed, the Gaussian tail is at most
\[
 2(2\pi)^{-1/2}p_{0,k}^{-1}\exp(-p_{0,k}^2/2);
\]
at
$p_{0,k}=1$ this is less than $1/2$, and
$-\log(1-t)\le2t$ for $t\le1/2$ suffices.

Set
\[
 \widetilde\Delta_k=C^T\Delta_k C,
 \qquad C_k=\widetilde\Delta_k^{-1}.
\]
The first of these definitions makes sense before invertibility
has been proved. From the axial estimate, the contraction of the
plaquette average, and $\mathbf Q_kH_k^{\mathrm{ph}}=1$, we get
\[
 \|z\|_2^2\le\|Cz\|_2^2
 \le9L^7\|dCz\|_2^2
 =9L^7\|\mathbf Q_k^{(2)}dH_k^{\mathrm{ph}}Cz\|_2^2
 \le9L^7\langle z,\widetilde\Delta_kz\rangle.
\]
The upper bound follows from
\eqref{eq:photon-form-kernel-decay} and
\eqref{eq:photon-coordinate-estimates}. Thus, with
\[
 c_P=\max\{9L^7,\ 5L^3 C_\Delta^{(2)}\},
 \qquad
 C_{\mathrm{Dt}}=c_1^2C_\Delta\exp(2\gamma_HL),
\]
we have
\begin{equation}\label{eq:photon-fluctuation-form-estimates}
 c_P^{-1}\le\widetilde\Delta_k\le c_P,
 \qquad
 |\widetilde\Delta_k(\Upsilon,\Upsilon')|
       \le C_{\mathrm{Dt}}\exp(-\gamma_Hd(\Upsilon,\Upsilon')).
\end{equation}
For the second estimate, expand
$\langle C\Upsilon,\Delta_kC\Upsilon'\rangle$ as a double bond
sum. Every contributing pair satisfies
$d(\Upsilon,\Upsilon')\le d(b,b')+2L$ by
\eqref{eq:photon-coordinate-support}; the two $\ell^1$ bounds
then give the asserted constant.

The distance between supports has a triangle defect at most $L$.
We record the weighted resolvent argument with this defect, since
it is responsible for the uniform decay of the square root.
Suppose that $T\ge t_0>0$ is a real symmetric matrix, that
$|T(a,b)|\le C_T\exp(-g d(a,b))$, and that
\[
 d(a,c)\le d(a,b)+d(b,c)+\phi.
\]
Choose
\[
 S\ge\exp(g\phi/2)\sup_a\sum_b(d(a,b)+\phi)
                                      \exp(-g d(a,b)/2),
 \qquad
 \theta=\min\{g/2,\ t_0/(2C_TS)\}.
\]
Fix $c$ and conjugate $T$ by the diagonal matrix
$W(a,a)=\exp(\theta d(a,c))$. The error $E=WTW^{-1}-T$ satisfies
\[
 |E(a,b)|\le\theta C_T\exp(g\phi/2)(d(a,b)+\phi)
                                      \exp(-g d(a,b)/2),
\]
because $|d(a,c)-d(b,c)|\le d(a,b)+\phi$ and
$|\exp(u)-1|\le|u|\exp(|u|)$. The Schur test gives
$\|E\|\le t_0/2$. Therefore
$\|(T+x+E)^{-1}\|\le(t_0/2+x)^{-1}$ for $x\ge0$.
Undoing the conjugation yields
\[
 |(T+x)^{-1}(a,c)|\le(t_0/2+x)^{-1}\exp(-\theta d(a,c)).
\]
The scalar integral
$\pi^{-1}\int_0^\infty x^{-1/2}(t+x)^{-1}\,\dd x=t^{-1/2}$
and finite-dimensional diagonalization now give
\begin{equation}\label{eq:photon-root-integral}
 T^{-1/2}=\frac1\pi\int_0^\infty x^{-1/2}(T+x)^{-1}\,\dd x,
 \qquad
 |T^{-1/2}(a,c)|\le(2/t_0)^{1/2}\exp(-\theta d(a,c)).
\end{equation}

Apply this to $T=\widetilde\Delta_k$, $\phi=L$,
$t_0=c_P^{-1}$ and $g=\gamma_H$. The basis support count gives
\[
 \sup_\Upsilon\sum_{\Upsilon'}\exp(-g d(\Upsilon,\Upsilon'))
       \le c_{\mathrm{bas}}(g),\qquad
 c_{\mathrm{bas}}(g)=2L^3\exp(2gL)(1+6/g)^3.
\]
For example, assign a basis vector to its block center or to an
end of its face bond. At most $2L^3$ vectors are assigned to any
site, and replacing support distance by point distance costs at
most $2L$; the remaining lattice sum is bounded by $(1+6/g)^3$.
Since $t\exp(-gt/4)\le2/g$, we may take
\[
 S_P=\exp(\gamma_HL/2)(L+2/\gamma_H)
                          c_{\mathrm{bas}}(\gamma_H/4),
 \quad
 \gamma_P=\min\{\gamma_H/2,\ (2c_PC_{\mathrm{Dt}}S_P)^{-1}\}.
\]
With
\[
 C_P^{\mathrm{ph}}=\max\{C_H(\alpha),C_{\mathrm{Dt}},(2c_P)^{1/2},4\exp(2)\},
\]
all the kernels in \eqref{eq:photon-kernel-regularity},
\eqref{eq:photon-fluctuation-form-estimates}, and the positive
square root satisfy their respective bounds at rate $\gamma_P$
with constant $C_P^{\mathrm{ph}}$. In particular,
\begin{equation}\label{eq:photon-square-root-decay}
 C_k^{1/2}=\widetilde\Delta_k^{-1/2},\qquad
 |C_k^{1/2}(\Upsilon,\Upsilon')|
                  \le C_P^{\mathrm{ph}}\exp(-\gamma_Pd(\Upsilon,\Upsilon')).
\end{equation}
At $k=0$, $H_0^{\mathrm{ph}}=1$; its first differences and their
H\"older differences are supported within distance two and have
entries of modulus at most two. Our choice $C_P^{\mathrm{ph}}\ge4\exp(2)$
therefore includes this level as well.

\subsection{The exact block decomposition}
\label{subsec:photon-block-decomposition}

For a real unit bond field $A'$ on the torus of side $\ell/L$,
put $A'_L(b)=L^{-1/2}A'(b/L)$. Let $A_k^{\min}(A')$ minimize
$\langle A,\Delta_kA\rangle$ under
\[
 \mathbf QA=A'_L,\qquad \tau A=0.
\]
This affine space is nonempty: set each face bond of the face $c$
equal to $L A'_L(c)$, and set every interior bond to zero.
Its direction is $C(V_k)$, and
\eqref{eq:photon-fluctuation-form-estimates} makes the quadratic
form positive definite on this direction. The minimizer is
therefore unique and is characterized by
\begin{equation}\label{eq:photon-block-minimizer-orthogonality}
 \langle A_k^{\min}(A'),\Delta_kCz\rangle=0
                    \quad(z\in V_k).
\end{equation}
This characterization is a real linear system with a unique
solution, so $A_k^{\min}$ is real and linear.

Define $\mathrm{Sq}_k=CC_k^{1/2}$. It maps $V_k$ bijectively onto
the axial fluctuations, has zero block average and mean, and
\[
 \langle\mathrm{Sq}_k\zeta,\Delta_k\mathrm{Sq}_k\zeta'\rangle
  =\langle\zeta,C_k^{1/2}\widetilde\Delta_kC_k^{1/2}\zeta'\rangle
  =\langle\zeta,\zeta'\rangle.
\]
Thus the exact Gaussian decomposition is
\begin{equation}\label{eq:photon-Gaussian-block-decomposition}
 \langle A_k^{\min}(A')+\mathrm{Sq}_k\zeta,
           \Delta_k(A_k^{\min}(A')+\mathrm{Sq}_k\zeta)\rangle
       =\langle A_k^{\min}(A'),\Delta_kA_k^{\min}(A')\rangle
                                                   +\|\zeta\|_2^2.
\end{equation}
The block integration uses the standard Gaussian in these
$\zeta$ coordinates, restricted to $\mathrm{Box}_k$; the loss of
normalization is exactly \eqref{eq:photon-Gaussian-box-mass}.

We also need the relation between the two backgrounds. With the
scalar scaling $\lambda_L(x)=L^{1/2}\lambda(x/L)$, there is a
unique real scalar $\omega_k(A')$ with
$\mathbf Q_{k+1}\omega_k(A')=0$ such that
\begin{equation}\label{eq:photon-background-one-step}
 H_k^{\mathrm{ph}}A_k^{\min}(A')
       =(H_{k+1}^{\mathrm{ph}}A')_L-d\omega_k(A'),
 \qquad
 \|dH_k^{\mathrm{ph}}A_k^{\min}(A')\|_2
       =\|dH_{k+1}^{\mathrm{ph}}A'\|_2.
\end{equation}
The harmonic part of $A'$ is arbitrary.
Here $\mathbf Q_{k+1}=\mathbf Q\mathbf Q_k$ is the average from
the mesh-$\eta_k$ lattice of side $\ell$ to its $L$-spaced
lattice. The operator $H_{k+1}^{\mathrm{ph}}$ is formed at mesh
$\eta_{k+1}$ and side $\ell/L$, before scaling by $L$.

Here is the minimizing argument. At mesh $\eta$, the scaled field
$(H_{k+1}^{\mathrm{ph}}A')_L$ minimizes the curvature norm under
the $(k+1)$-step constraint $\mathbf Q_{k+1}\mathscr A=A'_L$.
This follows by scaling the defining problem; the curvature
$L^2$ norm is unchanged. The field
$\mathscr B=H_k^{\mathrm{ph}}A_k^{\min}(A')$ has the same average.
If $\mathscr A$ is any competitor, set $A=\mathbf Q_k\mathscr A$.
Choose the one-step restricted gauge
\[
 \mu(x)=(\tau A)(y,x)
            -L^{-3}\sum_{x'\in B(y)}(\tau A)(y,x'),
                 \qquad x\in B(y).
\]
Then $\mathbf Q\mu=0$ and $\tau(A-d\mu)=0$, by telescoping
the exact-field path sums. Therefore
\[
 \|d\mathscr A\|_2^2
 \ge\|dH_k^{\mathrm{ph}}A\|_2^2
 =\|dH_k^{\mathrm{ph}}(A-d\mu)\|_2^2
 \ge\|d\mathscr B\|_2^2.
\]
Both backgrounds are minimizers of the same problem. The
Landau-minimizer argument shows that their difference is the
coboundary of a scalar with zero $(k+1)$-step average. This scalar
is unique: a closed field with zero average is exact, and the
only constant scalar with zero block average is zero. The two
backgrounds are real and linear in $A'$, hence so is $\omega_k$.
This proves \eqref{eq:photon-background-one-step} and identifies
the first term on the right of
\eqref{eq:photon-Gaussian-block-decomposition} with
$\|dH_{k+1}^{\mathrm{ph}}A'\|_2^2$.

\subsection{Curvature control in local gauges}
\label{subsec:photon-local-gauge}

Assume from now on that $\ell/(LM)$ is odd and at least $13$.
For an $M$-cube $\square$, let $\square^\dagger$ be its enlargement
by $LM$ in each direction. Its side is $s_d=(2L+1)M\le3\ell/13$.
For a bond field $B$ on a set $W$, write
\[
 \|B\|_{1,W}=\|B\|_{\infty,W},\qquad
 \|B\|_{2,W}=\|\partial B\|_{\infty,W},\qquad
 \|B\|_{3,W}=\|\delta_\alpha\partial B\|_{\infty,W}.
\]
Only bonds and differences contained in $W$ enter these norms.
They are ordinary differences of the photon coefficients; no
connection occurs in their definition.

\begin{proposition}\label{prop:photon-curvature-local-regularity}
For every real unit bond field $A$, including an arbitrary
constant part, and every $M$-cube $\square$, there is a real scalar
$\nu$ on $\square^\dagger$ such that
\begin{equation}\label{eq:photon-curvature-local-regularity}
 \|H_k^{\mathrm{ph}}A-d\nu\|_{i,\square^\dagger}
                         \le C_B\|dA\|_\infty,
                         \qquad i=1,2,3.
\end{equation}
The constant $C_B$ depends on $L,M,\alpha$ and is independent of
$k,\ell$ and the harmonic part of $A$.
\end{proposition}

\begin{proof}
Let $a$ be the mean of $A$, put $A^0=A-a$, and choose the center
$x_0$ of $\square$ as base point. On the canonical lifted period
cell
\[
 W=x_0+\{-(\ell-1)/2,\ldots,(\ell-1)/2\}^3
\]
define $\lambda(x)$ to be the line sum of $A^0$ along the
rectilinear path from $x_0$ to $x$, in coordinate order $1,2,3$.
This defines a scalar on the torus. Put $A''=A^0-d\lambda$ and
$p_A=\|dA\|_\infty$. The path sums of $A''$ vanish. For a bond
whose two ends lie in $W$ without crossing its boundary,
the strip identity used in the axial estimate gives
\begin{equation}\label{eq:photon-inner-basepoint-bound}
 |A''(x,x+e_\mu)|\le\sum_{\nu>\mu}|x_\nu-x_{0,\nu}|p_A
                                      \le2d(x,x_0)p_A.
\end{equation}
Indeed, commuting the final bond through the later portions of
the two base-point paths sweeps one strip in each direction
$\nu>\mu$. Each strip contributes its plaquette sum, and both
remaining path sums are zero.

A seam bond lies on a closed coordinate line with $\ell-1$
inner bonds. Let $\mathrm{hol}(v)$ be the line sum of $A^0$ in
that direction. Translating this line by one transverse step
changes its sum by at most $\ell p_A$, by summing the intervening
$\ell$ plaquettes. Its mean over the transverse two-torus is zero,
since $A^0$ has mean zero. The mean distance in each coordinate
circle is at most $\ell/4$, hence
$|\mathrm{hol}(v)|\le\ell^2p_A/2$. The line sum of $d\lambda$
around a closed line vanishes. Removing the inner bonds using
\eqref{eq:photon-inner-basepoint-bound} therefore gives
\begin{equation}\label{eq:photon-seam-basepoint-bound}
 |A''(b')|\le|\mathrm{hol}(v)|+(\ell-1)\ell p_A
                                  \le\tfrac32\ell^2p_A
\end{equation}
for the seam value. In particular this construction also gives a
global real gauge with
$\|A^0-d\lambda\|_\infty\le2\ell^2\|dA\|_\infty$.

Since $H_k^{\mathrm{ph}}d\lambda$ is closed and has mean zero,
it is exact, say $d\nu'$. On the nonwrapping cube
$\square^\dagger$, let $x$ denote lifted coordinates and set
$\nu=a\cdot x+\nu'$. Then
$H_k^{\mathrm{ph}}A-d\nu=H_k^{\mathrm{ph}}A''$ there.
If $b\subset\square^\dagger$ and $b'$ is inner, its initial
point is at distance at most $d(b,b')+2$ from an end of $b$.
Thus \eqref{eq:photon-inner-basepoint-bound} gives
\[
 |A''(b')|\le(2d(b,b')+4+s_d)p_A.
\]
For a seam bond,
$d(b,b')\ge(\ell-1-s_d)/2$.

Put $g=\gamma_P$ and
$c_{\mathrm{bd}}(g)=3\exp(3g)(1+6/g)^3$, which bounds the sum
of $\exp(-g d(b,b'))$ over unit bonds. Summing any of the three
kernel bounds against the inner bonds gives
\[
 C_P^{\mathrm{ph}}p_A\big[(s_d+4)c_{\mathrm{bd}}(g)
                         +(2/g)c_{\mathrm{bd}}(g/2)\big],
\]
using $2t\exp(-gt)\le(2/g)\exp(-gt/2)$.
For the seam contribution split off half the decay. Since
$s_d\le3\ell/13$,
\[
 \exp(-g d(b,b'))
 \le\exp(g/4)\exp(-5g\ell/26)\exp(-g d(b,b')/2),
 \qquad
 \ell^2\exp(-5g\ell/26)\le15/g^2.
\]
Together with \eqref{eq:photon-seam-basepoint-bound}, this bounds
the seam contribution by
$62C_P^{\mathrm{ph}}p_Ag^{-2}c_{\mathrm{bd}}(g/2)$.
The proposition follows with
\begin{equation}\label{eq:photon-local-regularity-constant}
 C_B=C_P^{\mathrm{ph}}\big[(s_d+4)c_{\mathrm{bd}}(\gamma_P)
 +(2/\gamma_P+62/\gamma_P^2)c_{\mathrm{bd}}(\gamma_P/2)\big].
\end{equation}
\end{proof}

\subsection{Weakened photon kernels}
\label{subsec:photon-weakened-kernels}

For $a=M$ or $LM$, let $\pi_a$ consist of the closed cubes of
half-side $a$ centered at $a\Z^3/\ell\Z^3$. These are overlapping
parameter cubes, distinct from the cubes that partition polymers.
Two parameter cubes are linked when they contain a common lattice
site; components below refer to chains of such links. For a set
$S$ of unit sites, let $\mathrm{Cub}_a(S)$ be all parameter cubes
meeting $S$.

Choose the shortest rectilinear path $\mathrm{Pth}(p,p')$ in
coordinate order $1,2,3$. Its signed coordinate offsets are
uniquely determined because $\ell$ is odd. A site belongs to at
most $27$ parameter cubes. If $S$ is connected by nearest-neighbor
steps, then $\mathrm{Cub}_a(S)$ is linked: consecutive sites lie
together in a cube, and all cubes containing either site are
linked to this cube. Along a coordinate segment, at most nine
transverse choices of centers occur. Counting the longitudinal
centers for the three segments of $\mathrm{Pth}(p,p')$ gives
\begin{equation}\label{eq:photon-parameter-path-count}
 \#\mathrm{Cub}_a(\mathrm{Pth}(p,p'))
                            \le9(3d(p,p')/a+9).
\end{equation}
An $L$-block, and the ends of a unit bond, meet at most $27$
cubes of either kind.

We smooth only in the row variable of the minimizer. Let
\[
 \vartheta(r)=
 \begin{cases}0&r\le0,\\3r^2-2r^3&0\le r\le1,\\1&r\ge1,
 \end{cases}
 \qquad
 \varphi(t)=\vartheta(2(3/4-|t|)),
 \qquad
 \chi_c(x)=\prod_{\nu=1}^3\varphi((x_\nu-c_\nu)/M).
\]
The coordinates in the last expression are represented with
modulus at most $\ell/2$. Put $\chi_c(b)=\chi_c(x_b)$, where
$x_b$ is the midpoint of $b$. The identity
$\vartheta(r)+\vartheta(1-r)=1$ implies
$\sum_{j\in\Z}\varphi(t-j)=1$, and therefore
$\chi_c\ge0$ and $\sum_c\chi_c=1$ on the torus.
Moreover $\varphi'$ is bounded by three and is Lipschitz with
constant $24$. Tensor products and the averaged derivative
formula for the discrete difference give, for $\tilde b=b+r$
with $|r|\le1$,
\begin{equation}\label{eq:photon-cutoff-differences}
 |\partial_\nu\chi_c(b)|\le3/M,\qquad
 |\chi_c(\tilde b)-\chi_c(b)|\le9|r|/M,\qquad
 |\partial_\nu\chi_c(\tilde b)-\partial_\nu\chi_c(b)|
                                                     \le42|r|/M^2.
\end{equation}
For the last bound the derivative factor contributes $24/M^2$,
and the two remaining factors contribute $2\cdot9/M^2$.
If $\chi_c(b)\ne0$, both ends of $b$ belong to the cube centered
at $c$.

For a unit bond $b'$ with ends $p'_1,p'_2$, define
\[
 P'(c,b')=\mathrm{Cub}_M
                 (\mathrm{Pth}(c,p'_1)\cup\mathrm{Pth}(c,p'_2)),
 \qquad
 w(b,b';s)=\sum_c\chi_c(b)\prod_{a\in P'(c,b')}s_a.
\]
For an interior basis vector set
$\mathrm{St}(e_i^y)=B(y)$ and $\mathrm{Pt}(e_i^y)=\{y\}$;
for $\delta_b$, both sets are the two ends of $b$. Put
\[
 P''(\Upsilon,\Upsilon')=\mathrm{Cub}_{LM}
 \left(\mathrm{St}(\Upsilon)\cup\mathrm{St}(\Upsilon')
 \cup\!\!\bigcup_{\substack{p\in\mathrm{Pt}(\Upsilon)\\
                           p'\in\mathrm{Pt}(\Upsilon')}}
                                  \mathrm{Pth}(p,p')\right).
\]
The weakened kernels are
\begin{equation}\label{eq:photon-weakened-definitions}
 \begin{split}
 H_k^{\mathrm{ph}}(s)(b,b')&=H_k^{\mathrm{ph}}(b,b')w(b,b';s),
                                                  \quad(k\ge1),\\
 C_k^{1/2}(s)(\Upsilon,\Upsilon')
    &=C_k^{1/2}(\Upsilon,\Upsilon')
                         \prod_{a\in P''(\Upsilon,\Upsilon')}s_a,
                                                  \quad(k\ge0).
 \end{split}
\end{equation}
The two families have parameter sets $\pi_M$ and $\pi_{LM}$,
respectively. When both occur in one expression, $s$ denotes
their pair. Each entry is a polynomial, independent of the
connection and real for real parameters. At $s=1$ the exact
operators are recovered. For general complex parameters, the
second family is defined by the displayed kernel; positivity
or the square-root identity is used only at $s=1$.

Fix $0<\alpha_0\le1/4$ and take
\begin{equation}\label{eq:photon-weakening-scale}
 M=L^m\ge L^2,\qquad
 2\alpha_0n_P\log M\le\gamma_PM,\qquad n_P=402.
\end{equation}
Parameters range over the closed polydisc $|s_a|\le M^{\alpha_0}$.
Such an $M$ exists because $(\log M)/M\to0$.
If $\chi_c(\hat b)\ne0$ for a translate of $b$ by distance at
most two, then $d(c,x_b)\le M+1$; applying
\eqref{eq:photon-parameter-path-count} to both paths gives
\[
 \#P'(c,b')\le234+54d(b,b')/M
                              \le n_P(1+d(b,b')/M).
\]
At most $27$ such centers occur: in each coordinate their
representatives lie in an interval of length $3M/2+5<3M$.
The set $P''$ has two endpoint sets and at most four paths.
Since their endpoints are within $L$ of the respective supports,
the same count gives
\[
 \#P''(\Upsilon,\Upsilon')
 \le378+216/M+108d(\Upsilon,\Upsilon')/(LM)
 \le n_P(1+d(\Upsilon,\Upsilon')/M).
\]
For either product, with corresponding distance $D$, we have
\begin{equation}\label{eq:photon-weakening-product-bound}
 \left|\prod_{a\in P}s_a\right|
 \le M^{\alpha_0n_P}\exp(\alpha_0n_P D\log M/M)
 \le M^{\alpha_0n_P}\exp(\gamma_PD/2).
\end{equation}
This is the loss that the exact exponential decay absorbs.

\begin{proposition}\label{prop:photon-weakened-bounds}
Under \eqref{eq:photon-weakening-scale}, put
$\gamma^P=\gamma_P/2$ and $C^P=116C_P^{\mathrm{ph}}M^{\alpha_0n_P}$.
The three kernels of $H_k^{\mathrm{ph}}(s)$ satisfy
\begin{equation}\label{eq:photon-weakened-regularity}
 |H_k^{\mathrm{ph}}(s)(b,b')|,
 |\partial H_k^{\mathrm{ph}}(s)(b,b')|,
 |\delta_\alpha\partial H_k^{\mathrm{ph}}(s)((b,\tilde b);b')|
                                  \le C^P\exp(-\gamma^Pd(b,b')),
\end{equation}
and
\begin{equation}\label{eq:photon-weakened-root-decay}
 |C_k^{1/2}(s)(\Upsilon,\Upsilon')|
                  \le C^P\exp(-\gamma^Pd(\Upsilon,\Upsilon')).
\end{equation}
These bounds hold on the closed polydiscs.
\end{proposition}

\begin{proof}
The root estimate follows at once from
\eqref{eq:photon-square-root-decay} and
\eqref{eq:photon-weakening-product-bound}.
For the minimizer, fix $b'$ and abbreviate
$H(b)=H_k^{\mathrm{ph}}(b,b')$, $w(b)=w(b,b';s)$ and
$b^+=b+\eta e_\nu$. Put
\[
 \mathrm{Pm}=M^{\alpha_0n_P}\exp(\gamma_Pd(b,b')/2).
\]
For $b,b^+,\tilde b,\tilde b^+$, and the intervening translates,
\eqref{eq:photon-cutoff-differences} and the $27$-center count
give
\[
 |w|\le\mathrm{Pm},\quad |\partial_\nu w|\le9\mathrm{Pm},\quad
 |w(\tilde b^+)-w(b^+)|\le27|r|\mathrm{Pm},\quad
 |\partial_\nu w(\tilde b)-\partial_\nu w(b)|
                                             \le14|r|\mathrm{Pm}.
\]
Here $\eta\le1/3$ and $M\ge9$; all translates are within two
of the initial point. A rectilinear path between $b$ and
$\tilde b=b+r$ has at most $3|r|/\eta$ steps, at distance at
most one from $b$. Summing the first-derivative bound for $H$
along this path gives
\[
 |H(\tilde b)-H(b)|
       \le3|r|C_P^{\mathrm{ph}}\exp(\gamma_P)\exp(-\gamma_Pd(b,b')).
\]
Use the discrete product rule
\[
 \partial_\nu(Hw)(b)
             =(\partial_\nu H)(b)w(b^+)+H(b)\partial_\nu w(b).
\]
It bounds the zeroth and first kernels with factors one and ten.
Subtracting this formula at $b$ and $\tilde b$ gives four terms:
the difference of $\partial_\nu H$, the difference of $w(b^+)$,
the difference of $H$, and the difference of $\partial_\nu w$.
Their bounds have factors $1,27,27\exp(\gamma_P),14$, respectively,
times
$|r|^\alpha C_P^{\mathrm{ph}}\mathrm{Pm}\exp(-\gamma_Pd(b,b'))$.
We used $|r|\le|r|^\alpha$ and $\gamma_P\le1$.
Their sum is at most $116$ times this expression, proving
\eqref{eq:photon-weakened-regularity}.
\end{proof}

At level zero the row smoothing is replaced by the diagonal
family
\begin{equation}\label{eq:photon-weakened-level-zero}
 H_0^{\mathrm{ph}}(s)(b,b')
               =\delta_{bb'}\prod_{a\in\pi_M:\,b\subset a}s_a.
\end{equation}
Every product has at most $27$ factors and at least one. The
first difference has two diagonal terms; the H\"older difference
has four, supported within distance two. Thus their entries are
bounded by $4\exp(2\gamma^P)M^{27\alpha_0}
\exp(-\gamma^Pd(b,b'))$, which is at most the right side of
\eqref{eq:photon-weakened-regularity} by $C_P^{\mathrm{ph}}\ge4\exp(2)$.

The weakening also separates components exactly. Suppose that
$s_a=0$ outside a set $\mathscr Y$ of parameter cubes, and let
$Y_1$ be one linked component of their union. A linked path set
$P$ meeting this component either lies entirely in it, or
contains a cube outside $\mathscr Y$ and its parameter product
vanishes. Both $P'$ and $P''$ are linked and contain every
parameter cube meeting the respective column support; $P''$
also contains every cube meeting the row support. In the first
family, every contributing row cube $c$ contains the fine bond
$b$ and belongs to $P'(c,b')$. It follows that a row or column
supported in $Y_1$ depends only on parameters in that component
and has no nonzero entry at an index outside it. A nonzero row
is supported in some component. The diagonal family has the
same properties because all cubes containing $b$ are linked
through an end of $b$.

Finally, coordinate sign changes preserve the ordered shortest
paths, and translations by the relevant center lattice preserve
the parameter cubes. The exact kernels commute with these
transformations: the averaging, axial and Landau constraints
are preserved, and uniqueness identifies the transformed
minimizers and covariance. The basis transforms by a signed
permutation. The parameter products and $\chi_c$ transform
accordingly. Hence $H_k^{\mathrm{ph}}(s)$ is covariant under
coordinate sign changes and $M$ translations, and
$C_k^{1/2}(s)$ under coordinate sign changes and $LM$
translations, with the parameters transformed as well.

\subsection{The domain bounds for block integration}
\label{subsec:photon-domain-bounds}

We make explicit the domain conditions used with these operators.
Take $0<\varepsilon\le1/200$, $e_k\in(0,1)$, and put
$t_{i,k}=e_k^{-3/4+i\varepsilon}$ for $i=1,2,3$.
The set $R_k$ consists of complex fine bond fields with
$\|B\|_i<t_{i,k}$ in the three norms. A field belongs to
$c\widetilde R_k$, $0<c\le1$, if it has a decomposition
$B=B_0+B_1$, with $B_1\in cR_k$ and $B_0$ real, such that on
each $\square^\dagger$ the field $B_0$ is gauge equivalent by a real scalar
to a real field with norms strictly below $ct_{i,k}$.
The decomposition and gauges need only exist. For
$\widetilde R_k^-$ the local real representative has the bounds
$t_{i,k}$, while the global term $B_1$ has the bounds
$(1-e_k^\varepsilon)t_{i,k}$. Finally,
\[
 \mathrm{Dom}_k=
 \{A\text{ real unit bond field}:\ H_k^{\mathrm{ph}}A
                                      \in\tfrac14\widetilde R_k\}.
\]
These definitions permit every real harmonic part.

Let $C_w\ge C^P$ and $0<\gamma_w\le\gamma^P$, and choose
$c_\gamma$ to bound both sums
\[
 \sup_\Upsilon\sum_{\Upsilon'}\exp(-\gamma_wd(\Upsilon,\Upsilon')),
 \qquad
 \sup_b\sum_{b'}\exp(-\gamma_wd(b,b'))
\]
over basis elements and unit bonds, respectively. The counts
above make $c_\gamma$ independent of $k,\ell$.
Let $p_{0,k}=(-\log e_k)^{p_0}$ and $p_k=(-\log e_k)^p$.
Assume
\begin{equation}\label{eq:photon-domain-small-charge}
 \begin{split}
 p_{0,k}&\le e_k^{-\varepsilon/4},\qquad
 2C_w^2c_\gamma^2c_Ce_k^{\varepsilon/4}\le1,\\
 8C_Bp_k&\le t_{3,k},\qquad
 8e_k^{3/4-7\varepsilon/2}\le1,\qquad
 L^{3/4}\ge8/3.
 \end{split}
\end{equation}
Since the constants are fixed before the charge, these are
conditions on an upper bound $e_0$, valid for every
$0<e_k\le e_0$.

\begin{proposition}\label{prop:photon-domain-preservation}
Under \eqref{eq:photon-domain-small-charge}, every real unit
bond field with $\|dA\|_\infty\le p_k$ belongs to
$\mathrm{Dom}_k$. For every $\zeta\in\mathrm{Box}_k$ and every
pair of parameters in the closed polydiscs, the field
\[
 \mathscr Z_k(s,\zeta)
       =H_k^{\mathrm{ph}}(s)C C_k^{1/2}(s)\zeta
\]
has all three norms bounded by $\frac12e_k^{-\varepsilon/2}$.
If $\ell/L$ is also an admissible side and
$e_{k+1}=L^{1/2}e_k$, then
\begin{equation}\label{eq:photon-domain-preserved-one-step}
 A_k^{\min}(A')+\mathrm{Sq}_k\zeta\in\mathrm{Dom}_k
 \quad(A'\in\mathrm{Dom}_{k+1},\ \zeta\in\mathrm{Box}_k).
\end{equation}
\end{proposition}

\begin{proof}
The first assertion follows from
\eqref{eq:photon-curvature-local-regularity}: its local real
representative has each norm at most
$C_Bp_k\le t_{3,k}/8<t_{i,k}/4$. Use $B_1=0$ in the domain
decomposition.

For the fluctuation bound, the root estimate and the first
summation bound give
\[
 \sup_\Upsilon|(C_k^{1/2}(s)\zeta)(\Upsilon)|
                         \le C_wc_\gamma e_k^{-\varepsilon/4}.
\]
The coordinate estimate for $C$ holds also for complex
coordinates: evaluating its real bound on the signs of the real
row coefficients gives $\sum_\Upsilon|C(b,\Upsilon)|\le c_C$.
Each norm of $H_k^{\mathrm{ph}}(s)v$ is bounded by
$C_wc_\gamma\|v\|_\infty$, by the respective kernel sum.
Consequently all three norms of $\mathscr Z_k(s,\zeta)$ are at most
\[
 C_w^2c_\gamma^2c_Ce_k^{-\varepsilon/4}
                                     \le\tfrac12e_k^{-\varepsilon/2}.
\]
This proves the complex weakened assertion at every level,
including the diagonal family at level zero.

For \eqref{eq:photon-domain-preserved-one-step}, write
$\mathscr A'=H_{k+1}^{\mathrm{ph}}A'$. It is real and belongs to
$\frac14\widetilde R_{k+1}$. On any coarse enlarged cube there
is a real gauge in which $\mathscr A'$ is represented by a real
field $A^c$ with $i$-th norm below $t_{i,k+1}/2$: add the real
global summand to the local representative. For any fine
$M$-cube $\square$, choose a coarse cube $\square'$ with
$\square\subset L\square'$. Then
$\square^\dagger\subset L(\square'^\dagger)$.
The scaling of the three norms is
$L^{-1/2},L^{-3/2},L^{-3/2-\alpha}$, respectively, and
\[
 t_{i,k+1}=L^{-3/8+i\varepsilon/2}t_{i,k}.
\]
Thus $(A^c)_L$ has $i$-th norm strictly below
$\frac12L^{-7/8+3\varepsilon/2}t_{i,k}
\le\frac12L^{-3/4}t_{i,k}\le3t_{i,k}/16$.
By \eqref{eq:photon-background-one-step}, the new background on
$\square^\dagger$ is gauge equivalent to
$(A^c)_L+\mathscr Z_k(1,\zeta)$. The latter fluctuation is real,
and its norms satisfy
\[
 \tfrac12e_k^{-\varepsilon/2}\le t_{3,k}/16\le t_{i,k}/16
\]
by \eqref{eq:photon-domain-small-charge}. The sum has norms
strictly below $t_{i,k}/4$. This is the required local
representative, again with global summand zero.
\end{proof}

The weakened fluctuation also gives the complex radius used in
the later Taylor expansions. If $A\in\widetilde R_k^-$ and
\begin{equation}\label{eq:photon-fluctuation-analytic-radius}
 |t|\le\mathrm{rad}_1:=\tfrac12e_k^{-3/4+9\varepsilon/2},
\end{equation}
then $A+t\mathscr Z_k(s,\zeta)\in\widetilde R_k$. Indeed the
three added norms are at most
$\frac14e_k^{-3/4+4\varepsilon}
<e_k^\varepsilon t_{3,k}$; this fits the margin in the global
summand of $\widetilde R_k^-$. The fluctuation is a polynomial
in the two parameter families and linear in $\zeta$.
This is the photon input to analyticity of the step integrand.

\subsection{Consistency of the exact operators on covers}
\label{subsec:photon-cover-consistency}

Let $\ell_2=J\ell_1$ with $J$ odd and with both sides admissible.
Write $\mathrm{pj}^*$ for periodic pullback, also on the coarse
unit lattices and on fluctuation coordinates. The common
reflection eigenbasis in each block makes the last pullback
well defined. Its image is the subspace of $\ell_1$-periodic
fields, and
\[
 \langle\mathrm{pj}^*f,\mathrm{pj}^*h\rangle_{\ell_2}
                              =J^3\langle f,h\rangle_{\ell_1}.
\]
All exact photon operators commute with this pullback:
\begin{equation}\label{eq:photon-exact-cover-consistency}
 O^{(2)}\mathrm{pj}^*=\mathrm{pj}^*O^{(1)},
 \qquad
 O=H_k^{\mathrm{ph}},\Delta_k,C,\widetilde\Delta_k,
                (\widetilde\Delta_k+x)^{-1},C_k^{1/2},\mathrm{Sq}_k.
\end{equation}
The superscript specifies the torus; $x\ge0$. Also
\begin{equation}\label{eq:photon-minimizer-cover-consistency}
 (A_k^{\min})^{(2)}(\mathrm{pj}^*A')
                         =\mathrm{pj}^*(A_k^{\min})^{(1)}(A').
\end{equation}

To prove these identities, first note that $d,\delta$,
$\Delta_{\mathrm{sc}},\mathbf Q_j,\mathbf Q_j^T,\tau,C$
have the same finite row formulas on the periodic lifts. They
commute with pullback and with translations by $\ell_1\Z^3$.
Pulling back the three Euler equations
\eqref{eq:photon-minimizer-Euler} therefore characterizes the
larger-torus minimizer with the pulled-back datum. Uniqueness
gives the identity for $H_k^{\mathrm{ph}}$.

The exact operators intertwine the orthogonal translations
by $\ell_1\Z^3$. Transposing these relations shows that their
adjoints also intertwine the translations and preserve the
periodic subspaces. If
$O^{(2)}\mathrm{pj}^*=\mathrm{pj}^*O^{(1)}$, pairing with two
pullbacks and canceling $J^3$ gives the same identity for $O^T$.
Products give the identities for $\Delta_k$ and
$\widetilde\Delta_k$. Invertibility gives the resolvent identity;
integrating it by \eqref{eq:photon-root-integral} gives the root
identity, and hence the identity for $\mathrm{Sq}_k$.

Finally let $A=(A_k^{\min})^{(1)}(A')$ and test its pullback
against a larger-torus axial fluctuation $\widetilde Z$. Sum the
translates of $\widetilde Z$ over
$\ell_1\Z^3/\ell_2\Z^3$ to obtain $\mathrm{pj}^*Z$. Its average
and axial sums vanish, so $Z\in C(V_k)$ on the smaller torus.
The field $\Delta_k^{(2)}\mathrm{pj}^*A$ is periodic; consequently
\[
 \langle\mathrm{pj}^*A,\Delta_k^{(2)}\widetilde Z\rangle_{\ell_2}
  =J^{-3}\langle\mathrm{pj}^*\Delta_k^{(1)}A,
                                         \mathrm{pj}^*Z\rangle_{\ell_2}
  =\langle A,\Delta_k^{(1)}Z\rangle_{\ell_1}=0.
\]
Its average and axial constraints also pull back. The
characterization \eqref{eq:photon-block-minimizer-orthogonality}
proves \eqref{eq:photon-minimizer-cover-consistency}.


\section{Local fermion integration}\label{sec:fermion-primitives}

The walk expansion also provides the localization required by the
renormalization step. A parameter is attached to each $M$-cube traversed
by a walk, together with its one-layer neighbourhood. Setting parameters
to zero then separates the operator into components. The coarse
covariance is obtained from this expansion by an exact resolvent
identity. At parameter value one it agrees with the exact covariance;
its local terms expand the determinant and the interaction.
We work at mesh $\eta_k=L^{-k}$ on the torus of side $\ell$; this is
the side denoted by $R$ in the local inverse construction.
Throughout this section the mass satisfies $0<\mu\le1$,
the range of the renormalization data.
Transposition is with respect to the bilinear lattice pairings:
for $O:E\to E'$, $\langle v,Ou\rangle_{E'}=
\langle O^Tv,u\rangle_E$. These pairings use the counting,
chamber, or coarse lattice measure of the respective space.
No complex conjugation enters this operation.

\subsection{Weakened walks and polymer propagators}\label{subsec:weakened-walks}

The resolvent parameter ranges over $Y_0=\R$ at level zero and
$Y_k=\R\cup\{\infty\}$ at $k\ge1$, with $-\infty=\infty$
as in the preceding resolvent construction.

Let $\pi_k$ be the partition into $M$-cubes and let $\mathcal P_k$ be the set of
nonempty unions of these cubes connected through faces. For $X\in\mathcal P_k$
write $|X|_M$ for its number of cubes. With the net and length
conventions of Appendix~\ref{app:polymer-geometry}, set
\[
 M d_M(X)=\inf\{\operatorname{length}(\mathcal N):
       \mathcal N\text{ is a net meeting every closed cube of }X\}.
\]
For a face-connected union of cubes, joining successive cube centres gives
$d_M(X)\le |X|_M-1$.

For a walk $\omega=(\square_0,\ldots,\square_n)$ from
\S\ref{subsec:propagator-walks}, put
\[
 X_\omega=\bigcup_{i=0}^n\square_i^{\sim1},\qquad
 s_\omega=\prod_{\square\subset X_\omega}s_\square,
 \qquad |s_\square|\le M^{\alpha_0},\quad \alpha_0=1/54.
\]
The union $X_\omega$ is a polymer and
\begin{equation}\label{eq:walk-polymer-size}
 |X_\omega|_M\le27(n+1),\qquad d_M(X_\omega)\le27(n+1).
\end{equation}
Indeed, each one-layer neighbourhood contains $27$ cubes. Successive
centres differ by at most $M$ in each coordinate, so their neighbourhoods
are joined through faces. The preceding bound on $d_M$ applies.
Every walk, including a walk of length zero, carries the parameters of
all cubes in $X_\omega$.

The left parametrix gives terms
\[
 S'_{y,\omega}=(h_nS_n)(-C_nh_{n-1}S_{n-1})\cdots
                 (-C_1h_0S_0)h_0.
\]
Here, as for $S_{y,\omega}$, $S_i$ is the local inverse on
$\square_i$ and $C_i=-[F_y,h_i]$. Define
\begin{align}
 S_y(s,A)&=\sum_\omega s_\omega S_{y,\omega}(A),&
 S_y(X,A)&=\sum_{X_\omega=X}S_{y,\omega}(A),\label{eq:weakened-walk-definitions}\\
 S'_y(s,A)&=\sum_\omega s_\omega S'_{y,\omega}(A),&
 S'_y(X,A)&=\sum_{X_\omega=X}S'_{y,\omega}(A).\notag
\end{align}
Their averages are denoted by $S_y^{\mathrm{sy}}$. The superscript
$\#$ below means any of these three choices. Put
$\kappa'=\kappa_{\mathrm w}+\kappa_0$, where
$\kappa_{\mathrm w},\kappa_0>0$ are fixed polymer decay rates.
Choose $M$ large enough that
\begin{equation}\label{eq:weakened-walk-contraction}
 q(M^{27\alpha_0})\le\tfrac12,\qquad
 q(e^{27\kappa'})\le\tfrac12,
 \qquad q(p)=54\pi C_Va'p/M.
\end{equation}

\begin{proposition}\label{prop:weakened-polymer-propagators}
For $A\in\widetilde{\mathcal G}_k$, $y\in Y_k$, and
$s\in\prod_{\pi_k}\{|s_\square|\le M^{\alpha_0}\}$, the series
\eqref{eq:weakened-walk-definitions} converge absolutely in $N_\gamma$,
uniformly in $(s,A)$, and
\begin{align}
 N_\gamma(S_y^\#(s,A))&\le32a'M^{27\alpha_0}w_k(y),
                                       \label{eq:weakened-propagator-bound}\\
 N_\gamma(\delta_{\alpha,A}S_y^\#(s,A))&\le304a'M^{27\alpha_0},\notag\\
 \|\mathbf1_{\Delta_z}S_y^\#(X,A)\mathbf1_{\Delta_{z'}}\|
   &\le32a'e^{27\kappa'}w_k(y)
       e^{-\gamma d(z,z')-\kappa'd_M(X)}.
                                       \label{eq:polymer-propagator-bound}
\end{align}
The norm in the last line is the sup operator norm. At $k=0$ it can
also be replaced by the supremum of the normalized kernel over the two
clusters, with $w_0(y)=(1+|y|)^{-1}$.
At parameter value one,
\begin{equation}\label{eq:polymer-propagator-sum}
 S_y^\#(1,A)=S_y(A)=\sum_{X\in\mathcal P_k}S_y^\#(X,A).
\end{equation}
The weakened operators are continuous on the closed parameter
polydisk, jointly analytic in its interior and in $A$, and analytic in
$A$ for each fixed parameter value, including boundary values.
For finite real $y$ the families are continuous in $y$.

The operator $S_y^\#(X,A)$ is supported on $X$ in both arguments and
depends on $A$ only on bonds with endpoints in $X$. If parameters vanish
off a union $Y$ of $M$-cubes, then each component of $Y$ is preserved by
$S_y^\#(s,A)$. On such a component the operator depends only on its
parameters and its connection. Its row at $x=(c,g)$ vanishes unless all
$M$-cubes meeting $K_c=\{c':|c'-c|_\infty\le2L\}$ have active
parameters; the corresponding assertion holds for columns.
\end{proposition}

\begin{proof}
By \eqref{eq:walk-polymer-size}, $|s_\omega|\le p_s^{n+1}$ with
$p_s=M^{27\alpha_0}$. Apply the walk row estimate with this weight.
The length-$n$ contribution is bounded by
$16a'p_sw_k(y)2^{-n}$; its H\"older row estimate is bounded by
$152a'p_s2^{-n}$. Summing these bounds proves the first two estimates.
For a fixed polymer use
\[
 \mathbf1[X_\omega=X]\le
       e^{-\kappa'd_M(X)}p_X^{n+1},\qquad p_X=e^{27\kappa'}.
\]
The same row estimate with $p_X$ proves
\eqref{eq:polymer-propagator-bound}. The left walk estimates follow
from the left parametrix by the same weighted product estimate; taking
the average preserves all bounds.

On a fixed torus, summing the row bounds over its finitely many unit
blocks also bounds the sum of the $N_\gamma$ norms of the individual
terms. The tails decay geometrically, uniformly in $(s,A)$. The series
may therefore be regrouped by $X_\omega$. At $s=1$ they are the right
and left inverse expansions, so both equal $S_y(A)$. This proves
\eqref{eq:polymer-propagator-sum}. Each finite sum is analytic in its
connection variables and polynomial in $s$. The uniform tail estimates
give the stated analyticity and continuity.
The local inverse is continuous at every real resolvent parameter.
Its walk terms inherit this property, and their tails are uniform in
$y$ since $w_k(y)\le1$. The sums are therefore continuous in real $y$.

A walk term is a product of local factors on the neighbourhoods
$\square_i^{\sim1}$. Its row and column cutoffs, its commutators, and
all connection variables consequently lie in $X_\omega$. If the first
row is nonzero, its photon site satisfies
$|c-z_{\square_0}|_\infty<3M/4$. Since $3M/4+2L<3M/2$, every
$M$-cube meeting $K_c$ is in $\square_0^{\sim1}$, and its parameter
occurs in $s_\omega$. The last cutoff gives the column assertion;
the order is reversed for a left walk. Finally, a surviving walk has
every cube of $X_\omega$ in $Y$. Since $X_\omega$ is face-connected,
it lies in a single component of $Y$. Its local factors use only the
connection in that component. This proves all the locality statements.
\end{proof}
The lattice symmetries used in this section have coordinate sign
changes and translations in $M\Z^3$. Their spin lifts lie in the
group generated by $\gamma_1,\gamma_2,\gamma_3,\gamma_5$; they act
on $(c,g)$ by $(rc+a,rg)$, where
$S\gamma_\mu S^{-1}=r_{\mu\mu}\gamma_\mu$ determines $r$.
Write $\widetilde g=(S,a)$ for such a symmetry.
The step group restricts $a$ to
$LM\Z^3$. A lift commuting with $\gamma_5$ preserves $y$; a lift
anticommuting with it sends $y$ to $-y$ at nonneutral values.
Gauge covariance follows term by term: the conjugating phases of
adjacent factors cancel. A lattice symmetry permutes cubes and walks,
and hence acts on the parameters by
$(\widetilde g s)_\square=s_{\widetilde g^{-1}\square}$.
Transposition reverses the order of
the local factors and interchanges the right and left expansions.
More precisely, with $\mathsf C=i\gamma_2$ one has
$F_y(A)^T=\mathsf C^{-1}F_{-y}(-A)\mathsf C$. Therefore
\begin{equation}\label{eq:weakened-transpose}
 S_y(s,A)^T=\mathsf C^{-1}S'_{-y}(s,-A)\mathsf C.
\end{equation}
Thus at the neutral values $y=0,\infty$ for $k\ge1$, and $y=0$
for $k=0$, the symmetrized family has the same transpose covariance as
the exact inverse. This is the reason for using it in the coarse
covariance and determinant.

\subsection{The fine and coarse resolvents}\label{subsec:fine-coarse-resolvent}

For $k\ge1$ let
\[
 \mathcal Q_k=Q_k(A)J^*,\qquad
 \mathcal Q_k^\sharp=JQ_k^T(-A),\qquad
 P_L=Q^T(-A)Q(A).
\]
Here $Q_k$ averages the fine lattice onto the unit lattice, and $Q$
averages the unit lattice onto the lattice of spacing $L$.
The notation $\sharp$ denotes the displayed algebraic transpose with
reversed connection; for complex $A$ it is not a Hermitian adjoint.
Here $B(z')=E_{z'}$ is the $L$-block from the preceding section.
To specify the phases, if $c\in\Delta_z$ and $z\in B(z')$, set
\[
 \varphi_A(z,c)=e^{ie_k\eta_k(\tau_kA)(z,c)},\qquad
 U_A(z',z)=e^{ie_k\eta_k(\tau A)(z',z)}.
\]
The path average $\tau A$ is the arithmetic mean of the oriented bond
sums in the six coordinate orders, and $\tau_kA$ is their sum along
the nested block centres. Then
\begin{align*}
 (\mathcal Q_kf)(z)&=\eta_k^3\sum_{c\in\Delta_z}
                \varphi_A(z,c)\frac1{48}\sum_gf(c,g),\\
 (\mathcal Q_k^\sharp v)(c,g)&=\varphi_A(z,c)^{-1}v(z),\\
 (P_Lv)(z)&=U_A(z',z)^{-1}L^{-3}
                    \sum_{z''\in B(z')}U_A(z',z'')v(z'').
\end{align*}
The nesting of the odd blocks and cancellation of the phases give
\begin{equation}\label{eq:block-algebra}
 \mathcal Q_k\mathcal Q_k^\sharp=1,\quad P_L^2=P_L,\quad
 \mathcal Q_k^\sharp\mathcal Q_k=\Pi_1(A),\quad
 \mathcal Q_k^\sharp P_L\mathcal Q_k=\Pi_L(A).
\end{equation}
All these operators commute with constant spin matrices.

Write $S_k=S_\infty(A)$ and define maps on the unit spinors by
\begin{align}
 H_k&=b_kS_k\mathcal Q_k^\sharp,&
 \mathfrak H_k&=b_k\mathcal Q_kS_k,\notag\\
 D_k&=b_k-b_k^2\mathcal Q_kS_k\mathcal Q_k^\sharp,&
 M_k&=b_k-D_k.\label{eq:coarse-dirac-definition}
\end{align}
For real $y$ put
\[
 \mathcal B_{k,y}=(b_k+i\gamma_5y)^{-1}(1-P_L)
                   +(b_k+b_L+i\gamma_5y)^{-1}P_L.
\]
These $4\times4$ inverses exist since
$(t+i\gamma_5y)^{-1}=(t-i\gamma_5y)/(t^2+y^2)$ for $t>0$.
The idempotent decomposition and \eqref{eq:block-algebra} give
\begin{equation}\label{eq:fine-resolvent-perturbation}
 \mathcal B_{k,y}^{-1}=b_k+b_LP_L+i\gamma_5y,\qquad
 F_y=F_\infty-b_k^2\mathcal Q_k^\sharp
                       \mathcal B_{k,y}\mathcal Q_k.
\end{equation}

\begin{proposition}\label{prop:fine-coarse-resolvents}
For $k\ge1$, $A\in\widetilde{\mathcal G}_k$, and real $y$, the
operator $D_k+b_LP_L+i\gamma_5y$ is invertible and its inverse is
\begin{equation}\label{eq:coarse-resolvent-formula}
 \Gamma_{k,y}=\mathcal B_{k,y}
       +b_k^2\mathcal B_{k,y}\mathcal Q_kS_y
                         \mathcal Q_k^\sharp\mathcal B_{k,y}.
\end{equation}
Moreover,
\begin{equation}\label{eq:resolved-fine-covariance}
 S_y=S_k+H_k\Gamma_{k,y}\mathfrak H_k.
\end{equation}
In particular, $\Gamma_k=\Gamma_{k,0}$ exists and
$\delta Z_k(A)=\det(D_k+b_LP_L)\ne0$.
\end{proposition}

\begin{proof}
Set $Z=b_k^2\mathcal Q_k^\sharp\mathcal B_{k,y}\mathcal Q_k$.
The two fine inverses exist by the preceding section, and
\eqref{eq:fine-resolvent-perturbation} gives both identities
\[
 S_y-S_k=S_kZS_y=S_yZS_k.
\]
Since $D_k+b_LP_L+i\gamma_5y
=\mathcal B_{k,y}^{-1}-b_k^2\mathcal Q_kS_k\mathcal Q_k^\sharp$,
multiplication by the right side of
\eqref{eq:coarse-resolvent-formula} gives the identity: the additional
term is
$b_k^2\mathcal Q_k(S_y-S_k-S_kZS_y)
\mathcal Q_k^\sharp\mathcal B_{k,y}=0$.
The product in the other order uses $S_y-S_k-S_yZS_k=0$.
This proves the inverse formula without assuming that the coarse
operator is invertible. Substituting the formula in
$H_k\Gamma_{k,y}\mathfrak H_k$ gives
\[
 S_kZS_k+S_kZS_yZS_k
     =S_kZ(S_k+S_yZS_k)=S_kZS_y=S_y-S_k.
\]
At $y=0$ the coarse operator is an invertible map on the finite
dimensional unit spinor space, so its determinant is nonzero.
\end{proof}

At level zero this algebra is replaced by a single chamber
integration. The operators $b_0$, $S_\infty$, $\mathcal B_{0,y}$, and
$M_0$ are not defined. Instead,
\begin{equation}\label{eq:level-zero-covariance}
 F_0=D_A+\mu+b_LJP_LJ^*,\qquad
 \Gamma_0=F_0^{-1},\qquad \delta Z_0(A)=\det F_0\ne0.
\end{equation}
Here $\Gamma_0$ acts on $192$ components per photon site, whereas
$\Gamma_k$ for $k\ge1$ acts on four components per unit site.

\subsection{Localized coarse covariances}\label{subsec:localized-coarse-covariance}

The weakened families for $k\ge1$ are
\begin{align}
 H_k(s,A)&=b_kS_\infty(s,A)\mathcal Q_k^\sharp,
                                          \label{eq:local-coarse-families}\\
 \Gamma_k(s,A)&=\mathcal B_{k,0}
    +b_k^2\mathcal B_{k,0}\mathcal Q_kS_0^{\mathrm{sy}}(s,A)
                 \mathcal Q_k^\sharp\mathcal B_{k,0},\notag\\
 M_k(X,A)&=b_k^2\mathcal Q_kS_\infty^{\mathrm{sy}}(X,A)
                  \mathcal Q_k^\sharp.\notag
\end{align}
The unprimed polymer propagators $S_{k,y}(X,A)=S_y(X,A)$ are used
in the determinant integral. At level zero let
$\operatorname{Cub}(c)$ be the $M$-cubes having a lattice site in
common with $K_c$, and set
\begin{equation}\label{eq:local-level-zero-families}
 \begin{aligned}
 (H_0(s,A)u)(c,g)&=
       \left(\prod_{\square\in\operatorname{Cub}(c)}s_\square\right)u(c,g),\\
 \Gamma_0(s,A)&=S_0^{\mathrm{sy}}(s,A),\qquad
 S_{0,y}(X,A)=S_y(X,A).
 \end{aligned}
\end{equation}
At $s=1$ these are the exact operators; also
$\sum_XM_k(X,A)=M_k(A)$. The exponential formulas for the averaging
phases and Proposition~\ref{prop:weakened-polymer-propagators} give
analyticity on the same domains.

For use in the estimates, let
$\mathfrak i(A)=e_k\|\operatorname{Im}A\|_\infty$.
The small-connection domain and the charge bound give
$\mathfrak i(A)\le1/(6L)$; reduce $\gamma$, if necessary, so that
$\gamma L\le1/6$. The path lengths imply
\begin{equation}\label{eq:coarse-phase-bounds}
 |\varphi_A(z,c)^{\pm1}|\le e^{3\mathfrak i(A)/2},\qquad
 \|\mathcal Q_k\|_{\infty\to\infty},
 \|\mathcal Q_k^\sharp\|_{\infty\to\infty}
       \le e^{3\mathfrak i(A)/2}.
\end{equation}
The sup norms here are formed with the Euclidean spinor norm.
Each unit block has chamber measure one. In particular, for a fine
operator $O$,
\begin{equation}\label{eq:block-compression-bound}
 |(\mathcal Q_kO\mathcal Q_k^\sharp)(z,z')|_2
       \le \exp(1)\,\|\mathbf1_{\Delta_z}O\mathbf1_{\Delta_{z'}}\|.
\end{equation}
Indeed, apply the two averaging bounds to a unit spinor supported at
$z'$; their product is $e^{3\mathfrak i(A)}\le \exp(1)$.

\begin{proposition}\label{prop:coarse-local-family-bounds}
The families \eqref{eq:local-coarse-families} satisfy
\begin{align}
 \|H_k(s,A)\|_{\infty\to\infty}
       &\le32\exp(1)b_ka'M^{27\alpha_0},\label{eq:weakened-H-bound}\\
 \|\delta_{\alpha,A}H_k(s,A)\|_{\infty\to\infty}
       &\le304\exp(1)b_ka'M^{27\alpha_0},\notag\\
 |\Gamma_k(s,A)(z,z')|_2
       &\le C_\Gamma M^{27\alpha_0}e^{-\gamma d(z,z')},
                                             \label{eq:weakened-Gamma-bound}\\
 |M_k(X,A)(z,z')|_2
       &\le32\exp(1)b_k^2a'e^{27\kappa'}
                     e^{-\gamma d(z,z')-\kappa'd_M(X)},\notag
\end{align}
where
\[
 C_\Gamma=\frac{2(1+2\exp(1))}{b_-}+128\exp(1)(1+2\exp(1))^2a'.
\]
The H\"older estimate is on short pairs, including the distinct
sites of one cluster. The families $H_k(s,A)$ and $\Gamma_k(s,A)$
preserve every component of the active $M$-cubes and depend there only
on its parameters and connection. The row of $H_k(s,A)$ at $(c,g)$
vanishes unless all cubes meeting $K_c$ are active. The polymer
$M_k(X,A)$ is supported on the unit sites of $X$ and depends only on
the connection in $X$.

At level zero the normalized kernel of $\Gamma_0(s,A)$ is bounded by
$32a'M^{27\alpha_0}e^{-\gamma d(c,c')}$, and
\eqref{eq:polymer-propagator-bound} bounds $S_{0,y}(X,A)$ with its
factor $(1+|y|)^{-1}$. The multiplier $H_0(s,A)$ has sup norm
at most $M^{8\alpha_0}$ and equals the identity at $s=1$.
Its row at $(c,g)$ vanishes unless every cube meeting $K_c$
is active. There are no short pairs of distance less than one
at this level.
\end{proposition}

\begin{proof}
The first two estimates follow by composing the weakened fine bounds
with \eqref{eq:coarse-phase-bounds}. The third and fourth use
\eqref{eq:block-compression-bound}. For the covariance it remains to
control the two factors $\mathcal B_{k,0}$. Its kernel is
\[
 \mathcal B_{k,0}(z,z')=b_k^{-1}\mathbf1[z=z']1_4
      -\frac{b_L}{b_k(b_k+b_L)}P_L(z,z').
\]
It vanishes unless $z,z'$ belong to one $L$-block. Since
$|P_L(z,z')|_2\le\exp(1)L^{-3}$ on that block, its row and column sums
are at most $(1+2\exp(1))/b_k$. In a nonzero term of
$\mathcal B_{k,0}G\mathcal B_{k,0}$,
where $G=\mathcal Q_kS_0^{\mathrm{sy}}(s,A)\mathcal Q_k^\sharp$,
the two intermediate sites are within $L$ of the exterior sites.
Therefore
\[
 e^{-\gamma d(z_1,z_2)}\le
      e^{2\gamma L}e^{-\gamma d(z,z')}
      \le4e^{-\gamma d(z,z')}.
\]
The fine bound and \eqref{eq:block-compression-bound} give
$|G(z_1,z_2)|_2\le32\exp(1)a'M^{27\alpha_0}e^{-\gamma d(z_1,z_2)}$.
Multiplying the row and column sums cancels the factor $b_k^2$ and
gives $128\exp(1)(1+2\exp(1))^2a'M^{27\alpha_0}$.
The local term $\mathcal B_{k,0}$ is bounded by
$2(1+2\exp(1))b_-^{-1}e^{-\gamma d(z,z')}$, proving
\eqref{eq:weakened-Gamma-bound}.

The odd block partitions are nested, and every unit block and every
$L$-block belongs to a single $M$-cube. Moreover, the paths entering
$\varphi_A(z,c)$ lie in $\Delta_z$, and the paths entering
$U_A(z',z)$ lie in the fine block of side $L$ centred at $z'$:
each coordinate of a rectilinear path remains between those of its
endpoints. Consequently the averaging maps and $\mathcal B_{k,0}$
preserve every union of $M$-cubes and depend only on its connection.
Composing them with the component-local fine families proves the
locality assertions. This applies also to the parameter-independent
term $\mathcal B_{k,0}$ in $\Gamma_k(s,A)$.
For the row of $H_k$, the terminal average does not change the first
row index of $S_\infty(s,A)$, so the stronger vanishing statement is
preserved. At level zero the assertions for $\Gamma_0$ follow
from the fine kernel bound and locality with $\eta_0=1$.

For $H_0$, lift $c$ and $K_c$ to the lattice. In each coordinate
$K_c$ is an interval of $4L+1$ sites. Since $M\ge L^2$ and $L\ge5$,
this interval meets at most two consecutive $M$-blocks. The cubes
meeting $K_c$ are the product of these three sets of blocks.
They number at most eight, form a face-connected set, and include
the cube of $c$. Thus the defining product is a monomial of degree
at most eight, bounded by $M^{8\alpha_0}$. If one of its cubes is
inactive, its factor is zero. If all are active, their connectedness
puts them in the active component containing $c$; hence changing
parameters outside that component has no effect on the row or
column at $(c,g)$. The multiplier is scalar in spin and chamber
labels and independent of the connection. These observations give
its analyticity, gauge covariance, and component locality.
\end{proof}

These families have the gauge and lattice covariance of their exact
counterparts. For $\Gamma_k$ and $M_k$, transposition follows from
\eqref{eq:weakened-transpose} at a neutral value and the scalar
averaging phases. The barred field map is
$\bar H_k(s,A)=\mathsf C^{-1}H_k(s,-A)\mathsf C$.
Its bounds follow from those for $H_k$ because $\mathsf C$ is unitary
and $\mathsf C\gamma_\mu\mathsf C^{-1}=-\gamma_\mu^T$.
When component sup norms are used in the Grassmann substitution,
the Euclidean operator bounds above acquire at most the factor two
from $|v|_c\le|v|_2\le2|v|_c$.

\subsection{Weighted Grassmann integration}\label{subsec:weighted-grassmann}

Let $\phi_i(t)$, $1\le i\le p$, be independent Grassmann families
on finite ordered index sets. If $\phi_i(I_i)$ is the ordered product
over $I_i$, write
\[
 F=\sum_{I_1,\ldots,I_p}c_F(I_1,\ldots,I_p)
                 \phi_1(I_1)\cdots\phi_p(I_p),\qquad
 \|F\|_{\mathsf w}=\sum_I|c_F(I)|\prod_i\mathsf w_i^{|I_i|}.
\]
The positive measures $\nu_i(t)$ on the index sets enter the kernels,
not this coefficient norm. Thus an antisymmetric kernel expansion
\[
 F=\sum_n\prod_i\frac1{n_i!}\int\dd t\,K_n(t)
                       \prod_i\prod_{j=1}^{n_i}\phi_i(t_j^i)
\]
has norm
\begin{equation}\label{eq:grassmann-kernel-norm}
 \|F\|_{\mathsf w}
     =\sum_n\prod_i\frac{\mathsf w_i^{n_i}}{n_i!}
                      \int\dd t\,|K_n(t)|.
\end{equation}
For arbitrary kernels the right side is an upper bound. Products of
ordered monomials vanish on repeated generators and otherwise multiply
their weights, so
\begin{equation}\label{eq:grassmann-product-bound}
 \|FG\|_{\mathsf w}\le\|F\|_{\mathsf w}\|G\|_{\mathsf w},
 \qquad \|e^F-1\|_{\mathsf w}\le e^{\|F\|_{\mathsf w}}-1
       \quad(F\text{ even}).
\end{equation}

For an additional finite nonempty mode set $\mathcal I$ with positive weights
$\nu(t)$, introduce $W(t),\bar W(t)$ and set
\[
 \langle\bar W,OW\rangle_\nu
    =\sum_{t,t'}\nu(t)\nu(t')\bar W(t)O(t,t')W(t'),
 \qquad \mathrm{Top}=\prod_{t\in\mathcal I}W(t)\bar W(t).
\]
The kernel represents $O$ with measure $\nu$; its matrix entry is
$O(t,t')\nu(t')$. Define $\int\mathrm DW$ as extraction of the
coefficient of $\mathrm{Top}$, multiplied by
$\prod_t\nu(t)^{-1}$, and define
\begin{align}
 \int\dd\mu_I(W)F
     &=\int\mathrm DW\,e^{-\langle\bar W,W\rangle_\nu}F,
                                            \label{eq:weighted-Berezin}\\
 \int\dd\mu_C(W)F
     &=(\det C)\int\mathrm DW\,
               e^{-\langle\bar W,C^{-1}W\rangle_\nu}F
                    \quad(C\text{ invertible}).\notag
\end{align}
These are algebra-valued functionals when other Grassmann generators
are present. The order of modes in $\mathrm{Top}$ is immaterial because
its factors are even.

\begin{proposition}\label{prop:weighted-Berezin-bounds}
Put $w_W=\max_t\nu(t)^{-1/2}$. Then
\begin{equation}\label{eq:weighted-integration-bound}
 \left\|\int\dd\mu_I(W)F\right\|_{\mathsf w}
                  \le\|F\|_{(\mathsf w,w_W)}.
\end{equation}
The integral factors for polynomials involving disjoint sets of modes.
For an invertible covariance,
\begin{align}
 \int\dd\mu_C(W)F(\bar W,W)
     &=\int\dd\mu_I(W)F(\bar W,CW),\label{eq:covariance-substitution}\\
 \int\dd\mu_C(W)W(t)\bar W(t')&=C(t,t'),\notag\\
 \int\mathrm DW\,e^{-\langle\bar W,OW\rangle_\nu}
     &=\det O\quad\text{for every }O.\label{eq:weighted-determinant}
\end{align}
Finally, consider a linear substitution
$\phi_i\mapsto\sum_jO_{ij}\theta_j$, and put
\[
 \|O_{ij}\|_{1,\infty}
     =\sup_t\sum_s|O_{ij}(t,s)|\nu'_j(s).
\]
If $\sum_j\|O_{ij}\|_{1,\infty}\mathsf w'_j\le\mathsf w_i$ for
each $i$, then its image $F'$ obeys
$\|F'\|_{\mathsf w'}\le\|F\|_{\mathsf w}$.
Only rows corresponding to generators occurring in $F$ need satisfy
this bound.
\end{proposition}

\begin{proof}
The identity density factors as
\[
 e^{-\langle\bar W,W\rangle_\nu}
            =\prod_t(1+\nu(t)W(t)\bar W(t)).
\]
An ordered monomial $W(J)$ contributes to its integral only if $J$
contains both generators of every mode it meets. In that case its
integral is its reordering sign times
$\prod_{(t,0)\in J}\nu(t)^{-1}$; otherwise it is zero. Its modulus
is therefore at most $\prod_{u\in J}\nu(u)^{-1/2}\le w_W^{|J|}$.
Apply this coefficientwise to obtain
\eqref{eq:weighted-integration-bound}. For disjoint mode sets a
nonzero integral requires each monomial separately to be paired.
Their degrees in $W$ are then even, so the signs from exchanging
external generators cancel. The paired-mode formula gives
factorization. It also shows that modes absent from the integrand may
be removed.

A linear change $W\mapsto OW$ multiplies the coefficient of the top
monomial by $\det O$, since the matrix of the change is
$O(t,t')\nu(t')$. Applied to the identity density this proves
\eqref{eq:weighted-determinant}, without an invertibility assumption.
For invertible $C$, use $W\mapsto CW$, leaving $\bar W$ fixed, in
the density with covariance $C$. This proves
\eqref{eq:covariance-substitution} and normalization. Its two-point
function is
\[
 \sum_s C(t,s)\nu(s)\int\dd\mu_I(W)W(s)\bar W(t')
       =\sum_s C(t,s)\nu(s)\frac{\mathbf1[s=t']}{\nu(t')}
       =C(t,t').
\]
Thus it is the kernel, rather than the matrix entry.

Under the stated substitution, the norm of the image of
$\phi_i(t)$ is at most
\[
 \sum_j\mathsf w'_j\sum_s|O_{ij}(t,s)|\nu'_j(s)\le\mathsf w_i
\]
by hypothesis. Submultiplicativity bounds
the image of each monomial by its original weight; summing its
coefficients proves the last assertion, with the same proof on the
subalgebra generated by the rows that occur.
\end{proof}

For linear combinations $v_1,\ldots,v_{2r}$ of the $W,\bar W$,
Wick's formula gives
\begin{equation}\label{eq:weighted-wick}
 \int\dd\mu_I(W)v_1\cdots v_{2r}
      =\operatorname{Pf}\left(\int\dd\mu_I(W)v_iv_j\right)_{i,j}.
\end{equation}
Odd products integrate to zero. To fix the signs, let $\partial_u$
be the left Grassmann derivative. It satisfies
$\partial_u(vG)=(\partial_uv)G-v\partial_uG$, and its Berezin
integral is zero. Differentiating the factored Gaussian density gives
\begin{align*}
 \int\dd\mu_I(W)W(t)G
      &=\nu(t)^{-1}\int\dd\mu_I(W)\partial_{\bar W(t)}G,\\
 \int\dd\mu_I(W)\bar W(t)G
      &=-\nu(t)^{-1}\int\dd\mu_I(W)\partial_{W(t)}G.
\end{align*}
Consequently, if $I(v)=\int\dd\mu_I(W)v$,
\[
 I(v_1\cdots v_q)=\sum_{j=2}^q(-1)^jI(v_1v_j)
                         I(v_2\cdots\widehat v_j\cdots v_q).
\]
This is the expansion of the Pfaffian along its first row and proves
\eqref{eq:weighted-wick} by induction. Formula
\eqref{eq:covariance-substitution} gives the corresponding Wick rule
for $\mu_C$.

In the renormalization step the modes have weight one for $k\ge1$.
At $k=0$ they are $(c,g,\beta)$ with weight $1/48$. Thus the
integration weight is
\begin{equation}\label{eq:first-step-Grassmann-weight}
 w_W=1\quad(k\ge1),\qquad w_W=\sqrt{48}\quad(k=0).
\end{equation}
In particular the first-step identity covariance has two-point
function $48\mathbf1[(c,g,\beta)=(c',g',\beta')]$.
The pair index used for a H\"older field has measure
$(\eta_k^3/48)^2$. In a substitution involving
$\delta_{\alpha,A}H_k$, its scalar component row sum is computed
with the input measure, as in
Proposition~\ref{prop:weighted-Berezin-bounds}. Internal-cluster
differences are therefore controlled by the H\"older bound with the
same field weight as other pairs.

\subsection{The Gaussian block identity}\label{subsec:Gaussian-block-identity}

For this identity use a common notation for the integration space:
at $k\ge1$ let $\mathcal D_k=D_k$ and
$\mathcal A_k(A)=Q(A)$; at $k=0$ let
$\mathcal D_0=D_A+\mu$ and $\mathcal A_0(A)=Q(A)J^*$.
The integration measure is respectively counting measure or the
chamber measure $1/48$. The target is the four-component spinor
space on the lattice of spacing $L$, with measure $L^3$.
Set
\[
 \mathcal A_k^\sharp(A)=\mathcal A_k(-A)^T,\qquad
 K_k=\mathcal D_k+b_L\mathcal A_k^\sharp\mathcal A_k,
 \qquad\Gamma_k=K_k^{-1}.
\]
Thus $K_k=D_k+b_LP_L$ for $k\ge1$ and $K_0=F_0$.
Define the quadratic operator on the coarse lattice by
\[
 D_{k+1}^0=b_L-b_L^2\mathcal A_k\Gamma_k
                                       \mathcal A_k^\sharp.
\]
Let $\Phi,\bar\Phi$ be external coarse Grassmann fields and put
\[
 \Psi^c=b_L\Gamma_k\mathcal A_k^\sharp\Phi,\qquad
 \bar\Psi^c=b_L\Gamma_k^T\mathcal A_k(A)^T\bar\Phi.
\]
For $\Psi=\Psi^c+W'$ and $\bar\Psi=\bar\Psi^c+\bar W'$,
completion of the square is the exact identity
\begin{multline}\label{eq:block-square-completion}
 \langle\bar\Psi,\mathcal D_k\Psi\rangle
   +b_L\langle\bar\Phi-\mathcal A_k(-A)\bar\Psi,
                    \Phi-\mathcal A_k(A)\Psi\rangle_L\\
   =\langle\bar\Phi,D_{k+1}^0\Phi\rangle_L
                    +\langle\bar W',K_kW'\rangle.
\end{multline}
Indeed, write $s=b_L\mathcal A_k^\sharp\Phi$ and
$\bar s=b_L\mathcal A_k(A)^T\bar\Phi$. The left side is
\[
 \langle\bar\Psi,K_k\Psi\rangle
      -\langle\bar\Psi,s\rangle-\langle\bar s,\Psi\rangle
      +b_L\langle\bar\Phi,\Phi\rangle_L.
\]
Substitute the means $\Gamma_ks,\Gamma_k^T\bar s$. Both cross
terms cancel, and the remaining constant is
$b_L\langle\bar\Phi,\Phi\rangle_L
-\langle\bar s,\Gamma_ks\rangle
=\langle\bar\Phi,D_{k+1}^0\Phi\rangle_L$.
All factors remain in their displayed order; no conjugation of
Grassmann fields is involved.

Translation by external odd fields preserves the top coefficient
in the integration variables. Applying this to
\eqref{eq:block-square-completion} and then using
\eqref{eq:covariance-substitution} gives, for any Grassmann
polynomial $F$, the block integration formula
\begin{multline}\label{eq:Gaussian-block-integration}
 \int\mathrm D\Psi\,
    e^{-\langle\bar\Psi,\mathcal D_k\Psi\rangle
      -b_L\langle\bar\Phi-\mathcal A_k(-A)\bar\Psi,
                         \Phi-\mathcal A_k(A)\Psi\rangle_L}
                     F(\bar\Psi,\Psi)\\
 =\delta Z_k(A)e^{-\langle\bar\Phi,D_{k+1}^0\Phi\rangle_L}
      \int\dd\mu_I(W)
          F(\bar\Psi^c+\bar W,\Psi^c+\Gamma_kW).
\end{multline}
At the first step this is a single integration over all chamber
components. The determinant and the integration weight are those of
\eqref{eq:level-zero-covariance} and
\eqref{eq:first-step-Grassmann-weight}; the heavy components are
included in both.

\subsection{An integral for the determinant}\label{subsec:determinant-integral}

The determinant in \eqref{eq:Gaussian-block-integration} has a local
logarithm obtained by inserting the polymer propagators into a
resolvent integral. We first derive the integral at real connections.
Put $n_k=\ell^3$ and
\[
 c_k^d=4\bigl((1-L^{-3})\log b_k+L^{-3}\log(b_k+b_L)\bigr)
       \quad(k\ge1),\qquad c_0^d=0.
\]
For real $A\in\widetilde{\mathcal G}_k$ and $0<\mu\le1$,
\begin{align}
 \delta Z_k(A)
  &=\exp\left\{c_k^dn_k-ib_k^2\int_0^\infty\dd y\,
       \operatorname{Tr}\bigl[\gamma_5\mathcal B_{k,y}
          \mathcal Q_kS_y\mathcal Q_k^\sharp\mathcal B_{k,y}\bigr]
                        \right\}\quad(k\ge1),
                                    \label{eq:determinant-integral-positive-level}\\
 \delta Z_0(A)
  &=\exp\left\{\int_1^\infty\frac{\dd y}{y}
                       \operatorname{Tr}[F_0S_y]
             -i\int_0^1\dd y\,\operatorname{Tr}[\gamma_5S_y]
                        \right\}.\label{eq:determinant-integral-level-zero}
\end{align}
The trace in the first line is on the unit spinor space. At level
zero it is the normalized chamber trace
$\operatorname{Tr}O=48^{-1}\sum_{c,g}\tr O((c,g),(c,g))$.

To prove these identities, write $K=K_k$ as in the preceding
subsection. At real $A$, the fine adjoint identity
$D_A^*=\gamma_5D_A\gamma_5$, the reality of $\mu,b_k,b_L$, and
the orthogonal averaging projections give
$K^*=\gamma_5K\gamma_5$. Hence $K\gamma_5$ is Hermitian and
invertible. Since $\det\gamma_5=1$ on each four-component spinor,
its determinant is $\det K$. Let $\lambda_j$ be its real nonzero
eigenvalues, counted with multiplicity. There are $4n_k$ at positive
levels and $192n_0$ at level zero. Then
\[
 \operatorname{Tr}\bigl[\gamma_5(K+i\gamma_5y)^{-1}\bigr]
        =\sum_j(\lambda_j+iy)^{-1},\qquad
 \operatorname{Tr}\bigl[K(K+i\gamma_5y)^{-1}\bigr]
        =\sum_j\lambda_j(\lambda_j+iy)^{-1}.
\]
For a real $\lambda\ne0$, choose along $y\ge0$ the logarithm
with value $\log|\lambda|+i\pi\mathbf1[\lambda<0]$ at zero.
Its derivative is $i/(\lambda+iy)$ and
$\operatorname{Log}(\lambda+iy)-\log y\to i\pi/2$.
Integration gives
\begin{align*}
 -i\int_0^Y\frac{\dd y}{\lambda+iy}
        &=\operatorname{Log}\lambda-
                          \operatorname{Log}(\lambda+iY),\\
 \int_1^\infty\frac{\dd y}{y}\frac{\lambda}{\lambda+iy}
       -i\int_0^1\frac{\dd y}{\lambda+iy}
        &=\operatorname{Log}\lambda-i\pi/2.
\end{align*}
Summing the second identity proves
\eqref{eq:determinant-integral-level-zero}: its constant phase is
$-96\pi i n_0\in2\pi i\Z$.

For $k\ge1$, the ranks of $P_L$ and $1-P_L$ are respectively
$4L^{-3}n_k$ and $4(1-L^{-3})n_k$. The explicit formula for
$\mathcal B_{k,y}$ therefore gives
\[
 -i\int_0^Y\dd y\,\operatorname{Tr}[\gamma_5\mathcal B_{k,y}]
          =-4n_k\log Y+c_k^dn_k+o(1)
                  \quad(Y\to\infty).
\]
Insert \eqref{eq:coarse-resolvent-formula} in the first scalar
integral. The terms $4n_k\log Y$ cancel. The remaining phase is
$2\pi i n_k$, which disappears on exponentiation, and this proves
\eqref{eq:determinant-integral-positive-level}.
The integral converges absolutely: block compression bounds
$\mathcal Q_kS_y\mathcal Q_k^\sharp$ uniformly, while each row
and column sum of $\mathcal B_{k,y}$ is at most
$(1+2\exp(1))(b_k^2+y^2)^{-1/2}$. Its trace is bounded by
\[
 128\exp(1)(1+2\exp(1))^2a'n_k(b_k^2+y^2)^{-1}.
\]
At level zero the spectral identities already give absolute
convergence. The uniform kernel bounds below will give a bound
independent of the eigenvalues.

\subsection{Local determinant terms}\label{subsec:local-determinant-terms}

Choose $\gamma_{\mathrm w}\le\gamma$ and
$\kappa_{\mathrm w}=\kappa+\kappa_0+3$ with $\kappa,\kappa_0>0$.
The propagator bounds above give a fixed constant $C_w$ such that,
in component sup norms, for $y\ge0$,
\begin{align*}
 \|\mathbf1_{\Delta_z}S_{k,y}(X,A)\mathbf1_{\Delta_{z'}}\|
       &\le C_we^{-\gamma_{\mathrm w}d(z,z')
                                  -\kappa_{\mathrm w}d_M(X)}
                                                    \quad(k\ge1),\\
 |S_{0,y}(X,A)(x,x')|_c
       &\le C_w(1+y)^{-1}e^{-\gamma_{\mathrm w}d(c,c')
                                  -\kappa_{\mathrm w}d_M(X)}.
\end{align*}
Here $|\cdot|_c$ is the matrix norm induced by the component sup
norm. We may use the same $C_w$ after reducing the rates from
$\gamma,\kappa'$.
For $X\in\mathcal P_k$ define
\begin{align}
 \widetilde E_k(X,A)
    &=-ib_k^2\int_0^\infty\dd y\,
       \operatorname{Tr}\bigl[\gamma_5\mathcal B_{k,y}\mathcal Q_k
           S_{k,y}(X,A)\mathcal Q_k^\sharp\mathcal B_{k,y}\bigr]
                   \quad(k\ge1),\label{eq:local-determinant-definition}\\
 \widetilde E_0(X,A)
    &=\int_1^\infty\frac{\dd y}{y}
                     \operatorname{Tr}[F_0S_{0,y}(X,A)]
       -i\int_0^1\dd y\,\operatorname{Tr}[\gamma_5S_{0,y}(X,A)],\notag\\
 E_k^d(X,A)&=\tfrac14\bigl(\widetilde E_k(X,A)
       +\widetilde E_k(X,-A)
       +\overline{\widetilde E_k(X,\overline A)}
       +\overline{\widetilde E_k(X,-\overline A)}\bigr),\notag\\
 E_k^{\mathrm{det}}(X,A)&=E_k^d(X,A)-E_k^d(X,0).\notag
\end{align}
Set $c_X=54$ and
\[
 C_d=16c_XC_w\max\left\{\tfrac\pi2b\exp(1)(1+2\exp(1)),\,
                                  2+58\exp(1)+b\exp(1)\right\}.
\]

\begin{proposition}\label{prop:local-determinant-expansion}
The integrals \eqref{eq:local-determinant-definition} converge
absolutely for $A\in\widetilde{\mathcal G}_k$. They define analytic
functions, depending only on bonds meeting $X$, with
\begin{equation}\label{eq:local-determinant-bound}
 |\widetilde E_k(X,A)|, |E_k^d(X,A)|
               \le C_dM^3e^{-(\kappa+1)d_M(X)}.
\end{equation}
They satisfy
\begin{equation}\label{eq:local-determinant-quotient}
 \delta Z_k(A)=e^{c_k^dn_k+\sum_X\widetilde E_k(X,A)},\qquad
 \frac{\delta Z_k(A)}{\delta Z_k(0)}
                         =e^{\sum_XE_k^{\mathrm{det}}(X,A)}.
\end{equation}
The functions $E_k^d$ are gauge invariant, even in $A$, and satisfy
$\overline{E_k^d(X,\overline A)}=E_k^d(X,A)$. They are covariant
under the lattice step group.
\end{proposition}

\begin{proof}
For $k\ge1$ write
$N_y(X)=\mathcal Q_kS_{k,y}(X,A)\mathcal Q_k^\sharp$.
It is supported on the unit sites of $X$, and each entry is bounded
by $\exp(1)C_we^{-\kappa_{\mathrm w}d_M(X)}$ by normalized block
compression. To estimate the trace, cycle one factor
$\mathcal B_{k,y}$ to the left and use
\[
 G_y=\mathcal B_{k,y}\gamma_5\mathcal B_{k,y}
     =\gamma_5\bigl((b_k+i\gamma_5y)^{-2}(1-P_L)
              +(b_k+b_L+i\gamma_5y)^{-2}P_L\bigr).
\]
Its column sum in Euclidean matrix norm is at most
$(1+2\exp(1))/(b_k^2+y^2)$: the identity contributes one and the two
projection coefficients contribute at most $2\exp(1)$. Since
$|\tr(mm')|\le8|m|_2\max_{i,j}|m'_{ij}|$ for four-component
matrices, the integrand is bounded by
\begin{equation}\label{eq:polymer-determinant-trace-bound}
 8\exp(1)(1+2\exp(1))C_wM^3|X|_M
           e^{-\kappa_{\mathrm w}d_M(X)}(b_k^2+y^2)^{-1}.
\end{equation}
Integrating and multiplying by $b_k^2$ leaves the factor
$\pi b_k/2\le\pi b/2$.

At $k=0$, the formula for $D_A$ and the averaging penalty give
the scalar component row bound
\begin{equation}\label{eq:first-step-Dirac-row-bound}
 \sup_{x,\beta}\sum_{\beta'}\int\dd x'\,
                    |F_0(A)(x,x')_{\beta\beta'}|
                      \le R_F:=58\exp(1)+1+b\exp(1).
\end{equation}
The diagonal and three internal hops contribute
at most $41+\mu$, the external hop at most
$6e^{\mathfrak i(A)}$, and the block penalty at most $b\exp(1)$.
The bound follows from $\mu\le1$ and
$e^{\mathfrak i(A)}\le \exp(1)$. Point evaluation has normalized
kernel $48\mathbf1[x=x']$; its factor $48$ cancels the chamber
measure. The two traces in the level-zero integral are consequently
bounded by
\[
 4R_FC_wM^3|X|_M(1+y)^{-1}e^{-\kappa_{\mathrm w}d_M(X)},
 \qquad
 4C_wM^3|X|_M(1+y)^{-1}e^{-\kappa_{\mathrm w}d_M(X)}.
\]
The second bound uses that $\gamma_5$ has four entries of modulus
one. Their integrals contribute respectively
$\int_1^\infty(y(1+y))^{-1}\dd y=\log2$ and
$\int_0^1(1+y)^{-1}\dd y=\log2$.

The geometric bound \eqref{eq:polymer-size-tree} gives
$|X|_M\le c_X(d_M(X)+1)\le c_Xe^{d_M(X)}$.
Applying this to the trace bounds,
and using $\kappa_{\mathrm w}-1\ge\kappa+1$, proves
\eqref{eq:local-determinant-bound} for $\widetilde E_k$ and hence
for its four-term mean. The majorants are independent of $A$ and
integrable in $y$. Each integrand is analytic in $A$; Cauchy's
formula on a compact polydisk and the majorant permit its contour
integrals to be exchanged with the $y$ integral. The result is
analytic.

Every phase in the compressed trace belongs to a unit block or an
$L$-block meeting $X$, and that block is contained in $X$ by
nesting. The fine propagator is supported on $X$ and uses only its
connection. At level zero a nonzero external-hop trace uses a bond
meeting $X$; the penalty uses the $L$-block of its row. This proves
locality. Gauge transformations conjugate each product under the
trace and hence leave it unchanged. A lattice transformation whose
spin matrix commutes with $\gamma_5$ does the same. Every spatial
action of the step group has such a lift: multiplication by
$\gamma_1\gamma_2\gamma_3$ changes its commutation sign with
$\gamma_5$ and leaves its spatial action fixed. This proves step
covariance. The mean in \eqref{eq:local-determinant-definition}
then makes $E_k^d$ even and real in the stated analytic sense.

For real connections, sum the polymer propagators and use
\eqref{eq:determinant-integral-positive-level} or
\eqref{eq:determinant-integral-level-zero}. This proves the first
identity of \eqref{eq:local-determinant-quotient} there. Both sides
are analytic on $\widetilde{\mathcal G}_k$, which is open and
star-shaped. They agree on a real neighbourhood of zero, so their
Taylor series agree there and analytic continuation proves the
identity throughout the domain.

To justify the symmetrized quotient, put
$\mathcal L(A)=\sum_X\widetilde E_k(X,A)$.
The charge-conjugation transpose and the complex-conjugation
identity give
$\delta Z_k(-A)=\delta Z_k(A)
=\overline{\delta Z_k(\overline A)}$.
For the latter, the spin matrix $\varkappa=\gamma_5\mathsf C$
conjugates $\overline{K_k(\overline A)}$ to $K_k(-A)$.
Thus $\mathcal L(-A)-\mathcal L(A)$ and
$\overline{\mathcal L(\overline A)}-\mathcal L(A)$ take values
in $2\pi i\Z$. They are continuous on the connected domain,
so they are constant. Subtracting their values at zero gives
\[
 \sum_X\bigl(E_k^d(X,A)-E_k^d(X,0)\bigr)
                            =\mathcal L(A)-\mathcal L(0).
\]
This proves the second identity in
\eqref{eq:local-determinant-quotient} without a choice of logarithm
branch for the complex determinant.
\end{proof}

\subsection{The small-field gain for the determinant}\label{subsec:determinant-small-field-gain}

For a photon connection write $A_\mu(c)=A(c,c+\eta_ke_\mu)$ and
\[
 \partial_\nu A_\mu(c)
       =\eta_k^{-1}(A_\mu(c+\eta_ke_\nu)-A_\mu(c)),\qquad
 \|\delta_\alpha\partial A\|_\infty
       =\sup_{\mu,\nu}\sup_{0<d(c,c')\le1}
          \frac{|\partial_\nu A_\mu(c')-\partial_\nu A_\mu(c)|}
               {d(c,c')^\alpha}.
\]
Local versions use only differences and bonds inside the given set.
In this section the curl is mesh normalized:
$(dA)_{\mu\nu}=\partial_\mu A_\nu-\partial_\nu A_\mu$.
Let $R_k$ be the complex bond functions whose three norms
$\|A\|_\infty$, $\|\partial A\|_\infty$, and
$\|\delta_\alpha\partial A\|_\infty$ are less than
$e_k^{-3/4+\varepsilon}$, $e_k^{-3/4+2\varepsilon}$, and
$e_k^{-3/4+3\varepsilon}$, respectively. Let $\widetilde R_k$
consist of $A_0+A_1$, with $A_1\in R_k$ and $A_0$ real, such that
on every $\square^\dagger=\square^{\sim L}$ one can write
$A_0=A'+d\lambda$, where $\lambda,A'$ are real and the same
three local bounds hold for $A'$. The external connection domain is
\[
 \mathfrak X_k=\{\mathscr C:\|\mathscr C\|_\infty
                  <e_k^{-1+2\varepsilon},\quad
             \|\partial\mathscr C\|_\infty<e_k^{-1+3\varepsilon}\}.
\]
If $e_k^{1/4}+e_k^\varepsilon<1$, both $\mathfrak X_k$ and
$\widetilde R_k+\mathfrak X_k$ lie in
$\widetilde{\mathcal G}_k$: on the local cubes the real part has
a representative of size $e_k^{-3/4+\varepsilon}$, and the
remaining part has size at most
$(e_k^{1/4}+e_k^\varepsilon)e_k^{-1+\varepsilon}$.

Put $v_k=-\log e_k$. Assume
\begin{equation}\label{eq:determinant-gain-conditions}
 v_k\ge9,\qquad e_k^\varepsilon\le\tfrac14,\qquad
 8M(v_k/4+4)e_k^\varepsilon\le1,\qquad
 \ell\ge Mv_k^2.
\end{equation}
The last condition follows from the stop-scale volume condition
$\ell\ge Mr_k$ with $r_k\ge v_k^2$.

\begin{proposition}\label{prop:determinant-small-field-gain}
For $\mathscr A\in\widetilde R_k$, $\mathscr C\in\mathfrak X_k$,
and $X\in\mathcal P_k$, under \eqref{eq:determinant-gain-conditions},
\begin{equation}\label{eq:determinant-small-field-gain}
 |E_k^d(X,\mathscr A+\mathscr C)-E_k^d(X,\mathscr C)|
       \le8C_dM^3e_k^{1/4}e^{-\kappa d_M(X)}.
\end{equation}
\end{proposition}

\begin{proof}
If $d_M(X)>v_k/4$, bound the two terms separately by
\eqref{eq:local-determinant-bound} and use
$e^{-d_M(X)}\le e_k^{1/4}$. Suppose therefore that
$d_M(X)\le v_k/4$. The polymer lift and diameter estimates give
a nonwrapping lattice box $W$ containing every bond meeting $X$,
with each side of length at most
$s_W=M(d_M(X)+3)$.
Their applicability follows from
$2M(d_M(X)+2)<\ell$ and $M(d_M(X)+4)<\ell$, which hold by
\eqref{eq:determinant-gain-conditions}.

Write $\mathscr A=A_0+A_1$ as in the domain definition.
On each plaquette, the local real representative gives
$\|dA_0\|_\infty<2e_k^{-3/4+2\varepsilon}$.
Choose a fixed-order rectilinear path $\Gamma(o,c)$ from a site
$o\in W$ and set $\lambda(c)=\eta_k A_0(\Gamma(o,c))$ on $W$.
Then
\[
 \|A_0-d\lambda\|_{\infty,W}
             \le2s_W\|dA_0\|_{\infty,W}
             <4s_We_k^{-3/4+2\varepsilon}.
\]
To see the factor $2s_W$, append one bond to $\Gamma(o,c)$ and
subtract $\Gamma(o,c')$. After cancelling common bonds, the
resulting loop bounds at most two rectangles of width $\eta_k$.
Together they contain at most $2s_W/\eta_k$ plaquettes, each
contributing $\eta_k dA_0$. The paths and rectangles remain in
$W$ because it is a product of intervals. Extend the real function
$\lambda$ by zero outside the projected box. On every bond meeting
$X$, the bound is unchanged. By
\eqref{eq:determinant-gain-conditions} it is at most
$\tfrac12e_k^{-3/4+\varepsilon}$.

Let $B=A_0-d\lambda+A_1$ on bonds meeting $X$ and zero elsewhere.
Then $\|B\|_\infty<\tfrac32e_k^{-3/4+\varepsilon}$.
For $|t|\le\tfrac13e_k^{-1/4}$,
\[
 \|tB+\mathscr C\|_\infty
       <(\tfrac12+e_k^\varepsilon)e_k^{-1+\varepsilon}
       \le\tfrac34e_k^{-1+\varepsilon}.
\]
Thus this connection belongs to the complex connection ball
$\mathcal G_k=\{A:\|A\|_\infty<e_k^{-1+\varepsilon}\}$,
which is contained in $\widetilde{\mathcal G}_k$, and
$f(t)=E_k^d(X,tB+\mathscr C)$ is analytic on this disk and
bounded by the right side of \eqref{eq:local-determinant-bound}.
Gauge invariance and locality identify the difference in
\eqref{eq:determinant-small-field-gain} with $f(1)-f(0)$.
On the circle of radius $r=\tfrac14e_k^{-1/4}>2$, Cauchy's
formula yields
\[
 f(1)-f(0)=\frac1{2\pi i}\oint_{|t|=r}
                         \frac{f(t)}{t(t-1)}\,\dd t,
 \qquad |f(1)-f(0)|\le\frac{\sup|f|}{r-1}
                          \le8e_k^{1/4}\sup|f|.
\]
This proves the claimed gain, retaining the stronger decay
$e^{-(\kappa+1)d_M(X)}$ in this case.
\end{proof}


\section{Small-field renormalization}
\label{sec:small-field-rg}

The renormalization step separates the vacuum and mass terms from an
activity, transports its remaining quadratic part to local averages,
and integrates the fluctuation fields. Two estimates are needed before
this integration. Extraction must preserve the chamber symmetries and
produce a scalar mass. Reconstruction of a coarse field must preserve
the support of an activity, including the bonds that cross its boundary.
The latter estimate uses cancellation in the chamber Dirac operator
and curvature control of the averaging phases.

We use $\gamma_\mu=\sigma_3\otimes\sigma_\mu$,
$\gamma_5=\sigma_1\otimes1_2$ and $\Sigma_3=\sigma_3\otimes1_2$.
The matrices acting on the Grassmann variables below are
\[
 \mathsf C=i\gamma_2,\qquad
 \varkappa=\gamma_5\mathsf C=\sigma_2\otimes\sigma_2.
\]
Both are real, $\mathsf C^T=\mathsf C^{-1}=-\mathsf C$,
and $\varkappa^T=\varkappa=\varkappa^{-1}$.
The regulator matrix in the reflection-positive construction
is $C=\sigma_2\oplus\sigma_2$; the relation in these spin
conventions is $\mathsf C=i\Sigma_3C$.

\subsection{Activities, norms and fields}
\label{subsec:rg-activity-spaces}

Fix a level $k$, mesh $\eta_k=L^{-k}$ and an admissible side
$\ell$, with $\ell/(LM)$ odd and at least $13$.
Let $\mathcal P_k$ be the collection of face-connected polymers
of $M$-cubes on this torus. A polymer is small when
$d_M(X)\le L$. A small polymer has an injective lift, unique up
to a period translation, by the diameter and lifting estimates
of Appendix~\ref{app:polymer-geometry}.
Write $\operatorname{Vol}(X)=M^3|X|_M$.

The independent Grassmann generators are $\psi$, $\bar\psi$ on
the chamber sites $x=(c,g)$ and $\chi$, $\bar\chi$ on ordered
pairs. For distinct chamber sites on the covering lattice, put
$d_X(x,x')=\max(|c-c'|_\infty,\eta_k)$ and restrict pairs to
$d_X\le M$. Pairs on the torus are simultaneous period classes;
their shortest lifts are unique since $2M<\ell$. A pair is
admissible for $X$ when its photon sites and all six rectilinear
paths between them lie in $X$, or when it is $(x,x^{(3)})$ and
one of its photon sites belongs to $X$.
Thus a pair may leave $X$ only along a bond of type $3$.
Pairs inside a cluster have $d_X=\eta_k$ and are admissible when
their common photon site belongs to $X$.

Let $\xi=(x,\beta,\omega)$ index an ordinary leg and
$\zeta=(x,x',\beta,\omega)$ a H\"older leg, where
$\omega=0$ is unbarred and $\omega=1$ is barred.
We use the index measures
\[
 \nu(\xi)=\eta_k^3/48,\qquad
 \nu(\zeta)=(\eta_k^3/48)^2.
\]
Integrals sum the spin and bar indices as well. An activity
has a unique expansion with separately antisymmetric kernels,
\begin{equation}\label{eq:rg-activity-kernels}
 \mathscr E(X;\mathscr A,\mathscr C)
 =\sum_{n,m}\frac1{n!m!}\int\dd\xi\,\dd\zeta\,
      \mathscr E_{nm}(X;\mathscr A,\mathscr C;\xi;\zeta)
                       \psi(\xi_1)\cdots\psi(\xi_n)
                       \chi(\zeta_1)\cdots\chi(\zeta_m).
\end{equation}
Barred generators are included in this notation.
The generator weights, as distinct from the index measures, are
\[
 h_{k,1}=e_k^{-1/4},\qquad h_{k,2}=e_k^{-1/4+\varepsilon}.
\]
For an element $F$ the weighted coefficient norm is
\begin{equation}\label{eq:rg-activity-coefficient-norm}
 \|F\|_{h_k}
       =\sum_{n,m}\frac{h_{k,1}^nh_{k,2}^m}{n!m!}
                                   \int\dd\xi\,\dd\zeta\,|F_{nm}|.
\end{equation}
This is the coefficient norm already used for weighted fermion
integration. In particular a generator has norm $h_{k,i}$;
the weight $\nu$ enters only when coefficients are written as
integral kernels.

An activity family in $K_k^{x,\C}$ is balanced: every monomial
has as many barred as unbarred legs. Its kernels are holomorphic
on $\widetilde R_k\times\mathfrak X_k$, vanish unless all ordinary
legs lie over $X$ and all pairs are admissible for $X$, and read
the two connections only on bonds meeting $X$. The norm is
\begin{equation}\label{eq:rg-activity-norm}
 \|\mathscr E\|_k^x
  =\sup_{X\in\mathcal P_k}\exp(\kappa d_M(X))
    \sup_{(\mathscr A,\mathscr C)\in\widetilde R_k\times\mathfrak X_k}
                          \|\mathscr E(X;\mathscr A,\mathscr C)\|_{h_k}.
\end{equation}
We require $\kappa\ge\kappa_0$.

The required symmetries act on both the connections and the
generators. For $\tilde g=(S,a)$ with the coordinate sign action
$r(S)$ and $a\in M\Z^3$, put
\[
 \tilde g(c,g)=(r(S)c+a,r(S)g),\qquad
 a_{\tilde g}\psi(x)=S\psi(\tilde g^{-1}x),\qquad
 a_{\tilde g}\bar\psi(x)=(S^{-1})^T\bar\psi(\tilde g^{-1}x).
\]
The same spin matrices act on pair generators after transforming
both endpoints. Local gauge transformations multiply an unbarred
ordinary or pair leg by $\exp(ie_k\Lambda(c))$, using the photon
site of its first point, and a barred leg by the inverse phase.
Write this automorphism as $G_\Lambda$.
Charge conjugation and the antilinear reality map act by
\[
 \mathrm{Cg}\,\psi=\mathsf C\bar\psi,\qquad
 \mathrm{Cg}\,\bar\psi=-\mathsf C\psi,\qquad
 \mathbf K\psi=\varkappa\psi,\qquad
 \mathbf K\bar\psi=\varkappa\bar\psi,
\]
and by the same rules on the H\"older generators.

The covariance requirements on a family are
\begin{equation}\label{eq:rg-activity-symmetries}
 \begin{split}
 a_{\tilde g}\mathscr E(\tilde gX;\tilde g\mathscr A,
                                      \tilde g\mathscr C)
                                  &=\mathscr E(X;\mathscr A,\mathscr C),\\
 \mathrm{Cg}\,\mathscr E(X;-\mathscr A,-\mathscr C)
                                  &=\mathscr E(X;\mathscr A,\mathscr C).
 \end{split}
\end{equation}
For real scalar functions $\lambda,\lambda'$ on the photon torus
and two pairs in the domain
which satisfy $\mathscr A^1=\mathscr A-d\lambda$ and
$\mathscr C^1=\mathscr C-d\lambda'$ on bonds at distance at
most one from $X$, the local gauge requirement is
\[
 G_{\lambda+\lambda'}\mathscr E(X;\mathscr A^1,\mathscr C^1)
                               =\mathscr E(X;\mathscr A,\mathscr C).
\]
The real subspace $K_k^x$ has in addition
\[
 \mathbf K\mathscr E(X;-\overline{\mathscr A},-\overline{\mathscr C})
                               =\mathscr E(X;\mathscr A,\mathscr C).
\]
All requirements refer to the generators before evaluating fields.
They are essential to the form of the counterterms.

The spaces $K_k^\C,K_k$ without the second connection argument
have the same definitions on $\widetilde R_k$. Restriction
$\mathscr E|_0(X;\mathscr A)=\mathscr E(X;\mathscr A,0)$ maps
$K_k^{x,\C}$ into $K_k^\C$ and has norm at most one. The real
subspaces also restrict. Analyticity, locality and the symmetries
follow by setting the second argument and its gauge function to
zero; the norm bound follows from the smaller supremum.

For $k\ge1$ the step is defined on balls with windows
\begin{equation}\label{eq:rg-windows}
 w_k^{\mathrm m}=e_k^{3/4-8\varepsilon},\qquad
 w_k^{\mathrm E}=e_k^{1/4-7\varepsilon},\qquad
 \mathrm{Ball}_k^x(c)=
 \{(m,\mathscr E):|m|\le c w_k^{\mathrm m},
                         \|\mathscr E\|_k^x\le c w_k^{\mathrm E}\}.
\end{equation}
At level zero the ball is instead
$\{(m,0):|m|\le c w_0^{\mathrm m}\}$. This restriction is part
of the first-step construction.

When fields are evaluated, the total connection is
$\mathscr T=\mathscr A+\mathscr C$ and
\[
 \psi=H_k(\mathscr T)\Psi,\qquad
 \chi=\delta_{\alpha,\mathscr T}\psi,\qquad
 \bar\psi=\bar H_k(\mathscr T)\bar\Psi,\qquad
 \bar\chi=\delta_{\alpha,-\mathscr T}\bar\psi.
\]
For an intra-cluster pair the phase is one, so
$\psi(x')-\psi(x)=\eta_k^\alpha\chi(x,x')$. Accordingly
\[
 (1-JJ^*)\psi(x)
               =-\frac{\eta_k^\alpha}{48}\sum_{x'\ne x}\chi(x,x').
\]
The sum is within the cluster of $x$. This identity expresses
the heavy part through the same ordinary and H\"older legs.
The norm of a H\"older generator remains $h_{k,2}$, including
for an intra-cluster pair; the factor $\eta_k^\alpha$ belongs
to this evaluated identity.

\subsection{The total connection and its inverse domain}
\label{subsec:rg-total-connection-domain}

The inverse domain $\widetilde G_k$ consists of fields $B_0+B_1$
with $B_0$ real and $\|B_1\|_\infty<e_k^{-1+\varepsilon}$,
where on each $\square^{\sim5}$ the field $B_0$ has a real-gauge
representative of supremum norm below $e_k^{-1+\varepsilon}$.
The external domain is
\[
 \mathfrak X_k=
 \{\mathscr C:\|\mathscr C\|_\infty<e_k^{-1+2\varepsilon},
                  \|\partial\mathscr C\|_\infty<e_k^{-1+3\varepsilon}\}.
\]
The sets $R_k,\widetilde R_k,\widetilde R_k^-$ were defined with
the photon estimates. Their real-gauge representatives may be
chosen independently on each enlarged cube.

Assume $L\ge5$, $0<\varepsilon\le1/200$ and
\begin{equation}\label{eq:rg-domain-charge-conditions}
 e_k^\varepsilon\le\tfrac14,\qquad
 e_k^{1/4}+e_k^\varepsilon\le1,\qquad
 4e_k^{1/4+2\varepsilon}+2e_k^{3\varepsilon}\le1.
\end{equation}
Then
\begin{equation}\label{eq:rg-total-domain-inclusion}
 \widetilde R_k+\mathfrak X_k\subset\widetilde G_k.
\end{equation}
Indeed, write $\mathscr A=A_0+A_1$ with $A_1\in R_k$ and
$A_0$ locally real-gauge equivalent on $\square^\dagger$ to
$A_\square^b$ with the three threshold bounds. Restrict the
representative to $\square^{\sim5}\subset\square^\dagger$.
Its supremum norm is below
$e_k^{-3/4+\varepsilon}\le e_k^{-1+\varepsilon}$, and
\[
 \|A_1+\mathscr C\|_\infty
 <e_k^{-3/4+\varepsilon}+e_k^{-1+2\varepsilon}
 =(e_k^{1/4}+e_k^\varepsilon)e_k^{-1+\varepsilon}
 \le e_k^{-1+\varepsilon}.
\]
This is the required decomposition of the total connection.
Real gauge transformations can absorb any real constant part,
locally on the nonwrapping cubes.

Put $\mathfrak i(B)=e_k\|\operatorname{Im}B\|_\infty$ and
$q_d(B)=e_k\|dB\|_\infty$. The same decomposition gives
\begin{equation}\label{eq:rg-total-imaginary-curvature-bounds}
 \mathfrak i(\mathscr A+\mathscr C)
       \le e_k^{1/4+\varepsilon}+e_k^{2\varepsilon}\le1,
 \qquad
 q_d(\mathscr A+\mathscr C)
       \le4e_k^{1/4+2\varepsilon}+2e_k^{3\varepsilon}\le1.
\end{equation}
For the curvature, every plaquette lies in an enlarged cube,
and $dA_0=dA_\square^b$ there. This costs at most
$2\|\partial A_\square^b\|_\infty$; the global small field
costs another $2\|\partial A_1\|_\infty$, and the external
field costs $2\|\partial\mathscr C\|_\infty$.
The imaginary bound uses only that $A_0$ is real.

The domains are open: the strict inequalities define open
sets for the global summand, which can absorb a sufficiently
small complex perturbation while the local real representative
is kept fixed. The norm-defined sets are invariant under sign
and conjugation; the local real-gauge definitions are also
invariant under real gauges and real constant additions.
Coordinate sign changes and $M$ translations preserve them.
Their derivative and H\"older norms are preserved, with a
possible reversal of a bond and translation of its base point.
Thus \eqref{eq:rg-total-domain-inclusion} applies also at the
signs, conjugates and real gauges used in the step.

\subsection{Vacuum and mass extraction}
\label{subsec:rg-extraction}

For a balanced activity, its part of bidegree $(2,0)$ can be
written uniquely as
\[
 \mathscr E^{[2,0]}(X;\mathscr A,\mathscr C)
   =\int\dd x\,\dd x'\,
             \bar\psi(x)^T K_X(\mathscr A,\mathscr C;x,x')\psi(x').
\]
The matrix kernel is the ordinary antisymmetric kernel with
the first bar index one and the second bar index zero.
Antisymmetry accounts for the other ordering; in particular
\begin{equation}\label{eq:rg-bilinear-norm}
 \|\mathscr E^{[2,0]}(X;\mathscr A,\mathscr C)\|_{h_k}
   =h_{k,1}^2\sum_{\beta,\beta'}\int\dd x\,\dd x'
                     |K_X(\mathscr A,\mathscr C;x,x')_{\beta\beta'}|.
\end{equation}
The kernel is antiperiodic in each argument separately.

For a small $X$, choose one lift $\widetilde X$ and define
\begin{equation}\label{eq:rg-local-extractions}
 a_0(\mathscr E,X)
       =\operatorname{Vol}(X)^{-1}\mathscr E_{00}(X;0,0),\qquad
 a_2(\mathscr E,X)
       =\operatorname{Vol}(X)^{-1}
             \int_{\widetilde X}\dd x\int_{\widetilde X}\dd x'\,
                                                    K_X(0,0;x,x').
\end{equation}
Both cluster labels are averaged with measure $\eta_k^3/48$.
The second expression is independent of the lift: simultaneous
translation by a period multiplies its two arguments by the
same antiperiodic sign, whose product is one.

For an $M$-cube $\square$ put
\begin{equation}\label{eq:rg-counterterm-maps}
 \varepsilon_{\mathrm v}(\mathscr E,\square)
          =-\sum_{\substack{X\supset\square\\X\text{ small}}}
                                                     a_0(\mathscr E,X),
 \qquad
 M(\mathscr E,\square)
          =-\sum_{\substack{X\supset\square\\X\text{ small}}}
                                                     a_2(\mathscr E,X).
\end{equation}
Translations by $M\Z^3$ identify the small-polymer sets over
different cubes. They leave the two extracted quantities
unchanged. Thus the two sums do not depend on $\square$.
Write their values as $\varepsilon_{\mathrm v}(\mathscr E)$ and
$M(\mathscr E)$.

\begin{proposition}\label{prop:rg-scalar-counterterms}
There is a complex-linear scalar map $m(\mathscr E)$ such that
$M(\mathscr E)=m(\mathscr E)1$. Both $m$ and
$\varepsilon_{\mathrm v}$ are real on $K_k^x$, and
\begin{equation}\label{eq:rg-counterterm-bounds}
 |\varepsilon_{\mathrm v}(\mathscr E)|
                  \le K_0M^{-3}\|\mathscr E\|_k^x,
 \qquad
 |m(\mathscr E)|\le K_0M^{-3}e_k^{1/2}\|\mathscr E\|_k^x.
\end{equation}
They commute with restriction of the second connection argument
to zero.
\end{proposition}

\begin{proof}
Substitution of the generator transformations in a bilinear,
followed by uniqueness of its kernel, gives
\begin{equation}\label{eq:rg-extraction-matrix-symmetries}
 \begin{split}
 a_2(\mathscr E,\tilde gX)&=S a_2(\mathscr E,X)S^{-1},\\
 a_2(\mathscr E,X)&=\mathsf C^{-1}a_2(\mathscr E,X)^T\mathsf C,\\
 a_2(\mathscr E,X)&=\varkappa^{-1}
                         \overline{a_2(\mathscr E,X)}\varkappa
                                                   \quad(\mathscr E\in K_k^x).
 \end{split}
\end{equation}
For the first relation, the two variables and their labels
are transformed together, and the measure is unchanged.
For charge conjugation, moving the unbarred leg past the barred
leg cancels the minus sign in its transformation, and exchanges
the two kernel variables. For the last relation, the reality
map conjugates the coefficient and multiplies both legs by
$\varkappa$; its inverse transpose is $\varkappa^{-1}$.
The constant kernel has the corresponding scalar invariance
and is real for a real family.

Let $z$ be the center of $\square$ and take
$S=\gamma_1,\gamma_2,\gamma_3$ or $\gamma_5$.
The symmetry $(S,(1-r(S))z)$ fixes $\square$ and permutes the
small polymers containing it. The first relation in
\eqref{eq:rg-extraction-matrix-symmetries} therefore shows that
$M(\mathscr E)$ commutes with all four matrices. Their commutant
is scalar. In the fixed Dirac basis, $\gamma_3$ and
$\Sigma_3=-i\gamma_1\gamma_2\gamma_3$ are diagonal with the
four distinct pairs of diagonal entries
$(1,1),(-1,1),(-1,-1),(1,-1)$. A commuting matrix is diagonal;
commutation with $\gamma_1$ identifies its first two and last
two entries, and commutation with $\gamma_5$ identifies these
pairs. This proves the scalar assertion. The last relation in
\eqref{eq:rg-extraction-matrix-symmetries} makes this scalar
real for a real family.

The coefficient norm gives, for each small $X$,
\begin{equation}\label{eq:rg-single-extraction-bounds}
 \begin{split}
 \operatorname{Vol}(X)|a_0(\mathscr E,X)|
                    &\le\|\mathscr E(X;0,0)\|_{h_k},\\
 \operatorname{Vol}(X)\sum_{\beta,\beta'}
                              |a_2(\mathscr E,X)_{\beta\beta'}|
                    &\le h_{k,1}^{-2}
                               \|\mathscr E^{[2,0]}(X;0,0)\|_{h_k}.
 \end{split}
\end{equation}
Use $h_{k,1}^{-2}=e_k^{1/2}$ and
$\operatorname{Vol}(X)\ge M^3$. Each contribution is bounded
by the corresponding right side of
\eqref{eq:rg-counterterm-bounds}, times $\exp(-\kappa d_M(X))$.
The polymer sum over a fixed cube is at most $K_0$ by
$\kappa\ge\kappa_0$. For the mass use one diagonal entry of
the matrix sum, now known to be scalar. This proves the bounds.
Linearity follows from uniqueness of the antisymmetric kernels.
Restriction gives the same kernels at $(0,0)$, proving the
last assertion.
\end{proof}

Subtract these local terms only on small polymers:
\begin{equation}\label{eq:rg-normalized-remainder}
 (R\mathscr E)(X;\mathscr A,\mathscr C)
 =\mathscr E(X;\mathscr A,\mathscr C)
       -a_0(\mathscr E,X)\operatorname{Vol}(X)
       -\int_X\dd x\,\bar\psi(x)^Ta_2(\mathscr E,X)\psi(x).
\end{equation}
For large $X$ set $R\mathscr E=\mathscr E$. Then
\begin{equation}\label{eq:rg-exact-extraction-identity}
 \sum_X\mathscr E(X;\mathscr A,\mathscr C)
  =-\varepsilon_{\mathrm v}(\mathscr E)\ell^3
     -m(\mathscr E)\langle\bar\psi,\psi\rangle
                           +\sum_X(R\mathscr E)(X;\mathscr A,\mathscr C).
\end{equation}
Indeed, decompose the integral and volume of each small polymer
over its $M$-cubes, interchange the finite sums, and use
\eqref{eq:rg-counterterm-maps} on each cube. Every chamber
component receives the same scalar mass coefficient:
\[
 m(\mathscr E)\langle\bar\psi,\psi\rangle
 =m(\mathscr E)\frac{\eta_k^3}{48}
         \sum_c\sum_g\sum_{\beta=1}^4
                                      \bar\psi_\beta(c,g)\psi_\beta(c,g).
\]

The diagonal bilinear of a constant matrix $A$ on $X$ has
kernel $48\eta_k^{-3}1[x=x'\text{ over }X]A$ when both
arguments are taken over the same lift. Its integral is
$\operatorname{Vol}(X)A$, and its norm is
\[
 h_{k,1}^2\operatorname{Vol}(X)
                         \sum_{\beta,\beta'}|A_{\beta\beta'}|.
\]
These identities and \eqref{eq:rg-single-extraction-bounds}
show that
\begin{equation}\label{eq:rg-remainder-properties}
 \begin{split}
 (R\mathscr E)_{00}(X;0,0)&=0\quad(X\text{ small}),\\
 \int_{\widetilde X}\dd x\,\dd x'\,
                  (R\mathscr E)_{20}(X;0,0;x,x')&=0\quad(X\text{ small}),\\
 \|R\mathscr E\|_k^x&\le2\|\mathscr E\|_k^x.
 \end{split}
\end{equation}
For the bound, the two subtracted terms have distinct bidegrees
and together have norm at most $\|\mathscr E(X;0,0)\|_{h_k}$.
Apply the triangle inequality to
\eqref{eq:rg-normalized-remainder}, then the activity norm.

The map $R$ preserves the activity spaces and restriction to
zero. The subtracted kernels are independent of the connections,
and their two legs lie over $X$, so analyticity and locality
remain. A diagonal bilinear is invariant under every local
gauge. Its lattice, charge and reality transformations follow
from \eqref{eq:rg-extraction-matrix-symmetries}; the volume
and constant term have the same symmetries. This verifies
each requirement on the remainder, without assuming that
an individual matrix $a_2(\mathscr E,X)$ is scalar.

For later comparisons, the extraction has a precise local
meaning across different sides. Fix a lifted $M$-cube and
consider all lifted small polymers containing it. Their
projections are bijective onto the small polymers containing
its image, and their $d_M$ and volumes are unchanged.
If two families have the same constant and quadratic kernels
on every such lifted polymer at zero connections, then their
vacuum and mass matrices over this cube agree term by term in
\eqref{eq:rg-counterterm-maps}. Volume independence of a flow
will require this equality of kernels from the step comparison.

\subsection{Reconstruction by interior and boundary legs}
\label{subsec:rg-field-reconstruction}

Let $k\ge1$, $0<\mu\le1$ and $b_-\le b_k$. With
$\mathcal Q_k(A)=Q_k(A)J^*$ and
$\mathcal Q_k^\sharp(A)=JQ_k^T(-A)$, set
\begin{equation}\label{eq:rg-reconstruction-definition}
 T_k(A)=b_k^{-1}\mathcal Q_k(A)
            \big(D_A+\mu+b_k\mathcal Q_k^\sharp(A)\mathcal Q_k(A)\big)
       =\mathcal Q_k(A)\big(b_k^{-1}D_A+1+\mu/b_k\big).
\end{equation}
The second form follows from
$\mathcal Q_k\mathcal Q_k^\sharp=1$.
The fine inverse used in the definition of $H_k(A)$ exists
on $\widetilde G_k$. Hence
\begin{equation}\label{eq:rg-reconstruction-left-inverse}
 T_k(A)H_k(A)
 =b_k^{-1}\mathcal Q_k(A)F_\infty(A)
               b_kF_\infty(A)^{-1}\mathcal Q_k^\sharp(A)=1.
\end{equation}
The barred identity follows at $-A$, by conjugation with
$\mathsf C$. At level zero all four field and reconstruction
operators are the identity on the $192$ chamber components
per photon site. This first step uses the unblocked chamber spinors.

For $c\in\Delta_y$, let $c_j(c)$ be the center of its block
of side $L^j\eta_k$, $0\le j\le k$, so that $c_0(c)=c$ and
$c_k(c)=y$. The averaging phase is
\begin{equation}\label{eq:rg-reconstruction-phase}
 \varphi_A(y,c)=\exp\left(ie_k\eta_k
                 \sum_{j=0}^{k-1}(\tau A)(c_{j+1}(c),c_j(c))\right).
\end{equation}
The sum uses the average of the six rectilinear path sums.
We write $(\tau_kA)(y,c)$ for this sum.
All of these paths lie in $\Delta_y$ and their total length
in bonds is at most $\frac32(L^k-1)$. Thus
\begin{equation}\label{eq:rg-reconstruction-phase-moduli}
 \begin{aligned}
 |\varphi_A(y,c)^{\pm1}|&\le\exp(3\mathfrak i(A)/2),\\
 |u_A(c,c')^{\pm1}|&\le\exp(\eta_k\mathfrak i(A))
                                             \le\exp(\mathfrak i(A)).
 \end{aligned}
\end{equation}
for a photon bond, with $u_A(c,c')=\exp(ie_k\eta_k A(c,c'))$.

The row form of \eqref{eq:rg-reconstruction-definition} is
\begin{equation}\label{eq:rg-reconstruction-row}
 \begin{aligned}
 (T_k(A)\psi)(y)
 &=\int_{\Delta_y}\dd x\,\varphi_A(y,c)
                  \Bigl\{(1+\mu/b_k)\psi(x)\\
 &\qquad-(b_k\eta_k)^{-1}\sum_{i=0}^3M_i(x)
           [U_A(x,x^{(i)})\psi(x^{(i)})-\psi(x)]\Bigr\}.
 \end{aligned}
\end{equation}
Write $c_3=c+\eta_kge_1$ for the photon base point of $x^{(3)}$.
Here $U_A=1$ for $i\le2$ and $U_A(x,x^{(3)})=u_A(c,c_3)$.
The normalized chamber integral over $\Delta_y$ has total mass
one. This row exists for every complex connection, regardless
of invertibility; the inverse domain was needed only for
\eqref{eq:rg-reconstruction-left-inverse}.

The rows map antiperiodic fields to antiperiodic coarse fields:
translation by a period preserves blocks, neighbors and phases,
and supplies the single antiperiodic sign from the source field.

If $Y$ is a union of $M$-cubes and $y\in Y$ is a unit site,
then $\Delta_y\subset Y$: both $M$ and $L^k$ are odd, and the
unit block is contained in the $M$-cube of its center.
Define the two rows
\begin{equation}\label{eq:rg-reconstruction-two-pieces}
 \begin{split}
 (T_Y^{(1)}(A)\psi)(y)
 &=\int_{\Delta_y}\dd x\,\varphi_A(y,c)
    \Bigl\{(1+\mu/b_k)\psi(x)\\
 &\qquad-(b_k\eta_k)^{-1}\!\!\sum_{i:\,x^{(i)}\text{ over }Y}
           M_i(x)[U_A(x,x^{(i)})\psi(x^{(i)})-\psi(x)]\Bigr\},\\
 (T_Y^{(2)}(A)\chi)(y)
 &=-b_k^{-1}\eta_k^{\alpha-1}\int_{\Delta_y}\dd x\,
                1[c_3\notin Y]\varphi_A(y,c)M_3(x)\chi(x,x^{(3)}).
 \end{split}
\end{equation}
Rows outside $Y$ are zero. The first row contains ordinary
legs only over $Y$. The second row contains precisely the
type-$3$ pairs leaving $Y$, which are admissible boundary
pairs of length $\eta_k$. Their coefficients read $A$ only
on bonds with an endpoint in $\Delta_y$.
If $\chi=\delta_{\alpha,A}\psi$ on these boundary pairs, then
\begin{equation}\label{eq:rg-reconstruction-local-identity}
 T_Y^{(1)}(A)\psi+T_Y^{(2)}(A)\chi=T_k(A)\psi
                                             \quad\text{on the unit sites of }Y.
\end{equation}
This follows by replacing the missing type-$3$ difference in
\eqref{eq:rg-reconstruction-row} by
$\eta_k^\alpha\chi(x,x^{(3)})$.
The intra-cluster differences remain in $T_Y^{(1)}$.

\begin{proposition}\label{prop:rg-reconstruction-estimate}
For any complex connection, write $\mathfrak i=\mathfrak i(A)$
and $q_d=q_d(A)$. In the norm obtained by taking the supremum
of the Euclidean spinor norm,
\begin{equation}\label{eq:rg-reconstruction-sup-estimate}
 \begin{split}
 \|T_Y^{(1)}(A)\|_{\mathrm{su}},\ \|T_k(A)\|_{\mathrm{su}}
 &\le\exp(3\mathfrak i/2)
       \left[1+b_-^{-1}
                 (54Lq_d\exp(6q_d)+12+6\exp(\mathfrak i)+\mu)\right],\\
 \|T_Y^{(2)}(A)\|_{\mathrm{su}}
 &\le6b_-^{-1}\exp(3\mathfrak i/2)\eta_k^\alpha.
 \end{split}
\end{equation}
In particular, if $A\in\widetilde R_k+\mathfrak X_k$ and
\eqref{eq:rg-domain-charge-conditions} holds, the weighted
coefficient row sums satisfy
\begin{equation}\label{eq:rg-reconstruction-coefficient-bound}
 |T_Y^{(i)}(A)|_{1,\infty},\ |T_k(A)|_{1,\infty},
 |\bar T_Y^{(i)}(A)|_{1,\infty},\ |\bar T_k(A)|_{1,\infty}
                                                        \le C_T,
\end{equation}
where
$C_T=2\exp(3/2)[1+b_-^{-1}(54L\exp(6)+12\exp(1)+1)]$.
The bound is independent of $k,\ell$ and $Y$.
\end{proposition}

\begin{proof}
Put $E=\Delta_y\times O_{\mathrm{cub}}$ and divide it into
\[
 E_1=\{x\in E:c_3\in\Delta_y\},\quad
 E_2=\{x\in E:c_3\in Y\setminus\Delta_y\},\quad
 E_3=\{x\in E:c_3\notin Y\}.
\]
For each coordinate direction $a\in\{\pm e_\mu\}$ there are
eight labels with $ge_1=a$. The face of $\Delta_y$ through
which $c+\eta_ka$ leaves the unit block has $\eta_k^{-2}$
photon sites. Thus
\begin{equation}\label{eq:rg-reconstruction-face-count}
 \#E=48\eta_k^{-3},\qquad \#(E_2\cup E_3)=48\eta_k^{-2}.
\end{equation}

The chamber hopping matrices satisfy
\[
 M_i(x^{(i)})-M_i(x)=w_i\gamma\cdot\widehat d_i(x),\qquad
 \sum_{i=0}^3w_i\gamma\cdot\widehat d_i(x)=0.
\]
Also $\|M_3(x)\|_2\le6$ and
$\|w_3\gamma\cdot\widehat d_3(x)\|_2=6$.
Reindex each intra-cluster sum by $x\mapsto x^{(i)}$ and
each type-$3$ sum in $E_1$ by $x\mapsto x^{(3)}$.
Both are involutions of their summation sets. The intra-cluster
sums become
$\sum_{x\in E}\varphi_A(y,c)\sum_{i\le2}
 (M_i(x^{(i)})-M_i(x))\psi(x)$, which equals
$-\sum_{x\in E}\varphi_A(y,c)w_3\gamma\cdot\widehat d_3(x)\psi(x)$.
The part over $E_1$ cancels with the corresponding matrix
difference from the reindexed type-$3$ sum. The resulting row is
$(T_Y^{(1)}(A)\psi)(y)=R_0+R_1+R_2+R_3$, with
\begin{equation}\label{eq:rg-reconstruction-summation-parts}
 \begin{split}
 R_0&=(1+\mu/b_k)\frac{\eta_k^3}{48}
                         \sum_{x\in E}\varphi_A(y,c)\psi(x),\\
 R_1&=b_k^{-1}\frac{\eta_k^2}{48}
       \sum_{x\in E_2\cup E_3}
                     \varphi_A(y,c)w_3\gamma\cdot\widehat d_3(x)\psi(x),\\
 R_2&=-b_k^{-1}\frac{\eta_k^2}{48}\sum_{x\in E_2}
       \varphi_A(y,c)M_3(x)[u_A(c,c_3)\psi(x^{(3)})-\psi(x)],\\
 R_3&=-b_k^{-1}\frac{\eta_k^2}{48}\sum_{x\in E_1}
       [\varphi_A(y,c_3)u_A(c_3,c)-\varphi_A(y,c)]
                                             M_3(x^{(3)})\psi(x).
 \end{split}
\end{equation}
The face count and phase moduli bound the first three terms by
\[
 \exp(3\mathfrak i/2)\|\psi\|_{\mathrm{su}}
       \big[1+\mu/b_k+6b_k^{-1}
                           +6b_k^{-1}(1+\exp(\mathfrak i))\big].
\]
The fourth term requires the curvature estimate.

For neighboring photon sites $c,c'=c+\eta_ka$ in $\Delta_y$,
let $j(c,a)$ be the first scale at which their block centers
agree. The centers at all larger scales then agree as well.
For $j\ge2$, such a $c$ lies on a face of a block of side
$L^{j-1}\eta_k$. There are $L^{3(k-j+1)}$ such blocks inside
$\Delta_y$, with $L^{2(j-1)}$ sites on the indicated face.
For $j=1$ use all sites. We obtain
\begin{equation}\label{eq:rg-first-common-block-count}
 \#\{c:c,c+\eta_ka\in\Delta_y,\ j(c,a)=j\}
                         \le\eta_k^{-3}L^{-(j-1)}.
\end{equation}

Write
\[
 \begin{aligned}
 \varphi_A(y,c')u_A(c',c)-\varphi_A(y,c)
    &=\varphi_A(y,c)(\exp(v(c,a))-1),\\
 v(c,a)&=ie_k\eta_k
      [(\tau_kA)(y,c')+A(c',c)-(\tau_kA)(y,c)].
 \end{aligned}
\]
The contributions above scale $j(c,a)$ cancel in this bracket.
The remaining expression is the mean of the line sums over
closed paths from their common scale-$j$ center down to $c'$,
then across the bond to $c$, and back up the reversed paths.
All these loops lie in the scale-$j$ block, of side at most
$L^j\eta_k$, and have at most $3L^j$ bonds.
For a closed path of $n$ bonds in a box of side $s$, the axial
path estimate gives $|A(\Gamma)|\le2ns\|dA\|_\infty$.
To recall the reason, subtract the gauge obtained by line sums
from one corner in coordinate order $1,2,3$. The difference on
each bond is a sum over at most two strips of width one mesh
and length at most $s$. Its modulus is at most
$2s\|dA\|_\infty$, and the gauge sum around the loop vanishes.
This argument applies to complex fields as well.
Consequently, with $j=j(c,a)$,
\begin{equation}\label{eq:rg-reconstruction-phase-defect}
 |v(c,a)|\le6L^{2(j-k)}q_d,\qquad
 |\exp(v(c,a))-1|
                    \le6L^{2(j-k)}q_d\exp(6q_d).
\end{equation}
The factor $L^{2(j-k)}$ must be kept until after summing the
locations of the first common block. By
\eqref{eq:rg-first-common-block-count},
\begin{align*}
 \eta_k^2\sum_{\substack{c\in\Delta_y\\c+\eta_ka\in\Delta_y}}
                                  |\exp(v(c,a))-1|
 &\le6q_d\exp(6q_d)L^{1-k}\sum_{j=1}^kL^j\\
 &\le9Lq_d\exp(6q_d).
\end{align*}
The six directions, eight labels per direction, and
$\|M_3\|_2\le6$ now give
\[
 |R_3|_2\le54Lq_d\exp(6q_d)b_k^{-1}
                         \exp(3\mathfrak i/2)\|\psi\|_{\mathrm{su}}.
\]
Together with the first three terms this proves the first
line of \eqref{eq:rg-reconstruction-sup-estimate} for $T_Y^{(1)}$.
Take $Y$ to be the full torus to obtain the bound for $T_k$.

The second row is supported on $E_3$. Its definition and
\eqref{eq:rg-reconstruction-face-count} give directly
\[
 |(T_Y^{(2)}(A)\chi)(y)|_2
 \le b_k^{-1}\eta_k^{\alpha-1}\frac{\eta_k^3}{48}
            48\eta_k^{-2}\exp(3\mathfrak i/2)6\|\chi\|_{\mathrm{su}},
\]
which is the second bound.

For the coefficient norm, fix a row and choose each source
component to have the conjugate phase of its row coefficient.
Its component supremum is at most one, and its Euclidean
spinor supremum is at most two. The row coefficient sum is
then at most twice the corresponding bound in
\eqref{eq:rg-reconstruction-sup-estimate}. At a total connection,
\eqref{eq:rg-total-imaginary-curvature-bounds} gives
$\mathfrak i,q_d\le1$; also $\mu\le1$ and $\eta_k^\alpha\le1$.
The resulting bound is $C_T$ as stated. The normalized index
measure is already in the row coefficients, so no extra
factor $48$ enters. The domain is invariant under sign, and
the matrix $\mathsf C$ is monomial; conjugating $T_Y^{(i)}(-A)$
or $T_k(-A)$ by it permutes row and column moduli.
This proves the barred bounds.
\end{proof}

The local reconstruction respects the symmetries of an activity.
Telescoping the scalar gauge sum in the averaging paths gives
\[
 \begin{aligned}
 \varphi_{A-d\lambda}(y,c)\exp(ie_k\lambda(c))
       &=\exp(ie_k\lambda(y))\varphi_A(y,c),\\
 u_{A-d\lambda}(c,c')\exp(ie_k\lambda(c'))
       &=\exp(ie_k\lambda(c))u_A(c,c').
 \end{aligned}
\]
Insert these two identities in
\eqref{eq:rg-reconstruction-two-pieces}. Every term acquires
the coarse factor $\exp(ie_k\lambda(y))$, where the boundary
leg transforms with the phase of its first point. Thus, with
$\lambda^{\mathrm u}$ the restriction to unit sites,
\begin{equation}\label{eq:rg-reconstruction-gauge-covariance}
 \begin{split}
 T_Y^{(1)}(A-d\lambda)\exp(ie_k\lambda)\psi
      &=\exp(ie_k\lambda^{\mathrm u})T_Y^{(1)}(A)\psi,\\
 T_Y^{(2)}(A-d\lambda)\chi^\lambda
      &=\exp(ie_k\lambda^{\mathrm u})T_Y^{(2)}(A)\chi,
 \qquad\chi^\lambda(x,x')=\exp(ie_k\lambda(c))\chi(x,x').
 \end{split}
\end{equation}
The identity also holds for gauges specified locally on the
bonds and sites read by the row.

A scheme symmetry $\tilde g=(S,a)$ maps nested blocks and their
centers to their images, preserves the order of each coordinate
path, and permutes the $48$ labels. Moreover
$M_i(\tilde gx)S=SM_i(x)$ and
$(\tilde gx)^{(i)}=\tilde g(x^{(i)})$.
Changing variables in the two rows therefore proves
\begin{equation}\label{eq:rg-reconstruction-lattice-covariance}
 T_{\tilde gY}^{(i)}(\tilde gA)\tilde g
                           =\tilde gT_Y^{(i)}(A).
\end{equation}
Complex conjugation changes the phases to those at $-A$, and
$\overline{M_i}=\varkappa^{-1}M_i\varkappa$, giving
\begin{equation}\label{eq:rg-reconstruction-reality}
 \overline{T_Y^{(i)}(\overline A)}
                         =\varkappa^{-1}T_Y^{(i)}(-A)\varkappa.
\end{equation}

The barred rows are $\bar T_Y^{(i)}(A)=\mathsf C^{-1}T_Y^{(i)}(-A)\mathsf C$.
Their gauges have the opposite sign. For the lattice covariance,
$(S^{-1})^T=\sigma(S)\mathsf C^{-1}S\mathsf C$ with
$\sigma(S)\in\{\pm1\}$: both matrix maps are homomorphisms,
and on each generating Dirac matrix they differ by a minus
sign. This same sign occurs in the source and target spin
actions and cancels in the barred intertwining identity.
For reality, $\mathsf C$ is real and
$\varkappa\mathsf C=-\mathsf C\varkappa$; the two signs from
commuting the conjugations cancel. Thus the barred versions
of \eqref{eq:rg-reconstruction-lattice-covariance} and
\eqref{eq:rg-reconstruction-reality} hold with the inverse
transpose spin action and the same $\varkappa$.

For $B\in\widetilde R_k+\mathfrak X_k$, equations
\eqref{eq:rg-reconstruction-left-inverse} and
\eqref{eq:rg-reconstruction-local-identity} now define the
local substitution used after fluctuation integration:
\[
 \Psi(y)\longmapsto(T_Y^{(1)}(B)\psi)(y)+(T_Y^{(2)}(B)\chi)(y),\qquad
 \bar\Psi(y)\longmapsto(\bar T_Y^{(1)}(B)\bar\psi)(y)
                                      +(\bar T_Y^{(2)}(B)\bar\chi)(y).
\]
After setting $\psi=H_k(B)\Psi$ and
$\chi=\delta_{\alpha,B}\psi$, and their barred analogues,
this substitution is the identity
on the coarse generators over $Y$. Its coefficients are
holomorphic in $B$ and satisfy the uniform row bounds above.
The new legs lie over $Y$ and its admissible boundary pairs.
The substitution fixes any old generators, so the full polynomial
has these support properties when its old legs already have them.
The weighted linear-substitution
bound therefore applies with these two row operators. The
resulting degree weights will be fixed with the full step.

\subsection{Adjustment of the quadratic kernel}
\label{subsec:rg-adjustment}

The extraction makes the integral of the zero-field quadratic
kernel vanish. For contraction we need the kernel itself to
vanish. We achieve this by writing each ordinary leg as a
transported average and a sum of H\"older differences. This
changes the activity before evaluation and leaves its evaluated
value unchanged.

Let $X$ be small, choose a lift $\widetilde X$, and let
$x=(c,g)$ and $\tilde x=(\tilde c,\tilde g)$ lie over this lift.
Take all shortest face chains of its $M$-cubes between their
two cubes. For a chain $(\square_0,\ldots,\square_n)$ with
centers $z_i$, use the points
\[
 p_0=x,\qquad p_{i+1}=(z_i,g)\quad(0\le i\le n),\qquad
 p_{n+2}=\tilde x.
\]
The initial label is kept at each intermediate center.
Each distinct consecutive pair is admissible for $X$ of the
first kind. The paths in the first and last cube remain in
those cubes; between adjacent centers all six paths are the
straight segment across their common face. The last pair
can have different labels, and can lie inside one cluster.
For a coincident pair there is no difference term.

With $u_B(p,p')=U_B(p,p')$ and $u_B(p,p)=1$, set
\[
 U_i^\pm=\prod_{j<i}u_B(p_j,p_{j+1})^{\pm1}.
\]
The chain has no repeated cube, and
$n+1\le|X|_M\le c_X(L+1)$. Hence
\begin{equation}\label{eq:rg-chain-length}
 \sum_{i=0}^{n+1}|\pi p_i-\pi p_{i+1}|_1\le(n+3)M
          \le3c_XLM,\qquad
 \sum_{i:p_i\ne p_{i+1}}d_X(p_i,p_{i+1})^\alpha\le3c_XLM.
\end{equation}
The second estimate uses $d_X^\alpha\le\max(1,d_X)\le M$.
The transport bound is
$|u_B(p,p')^{\pm1}|\le\exp(e_k\|\operatorname{Im}B\|_\infty
|\pi p-\pi p'|_1)$, obtained by counting the bonds in each path.

For $B=\mathscr A+\mathscr C$, the domain decomposition and
$e_k^\varepsilon\le1/4$ give
$e_k\|\operatorname{Im}B\|_\infty\le e_k^\varepsilon/2$.
Assume in addition
\begin{equation}\label{eq:rg-adjustment-charge-condition}
 90c_XLM e_k^\varepsilon\le1.
\end{equation}
The modulus of every partial product $U_i^\pm$ is then at most
$\exp(1)$, uniformly in the chain and mesh.

Let $\operatorname{Mn}$ be the mean over the shortest chains.
Define the homomorphism $\Theta_{X,B}$ on ordinary generators by
\begin{equation}\label{eq:rg-chain-homomorphism}
 \begin{split}
 \Theta_{X,B}\psi(x)
  =\operatorname{Vol}(X)^{-1}\int_{\widetilde X}\dd\tilde x\,
    \operatorname{Mn}\left[
       U_{n+2}^+\psi(\tilde x)
       -\sum_{i:p_i\ne p_{i+1}}
           U_i^+d_X(p_i,p_{i+1})^\alpha\chi(p_i,p_{i+1})\right],\\
 \Theta_{X,B}\bar\psi(x)
  =\operatorname{Vol}(X)^{-1}\int_{\widetilde X}\dd\tilde x\,
    \operatorname{Mn}\left[
       U_{n+2}^-\bar\psi(\tilde x)
       -\sum_{i:p_i\ne p_{i+1}}
           U_i^-d_X(p_i,p_{i+1})^\alpha\bar\chi(p_i,p_{i+1})\right].
 \end{split}
\end{equation}
The spin components are unchanged. Both displayed images are
linear combinations of odd independent generators, so they
anticommute and determine a unique homomorphism on the
subalgebra of ordinary legs over $X$. Translation of the lift
by a period preserves every transport and chain and multiplies
each generator by the same antiperiodic sign. Thus the
homomorphism is independent of this choice of lift.

If the generators are evaluated at fields satisfying their
covariant H\"older relations, repeated substitution of
\[
 \psi(p_i)=u_B(p_i,p_{i+1})\psi(p_{i+1})
                    -d_X(p_i,p_{i+1})^\alpha\chi(p_i,p_{i+1})
\]
along a chain gives exactly the first expression in
\eqref{eq:rg-chain-homomorphism}; the barred relation uses
$u_B^{-1}$. Taking the two means therefore reproduces the
original leg. This proves equality after evaluation for
every polynomial in the ordinary legs over $X$.

The adjusted family is
\begin{equation}\label{eq:rg-adjustment-definition}
 \mathscr E^{\mathrm{nat}}(X)
   =\mathscr E(X)-\mathscr E^{[2,0]}(X)
           +\Theta_{X,\mathscr A+\mathscr C}
                                      \mathscr E^{[2,0]}(X)
                                      \qquad(X\text{ small}),
\end{equation}
and is unchanged on large polymers. The equality just proved
shows that the adjusted and original families have equal
evaluated values. Only bidegrees $(2,0),(1,1),(0,2)$ can change.

At zero connection all $U_i^\pm$ are one. The ordinary part
of an adjusted bilinear with kernel $K$ is therefore
\[
 \bar\psi_{\mathrm{av}}^T
    \left[\int_{\widetilde X}\dd x\,\dd x'\,K(x,x')\right]
                                                     \psi_{\mathrm{av}},
 \qquad
 \psi_{\mathrm{av}}=\operatorname{Vol}(X)^{-1}
                                        \int_{\widetilde X}\dd x\,\psi(x).
\]
Apply this to $R\mathscr E$, whose kernel integral vanishes.
With $E=(R\mathscr E)^{\mathrm{nat}}$, we obtain the stronger
normalization
\begin{equation}\label{eq:rg-adjusted-kernel-normalization}
 E_{00}(X;0,0)=0,\qquad E_{20}(X;0,0;\cdot,\cdot)=0
                                                     \quad(X\text{ small}).
\end{equation}

The same chain estimate controls the norm of adjustment.
Decompose $\Theta\psi=L_\psi+L_\chi$ by bidegree. For arbitrary
generator weights $(w_1,w_2)$,
\[
 \|L_\psi\|_w\le\exp(1)w_1,\qquad
 \|L_\chi\|_w\le\exp(1)(3c_XLM)w_2.
\]
The normalized site and label mean has total mass one.
Writing $s_K=\sum_{\beta,\beta'}\int|K_{\beta\beta'}|$ for
the balanced quadratic kernel and $F=\Theta\mathscr E^{[2,0]}$
gives the three bounds
\[
 \begin{aligned}
 \|F^{[2,0]}\|_w&\le\exp(2)w_1^2s_K,\\
 \|F^{[1,1]}\|_w&\le2\exp(2)(3c_XLM)w_1w_2s_K,\\
 \|F^{[0,2]}\|_w&\le\exp(2)(3c_XLM)^2w_2^2s_K.
 \end{aligned}
\]
Since $h_{k,2}=e_k^\varepsilon h_{k,1}$ and
$3c_XLM e_k^\varepsilon\le1/30$, summing these bounds yields
\begin{equation}\label{eq:rg-adjustment-norm}
 \|\mathscr E^{\mathrm{nat}}(X)\|_{h_k}
 \le\exp(2)(1+3c_XLM e_k^\varepsilon)^2\|\mathscr E(X)\|_{h_k}
                                               \le10\|\mathscr E(X)\|_{h_k}.
\end{equation}
Subtracting the original quadratic part before adding its
image uses additivity of the norm over bidegrees.

All new legs lie over $X$ or at its admissible chain pairs,
and all new transports read bonds with both ends in $X$.
Thus adjustment preserves locality. It preserves analyticity
because these transports are entire in the connection.
For the lattice covariance, symmetries carry shortest chains,
their centers and labels to the corresponding image chains.
For a real gauge the factors along a partial chain telescope
to the phase difference at its endpoints, giving the original
leg phase. Charge conjugation replaces $u_B$ by $u_B^{-1}$;
the reality map uses $\overline{u_{-\overline B}}=u_B$.
These identities on generators prove all four covariance
identities for $\Theta$, and hence the activity symmetries
for adjustment. Restriction to $\mathscr C=0$ commutes with
adjustment because it leaves the same total connection and
the same ordinary quadratic kernel. Therefore adjustment
preserves the complex and real activity spaces.

\subsection{Scaling and reblocking}\label{subsec:rg-scaling}

For the rest of this group let both $\ell$ and $\ell/L$ be
admissible sides. In particular, $\ell\ge13L^2M$.
Set $\bar\varepsilon=1/200$ and assume
\begin{equation}\label{eq:rg-contraction-parameters}
 \begin{gathered}
 0<\varepsilon\le\bar\varepsilon,\qquad 1-\varepsilon<\alpha<1,\qquad
 L^{1/8-3\bar\varepsilon/2}\ge2400,\qquad e_k^\varepsilon\le1/4,\\
 160L^{1+3\varepsilon/2}M e_k^\varepsilon\le1,\qquad
 9ML^\varepsilon e_k^{2\varepsilon}\le1.
 \end{gathered}
\end{equation}
The adjustment condition \eqref{eq:rg-adjustment-charge-condition}
is also retained. These inequalities will control the analytic
disk and the phases on a small polymer.

Write the three thresholds of $R_k$ as
\[
 \theta_1=e_k^{-3/4+\varepsilon},\qquad
 \theta_2=e_k^\varepsilon\theta_1,\qquad
 \theta_3=e_k^\varepsilon\theta_2,
\]
and the thresholds of $\mathfrak X_k$ as
$d_1=e_k^{-1+2\varepsilon}$ and $d_2=e_k^\varepsilon d_1$.
For a level $k+1$ bond function and gauge function define
\begin{equation}\label{eq:rg-connection-scaling}
 A_L(b)=L^{-1/2}A(b/L),\qquad
 \lambda_L(c)=L^{1/2}\lambda(c/L).
\end{equation}
Since $\eta_{k+1}=\eta_k/L$, the finite differences obey
$(d\lambda)_L=d\lambda_L$ and
\[
 \|A_L\|_\infty=L^{-1/2}\|A\|_\infty,\qquad
 \|\partial A_L\|_\infty=L^{-3/2}\|\partial A\|_\infty,\qquad
 \|\delta_\alpha\partial A_L\|_\infty
                \le L^{-3/2-\alpha}\|\delta_\alpha\partial A\|_\infty.
\]
The last inequality uses only the coarse pairs at distance
$\le1/L$, a subset of those in the coarse seminorm.
The charge relation $e_{k+1}=L^{1/2}e_k$ therefore gives
the following fine-scale bounds for the scaled complex
small field, and for each scaled local real representative:
\begin{equation}\label{eq:rg-scaled-field-thresholds}
 \begin{split}
 b_1&=L^{-7/8+\varepsilon/2}\theta_1,\\
 b_2&=L^{-15/8+\varepsilon}\theta_2,\\
 b_3&=L^{-15/8+3\varepsilon/2}\theta_3.
 \end{split}
\end{equation}
Here the additional factor $L^{-\alpha}$ in the third bound
has been discarded. For the external argument the bounds are
\begin{equation}\label{eq:rg-scaled-external-thresholds}
 \|\mathscr C_L\|_\infty<L^{-1+\varepsilon}d_1,\qquad
 \|\partial\mathscr C_L\|_\infty<L^{-2+3\varepsilon/2}d_2.
\end{equation}
The enlarged coarse cube used to define a local representative,
after multiplication by $L$, contains the enlarged fine cube.
Scaling its real gauge thus gives the corresponding fine
representative. All the ratios in
\eqref{eq:rg-scaled-field-thresholds} and
\eqref{eq:rg-scaled-external-thresholds} are at most $L^{-3/4}$,
by \eqref{eq:rg-contraction-parameters}. It follows that
\begin{equation}\label{eq:rg-scaled-domain-inclusion}
 \mathscr A_L\in L^{-3/4}\widetilde R_k
                     \subset\widetilde R_k^-,\qquad
 \mathscr C_L\in L^{-3/4}\mathfrak X_k.
\end{equation}
For the second inclusion in the first statement use
$L^{-3/4}\le1/2\le1-e_k^\varepsilon$.

The Grassmann scaling homomorphism is
\begin{equation}\label{eq:rg-leg-scaling}
 S_L\psi(x)=L^{-1}\psi'(x/L),\qquad
 S_L\chi(x,x')=L^{-1-\alpha}\chi'(x/L,x'/L),
\end{equation}
with the same powers for the barred generators.
Only the photon coordinate is divided by $L$; the chamber
label and spin index are kept. Ordinary indices are mapped
bijectively onto the coarse ordinary indices. Pair indices
are mapped bijectively onto the coarse pair indices with
$d_X\le M/L$, since $d_X$ is divided by $L$.
This includes a pair inside one cluster:
$d_X=\eta_k$ becomes $d_X=\eta_{k+1}$.
The antiperiodic signs agree under this bijection.

Thus distinct basis monomials have distinct images, up to
their ordering signs, and every monomial of bidegree $(n,m)$
is multiplied by $L^{-n-(1+\alpha)m}$. There is no cancellation
of coefficients. For arbitrary positive weights,
\begin{equation}\label{eq:rg-scaling-norm-identity}
 \|S_LF\|_{w_1,w_2}=\|F\|_{L^{-1}w_1,L^{-1-\alpha}w_2}.
\end{equation}
As $h_{k+1,1}/h_{k,1}=L^{-1/8}$ and
$h_{k+1,2}/h_{k,2}=L^{-1/8+\varepsilon/2}$, this gives
\begin{equation}\label{eq:rg-degree-scaling}
 \|(S_LF)^{[n,m]}\|_{h_{k+1}}
       =L^{-p(n,m)}\|F^{[n,m]}\|_{h_k},\qquad
 p(n,m)=n+(1+\alpha)m+\frac{n+m}{8}-\frac{m\varepsilon}{2}.
\end{equation}
The powers are identical for all admissible pair types.

In kernel notation the same calculation reads
\begin{equation}\label{eq:rg-scaled-kernel}
 (S_LF)_{nm}(\xi';\zeta')=
 \begin{cases}
 L^{3n+6m-n-(1+\alpha)m}F_{nm}(L\xi';L\zeta'),
                       &d_X(\zeta'_j)\le M/L\text{ for every }j,\\
 0,                    &\text{otherwise}.
 \end{cases}
\end{equation}
Indeed, the index weights satisfy
$\eta_k^3/48=L^3\eta_{k+1}^3/48$ and hence contribute
$L^{3n+6m}$ to the kernel. Antisymmetry and antiperiodicity
are preserved, so uniqueness of the kernel expansion gives
\eqref{eq:rg-scaled-kernel}.

Let $\bar X^L$ be the union of $LM$-cubes containing the cubes
of $X$. Define the reblocked family by
\begin{equation}\label{eq:rg-reblocking-definition}
 (\mathcal L F)(Y;\mathscr A,\mathscr C)
       =\sum_{X:\bar X^L=LY}S_L
                          F(X;\mathscr A_L,\mathscr C_L).
\end{equation}
The sum is finite on a fixed torus. It is local in $Y$ if
$F$ is local in each $X$. Ordinary legs over $X$ scale to
legs over $Y$; the six paths for a pair of the first kind
scale to paths in $Y$. For a boundary pair,
$x^{(3)}/L=(x/L)^{(3)}$, and an endpoint in $X$ scales to
an endpoint in $Y$. Finally, every fine bond meeting $X$
scales to a coarse bond meeting $Y$. Equality of both
coarse connection arguments on these bonds consequently
gives equality of the fine arguments read by each term.
This proves both the support and the coefficient locality
of \eqref{eq:rg-reblocking-definition}.

\subsection{Global representatives on a small polymer}
\label{subsec:rg-small-polymer-gauge}

The scaled fields belong to the fine domain, but this alone
does not give an analytic gain for their vacuum and ordinary
quadratic kernels. On a small polymer they admit a stronger
representation, with a real gauge, two constant fields and
global fields small enough for an analytic disk of radius
$L^{7/4}$.

\begin{proposition}\label{prop:rg-small-polymer-representatives}
Assume \eqref{eq:rg-contraction-parameters}. Let $X$ be small
and let $(\mathscr A,\mathscr C)\in
\widetilde R_{k+1}\times\mathfrak X_{k+1}$.
There are a real function $\lambda_*$, constant bond
functions $\mathscr K_1,\mathscr K_2$, and global bond
functions $\mathscr A^{\mathrm b},\mathscr C^{\mathrm b}$
on the fine torus such that, on every bond with an endpoint
at distance at most $1$ from $X$,
\begin{equation}\label{eq:rg-small-polymer-gauge-identity}
 \mathscr A_L-d\lambda_*=\mathscr A^{\mathrm b}+\mathscr K_1,
 \qquad
 \mathscr C_L=\mathscr C^{\mathrm b}+\mathscr K_2.
\end{equation}
They satisfy
\begin{equation}\label{eq:rg-small-polymer-global-bounds}
 \begin{split}
 \|\mathscr A^{\mathrm b}\|_\infty
          &\le64L^{-15/8+3\varepsilon/2}\theta_1,\\
 \|\partial\mathscr A^{\mathrm b}\|_\infty
          &\le137L^{-15/8+3\varepsilon/2}\theta_2,\\
 \|\delta_\alpha\partial\mathscr A^{\mathrm b}\|_\infty
          &\le1382L^{-15/8+3\varepsilon/2}\theta_3,\\
 \|\mathscr C^{\mathrm b}\|_\infty&<L^{-7/4}d_1,\qquad
 \|\partial\mathscr C^{\mathrm b}\|_\infty<L^{-7/4}d_2,
 \end{split}
\end{equation}
and
\begin{equation}\label{eq:rg-small-polymer-constant-bounds}
 \|\mathscr K_1\|_\infty<b_1,\qquad
 \|\mathscr K_2\|_\infty<L^{-1+\varepsilon}d_1.
\end{equation}
Consequently
\begin{equation}\label{eq:rg-small-polymer-analytic-disk}
 (t\mathscr A^{\mathrm b},t\mathscr C^{\mathrm b})
                       \in R_k\times\mathfrak X_k
                             \qquad(|t|\le L^{7/4}).
\end{equation}
The construction uses a base point of $X$ and a decomposition
and local real representative of $\mathscr A$, and is
independent of the activity. When $\mathscr C=0$, it gives
$\mathscr C^{\mathrm b}=\mathscr K_2=0$.
\end{proposition}

\begin{proof}
Choose a lift of $X$ and a photon site $c_*$ in it.
Set $r=2LM\ge50$ and $K_{\mathrm{loc}}=K(c_*,2r+5)$, with
sup balls understood through their nonwrapping lifts.
Choose the coarse $M$-cube containing $c_*/L$.
If $\mathscr A=A_0+A_1$ is its domain decomposition, its
real representative $A'$ and gauge $\lambda$ on the
enlarged cube give
\[
 a_0=(A_0)_L=a'_0+d\lambda_L,\qquad
 a'_0=(A')_L,\qquad a_1=(A_1)_L
\]
on $K_{\mathrm{loc}}$. To see that these functions are
defined there, the scaled enlarged cube has side
$(2L+1)LM$ and center $z$ with $|c_*-z|<LM/2$.
The inequality $4LM+6\le L^2M$ puts
$K(c_*,2r+6)$ inside this cube. Moreover
$4r+11<\ell$, since $\ell\ge13L^2M$.
Thus the boxes and all derivative cells used below do
not wrap.

We first record the cutoff estimate with its derivative
support. Put
\[
 f(t)=
 \begin{cases}
 1,&|t|\le1,\\
 1-2(|t|-1)^2,&1\le|t|\le3/2,\\
 2(2-|t|)^2,&3/2\le|t|\le2,\\
 0,&|t|\ge2.
 \end{cases}
\]
For a bond $(c,c+\eta_ke_\mu)$ use the product of the three
functions $f((c_\nu+\eta_k\delta_{\mu\nu}/2-c_{*,\nu})/r)$
as its cutoff multiplier. Denote this operation by
$\operatorname{cut}_r$. The function $f$ is continuously
differentiable, $|f'|\le2$, and $f'$ is Lipschitz with
constant $4$. If a bond function $b$ has local norms
$B_0,B_1,B_\alpha$ on $K_{\mathrm{loc}}$, extension by zero
after multiplication gives
\begin{equation}\label{eq:rg-cutoff-estimate}
 \begin{split}
 \|\operatorname{cut}_r b\|_\infty&\le B_0,\\
 \|\partial\operatorname{cut}_r b\|_\infty
                                      &\le B_1+2B_0/r,\\
 \|\delta_\alpha\partial\operatorname{cut}_r b\|_\infty
                             &\le B_\alpha+12B_1/r+12B_0/r^2.
 \end{split}
\end{equation}
Indeed, the multiplier has oscillation at most
$6|c-c'|/r$, first difference at most $2/r$, and
first-difference oscillation at most $12|c-c'|/r^2$.
Apply the product rule to the first difference.
In its H\"older difference the four terms are bounded
by $B_\alpha |c-c'|^\alpha$, twice
$6B_1|c-c'|/r$, and $12B_0|c-c'|/r^2$.
Here $|c-c'|\le1$ implies
$|c-c'|\le|c-c'|^\alpha$.
If any such term is nonzero, all its cells lie in
$K(c_*,2r+5)$: a nonzero multiplier or its difference
requires a site within $2r+1$, the first-difference
cell adds one mesh length, and the pair adds at most
$1$. The same bound covers the two-step cell at
$k=0$. Thus the local estimates suffice for the global
ones. The cutoff equals one on bonds with both ends in
$K(c_*,r)$.

For $b=a'_0$ or $a_1$, write $v=c-c_*$ and define
\[
 V[b]_{\nu\mu}=(\partial_\nu b_\mu)(c_*),\qquad
 \operatorname{Lin}[b]_\mu(c)=\sum_\nu V[b]_{\nu\mu}v_\nu,
 \qquad
 \operatorname{rem}[b]=b-b(c_*)-\operatorname{Lin}[b].
\]
Its first difference is
$\partial_\nu b_\mu(c)-\partial_\nu b_\mu(c_*)$,
and its H\"older first-difference seminorm is at most
$b_3$. Join the two sites by at most $2r+5$ simultaneous
coordinate steps of length at most $1$. Their difference
is therefore at most $(2r+5)b_3\le3rb_3$. A rectilinear
path to $c_*$ has length at most $3(2r+5)$, and the
remainder vanishes at $c_*$. Keeping the sharper derivative
bound along this path and using $r\ge50$ gives
\[
 \|\operatorname{rem}[b]\|_{K_{\mathrm{loc}}}
              \le3(2r+5)^2b_3\le18r^2b_3,\qquad
 \|\partial\operatorname{rem}[b]\|_{K_{\mathrm{loc}}}
              \le3rb_3.
\]
Equation \eqref{eq:rg-cutoff-estimate} consequently bounds
the three global norms of its cutoff by
$18r^2b_3$, $39rb_3$, and $253b_3$, respectively.

Only a real gauge is available. We remove the constant
and symmetric linear part of $a'_0$, leaving its curvature
and the complex linear part of $a_1$. Set
\[
 V^0=V[a'_0],\qquad
 V^{\mathrm s}=(V^0+(V^0)^T)/2,\qquad
 F=V^0-(V^0)^T,
\]
and use the real quadratic function
\begin{equation}\label{eq:rg-real-quadratic-gauge}
 q_*(c)=\sum_\mu\left[a'_{0,\mu}(c_*)-
                    \frac{\eta_k}{2}V^{\mathrm s}_{\mu\mu}\right]v_\mu
              +\frac12\sum_{\nu,\mu}V^{\mathrm s}_{\nu\mu}v_\nu v_\mu.
\end{equation}
Its finite difference is exactly
$(dq_*)_\mu(c)=a'_{0,\mu}(c_*)+
\sum_\nu V^{\mathrm s}_{\nu\mu}v_\nu$.
The mesh term in \eqref{eq:rg-real-quadratic-gauge} cancels
the diagonal term from differencing the quadratic polynomial.
With $\mathscr K_1=a_1(c_*)$ and
\[
 P_\mu(c)=\sum_\nu
                    \left[\frac12F_{\nu\mu}+V[a_1]_{\nu\mu}\right]v_\nu,
\]
we have the local identity
\begin{equation}\label{eq:rg-local-taylor-gauge}
 \mathscr A_L-d(\lambda_L+q_*)=
       \mathscr K_1+P+\operatorname{rem}[a'_0]+\operatorname{rem}[a_1].
\end{equation}

To cut off $P$ we need a bound that also uses the side
$\ell$. A global mean-zero periodic function whose
oscillation between sites at distance at most $1$ is at
most $B$ has sup norm at most $\ell B$. Any two sites
can be joined by at most $\lceil\ell/2\rceil\le\ell$
simultaneous coordinate steps of length at most $1$;
comparison with the mean gives the claim.
Every $\partial_\nu a_{1,\mu}$ has mean zero, sup norm
at most $b_2$, and this oscillation at most $b_3$.
For $a_0$, the global curl agrees on each derivative
cell with the curl of a local real representative.
Its sup norm is at most $2b_2$, its oscillation at most
$2b_3$, and it has mean zero. The cells of a pair at
distance at most $1$ lie in one enlarged cube, so that
the same representative can be used for their comparison.
It follows that every coefficient of $P$ is bounded by
\begin{equation}\label{eq:rg-linear-slope-bound}
 s=2\min(b_2,\ell b_3).
\end{equation}

Use a second radius
\[
 r'=\min(e_k^{-\varepsilon},\ell/5).
\]
The charge condition in \eqref{eq:rg-contraction-parameters}
gives $r'\ge r$, and $4r'+11<\ell$. Thus its cutoff is
also defined without wrapping. The three norms of $P$
on its support are at most $6.3r's$, $s$, and $0$.
Equation \eqref{eq:rg-cutoff-estimate}, with $r'$ in
place of $r$, bounds the cutoff norms by
$6.3r's$, $13.6s$, and $87.6s/r'$.
If $r'=e_k^{-\varepsilon}$, use $s\le2b_2$;
in units of $(\theta_1,\theta_2,\theta_3)$ these are
at most $(12.6,27.2,175.2)L^{-15/8+\varepsilon}$.
If $r'=\ell/5<e_k^{-\varepsilon}$, use
$s\le2\ell b_3$. Since $\ell e_k^\varepsilon<5$,
the corresponding three bounds are at most
$(63,136,876)L^{-15/8+3\varepsilon/2}$.

Define the global field and real gauge by
\begin{equation}\label{eq:rg-global-representative-construction}
 \begin{split}
 \mathscr A^{\mathrm b}
   &=\operatorname{cut}_r
             \big(\operatorname{rem}[a'_0]+\operatorname{rem}[a_1]\big)
                  +\operatorname{cut}_{r'}P,\\
 \lambda_*&=\lambda_L+q_*\quad\text{on }K_{\mathrm{loc}},
                  \qquad \lambda_*=0\quad\text{elsewhere}.
 \end{split}
\end{equation}
The two remainder cutoffs contribute
$(36r^2,78r,506)b_3$. As $r=2LM$, the charge condition
$160L^{1+3\varepsilon/2}Me_k^\varepsilon\le1$ bounds
the first two by $L^{-15/8+3\varepsilon/2}\theta_1$
and $L^{-15/8+3\varepsilon/2}\theta_2$, respectively.
Adding the linear contribution gives the constants
$64$, $137$, and $1382$ in
\eqref{eq:rg-small-polymer-global-bounds}.

In the external argument take
\[
 \mathscr K_2=\mathscr C_L(c_*),\qquad
 \mathscr C^{\mathrm b}
                   =\operatorname{cut}_r(\mathscr C_L-\mathscr K_2).
\]
A rectilinear path to $c_*$ bounds the undifferentiated
field on $K_{\mathrm{loc}}$ by
$6.3r\|\partial\mathscr C_L\|_\infty$.
The cutoff bounds therefore give
\[
 \|\mathscr C^{\mathrm b}\|_\infty
       \le6.3r\|\partial\mathscr C_L\|_\infty,\qquad
 \|\partial\mathscr C^{\mathrm b}\|_\infty
       \le13.6\|\partial\mathscr C_L\|_\infty.
\]
Use \eqref{eq:rg-scaled-external-thresholds} and
$d_2=e_k^\varepsilon d_1$. The first charge condition
involving $160$ gives the desired bound for the sup norm;
the large-$L$ condition gives
$13.6L^{-2+3\varepsilon/2}\le L^{-7/4}$ for the derivative.
This proves the external bounds and the constant bounds.

The diameter estimate in Appendix~\ref{app:polymer-geometry}
puts $X$ in $K(c_*,M(L+2))$. Every bond with an endpoint
within distance $1$ of $X$ has both ends in $K(c_*,r)$,
since $M(L+2)+2\le2LM$. All cutoffs in the local identity
are one there, and the gauge is $\lambda_L+q_*$ there.
This proves \eqref{eq:rg-small-polymer-gauge-identity}.
Finally, the large-$L$ condition implies
$1382L^{-15/8+3\varepsilon/2}<L^{-7/4}$.
Multiplication by any $|t|\le L^{7/4}$ consequently leaves
all five norms strictly below their domain thresholds,
proving \eqref{eq:rg-small-polymer-analytic-disk}.
\end{proof}

\begin{proposition}\label{prop:rg-small-polymer-kernel-comparison}
Let $F(X)$ be a level-$k$ Grassmann activity holomorphic on
$\widetilde R_k\times\mathfrak X_k$, with the kernel support
and coefficient locality specified in
\S\ref{subsec:rg-activity-spaces}, and covariant under the
independent real gauges of its two connection arguments.
For the fields of Proposition~\ref{prop:rg-small-polymer-representatives},
its kernels satisfy
\begin{equation}\label{eq:rg-small-polymer-kernel-comparison}
 |F_{nm}(X;\mathscr A_L,\mathscr C_L;\xi;\zeta)|
   \le\exp(2(n+m))
        |F_{nm}(X;\mathscr A^{\mathrm b},\mathscr C^{\mathrm b};\xi;\zeta)|.
\end{equation}
For the vacuum kernel the two values are equal.
\end{proposition}

\begin{proof}
For each kernel argument $a$ let $c(a)$ be the photon site
of its ordinary leg, or the first photon site of its pair
leg, and let $\omega(a)\in\{0,1\}$ distinguish unbarred and
barred legs. On kernels the gauge homomorphism has the
form
\begin{equation}\label{eq:rg-kernel-gauge-phase}
 (G_\lambda F)_{nm}(\xi;\zeta)
       =F_{nm}(\xi;\zeta)
          \prod_a\exp\{(-1)^{\omega(a)}ie_k\lambda(c(a))\}.
\end{equation}
This follows by applying the homomorphism to the kernel
expansion; the multiplying product is symmetric in each
set of arguments, and hence preserves their antisymmetry.
The gauge relation \eqref{eq:rg-small-polymer-gauge-identity}
first gives
\[
 G_{\lambda_*}F(X;\mathscr A^{\mathrm b}+\mathscr K_1,
                         \mathscr C^{\mathrm b}+\mathscr K_2)
                       =F(X;\mathscr A_L,\mathscr C_L).
\]
Both pairs lie in the domains: a constant changes only
the sup norm, and
$L^{-7/4}+L^{-7/8+\varepsilon/2}<1$ and
$L^{-7/4}+L^{-1+\varepsilon}<1$.
Every phase from this real gauge has modulus one.

We remove the constant fields by continuation of a real
gauge identity. Write $\mathscr K_j=K_j^{\mathrm R}+iK_j^{\mathrm I}$
with real constant fields. On the connected domain
\[
 \Omega=\{z\in\C^4:|\operatorname{Im}z_1|<2,\,
       |\operatorname{Im}z_2|<2,\,|z_3|<2,\,|z_4|<2\},
\]
set
\[
 A(z)=\mathscr A^{\mathrm b}+z_1K_1^{\mathrm R}+z_2K_1^{\mathrm I},
 \qquad
 C(z)=\mathscr C^{\mathrm b}+z_3K_2^{\mathrm R}+z_4K_2^{\mathrm I},
\]
and set
\[
 \Lambda_z(c)=-
       (z_1K_1^{\mathrm R}+z_2K_1^{\mathrm I}
                  +z_3K_2^{\mathrm R}+z_4K_2^{\mathrm I})\cdot(c-c_*)
                    \quad(c\in K(c_*,r)),
\]
with value zero elsewhere. The first connection belongs to
$\widetilde R_k$: its complex small summand has sup norm
below $(L^{-7/4}+4L^{-7/8+\varepsilon/2})\theta_1<\theta_1$
and derivatives those of $\mathscr A^{\mathrm b}$, and its
remaining constant summand is real. Invariance under real
constant additions handles its unrestricted real part.
The external connection has sup norm below
$(L^{-7/4}+4L^{-1+\varepsilon})d_1<d_1$, and derivative
that of $\mathscr C^{\mathrm b}$. The large-$L$ condition
implies these two strict inequalities.

For fixed kernel arguments the function
\[
 H(z)=F_{nm}(X;A(z),C(z);\xi;\zeta)
             \prod_a\exp\{(-1)^{\omega(a)}ie_k\Lambda_z(c(a))\}
\]
is holomorphic on $\Omega$. When $z$ is real, split
$\Lambda_z$ into its first two and last two terms.
These are independent real gauge functions for the two
connections, with differentials equal to the negatives
of their respective constants on $K(c_*,r)$. Every bond
read by the local gauge condition lies in this ball.
The real gauge identity and
\eqref{eq:rg-kernel-gauge-phase} therefore give
$H(z)=H(0)$ for $z\in\R^2\times(-2,2)^2$.
All derivatives of $H-H(0)$ at zero vanish, so its
holomorphic power series is zero there. The identity
theorem extends this equality to the connected domain
$\Omega$. Evaluating at $z=(1,i,1,i)$ gives the identity
for the two original complex constants.

It remains to bound the inverse phases. By locality,
an ordinary argument lies over $X$. For a pair of the
first kind its first photon site also lies in $X$.
For a boundary pair that site has distance at most
$\eta_k\le1$ from $X$. Hence in all nonzero kernel terms
\[
 |c(a)-c_*|_1\le3(M(L+2)+1)\le9LM.
\]
The constant bounds and the charge conditions give
\[
 \begin{split}
 e_k|K_1^{\mathrm I}\cdot(c(a)-c_*)|
       &\le9LM L^{-7/8+\varepsilon/2}e_k^{1/4+\varepsilon}
                                                    \le1,\\
 e_k|K_2^{\mathrm I}\cdot(c(a)-c_*)|
       &\le9ML^\varepsilon e_k^{2\varepsilon}\le1.
 \end{split}
\]
Each leg phase, and its inverse, thus has modulus at most
$\exp(2)$. Combining this bound with the unit modulus
phases from $\lambda_*$ proves
\eqref{eq:rg-small-polymer-kernel-comparison}.
The vacuum has no leg phases, which proves its equality.
\end{proof}

\subsection{Contraction of the normalized activity}
\label{subsec:rg-activity-contraction}

The ordinary quadratic kernel now vanishes at zero
connection. The vacuum is also even in the connection,
so its first derivative vanishes there. The analytic
disk gives an additional gain for precisely these two
terms; every other term is controlled directly by its
Grassmann degree.

\begin{proposition}\label{prop:rg-activity-contraction}
Assume \eqref{eq:rg-adjustment-charge-condition},
\eqref{eq:rg-contraction-parameters}, and
\begin{equation}\label{eq:rg-contraction-polymer-parameters}
 \kappa\ge7\kappa_0+7,\qquad
             L(\kappa-6\kappa_0-4)-1\ge\kappa.
\end{equation}
For $\mathscr E\in K_k^{x,\C}(\ell)$ set
$E=(R\mathscr E)^{\mathrm{nat}}$.
The linear map $\mathscr E\mapsto\mathcal L E$ takes
$K_k^{x,\C}(\ell)$ into $K_{k+1}^{x,\C}(\ell/L)$,
preserves the real subspaces, and satisfies
\begin{equation}\label{eq:rg-contraction-norm}
 \|\mathcal L E\|_{k+1}^x
       \le c_{\mathrm{fl}}L^{-1/4+3\varepsilon/2}\|\mathscr E\|_k^x
       \le c_{\mathrm{fl}}L^{-1/4+2\varepsilon}\|\mathscr E\|_k^x,
       \qquad c_{\mathrm{fl}}=61c_{\mathrm{rb}}.
\end{equation}
Relative to the activity windows,
\begin{equation}\label{eq:rg-contraction-window}
 \frac{\|\mathcal L E\|_{k+1}^x}{w_{k+1}^{\mathrm E}}
       \le c_{\mathrm{fl}}L^{-3/8+11\varepsilon/2}
                       \frac{\|\mathscr E\|_k^x}{w_k^{\mathrm E}}.
\end{equation}
Restriction to the zero external argument commutes with
the map.
\end{proposition}

\begin{proof}
Write
\[
 N(X)=\sup_{\widetilde R_k\times\mathfrak X_k}
                                      \|E(X;\cdot,\cdot)\|_{h_k}.
\]
Extraction and adjustment give
\begin{equation}\label{eq:rg-adjusted-polymer-bound}
 N(X)\le
 \begin{cases}
 20\|\mathscr E\|_k^x\exp(-\kappa d_M(X)),&X\text{ small},\\
 \|\mathscr E\|_k^x\exp(-\kappa d_M(X)),&X\text{ large}.
 \end{cases}
\end{equation}
For a small polymer this uses
$\|R\mathscr E(X;\mathscr A,\mathscr C)\|_{h_k}
\le\|\mathscr E(X;\mathscr A,\mathscr C)\|_{h_k}
 +\|\mathscr E(X;0,0)\|_{h_k}$ and
\eqref{eq:rg-adjustment-norm}. For a large polymer both
operations leave the activity unchanged.

Fix a small $X$ and a coarse connection pair, and use
the global representatives above. Put $\rho=L^{7/4}$.
The inverse image of the open set $R_k\times\mathfrak X_k$
under $t\mapsto(t\mathscr A^{\mathrm b},t\mathscr C^{\mathrm b})$
contains the closed disk $|t|\le\rho$, hence a slightly
larger disk. The functions
\[
 f_{nm}(t)=E_{nm}(X;t\mathscr A^{\mathrm b},
                              t\mathscr C^{\mathrm b})
\]
are holomorphic there, with values in a finite-dimensional
kernel space. Their weighted norms obey
\[
 \frac{h_{k,1}^nh_{k,2}^m}{n!m!}\|f_{nm}(t)\|_1
                                      \le N(X)\qquad(|t|\le\rho).
\]
Here $\|\cdot\|_1$ is the sum of kernel moduli with their
index weights.

For the vacuum, \eqref{eq:rg-adjusted-kernel-normalization}
gives $f_{00}(0)=0$. Charge conjugation fixes a scalar
kernel and reverses both connections, so
$f_{00}(-t)=f_{00}(t)$ and $f_{00}'(0)=0$.
Cauchy's formula gives
\[
 f_{00}(1)=\frac1{2\pi i}\oint_{|t|=\rho}
                                   \frac{f_{00}(t)}{t^2(t-1)}\,\dd t.
\]
The vacuum equality in
Proposition~\ref{prop:rg-small-polymer-kernel-comparison}
and $\rho\ge2$ therefore give
\begin{equation}\label{eq:rg-vacuum-cauchy-gain}
 \|(S_LE)^{[0,0]}(X;\mathscr A_L,\mathscr C_L)\|_{h_{k+1}}
              \le\frac{N(X)}{\rho(\rho-1)}
              \le2L^{-7/2}N(X).
\end{equation}

For the ordinary quadratic kernel, normalization gives
$f_{20}(0)=0$ as a full antisymmetric kernel. Its balanced
matrix components vanish by
\eqref{eq:rg-adjusted-kernel-normalization}; their
oppositely ordered components vanish by antisymmetry,
and the two equally barred or equally unbarred components
vanish by balance. Hence
\[
 f_{20}(1)=\frac1{2\pi i}\oint_{|t|=\rho}
                                    \frac{f_{20}(t)}{t(t-1)}\,\dd t,
 \qquad
 \|f_{20}(1)\|_1\le\frac{\sup_{|t|=\rho}\|f_{20}(t)\|_1}{\rho-1}.
\]
The comparison of kernels costs $\exp(4)$ in this
bidegree. Its scaling exponent is $p(2,0)=9/4$.
Multiplying the last inequality by $h_{k,1}^2/2$ and
using \eqref{eq:rg-degree-scaling} gives
\begin{equation}\label{eq:rg-quadratic-cauchy-gain}
 \|(S_LE)^{[2,0]}(X;\mathscr A_L,\mathscr C_L)\|_{h_{k+1}}
                       \le2\exp(4)L^{-4}N(X).
\end{equation}

For all remaining bidegrees use the original scaled
connection pair, which belongs to the fine domain.
Evenness leaves only $(1,1)$, $(0,2)$, and total degrees
at least $4$. The exponents satisfy
\[
 p(1,1)>\frac{13}{4}-\frac{3\varepsilon}{2},\qquad
 p(0,2)>\frac{13}{4}-\frac{3\varepsilon}{2},\qquad
 p(n,m)\ge n+m.
\]
Consequently
\[
 \begin{split}
 \sum_{(n,m)\ne(0,0),(2,0)}
       \|(S_LE)^{[n,m]}(X;\mathscr A_L,\mathscr C_L)\|_{h_{k+1}}
 &\le L^{-13/4+3\varepsilon/2}
       \sum_{n,m}\|E^{[n,m]}(X;\mathscr A_L,\mathscr C_L)\|_{h_k}\\
 &\le L^{-13/4+3\varepsilon/2}N(X).
 \end{split}
\]
The sum is taken at one connection pair before bounding
it by $N(X)$. The large-$L$ condition ensures
$2L^{-7/2}\le L^{-13/4}$ and
$2\exp(4)L^{-4}\le L^{-13/4}$. Adding the three
degree contributions yields
\begin{equation}\label{eq:rg-small-polymer-contraction}
 \|S_LE(X;\mathscr A_L,\mathscr C_L)\|_{h_{k+1}}
                       \le3L^{-13/4+3\varepsilon/2}N(X).
\end{equation}

We now sum over the fine polymers in a coarse polymer
$Y$. The hypotheses \eqref{eq:rg-contraction-polymer-parameters}
imply
$\kappa-1\ge\kappa_0$ and
$L(\kappa-\kappa_0-2)\ge\kappa$.
The reblocking bound
\eqref{eq:polymer-reblock-sum} at rate $\kappa$
and \eqref{eq:rg-adjusted-polymer-bound} therefore give
\[
 \begin{split}
 \sum_{\substack{X\text{ small}\\\bar X^L=LY}}
       \|S_LE(X;\mathscr A_L,\mathscr C_L)\|_{h_{k+1}}
 &\le60c_{\mathrm{rb}}L^{-1/4+3\varepsilon/2}
       \|\mathscr E\|_k^x\exp(-\kappa d_M(Y)).
 \end{split}
\]
For large $X$, scaling decreases the norm by
\eqref{eq:rg-scaling-norm-identity} and the ratios of
the two weights. Since $d_M(X)>L$,
\[
 \exp(-\kappa d_M(X))
       \le\exp(-L)\exp(-(\kappa-1)d_M(X)).
\]
Use \eqref{eq:polymer-reblock-sum} at rate $\kappa-1$
to obtain
\[
 \sum_{\substack{X\text{ large}\\\bar X^L=LY}}
       \|S_LE(X;\mathscr A_L,\mathscr C_L)\|_{h_{k+1}}
 \le c_{\mathrm{rb}}L^3\exp(-L)
             \|\mathscr E\|_k^x\exp(-\kappa d_M(Y)).
\]
As $L^3\exp(-L)\le L^{-1/4}$, these estimates prove
\eqref{eq:rg-contraction-norm} after taking the two
suprema in the activity norm. Division by
$w_{k+1}^{\mathrm E}=L^{1/8-7\varepsilon/2}w_k^{\mathrm E}$
gives the relative factor $c_{\mathrm{fl}}L^{-3/8+5\varepsilon}$,
which is at most the factor in \eqref{eq:rg-contraction-window}.
The additional parameter choice
$c_{\mathrm{fl}}L^{-3/8+6\bar\varepsilon}\le1/16$
makes its relative factor at most $1/16$.

The map is linear. Its kernels are finite sums of
holomorphic fine kernels composed with the complex
linear connection scaling. The locality was proved
after \eqref{eq:rg-reblocking-definition}.
For lattice covariance a coarse group element
$g'=(S,a')$ has the fine representative $g=(S,La')$.
It gives a bijection of the polymers in the two
reblocking sums, satisfies
$(g'\mathscr A)_L=g\mathscr A_L$, and intertwines the
spin actions by $a_{g'}S_L=S_La_g$.
The same identity holds for $\mathscr C$.
Thus covariance of each fine activity gives covariance
of its sum. For the real gauges, every fine bond within
distance $1$ of $X$ scales to a bond within distance
$1/L\le1$ of $Y$, while
$G_\lambda S_L=S_LG_{\lambda_L}$ follows from
$e_k\lambda_L(c)=e_{k+1}\lambda(c/L)$.
Independent real gauges of the two arguments therefore
pass through each term. Finally, scaling has real
coefficients, preserves the bar index, and commutes with
charge conjugation and the antilinear reality map.
These facts prove membership in both stated activity
spaces. Since extraction and adjustment commute with
restriction and $0_L=0$, reblocking also commutes with
restriction to the zero external argument.
\end{proof}

The reblocked sum retains the evaluated normalized
activity. Indeed, each fine polymer belongs to exactly
one reblocking class, so
\begin{equation}\label{eq:rg-reblocked-sum}
 \sum_Y(\mathcal LE)(Y;\mathscr A,\mathscr C)
      =S_L\sum_X E(X;\mathscr A_L,\mathscr C_L).
\end{equation}
If the coarse generators are evaluated at fields
with $\chi'=\delta_{\alpha,\mathscr A+\mathscr C}\psi'$
and its barred counterpart, their scaled fields satisfy
the fine relations. This follows from
$d_X(x/L,x'/L)=d_X(x,x')/L$ and
$U_{(\mathscr A+\mathscr C)_L}(x,x')
 =U_{\mathscr A+\mathscr C}(x/L,x'/L)$:
the path bijection and
$e_k\eta_kL^{-1/2}=e_{k+1}\eta_{k+1}$ give the latter.
Substitution in the difference quotient produces
$L^{-1-\alpha}\chi'(x/L,x'/L)$, and the barred quotient
uses the inverse transport. Adjustment consequently
has its original value under this evaluation. Thus the
evaluation of \eqref{eq:rg-reblocked-sum} is also the
evaluation of
\[
 S_L\sum_X(R\mathscr E)(X;\mathscr A_L,\mathscr C_L).
\]

\subsection{Lifted formulas and comparison of local data}
\label{subsec:rg-local-pullback}

For the later comparison of different torus volumes we
need formulas that depend on a lift and local data alone.
The required margin also distinguishes the pairs produced
by scaling from the full set of coarse pairs.

For a finite photon set $\widetilde U$ on the covering
lattice, let $\mathfrak G_k(\widetilde U)$ be the algebra
generated by the ordinary indices over $\widetilde U$
and pair indices whose two photon sites lie in
$\widetilde U$. Its ordinary and pair measures remain
$\eta_k^3/48$ and $(\eta_k^3/48)^2$. For a system of
separately antisymmetric kernels $K=(K_{nm})$ define
\[
 \operatorname{El}_k(K)
   =\sum_{n,m}\frac1{n!m!}\int\dd\xi\,\dd\zeta\,
                 K_{nm}(\xi;\zeta)\psi(\xi_1)\cdots\psi(\xi_n)
                                    \chi(\zeta_1)\cdots\chi(\zeta_m).
\]
The integrals range over these two finite index sets.
The kernel expansion is a bijection between the systems
of such kernels and this algebra: at increasing arguments
its coefficient is the kernel times the positive product
of index weights, as in \eqref{eq:rg-activity-kernels}.

Map each covering generator to the torus generator at
its simultaneous period class, with the antiperiodic
convention, and denote the homomorphism by
$\operatorname{cov}_\ell$. If
\begin{equation}\label{eq:rg-generator-pullback-margin}
 \operatorname{diam}_\infty(\widetilde U)+\eta_k<\ell,
\end{equation}
it is injective and preserves kernels when read on this
lift. In fact two congruent photon sites in
$\widetilde U$ must coincide. Choose a fundamental
domain containing $\widetilde U$ and order its
independent ordinary and pair indices so that they
extend the covering orders. The homomorphism maps
distinct covering basis monomials to distinct torus
basis monomials with the same coefficients and weights.
This proves injectivity, preservation of bidegrees,
and equality of the kernels at the lifted arguments.
All torus kernels at indices outside these period
classes vanish.

Now let $\widetilde X$ be a face-connected lift and put
$\widetilde U=\widetilde X^{+\eta_k}$, its mesh
neighborhood. The margin
\begin{equation}\label{eq:rg-activity-pullback-margin}
 \operatorname{diam}_\infty(\widetilde X)+2\eta_k+M<\ell
\end{equation}
gives a bijection of ordinary indices over
$\widetilde X$ onto those over $X$, and of admissible
pairs for $\widetilde X$ onto those for $X$.
For a pair of the first kind, its endpoints and every
point of its shortest paths lift inside $\widetilde X$:
a lifted path point within distance $M$ of its first
endpoint differs from its representative in
$\widetilde X$ by a period of norm at most
$M+\operatorname{diam}_\infty(\widetilde X)<\ell$.
That period is zero. For a boundary pair
$x'=x^{(3)}$, lift the endpoint in $X$; the other is
its prescribed nearest neighbor in $\widetilde U$.
Two lifted pairs in the same class coincide because
their first photon sites coincide.
These arguments prove surjectivity, injectivity and
preservation of pair kind.
The margin also implies
\eqref{eq:rg-generator-pullback-margin} for $\widetilde U$.
By the nonwrapping lifting result, $X$ is small exactly
when $d_M(\widetilde X)\le L$: for either a small torus
polymer or a small covering polymer the unique
nonwrapping lift has the same net length.

Take data $(K^0,K,\widetilde B)$ consisting of two
kernel systems over $\widetilde U$ and a connection
on bonds with both ends in $\widetilde X$.
For $d_M(\widetilde X)\le L$, define
\begin{equation}\label{eq:rg-lifted-extraction}
 \begin{split}
 M^0&=\operatorname{Vol}(\widetilde X)^{-1}
      \int_{\widetilde X}\dd x'\int_{\widetilde X}\dd x''\,
                                                    K^0_{20}(x',x''),\\
 \Phi^R&=\operatorname{El}_k(K)-K^0_{00}
               -\int_{\widetilde X}\dd x\,\bar\psi(x)^TM^0\psi(x),\\
 \Phi^E&=\Phi^R-(\Phi^R)^{[2,0]}_{\widetilde X}
               +\Theta_{\widetilde X,\widetilde B}
                                      (\Phi^R)^{[2,0]}_{\widetilde X}.
 \end{split}
\end{equation}
Here $K^0_{20}$ is the balanced matrix component in the
earlier convention, and the subscript $\widetilde X$
retains only quadratic terms whose two sites lie over
$\widetilde X$. Thus the transport homomorphism is
applied within its domain. For $d_M(\widetilde X)>L$,
set $\Phi^R=\Phi^E=\operatorname{El}_k(K)$. In both
cases define
\[
 \Phi^S=\widetilde S_L\Phi^E,
\]
where $\widetilde S_L$ is the covering version of
\eqref{eq:rg-leg-scaling}. These formulas use no side
$\ell$. Translation by $v\in M\Z^3$ preserves the
cube geometry, measures and chain transports. Therefore
\begin{equation}\label{eq:rg-lifted-translation}
 \Phi^R(T_v\mathrm{data})=T_v\Phi^R(\mathrm{data}),\qquad
 \Phi^E(T_v\mathrm{data})=T_v\Phi^E(\mathrm{data}),\qquad
 \Phi^S(T_v\mathrm{data})=T_{v/L}\Phi^S(\mathrm{data}),
\end{equation}
with the translated polymer understood on the left.

For an activity $\mathscr E$, let $K^0$ and $K$ be its
kernels at $(0,0)$ and at $(\mathscr A_L,\mathscr C_L)$,
restricted to the lift, and let $\widetilde B$ be the
restriction of $(\mathscr A+\mathscr C)_L$ to the
internal bonds. Then
\begin{equation}\label{eq:rg-lifted-activity-identities}
 \begin{split}
 (R\mathscr E)(X;\mathscr A_L,\mathscr C_L)
                              &=\operatorname{cov}_\ell\Phi^R,\\
 (R\mathscr E)^{\mathrm{nat}}(X;\mathscr A_L,\mathscr C_L)
                              &=\operatorname{cov}_\ell\Phi^E,\\
 S_L[(R\mathscr E)^{\mathrm{nat}}(X;\mathscr A_L,\mathscr C_L)]
                              &=\operatorname{cov}_{\ell/L}\Phi^S.
 \end{split}
\end{equation}
To prove the first identity, locality and the index
bijection reproduce the original activity from $K$.
The vacuum subtraction is exactly $K^0_{00}$, and
the two label-averaged mass integrals have the same
weights and volume in the lift and on the torus.
The diagonal bilinear is period invariant, since its
two antiperiodic signs cancel. The subtraction is
therefore also the same.
For the second identity the ordinary quadratic terms
of $\Phi^R$ lie over $\widetilde X$. Cubes, their
shortest chains and their transports are the lifted
versions of the torus construction. On each ordinary
generator this gives
$\operatorname{cov}_\ell\Theta_{\widetilde X,\widetilde B}
=\Theta_{X,(\mathscr A+\mathscr C)_L}\operatorname{cov}_\ell$.
The homomorphism identity proves the adjusted formula.
Finally, on each generator
\[
 \operatorname{cov}_{\ell/L}\widetilde S_L
                                  =S_L\operatorname{cov}_\ell.
\]
Both sides have the same scaled index and coefficient,
with compatible antiperiodic signs. This proves the
third identity.

The margin \eqref{eq:rg-activity-pullback-margin}
suffices for the image of $\widetilde S_L$, whose
pairs have $d_X\le M/L$. To identify
$\operatorname{cov}_{\ell/L}$ with the quotient map
on the full coarse local algebra, including pairs
with $M/L<d_X\le M$, we use the stronger margin
\begin{equation}\label{eq:rg-full-coarse-pullback-margin}
 \operatorname{diam}_\infty(\widetilde U)+LM<\ell.
\end{equation}
It is equivalent to
$\operatorname{diam}_\infty(\widetilde U/L)+M<\ell/L$.
For a small polymer both margins hold, since
\[
 \operatorname{diam}_\infty(\widetilde X)+2\eta_k+LM
                    \le2LM+4M<\ell.
\]
The torus expressions in
\eqref{eq:rg-lifted-activity-identities} do not depend
on the chosen lift. Translation of that lift by $\ell w$
translates the data; their total even degree makes the
product of their argument signs one.
Equation \eqref{eq:rg-lifted-translation} and the same
evenness then cancel the generator signs under the
quotient maps at both levels.
Thus identical lifted data produce identical lifted
extraction, adjustment and scaled kernels on different
admissible tori, subject to the stated margins.


\section{The exact block density}\label{sec:rg-step}

The block transformation has three factors: a new quadratic
fermion density, a determinant quotient, and a finite integral
over photon and fermion fluctuations. We derive this identity
before expanding its last factor. The integration at level
zero includes every chamber component and uses the chamber
Berezin measure.

\subsection{Integration spaces and real backgrounds}
\label{subsec:rg-step-real-background}

Fix $k\ge0$, $\eta_k=L^{-k}$, odd $L\ge5$, and
$M=L^m\ge L^2$, with $\ell/(LM)$ odd and at least $13$.
The regulator parameter satisfies $0<\rho<1$.
Keep $b=1$ as a letter, and write
$b_L=b/L$ and
$b_k=b(1-L^{-1})/(1-L^{-k})$ for $k\ge1$.
Write $\mathscr V_k=\Sspace(1,\ell)$ for $k\ge1$, with
four spin components and counting measure, and
$\mathscr V_0=\Sspace^{\mathrm{ch}}(1,\ell)$, with
$192$ components per photon site and measure $1/48$.
The one-step averaging operator into $\Sspace(L,\ell)$ is
\[
 \mathcal A_k(B)=
 \begin{cases}
 Q(B),&k\ge1,\\
 Q(B)J^*,&k=0,
 \end{cases}
 \qquad
 \mathcal A_k^\sharp(B)=\mathcal A_k(-B)^T.
\]
The target measure is $L^3$. On these spaces set
\begin{equation}\label{eq:rg-step-quadratic-operators}
 \mathcal D_k(B)=
 \begin{cases}D_k(B),&k\ge1,\\D_B+\mu,&k=0,\end{cases}
 \qquad
 K_k(B)=\mathcal D_k(B)+b_L\mathcal A_k^\sharp(B)\mathcal A_k(B).
\end{equation}
Thus $K_k=D_k+b_LP_L$ for $k\ge1$ and
$K_0=F_0=D_B+\mu+b_LJP_LJ^*$.
We write $D_0(B)=D_B+\mu$, so that
$\mathcal D_k(B)=D_k(B)$ at every level.
When it exists, $\Gamma_k=K_k^{-1}$, and
$\delta Z_k(B)=\det K_k(B)$.
The transpose refers to the displayed bilinear measures.
For a real $B$, the phases have modulus one and
$\mathcal A_k^\sharp=\mathcal A_k^*$.

The normalized averages give
$Q(B)Q^T(-B)=1$. On a coarse block $B(z')$ the product
is
$L^{-3}\sum_{z\in B(z')}U_B(z',z)U_B(z',z)^{-1}=1$.
At level zero the additional cluster mean contributes
$48/48=1$. Similarly
$\mathcal Q_k\mathcal Q_k^\sharp=1$ for $k\ge1$.
In a unit block the photon count gives
$\eta_k^3\eta_k^{-3}=1$ and the cluster count again
gives $48/48=1$.
For real $B$ the inverse phases in the embeddings are
the complex conjugates of the averaging phases.
The products $\Pi_1=\mathcal Q_k^\sharp\mathcal Q_k$
and $\Pi_L=\mathcal Q_k^\sharp P_L\mathcal Q_k$ are
therefore orthogonal projections; so are $P_L$ and,
at level zero, $JP_LJ^*$. Their scalar spin kernels
commute with every constant spin matrix.

\begin{proposition}\label{prop:rg-step-real-inverses}
Let $B$ be real and $\mu>0$. The operators
$D_B+\mu$, $\mathcal D_k(B)$ and $K_k(B)$ are
invertible at every level at which they are defined.
For $k\ge1$,
\begin{equation}\label{eq:rg-step-real-coarse-inverse}
 D_k(B)^{-1}=b_k^{-1}
           +\mathcal Q_k(B)(D_B+\mu)^{-1}\mathcal Q_k^\sharp(B),
 \qquad \operatorname{Re}D_k(B)^{-1}\ge b_k^{-1}.
\end{equation}
Moreover $\operatorname{Re}D_k>0$ and
$\operatorname{Re}K_k>0$ as quadratic forms for
$k\ge1$, whereas $\operatorname{Re}K_0\ge\mu$.
Each $K_k$ satisfies
$K_k^*=\gamma_5K_k\gamma_5$.
For every real $y$, $K_k+i\gamma_5y$ is invertible,
and for $y\ne0$,
\begin{equation}\label{eq:rg-step-real-resolvent}
 \|(K_k+i\gamma_5y)^{-1}\|_{2\to2}\le |y|^{-1}.
\end{equation}
For $k\ge1$ the fine inverse $S_{k,y}$ also
exists for every real $y$.
In particular $\det D_k(0)$ and $\delta Z_k(0)$
are nonzero, including the determinants on
$\mathscr V_0$.
\end{proposition}

\begin{proof}
At a real connection the chamber hops are self-adjoint
unitary involutions. In the proof of
\eqref{eq:dirac-realpart}, reindexing $x$ by $x^{(i)}$
uses
$U_B(x^{(i)},x)=\overline{U_B(x,x^{(i)})}$.
It gives
\[
 (f,\mathcal L_Bf)=\frac12\sum_iw_i
                  \int\dd x\,|R_i(B)f(x)-f(x)|_2^2\ge0.
\]
The direction vector changes sign under the same
involution; hence
$\sum_iw_i(\gamma\cdot\widehat d_i)R_i(B)$
is anti-self-adjoint. The cancellation
$\sum_iw_i\widehat d_i=0$ yields
$D_B^*=\gamma_5D_B\gamma_5$ and
$\operatorname{Re}D_B=(\rho/(2\eta_k))\mathcal L_B\ge0$.
Thus $\operatorname{Re}(D_B+\mu)\ge\mu$.
For $k\ge1$ adding the positive penalty
$b_k\Pi_1$ gives $\operatorname{Re}F_\infty\ge\mu$.
These inequalities imply injectivity and then
invertibility in the finite-dimensional spaces.

Put $G=D_B+\mu$ and $S_k=(G+b_k\mathcal Q_k^\sharp
\mathcal Q_k)^{-1}$. Multiplication of
$b_k^{-1}+\mathcal Q_kG^{-1}\mathcal Q_k^\sharp$
by $D_k=b_k-b_k^2\mathcal Q_kS_k\mathcal Q_k^\sharp$
gives
\[
 1+b_k\mathcal Q_kG^{-1}
      [1-(G+b_k\mathcal Q_k^\sharp\mathcal Q_k)S_k]
                                               \mathcal Q_k^\sharp=1.
\]
This proves \eqref{eq:rg-step-real-coarse-inverse}.
For an invertible operator $O$,
$\operatorname{Re}(u,Ou)=
\operatorname{Re}(Ou,O^{-1}Ou)$.
In particular $\operatorname{Re}O\ge0$ implies
$\operatorname{Re}O^{-1}\ge0$.
As $\mathcal Q_k^\sharp=\mathcal Q_k^*$ here,
the second term of $D_k^{-1}$ has nonnegative real
part. For a nonzero $v$, consequently,
\[
 \operatorname{Re}(v,D_kv)
       =\operatorname{Re}(D_kv,D_k^{-1}D_kv)
                           \ge b_k^{-1}\|D_kv\|_2^2>0.
\]
Adding $b_LP_L$ preserves strict positivity.
At $k=0$ the same conclusion follows directly by
adding $b_LJP_LJ^*$ to $D_B+\mu$.

The projections commute with $\gamma_5$.
It follows first that
$F_\infty^*=\gamma_5F_\infty\gamma_5$ and
$S_k^*=\gamma_5S_k\gamma_5$.
Taking the adjoint in the formula for $D_k$ and using
$\mathcal Q_k^\sharp=\mathcal Q_k^*$ gives
$D_k^*=\gamma_5D_k\gamma_5$. The penalty has the
same symmetry, proving it for $K_k$.
For $k=0$ this is already the symmetry of its three
terms. Hence
\[
 (K_k+i\gamma_5y)^*(K_k+i\gamma_5y)
                                      =K_k^*K_k+y^2,
 \qquad
 \|(K_k+i\gamma_5y)f\|_2^2
                         =\|K_kf\|_2^2+y^2\|f\|_2^2.
\]
For $y\ne0$ this proves injectivity and the inverse
bound; $y=0$ is covered by strict positivity.
For the fine inverse at $k\ge1$, use the block
matrix with diagonal blocks
$b_k+b_LP_L+i\gamma_5y$ and $F_\infty$,
and off-diagonal blocks
$-b_k\mathcal Q_k$ and $-b_k\mathcal Q_k^\sharp$.
The first diagonal block is invertible by its
idempotent decomposition
\eqref{eq:fine-resolvent-perturbation}.
Its two Schur complements are $F_y$ and
$K_k+i\gamma_5y$. The latter is invertible,
so the block matrix and then $F_y$ are invertible.
This proves existence of $S_{k,y}$.
The asserted normalizing determinants are therefore
nonzero at the real zero connection.
\end{proof}

\subsection{The next quadratic density}
\label{subsec:rg-step-scale-identification}

Assume now that $\ell/L$ is also admissible, and let
$B=B'_L$. At the outgoing level the mass is $L\mu$ and
the charge is $e_{k+1}=L^{1/2}e_k$.
The following identities hold whenever $S_k(B)$ and
$\Gamma_k(B)$ exist for $k\ge1$, or $\Gamma_0(B)$
exists for $k=0$. This invertibility domain is open
and invariant under $B\mapsto-B$: all elementary
entries are entire functions of the bond values,
and the transpose identities are
\begin{equation}\label{eq:rg-step-transpose}
 \mathcal D_k(B)^T=\mathsf C^{-1}\mathcal D_k(-B)\mathsf C,\qquad
 K_k(B)^T=\mathsf C^{-1}K_k(-B)\mathsf C,\qquad
 \Gamma_k(B)^T=\mathsf C^{-1}\Gamma_k(-B)\mathsf C.
\end{equation}
For the first identity reindex each hopping term by
$x^{(i)}$, using
$M_i(x^{(i)})^T=\mathsf C^{-1}M_i(x)\mathsf C$
and the reverse phase at $-B$. Transposition of the
averages reverses their connection in the product
$P_L$, and the scalar kernels commute with
$\mathsf C$. The remaining identities follow by
products and inversion. Invertible matrices have
analytic inverses on their open determinant domains.

Define the fine-to-$L$ average and its embedding by
\[
 \mathcal Q_k^+(B)=
 \begin{cases}Q(B)\mathcal Q_k(B),&k\ge1,\\Q(B)J^*,&k=0,\end{cases}
 \qquad
 \mathcal Q_k^{+,\sharp}(B)=\mathcal Q_k^+(-B)^T.
\]
They obey
$\mathcal Q_k^+\mathcal Q_k^{+,\sharp}=1$.
Set $\bar b_k=b_{k+1}/L$ and
\begin{equation}\label{eq:rg-step-fine-plus}
 F_k^+(B)=D_B+\mu+\bar b_k\mathcal Q_k^{+,\sharp}\mathcal Q_k^+,
 \qquad \Sigma_k^+(B)=F_k^+(B)^{-1}.
\end{equation}
For $k=0$ this is $F_0$ and its inverse $\Gamma_0$.
For $k\ge1$ it is the fine resolvent operator at
$y=0$, since
$\bar b_k=b_kb_L/(b_k+b_L)$.

To verify existence and the identities needed below,
write $\mathcal B_k=(b_k+b_LP_L)^{-1}$ for $k\ge1$
and consider
\[
 \begin{pmatrix}
 b_k+b_LP_L&-b_k\mathcal Q_k\\
 -b_k\mathcal Q_k^\sharp&F_\infty
 \end{pmatrix}.
\]
Its lower-right Schur complement is $K_k$.
Since $P_L$ is idempotent and scalar in spin,
$\mathcal B_k=b_k^{-1}(1-P_L)+(b_k+b_L)^{-1}P_L$.
The other Schur complement is
\[
 F_\infty-b_k^2\mathcal Q_k^\sharp\mathcal B_k\mathcal Q_k
    =D_B+\mu+\frac{b_kb_L}{b_k+b_L}
                                  \mathcal Q_k^\sharp P_L\mathcal Q_k
    =F_k^+.
\]
The two assumed inverses make the block matrix
invertible, and hence $F_k^+$ is invertible.
Equating the upper-left and lower-left blocks of
its inverse gives
\begin{equation}\label{eq:rg-step-block-inverse-identities}
 \Gamma_k=\mathcal B_k+
       b_k^2\mathcal B_k\mathcal Q_k\Sigma_k^+
                                     \mathcal Q_k^\sharp\mathcal B_k,
 \qquad
 H_k\Gamma_k=b_k\Sigma_k^+\mathcal Q_k^\sharp\mathcal B_k.
\end{equation}
This is the same Schur calculation as
Proposition~\ref{prop:fine-coarse-resolvents}, retaining
the off-diagonal block for the mean field.

Let
\begin{equation}\label{eq:rg-step-unscaled-coarse}
 D_{k+1}^0(B)=b_L-
                       b_L^2\mathcal A_k(B)\Gamma_k(B)\mathcal A_k^\sharp(B).
\end{equation}
For $k\ge1$, $QP_L=Q$ and $P_LQ^T(-B)=Q^T(-B)$
give
$Q\mathcal B_k=(b_k+b_L)^{-1}Q$ and
$\mathcal B_kQ^T(-B)=(b_k+b_L)^{-1}Q^T(-B)$.
Insert these in \eqref{eq:rg-step-block-inverse-identities}
to obtain
\begin{equation}\label{eq:rg-step-compressed-identities}
 \begin{split}
 D_{k+1}^0(B)
       &=\bar b_k-\bar b_k^2
                            \mathcal Q_k^+\Sigma_k^+\mathcal Q_k^{+,\sharp},\\
 b_LH_k(B)\Gamma_k(B)\mathcal A_k^\sharp(B)
       &=\bar b_k\Sigma_k^+\mathcal Q_k^{+,\sharp}.
 \end{split}
\end{equation}
For $k=0$ both formulas follow directly from
$\bar b_0=b_L$, $H_0=1$,
$\Sigma_0^+=\Gamma_0$, and
$\mathcal Q_0^+=\mathcal A_0$.

Let $\mathsf R_L$ denote the scaling
$f_L(c,g)=L^{-1}f(c/L,g)$ on chamber spinors and
$v_L(z')=L^{-1}v(z'/L)$ on unit spinors.
These are bijections between the corresponding
incoming and outgoing spaces. Scaling preserves
the label and the nested block centers, and
$e_k\eta_kL^{-1/2}=e_{k+1}\eta_{k+1}$ preserves
each transport phase. Consequently
\[
 (D_B+\mu)\mathsf R_L
       =L^{-1}\mathsf R_L(D_{B'}+L\mu),\qquad
 \mathcal Q_k^+(B)\mathsf R_L
       =\mathsf R_L\mathcal Q_{k+1}(B'),
\]
and the corresponding embedding identity holds.
For the average, its $L^{-3}$ factor cancels the
$L^3$ change of the chamber index measure.
Thus
\begin{equation}\label{eq:rg-step-fine-scaling}
 F_k^+(B)\mathsf R_L
                   =L^{-1}\mathsf R_LF_\infty(B'),\qquad
 \Sigma_k^+(B)\mathsf R_L
                         =L\mathsf R_LS_{k+1}(B').
\end{equation}
In particular the outgoing inverse exists.
Equations \eqref{eq:rg-step-compressed-identities} and
\eqref{eq:rg-step-fine-scaling} now imply
\begin{equation}\label{eq:rg-step-coarse-scaling}
 D_{k+1}^0(B)\mathsf R_L
                       =L^{-1}\mathsf R_LD_{k+1}(B'),\qquad
 b_LH_k(B)\Gamma_k(B)\mathcal A_k^\sharp(B)\mathsf R_L
                       =\mathsf R_LH_{k+1}(B').
\end{equation}
All expressions in the second identity act on the
outgoing four-component field.

For independent outgoing Grassmann fields define
\begin{equation}\label{eq:rg-step-mean-fields}
 \Psi^c=b_L\Gamma_k(B)\mathcal A_k^\sharp(B)\Psi'_L,\qquad
 \bar\Psi^c=b_L\Gamma_k(B)^T\mathcal A_k(B)^T\bar\Psi'_L.
\end{equation}
Then, with $\psi'=H_{k+1}(B')\Psi'$ and
$\bar\psi'=\bar H_{k+1}(B')\bar\Psi'$,
\begin{equation}\label{eq:rg-step-induced-fields}
 H_k(B)\Psi^c=\psi'_L,\qquad
 \bar H_k(B)\bar\Psi^c=\bar\psi'_L.
\end{equation}
The first relation is \eqref{eq:rg-step-coarse-scaling}.
For the second use \eqref{eq:rg-step-transpose}
and apply the first relation at $-B$ to
$\mathsf C\bar\Psi'$; the constant matrix
$\mathsf C$ commutes with scaling and the averages.
The covariant H\"older quotients of these induced
fields are also the scaled outgoing quotients:
the transports agree and $d_X$ is divided by $L$,
so the quotient carries the factor $L^{-1-\alpha}$.
This holds for every pair, including a pair inside
one cluster.
Applying $T_kH_k=1$ and its barred version shows
that the means in \eqref{eq:rg-step-mean-fields}
are also $T_k(B)\psi'_L$ and $\bar T_k(B)\bar\psi'_L$.
At level zero these maps are identities on
the chamber field.

The quadratic term in the Gaussian block formula is
therefore exactly
\[
 \langle\bar\Psi'_L,D_{k+1}^0(B)\Psi'_L\rangle_L
                         =\langle\bar\Psi',D_{k+1}(B')\Psi'\rangle_1,
\]
since the pairings of scaled fields multiply by $L$.
Together with \eqref{eq:block-square-completion},
this identifies the entire outgoing quadratic density.

\begin{proposition}\label{prop:rg-step-determinant-normalization}
On the invertibility domain just specified,
\begin{equation}\label{eq:rg-step-determinant-normalization}
 \det D_{k+1}(B')=
           b^{4n_{k+1}}\frac{\det\mathcal D_k(B)}{\delta Z_k(B)},
 \qquad n_k=\ell^3,\quad n_{k+1}=(\ell/L)^3.
\end{equation}
In particular, with $c_k=b^{-4n_{k+1}}$,
\begin{equation}\label{eq:rg-step-zero-normalization}
 \frac{c_k}{\det D_k(0)}
              =\frac{1}{\det D_{k+1}(0)\delta Z_k(0)}.
\end{equation}
\end{proposition}

\begin{proof}
Apply Sylvester's determinant identity to
$b_L\mathcal A_k\Gamma_k$ and $\mathcal A_k^\sharp$
on the two different spaces. It gives
\[
 b_L^{-4n_{k+1}}\det D_{k+1}^0(B)
       =\det(1-b_L\mathcal A_k^\sharp\mathcal A_k\Gamma_k)
       =\det(\mathcal D_k\Gamma_k)
       =\frac{\det\mathcal D_k(B)}{\delta Z_k(B)}.
\]
The incoming determinant at $k=0$ has size $192n_0$;
the exponent of $b_L$ is the outgoing dimension
$4n_1$, as in every other step.
Equation \eqref{eq:rg-step-coarse-scaling} gives
$\det D_{k+1}^0(B)=L^{-4n_{k+1}}\det D_{k+1}(B')$.
Using $Lb_L=b$ proves
\eqref{eq:rg-step-determinant-normalization}.
At the zero connection all denominators are nonzero
by Proposition~\ref{prop:rg-step-real-inverses}, so
\eqref{eq:rg-step-zero-normalization} follows.
\end{proof}

\subsection{The density with an external connection}
\label{subsec:rg-step-density}

For this subsection assume $0<\mu\le1/L$ and
$e_{k+1}=L^{1/2}e_k\le e_0$, with the hypotheses
of Propositions~\ref{prop:walk-propagator-bounds},
\ref{prop:fine-coarse-resolvents} and
\ref{prop:photon-domain-preservation}. We also impose
\eqref{eq:photon-domain-small-charge} and
\eqref{eq:rg-domain-charge-conditions}.
The inverse domain $\widetilde G_k$ is the domain
$\widetilde{\mathcal G}_k$ of the propagator estimates.
Let $A'\in\mathrm{Dom}_{k+1}$ be a real block field,
with arbitrary harmonic part, and let
$\mathscr C'\in\mathfrak X_{k+1}$.
Put
\[
 \mathscr A'=H_{k+1}^{\mathrm{ph}}A',\qquad
 \mathscr A=\mathscr A'_L,\qquad
 \mathscr C=\mathscr C'_L,\qquad
 \mathscr T'=\mathscr A'+\mathscr C',\qquad
 \mathscr T=\mathscr T'_L.
\]
The photon block objects are
\begin{equation}\label{eq:rg-step-photon-fields}
 \begin{split}
 A_k(\zeta)&=A_k^{\min}(A')+\mathrm{Sq}_k\zeta,\\
 \mathscr Z(\zeta)&=H_k^{\mathrm{ph}}\mathrm{Sq}_k\zeta,\qquad
 \mathscr B_k=H_k^{\mathrm{ph}}A_k^{\min}(A')
                                      =\mathscr A-d\omega_k.
 \end{split}
\end{equation}
Here $\omega_k$ is real and periodic, as in
\eqref{eq:photon-background-one-step}.
The truncated photon measure has the normalization
\begin{equation}\label{eq:rg-step-box-normalization}
 \mu_I(\mathrm{Box}_k)=\exp(-n_k\delta\varepsilon_k^0),\qquad
 \widehat\mu_k=
              \frac{\mathbf1_{\mathrm{Box}_k}}{\mu_I(\mathrm{Box}_k)}\mu_I,
\end{equation}
proved in \eqref{eq:photon-Gaussian-box-mass}.
Its denominator is positive: $e_k<1$ makes
$p_{0,k}=(-\log e_k)^{p_0}>0$, and the box has
nonempty interior in the photon coordinate space.

All connections read by the step have their required
domains. Indeed, scaling gives
$\mathscr A\in\frac14L^{-3/4}\widetilde R_k
\subset\widetilde R_k^-$ and
$\mathscr C\in\mathfrak X_k$.
Proposition~\ref{prop:photon-domain-preservation}
gives $A_k(\zeta)\in\mathrm{Dom}_k$ for every
$\zeta\in\mathrm{Box}_k$.
The real-gauge identity in
\eqref{eq:rg-step-photon-fields} yields
\[
 H_k^{\mathrm{ph}}A_k(\zeta)
                 =\mathscr A+\mathscr Z(\zeta)-d\omega_k.
\]
Since $\frac14\widetilde R_k$ is invariant under
real gauges, $\mathscr A+\mathscr Z(\zeta)$ also
lies in this domain. Thus both connection pairs
$(\mathscr A,\mathscr C)$ and
$(\mathscr A+\mathscr Z,\mathscr C)$ lie in
$\widetilde R_k\times\mathfrak X_k$.
By \eqref{eq:rg-total-domain-inclusion},
$\mathscr T$, $\mathscr T+\mathscr Z$ and
$\mathscr T+\mathscr Z-d\omega_k$, and their negatives,
belong to $\widetilde G_k$.
The fine and coarse inverse estimates consequently
apply to them. The outgoing operators at
$\mathscr T'$ then exist by
\eqref{eq:rg-step-fine-scaling}.

For a total connection $B$ let $\operatorname{ev}_B$
insert
\[
 \psi=H_k(B)\Psi,\qquad
 \bar\psi=\bar H_k(B)\bar\Psi,\qquad
 \chi=\delta_{\alpha,B}H_k(B)\Psi,\qquad
 \bar\chi=\delta_{\alpha,-B}\bar H_k(B)\bar\Psi.
\]
At $k=0$ the two field maps are identities.
For complex $m,\varepsilon^{\mathrm v}$ and
$\mathscr E\in K_k^{x,\C}$, with $\mathscr E=0$
at $k=0$, define
\begin{equation}\label{eq:rg-step-interaction-action}
 \begin{split}
 I_k^x[m,\mathscr E](A,C;\Psi)
     &=\operatorname{ev}_{A+C}
         \left(-m\langle\bar\psi,\psi\rangle
                               +\sum_X\mathscr E(X;A,C)\right),\\
 \mathrm{Act}_k^x[m,\mathscr E](A,C;\Psi)
     &=-\langle\bar\Psi,D_k(A+C)\Psi\rangle
                                  +I_k^x[m,\mathscr E](A,C;\Psi).
 \end{split}
\end{equation}
The mass pairing is always the chamber pairing
$\eta_k^3\,48^{-1}\sum_{c,g,\beta}\bar\psi_\beta(c,g)\psi_\beta(c,g)$.
The kinetic pairing uses the measure of $\mathscr V_k$.
For an arbitrary finite family of scalar functions
$\mathfrak P_k(X,C)$ on $\mathfrak X_k$, set
\begin{equation}\label{eq:rg-step-input-density}
 \begin{split}
 \mathscr F_k^x[\varepsilon^{\mathrm v},m,\mathscr E](A,C;\Psi)
   =\det D_k(0)^{-1}\exp\left\{
         \mathrm{Act}_k^x[m,\mathscr E](H_k^{\mathrm{ph}}A,C;\Psi)
                   -\varepsilon^{\mathrm v}n_k
                              +\sum_X\mathfrak P_k(X,C)\right\}.
 \end{split}
\end{equation}
The arguments here are $A\in\mathrm{Dom}_k$ and
$C\in\mathfrak X_k$. The fields $\Psi$ and $\bar\Psi$
are independent Grassmann variables.

Define the block transformation on an even density
$F$ whose coefficients are integrable along
$A_k(\zeta)$ by
\begin{equation}\label{eq:rg-step-block-transformation}
 \begin{split}
 (\mathbf T_k^xF)(A',\mathscr C';\Psi')
  &=c_k\int_{\mathrm{Box}_k}\dd\mu_I(\zeta)
                       \int\mathrm D\Psi\,e^{\mathcal Q_{\mathrm{bl}}}
                                         F(A_k(\zeta),\mathscr C;\Psi),\\
 \mathcal Q_{\mathrm{bl}}
  &=-b_L\big\langle
       e^{-ie_k\omega_k}\bar\Psi'_L
                    -\mathcal A_k(-\mathscr B_k-\mathscr C)\bar\Psi,\,
       e^{ie_k\omega_k}\Psi'_L
                    -\mathcal A_k(\mathscr B_k+\mathscr C)\Psi
                                                \big\rangle_L.
 \end{split}
\end{equation}
The gauge function in the two outgoing phases is
evaluated at the $L$-spaced sites.
The Berezin functional is the weighted functional
of \eqref{eq:weighted-Berezin}; its incoming mode
measure is $1$ for $k\ge1$ and $1/48$ for $k=0$.

Its normalization is
\begin{equation}\label{eq:rg-step-block-marginal}
 \int\mathrm D\Psi'\,(\mathbf T_k^xF)(A',\mathscr C';\Psi')
       =\int_{\mathrm{Box}_k}\dd\mu_I(\zeta)
                       \int\mathrm D\Psi\,F(A_k(\zeta),\mathscr C;\Psi).
\end{equation}
To prove this, integrate the outgoing family first.
Reciprocal multiplication of its barred and unbarred
generators removes the two phases and leaves the top
coefficient unchanged. Translation by the two averaged
incoming fields likewise leaves it unchanged: every
translated term that contains an external odd field
has lower degree in the outgoing family.
The remaining outgoing Gaussian exponent is
$-b_L\langle\bar\Psi'_L,\Psi'_L\rangle_L
=-b\langle\bar\Psi',\Psi'\rangle_1$.
Its integral is $b^{4n_{k+1}}=c_k^{-1}$ by
\eqref{eq:weighted-determinant}.
The two Berezin integrations commute, since both
extract an even top-degree coefficient; integration
of the finite-dimensional coefficient functions in
$\zeta$ also commutes with them. This proves
\eqref{eq:rg-step-block-marginal}.

For the density \eqref{eq:rg-step-input-density}
the coefficient integrability follows from the domain
statements above. The bond entries of all inverse
operators are analytic on their invertibility domains,
and the activity kernels are holomorphic. Composing
them with the affine map $\zeta\mapsto
H_k^{\mathrm{ph}}A_k(\zeta)$ gives continuous
coefficients on the compact photon box. Their even
Grassmann exponentials are also continuous there.

\begin{proposition}\label{prop:rg-step-density-identity}
Let the data and domains be as above. Put
\[
 m'=m+m(\mathscr E),\qquad \mathscr E^r=R\mathscr E,\qquad
 m^+=Lm',\qquad
 \varepsilon^{\mathrm v,+}
       =L^3\big(\varepsilon^{\mathrm v}
              +\varepsilon_{\mathrm v}(\mathscr E)+\delta\varepsilon_k^0\big).
\]
At $k=0$ the activity and its extracted terms vanish.
Define the outgoing inserted fields
$\psi'=H_{k+1}(\mathscr T')\Psi'$,
$\bar\psi'=\bar H_{k+1}(\mathscr T')\bar\Psi'$ and
their covariant H\"older quotients, and let
$\psi^{0,\#}$ denote their scaled fine fields,
including the barred ones.
Let the means $\Psi^c,\bar\Psi^c$ be those of
\eqref{eq:rg-step-mean-fields} at $B=\mathscr T$,
and put
\[
 \Psi^f=\Psi^c+\Gamma_k(\mathscr T)W,\qquad
 \bar\Psi^f=\bar\Psi^c+\bar W.
\]
With $\mathscr Z=H_k^{\mathrm{ph}}\mathrm{Sq}_k\zeta$,
set
\begin{equation}\label{eq:rg-step-fluctuation-exponent}
 \begin{split}
 V_k^x={}&
       \mathrm{Act}_k^x[m',\mathscr E^r]
                    (\mathscr A+\mathscr Z,\mathscr C;\Psi^f)
       -\mathrm{Act}_k^x[m',\mathscr E^r]
                    (\mathscr A,\mathscr C;\Psi^f)\\
       &+I_k^x[m',\mathscr E^r](\mathscr A,\mathscr C;\Psi^f)
       -I_k^x[m',\mathscr E^r](\mathscr A,\mathscr C;\Psi^c),\\
 \Xi_k^x={}&\int\dd\mu_I(W)\,\dd\widehat\mu_k(\zeta)\,e^{V_k^x}.
 \end{split}
\end{equation}
Then the exact identity is
\begin{equation}\label{eq:rg-step-density-identity}
 \begin{split}
 &(\mathbf T_k^x\mathscr F_k^x
                [\varepsilon^{\mathrm v},m,\mathscr E])
                                      (A',\mathscr C';\Psi')\\
 &\quad=\det D_{k+1}(0)^{-1}
       \exp\Big\{-\langle\bar\Psi',D_{k+1}(\mathscr T')\Psi'\rangle_1
                  -m^+\langle\bar\psi',\psi'\rangle
                  -\varepsilon^{\mathrm v,+}n_{k+1}\\
 &\hspace{48mm}
          +\sum_X\mathscr E^r(X;\mathscr A,\mathscr C;\psi^{0,\#})
          +\sum_X\mathfrak P_k(X,\mathscr C)\Big\}
                     \frac{\delta Z_k(\mathscr T)}{\delta Z_k(0)}\Xi_k^x.
 \end{split}
\end{equation}
All integrals in this formula are finite coefficient
integrals. Its right side depends on $A'$ through
$H_{k+1}^{\mathrm{ph}}A'$.
\end{proposition}

\begin{proof}
Extraction gives
\[
 I_k^x[m,\mathscr E](A,C;\Psi)
      =-\varepsilon_{\mathrm v}(\mathscr E)n_k
                                      +I_k^x[m',\mathscr E^r](A,C;\Psi).
\]
This is the image of
\eqref{eq:rg-exact-extraction-identity} under the
field evaluation homomorphism. At $k=0$ it holds
with zero activity and $m'=m$.

For a fixed $\zeta$, change the integrated variables
by $\Psi\mapsto e^{ie_k\omega_k}\Psi$ and
$\bar\Psi\mapsto e^{-ie_k\omega_k}\bar\Psi$.
This leaves the weighted Berezin functional unchanged,
since each mode's factors cancel in its top monomial,
also for the $192n_0$ modes of weight $1/48$.
The transport identity
\[
 U_{B-d\omega_k}(c,c')e^{ie_k\omega_k(c')}
                       =e^{ie_k\omega_k(c)}U_B(c,c')
\]
intertwines the Dirac operator and every average.
Products and inverses preserve this intertwining,
so it also holds for $H_k$ and $D_k$.
Its covariant difference identity uses the phase
at the first photon site, including for a pair
inside one cluster. The barred identity has the
inverse phase. Gauge covariance of the activity
and invariance of the mass pairing consequently
turn its first connection argument
$\mathscr A+\mathscr Z-d\omega_k$ into
$\mathscr A+\mathscr Z$.
The same change removes the gauge from the block
Gaussian factor. The integral in
\eqref{eq:rg-step-block-transformation} is now
\[
 \int\mathrm D\Psi\,
       \exp\left\{-b_L
          \langle\bar\Psi'_L-\mathcal A_k(-\mathscr T)\bar\Psi,\,
                  \Psi'_L-\mathcal A_k(\mathscr T)\Psi\rangle_L
            +\mathrm{Act}_k^x[m',\mathscr E^r]
                    (\mathscr A+\mathscr Z,\mathscr C;\Psi)\right\},
\]
multiplied by
$\det D_k(0)^{-1}
\exp\{-(\varepsilon^{\mathrm v}+
\varepsilon_{\mathrm v}(\mathscr E))n_k+
\sum_X\mathfrak P_k(X,\mathscr C)\}$.

Separate the background quadratic term by writing
\[
 \begin{split}
 \mathcal R(\Psi)={}&I_k^x[m',\mathscr E^r](\mathscr A,\mathscr C;\Psi)\\
  &+\mathrm{Act}_k^x[m',\mathscr E^r]
                    (\mathscr A+\mathscr Z,\mathscr C;\Psi)
   -\mathrm{Act}_k^x[m',\mathscr E^r]
                    (\mathscr A,\mathscr C;\Psi).
 \end{split}
\]
The last exponent above is then the background
Gaussian exponent plus $\mathcal R(\Psi)$.
Translation by the means \eqref{eq:rg-step-mean-fields}
and the block formula \eqref{eq:Gaussian-block-integration}
give
\[
 \delta Z_k(\mathscr T)
   e^{-\langle\bar\Psi',D_{k+1}(\mathscr T')\Psi'\rangle_1}
          \int\dd\mu_I(W)\,e^{\mathcal R(\Psi^f)}.
\]
The change of variables in this formula uses the
same mode measure for $\Psi$ and $W$.
At $k=0$ it is the identity
\[
 \int\mathrm D\Psi\,e^{-\langle\bar\Psi,F_0(\mathscr T)\Psi\rangle^{\mathrm{ch}}}
       G(\bar\Psi,\Psi)
   =\det F_0(\mathscr T)
          \int\dd\mu_I(W)\,
                    G(\bar W,\Gamma_0(\mathscr T)W),
\]
on the entire $192n_0$-mode space.
The Gaussian measure is
$e^{-\langle\bar W,W\rangle^{\mathrm{ch}}}\mathrm DW$,
with the weight $\sqrt{48}$ in its coefficient bound.

By definition of $V_k^x$,
$\mathcal R(\Psi^f)=
I_k^x[m',\mathscr E^r](\mathscr A,\mathscr C;\Psi^c)+V_k^x$.
Its first summand is independent of $W$ and $\zeta$,
so its even exponential factors out.
The induced-field identity
\eqref{eq:rg-step-induced-fields} gives
\[
 I_k^x[m',\mathscr E^r](\mathscr A,\mathscr C;\Psi^c)
      =-m'\langle\bar\psi'_L,\psi'_L\rangle
                 +\sum_X\mathscr E^r(X;\mathscr A,\mathscr C;\psi^{0,\#}).
\]
The chamber measure changes by $L^3$ and each field
by $L^{-1}$. Hence
$\langle\bar\psi'_L,\psi'_L\rangle
=L\langle\bar\psi',\psi'\rangle$, giving $m^+=Lm'$.
Normalizing the photon-box restriction in
\eqref{eq:rg-step-box-normalization} adds
$-\delta\varepsilon_k^0 n_k$ to the vacuum term.
Since $n_k=L^3n_{k+1}$, this term has exactly
the stated outgoing vacuum coefficient.
Finally \eqref{eq:rg-step-zero-normalization}
turns the remaining constant factor into
$\det D_{k+1}(0)^{-1}\delta Z_k(0)^{-1}$,
proving \eqref{eq:rg-step-density-identity}.

The fluctuation exponent is even and continuous
on the photon box, by the same coefficient argument
as for the input density. Thus its integral is defined,
and its two finite integrations commute.
Every background in the final expression is formed
from $\mathscr A'$ and $\mathscr C'$; the remaining
photon maps and measure are independent of $A'$.
The minimum and the real gauge have disappeared
from that expression, proving the last assertion.
\end{proof}

\subsection{Local substitution of the mean field}
\label{subsec:rg-step-local-substitution}

For $k\ge1$, let $Y$ be a union of partition $M$-cubes and allow independent
unit generators $\Psi,\bar\Psi$ over $Y$, together
with the ordinary and H\"older generators.
Using the two operators from
\eqref{eq:rg-reconstruction-two-pieces}, define
the substitution that fixes the latter generators and
sets, at the unit sites of $Y$,
\begin{equation}\label{eq:rg-step-local-substitution}
 \begin{split}
 \operatorname{sub}_{Y,B}\Psi
       &=T_Y^{(1)}(B)\psi+T_Y^{(2)}(B)\chi,\\
 \operatorname{sub}_{Y,B}\bar\Psi
       &=\bar T_Y^{(1)}(B)\bar\psi+\bar T_Y^{(2)}(B)\bar\chi.
 \end{split}
\end{equation}
The images are odd linear combinations of independent
generators and therefore anticommute with one another
and with the fixed generators. They determine a
unique Grassmann homomorphism.
Its new ordinary legs lie over $Y$ and its new pair
legs are boundary pairs of the second kind with first
point over $Y$. Its coefficients read $B$ on bonds
meeting $Y$, by the row support proved in
Subsection~\ref{subsec:rg-field-reconstruction}.
Since a bar index is preserved, the substitution
preserves balance and parity.

If an evaluation has
$\chi=\delta_{\alpha,B}\psi$ and the barred
relation on these boundary pairs, then
\eqref{eq:rg-reconstruction-local-identity}
turns \eqref{eq:rg-step-local-substitution}
into $T_k(B)\psi$ and $\bar T_k(B)\bar\psi$.
If in addition $\psi=H_k(B)\Phi$ and
$\bar\psi=\bar H_k(B)\bar\Phi$, the left-inverse
identity gives $\Phi$ and $\bar\Phi$ over $Y$.
Thus this local substitution represents the exact
mean field under evaluation.

The row covariances yield
\begin{equation}\label{eq:rg-step-substitution-covariances}
 \begin{split}
 a_g\operatorname{sub}_{gY,gB}
                        &=\operatorname{sub}_{Y,B}a_g,\\
 G_\lambda\operatorname{sub}_{Y,B^1}
                        &=\operatorname{sub}_{Y,B}G_\lambda,
                 \qquad B^1=B-d\lambda\text{ on bonds meeting }Y,\\
 \mathrm{Cg}\operatorname{sub}_{Y,-B}
                        &=\operatorname{sub}_{Y,B}\mathrm{Cg},\\
 \mathbf K\operatorname{sub}_{Y,-\overline B}
                        &=\operatorname{sub}_{Y,B}\mathbf K.
 \end{split}
\end{equation}
The final relation uses real $\mu$.
Both sides fix the old generators according to the
stated symmetry, so it suffices to compare their
images of $\Psi,\bar\Psi$. The lattice and real-gauge
relations follow from
\eqref{eq:rg-reconstruction-lattice-covariance} and
\eqref{eq:rg-reconstruction-gauge-covariance}; for
a local gauge equality the rows read only the
bonds on which that equality is assumed.
For charge conjugation use
$T_Y^{(i)}(-B)\mathsf C=\mathsf C\bar T_Y^{(i)}(B)$
and its barred counterpart. For the antilinear map,
\eqref{eq:rg-reconstruction-reality} with $-B$ gives
$\overline{T_Y^{(i)}(-\overline B)}
=\varkappa^{-1}T_Y^{(i)}(B)\varkappa$, and likewise
for the barred row. These identities give
\eqref{eq:rg-step-substitution-covariances} on
each new generator.

Assume the constants and charge conditions of
Proposition~\ref{prop:rg-reconstruction-estimate}, including
\eqref{eq:rg-domain-charge-conditions}, with
$0<\mu\le1$ and $0<e_k\le e_0$.
For $B\in\widetilde R_k+\mathfrak X_k$, the row
estimate gives the following norm statement. Assign weights
$w_1,w_2$ to the ordinary and pair legs and weight $w_3$
to the independent unit generators:
\begin{equation}\label{eq:rg-step-substitution-bound}
 \|\operatorname{sub}_{Y,B}F\|_{w_1,w_2}
                         \le\|F\|_{w_1,w_2;w_3}
             \quad\text{if }C_T(w_1+w_2)\le w_3.
\end{equation}
Indeed the image of a $\Psi$ generator has norm
at most $C_Tw_1+C_Tw_2$, and the original generators
retain their weights. The weighted linear-substitution
bound in Proposition~\ref{prop:weighted-Berezin-bounds}
applies on the subalgebra whose unit generators lie
over $Y$. The measures for the two row sums are
$\eta_k^3/48$ and $(\eta_k^3/48)^2$, respectively;
their identity kernels include the reciprocal
column measure. Thus this estimate has the same
normalization as the reconstruction bound.

\subsection{The first-step vertex}\label{subsec:rg-step-first-vertex}

At $k=0$ write the charge as $q$.
There is no input activity, and $H_0(B)=\bar H_0(B)=1$.
Consequently the exponent in
\eqref{eq:rg-step-fluctuation-exponent} is
\begin{equation}\label{eq:rg-step-first-exponent}
 V_0^x=
  -\langle\bar\Psi^f,(D_{\mathscr T+\mathscr Z}
                                    -D_{\mathscr T})\Psi^f\rangle^{\mathrm{ch}}
  -m\left(\langle\bar\Psi^f,\Psi^f\rangle^{\mathrm{ch}}
                   -\langle\bar\Psi^c,\Psi^c\rangle^{\mathrm{ch}}\right).
\end{equation}
The terms with $\mu$ and $m$ cancel in the connection
difference of the actions; the second difference
of interactions gives the displayed mass term.
Here $\Psi^c=\psi'_L$, $\bar\Psi^c=\bar\psi'_L$,
$\Gamma_0(\mathscr T)=F_0(\mathscr T)^{-1}$,
and $\Psi^f=\Psi^c+\Gamma_0(\mathscr T)W$.
All these are chamber fields on $T_0^{\mathrm f}$.
The identity measure in $W$ has mode measure $1/48$,
and its two-point kernel is $48$ times the identity
matrix of the modes.

For a bond fluctuation $Z$ and an $M$-cube $\square$,
define its local vertex by
\begin{equation}\label{eq:rg-step-first-vertex}
 \begin{split}
 \mathcal V_0(Z,\square)=\frac12\int_\square\dd x\,
   \Big\{&
    \bar\psi(x)^TM_3(x)(e^{iqZ(c,c_3)}-1)
                                  [\psi(x)+\chi(x,x^{(3)})]\\
    &+[\bar\psi(x)+\bar\chi(x,x^{(3)})]^T
               M_3(x^{(3)})(e^{-iqZ(c,c_3)}-1)\psi(x)\Big\},
 \end{split}
\end{equation}
where $x=(c,g)$, $c_3=c+ge_1$ and
$\int_\square\dd x=48^{-1}\sum_{c\in\square}\sum_g$.
The pair $(x,x^{(3)})$ has $d_X=1$ and is admissible
for $\square$ of the second kind. Thus the vertex
has ordinary legs over the cube and precisely the
boundary pair legs allowed in the activity algebra.
Its coefficients are entire functions of $Z$ on
the bonds meeting the cube.

\begin{proposition}\label{prop:rg-step-first-vertex}
Evaluate $\chi=\delta_{\alpha,\mathscr T}f$,
$\bar\chi=\delta_{\alpha,-\mathscr T}\bar f$,
$\psi=f$ and $\bar\psi=\bar f$ at any pair of
antiperiodic odd chamber fields. Then
\begin{equation}\label{eq:rg-step-first-vertex-evaluation}
 \operatorname{ev}\sum_\square\mathcal V_0(Z,\square)
          =-\langle\bar f,(D_{\mathscr T+Z}-D_{\mathscr T})
                                                    f\rangle^{\mathrm{ch}}.
\end{equation}
The two halves have equal images after summing
over cubes. The vertex has the lattice, real-gauge,
charge-conjugation and reality symmetries
\begin{equation}\label{eq:rg-step-first-vertex-symmetries}
 \begin{split}
 a_g\mathcal V_0(gZ,g\square)&=\mathcal V_0(Z,\square),\\
 G_\lambda\mathcal V_0(Z,\square)&=\mathcal V_0(Z,\square),\\
 \mathrm{Cg}\mathcal V_0(-Z,\square)&=\mathcal V_0(Z,\square),\\
 \mathbf K\mathcal V_0(-\overline Z,\square)&=\mathcal V_0(Z,\square).
 \end{split}
\end{equation}
For any generator weights $(w_1,w_2)$,
\begin{equation}\label{eq:rg-step-first-vertex-bound}
 \|\mathcal V_0(Z,\square)\|_{w_1,w_2}
       \le12(1+\rho)M^3w_1(w_1+w_2)
       \sup_{b\text{ oriented and meeting }\square}
                                                      |e^{iqZ(b)}-1|.
\end{equation}
If $q\le1$ and
$q|Z(b)|\le\frac12q^{1-\varepsilon/2}$ on these bonds,
then
\begin{equation}\label{eq:rg-step-first-vertex-small-bound}
 \|\mathcal V_0(Z,\square)\|_{h_0/2}
                 \le6(1+\rho)M^3q^{1/2-\varepsilon/2}
                 \le12M^3q^{1/2-\varepsilon/2}.
\end{equation}
\end{proposition}

\begin{proof}
At a nearest-neighbor pair the six paths reduce to
the bond, so the quotient relations read
\[
 f(x)+\chi(x,x^{(3)})=u_{\mathscr T}(c,c_3)f(x^{(3)}),
 \qquad
 \bar f(x)+\bar\chi(x,x^{(3)})
                         =u_{\mathscr T}(c,c_3)^{-1}\bar f(x^{(3)}).
\]
The first half of the summed vertex is therefore
one half of the integral of
\[
 \bar f(x)^TM_3(x)(e^{iqZ(c,c_3)}-1)
                                  u_{\mathscr T}(c,c_3)f(x^{(3)}).
\]
Reindexing the second half by the measure-preserving
involution $x\mapsto x^{(3)}$ gives the same
integrand: both $Z$ and $\mathscr T$ change sign
on the reversed bond. Only the type-$3$ hop
depends on the connection, and
$u_{\mathscr T+Z}=e^{iqZ}u_{\mathscr T}$.
The difference of the two Dirac operators is
thus the negative of this hopping coefficient,
proving \eqref{eq:rg-step-first-vertex-evaluation}.
Every product contains two antiperiodic fields,
so the reindexing is legitimate on the torus.

For lattice symmetry use
$g^{-1}(x^{(3)})=(g^{-1}x)^{(3)}$ and
$S^{-1}M_3(x)S=M_3(g^{-1}x)$ for the spin
lift $S$. Each half transforms separately, and
reindexing maps $g\square\times O_{\mathrm{cub}}$
onto $\square\times O_{\mathrm{cub}}$ with the
same measure. For a real gauge, the two generators
in each term have the same first photon site and
opposite bar indices, so their phases cancel.
Charge conjugation interchanges the two halves.
Its algebraic rule on the first half gives
\[
 \mathrm{Cg}[\bar\psi(x)^Tm(\psi(x)+\chi(x,x^{(3)}))]
       =(\bar\psi(x)+\bar\chi(x,x^{(3)}))^T
                              (\mathsf C^{-1}m^T\mathsf C)\psi(x).
\]
The sign from exchanging the odd factors cancels
the minus sign in the image of $\bar\psi$.
As
\[
 \mathsf C^{-1}M_3(x)^T\mathsf C=M_3(x^{(3)}),
\]
the first half at $-Z$ becomes the second at $Z$;
the reverse relation gives the other half.
Finally
\[
 \varkappa^T\overline{M_3(x)}\varkappa=M_3(x),\qquad
 \overline{e^{\pm iq(-\overline Z)}-1}=e^{\pm iqZ}-1.
\]
The antilinear map consequently preserves each
half in its reality identity. This proves
\eqref{eq:rg-step-first-vertex-symmetries}.

Put $\delta_x=e^{iqZ(c,c_3)}-1$ and
$\delta'_x=e^{-iqZ(c,c_3)}-1$. Collecting the
ordinary bilinear in \eqref{eq:rg-step-first-vertex},
its coefficient is
$[M_3(x)\delta_x+M_3(x^{(3)})\delta'_x]/96$.
The other two products have coefficients
$M_3(x)\delta_x/96$ and
$M_3(x^{(3)})\delta'_x/96$, and weights
$w_1w_2$.
The ordinary product has weight $w_1^2$.
The chamber sites in the cube and the ordered
pairs with these first sites are distinct, so
these are exactly the coefficients, up to
ordering signs. Each endpoint matrix satisfies
$\sum_{\beta,\beta'}|M_3(\cdot)_{\beta\beta'}|
\le12(1+\rho)$, since
$M_3=3(\rho\mp\gamma_\mu)$ and $\gamma_\mu$
has four entries of modulus one.
There are $48M^3$ chamber sites in the cube.
The triangle inequality for the collected
bilinear therefore gives
\eqref{eq:rg-step-first-vertex-bound}.

Under the smallness condition,
$|e^{iqZ(b)}-1|\le2q|Z(b)|\le q^{1-\varepsilon/2}$.
For $w=h_0/2$,
$w_1(w_1+w_2)=\frac14q^{-1/2}(1+q^\varepsilon)
\le\frac12q^{-1/2}$. Inserting these bounds proves
\eqref{eq:rg-step-first-vertex-small-bound}.
\end{proof}

For $Z=\mathscr Z$ and the evaluation at
$\Psi^f,\bar\Psi^f$, the first term of
\eqref{eq:rg-step-first-exponent} is the image of
the summed vertex. Under the additional charge
condition $96M^3q^{1/4}\le1$, its bound becomes
$\|\mathcal V_0(\mathscr Z,\square)\|_{h_0/2}
\le q^{1/4-5\varepsilon}$ whenever the displayed
bond smallness condition holds. These statements
apply before any expansion of the finite integral.

The first-step density identity retains every
vacuum and determinant factor:
\begin{equation}\label{eq:rg-step-first-density}
 \begin{split}
 (\mathbf T_0^x\mathscr F_0^x[\varepsilon^{\mathrm v},m,0])
                                    (A',\mathscr C';\Psi')\\
  ={}&\det D_1(0)^{-1}
       \exp\Big\{-\langle\bar\Psi',D_1(\mathscr T')\Psi'\rangle_1
                   -Lm\langle\bar\psi',\psi'\rangle\\
      &\hspace{14mm}
                   -L^3(\varepsilon^{\mathrm v}
                        +\delta\varepsilon_0^0)n_1
                   +\sum_X\mathfrak P_0(X,\mathscr C)\Big\}\\
       &\times
              \frac{\det F_0(\mathscr T)}{\det F_0(0)}
        \int\dd\mu_I(W)\,\dd\widehat\mu_0(\zeta)\,e^{V_0^x}.
 \end{split}
\end{equation}
The two determinants of $F_0$ have size $192n_0$.
The original level-zero density, integrated before
the block transformation, is
\begin{equation}\label{eq:rg-step-initial-integrated-density}
 \int\mathrm D\Psi\,\mathscr F_0^x
                     [\varepsilon^{\mathrm v},m,0](A,C;\Psi)
     =e^{-\varepsilon^{\mathrm v}n_0+\sum_X\mathfrak P_0(X,C)}
               \frac{\det(D_{A+C}+\mu+m)}{\det D_0(0)}.
\end{equation}
This follows directly from the weighted determinant
identity. The numerator need not be invertible;
the nonzero denominator uses $\mu>0$.


\section{The nonlinear small-field step}\label{sec:rg-nonlinear-step}

The finite fluctuation integral in
\eqref{eq:rg-step-fluctuation-exponent} has a logarithm which is a
sum of local activities. Its construction has two localizations.
The first resolves the change of background and reconstructs
the coarse mean; the second resolves the fermion fluctuations.
After these finite operations, a connected-polymer expansion
gives the logarithm. The estimates below apply also to the first
step, where the fluctuation generators have chamber measure
$1/48$.

\subsection{Contour domains and localization}
\label{subsec:rgnl-domains}

Keep the admissible sides, mass and charge of
Section~\ref{sec:rg-step}, and assume in addition
$\ell\ge Mr_k$. Put $e=e_k$, $h=(h_{k,1},h_{k,2})$.
The weakened photon and fermion primitives are taken on the
polydisk
\[
 \mathrm{Pd}_k=\{s=(s',s''):|s_{\mathrm{pc}}|\le M^{\alpha_0}\}.
\]
Here $s'=(s'_M,s'_{LM})$ contains the photon reconstruction
and square-root parameters, of kinds $\pi_M,\pi_{LM}$ in
\eqref{eq:photon-weakened-definitions}; $s''$ contains the
fermion $M$-cube parameters $\pi_k$ of
Section~\ref{subsec:weakened-walks}.
We assume their analytic kernel and H\"older estimates,
component restriction, the row-vanishing properties stated for
each family, and coefficient support
with common constants $(\alpha_0,\gamma_w,\kappa_w;C_w)$,
where $\kappa_w=\kappa+\kappa_0+3$.
The photon hypotheses are those of
Proposition~\ref{prop:photon-weakened-bounds},
\eqref{eq:photon-weakened-level-zero} and the component
restriction following them. The fermion hypotheses are those of
Propositions~\ref{prop:weakened-polymer-propagators} and
\ref{prop:coarse-local-family-bounds}, including
\eqref{eq:local-coarse-families},
\eqref{eq:local-level-zero-families}, and their covariance
identities. At parameter value one all families equal the
exact operators. Their gauge, lattice, transpose and reality
actions are the actions of those exact operators with the
parameters transformed accordingly. At level zero we retain
the stated diagonal families and the chamber covariance
with measure $1/48$.

The fixed parameters satisfy
\eqref{eq:rg-contraction-polymer-parameters},
\eqref{eq:rg-adjustment-charge-condition} and
\eqref{eq:rg-contraction-parameters},
$M>4L$, $L\ge4^8$, $0<1-\alpha<\varepsilon\le1/200$,
$C_w,c_\gamma\ge1$, and
$C_w\ge4\exp(2\gamma_w)M^{27\alpha_0}$.
The parameter radius satisfies
$M^{\alpha_0}-1\ge\exp(N_c\kappa+\log(6c_X))$,
where $N_c=(2L+1)^3$.
For later use put $c_S=10^5c_X^5K_0^5\exp(6\kappa)$.
The running charge lies in the fixed range on which
\begin{align}
 &p_{0,k}\le e^{-\varepsilon/4},\qquad
       12\exp(-p_{0,k}^2/2)\le e^7,\qquad
       2e^{5\varepsilon/2}\le1,\qquad
       c_{\mathrm{fl}}e^\varepsilon\le1,\qquad3Me^\varepsilon\le1,
                                                        \notag\\
 &2c_CC_w^2c_\gamma^2e^{\varepsilon/4}\le1,\qquad
       16c_Xc_\gamma C_wC_T^2M^3e^{\varepsilon/2}\le1,\qquad
       96M^3e^{1/4}\le1,                                \notag\\
 &32C_wC_Te^{3\varepsilon/2}\le1,\qquad
       88M^3e^{1/4-6\varepsilon}\le1,\qquad
       112C_w^2c_\gamma e^{2\varepsilon}\le3,             \notag\\
 &c_Se^{1/4-5\varepsilon}\le1,\qquad
       c_SL^3e^\varepsilon\le1,\qquad e^\varepsilon\le\tfrac14,\qquad
       8c_{\mathrm{rb}}L^3C_dM^3e^{2\varepsilon}\le1.
                                      \label{eq:rgnl-fixed-smallness}
\end{align}
The density identity also uses the photon domain-preservation
conditions of \eqref{eq:photon-domain-small-charge} and the
determinant conditions of \eqref{eq:determinant-gain-conditions}.
All are imposed by the fixed choice of the charge ceiling.

For $\zeta\in\mathrm{Box}_k$, write
\begin{equation}\label{eq:rgnl-weakened-photon}
 Z(s')=H_k^{\mathrm{ph}}(s')C C_k^{1/2}(s')\zeta .
\end{equation}
In this formula $C$ is the real photon coordinate map of
\eqref{eq:photon-coordinate-bijection}. The diagonal photon
family is used at $k=0$. For a bond field $Z$ set
\[
 |Z|_*=\max\{\|Z\|_\infty,\|\partial Z\|_\infty,
                              \|\delta_\alpha\partial Z\|_\infty\}.
\]
The three thresholds defining $R_k$ are
$\theta_i=e^{-3/4+i\varepsilon}$, $i=1,2,3$.
Thus $\widetilde R_k^-$ consists of the permitted real-gauge
part plus a field satisfying the three strict bounds
$(1-e^\varepsilon)\theta_i$.
Introduce the radii
\begin{equation}\label{eq:rgnl-contour-radii}
 r_1=\tfrac12e^{-3/4+9\varepsilon/2},\qquad
 r_2=e^{-3/4+7\varepsilon},\qquad
 r_4=e^{-1/4+3\varepsilon}.
\end{equation}

\begin{proposition}\label{prop:rgnl-contour-domains}
For every $s\in\mathrm{Pd}_k$ and
$\zeta\in\mathrm{Box}_k$,
\[
 |Z(s')|_*\le\tfrac12e^{-\varepsilon/2}.
\]
The radii satisfy $2\le r_2,r_4\le r_1$ and
$(r_i-1)^{-1}\le2/r_i$, $i=2,4$.
If $A\in\widetilde R_k^-$, $\mathscr C\in\mathfrak X_k$,
and $|Z|_*\le\tfrac12e^{-\varepsilon/2}$, then
\[
 A+tZ\in\widetilde R_k,\qquad
 A+\mathscr C+tZ\in\widetilde{\mathcal G}_k
                  \quad (|t|\le r_1).
\]
\end{proposition}

\begin{proof}
The square-root kernel is bounded by
$C_w\exp(-\gamma_w d(\Upsilon,\Upsilon'))$.
The summability constant $c_\gamma$ and
the coordinate bound $|\zeta(\Upsilon')|\le p_{0,k}
\le e^{-\varepsilon/4}$ give
\[
 \sup_\Upsilon|(C_k^{1/2}(s')\zeta)(\Upsilon)|
                                      \le C_wc_\gamma e^{-\varepsilon/4}.
\]
The coordinate map is real. Its bound
$\|Cv\|_\infty\le c_C\sup_\Upsilon|v(\Upsilon)|$
also gives the absolute row bound
$\sum_\Upsilon|C(b,\Upsilon)|\le c_C$:
apply the real bound to the signs of the entries of a fixed
row. Therefore
\[
 \|CC_k^{1/2}(s')\zeta\|_\infty
                                  \le c_CC_wc_\gamma e^{-\varepsilon/4}.
\]
Each of the reconstruction, derivative and H\"older derivative
kernels has the same summable exponential bound. Consequently
\[
 |Z(s')|_*\le c_CC_w^2c_\gamma^2 e^{-\varepsilon/4}
                                          \le\tfrac12 e^{-\varepsilon/2},
\]
where the last inequality is
$2c_CC_w^2c_\gamma^2e^{\varepsilon/4}\le1$.
At level zero the polynomial diagonal family has the same
row bounds: its constant is
$4\exp(2\gamma_w)M^{27\alpha_0}\le C_w$.

Since $\varepsilon\le1/200$,
\[
 \frac{r_2}{r_1}=2e^{5\varepsilon/2}\le1,\qquad
 \frac{r_4}{r_1}=2e^{1/2-3\varepsilon/2}
                                      \le2e^{5\varepsilon/2}\le1.
\]
The small-charge condition
$88M^3e^{1/4-6\varepsilon}\le1$ gives
$r_4^{-1}=e^{1/4-3\varepsilon}\le1/88$ and
$r_2^{-1}\le r_4^{-1}$.
In particular both radii are at least $2$, and their stated
inverse bounds follow.

For $|t|\le r_1$ the three norms of $tZ$ are at most
$\frac14e^{-3/4+4\varepsilon}$, strictly less than
$e^\varepsilon\theta_3$. Adding this field fills less than the
margin in each of the three defining bounds of
$\widetilde R_k^-$. The permitted gauge part is unchanged.
Hence $A+tZ\in\widetilde R_k$.
Finally $\widetilde R_k+\mathfrak X_k
\subset\widetilde{\mathcal G}_k$ gives the second inclusion.
\end{proof}
The first-step vertex estimate
\eqref{eq:rg-step-first-vertex-small-bound} now applies to
$Z(s')$ for every weakening parameter. In particular,
$96M^3e^{1/4}\le1$ gives its bound
$e^{1/4-5\varepsilon}$ in the $h/2$ norm.

Let $B=A+\mathscr C$.
The arguments of $H_k^\#$ are independent pairs
$(\Phi,\bar\Phi)$; in $H_k^\#\Phi$ the single symbol denotes
this pair. Write
$\psi^\#=(\psi,\chi;\bar\psi,\bar\chi)$ for the ordinary,
pair, barred ordinary and barred pair generators, respectively.
Then
\[
 \begin{aligned}
 H_k^\#(s'',B)\Phi
   ={}&(H_k(s'',B)\Phi,\delta_{\alpha,B}H_k(s'',B)\Phi;\\
       &
       \bar H_k(s'',B)\bar\Phi,
                       \delta_{\alpha,-B}\bar H_k(s'',B)\bar\Phi),\\
 \bar H_k(s'',B)&=\mathsf C^{-1}H_k(s'',-B)\mathsf C.
 \end{aligned}
\]
Define
$J_k^\#(s'',B,Z)=H_k^\#(s'',B+Z)-H_k^\#(s'',B)$ and
\[
 \widetilde\Gamma_k(s'',B)(W,\bar W)
                         =(\Gamma_k(s'',B)W,\bar W),\qquad
 W_k^\#(s'',B)=H_k^\#(s'',B)
                         \big(\Gamma_k(s'',B)W,\bar W\big).
\]
The notation $J_k,\bar J_k$ denotes the components without
H\"older quotient; $\mathrm{Wc},\overline{\mathrm{Wc}}$ denotes
the corresponding components of $W_k^\#$.
The auxiliary generators $W$ have weight
$\mathsf w_W=1$ for $k\ge1$ and $\sqrt{48}$ for $k=0$.

These fields are continuous in all closed polydisk and box
variables and analytic in the open parameter polydisk and the
connection variables. For the H\"older rows this follows
directly from
\[
 (\delta_{\alpha,B}H_k)((x,x'),u)
   =d_X(x,x')^{-\alpha}
       \big(U_B(x,x')H_k(x',u)-H_k(x,u)\big).
\]
The transport is entire in $B$. Proposition~\ref{prop:rgnl-contour-domains}
places every connection on the closed contours inside the
open domain of the primitives. Thus compositions at
$B+tZ(s')$ have the stated regularity also when the contour
variable is on its circle.

In the following localization, cubes and polymers are sets
of photon sites. Containment refers to all their lattice sites,
and two cubes meet when they share a lattice site.
Let $P_0(X)$ be the parameter cubes contained in $X$.
For $P$ outside $P_0(X)$ let $Y'(X,P)$ be the union of
the $M$-cubes meeting $X$ or a cube of $P_0(X)\cup P$.
For a continuous analytic family $G(s,X)$ define
\begin{align}
 \mathcal T_PG(X)
  &=\int_{[0,1]^P}\dd s_P\,
       \partial_PG(s,X)
       \big|_{s=1\text{ on }P_0(X),\,
                    s=0\text{ outside }P_0(X)\cup P},
                          \label{eq:rgnl-interpolation}\\
 \operatorname{Loc}G(Y)
  &=\sum_{(X,P):Y'(X,P)=Y}\mathcal T_PG(X).
                          \label{eq:rgnl-localization}
\end{align}
Repeated one-variable integration gives the useful identity
\begin{equation}\label{eq:rgnl-finite-difference}
 \mathcal T_PG(X)=
    \sum_{Q\subseteq P}(-1)^{|P|-|Q|}
                                 G(\mathbf1_{P_0(X)\cup Q},X).
\end{equation}
In particular, coefficient locality can be checked at
indicator families; no statement about the convergence of
the later logarithm is used here.

\begin{lemma}\label{lem:rgnl-resummation}
Suppose
$\|G(s,X)\|\le B_0\exp(-\kappa'd_M(X))$ on the parameter
polydisk, where $\kappa_0\le\kappa'\le\kappa$.
Assume $\mathcal T_PG(X)=0$ whenever $Y'(X,P)$ is
disconnected. Then
\begin{gather}
 \sum_Y\operatorname{Loc}G(Y)=\sum_XG(1,X),
                         \label{eq:rgnl-localization-sum}\\
 \|\operatorname{Loc}G(Y)\|
   \le4c_XK_0B_0
             \exp\big(-(\kappa'-\kappa_0-1)d_M(Y)\big).
                         \label{eq:rgnl-resummation-bound}
\end{gather}
\end{lemma}

\begin{proof}
Applying $f(1)=f(0)+\int_0^1f'(t)\dd t$ in each parameter
outside $P_0(X)$ gives $G(1,X)=\sum_P\mathcal T_PG(X)$.
The sums are finite. After removing the vanishing disconnected
terms, each term belongs to exactly one polymer $Y'(X,P)$.
This proves \eqref{eq:rgnl-localization-sum}.

Put $R=M^{\alpha_0}$. Cauchy's estimate on a polydisk of
radius less than $R-1$ about each integration point, followed
by the limit of these radii to $R-1$, gives
\[
 \|\mathcal T_PG(X)\|
       \le B_0\exp(-\kappa'd_M(X))(R-1)^{-|P|}.
\]
If $Y'(X,P)=Y$, then $X\subseteq Y$ and
$|Y\setminus X|_M\le N_c|P|$, since every added $M$-cube
meets a parameter cube and each parameter cube meets at
most $N_c=(2L+1)^3$ such cubes.
To compare their tree lengths, order the cubes of $Y\setminus X$
by graph distance from $X$. Each added cube is adjacent
to a preceding one. Extend a net by a segment from
a point of that preceding closed cube to their common
face, of length at most $M$. Starting from an
arbitrarily near-minimal net for $X$ gives
$d_M(Y)\le d_M(X)+|Y\setminus X|_M$.
Conversely every net for $Y$ is a net for $X$.
Thus $0\le d_M(Y)-d_M(X)\le N_c|P|$.
Writing $a_X=\log(6c_X)$, the condition
$R-1\ge\exp(N_c\kappa+a_X)$ implies
\[
 \|\mathcal T_PG(X)\|
 \le B_0\exp(-(\kappa'-\kappa_0)d_M(Y))
             \exp(-\kappa_0d_M(X))\exp(-a_X|P|).
\]
At most $3|Y|_M$ parameter cubes are contained in $Y$.
Hence
\[
 \sum_P\exp(-a_X|P|)
       \le(1+(6c_X)^{-1})^{3|Y|_M}
       \le \exp((d_M(Y)+1)/2).
\]
The polymer counting and size estimates give
\[
 \sum_{X\subseteq Y}\exp(-\kappa_0d_M(X))
   \le K_0|Y|_M
   \le2c_XK_0\exp(d_M(Y)/2).
\]
The product of the last two bounds is at most
$4c_XK_0\exp(d_M(Y))$. Summing the estimate for
$\mathcal T_PG(X)$ proves
\eqref{eq:rgnl-resummation-bound}.
\end{proof}

\subsection{The two finite localizations}
\label{subsec:rgnl-finite-localizations}

For a function analytic near the closed disk $|t|\le r$, $r>1$,
write
\[
 \mathcal C_rf=\frac1{2\pi i}\int_{|t|=r}
                                      \frac{f(t)}{t(t-1)}\,\dd t .
\]
This is $f(1)-f(0)$ and satisfies
$\|\mathcal C_rf\|\le(r-1)^{-1}\sup_{|t|=r}\|f(t)\|$.
Let $(m,\mathscr E)\in\operatorname{Ball}_k^x(4)$ and set
$m'=m+m(\mathscr E)$, $\mathscr E^r=R\mathscr E$.
For $k\ge1$ define
\begin{align}
 G_1(s,X)&=\mathcal C_{r_1}
                \langle\bar\Psi,M_k(X,B+tZ(s'))\Psi\rangle_1,
                                      \label{eq:rgnl-G1}\\
 G_2(s,X)&=\mathcal C_{r_2}\,
       \mathscr E^r(X;A+tZ(s'),\mathscr C,
                    \psi^\#+tJ_k^\#(s'',B,Z(s'))\Psi),
                                      \label{eq:rgnl-G2}\\
 G_3(s,\square)&=\mathcal C_{r_2}(-m')
       \langle\bar\psi+t\bar J_k(s'',B,Z(s'))\bar\Psi,
                         \psi+tJ_k(s'',B,Z(s'))\Psi\rangle_\square .
                                      \label{eq:rgnl-G3}
\end{align}
The symbols $M_k(X,\cdot)$ denote the local Dirac primitives.
A function defined for one cube is zero on the other polymers.
Put
\[
 E_i^*(Y)=\operatorname{Loc}G_i(Y),\qquad
 E_i(Y)=\operatorname{sub}_{Y,B}E_i^*(Y),\qquad
 E_{\le3}=\sum_{i=1}^3E_i.
\]
The substitution is \eqref{eq:rg-step-local-substitution}.
At $k=0$ there is no $\Psi$ and no substitution:
$G_1(s,\square)=\mathcal V_0(Z(s'),\square)$,
$G_2=G_3=0$, and $E_{\le3}=\operatorname{Loc}G_1$.
All parameter kinds occur in this first localization.

For the second localization, which uses only $\pi_k$, let
$v=W_k^\#(s'',B)$ and define
\begin{align}
 G_{\le3}^{\,W}(s'',X)&=E_{\le3}(X;A,\mathscr C,\zeta,
                                                      \psi^\#+v),
                                      \label{eq:rgnl-GW}\\
 G_4(s'',X)&=\mathcal C_{r_4}
                  \mathscr E^r(X;A,\mathscr C,\psi^\#+tv),\\
 G_5(s'',\square)&=\mathcal C_{r_4}(-m')
                 \langle\bar\psi+t\overline{\mathrm{Wc}},
                                   \psi+t\mathrm{Wc}\rangle_\square,\\
 F(Y)&=\operatorname{Loc}(G_{\le3}^{\,W}+G_4+G_5)(Y).
                                      \label{eq:rgnl-finite-activity}
\end{align}
At $k=0$ the term $G_4$ is zero and $m'=m$.

\begin{proposition}\label{prop:rgnl-finite-locality}
All finite expressions just defined exist on
$\widetilde R_k^-\times\mathfrak X_k\times\mathrm{Box}_k$.
The element $F(Y)$ is even, has ordinary legs over $Y$ and
pair legs admissible for $Y$, and has fluctuation generators
only at unit indices in $Y$. Its coefficients read $A,\mathscr C$
only on bonds meeting $Y$ and read $\zeta$ only through the
coordinates $\zeta(\Upsilon)$ with $\supp\Upsilon\subset Y$.
They are continuous in these arguments and analytic in the
connection variables. The same assertions hold for $E_{\le3}(Y)$,
which has no fluctuation generators.

The disconnected terms in both localizations vanish, and
\begin{equation}\label{eq:rgnl-two-localization-sums}
 \sum_YE_i^*(Y)=\sum_XG_i(1,X),\qquad
 \sum_YF(Y)=\sum_X(G_{\le3}^{\,W}+G_4+G_5)(1,X).
\end{equation}
The finite expressions on $Y$ depend on the datum only through
$m'$ and the $\mathscr E^r(X)$ with $X\subseteq Y$.
\end{proposition}

\begin{proof}
We give the component argument which is responsible for
locality. Let $\mathrm{Ps}=P_0(X)\cup P$, let
$Y'=Y'(X,P)$ and let $Y'_1$ be its component containing $X$.
Set the weakening parameters to zero outside $\mathrm{Ps}$.
For each of the three parameter kinds, any link-component
$\mathrm{Lnk}$ of $\mathrm{Ps}$ is contained in the polymer
formed by the $M$-cubes meeting it. This polymer lies in one
component of $Y'$. If a site of $Y'_1$ lies in $\mathrm{Lnk}$,
or in the enlarged parameter cube of a photon kind, that
polymer has a cube in common with $Y'_1$, and hence lies
entirely in $Y'_1$.

The row-vanishing property of the weakened photon families
and their component restriction give, on any fine bond $b$
meeting $Y'_1$,
\begin{equation}\label{eq:rgnl-restricted-photon}
 Z(s')(b)=
  \sum_{\substack{b',\Upsilon,\Upsilon'\subset Y'_1}}
   H_k^{\mathrm{ph}}(s')(b,b')C(b',\Upsilon)
          C_k^{1/2}(s')(\Upsilon,\Upsilon')\zeta(\Upsilon').
\end{equation}
Indeed, the reconstruction row is zero unless $b$ lies in
one link-component; its columns lie in that component.
The coordinate map has columns supported in the $L$-blocks
of the ends of $b'$, and each block lies in the $M$-cube of
the corresponding end. The square-root row then has columns
in its own link-component. Each component involved is
attached to $Y'_1$ by the preceding argument.
These rows are unchanged if the parameters outside $Y'_1$
are changed while their common zero set outside $\mathrm{Ps}$
is kept.

For a chamber site $x=(c,g)$ at distance at most $1$ from
$Y'_1$, the fermion row is zero unless every parameter cube
meeting $K_c$ is present. These cubes form one
link-component. When the row is nonzero, that component
contains $c$ and a site of $Y'_1$, so it lies in $Y'_1$.
The row of $H_k(s'',B)$, and its barred counterpart,
then has columns in that component and reads $B$ only on
bonds meeting it. At a column of this row, the corresponding
row of $\Gamma_k(s'',B)$ has columns in the same component.
These conclusions also hold at level zero: the diagonal
factor is
\[
 \prod_{\mathrm{pc}\in\pi_0:\mathrm{pc}\cap K_c\ne\varnothing}
                                                       s''_{\mathrm{pc}},
\]
and the kernel of that multiplication operator is
$48\mathbf1_{x=x'}$ times this factor. Its composition with
$\Gamma_0$ uses the measure $1/48$.

An endpoint of a pair admissible for $X$ is at distance
at most $\eta_k\le1$ from $X$. Its transport reads only
bonds contained in $X$, or the one type-$3$ bond meeting $X$.
Thus the preceding row argument applies to every row
used in $J_k^\#$ and $W_k^\#$.
The H\"older row has only the two columns $x,x'$ and
the entire transport coefficient displayed above.
Finite compositions preserve this dependence: the columns
of $O_2O_1$ are the union of the columns of $O_1$ at the
finitely many columns of $O_2$, and its read bonds are the
union of their read bonds. Entries of the transport and
reconstruction rows are entire in the connections, and finite
compositions of these algebraic rows retain this property.
Compositions involving weakened primitives are holomorphic
on their stated primitive domains. These are the finite-row
properties used here; the weakened operators themselves
are not asserted to have finite range.

For $G_1$ at $k\ge1$, the local primitive $M_k(X,\cdot)$
has both indices in $X$ and reads only bonds meeting $X$.
For $G_2,G_3$, the incoming activity or mass form has ordinary
legs in $X$ and admissible pair legs; the substitution by
$J_k^\#\Psi$ introduces $\Psi$ only in $Y'_1$.
At $k=0$ the vertex has precisely these ordinary and
type-$3$ pair legs and reads $Z$ only on bonds meeting
its cube. Equation~\eqref{eq:rgnl-restricted-photon}
therefore proves the same component dependence for every
first-localization function.

If $Y'$ is disconnected, choose an $M$-cube outside $Y'_1$.
It cannot meet $X$ or a parameter cube contained in $X$,
so it meets a cube of $P$ outside $Y'_1$. The integrand is
independent of that parameter. Pairing the subsets in
\eqref{eq:rgnl-finite-difference} according to whether they
contain this cube proves that the term is zero.
For a term assigned to a connected $Y$, all parameter cubes
in $P_0(X)\cup P$ and $X$ lie in $Y$.
Equation~\eqref{eq:rgnl-finite-difference} now gives the
claimed locality of $E_i^*$.

The reconstruction substitution adds only ordinary legs
over $Y$ and boundary pairs $(x,x^{(3)})$ admissible for
$Y$. Its finite rows at a unit site $y$ read bonds contained
in $\Delta_y$ or type-$3$ bonds from this block. These bonds
meet $Y$ since $\Delta_y\subset Y$. The barred rows have
the same support. Thus it preserves the claimed coefficient
locality and removes $\Psi$.
For the second localization, the row argument for
$H_k^\#\widetilde\Gamma_kW$ gives fluctuation generators
only in $Y'_1$. It applies to each of the three functions
in \eqref{eq:rgnl-finite-activity}. The same disconnected
parameter cancellation proves the second assertion.
Lemma~\ref{lem:rgnl-resummation} gives both sum identities.

Regularity follows from the established contour domains,
the primitive regularity, and finite sums and substitutions.
Contour integrals preserve analyticity: their Riemann sums
converge uniformly on compact sets. The same argument
applies to finite parameter integrals. The incoming
activities and mass forms are even, and each substitution
sends generators to odd linear combinations, so every
expression remains even.
Finally $G_1$ is independent of the datum; $G_2,G_4$
use $\mathscr E^r(X)$ and $G_3,G_5$ use $m'$.
Tracing the two localizations shows that no incoming
polymer outside $Y$ enters the result.
\end{proof}

The proof identifies the primitive entries determining a
finite expression on $Y$. With parameters zero outside $Y$,
the photon field reads only
$H_k^{\mathrm{ph}}(s')(b,b')$ for $b$ meeting $Y$,
$b'\subset Y$, and
$C_k^{1/2}(s')(\Upsilon,\Upsilon')$ with both supports
in $Y$. Its remaining factor is the fixed coordinate map.
The fermion fields use $H_k,\bar H_k$ at chamber sites
over $Y$, with unit columns in $Y$, at $B+Z(s')$ and
$B$, and $\Gamma_k(s'',B)$ with both unit indices in $Y$.
At zero, the kernel sums have measure $1/48$ and the
diagonal rows have the factor $48$ already recorded.
The Dirac primitive has both indices in $X\subseteq Y$
and connection $B+tZ(s')$ on $|t|=r_1$.
H\"older transports read only bonds meeting $Y$.
The reconstruction substitution contains only the fixed
finite $T_k^{(i)},\bar T_k^{(i)}$ rows at unit sites of its
own subpolymer.

Consequently two systems with the stated hypotheses on this torus give the
same finite expressions on $Y$ if these entries agree
for every parameter family zero outside $Y$, at the
connections formed with their respective photon fields.
The restricted photon sum first gives equality of the
fields on bonds meeting $Y$; the transports and incoming
activities there then agree by their locality.
The listed fermion and Dirac entries give equal first
functions, their localizations and reconstruction
substitutions. They give equal second functions and
localizations in turn, and then equal activities $K(Y)$ and
$H(Y)$ in \eqref{eq:rgnl-Mayer-activity} and
\eqref{eq:rgnl-integrated-activity}, and each term of the
connected series \eqref{eq:rgnl-log-activity}.
Agreement of the photon
entries and agreement of the fermion entries as functions
on the full operator domain is sufficient: if the two
composed connections differ outside $Y$, local row
dependence removes that difference first.
The same comparison passes to the convergent logarithm
once its convergence has been established.

\subsection{Bounds before fluctuation integration}
\label{subsec:rgnl-finite-bounds}

Set
\[
 \begin{gathered}
 \kappa_w=\kappa+\kappa_0+3,\qquad
 a_1=\kappa-\kappa_0-1,\qquad a_2=\kappa-2\kappa_0-2,\\
 C_W=28C_w^2c_\gamma,
 \qquad H_{\mathrm{cl}}=36c_X^2K_0^2e^{1/4-5\varepsilon}.
 \end{gathered}
\]
The notation $H_{\mathrm{cl}}$ is a scalar cluster bound.
The finite primitive bounds imply
\begin{align}
 \|\langle\bar\Psi,M_k(X,B)\Psi\rangle_1\|_{w_3}
 &\le4c_Xc_\gamma C_wM^3w_3^2
                   \exp(-(\kappa_w-1)d_M(X)),\quad k\ge1,
                                      \label{eq:rgnl-M-bound}\\
 |H_k^\#(s'',B)|_{1,\infty}&\le4C_w,\\
 |J_k^\#(s'',B,Z)|_{1,\infty}
                  &\le16C_we^{3/4-9\varepsilon/2},\\
 \mathsf w_W|H_k^\#(s'',B)\widetilde\Gamma_k(s'',B)|_{1,\infty}
                  &\le C_W.           \label{eq:rgnl-W-bound}
\end{align}
Here $B=A+\mathscr C$ and $|Z|_*\le\frac12e^{-\varepsilon/2}$.
To see the first bound, count $4M^3|X|_M$ rows and sum
the exponential kernel in its second index, giving
$c_\gamma C_w \exp(-\kappa_wd_M(X))$ per row.
Then $|X|_M\le c_X\exp(d_M(X))$ gives
\eqref{eq:rgnl-M-bound}.
For $H_k^\#$ the primitive bound controls the ordinary
and short-pair rows by $C_w$. A long pair, $d_X\ge1$,
has its H\"older quotient bounded by four times the
ordinary sup norm. At level zero every pair is long.
The monomial matrices $\mathsf C^{\pm1}$ preserve row norms,
so the barred rows have the same bounds.
Apply $\mathcal C_{r_1}$ to
$H_k^\#(s'',B+tZ)$ to obtain the bound for $J_k^\#$.
Finally $|\Gamma_k|_{1,\infty}\le C_wc_\gamma$ by kernel
summation, using $1/48$ at level zero. Submultiplicativity
and $\mathsf w_W\le7$ give \eqref{eq:rgnl-W-bound}.

\begin{proposition}\label{prop:rgnl-finite-bounds}
For the data of Proposition~\ref{prop:rgnl-finite-locality},
\begin{align}
 \|E_{\le3}(Y)\|_{h/2}
       &\le8c_XK_0e^{1/4-5\varepsilon}\exp(-a_1d_M(Y)),
                                      \label{eq:rgnl-E3-bound}\\
 \|F(Y)\|_{(h/4;\mathsf w_W)}
       &\le H_{\mathrm{cl}}\exp(-a_2d_M(Y)).
                                      \label{eq:rgnl-F-bound}
\end{align}
The element $F(Y)$ is the sum of a term independent of
$(m,\mathscr E)$ and a term complex linear in these data.
\end{proposition}

\begin{proof}
Extraction gives
\[
 \|\mathscr E^r(X)\|_h\le2\|\mathscr E\|_k^x\exp(-\kappa d_M(X))
                   \le12e^{1/4-7\varepsilon}\exp(-\kappa d_M(X)),
 \qquad |m'|\le8e^{3/4-8\varepsilon}.
\]
The mass bound uses
$c_{\mathrm{fl}}e^\varepsilon\le1$.
With chamber measure $\eta_k^3/48$, direct counting gives
\[
 \|\langle\bar\psi,\psi\rangle_\square\|_{(w_1,w_2)}
                                                    =4M^3w_1^2.
\]
The number of monomials is
$4\cdot48\eta_k^{-3}M^3$ and their common coefficient
is $\eta_k^3/48$.

In the first localization give $\Psi$ the weight
$w_3=C_Th_{k,1}$ and the old generators the weights $h/2$.
The contour bound and \eqref{eq:rgnl-M-bound} give
\[
 \|G_1(s,X)\|\le
 16c_Xc_\gamma C_wC_T^2M^3e^{1/4-9\varepsilon/2}
                 \exp(-(\kappa_w-1)d_M(X))
 \le e^{1/4-5\varepsilon}\exp(-\kappa d_M(X)).
\]
The last step uses
\[
 16c_Xc_\gamma C_wC_T^2M^3e^{\varepsilon/2}\le1.
\]
At $k=0$ the weighted vertex instead gives
$12M^3e^{1/2-\varepsilon/2}\le e^{1/4-5\varepsilon}$.
Moreover,
\[
 r_2|J_k^\#|_{1,\infty}w_3
     \le16C_wC_Te^{3\varepsilon/2}h_{k,2}
     \le\tfrac12h_{k,2}.
\]
The substitution bound therefore takes the integrand norm
of $G_2$ to the $h$ norm of $\mathscr E^r$, and that of
$G_3$ to the $h$ norm of the mass form. It gives
\[
 \|G_2(s,X)\|\le24e^{1-14\varepsilon}\exp(-\kappa d_M(X)),
 \qquad \|G_3(s,\square)\|\le64M^3e^{1-15\varepsilon}.
\]
The condition $88M^3e^{1/4-6\varepsilon}\le1$ implies
\[
 24e^{3/4-9\varepsilon}
                   +64M^3e^{3/4-10\varepsilon}\le1.
\]
Thus $\sum_{i\le3}G_i$ is bounded by
$2e^{1/4-5\varepsilon}\exp(-\kappa d_M(X))$.
Apply Lemma~\ref{lem:rgnl-resummation}, using the disconnected
cancellation in Proposition~\ref{prop:rgnl-finite-locality}.
The reconstruction substitution costs no norm, since
$C_T(h_{k,1}+h_{k,2})/2\le w_3$.
This proves \eqref{eq:rgnl-E3-bound}.

For the second localization, the condition
$112C_w^2c_\gamma e^{2\varepsilon}\le3$ gives
\[
 r_4C_W\le\tfrac34h_{k,2},\qquad C_W\le\tfrac14h_{k,2}.
\]
Consequently addition of $W_k^\#$ maps the $h/4$ norm
to the $h/2$ norm at $t=1$ and to the $h$ norm on
the contour $|t|=r_4$. The three input bounds are
\begin{align*}
 \|G_{\le3}^{\,W}(s'',X)\|
       &\le8c_XK_0e^{1/4-5\varepsilon}\exp(-a_1d_M(X)),\\
 \|G_4(s'',X)\|&\le24e^{1/2-10\varepsilon}\exp(-\kappa d_M(X)),\\
 \|G_5(s'',\square)\|&\le64M^3e^{1/2-11\varepsilon}.
\end{align*}
Their sum is at most
$9c_XK_0e^{1/4-5\varepsilon}\exp(-a_1d_M(X))$,
again by $88M^3e^{1/4-6\varepsilon}\le1$.
Since $\kappa_0\le a_1\le\kappa$, resummation gives
\eqref{eq:rgnl-F-bound}.
At level zero $G_4=0$; all the other inequalities include
the factor $\mathsf w_W$ already, so they apply unchanged.
Linearity follows from the definitions: $G_1$ is fixed,
$G_2,G_4$ are linear in $R\mathscr E$, $G_3,G_5$ in $m'$,
and localization, the reconstruction substitution and the
contour functionals are linear.
\end{proof}

\subsection{Connected polymers and the logarithm}
\label{subsec:rgnl-cluster-logarithm}

Two polymers are incompatible when their union is connected
through faces. In particular a polymer is incompatible with
itself. For an ordered tuple $(X_1,\ldots,X_n)$ define
$u^T(X_1)=1$ and
\[
 u^T(X_1,\ldots,X_n)=\sum_{G\text{ connected}}
                                      (-1)^{|G|},\quad n\ge2,
\]
where $G$ ranges over spanning subgraphs of the incompatibility
graph of the tuple. The coefficient is zero when that graph
is disconnected.

We will use the formal connected-graph identity
\begin{equation}\label{eq:rgnl-connected-graph-identity}
 \sum_{\{X_i\}\text{ compatible}}\prod_i z_{X_i}
 =\exp\left\{\sum_{n\ge1}\frac1{n!}
       \sum_{(X_1,\ldots,X_n)}u^T(X_1,\ldots,X_n)
                                           \prod_i z_{X_i}\right\}
\end{equation}
and the tree bound
\begin{equation}\label{eq:rgnl-tree-domination}
 |u^T(X_1,\ldots,X_n)|
       \le\#\{\text{spanning trees of the incompatibility graph}\}.
\end{equation}
Here the identity is one of formal series in commuting
variables, before substitution into a normed algebra.
For its proof, express the left side as an ordered sum with
factor $1/n!$ and multiply by
$\prod_{i<j}(1-\mathbf1_{\{X_i,X_j\}\text{ incompatible}})$.
Expanding this product sums over all graphs. A graph is
uniquely the union of its connected components. Grouping by
their vertex sets gives the exponential on the right:
an ordered list of component sizes $n_1,\ldots,n_r$ has
$n!/(n_1!\cdots n_r!)$ vertex assignments, and division
by $r!$ removes the order of its components.
Only finitely many terms contribute to a fixed degree.

For \eqref{eq:rgnl-tree-domination}, let $\mathbf G$ be
the incompatibility graph and suppose it is connected.
For a connected spanning subgraph $G$, root at $1$ and,
for $i\ne1$, choose as its parent the least labelled
neighbor at distance $d_G(i)-1$ from the root.
These edges form a tree $T(G)$ with the same root distances
as $G$. Given a spanning tree $\tau$, let $R(\tau)$ contain
its edges, all edges of $\mathbf G$ within one distance
level, and all edges from $i$ to the preceding level whose
other endpoint exceeds the parent of $i$ in $\tau$.
Then
\[
 T(G)=\tau\quad\Longleftrightarrow\quad
                                      \tau\subseteq G\subseteq R(\tau).
\]
In one direction the distances and least parents force
every edge of $G$ to belong to $R(\tau)$.
Conversely, along any edge of $R(\tau)$ the tree distance
changes by at most one. Thus a graph between $\tau$ and
$R(\tau)$ has the tree's distances, and the least neighbor
in the preceding level is its tree parent.
The connected subgraphs are therefore partitioned into
these intervals, and their signed sum is
\[
 \sum_\tau(-1)^{n-1}
           \sum_{A\subseteq R(\tau)\setminus\tau}(-1)^{|A|}
 =(-1)^{n-1}\#\{\tau:R(\tau)=\tau\}.
\]
This proves the tree bound, including the cases
of a disconnected graph and $n=1$.

\begin{proposition}\label{prop:rgnl-cluster-logarithm}
Let $V(X)$ be even elements in the old and fluctuation
generators, continuous on $\mathrm{Box}_k$, with
\[
 \|V(X)\|_{(w;\mathsf w_W)}
                       \le H_{\mathrm{cl}}\exp(-a_2d_M(X)).
\]
Assume that $V(X)$ has fluctuation indices only in $X$ and
depends on $\zeta$ only through coordinates supported in $X$.
For any union $Y$ of $M$-cubes put
\begin{align}
 K(Y)&=\sum_{\substack{\{X_i\}\text{ distinct}\\\bigcup_iX_i=Y}}
                    \prod_i(\exp V(X_i)-1),\qquad K(\varnothing)=1,
                                      \label{eq:rgnl-Mayer-activity}\\
 H(Y)&=\int\dd\mu_I(W)\,\dd\widehat\mu(\zeta)\,K(Y).
                                      \label{eq:rgnl-integrated-activity}
\end{align}
For a connected $Y$ define
\begin{equation}\label{eq:rgnl-log-activity}
 E^\#(Y)=\sum_{n\ge1}\frac1{n!}
   \sum_{\bigcup_iY_i=Y}u^T(Y_1,\ldots,Y_n)\prod_iH(Y_i).
\end{equation}
The last series converges absolutely in the $w$ norm and
\begin{align}
 \int\dd\mu_I(W)\,\dd\widehat\mu(\zeta)
                         \exp\left\{\sum_XV(X)\right\}
       &=\exp\left\{\sum_YE^\#(Y)\right\},
                                      \label{eq:rgnl-log-identity}\\
 \|E^\#(Y)\|_w
       &\le c_S'H_{\mathrm{cl}}\exp(-a_3d_M(Y)),
                                      \label{eq:rgnl-log-bound}
\end{align}
Here $a_3=\kappa-5\kappa_0-4$ and
$c_S'=14c_X^2K_0^2\exp(3\kappa)$.
The sum of the norms of the terms with $n>N$ in
\eqref{eq:rgnl-log-activity} is at most $2^{-N}$ times the
right side of \eqref{eq:rgnl-log-bound}.
\end{proposition}

\begin{proof}
Set
\[
 \begin{gathered}
 a_4=\kappa-3\kappa_0-3,\quad
 a_5=\kappa-3\kappa_0-2,\quad
 a_6=\kappa-5\kappa_0-3,\\
 H_1=7c_XK_0\exp(3a_5)H_{\mathrm{cl}},\qquad
 z_S=28H_1\exp(\kappa_0)K_0\exp(3a_6).
 \end{gathered}
\]
The fixed small-charge choice
\[
 c_S e^{1/4-5\varepsilon}\le1,\qquad
 c_S=10^5c_X^5K_0^5\exp(6\kappa)
\]
implies
\begin{equation}\label{eq:rgnl-cluster-smallness}
 H_{\mathrm{cl}}\le1,\qquad
 4H_{\mathrm{cl}}K_0c_X\exp(3a_5)\le1,\qquad
                                      z_S\le\tfrac12.
\end{equation}
Indeed the second left side is
$144c_X^3K_0^3\exp(3a_5)e^{1/4-5\varepsilon}$.
For the third,
$2z_S=14112c_X^3K_0^4
\exp(3a_5+3a_6+\kappa_0)e^{1/4-5\varepsilon}$, and
$3a_5+3a_6+\kappa_0\le6\kappa$.

Even elements commute, so expanding their exponential
gives $\exp(\sum_XV(X))=\sum_YK(Y)$.
Moreover $K(Y)$ factors over the connected components
of $Y$: every polymer in its defining set belongs to
exactly one component.
If distinct polymers $X_1,\ldots,X_n$ have connected
union $Y$, their incompatibility graph is connected.
Otherwise two nonempty groups would partition their
union; connectedness of $Y$ would give coincident or
adjacent cubes in different groups and hence an edge
between them.
For each incompatibility-tree edge, choose coincident or
face-adjacent cubes in the two polymers and points
where their nets meet these cubes. The two points
lie in a box of side lengths $M,M,2M$, so they
can be joined by a segment of length at most
$\sqrt6M<3M$. Joining near-minimal nets along
the tree proves
\[
 d_M(Y)\le\sum_i d_M(X_i)+3(n-1).
\]
Using $\|\exp V-1\|\le2\|V\|$ and $a_2=a_5+\kappa_0$
therefore bounds the product by
\[
 \exp(-a_5d_M(Y))
       \big(2H_{\mathrm{cl}}\exp(3a_5)\big)^n
                                  \prod_i \exp(-\kappa_0d_M(X_i)).
\]
Replace distinct sets by ordered tuples with factor $1/n!$,
and use the counting bound for subpolymers of $Y$. The result is
\[
 \|K(Y)\|\le \exp(-a_5d_M(Y))x_Y\exp(x_Y),\qquad
 x_Y=2H_{\mathrm{cl}}K_0\exp(3a_5)|Y|_M.
\]
By \eqref{eq:rgnl-cluster-smallness},
\[
 \begin{gathered}
 x_Y\le(d_M(Y)+1)/2,\\
 x_Y\le4H_{\mathrm{cl}}K_0c_X\exp(3a_5)\exp(d_M(Y)/2).
 \end{gathered}
\]
Since $4\exp(1/2)\le7$ and $a_5-1=a_4$, we obtain
\begin{equation}\label{eq:rgnl-Mayer-bound}
 \|K(Y)\|_{(w;\mathsf w_W)}
                         \le H_1\exp(-a_4d_M(Y))=:\mu(Y).
\end{equation}

Weighted Berezin integration costs no norm here:
each paired moment at level zero is
$\pm48^{|J|/2}$, which is exactly covered by the
weight $(\sqrt{48})^{|J|}$ of its generators.
The corresponding moment is $\pm1$ at positive levels.
The photon measure is a probability measure on its compact
box. Hence $\|H(Y)\|_w\le\mu(Y)$.
These integrals produce even elements, since a nonzero
Berezin moment uses an even number of generators.
For compatible polymers the fluctuation modes and
supported photon coordinates are disjoint. The two
product measures therefore factor, with the external
old generators left as coefficients. It follows that
\begin{equation}\label{eq:rgnl-compatible-partition}
 \int\dd\mu_I(W)\,\dd\widehat\mu(\zeta)e^{\sum_XV(X)}
             =\sum_{\{Y_i\}\text{ compatible}}\prod_iH(Y_i).
\end{equation}

We now bound the infinite connected series. Let
\[
 S_n(Y)=\sum_{\bigcup_iY_i=Y}
                   |u^T(Y_1,\ldots,Y_n)|\prod_i\mu(Y_i).
\]
For $n\ge2$, apply \eqref{eq:rgnl-tree-domination}.
For every contributing tree the union estimate and
$a_4=a_6+2\kappa_0$ give
\[
 \prod_i \exp(-a_4d_M(Y_i))
 \le \exp(-a_6d_M(Y))\exp(3a_6(n-1))
                                  \prod_i \exp(-2\kappa_0d_M(Y_i)).
\]
Drop the prescribed union, retaining $Y_i\subseteq Y$,
and root the tree at $1$. Write $q_j$ for the number
of children at $j$. A polymer incompatible with $Y'$
contains one of at most $7|Y'|_M$ cubes consisting of
its cubes and their face neighbors. Since
$|Z|_M\le c_X(d_M(Z)+1)\le\kappa_0(d_M(Z)+1)$ and
$a^q\le q!\exp(a)$, the counting estimate gives
\begin{align*}
 \sum_{Z\text{ incompatible with }Y'}
      \exp(-2\kappa_0d_M(Z))|Z|_M^{q}
            &\le7q!\exp(\kappa_0)K_0|Y'|_M,\\
 \sum_{Z\subseteq Y}\exp(-2\kappa_0d_M(Z))|Z|_M^q
            &\le q!\exp(\kappa_0)K_0|Y|_M.
\end{align*}
Summing from leaves to root thus costs at most
$(7\exp(\kappa_0)K_0)^n|Y|_M\prod_jq_j!$.
For prescribed children counts, Cayley's degree formula
\cite[Proposition~4]{PostnikovLin2019}
counts
$(n-2)!/((q_1-1)!\prod_{j\ge2}q_j!)$ trees, where
$q_1\ge1$ and $\sum_jq_j=n-1$.
Consequently
\[
 \sum_{\text{trees}}\prod_jq_j!
   =(n-2)!\sum_{(q_j)}q_1
   \le(n-1)!\binom{2n-2}{n-1}
   \le(n-1)!4^{n-1}.
\]
Collecting the factors, also for $n=1$, gives
\begin{equation}\label{eq:rgnl-connected-n-bound}
 \frac{S_n(Y)}{n!}
 \le7H_1\exp(\kappa_0)K_0|Y|_M
                                      \exp(-a_6d_M(Y))z_S^{n-1}.
\end{equation}
Since $z_S\le1/2$, sum the geometric series and use
$|Y|_M\le c_X\exp(d_M(Y))$ and $a_6-1=a_3$.
The bound becomes
\[
 98c_X^2K_0^2\exp(\kappa_0+3a_5)
                                 H_{\mathrm{cl}}\exp(-a_3d_M(Y))
 \le c_S'H_{\mathrm{cl}}\exp(-a_3d_M(Y)),
\]
because $\kappa_0+3a_5=3\kappa-8\kappa_0-6$.
Summing only $n>N$ gives the additional factor $2^{-N}$.
This proves absolute convergence and the claimed estimates.

It remains to justify the logarithm identity after
substitution in \eqref{eq:rgnl-connected-graph-identity}.
The even subalgebra with the coefficient norm is a
commutative Banach algebra, and the polymer set is finite.
The total norm sum of the connected terms is finite by
\eqref{eq:rgnl-connected-n-bound}.
In the exponential of their sum, the norm sum of all
products of connected terms is bounded by the exponential
of that total norm sum. Regrouping is therefore permitted.
Regroup by monomials in the formal variables $z_Y$.
The coefficients are those of
\eqref{eq:rgnl-connected-graph-identity}, namely one
for compatible sets and zero for all other monomials.
Its substitution $z_Y=H(Y)$ gives
\eqref{eq:rgnl-compatible-partition}. Grouping the
connected terms by their union gives
$\sum_YE^\#(Y)$ and proves \eqref{eq:rgnl-log-identity}.
\end{proof}

There are also polynomial approximants to this logarithm
which converge uniformly for all families satisfying the
preceding bound, at fixed side. Replace $\exp V-1$ by
$p_N(V)=\sum_{1\le j\le N}V^j/j!$ in $K,H$, and stop
the connected series at $n=N$; denote the result $E_N^\#$.
It is a polynomial in integrals of products of $V(X_i)$,
$X_i\subseteq Y$, and has the same bound
\eqref{eq:rgnl-log-bound}. To verify uniform convergence,
write $D=|\mathcal P_k|$. Since $\|V(X)\|\le1$,
\[
 \|\exp V(X)-1-p_N(V(X))\|
                                  \le\frac{\exp(1)}{(N+1)!}.
\]
Telescoping each finite product in $K$ gives
\[
 \|H(Y)-H_N(Y)\|\le
                             \frac{4^DD\exp(1)}{(N+1)!}.
\]
Thus this difference is at most $\delta_N\mu(Y)$, where
\[
 \delta_N=\frac{4^DD\exp(1)}
                      {(N+1)!\min_{Y\in\mathcal P_k}\mu(Y)}
                                                   \longrightarrow0.
\]
Replacing the $n$ factors in a connected term costs at
most $n\delta_N\prod_i\mu(Y_i)$. Summing
\eqref{eq:rgnl-connected-n-bound} with this factor, using
$\sum_{n\ge1}n2^{-(n-1)}=4$, and adding the tail yields
\[
 \sup_Y\|E^\#(Y)-E_N^\#(Y)\|
 \le28\delta_NH_1\exp(\kappa_0)K_0(\ell/M)^3
                                  +2^{-N}c_S'H_{\mathrm{cl}}.
\]
This convergence is uniform in the connections and the
allowed data. The displayed error depends on the side;
no volume-independent approximation is claimed.

\subsection{Symmetries of the local logarithm}
\label{subsec:rgnl-symmetries}

Apply Proposition~\ref{prop:rgnl-cluster-logarithm} with
$V=F$, $w=h/4$ and denote its output by
$E_k^\#(m,\mathscr E)$. Each connected term assigned to
$Y$ is a product of integrated activities on subpolymers
of $Y$. Thus ordinary and admissible pair support, local
connection dependence, and dependence only on $m'$ and
$\mathscr E^r(X)$, $X\subseteq Y$, pass to the sum.
The uniform bound and the polynomial approximants prove
analyticity of its coefficients in $(A,\mathscr C)$.

Let $\mathbb G_k^L$ denote the chamber symmetry group with
translations in $LM\Z^3$, so that it preserves every
parameter partition. Its actions and the gauge,
charge-conjugation and reality actions are those of
Section~\ref{subsec:rg-activity-spaces}.

\begin{proposition}\label{prop:rgnl-symmetries}
The local logarithm satisfies
\begin{align}
 a_{\tilde g}E_k^\#(\tilde gY;\tilde gA,\tilde g\mathscr C)
          &=E_k^\#(Y;A,\mathscr C),\qquad \tilde g\in\mathbb G_k^L,
                                      \label{eq:rgnl-lattice-symmetry}\\
 G_{\lambda+\lambda'}E_k^\#(Y;A^1,\mathscr C^1)
          &=E_k^\#(Y;A,\mathscr C),    \label{eq:rgnl-gauge-symmetry}\\
 \mathrm{Cg}\,E_k^\#(Y;-A,-\mathscr C)
          &=E_k^\#(Y;A,\mathscr C).    \label{eq:rgnl-charge-symmetry}
\end{align}
In the gauge identity both pairs belong to the domain and
 $A^1=A-d\lambda$, $\mathscr C^1=\mathscr C-d\lambda'$
on all bonds at distance at most $1$ from $Y$, with
real $\lambda,\lambda'$. Here $d\lambda$ denotes the
lattice gradient, as in the preceding gauge conventions.
For real $m$ and $\mathscr E\in K_k^x$ one also has
\begin{equation}\label{eq:rgnl-reality-symmetry}
 \mathbf K E_k^\#(Y;-\overline A,-\overline{\mathscr C})
                                        =E_k^\#(Y;A,\mathscr C).
\end{equation}
Every monomial is balanced.
\end{proposition}

\begin{proof}
We first record how transformations pass through the
finite operations. Each of the four algebra maps is a
continuous ring automorphism, linear for the first three
and antilinear for $\mathbf K$. If an odd field transforms
like the generators it replaces, the two compositions of
the field substitution with the automorphism agree on
each generator; hence they agree on the algebra.
The same observation applies to adding a field.
Real-coefficient finite differences
\eqref{eq:rgnl-finite-difference} then commute with these
actions. A lattice symmetry also permutes the parameter
cubes and localization summands. It preserves face
connectedness, incidence relations and the Ursell coefficient.

The contour functional commutes with the linear actions.
For the antilinear action, parametrize $t=r\exp(i\theta)$:
\[
 \mathcal C_rf=\frac1{2\pi}\int_0^{2\pi}
             \frac{f(r\exp(i\theta))}{r\exp(i\theta)-1}\,\dd\theta.
\]
Conjugation of coefficients, together with
$t\mapsto\bar t$ and $\theta\mapsto-\theta$, gives the
same integral at the transformed other arguments.
The integration over $W$ commutes with every linear or
antilinear action which sends
$W\mapsto UW$, $\bar W\mapsto(U^{-1})^T\bar W$.
Indeed Wick's rule preserves every paired moment under
this change, and all these moments are real.
This applies to the lattice and gauge actions, to
$\mathbf K$ with $U=\varkappa$, and to charge conjugation
extended by the identity on $W$.
The last extension will require a separate covariance
argument below.

For a lattice symmetry the primitive covariance gives
\[
 Z(\tilde gs')[\tilde g\zeta]=\tilde gZ(s')[\zeta],\qquad
 H_k(\tilde gs'',\tilde gB)\tilde g=\tilde gH_k(s'',B),
 \qquad
 \Gamma_k(\tilde gs'',\tilde gB)\tilde g=\tilde g\Gamma_k(s'',B).
\]
The barred action is $(S^{-1})^T$ on spinors;
$(S^{-1})^T=\pm\mathsf C^{-1}S\mathsf C$ gives the
same covariance for $\bar H_k$.
H\"older quotients commute with the simultaneous
transformation of their endpoints and connection.
Thus $J_k^\#\Psi$ and $W_k^\#$ transform like their
ordinary and pair generators.
The quadratic primitive satisfies
$M_k(\tilde gX,\tilde gB)=\tilde gM_k(X,B)\tilde g^{-1}$,
so its bilinear form is invariant.
The mass form is invariant by the change of chamber
site and spin indices, and $R\mathscr E$ has the
incoming lattice symmetry. The reconstruction substitution
commutes with the action by
\eqref{eq:rg-step-substitution-covariances}.
The box measure is invariant under $\zeta\mapsto\tilde g\zeta$.
Apply these identities to $G_1,G_2,G_3$, the first
localization and substitution, then to the second
localization, $K,H$, and the convergent connected sum.
They prove \eqref{eq:rgnl-lattice-symmetry}.

For the gauge identity set $\Lambda=\lambda+\lambda'$.
In a summand assigned to $Y$, all parameters outside $Y$
vanish. The local row result proves that $B^1$ may be
replaced by $B-d\Lambda$ in every row used in that
summand. The latter field is still in the operator domain,
since real gauge shifts preserve it.
Telescoping the gradient along each path gives
\[
 U_{B-d\Lambda}(x,x')
       =U_B(x,x')\exp\big(-ie(\Lambda(c')-\Lambda(c))\big).
\]
The ordinary primitive row obeys
\[
 H_k(s'',B^1)(x,\cdot)
       =\exp(ie\Lambda(c))H_k(s'',B)(x,\cdot)
                                                 \exp(-ie\Lambda^u),
\]
when it is nonzero, where $\Lambda^u$ is the function
on the unit indices. The barred row has the inverse
phases. In the H\"older row the phase at its second
point cancels that from the transport, leaving exactly
$\exp(ie\Lambda(c))$, the phase of the first point.
The row of $\Gamma_k$ has the corresponding two-sided
gauge covariance. Its composition with $H_k^\#$ thus
gives the same covariance for $W_k^\#$, including the
rows which vanish outside $Y$.
The primitive bilinear and the local mass form are
gauge invariant, and so is $R\mathscr E$ under the
separate shifts of its two slots. The reconstruction
substitution also commutes with this local gauge action.
The photon fields and photon measure are unchanged.
The finite operations and integration therefore prove
\eqref{eq:rgnl-gauge-symmetry}.

For reality, the photon primitives are real on real
parameter families. Their analytic dependence extends
this identity to
$\overline{O(\bar s')}=O(s')$ on the full polydisk.
Since $C$ is real,
\[
 Z(\bar s')[-\zeta]=-\overline{Z(s')[\zeta]}.
\]
Transform $s\mapsto\bar s$, $t\mapsto\bar t$,
$\zeta\mapsto-\zeta$ and
$(A,\mathscr C)\mapsto(-\bar A,-\overline{\mathscr C})$.
The fermion identities are
\[
 \overline{H_k(\bar s'',-\bar B)}
                         =\varkappa^{-1}H_k(s'',B)\varkappa,\qquad
 \overline{\Gamma_k(\bar s'',-\bar B)}
                         =\varkappa^{-1}\Gamma_k(s'',B)\varkappa.
\]
For the barred field use
$\varkappa\mathsf C=-\mathsf C\varkappa$.
The transport relation
$\overline{U_{-\bar B}(x,x')}=U_B(x,x')$
then gives the same identities for the H\"older rows.
The local quadratic primitive obeys the corresponding
identity, and $\varkappa^T\varkappa=1$ preserves the
mass form. The number $m'$ is real for the present data,
and $R\mathscr E$ and the reconstruction substitution
have the reality symmetry. The contour argument just
given, invariance of the box under sign reversal, and
weighted integration prove
\eqref{eq:rgnl-reality-symmetry}.

Charge conjugation gives the first-localization symmetry
by
\[
 \begin{gathered}
 H_k(s'',-B)\mathsf C=\mathsf C\bar H_k(s'',B),\qquad
 \bar H_k(s'',-B)\mathsf C=\mathsf C H_k(s'',B),\\
 M_k(X,-B)^T=\mathsf C^{-1}M_k(X,B)\mathsf C .
 \end{gathered}
\]
Anticommuting the two transformed legs verifies the
bilinear and mass-form identities. The first-step
vertex has this symmetry by its explicit two-half formula.
It follows that
\[
 \mathrm{Cg}\,E_{\le3}(X;-A,-\mathscr C,-\zeta)
                                        =E_{\le3}(X;A,\mathscr C,\zeta).
\]
Direct covariance of the second substitution is not
available, since its unbarred and barred fluctuations
use different factors. We instead compare their Wick
covariances.

Introduce independent generators $\phi^\#$ indexed like
$\psi^\#$, with their own charge-conjugation action,
and let $\operatorname{ad}$ add them to the old generators.
Put
\begin{align*}
 \mathcal V(X)
   ={}&\operatorname{ad}E_{\le3}(X)
      +\operatorname{ad}\mathscr E^r(X)-\mathscr E^r(X)\\
     &-\mathbf1_{X=\square}m'
          \big(\operatorname{ad}
                         \langle\bar\psi,\psi\rangle_\square
                                -\langle\bar\psi,\psi\rangle_\square\big).
\end{align*}
It has charge-conjugation symmetry under
$(A,\mathscr C,\zeta)\mapsto(-A,-\mathscr C,-\zeta)$.
Let $I_{s'',B}$ fix the old generators and replace
$\phi^\#$ by $W_k^\#(s'',B)$. The second-localization
input is exactly $I_{s'',B}\mathcal V(X)$, since
the polynomial contours in $G_4,G_5$ give their value
at $1$ minus their value at $0$.
For any union $U$, interpolation restricted to cubes
of $U$, with the disconnected terms removed, gives
\begin{equation}\label{eq:rgnl-subvolume-substitution}
 \sum_{Y\subseteq U}F(Y)=
                        I_{\mathbf1_U,B}\sum_{X\subseteq U}\mathcal V(X).
\end{equation}

At $s''=\mathbf1_U$ define
\[
 \phi_\pm=H_k(\mathbf1_U,\pm B)\Gamma_k(\mathbf1_U,\pm B)W,
 \qquad
 \bar\phi_\pm=\bar H_k(\mathbf1_U,\pm B)\bar W,
 \qquad
 \phi'=\mathsf C\bar\phi_-,\quad
 \bar\phi'=\mathsf C^{-1}\phi_-.
\]
The substitution by $(\phi',\bar\phi')$ and their
H\"older quotients, denoted $I'$, satisfies
$\mathrm{Cg}\,I_{\mathbf1_U,-B}=I'\mathrm{Cg}$
on every generator, when $\mathrm{Cg}$ fixes $W$.
The only nonzero ordinary pair expectations for the
positive substitution are the kernel of
\[
 \mathcal G_U(B)
        =H_k(\mathbf1_U,B)\Gamma_k(\mathbf1_U,B)
                                            \bar H_k(\mathbf1_U,B)^T.
\]
The identity moment has kernel
$\mathbf1_{u=u'}/\nu(u)$, with $\nu=1/48$ at zero.
This precisely cancels one of the two measures in
the kernel composition, so the same formula holds
at every level. The primitive transpose relation gives
\[
 \mathcal G_U(-B)^T
                      =-\mathsf C^{-1}\mathcal G_U(B)\mathsf C^{-1}.
\]
Anticommuting $\bar\phi_-$ and $\phi_-$ therefore gives
\[
 \int\dd\mu_I(W)\,\phi'(x)\bar\phi'(x')^T
     =\mathsf C[-\mathcal G_U(-B)^T(x,x')]\mathsf C
     =\mathcal G_U(B)(x,x').
\]
Same-bar pair expectations vanish in both substitutions.
The pair generators are obtained by the identical
linear H\"older maps at $B,-B$, so their pair expectations
also agree. Wick's theorem now proves
$\int\dd\mu_I(W)I'Q=\int\dd\mu_I(W)I_{\mathbf1_U,B}Q$
for every polynomial $Q$ in the old and $\phi^\#$
generators: each monomial integral is the Pfaffian of
its pair expectations, with old generators as coefficients.

Apply this equality to the exponential of
$\sum_{X\subseteq U}\mathcal V(X)$ and use
$\zeta\mapsto-\zeta$. By
\eqref{eq:rgnl-subvolume-substitution}, the integrated
subvolume expression
$\Xi_U=\int\dd\mu_I(W)\dd\widehat\mu(\zeta)
\exp(\sum_{Y\subseteq U}F(Y))$
has charge-conjugation symmetry.
But $\Xi_U=\sum_{Z\subseteq U}H(Z)$, by the finite
Mayer expansion. Inverting this finite subset sum gives
\[
 H(Z)=\sum_{U\subseteq Z}(-1)^{|Z\setminus U|_M}\Xi_U.
\]
Thus each $H(Z)$ has the symmetry, and the real
coefficients and absolute convergence of the connected
series prove \eqref{eq:rgnl-charge-symmetry}.

Finally take a constant $\lambda$ and $\lambda'=0$
in \eqref{eq:rgnl-gauge-symmetry}. A monomial with
$n_u$ unbarred and $n_b$ barred legs acquires the factor
$\exp(ie\lambda(n_u-n_b))$. Invariance for every real
$\lambda$ and independence of the basis monomials force
its coefficient to vanish unless $n_u=n_b$.
\end{proof}

\subsection{The scaled correction}
\label{subsec:rgnl-scaled-correction}

Let $S_L$ be the field scaling of
\eqref{eq:rg-leg-scaling}. The nonlinear correction is
\begin{equation}\label{eq:rgnl-scaled-correction}
 E_k^*(m,\mathscr E)(Y;\mathscr A,\mathscr C)
  =\sum_{\bar X^L=LY}S_L
        E_k^\#(m,\mathscr E)(X;\mathscr A_L,\mathscr C_L).
\end{equation}
This formula uses field scaling and summation over fine
polymers. Quadratic adjustment enters the incoming
activity term in the full step below.
The domain inclusions
\eqref{eq:rg-scaled-domain-inclusion} ensure
$(\mathscr A_L,\mathscr C_L)
\in\widetilde R_k^-\times\mathfrak X_k$.

\begin{proposition}\label{prop:rgnl-scaled-correction}
For $(m,\mathscr E)\in\operatorname{Ball}_k^x(4)$,
\[
 E_k^*(m,\mathscr E)\in K_{k+1}^{x,\C},\qquad
                  \|E_k^*(m,\mathscr E)\|_{k+1}^x
                                      \le e^{1/4-6\varepsilon}.
\]
For real data its value belongs to $K_{k+1}^x$.
On any affine complex line in the data space its value
is norm continuous on the part in the closed ball and
norm analytic in its interior.
For two data in $\operatorname{Ball}_k^x(1)$,
\begin{equation}\label{eq:rgnl-coupling-Lipschitz}
 \|E_k^*(m_1,\mathscr E_1)-E_k^*(m_2,\mathscr E_2)\|_{k+1}^x
 \le e^{1/4-6\varepsilon}
       \left(\frac{|m_1-m_2|}{w_k^{\mathrm m}}
              +\frac{\|\mathscr E_1-\mathscr E_2\|_k^x}
                                                    {w_k^{\mathrm E}}\right).
\end{equation}
The second quotient is absent at $k=0$.
Restriction to $\mathscr C=0$ gives the correction formed
without an external slot from $(m,\mathscr E|_0)$.
\end{proposition}

\begin{proof}
The scaling identity gives
\[
 \|S_LQ\|_{h_{k+1}}
   =\|Q\|_{(L^{-1}h_{k+1,1},L^{-1-\alpha}h_{k+1,2})}
   \le\|Q\|_{h_k/4},
\]
since $h_{k+1,1}=L^{-1/8}h_{k,1}$,
$h_{k+1,2}=L^{-1/8+\varepsilon/2}h_{k,2}$ and $L\ge4$.
The polymer sum in \eqref{eq:rgnl-scaled-correction}
uses $Ld_M(Y)\le d_M(X)$ and $X\subseteq LY$.
Splitting $a_3=(a_3-\kappa_0)+\kappa_0$ therefore yields
\begin{align*}
 \sum_{\bar X^L=LY}\exp(-a_3d_M(X))
  &\le \exp(-L(a_3-\kappa_0)d_M(Y))
                         \sum_{X\subseteq LY}\exp(-\kappa_0d_M(X))\\
  &\le c_XK_0L^3
                  \exp(-[L(\kappa-6\kappa_0-4)-1]d_M(Y))
   \le c_XK_0L^3\exp(-\kappa d_M(Y)).
\end{align*}
Together with \eqref{eq:rgnl-log-bound} this gives
\[
 \|E_k^*\|_{k+1}^x
  \le c_XK_0L^3c_S'H_{\mathrm{cl}}
  =504c_X^5K_0^5\exp(3\kappa)L^3e^{1/4-5\varepsilon}
  \le e^{1/4-6\varepsilon},
\]
using $c_SL^3e^\varepsilon\le1$.

All generators of a summand lie over $X$ or at a pair
admissible for $X$. Scaling maps the ordinary sites
into $Y$; it maps each of the six paths into the
corresponding scaled path and maps
$(x,x^{(3)})$ to $(x/L,(x/L)^{(3)})$.
These pairs are admissible for $Y$, including
intracluster pairs. A bond meeting $X$ scales to
one meeting $Y$, giving local connection dependence.
The scaled-domain inclusion and coefficient analyticity
give analytic kernels.
The map $(S,a)\mapsto(S,a/L)$ identifies
$\mathbb G_k^L$ with $\mathbb G_{k+1}$ and intertwines
its generator actions with $S_L$. For gauge transformations
use $\Lambda_L(c)=L^{1/2}\Lambda(c/L)$ and
$(d\Lambda)_L=d(\Lambda_L)$.
A gauge relation on the unit neighborhood of $Y$
then holds on the unit neighborhood of every
$X\subseteq LY$. Charge conjugation and reality also
commute with the real scaling factors.
Proposition~\ref{prop:rgnl-symmetries} thus gives all the
defining symmetries of $K_{k+1}^{x,\C}$, and the real
space assertion follows from its reality clause.

For norm analyticity, take data
$(m_a+t m_b,\mathscr E_a+t\mathscr E_b)$ in the closed
ball and denote their parameter set by $T$.
It is closed and convex. The finite activity $F_t(X)$
is affine in $t$ by Proposition~\ref{prop:rgnl-finite-bounds}.
The polynomial approximants to $E_k^\#$ are therefore
polynomials in $t$. Reblock them by
\eqref{eq:rgnl-scaled-correction}; the resulting polynomials
$f_N(t)$ converge uniformly on $T$ in the full activity
norm, since both finite polymer sets are fixed and the
error bound of the previous subsection is uniform.
Their values have bounded full norm.
If $T$ has more than one point, it contains infinitely
many points. The coefficients of each polynomial are
finite numerical linear combinations of its values at
distinct points of $T$, by polynomial interpolation.
Hence these coefficients have finite full norm, and the
uniform limit is norm continuous on $T$.
If $T$ is empty or a singleton there is no continuity
or interior analyticity assertion to check.

For a closed disk $|t-t_0|\le r$ in $T$, take its
coefficientwise Cauchy integrals
\[
 f_j=\frac1{2\pi i}\int_{|t-t_0|=r}
                             f(t)(t-t_0)^{-j-1}\,\dd t.
\]
They have analytic connection kernels, and the support,
symmetry and balance conditions pass through this
complex linear integral. Moreover
$\|f_j\|_{k+1}^x\le e^{1/4-6\varepsilon}r^{-j}$.
Coefficientwise Cauchy's formula, with this summable
norm bound, gives
$f(t)=\sum_{j\ge0}f_j(t-t_0)^j$ in activity norm for
$|t-t_0|<r$. This proves the stated analyticity.

To obtain \eqref{eq:rgnl-coupling-Lipschitz}, vary first
the mass, keeping $\mathscr E_1$ fixed. With
$\Delta m=m_1-m_2\ne0$, the line
$(m_2+t\Delta m,\mathscr E_1)$ remains in the large
ball for $|t|\le r=3w_k^{\mathrm m}/|\Delta m|$.
Here $r\ge3/2$. The contour formula on any circle
$1<r'<r$ gives
\[
 \|f(1)-f(0)\|_{k+1}^x
        \le\frac{e^{1/4-6\varepsilon}}{r'-1}.
\]
Let $r'\uparrow r$ and use $r-1\ge r/3$.
The bound is the first quotient on the right of
\eqref{eq:rgnl-coupling-Lipschitz}.
Vary the activity with mass $m_2$ fixed and radius
$3w_k^{\mathrm E}/\|\mathscr E_1-\mathscr E_2\|_k^x$
to obtain the second quotient.
The zero differences require no contour.

Finally restriction commutes with extraction:
$m(\mathscr E)=m(\mathscr E|_0)$ and
$(R\mathscr E)|_0=R(\mathscr E|_0)$.
Every fermion operator uses the sum of its first
connection and $\mathscr C$; photon operations and
measures do not use $\mathscr C$.
At $\mathscr C=0$ the two constructions therefore have
identical finite expressions, connected series and
scaled corrections. Analyticity without a slot is
understood in the first connection alone.
\end{proof}

The possible H\"older legs can also be read directly.
In $G_2,G_4$ an existing pair generator is changed by
addition of a fluctuation field; this creates no new
pair index. The mass contours contain no pair
generators. The reconstruction substitution and the
first-step vertex introduce only pairs $(x,x^{(3)})$.
Every later sum, product and integral preserves this
set of available indices. Thus the correction has
only scaled incoming pair indices and these scaled
type-$3$ indices. At $k=0$ it has only the latter,
of length $\eta_1$, since the incoming activity is zero.

\subsection{The step identity and bounds}
\label{subsec:rgnl-current-step}

For the local determinant terms of
Proposition~\ref{prop:local-determinant-expansion}, put
\[
 \delta E_k^{\mathrm{det}}(X;\mathscr A,\mathscr C)
       =E_k^d(X,\mathscr A+\mathscr C)-E_k^d(X,\mathscr C),
 \qquad
 E_k^{\mathrm{det}}(X,\mathscr C)
       =E_k^d(X,\mathscr C)-E_k^d(X,0).
\]
These are scalar activities. We use reblocking
$\mathcal L$ from \eqref{eq:rg-reblocking-definition}.
The adjustment acts on the incoming extracted activity.
Define
\begin{align}
 \mathscr E^+&=\mathcal L(R\mathscr E)^{\mathrm{nat}}
                        +\mathcal L\delta E_k^{\mathrm{det}}
                        +E_k^*(m,\mathscr E), \label{eq:rgnl-full-step}\\
 m^+&=L(m+m(\mathscr E)),\\
 \varepsilon^{\mathrm v,+}
       &=L^3(\varepsilon^{\mathrm v}
                  +\varepsilon_{\mathrm v}(\mathscr E)
                                      +\delta\varepsilon_k^0),\\
 \mathfrak P_{k+1}(Y,\mathscr C')
       &=\sum_{\bar X^L=LY}
           \big(\mathfrak P_k(X,\mathscr C'_L)
                          +E_k^{\mathrm{det}}(X,\mathscr C'_L)\big).
                                      \label{eq:rgnl-external-scalar-map}
\end{align}
Here $\mathfrak P_k$ is any given family of scalar
functions of the external slot.

The reblocked determinant term belongs to $K_{k+1}^x$.
Indeed the scaled arguments lie in its inverse domain,
and each summand is analytic and reads only fine
bonds meeting $X$. These scale to bonds meeting $Y$.
There are no legs. The lattice symmetry permutes the
reblocking summands; local gauge shifts of the two
slots become the corresponding real shifts on these
fine read bonds, where determinant locality and gauge
invariance apply. Evenness of $E_k^d$ gives charge
conjugation, and its conjugation identity gives reality.
The pointwise bound
\eqref{eq:determinant-small-field-gain} and
\eqref{eq:polymer-reblock-sum} show
\begin{align}
 \|\mathcal L\delta E_k^{\mathrm{det}}\|_{k+1}^x
    &\le8c_{\mathrm{rb}}L^3C_dM^3e_k^{1/4}
      \le e_k^{1/4-2\varepsilon}
      \le e_{k+1}^{1/4-2\varepsilon}.
                                      \label{eq:rgnl-reblocked-determinant}
\end{align}
Here $L(\kappa-\kappa_0-1)\ge\kappa$ controls the
remaining exponential weight, and the fixed charge
condition is
$8c_{\mathrm{rb}}L^3C_dM^3e_k^{2\varepsilon}\le1$.
Restriction to $\mathscr C=0$ gives the same result
for $\mathcal LE_k^{\mathrm{det}}$.

\begin{theorem}\label{thm:rgnl-current-step}
Let $\varepsilon^{\mathrm v},m$ be real,
$\mathscr E\in K_k^x$, and
$(m,\mathscr E)\in\operatorname{Ball}_k^x(1)$, with
$\mathscr E=0$ at $k=0$. For $A'\in\mathrm{Dom}_{k+1}$
with arbitrary harmonic part and
$\mathscr C'\in\mathfrak X_{k+1}$, the exact density satisfies
\begin{equation}\label{eq:rgnl-exact-step}
 \mathbf T_k^x\mathscr F_k^x[\varepsilon^{\mathrm v},m,\mathscr E]
   =\mathscr F_{k+1}^x
            [\varepsilon^{\mathrm v,+},m^+,\mathscr E^+],
\end{equation}
with external scalar families $\mathfrak P_k$ and
$\mathfrak P_{k+1}$ on the respective sides.
Its integral satisfies
\begin{equation}\label{eq:rgnl-integrated-step}
 y_{k+1}^x(A';\mathscr C')
  =\int_{\mathrm{Box}_k}\dd\mu_I(\zeta)\,
        y_k^x(A_k^{\min}(A')+\mathrm{Sq}_k\zeta;\mathscr C'_L),
\end{equation}
where $y_j^x=\int\mathrm D\Psi\,\mathscr F_j^x$ with the
data and scalar family of that level.
Moreover $\mathscr E^+\in K_{k+1}^x$ and
\begin{align}
 \|\mathscr E^+\|_{k+1}^x
 &\le c_{\mathrm{fl}}L^{-1/4+2\varepsilon}\|\mathscr E\|_k^x
           +e_{k+1}^{1/4-2\varepsilon}+e_k^{1/4-6\varepsilon},
                                      \label{eq:rgnl-activity-step-bound}\\
 |m^+-Lm|
 &\le Lc_{\mathrm{fl}}e_k^{1/2}\|\mathscr E\|_k^x,
                                      \label{eq:rgnl-mass-step-bound}\\
 |\varepsilon^{\mathrm v,+}-L^3\varepsilon^{\mathrm v}|
 &\le L^3(c_{\mathrm{fl}}\|\mathscr E\|_k^x+e_k^7).
                                      \label{eq:rgnl-vacuum-step-bound}
\end{align}
At level zero this formula integrates all $192n_0$
incoming pairs of generators, with chamber measure $1/48$,
and leaves $4n_1$ outgoing pairs. Its data are
\[
 \mathscr E^+=\mathcal L\delta E_0^{\mathrm{det}}+E_0^*(m,0),
 \quad m^+=Lm,\quad
 \varepsilon^{\mathrm v,+}
                         =L^3(\varepsilon^{\mathrm v}+\delta\varepsilon_0^0).
\]
In particular
$0\le\varepsilon^{\mathrm v,+}-L^3\varepsilon^{\mathrm v}
\le L^3q^7$.
\end{theorem}

\begin{proof}
Use the notation of Proposition~\ref{prop:rg-step-density-identity},
with $A=\mathscr A'_L$, $B=A+\mathscr C$ and
$Z=Z(1)$ the full-strength photon fluctuation.
The primitive families satisfy their exact identities
at full strength:
$H_k(1,B)=H_k(B)$, $\Gamma_k(1,B)=\Gamma_k(B)$, and
$\sum_XM_k(X,B)=b_k-D_k(B)$ for $k\ge1$.
The induced fields obey
\[
 \psi^{0,\#}=H_k^\#(B)\Psi^c,\qquad
 H_k^\#(B)\Psi^f=\psi^{0,\#}+W_k^\#(1,B).
\]
Define three evaluation homomorphisms. Let $\operatorname{ev}$
send old generators to $\psi^{0,\#}$ and fix $W$;
let $\operatorname{ev}_1$ send them to $H_k^\#(B)\Psi^f$;
for $k\ge1$, let $\operatorname{ev}_f$ additionally send
$\Psi,\bar\Psi$ to $\Psi^f,\bar\Psi^f$.
At zero take $\operatorname{ev}_f=\operatorname{ev}_1$.
These maps exist since their images are odd and have
the required antiperiodicity. They commute with the
finite coefficient integrals; $\operatorname{ev}$ also
commutes with Berezin integration, since it fixes $W$.

On the first contours $\mathcal C_rf=f(1)-f(0)$.
For $G_1$, summing the local primitives gives the
Dirac difference
$-\langle\bar\Psi,(D_k(B+Z)-D_k(B))\Psi\rangle_1$.
For $G_2,G_3$, evaluating
$\psi^\#+J_k^\#(1,B,Z)\Psi$ gives
$H_k^\#(B+Z)\Psi^f$.
The evaluations of their contour differences therefore give
\begin{equation}\label{eq:rgnl-first-contour-evaluation}
 \operatorname{ev}_f\sum_X\sum_{i\le3}G_i(1,X)
  =\mathrm{Act}_k^x[m',\mathscr E^r](A+Z,\mathscr C;\Psi^f)
       -\mathrm{Act}_k^x[m',\mathscr E^r](A,\mathscr C;\Psi^f).
\end{equation}
At zero the explicit first vertex gives this same
identity with the chamber pairing: $D_0(B)=D_B+\mu$
and the mass form has no connection dependence.

The reconstruction left inverse gives, on every local
first expression,
$\operatorname{ev}_1\operatorname{sub}_{Y,B}
=\operatorname{ev}_f$.
This follows on the $\Psi$ generators from $T_k(B)H_k(B)=1$
and on the old generators from their definition.
Adding $W_k^\#(1,B)$ to $\psi^\#$ in a polynomial and then
applying $\operatorname{ev}$ gives
$\operatorname{ev}_1$ applied to that polynomial.
Apply these identities to
\eqref{eq:rgnl-two-localization-sums}.
The evaluated $G_{\le3}^{\,W}$ sum is
\eqref{eq:rgnl-first-contour-evaluation}.
The evaluated polynomial contours $G_4,G_5$ give
\[
 I_k^x[m',\mathscr E^r](A,\mathscr C;\Psi^f)
                      -I_k^x[m',\mathscr E^r](A,\mathscr C;\Psi^c).
\]
Consequently
\[
 \operatorname{ev}\sum_YF(Y)=V_k^x,\qquad
 \Xi_k^x
     =\exp\left\{\sum_X
                  \operatorname{ev}E_k^\#(m,\mathscr E)(X;A,\mathscr C)\right\}.
\]
The second equality is
\eqref{eq:rgnl-log-identity} evaluated by
$\operatorname{ev}$; its evaluation commutes with the
integrals as just noted.

The other remaining factor is the determinant quotient.
Equation~\eqref{eq:local-determinant-quotient} gives
\[
 \frac{\delta Z_k(B)}{\delta Z_k(0)}
      =\exp\left\{\sum_X\delta E_k^{\mathrm{det}}(X;A,\mathscr C)
                               +\sum_XE_k^{\mathrm{det}}(X,\mathscr C)\right\}.
\]
On inserted fields adjustment changes no value:
$\operatorname{ev}R\mathscr E(X)
=\operatorname{ev}(R\mathscr E)^{\mathrm{nat}}(X)$.
The equality is the telescoping reconstruction identity
for a small polymer and the definition for a large one.
Finally $\operatorname{ev}=\operatorname{ev}'S_L$,
where $\operatorname{ev}'$ inserts the outgoing fields.
Each fine polymer has exactly one outgoing reblocking
polymer. Thus the three activity sums combine into the
outgoing $\mathscr E^+$, and the two scalar external sums
combine into \eqref{eq:rgnl-external-scalar-map}.
Insert these identities in
\eqref{eq:rg-step-density-identity}. All exponents are
even and commute, so their sums give
\eqref{eq:rgnl-exact-step}, with exactly its previously
derived mass, vacuum and determinant normalization.

Integrate over the outgoing generators and apply the exact
marginal \eqref{eq:rg-step-block-marginal}. This gives
\eqref{eq:rgnl-integrated-step}.
The integrand is admissible throughout the photon box:
Proposition~\ref{prop:photon-domain-preservation} places
$A_k^{\min}(A')+\mathrm{Sq}_k\zeta$ in
$\mathrm{Dom}_k$, and its reconstructed connection is
an affine, continuous function of $\zeta$.
The density coefficients are analytic in this connection
on its operator and activity domains, hence continuous
and integrable on the compact box.
The external scalar factor is fixed during this integral.
No condition on the harmonic part of $A'$ enters.

The three terms of \eqref{eq:rgnl-full-step} belong to
$K_{k+1}^x$ by Proposition~\ref{prop:rg-activity-contraction},
\eqref{eq:rgnl-reblocked-determinant}, and
Proposition~\ref{prop:rgnl-scaled-correction}.
Their respective norms are bounded by the three terms
in \eqref{eq:rgnl-activity-step-bound}.
Equation~\eqref{eq:rg-contraction-norm} bounds the
first term with the stronger exponent $3\varepsilon/2$.
The mass and vacuum extraction bounds give
\eqref{eq:rgnl-mass-step-bound} and
\eqref{eq:rgnl-vacuum-step-bound}, since
\[
 0\le\delta\varepsilon_k^0
        \le12\exp(-p_{0,k}^2/2)\le e_k^7.
\]
At zero all incoming extracted terms vanish.
The exact algebra uses all $192n_0$ pairs with the
weight $1/48$ and integrates to the $4n_1$-pair
outgoing density. There is no removed multiplicity
or additional vacuum normalization. The displayed
zero-level formulas and bounds follow.
\end{proof}

The restriction to zero external connection commutes
with all three terms of \eqref{eq:rgnl-full-step}.
The determinant difference there becomes
$E_k^d(X,\mathscr A)-E_k^d(X,0)$, and the external scalar
map preserves $\mathfrak P_k(\cdot,0)=0$.
Thus \eqref{eq:rgnl-exact-step} and
\eqref{eq:rgnl-integrated-step} also give the step without
a slot, with incoming activity $\mathscr E|_0$.

The displayed norm estimates depend on the running charge
and the fixed constants, and on no side or harmonic part.
They were proved for any primitive system with the stated hypotheses on
the present torus. Identifying the outputs from different
systems or different volumes requires the separate
comparison arguments; it does not follow from these bounds.


\section{Covers and comparison of local steps}\label{sec:rg-cover-comparison}

We construct the primitive systems on one large torus and
project them to smaller tori through odd covering maps. Projection preserves the
estimates with a fixed loss in the common constant and
one loss in polymer decay. The comparison of local steps
then uses supported rows near a common lift, where the
signed image sum has only one surviving term.

\subsection{Odd covers and projected kernels}
\label{subsec:rgcov-projection}

Fix $k\ge0$, $0<\mu\le1$, $0<e_k\le e_0$, and admissible
sides $\ell_2=J_c\ell_1$, where $J_c$ is odd.
Write $T_i=\mathbb L_k/\ell_i\Z^3$ and
$p:T_2\to T_1$ for the cover. Let
$\operatorname{pr}_i$ denote the covering-lattice projections.
Pullback $p^*$ is inclusion of functions on the covering
index sets. They have character
$\chi(w)=1$ for photon fields, coordinates and parameters,
and $\chi(w)=(-1)^{w_1+w_2+w_3}$ for spinors and pair
fields. Oddness gives
$\chi(J_cw)=\chi(w)$, so a field on $T_1$ is a field
on $T_2$ with the same character.
The image is
\[
 \mathrm{Per}=\{f:f(t+\ell_1w)=\chi(w)f(t),\ w\in\Z^3\}.
\]
Translations move the photon site and preserve the chamber
label. In particular pullback commutes with $J,J^*$,
complex conjugation and constant spin matrices.

The index weight is $1$ on unit sites, unit bonds and
photon coordinate indices, $\eta_k^3$ on fine photon
sites and bonds, and $\eta_k^3/48$ on chamber sites.
For fields with a common index set and character,
\begin{equation}\label{eq:rgcov-pullback-pairing}
 \|p^*f\|_\infty=\|f\|_\infty,\qquad
 \langle p^*f',p^*f\rangle_{T_2}
                                      =J_c^3\langle f',f\rangle_{T_1}.
\end{equation}
Indeed a fundamental set for $T_2$ is the union of
$\mathrm{Fnd}_1+\ell_1w$, over representatives
$w\in(\Z/J_c\Z)^3$. The two characters cancel in
each paired summand, and the weight is the same on both
tori. At zero the chamber weight remains $1/48$.

Pullback maps the parameter polydisk, including its interior,
into the corresponding larger polydisk, since
$(p^*s)_{\mathrm{pc}}=s_{p\,\mathrm{pc}}$.
It also maps
$\widetilde{\mathcal G}_k(T_1)$ into
$\widetilde{\mathcal G}_k(T_2)$.
To check the latter assertion, write the domain decomposition
$A=A_0+A_1$. Pullback preserves reality of $A_0$ and
the sup norm of $A_1$. Each five-cube neighborhood has
a common nonwrapping lift of side $11M<\ell_1$.
The real local gauge representative of $A_0$ on its
image in $T_1$ pulls back on this lift to the
corresponding neighborhood in $T_2$, with its same
sup bound and gradient identity. This is precisely
the local real-gauge condition in the inverse domain.

Every exact finite fermion operator used below intertwines
with pullback. For the chamber Dirac operator this follows
term by term from its four neighbors and bond phases.
For $\mathcal Q_k=Q_kJ^*$ and its reverse it follows from
the same block sums and the same six paths on the cover;
$J^*$ averages the same $48$ labels.
The penalty and the finite matrices multiplying the
resolvent have the same finite block formulas.
Sums and products preserve intertwining.
If $O_2p^*=p^*O_1$ and both operators are invertible,
then
$O_2p^*O_1^{-1}=p^*$ and hence
$O_2^{-1}p^*=p^*O_1^{-1}$.
The domain preservation just proved makes the already
established inverse bounds applicable on both tori.
It follows, from the formulas in Section~\ref{sec:fermion-primitives},
that the fine resolvents, $H_k,\bar H_k,D_k,M_k,\Gamma_k$
intertwine at the same connection:
\begin{equation}\label{eq:rgcov-exact-fermion-intertwining}
 O_2(p^*A)p^*=p^*O_1(A).
\end{equation}
At zero use the entire chamber Dirac stencil and the
inverse of $F_y$, on all chamber components.
The exact photon operators have the corresponding
identity \eqref{eq:photon-exact-cover-consistency}.

Let a primitive system on $T_2$ be given.
For a parameter family $O_2(s,A)$ put
$\widetilde O(s,A)=O_2(p^*s,p^*A)$.
For a polymer family put
\[
 \widetilde O(X,A)=\sum_{X'\in\mathcal P_k(T_2):pX'=X}
                                                    O_2(X',p^*A).
\]
This is applied to $M_k$ for $k\ge1$ and to $S_{k,y}$
at every level. A projected polymer is connected:
adjacent $M$-cubes have equal or adjacent images.
Moreover
\begin{equation}\label{eq:rgcov-projected-polymer-length}
                           d_M(pX')\le d_M(X').
\end{equation}
A projected net meets all image cubes. Projecting its chosen
finite segment representation preserves the lengths of the segments;
taking the infimum gives \eqref{eq:rgcov-projected-polymer-length}.
There are only finitely many polymers on either torus.

\begin{proposition}\label{prop:rgcov-kernel-projection}
If the system has translation covariance under
$\ell_1\Z^3$, its projection is well defined by
$p^*O_1=\widetilde Op^*$ and has kernel
\begin{equation}\label{eq:rgcov-signed-image-kernel}
 O_1(\cdot)(t,t')
  =\sum_{w\in(\Z/J_c\Z)^3}\chi(w)
                       \widetilde O(\cdot)(t,t'+\ell_1w).
\end{equation}
This formula uses the same index weights on both tori.
The signed summand depends only on the class of $w$.
Transposition of an intertwining operator also intertwines:
\begin{equation}\label{eq:rgcov-transpose-intertwining}
                          O_2^Tp^*=p^*O_1^T.
\end{equation}
\end{proposition}

\begin{proof}
The pulled-back parameters and connection are invariant
under $\ell_1$ translations. Parameter-family covariance
therefore says that $\widetilde O$ commutes with those
translations. For polymer families, translation permutes
the polymers $X'$ with $pX'=X$, so the same commutation
holds after their finite sum.
An operator commuting with these translations maps the
character eigenspace $\mathrm{Per}$ into itself.
Since $p^*$ is injective onto that space, this defines
exactly one projected operator.

Sum its action over the fundamental sets used in
\eqref{eq:rgcov-pullback-pairing} and insert
$f(t'+\ell_1w)=\chi(w)f(t')$.
The coefficient of $\nu(t')f(t')$ is
\eqref{eq:rgcov-signed-image-kernel}.
Replacing $w$ by $w+J_cv$ multiplies the character
by $\chi(v)$ and the large-torus kernel by the same
character, so the product is unchanged.
This proves the asserted kernel formula, including
the weights $\eta_k^3/48$ or $1/48$ on chamber columns.

Translations have transpose $T_v^T=T_{-v}$ for the
displayed bilinear pairings. Hence $O_2^T$ also commutes
with $\ell_1$ translations and maps $\mathrm{Per}$
into itself. Write $O_2^Tp^*f'=p^*g$.
For every source field $f$, intertwining and
\eqref{eq:rgcov-pullback-pairing} give
\[
 J_c^3\langle g,f\rangle_{T_1}
 =\langle O_2^Tp^*f',p^*f\rangle_{T_2}
 =J_c^3\langle f',O_1f\rangle_{T_1}.
\]
Nondegeneracy gives $g=O_1^Tf'$.
\end{proof}

\subsection{The uniform projection estimate}
\label{subsec:rgcov-projection-estimates}

The supports of primitive row and column indices have
sup diameter at most $L$. Assume
$\gamma_wM\ge1$. Then
\begin{equation}\label{eq:rgcov-exponential-image-sum}
 \sum_{w\in(\Z/J_c\Z)^3}
     \exp\big(-\gamma_wd_2(t,t'+\ell_1w)\big)
                     \le64\exp\big(-\gamma_wd_1(t,t')\big).
\end{equation}
To prove this, translate the column lift by an
$\ell_1$ period so that the distance of the two support
sets equals $d_1\le\ell_1/2$.
For each class of $w$, choose a representative $w'$
which realizes its distance on $T_2$. Distinct classes
have distinct representatives. Their distances are
at least $d_1$; for $|w'|_\infty=j\ge2$ they are at least
\[
 \ell_1j-d_1-2L
                         \ge d_1+(\ell_1-2L)(j-1).
\]
This uses the support diameters and $d_1\le\ell_1/2$.
Since $\ell_1\ge13LM$, we have
$\gamma_w(\ell_1-2L)\ge6$.
There are $27$ lattice vectors of norm at most $1$
and at most $26j^2$ of norm $j$.
The sum is therefore bounded by
\[
 \exp(-\gamma_wd_1)
                 \left(27+\sum_{j\ge2}26j^2\exp(-6(j-1))\right).
\]
The latter series starts below $0.269$ and its
successive ratios are at most
$(9/4)\exp(-6)<1/100$. Thus the bracket is below
$27.3<64$, proving \eqref{eq:rgcov-exponential-image-sum}.

There is a second, polymer summation estimate. For a
unit or chamber row index $t$, let $\square_t$ be its
$M$-cube on $T_2$. By
\eqref{eq:rgcov-projected-polymer-length} and counting,
\begin{equation}\label{eq:rgcov-polymer-image-sum}
 \sum_{\substack{pX'=X\\X'\supseteq\square_t}}
    \exp\big(-(\kappa_w+\kappa_0)d_M(X')\big)
                 \le K_0\exp\big(-\kappa_wd_M(X)\big).
\end{equation}
Indeed the factor with rate $\kappa_w$ is at most
the factor on the right; the remaining sum is bounded
by $\sum_{X'\supseteq\square_t}\exp(-\kappa_0d_M(X'))\le K_0$.
The additional top-system decay $\kappa_0$ pays for the
polymer image sum.

\begin{proposition}\label{prop:rgcov-projected-system}
Suppose the system on $T_2$ has the primitive properties
with parameters
\[
 (\alpha_0,\gamma_w,\kappa_w+\kappa_0;C_w^{(1)}),
 \qquad C_w^{(1)}\ge4\exp(2\gamma_w)M^{27\alpha_0}.
\]
Its projection has the corresponding properties on $T_1$
with
\begin{equation}\label{eq:rgcov-common-constant}
                          C_w=64K_0C_w^{(1)}
\end{equation}
and parameters $(\alpha_0,\gamma_w,\kappa_w;C_w)$.
In particular it satisfies the primitive hypotheses of
Theorem~\ref{thm:rgnl-current-step}.
At zero the resolvent estimate retains $(1+y)^{-1}$.
\end{proposition}

\begin{proof}
We verify the identities, estimates and local structure
needed for the step.
At full strength the exact intertwining identities and
the primitive sum identities on $T_2$ give
\[
 H_{k,1}^{\mathrm{ph}}(1)=H_{k,1}^{\mathrm{ph}},\quad
 C_{k,1}^{1/2}(1)=C_{k,1}^{1/2},\quad
 H_{k,1}(1,A)=H_{k,1}(A),\quad
 \Gamma_{k,1}(1,A)=\Gamma_{k,1}(A).
\]
For the two polymer families, sum first over all $X$:
each $X'$ occurs once, under $X=pX'$.
Equation~\eqref{eq:rgcov-exact-fermion-intertwining}
then identifies the sum with the exact operator on $T_1$.
All projected entries are finite sums of original entries
composed with the linear pullback. Continuity on the
closed parameter polydisk, holomorphy in its interior and
in the connection for any fixed parameter, and Borel
dependence on the resolvent parameter follow from these
finite sums.

For the photon reconstruction, its first and H\"older
derivative kernels have the same image-sum formula:
the difference acts in the row variable, and short
lifted differences agree on the two tori.
Equation~\eqref{eq:rgcov-exponential-image-sum} bounds
each by $64C_w^{(1)}\exp(-\gamma_wd_1)$.
It gives the same bound for the square-root kernel
and for $\Gamma_k$.
The $H_k$ sup and short-pair H\"older bounds instead
follow by the isometric pullback in
\eqref{eq:rgcov-pullback-pairing}; their constant is
$C_w^{(1)}$. A short pair and its six paths lift
identically. An intracluster quotient is
$\eta_k^{-\alpha}(f(c,g')-f(c,g))$ on each torus.

For $M_k(X,A)$, only $X'$ containing $\square_t$
contribute to its row, by original polymer support.
Combine \eqref{eq:rgcov-signed-image-kernel},
\eqref{eq:rgcov-exponential-image-sum} and
\eqref{eq:rgcov-polymer-image-sum} to get
\[
 |M_{k,1}(X,A)(t,t')|_c
     \le C_w\exp(-\gamma_wd_1(t,t')-\kappa_wd_M(X)).
\]
For a positive-level fine resolvent, the pullback of
a unit-block indicator is the disjoint sum of its
$J_c^3$ translated large-torus block indicators.
The same two estimates bound its block operator
norm. At zero use its chamber kernel directly:
the original factor $(1+y)^{-1}$ is common to all
terms of both image sums and remains on the right.
All column weights are unchanged, so no further
factor $48$ appears.

For component locality, an index of one parameter kind
whose support is contained in its active union belongs
to exactly one link-component. Here is the geometric
point. Site supports are singletons. For a bond of
length $a\in\{\eta_k,1\}$, cubes of radius $M$ meeting
its support either share an endpoint or would have
centers differing in that direction by $2M+a$.
The latter is impossible since the centers differ by
a multiple of $M$ and $0<a<M$.
For a coordinate index, its support has diameter at
most $L$. The centers of two cubes of radius $LM$
meeting it differ coordinatewise by less than $3LM$,
and are multiples of $LM$; choosing the coordinate
midpoint when the difference is $2LM$, and either
center when it is $LM$, exhibits a lattice site in
both cubes. Thus all cubes meeting the support are
linked. The same argument applies on either torus
by lifting the support and these cubes into a
nonwrapping neighborhood.

Pull back the active parameter cubes on $T_1$.
A link-component on $T_2$ maps into one component
on $T_1$, since the covering map preserves links.
For a row in a small-torus component, the unique
large-torus component containing its lift maps into
that component. Original component restriction
confines all surviving columns there.
Changing the parameters and connection outside the
small component leaves these original rows unchanged;
the signed image sum therefore preserves their row
dependence. The identical argument applied to each
lifted column preserves column component restriction.
Each row-vanishing property stated for an original family
passes to its projection. In particular, if a cube meeting
$K_c$ is missing, its unique lift meeting the lifted $K_c$
is missing too, so the $H_k$ row vanishes. For $k\ge1$,
the parameter-independent term $\mathcal B_{k,0}$ in
$\Gamma_k$ is retained; it preserves every union of partition
$M$-cubes, by Proposition~\ref{prop:coarse-local-family-bounds}
and its proof.
For a projected polymer family a nonzero term has
both original indices in $X'$, so both projected
indices lie in $X$. Equality of connections on bonds
meeting $X$ gives equality on bonds meeting every
such $X'$, preserving the polymer's local dependence.

Gauge covariance passes through pullback because
the local gradients, restriction to unit sites and
phases commute with it. Lattice covariance follows
by lifting a symmetry $(S,a)$ with one representative
$a\in LM\Z^3$ to the same affine map on both tori;
it permutes the image polymers. A change of representative
multiplies source and target spinors by the same sign,
which cancels in the operator relation.
Transposition passes by
\eqref{eq:rgcov-transpose-intertwining};
the constant $\mathsf C$ commutes with pullback.
Complex conjugation and $\varkappa$ also commute
with it. These facts prove the charge and reality
relations for the projected primitives.

At zero, the photon and fermion diagonal factors
are unchanged by projection. Cubes containing a
given bond, and cubes meeting $K_c$, have unique
corresponding lifts, so their parameter products
are the same on both tori. The fermion factor is
the same for all $48$ labels at $c$.
Their constants are bounded by
$4\exp(2\gamma_w)M^{27\alpha_0}\le C_w^{(1)}$.
The remaining zero-level properties are those
just proved for the square root, chamber covariance
and local resolvent. This completes the system
verification with the common constant
\eqref{eq:rgcov-common-constant}.
\end{proof}

\subsection{A common top-torus system}
\label{subsec:rgcov-top-system}

On a chosen top torus use the photon families
\eqref{eq:photon-weakened-definitions} and the local
fermion families \eqref{eq:local-coarse-families}.
In the fermion walk construction take the polymer
decay rate $\kappa'=\kappa_w+\kappa_0$.
The formulas do not depend on this rate; it specifies
the parameter restrictions and the constants in their
estimates. We write $\gamma^F=\gamma$ for the fermion
decay rate in \eqref{eq:zero-row-bounds}. Define
\begin{gather}
 \alpha_0=1/54,\qquad
 \gamma_w=\min(\gamma^F,\gamma^P),\quad
                                \gamma^P=\gamma_P/2, \notag\\
 C^P=116C_P^{\mathrm{ph}}M^{402\alpha_0},\qquad
 C_w^{(1)}=\max\{C^F,C^P,4\exp(2\gamma_w)M^{27\alpha_0}\},
 \qquad C_w=64K_0C_w^{(1)}.             \label{eq:rgcov-top-constants}
\end{gather}
Here $C^F,\gamma^F$ are the constants of the fermion
construction at the indicated polymer rate.
The fixed $M$ satisfies both its fermion condition and
\eqref{eq:photon-weakening-scale}; the charge satisfies
the inverse and walk conditions of those constructions.

These families form a system with
$(\alpha_0,\gamma_w,\kappa_w+\kappa_0;C_w^{(1)})$.
This follows by matching their proved properties.
The photon kernel multipliers are polynomials in their
parameter cubes; at full strength the partition of unity
and the products give the exact kernels. Their path
products confine rows and columns to their active
link-components. The photon derivative and square-root
bounds have constants $C^P$ and rate $\gamma^P$.
The fermion families are the convergent weakened walks
and the localized block expressions, with their proved
support, covariance, analyticity and H\"older estimates.
Replacing their rates $\gamma^F,\gamma^P$ by the smaller
$\gamma_w$ and their constants by the larger $C_w^{(1)}$
preserves every inequality. The row and column
restriction identities contain no decay constant and
are unchanged.

At level zero the two diagonal families are used.
The photon derivative bound has constant
$4\exp(2\gamma_w)M^{27\alpha_0}$, and the fermion
diagonal factor is bounded by
$M^{8\alpha_0}$, below that constant.
There is no primitive short pair with $0<d_X<1$ at $\eta_0=1$.
The covariance and local resolvent are chamber kernels
with measure $1/48$, and the resolvent bound includes
$(1+y)^{-1}$. A unit block has total chamber measure
$1$, since
$(\eta_k^3/48)(48\eta_k^{-3})=1$.
Consequently the positive-level block estimates and
zero-level chamber estimates use precisely the common
constants \eqref{eq:rgcov-top-constants}.
Projecting this system by Proposition~\ref{prop:rgcov-projected-system}
gives a system with these primitive bounds and the same $C_w$ on every
odd smaller torus. The top system itself has these
weaker bounds as well. None of these constants depends
on the level, side, mass, charge or harmonic background.

\subsection{Common lifts and local algebras}
\label{subsec:rgcov-common-lifts}

For the rest of the comparison assume
\begin{equation}\label{eq:rgcov-comparison-side}
                      \ell_1>(31L+16)M.
\end{equation}
Put $d_U=2L+1$, $c_l=2L$.
For a finite nonempty union $\widetilde X$ of covering
$M$-cubes, let $\operatorname{box}(\widetilde X)$ be
its smallest rectangular union of cubes and let
$\operatorname{diam}(\widetilde X)$ be that box's
largest side. Define
\[
 N(\widetilde X)=\operatorname{box}(\widetilde X)^{\sim6d_U}.
\]
We call $\widetilde X$ admissible
when $\operatorname{diam}(\widetilde X)
\le\ell_1-14d_UM$; a polymer is flat when it is the
image of a connected admissible $\widetilde X$.
These are the sets on which the two torus algebras
will be identified.

\begin{proposition}\label{prop:rgcov-common-lifts}
For admissible $\widetilde X$, projection is injective
on $N(\widetilde X)$ and on every supported index
contained there. An index of either torus supported
in the image of a subset of $N(\widetilde X)$ has
exactly one lift contained in that subset.
Pairs have their shortest lifts, with the same
$d_X$ and six paths.
Local incidences of total support diameter below
$2d_UM$ are preserved in both directions.

There is a norm-preserving algebra isomorphism
$\iota_{\widetilde X}$ between the two subalgebras
generated by ordinary fields over $N(\widetilde X)$,
pair fields with both endpoints there, and their
unit fluctuation and auxiliary fields there.
It identifies covering generator indices.
At zero it includes the chamber $W$ fields and
no auxiliary $\Psi$.
\end{proposition}

\begin{proof}
Every side of $N(\widetilde X)$ is at most
$\ell_1-2d_UM$. Since faces contain no fine lattice
sites, any two of its sites satisfy
\[
 |u-u'|_\infty\le\ell_1-2d_UM-\eta_k.
\]
For $w\ne0$ this implies
\begin{equation}\label{eq:rgcov-lift-separation}
              |u-u'-\ell_iw|_\infty\ge2d_UM+\eta_k.
\end{equation}
It proves disjointness of $N(\widetilde X)$ from
its nonzero $\ell_1$ translates, injectivity, and
uniqueness of any representative with distance
below $2d_UM$.

To lift an index supported in the projected subset,
translate one of its lifts so that one support point
lies in that subset. Every remaining support point
is congruent to a point of the subset and is within
$2LM<2d_UM$ of the first. Equation~
\eqref{eq:rgcov-lift-separation} forces the congruence
translation for each point to be zero.
Thus the whole index lies in the subset; injectivity
gives uniqueness. In particular the shortest pair
has its endpoints there, distance at most $M$, and
its same covering paths.

For an incidence relation, lift a tuple realizing
the projected relation and translate its first index
to the prescribed lift. Each other lifted index
differs from the prescribed one by an $\ell_i$
period. Its support is at distance below $2d_UM$
from the first one. Equation~
\eqref{eq:rgcov-lift-separation} again forces the
period to vanish. This proves the reverse incidence
implication; the forward one is projection.
The relations used below meet this range condition:
nearest neighbors and chamber neighbors, short
pairs and their paths, block and cube membership,
common faces, linked parameter cubes, cube meeting,
and coordinate support. The largest diameter is
$4LM<2d_UM$, for two linked photon parameter cubes.
They are all invariant under $LM\Z^3$ translations.

No two covering generator indices in $N(\widetilde X)$
have the same period class on either torus.
Relative to any fundamental generator system their
images are therefore independent generators, up
to their antiperiodic signs.
Identifying the same covering indices extends
uniquely to a free Grassmann algebra isomorphism.
It preserves bar indices and degrees and sends
basis monomials to basis monomials, up to the
same coefficient signs. All generator weights and
index measures are common, so it preserves norms.
The maps for nested admissible sets restrict to
one another.
\end{proof}

A local activity on $X_i=\operatorname{pr}_i\widetilde X$
belongs to this identified algebra. Ordinary sites
lift to $\widetilde X$; an admissible pair lifts
to $\widetilde X^{\sim1}$. For its first kind the
two sites and paths lie in $\widetilde X$.
For its second kind the pair is $(x,x^{(3)})$
with an endpoint over $\widetilde X$.
Each assertion follows from the incidence comparison.
Two local families $F_1,F_2$ agree on flat polymers
when
\begin{equation}\label{eq:rgcov-local-agreement}
 \iota_{\widetilde X}F_1(X_1;A_1,\mathscr C_1)
                             =F_2(X_2;A_2,\mathscr C_2)
\end{equation}
whenever the two admissible connection pairs have
equal lifts on $N(\widetilde X)$.
It suffices to check $A_2=p^*A_1$,
$\mathscr C_2=p^*\mathscr C_1$:
for other pairs with equal lifts, locality of
$F_2$ identifies them with these pullbacks on all
bonds meeting $\widetilde X$.
Thus once agreement is proved, equality holds
already under equality on those read bonds.

Every small polymer is flat. The net diameter bound
gives
$\operatorname{diam}(\widetilde X)
\le M(d_M(\widetilde X)+3)\le M(L+3)$ for a small
connected lift. By \eqref{eq:rgcov-comparison-side},
$M(L+3)+14d_UM<\ell_1$.
The nonwrapping lift proposition then identifies
small torus polymers containing a chosen cube
with small covering polymers containing its
chosen lift. This identification is the same
on both tori. Conversely, if an admissible
connected lift projects to a small polymer,
the small nonwrapping lift is its period translate:
start at one common cube and propagate equality
along adjacent cubes using the incidence result.
Hence smallness, tree length and volume coincide
for its two images.

Subpolymers and parameter cubes contained in an
admissible $\widetilde Y$ correspond bijectively
on the two tori. Common faces and links are
preserved, so incompatibility is preserved.
The definitions by cube meeting give
\[
 P_0(\operatorname{pr}_i\widetilde X)
                 =\operatorname{pr}_iP_0(\widetilde X),\qquad
 \operatorname{pr}_iY'(\widetilde X,\widetilde P)
          =Y'(\operatorname{pr}_i\widetilde X,\operatorname{pr}_i\widetilde P).
\]
When $\widetilde X$ and all cubes of $\widetilde P$
are contained in $\widetilde Y$, the resulting
$Y'$ is contained there too, since the $M$-cubes
partition sites. Reblocking satisfies
$\overline{\widetilde X}^{\,L}
\subset\widetilde Y^{\sim(L-1)}$ and commutes with
projection. These larger sets still lie in
$N(\widetilde Y)$, so their incidences have the
same unique lifts.

\subsection{Supported rows near a flat set}
\label{subsec:rgcov-supported-rows}

The preceding identification applies to the weakened
primitives because setting parameters to zero separates
their components. We first record the component argument.
Let $s$ vanish at every parameter cube not contained
in a union $Z$ of $M$-cubes. If $Z_1$ is a face-connected
component of $Z$ and
$\square(t)^{\sim2L}\subset Z_1$, then
\begin{equation}\label{eq:rgcov-component-row}
 \begin{split}
 O(s,A)(t,t')&=0 \quad\text{unless }t'\text{ is contained in }Z_1,\\
 O(s,A)(t,\cdot)&=O(s\mathbf1_{Z_1},A')(t,\cdot).
 \end{split}
\end{equation}
Here $O$ is a parameter primitive, and $A'$ is any
connection in the primitive's domain equal to $A$ on bonds meeting $Z_1$;
the photon primitives have no connection argument.
The set $\square(t)$ is the union of the $M$-cubes
meeting the support of $t$.

To prove \eqref{eq:rgcov-component-row}, choose a parameter
cube containing that support inside
$\square(t)^{\sim2L}$. For a chamber site or a unit site
it is its $M$-cube. For a fine bond it is the parameter
cube of radius $M$ about the $M$-cube center of one end:
both ends lie in it, and it is contained in
$\square(t)^{\sim1}$. For a photon coordinate index choose
the parameter cube of radius $LM$ about the $LM$-cube
center of one point of its support.
The support diameter is at most $L$ and $M\ge2$,
so this cube contains the entire support.
If $z_c$ is the $M$-cube center of that point and $z''$
is the $M$-cube center of any site in the parameter cube,
then
\[
 |z''-z_c|_\infty
 <(3L/2+1)M\le(2L+1)M.
\]
These centers differ by a multiple of $M$, hence the
cube is contained in $\square(t)^{\sim2L}$.

Let $\mathscr Y$ be all parameter cubes contained in $Z$,
and let $\mathscr Y_1$ be their link component containing
the cube just chosen. The $M$-cubes meeting
$\bigcup\mathscr Y_1$ form a face-connected polymer
contained in $Z$. It meets $\square(t)\subset Z_1$,
and therefore is contained in the component $Z_1$.
The component restriction property of the weakened
primitives, proved in the photon and fermion constructions,
now confines the row to $\bigcup\mathscr Y_1$.
Replacing $s$ by $s\mathbf1_{Z_1}$ leaves all its
parameters there unchanged; replacing $A$ by $A'$
leaves all its read bonds unchanged. This proves
\eqref{eq:rgcov-component-row}.

\begin{proposition}\label{prop:rgcov-supported-rows}
Let $\widetilde Y$ be admissible, $Y_i=\operatorname{pr}_i
\widetilde Y$, and $0\le j\le5d_U$. Let $s_1$ vanish
outside $Y_1$. Define $s_2$ by giving the corresponding
parameter cubes contained in $\widetilde Y$ the same
values and setting all other entries to zero.
Suppose $A_i\in\widetilde{\mathcal G}_k(T_i)$ have equal
lifts on the bonds contained in
$\widetilde Y^{\sim(j+d_U)}$.
For each parameter primitive and each covering row $t$
contained in $\widetilde Y^{\sim j}$, the two rows,
written with covering indices in $N(\widetilde Y)$,
coincide:
\begin{equation}\label{eq:rgcov-local-primitive-row}
                 O_1(s_1,A_1)(t,t')=O_2(s_2,A_2)(t,t').
\end{equation}
Both rows have columns contained in
$Y_i^{\sim(j+c_l)}$, and read the connection only on
bonds contained in $Y_i^{\sim(j+d_U)}$.

For a connected nonempty $\widetilde X\subset\widetilde Y$
and connections with equal lifts on bonds meeting
$\widetilde X$, the local polymer families also satisfy
\begin{equation}\label{eq:rgcov-local-polymer-kernel}
 O_1(\operatorname{pr}_1\widetilde X,A_1)(t,t')
       =O_2(\operatorname{pr}_2\widetilde X,A_2)(t,t')
\end{equation}
at rows and columns over $\widetilde X$.
Here $O=M_k$ for $k\ge1$, or $O=S_{k,y}$ for any $k$
and $y\ge0$.
\end{proposition}

\begin{proof}
Put $\widetilde B=\widetilde Y^{\sim(j+c_l)}$.
Its one-cube enlargement is
$\widetilde Y^{\sim(j+d_U)}\subset N(\widetilde Y)$.
For the row $t$, the set $\square(t)$ is one $M$-cube
or two face-adjacent $M$-cubes. This follows directly
for site and bond indices; a coordinate support lies
in one $L$-block or is the support of one bond.
Thus $\square(t)^{\sim c_l}$ is a connected box
contained in $\widetilde B$.
Let $\widetilde B_1$ be the component of $\widetilde B$
containing this box.

Its image is a component of
$\operatorname{pr}_1\widetilde B$.
On $T_2$ its image is a component both of
$\operatorname{pr}_2\widetilde B$ and of
$p^{-1}(\operatorname{pr}_1\widetilde B)$.
Indeed a cube adjacent to its image has a lift adjacent
to a cube of $\widetilde B_1$. This lift lies in
$\widetilde B^{\sim1}\subset N(\widetilde Y)$.
If it belongs to a period translate of $\widetilde B$,
separation \eqref{eq:rgcov-lift-separation} forces that
translate to be zero. It then lies in $\widetilde B$
and in the same component. This proves maximality
of the images as components.

Apply \eqref{eq:rgcov-component-row} on $T_1$.
It confines the row to $\operatorname{pr}_1\widetilde B_1$
and lets the connection be changed off the bonds meeting
that component. Such bonds lift into
$\widetilde B^{\sim1}$, proving the claimed read set.
On the parameter cubes contained in
$\operatorname{pr}_2\widetilde B_1$, the families
$p^*s_1$ and $s_2$ agree. To see this even at cubes
not contained in $\widetilde Y$, observe that each
such cube in $N(\widetilde Y)$ has a unique lift;
its image cannot also be the image of a cube in
$\widetilde Y$. Both entries are therefore zero.
The connections $p^*A_1$ and $A_2$ agree on bonds
meeting $\operatorname{pr}_2\widetilde B_1$.
Applying \eqref{eq:rgcov-component-row} twice on $T_2$
gives
\[
 O_2(p^*s_1,p^*A_1)(t,\cdot)
       =O_2(s_2,A_2)(t,\cdot),
\]
with columns confined to $\operatorname{pr}_2\widetilde B_1$.

In the signed image formula
\eqref{eq:rgcov-signed-image-kernel}, a nonzero summand
at $t'+\ell_1w$ would have a lift in $\widetilde B_1$.
Both this lift and $t'$ belong to $N(\widetilde Y)$;
their difference is an $\ell_1$ period.
Separation forces $w=0$ in $(\Z/J_c\Z)^3$.
That summand has character $1$, proving
\eqref{eq:rgcov-local-primitive-row}.

For the polymer families, every polymer $X'$ on $T_2$
with $pX'=\operatorname{pr}_1\widetilde X$ is
$\operatorname{pr}_2(\widetilde X+v)$ for one
$v\in\ell_1\Z^3/\ell_2\Z^3$.
Distinct classes give disjoint cube sets.
To verify the assertion for an arbitrary $X'$, lift
two adjacent cubes as cubes of $\widetilde X+v_1$
and $\widetilde X+v_2$. Their adjacency places a
translate of the second cube in
$\widetilde X^{\sim1}\subset N(\widetilde Y)$.
Separation forces $v_1=v_2$ modulo $\ell_2\Z^3$.
Connectedness gives one class for every cube of $X'$,
and projection onto the whole image gives equality
with that entire copy.

The projected kernel is consequently a double sum over
$v$ and the column image $w$. The local polymer support
requires both its row and column to be over
$\operatorname{pr}_2(\widetilde X+v)$.
The common row lift forces $v=0$; the common column
lift then forces $w=0$. The remaining term is
$O_2(\operatorname{pr}_2\widetilde X,p^*A_1)(t,t')$.
Locality on the bonds meeting this polymer permits
$p^*A_1$ to be replaced by $A_2$.

The barred mean field satisfies
$\bar H_k(s,A)=\mathsf C^{-1}H_k(s,-A)\mathsf C$,
so its comparison follows at the connections $-A_i$.
At zero the mean-field kernels are
$48\delta_{x,x'}\mathbf1_4$ times their parameter products;
the coordinate and chamber measures are the same
on the two tori, and the chamber label is preserved.
There is no $M_0$ family. The proof applies to
$\Gamma_0$ and $S_{0,y}$ on the full chamber indices.
\end{proof}

\subsection{Finite rows and the local operations}
\label{subsec:rgcov-finite-rows}

For finite operators, equality of rows follows directly
from their covering formulas. Suppose a row has all its
columns in $N(\widetilde Y)$, and its connection and
polymer membership data agree on the bonds and sites
which that row reads. The two formulas are then the
same finite sum on either torus. There are no repeated
column classes by Proposition~\ref{prop:rgcov-common-lifts}.
Writing the row as
$\sum_{t'}r(t,t')f(t')$, its kernel is
\begin{equation}\label{eq:rgcov-row-kernel-weight}
                           O(t,t')=\nu(t')^{-1}r(t,t').
\end{equation}
Thus equality holds for the kernels as well, with the
same column weight.

We need this observation at connections $B_i$ having
equal lifts on bonds meeting a connected
$\widetilde U\subset\widetilde Y$.
The photon coordinate map $C$ reads no connection:
a noncentral face bond is its own coordinate, and
the central bond is the negative sum of the other
face bonds and the interior straight-path contributions
$n_c(b)/L$. All columns are supported in the one or
two $L$-blocks of the row bond's ends. Hence its rows
at bonds in $\widetilde U^{\sim(1+2d_U)}$ agree and
have range at most $2L-1$.

The covariant difference row at an admissible pair is
\[
 d_X(x,x')^{-\alpha}
               \big(U_B(x,x')f(x')-f(x)\big).
\]
Its only columns are $x,x'$, and it reads exactly
the six pair paths. These paths are contained in
$\widetilde U$ for the first kind, or are the bond
meeting $\widetilde U$ for the second kind.
The range is $M$; inside a chamber cluster the
transport is $1$ and no bond is read.
These rows therefore agree, also at $-B_i$.

For $k\ge1$ the reconstruction rows $T^{(1)},T^{(2)}$
and their barred versions are the two-piece formulas
of the reconstruction estimate. At $y\in\widetilde U$
they have columns over $\Delta_y$ and at $x^{(3)}$,
or at the pair $(x,x^{(3)})$ across its boundary.
They read the bonds in $\Delta_y$ and the bonds
$(c,c_3)$ starting there. All meet $\widetilde U$.
Their membership indicators are the same because
\[
 u\in\widetilde U+\ell_i\Z^3
       \quad\Longleftrightarrow\quad u\in\widetilde U
             \qquad(u\in\widetilde U^{\sim1}),
\]
again by separation. Every site involved is at
distance at most $(1+\eta_k)/2<1$ from $y$.
This proves equality of both pieces, including
the boundary term.

The finite rows of $Q_kJ^*$ and $JQ_k^T(-B)$ have
their columns and read bonds in a single $\Delta_y$.
The row of $B_{k,r}$, where $r\ge0$ is the resolvent
parameter, has its columns in one $L$-block and reads
the paths from its center to its sites.
These paths lie in that block. Their kernels are
\begin{align*}
 (Q_kJ^*)(y,x)&=\mathbf1[c\in\Delta_y]\varphi_B(y,c)\mathbf1_4,\\
 (JQ_k^T(-B))(x,y)&=\mathbf1[c\in\Delta_y]\varphi_B(y,c)^{-1}\mathbf1_4,\\
 B_{k,r}
   &=(b_k+\mathrm i\gamma_5r)^{-1}I\\
   &\quad+\big((b_k+b_L+\mathrm i\gamma_5r)^{-1}
                    -(b_k+\mathrm i\gamma_5r)^{-1}\big)P_B .
\end{align*}
The projection $P_B$ has kernel $L^{-3}$ times the two
block transports when its indices lie in the same
$L$-block, and zero otherwise. The first two rows
have range below $1$, the third range at most $L-1$.
At rows in $\widetilde U$ their columns and read
bonds stay in its $M$-cube. They agree on both tori.

At zero, $F_0$ has only its diagonal, the four chamber
neighbors, and its block-penalty columns. The hopping
part of its kernel is
\[
 48\Big(
  \bigl(\mu+\tfrac{\rho}{2}\sum_{i=0}^3w_i\bigr)\delta_{x,x'}\mathbf1_4
       -\sum_{j=0}^3\delta_{x',x^{(j)}}M_j(x)U_B(x,x^{(j)})
       \Big).
\]
The factor $48$ is the inverse column weight.
Only the fourth neighbor reads a bond meeting
$\widetilde U$; the penalty reads paths inside the
$L$-block of $c$. All columns lie over
$\widetilde U^{\sim1}$ and its range is at most $L-1$.
Thus this entire row agrees at zero.

These equalities also identify the mass form
with the covering sum:
\[
 \langle\bar\psi,\psi\rangle_{\operatorname{pr}_i\widetilde\square}
 =(\eta_k^3/48)\sum_{c\in\widetilde\square,g,\beta}
                          \bar\psi_\beta(c,g)\psi_\beta(c,g).
\]
At zero the two halves of the first-step vertex have
the common coefficient $1/96$, the same two boundary
transports, and the same generators over the cube
and at its outward pairs. Hence the first vertex
also agrees under $\iota_{\widetilde Y}$.

The remaining local operations can be compared in
this common free algebra. Equality of elements gives
equality of their constant coefficients and of their
antisymmetric quadratic kernels: expand in ordered
covering monomials and use the common measure
$\eta_k^3/48$. For a small polymer, the adjustment
$\Theta_{X,B}$ uses averages over shortest chains
of $M$-cubes in $X$. Common faces, path membership
and transports are identical on its lift, as are
its volume and the coefficients $d_X^\alpha$.
Therefore
\begin{equation}\label{eq:rgcov-adjustment-intertwining}
 \iota_{\widetilde X}\Theta_{X_1,B_1}
                      =\Theta_{X_2,B_2}\iota_{\widetilde X}.
\end{equation}

If $\widetilde U$ is admissible at level $k+1$,
then $L\widetilde U$ is admissible at level $k$.
Its box diameter is $L\operatorname{diam}(\widetilde U)$,
and $N(L\widetilde U)\subset L N(\widetilde U)$.
The scaling rules on the common generators are
$\psi(x)\mapsto L^{-1}\psi(x/L)$ and
$\chi(x,x')\mapsto L^{-1-\alpha}\chi(x/L,x'/L)$.
They preserve the chamber label and the period signs.
Thus
\begin{equation}\label{eq:rgcov-scaling-intertwining}
 S_L\iota_{L\widetilde U}=\iota_{\widetilde U}S_L.
\end{equation}

Finally, the integrations agree on any local element.
For $W$, expand in monomials. A monomial moment is
zero unless its modes are paired; otherwise it is
its ordering sign times $\nu^{-\#\mathrm{pairs}}$.
The mode indices and the weight $\nu$ agree.
At zero this weight is $1/48$, so the factor is
$48^{\#\mathrm{pairs}}$; other modes integrate to $1$.
This proves
\begin{equation}\label{eq:rgcov-W-integration}
 \iota_{\widetilde Y}\int\dd\mu_I(W)\,F
                    =\int\dd\mu_I(W)\,\iota_{\widetilde Y}F.
\end{equation}
The restricted photon law is a product of the same
normalized Gaussian on $[-p_{0,k},p_{0,k}]$ in each
coordinate. Consequently its marginal on common
coordinates is identical. With $d_V=2-2L^{-3}$,
its vacuum contribution on either torus is
\begin{equation}\label{eq:rgcov-box-vacuum}
 \delta\varepsilon_k^0
 =-d_V\log\left((2\pi)^{-1/2}
             \int_{-p_{0,k}}^{p_{0,k}}\exp(-z^2/2)\dd z\right).
\end{equation}
This number depends on $e_k$ and $L$, and is independent
of the side.

\subsection{Extraction and determinant comparison}
\label{subsec:rgcov-determinants}

Suppose $\mathscr E_i\in K_k^x(T_i)$ agree in the sense
of \eqref{eq:rgcov-local-agreement}. For a small
covering polymer $\widetilde X$, their values at $(0,0)$
have the same constant and quadratic coefficients.
Their local extractions
\eqref{eq:rg-local-extractions} are therefore the same:
their integrated kernels and their common volume are
the same. Fixing a covering cube identifies all small
polymers containing its image on either torus. The sums
\eqref{eq:rg-counterterm-maps} consequently give
\begin{equation}\label{eq:rgcov-counterterms}
 m(\mathscr E_1)=m(\mathscr E_2),\qquad
 \varepsilon_{\mathrm v}(\mathscr E_1)
                         =\varepsilon_{\mathrm v}(\mathscr E_2).
\end{equation}
Subtracting the common constant and mass forms preserves
local agreement of $R\mathscr E_i$.
On a small polymer, its quadratic part also agrees;
\eqref{eq:rgcov-adjustment-intertwining} proves agreement
of $(R\mathscr E_i)^{\mathrm{nat}}$.
On other polymers neither of these operations changes
the relevant element.

The determinant terms agree on every flat polymer,
not just on small ones. Let $X_i=\operatorname{pr}_i
\widetilde X$ and suppose $B_i$ have equal lifts on
bonds meeting $\widetilde X$.
For $k\ge1$, set
\[
 O_1=B_{k,r}\gamma_5B_{k,r},\qquad
 O_2=Q_kJ^*S_{k,r}(X_i,B_i)JQ_k^T(-B_i).
\]
The trace in the determinant integral is
\begin{equation}\label{eq:rgcov-local-determinant-trace}
 \operatorname{Tr}(O_1O_2)
       =\sum_{z_1,z_2\in X_i\cap\Z^3}
                      \tr\big(O_1(z_2,z_1)O_2(z_1,z_2)\big).
\end{equation}
Here and below $\operatorname{Tr}$ is the trace on the whole finite
function space, while $\tr$ is the spin trace.
Indeed
\begin{align*}
 O_2(z_1,z_2)
   &=\int_{\Delta_{z_1}}\dd x\int_{\Delta_{z_2}}\dd x'\,
       \varphi_{B_i}(z_1,c)S_{k,r}(X_i,B_i)(x,x')
                            \varphi_{B_i}(z_2,c')^{-1}.
\end{align*}
Local polymer support confines $x,x'$ to $X_i$,
and $\Delta_z$ lies in the $M$-cube of $z$.
Thus both unit sites in this trace lie in $X_i$.
The rows of $B_{k,r}$ stay in their own $L$-block,
so the intermediate sum in $O_1(z_2,z_1)$ also stays
in $X_i$. All sums may now be formed on the lift
$\widetilde X$; the two antiperiodic kernel signs
cancel in each trace summand.
The finite-row comparison and
\eqref{eq:rgcov-local-polymer-kernel} make every factor
the same on both tori.

At zero the required traces are
\begin{align}
 \operatorname{Tr}(F_0(B_i)S_{0,r}(X_i,B_i))
   &=\int_{X_i}\dd x\int_{X_i}\dd x'\,
          \tr\big(F_0(B_i)(x,x')S_{0,r}(X_i,B_i)(x',x)\big),
                                      \label{eq:rgcov-zero-trace}\\
 \operatorname{Tr}(\gamma_5S_{0,r}(X_i,B_i))
   &=\int_{X_i}\dd x\,\tr\big(\gamma_5S_{0,r}(X_i,B_i)(x,x)\big).
                                      \label{eq:rgcov-zero-trace-gamma}
\end{align}
Support again restricts all arguments to $X_i$.
Both are the same covering sums with measure $1/48$
and all $48$ chamber labels. The local $F_0$ and
$S_{0,r}$ kernels agree, so the traces agree for
every $r\ge0$.

These pointwise identities can be integrated in $r$.
The required domination follows
from the primitive estimates. Since $P_B^2=P_B$,
\[
 B_{k,r}\gamma_5B_{k,r}
  =\gamma_5\big((b_k+\mathrm i\gamma_5r)^{-2}(I-P_B)
                +(b_k+b_L+\mathrm i\gamma_5r)^{-2}P_B\big).
\]
In the component matrix norm,
$|(a+\mathrm i\gamma_5r)^{-1}|_c
\le\sqrt2(a^2+r^2)^{-1/2}$ for $a>0$.
Hence the first trace is bounded by a finite
$B_i,X_i$ dependent constant times
$(b_k^2+r^2)^{-1}$.
The other factor is uniformly bounded by the local
resolvent estimate and the two finite block transports.
This majorant is integrable because $b_k>0$.
At zero, $|S_{0,r}(X_i,B_i)(x,x')|_c
\le C_w(1+r)^{-1}$.
The first trace in \eqref{eq:rgcov-zero-trace} is
integrated against $\dd r/r$ on $[1,\infty)$,
and \eqref{eq:rgcov-zero-trace-gamma} against $\dd r$
on $[0,1]$; both are integrable.
The kernels are Borel in $r$, so these are
ordinary absolutely convergent scalar integrals.

The two unsymmetrized determinant primitives agree.
Applying the argument at $-B_i,\bar B_i,-\bar B_i$
identifies all four terms in its symmetrization.
It gives
\begin{equation}\label{eq:rgcov-determinant-agreement}
              E_k^{d,(1)}(X_1,B_1)=E_k^{d,(2)}(X_2,B_2).
\end{equation}
Subtracting the values at $\mathscr C_i$ or at zero
proves the corresponding equality for
$\delta E_k^{\mathrm{det}}$ and $E_k^{\mathrm{det}}$.

\subsection{Comparison of the local fluctuation sums}
\label{subsec:rgcov-fluctuation-comparison}

Let $m$ be real with $|m|\le w_k^{\mathrm m}$, let
$\|\mathscr E_i\|_k^x\le w_k^{\mathrm E}$, and assume
$\mathscr E_i=0$ at zero and local agreement as above.
Fix $A_1\in\widetilde R_k^-(T_1)$ and
$\mathscr C_1\in\mathfrak X_k(T_1)$, and put
$(A_2,\mathscr C_2)=p^*(A_1,\mathscr C_1)$.
All small-charge and primitive hypotheses of the
nonlinear step are imposed on both tori.
In particular the contour radii satisfy
$1<r_2\le r_1$, $r_4>1$, and the domain inclusions
of Proposition~\ref{prop:rgnl-contour-domains} hold.
We compare the terms of that step on a connected
nonempty $\widetilde U\subset\widetilde Y$.

Assign common coordinate values
$z_\Upsilon\in[-p_{0,k},p_{0,k}]$ to the photon basis
elements contained in $\widetilde U$, and set all
other coordinate values to zero on each torus.
Assign common parameter values of modulus at most
$M^{\alpha_0}$ to the parameter cubes contained in
$\widetilde U$, and set all other parameters to zero.
These data will be denoted by $\zeta_i,s_i$.
They belong to the respective box and polydisk
because indices on the lift are injective.

The fine fluctuation fields have the same lifts on
$\widetilde U^{\sim(1+d_U)}$:
\begin{equation}\label{eq:rgcov-fine-fluctuation}
                         Z_1(s'_1)(b)=Z_2(s'_2)(b).
\end{equation}
To check the successive supports in their composition,
write it in kernels,
\[
 Z_i(s'_i)(b)
  =\sum_{b',\Upsilon,\Upsilon'}
    H_k^{\mathrm{ph},(i)}(s'_i)(b,b')
          C_i(b',\Upsilon)
            C_k^{1/2,(i)}(s'_i)(\Upsilon,\Upsilon')\zeta_i(\Upsilon').
\]
At such a row $b$, Proposition~\ref{prop:rgcov-supported-rows}
with $j=1+d_U$ confines $b'$ to
$\widetilde U^{\sim2d_U}$ and identifies the row kernels.
The coordinate row has columns in the $M$-cubes of
the two ends of $b'$, hence in
$\widetilde U^{\sim(1+2d_U)}$; those finite rows agree.
The square-root row, with $j=1+2d_U$, confines
$\Upsilon'$ to $\widetilde U^{\sim3d_U}$.
All these sets lie in $N(\widetilde U)$, and the
last coordinate values agree there.
The two compositions are therefore the same finite
covering sum. At zero the first row is diagonal
and the sum over $b'$ has one term.

Put $B_i=A_i+\mathscr C_i$ and
$B_i(t)=B_i+tZ_i(s'_i)$.
For $|t|\le r_1$ these fields lie in
$\widetilde{\mathcal G}_k(T_i)$, and
$A_i+tZ_i(s'_i)\in\widetilde R_k(T_i)$.
They have equal lifts on the bonds just compared.
The fermion mean-field rows at sites at distance
at most $\eta_k$ from $\widetilde U$ consequently agree,
as do the $\Gamma_k$ rows at unit sites in
$\widetilde U$ and the $M_k(X_i,B_i(t))$ kernels
for lifted subpolymers $\widetilde X\subset\widetilde U$.
Here the parameters vanish outside $U_i$.
The row-vanishing and component argument in the proof
of Proposition~\ref{prop:rgnl-finite-locality} confines
their columns to $U_i$ itself: a nonzero mean-field
row has its active parameter component attached to
that set; the $\Gamma_k$ row starting in this component
has the same confined columns. For $k\ge1$, the
parameter-independent term $\mathcal B_{k,0}$ in $\Gamma_k$
also preserves $U_i$, a union of partition $M$-cubes.
Thus all these fields are polynomials in the common
generators over $\widetilde U$.

For example, at a relevant chamber site,
\[
 (H_k^{(i)}(s_i'',B_i(t))\Psi)_\beta(x)
   =\sum_{y\in\widetilde U\cap\Z^3,\beta'}
         H_k^{(i)}(s_i'',B_i(t))(x,y)_{\beta\beta'}
                                      \Psi(y,\beta',0)
\]
for $k\ge1$, with the analogous barred sum.
Their coefficients agree.
At every pair admissible for $\widetilde X$, applying
the common finite difference row preserves this equality.
The differences between these fields at $t=1$ and $t=0$
are precisely the ordinary and pair components of
$J_k^\#\Psi$ used in $G_2,G_3$.
Both values are in the disk because $r_1>1$.
Moreover locality of $R\mathscr E_i$ gives equality
of its values at the pairs
$(A_i+tZ_i,\mathscr C_i)$ on each subpolymer.
These pairs need only agree on its read bonds,
by the criterion following
\eqref{eq:rgcov-local-agreement}.

It follows that the three first-localization contour
integrands agree under $\iota_{\widetilde U}$.
For $G_1$, this is the common local $M_k$ bilinear.
For $G_2$, addition of the common odd field $tJ_k^\#\Psi$
to each generator commutes with an algebra homomorphism;
apply it to the equal local activity.
For $G_3$, the same substitution applies to the equal
mass form, and $m'=m+m(\mathscr E_i)$ is common by
\eqref{eq:rgcov-counterterms}.
Their contour radii and scalar contour measures
are the same. Integration thus gives
\begin{equation}\label{eq:rgcov-first-G}
              \iota_{\widetilde U}G_a^{(1)}(s_1,X_1)
                                      =G_a^{(2)}(s_2,X_2),
                    \qquad a=1,2,3.
\end{equation}
At zero $G_2=G_3=0$ and no auxiliary field occurs.
$G_1$ is the first-step vertex; its two factors
$\exp(\pm\mathrm i qZ_i(c,c_3))-1$ agree by
\eqref{eq:rgcov-fine-fluctuation}, with its common
coefficient $1/96$ and the common outward pair.
This proves \eqref{eq:rgcov-first-G} at that level too.

The finite localization commutes with this equality.
Indeed the pairs $(X,P)$ with $Y'(X,P)=U_i$ correspond
bijectively to lifted pairs $(\widetilde X,\widetilde P)$
inside $\widetilde U$ having
$Y'(\widetilde X,\widetilde P)=\widetilde U$.
If a torus pair has that union, its polymer and all
its parameter cubes are contained in $U_i$, so the
unique lifts give the inverse correspondence.
The identities for $P_0,Y'$ proved above show that
it preserves the parameter prescriptions.
Its summand is therefore
\begin{equation}\label{eq:rgcov-localization-term}
 \int_{[0,1]^{\widetilde P}}\dd u\,
       \prod_{\mathrm{pc}\in\widetilde P}\partial_{u_{\mathrm{pc}}}
                       G_i(s_i[u],\operatorname{pr}_i\widetilde X),
\end{equation}
where $s_i[u]$ is $1$ on $P_0(\widetilde X)$, $u$ on
$\widetilde P$, and zero elsewhere.
Identity \eqref{eq:rgcov-first-G} holds throughout
the parameter polydisk. Since $M^{\alpha_0}>1$,
the integration cube is in its interior, and
coefficient differentiation and integration preserve
the equality. This proves equality of
$\operatorname{Loc}G_a(U_i)$.

The reconstruction substitution replaces each $\Psi(y)$
by the same two-piece row $T^{(1)}\psi+T^{(2)}\chi$.
Both rows, including their boundary membership indicators,
have already been compared. Their resulting ordinary
generators are over $\widetilde U$, and their boundary
pairs are admissible for it. Equality on generators
proves equality of the substitution homomorphisms.
We obtain
\begin{equation}\label{eq:rgcov-first-localization}
              \iota_{\widetilde U}E_{\le3}^{(1)}(U_1)
                                      =E_{\le3}^{(2)}(U_2).
\end{equation}
At zero no reconstruction substitution is made.

For the second localization, the common fermion field is
\[
 \mathrm{Wc}=H_k(s'',B)\Gamma_k(s'',B)W,\qquad
 \overline{\mathrm{Wc}}=\bar H_k(s'',B)\bar W .
\]
For $k\ge1$ both are sums over the common unit sites
of $\widetilde U$. For $k=0$ they are
\[
 \mathrm{Wc}(x)=d(c)\frac1{48}
            \sum_{x'\text{ over }\widetilde U}
                                      \Gamma_0(s'',B)(x,x')W(x'),
 \qquad
 \overline{\mathrm{Wc}}(x)=d(c)\bar W(x),
\]
where the common diagonal factor $d(c)$ vanishes
outside $U_i$. These formulas preserve all chamber
labels and the weight $1/48$.
The kernel equalities proved above, followed by the
common pair difference rows, identify $W_k^\#$ at
every index used by a subpolymer.

The first-localized activity on a subpolymer reads
only photon coordinates supported there. Hence
\eqref{eq:rgcov-first-localization}, proved with coordinates
supported on that subpolymer, still holds when coordinates
are supported on the larger $\widetilde U$.
Adding the common $W_k^\#$ to its generators proves
equality of $G_{\le3}^{\,W}$.
The same addition to $R\mathscr E_i$ and to the mass
form proves equality of the contour integrands of
$G_4,G_5$. The second localization uses only $\pi_k$;
the correspondence and derivative argument
\eqref{eq:rgcov-localization-term} apply with this
one kind. Consequently
\begin{equation}\label{eq:rgcov-finite-F}
                       \iota_{\widetilde U}F_1(U_1)=F_2(U_2).
\end{equation}

We now pass to the local connected logarithm.
Its finite Mayer activity is
\[
 K_i(X_i)=\sum_{\{\widetilde X_a\}:\,\bigcup_a\widetilde X_a=\widetilde X}
                \prod_a\big(\exp(F_i(\operatorname{pr}_i
                                                    \widetilde X_a))-1\big),
\]
where the $\widetilde X_a$ are distinct connected
subpolymers of $\widetilde X$.
The products are even. Their index sets and factors
are the same on both tori by
\eqref{eq:rgcov-finite-F}; coordinate locality allows
their evaluation with common coordinates on
$\widetilde X$.
For the integrated activity $H_i(X_i)$,
\eqref{eq:rgcov-W-integration} identifies the Wick
integrals. The other photon coordinates integrate
to $1$, leaving the same product of restricted normal
laws over the coordinates supported in
$\widetilde X$. Hence $\iota_{\widetilde X}H_1(X_1)=H_2(X_2)$.

The logarithm has the expansion
\[
 E_k^{\#,(i)}(X_i)=
   \sum_{n\ge1}\frac1{n!}
    \sum_{\bigcup_{a=1}^n\widetilde X_a=\widetilde X}
       u^T(\operatorname{pr}_i\widetilde X_1,\ldots,
                    \operatorname{pr}_i\widetilde X_n)
                      \prod_{a=1}^n H_i(\operatorname{pr}_i\widetilde X_a).
\]
Projection preserves incompatibility inside this lift,
so the graph determining $u^T$ is the same.
Every summand agrees under the common algebra map.
Both series converge absolutely in the weighted norm
by Proposition~\ref{prop:rgnl-cluster-logarithm}.
The map $\iota_{\widetilde X}$ preserves that norm,
so equality passes to their sums:
\begin{equation}\label{eq:rgcov-logarithm-agreement}
       \iota_{\widetilde X}E_k^{\#,(1)}(X_1;A_1,\mathscr C_1)
                      =E_k^{\#,(2)}(X_2;A_2,\mathscr C_2).
\end{equation}

\subsection{Agreement of the step after reblocking}
\label{subsec:rgcov-step-comparison}

\begin{theorem}\label{thm:rgcov-step-comparison}
Suppose the two sides at levels $k$ and $k+1$ are
admissible and satisfy
$\ell_1/L>(31L+16)M$. Use the projected and top
systems constructed above, with the hypotheses of
Theorem~\ref{thm:rgnl-current-step} on both tori.
Let the real mass $m$ be common and let the incoming
activities agree on flat polymers, with
$\mathscr E_i=0$ at zero.
Then their extracted mass and vacuum contributions
are equal. Their outgoing activities agree on flat
polymers at level $k+1$.
The external scalar recursion
\eqref{eq:rgnl-external-scalar-map} preserves agreement
of scalar families on such polymers.
All these assertions hold for the activities without
an external slot as well.
\end{theorem}

\begin{proof}
The scalar conclusions are
\eqref{eq:rgcov-counterterms} and
\eqref{eq:rgcov-box-vacuum}.
For the outgoing activity it suffices to compare a
level $k+1$ connection pair with its pullback.
Let $\widetilde U$ be connected and admissible there,
put $Y_i=\operatorname{pr}_i'\widetilde U$,
and set $\widetilde Y=L\widetilde U$.
The latter is a connected union of $LM$-cubes,
hence also of $M$-cubes, admissible at level $k$.
Its neighborhood is contained in $L N(\widetilde U)$.

If $(\mathscr A',\mathscr C')$ satisfies the outgoing-domain hypotheses at level
$k+1$, its scaled pair lies in
$\widetilde R_k^-\times\mathfrak X_k$ by
\eqref{eq:rg-scaled-domain-inclusion}.
Scaling and pullback commute on covering bonds:
both give $L^{-1/2}B'(b/L)$.
Denote the scaled pairs by $(A_i,\mathscr C_i)$.
They satisfy exactly the hypotheses used in
\eqref{eq:rgcov-logarithm-agreement}.

The polymers with $\bar X^L=LY_i$ are precisely
$\operatorname{pr}_i\widetilde X$, where
$\widetilde X$ is connected in $\widetilde Y$ and
its $LM$-closure is $\widetilde Y$.
Indeed $X\subset\bar X^L$ gives its unique lift
inside $\widetilde Y$. Closure commutes with projection,
and lies in $\widetilde Y^{\sim(L-1)}$.
Projection is injective on that enlargement,
so equality of its projected closure with
$\operatorname{pr}_i\widetilde Y$ forces equality
of the lifted closure with $\widetilde Y$.
This also proves the converse and shows that the
reblocking index set is common.

In that index set the outgoing activity is
\begin{align}
 \mathscr E_i^+(Y_i;\mathscr A_i',\mathscr C_i')
  =\sum_{\substack{\widetilde X\subset\widetilde Y\text{ connected}\\
                         \overline{\widetilde X}^{\,L}=\widetilde Y}}
 S_L\Big[&
     (R\mathscr E_i)^{\mathrm{nat}}
                  (\operatorname{pr}_i\widetilde X;A_i,\mathscr C_i)
       +\delta E_k^{\mathrm{det},(i)}
                  (\operatorname{pr}_i\widetilde X;A_i,\mathscr C_i)
       \nonumber\\
       &+E_k^{\#,(i)}
                  (\operatorname{pr}_i\widetilde X;A_i,\mathscr C_i)\Big].
                                      \label{eq:rgcov-reblocked-comparison}
\end{align}
Each of the three entries agrees under
$\iota_{\widetilde Y}$, by extraction and adjustment,
\eqref{eq:rgcov-determinant-agreement}, and
\eqref{eq:rgcov-logarithm-agreement}, respectively.
The entries belong to its ordinary and pair subalgebra.
Equation~\eqref{eq:rgcov-scaling-intertwining} identifies
their scaled values, so the sums agree under
$\iota_{\widetilde U}$.
Locality of the outgoing family, already proved in
the step theorem, then extends the comparison
from pullback pairs to pairs with equal local lifts.

The external scalar recursion has the same reblocking
index set. Its incoming scalar term agrees by hypothesis,
and its local determinant term agrees by
\eqref{eq:rgcov-determinant-agreement} at
$(\mathscr C_1,\mathscr C_2)$ and at $(0,0)$.
Its two sums therefore agree.

The argument without an external slot uses precisely
these local operations with $B_i=A_i$.
The determinant difference is then
$E_k^{\mathrm{det}}(X_i,A_i)=E_k^d(X_i,A_i)-E_k^d(X_i,0)$.
Extraction, addition of odd fields, finite localization
and integration do not require that slot. Their
same domain and convergence assertions were established
for this case in the nonlinear-step construction.
Thus the three bracketed terms in
\eqref{eq:rgcov-reblocked-comparison} agree in that
case too, and scaling gives the conclusion.
Restriction of a family with a slot to
$\mathscr C=0$ also commutes with every operation
in the displayed formula.

At $k=0$ the incoming activity and its extracted
terms are zero. The comparison is for the local
determinant difference and the first correction,
with chamber generators, traces and Wick moments
of weight $1/48$, as explicitly included above.
\end{proof}


\section{The tuned flow}\label{sec:rg-flow}

The mass recursion expands by a factor $L$, while the
normalized activity recursion contracts. We determine
the masses backwards from a zero terminal mass and
the activities forwards from zero initial activity.
This gives a contraction on the whole finite sequence.
The vacuum is then fixed by a scalar sum. Comparing
local steps transfers the pair chosen on the top torus
to every smaller torus used in the construction.

\subsection{Physical scales and stopping levels}
\label{subsec:rgfl-scales}

Fix $\bar m\in(0,1]$ and put
$h=L^{\lfloor\log_L(1/\bar m)\rfloor}$.
For a cutoff $N\ge0$, define
\begin{equation}\label{eq:rgfl-physical-scales}
 \eta=L^{-N},\qquad K=N+\log_Lh,\qquad
 a_j=L^{j-N},\qquad e_j=e\sqrt{a_j},\qquad
 \mu_j=\bar m a_j.
\end{equation}
Thus $a_j$ is the physical length of one unit at
level $j$, and $a_K=h$.
At zero the charge is written $q=e\sqrt\eta$;
the fixed charge ceilings $e_0,e_1$ retain their
earlier meanings. A formula with variable index $j$
uses $e_j$ for the running charge, including $q$
when $j=0$.
For $0\le j\le K$ set
\[
 v_j=-\log e_j,\qquad
 r_j=\min\{L^i:i\in\N\cup\{0\},\ L^i\ge v_j^r\},
 \qquad p_j=v_j^p,\quad p_{0,j}=v_j^{p_0}.
\]
The integers $r,p,p_0$ and the constants of the
preceding estimates are fixed.
Here $r\ge2$, $L$ is odd with $L\ge13$, and $M$
is a power of $L$ with $M\ge L^2$.
We assume the hypotheses of the primitive and
step estimates, and use a ceiling
$e_1\le\min(e_0,1/3)$ such that, for all
$v\ge-\log e_1$,
\begin{equation}\label{eq:rgfl-scale-smallness}
 \begin{gathered}
 v\ge8,\qquad
 \left(1+\frac{3\log L}{2v}\right)^r<L,\qquad
 v^r>31L+16,\\
 (c_{\mathrm{fl}}+2)\exp(-\varepsilon v)
               \le\frac1{16}(1-L^{-\varepsilon/2}).
 \end{gathered}
\end{equation}
In addition
$0<\varepsilon\le1/200$,
$c_{\mathrm{fl}}\ge1$,
$e_0^\varepsilon\le1/4$,
$c_{\mathrm{fl}}L^{-3/8+6/200}\le1/16$, and
$L^{-5/8}\le1/4$.
The construction below applies whenever
$0<e\sqrt h\le e_1$.

The physical top side and reference side are
\begin{equation}\label{eq:rgfl-reference-sides}
 R_{\mathrm{top}}=hMr_K,\qquad
 R^*=R_{\mathrm{top}}/L=hL^{T^*(e)},\qquad
 T^*(e)=\log_L(Mr_K)-1.
\end{equation}
They are independent of $N$, since
$v_K=-\log(e\sqrt h)$.
The admissible physical sides for the flow are
\[
 \mathcal T(e)=\{hL^T:T\in\N\cup\{0\},\ hL^T\le R_{\mathrm{top}}\}.
\]
For $R\in\mathcal T(e)$ put
$\ell_j(R)=R/a_j$,
$T_j(R)=\mathbb L_j/\ell_j(R)\Z^3$ and
$\mathfrak n_j=\log_L(\ell_j/(Mr_j))$.
The allowed unit sides are
$\mathrm{Cl}=\{\ell:\ell/(LM)\text{ is an odd integer at least }13\}$.
Their occurrence at every required scale is part
of the following arithmetic argument.

\begin{proposition}\label{prop:rgfl-scales}
For each $R\in\mathcal T(e)$ there is a finite integer
$N_{\min}(R,e)\ge0$ such that, when $N\ge N_{\min}(R,e)$,
the first level with $\mathfrak n_j=0$ exists and
satisfies $1\le k_*(R)\le K$.
Its physical quantities $a_*,e_*,\mu_*,\ell_*,r_*$
are independent of $N$ in this range.
For $j\le k_*(R)$,
\begin{equation}\label{eq:rgfl-scale-ranges}
 \begin{gathered}
 \ell_j\in\mathrm{Cl},\qquad
 \ell_j\ge Mr_j,\qquad r_j\ge r_*>31L+16,\\
 v_j\ge v_*\ge v_K,\qquad
 0<e_j\le e_*\le e_K\le e_1,\qquad
 0<\mu_j\le\mu_*\le\bar m h\le1 .
 \end{gathered}
\end{equation}
For $j<k_*(R)$ one has $\ell_j\ge LMr_j$,
$\ell_{j+1}=\ell_j/L\in\mathrm{Cl}$,
$e_{j+1}=\sqrt L e_j\le e_0$ and
$0<\mu_j=\mu_{j+1}/L\le1/L$.
At the stop,
\begin{equation}\label{eq:rgfl-terminal-scales}
 \ell_*=Mr_*,\qquad
 Mv_*^r\le\ell_*<LMv_*^r,\qquad
 \mu_*^{-1}<L\ell_*,\qquad e_*^2\ell_*=e^2R.
\end{equation}
If $R\le R'$ in $\mathcal T(e)$, then
$N_{\min}(R',e)\le N_{\min}(R,e)$ and
$k_*(R)\le k_*(R')$ for a common cutoff above both stated thresholds.
\end{proposition}

\begin{proof}
Use the depth $s=K-j$. Its quantities are
\[
 \widehat v_s=v_K+\tfrac{s}{2}\log L,\qquad
 \widehat r_s=\min\{L^i:i\ge0,\ L^i\ge\widehat v_s^r\},\qquad
 \widehat n_s(R)=T+s-\log_L(M\widehat r_s).
\]
They do not depend on $N$. At depth $s=K-j$ they
equal $v_j,r_j,\mathfrak n_j$, while
$a_j=hL^{-s}$ and $\ell_j=L^{T+s}$.
Since $\widehat v_s\ge v_K$, the scalar condition gives
\[
 \widehat v_{s+1}^r
   =\widehat v_s^r\left(1+\frac{\log L}{2\widehat v_s}\right)^r
                    \le L\widehat v_s^r\le L\widehat r_s.
\]
Thus
$\widehat r_{s+1}\in\{\widehat r_s,L\widehat r_s\}$
and $\widehat n_{s+1}-\widehat n_s\in\{0,1\}$.
Also $\widehat n_0\le0$ because $R\le R_{\mathrm{top}}$.
Minimality of the rounded power gives
$\widehat r_s<L\widehat v_s^r$ and hence
\[
 \widehat n_s>T+s-1-\log_LM-r\log_L(v_K+\tfrac{s}{2}\log L)
                                                    \longrightarrow\infty .
\]
The integer sequence therefore has a largest zero.
Denote its index by $s_*(R,e)$.
It is nonpositive up to that index and at least $1$
afterwards.

The condition $\ell_0\ge LMr_0$ is
$\widehat n_K\ge1$, equivalently $K>s_*$.
Its exact threshold is
\begin{equation}\label{eq:rgfl-mesh-threshold}
 N_{\min}(R,e)=\max\{0,s_*(R,e)+1-\log_Lh\},
             \qquad k_*(R)=K-s_*(R,e).
\end{equation}
This gives $1\le k_*\le K$ and
$a_*=hL^{-s_*}$, proving cutoff independence of
the terminal quantities.
Increasing $T$ increases every $\widehat n_s$
by the same integer; thus it decreases $s_*$,
decreases $N_{\min}$ and increases $k_*$.

For $j\le k_*$, the depth satisfies $K-j\ge s_*$.
Monotonicity gives the asserted ranges of
$\mathfrak n_j,r_j,v_j,e_j$, and the running mass
increases with $j$.
Moreover
$\bar m h=L^{\lfloor\log_L(1/\bar m)\rfloor-\log_L(1/\bar m)}
\in(1/L,1]$.
Since $\mu_*\ell_*=\bar mR\ge\bar m h>1/L$,
we obtain $\mu_*^{-1}<L\ell_*$.
The rounded-power inequality at $v_*$ gives
the two bounds for $\ell_*$ in
\eqref{eq:rgfl-terminal-scales}; its last identity
uses $e_*^2=e^2a_*$.

Finally $r_*\ge v_*^r>31L+16>L$ implies
$r_*\ge L^2$, as it is a power of $L$.
At every required level,
$\ell_j/(LM)=r_jL^{\mathfrak n_j-1}$
is a positive power of the odd integer $L\ge13$.
This proves $\ell_j\in\mathrm{Cl}$.
Before the stop $\mathfrak n_j\ge1$, giving the
stronger side bound. The mass and charge bounds
for the step follow at $j+1$.
\end{proof}

On $R^*$ define the reference mesh threshold
\begin{equation}\label{eq:rgfl-reference-threshold}
                    N_{\mathrm{ref}}(e)=N_{\min}(R^*(e),e).
\end{equation}
The stop on this torus has depth $s_*\in\{1,2\}$.
Indeed, for $s=1,2,3$ the stronger inequality in
\eqref{eq:rgfl-scale-smallness} gives
$\widehat r_s\in\{r_K,Lr_K\}$.
On $R^*$,
$\widehat n_s=s-1-\log_L(\widehat r_s/r_K)$.
Its values at $s=0,1,2,3$ are respectively
$-1$, an element of $\{-1,0\}$, an element of
$\{0,1\}$, and an element of $\{1,2\}$.
The zero-step-or-one-step rule forces the last zero
to be $1$ or $2$. On $R_{\mathrm{top}}$ the sequence
is larger by one, so its last zero is $0$ or $1$.
Consequently
\begin{equation}\label{eq:rgfl-reference-stop}
 \begin{gathered}
 k_*(R^*)\in\{K-1,K-2\},\qquad
 a_*\in\{h/L,h/L^2\},\qquad r_*\in\{r_K,Lr_K\},\\
 k_{\mathrm{top}}\in\{K,K-1\},\qquad
 L^{-3}<\mu_*\le L^{-1}\quad\text{on }R^* .
 \end{gathered}
\end{equation}
All these physical terminal values are independent
of $N\ge N_{\mathrm{ref}}(e)$.

The scale proposition verifies every side, mass
and charge hypothesis of the step theorem
on $R_{\mathrm{top}}$ and on the smaller tori it covers with odd degree.
At each level use the top system of
Section~\ref{subsec:rgcov-top-system}; on a smaller
torus use its projection from
Proposition~\ref{prop:rgcov-projected-system}.
These are the actual mean-field, covariance,
local resolvent and photon families constructed
earlier, with the same common estimates.
The ratio $R_{\mathrm{top}}/R$ is a power of the
odd integer $L$, and $k_*(R)\le k_{\mathrm{top}}$.
Thus the required top families exist whenever
a smaller-volume step is used.
Furthermore
$\ell_{j+1}(R)\ge Mr_{j+1}>(31L+16)M$;
the next-side hypothesis of the cover comparison
holds at every such step.

\subsection{Continuation within the mass range}
\label{subsec:rgfl-continuation}

Write the mass and activity bounds as
\[
 w_j^{\mathrm m}=e_j^{3/4-8\varepsilon},\qquad
 w_j^{\mathrm E}=e_j^{1/4-7\varepsilon}.
\]
For a real pair $(m_{\mathrm{in}},\varepsilon_{\mathrm{in}}^{\mathrm v})$,
the initial vacuum is $\varepsilon_{\mathrm{in}}^{\mathrm v}$,
the initial member of the mass sequence is $m_{\mathrm{in}}$,
and $E_0=\mathscr E_0=\mathfrak P_0=0$.
The variable-index notation $m_j$ includes this initial member.
The recursion is
\begin{align}
 m_{j+1}&=L(m_j+m(E_j)),&
 \varepsilon_{j+1}^{\mathrm v}
       &=L^3(\varepsilon_j^{\mathrm v}
                 +\varepsilon_{\mathrm v}(E_j)+\delta\varepsilon_j^0),
                                      \label{eq:rgfl-coupling-recursion}\\
 E_{j+1}&=\mathcal L(RE_j)^{\mathrm{nat}}
                          +\mathcal LE_j^{\mathrm{det}}+E_j^*(m_j,E_j),
                                      \label{eq:rgfl-activity-recursion}\\
 \mathscr E_{j+1}
       &=\mathcal L(R\mathscr E_j)^{\mathrm{nat}}
               +\mathcal L\delta E_j^{\mathrm{det}}
                                      +E_j^*(m_j,\mathscr E_j).
                                      \label{eq:rgfl-slot-recursion}
\end{align}
The scalar family $\mathfrak P_j$ follows
\eqref{eq:rgnl-external-scalar-map}.
In the last line the correction has the external slot;
in \eqref{eq:rgfl-activity-recursion} it has none.
These are the corresponding maps proved earlier.
The trajectory exists up to a level if every preceding
$(m_j,\mathscr E_j)$ lies in $\operatorname{Ball}_j^x(4)$.
It is controlled when it exists to $k_*(R)$ and
$|m_j|\le w_j^{\mathrm m}$ throughout, including the stop.

The normalized estimates used below are
\begin{equation}\label{eq:rgfl-window-estimates}
 \begin{gathered}
 L^{-1}\frac{w_{j+1}^{\mathrm m}}{w_j^{\mathrm m}}
                  =L^{-5/8-4\varepsilon}\le\tfrac14,\qquad
 c_{\mathrm{fl}}e_j^{1/2}
          \frac{w_j^{\mathrm E}}{w_j^{\mathrm m}}
                           =c_{\mathrm{fl}}e_j^\varepsilon,\\
 c_{\mathrm{fl}}L^{-1/4+2\varepsilon}
          \frac{w_j^{\mathrm E}}{w_{j+1}^{\mathrm E}}
                     =c_{\mathrm{fl}}L^{-3/8+11\varepsilon/2}
                                                        \le\tfrac1{16},\\
 \frac{e_{j+1}^{1/4-2\varepsilon}}{w_{j+1}^{\mathrm E}}
                                      \le e_{j+1}^\varepsilon,\qquad
 \frac{e_j^{1/4-6\varepsilon}}{w_{j+1}^{\mathrm E}}
             =L^{-1/8+3\varepsilon}e_{j+1}^\varepsilon
                                      \le e_{j+1}^\varepsilon .
 \end{gathered}
\end{equation}
Their first and third bounds use the fixed inequalities
stated above; their last two use $e_{j+1}<1$ and
$\varepsilon\le1/200$.
The charge condition also gives
\begin{equation}\label{eq:rgfl-summable-charge}
 \sum_{j=0}^{k_*}(c_{\mathrm{fl}}+2)e_j^\varepsilon
             \le(c_{\mathrm{fl}}+2)e_*^\varepsilon
                          (1-L^{-\varepsilon/2})^{-1}
                                                        \le\tfrac1{16}.
\end{equation}
In particular
$c_{\mathrm{fl}}e_j^\varepsilon\le1/16$,
$2e_j^\varepsilon\le1/16$, and
$e_j\le e_j^\varepsilon\le1/48$.

\begin{proposition}\label{prop:rgfl-continuation}
If the recursively produced masses satisfy
$|m_i|\le w_i^{\mathrm m}$ for $i<j\le k_*(R)$,
then the trajectory is defined to $j$ and
\begin{equation}\label{eq:rgfl-activity-bound}
 \|E_j\|_j\le\|\mathscr E_j\|_j^x
       \le\frac{32}{15}e_j^\varepsilon w_j^{\mathrm E}
       \le4e_j^{1/4-6\varepsilon}\le w_j^{\mathrm E}.
\end{equation}
The data satisfy
$E_j=\mathscr E_j|_{\mathscr C=0}$,
$\mathfrak P_j(X,0)=0$, and belong to their real
mass, vacuum and activity spaces.
\end{proposition}

\begin{proof}
At zero all activities and scalar families vanish.
Assume the assertion through $i<j$.
The mass condition and
$4e_i^\varepsilon\le1$ put the input in the unit ball,
so Theorem~\ref{thm:rgnl-current-step} applies at $i$.
The scaled connection pair lies in
$\widetilde R_i^-\times\mathfrak X_i$, so every local
determinant and correction in the recursion is defined.
The scalar external sum is finite and uses
$\mathscr C_L\in\mathfrak X_i$.
The step theorem preserves the real spaces
and commutes with restriction to $\mathscr C=0$.
Consequently the coupling maps computed from
$E_i$ or $\mathscr E_i$ coincide.
The scalar recursion vanishes at zero because each
$E_i^{\mathrm{det}}(X,0)$ vanishes.

The activity bound of the current step, divided by
$w_{i+1}^{\mathrm E}$ and combined with
\eqref{eq:rgfl-window-estimates}, gives
\[
 \frac{\|\mathscr E_{i+1}\|_{i+1}^x}{w_{i+1}^{\mathrm E}}
      \le\frac1{16}
                   \frac{\|\mathscr E_i\|_i^x}{w_i^{\mathrm E}}
                                     +2e_{i+1}^\varepsilon.
\]
Starting from zero, the bound
$(32/15)e_i^\varepsilon$ is preserved since
$e_i\le e_{i+1}$ and $(32/15)/16+2=32/15$.
Restriction does not increase the activity norm.
This proves \eqref{eq:rgfl-activity-bound} and the induction.

At the first step the formulas use $192\ell_0^3$
incoming pairs with weight $1/48$ and
$4\ell_1^3$ outgoing pairs, as proved in the
exact density identity.
Its coupling maps are
$m_1=Lm_{\mathrm{in}}$ and
$\varepsilon_1^{\mathrm v}
 =L^3(\varepsilon_{\mathrm{in}}^{\mathrm v}+\delta\varepsilon_0^0)$;
its activity is the local determinant term plus
the first correction. The normalized estimate has
zero incoming-activity term and the same two
source bounds. Thus the induction includes that step.
\end{proof}

\subsection{The terminal conditions and the fixed point}
\label{subsec:rgfl-fixed-point}

\begin{theorem}\label{thm:rgfl-terminal-zero-pair}
On any $R\in\mathcal T(e)$ with $N\ge N_{\min}(R,e)$,
there is exactly one real controlled pair whose
trajectory has $m_{k_*}=0$ and
$\varepsilon_{k_*}^{\mathrm v}=0$.
Along its trajectory,
\begin{equation}\label{eq:rgfl-tuned-bounds}
 |m_j|\le\tfrac18w_j^{\mathrm m},\qquad
 \|E_j\|_j\le\|\mathscr E_j\|_j^x
                                  \le4e_j^{1/4-6\varepsilon},\qquad
 |\varepsilon_j^{\mathrm v}|
                          \le5c_{\mathrm{fl}}e_j^{1/4-6\varepsilon}.
\end{equation}
The vacuum at every level is
\begin{equation}\label{eq:rgfl-backward-vacuum}
 \varepsilon_j^{\mathrm v}
     =-\sum_{i=j}^{k_*-1}L^{-3(i-j)}
                      \big(\varepsilon_{\mathrm v}(E_i)
                                              +\delta\varepsilon_i^0\big).
\end{equation}
\end{theorem}

\begin{proof}
Take the finite product of the real mass coordinates
and real activity Banach spaces, with norm
\[
 \|\mathbf z\|_{\mathrm{fl}}
   =\max_{0\le j\le k_*}
       \max\left(\frac{|m_j|}{w_j^{\mathrm m}},
                              \frac{\|E_j\|_j}{w_j^{\mathrm E}}\right),
                  \qquad \mathbf z=(m_j,E_j)_{0\le j\le k_*}.
\]
The closed subset $\mathcal B$ defined by these
quotients being at most $1$ and $E_0=0$
is nonempty and complete. Define a self-map by
\begin{align}
 \widehat m_j&=L^{-1}m_{j+1}-m(E_j)\quad(j<k_*),
                 &\widehat m_{k_*}&=0, \nonumber\\
 \widehat E_0&=0,\qquad
 \widehat E_{j+1}
       =\mathcal L(RE_j)^{\mathrm{nat}}
                    +\mathcal LE_j^{\mathrm{det}}+E_j^*(m_j,E_j).
                                      \label{eq:rgfl-fixed-point-map}
\end{align}
All maps here are the actual maps without an external
slot; their inputs belong to the unit ball.
Extraction, adjustment followed by reblocking,
the determinant activity and the nonlinear correction
are defined on that ball and preserve the real space.

We verify that the map preserves $\mathcal B$.
The extraction estimate implies
$|m(E_j)|\le c_{\mathrm{fl}}e_j^{1/2}\|E_j\|_j$.
Equations~\eqref{eq:rgfl-window-estimates} and
\eqref{eq:rgfl-summable-charge} give
\[
 \frac{|\widehat m_j|}{w_j^{\mathrm m}}
                    \le\tfrac14+c_{\mathrm{fl}}e_j^\varepsilon
                    \le\tfrac5{16}.
\]
For the activity, the contraction, determinant and
correction estimates give
\[
 \frac{\|\widehat E_{j+1}\|_{j+1}}{w_{j+1}^{\mathrm E}}
                         \le\tfrac1{16}+2e_{j+1}^\varepsilon
                         \le\tfrac18 .
\]
The fixed end components are zero. Hence the map
takes $\mathcal B$ into itself.

For two sequences, let $d_j^{\mathrm m},d_j^{\mathrm E}$
be the corresponding normalized differences.
Linearity of extraction gives
\[
 \widehat d_j^{\mathrm m}
      \le\tfrac14d_{j+1}^{\mathrm m}
                            +c_{\mathrm{fl}}e_j^\varepsilon d_j^{\mathrm E}
      \le\tfrac5{16}\|\mathbf z-\widetilde{\mathbf z}\|_{\mathrm{fl}} .
\]
Adjustment followed by reblocking is linear in
the activity at fixed connection.
The local determinant term is independent of the
sequence. The correction's coupling estimate
\eqref{eq:rgnl-coupling-Lipschitz} therefore gives
\[
 \widehat d_{j+1}^{\mathrm E}
      \le\tfrac1{16}d_j^{\mathrm E}
                         +e_{j+1}^\varepsilon
                                      (d_j^{\mathrm m}+d_j^{\mathrm E})
      \le\tfrac18\|\mathbf z-\widetilde{\mathbf z}\|_{\mathrm{fl}} .
\]
The whole map contracts by $5/16$.
Banach's fixed point theorem
\cite[Theorem~8]{BrightMetric2023} gives a unique
fixed point in $\mathcal B$.
Its equations are precisely the mass and activity
recursions with $m_{k_*}=0$.

Let $m_{\mathrm{in}}$ be its initial mass coordinate
and choose
\[
 \varepsilon_{\mathrm{in}}^{\mathrm v}
      =-\sum_{i=0}^{k_*-1}L^{-3i}
                     \big(\varepsilon_{\mathrm v}(E_i)
                                           +\delta\varepsilon_i^0\big).
\]
This pair is real. Its full trajectory, including
the external slot, agrees with the fixed-point mass
and activity at every level: agreement at zero is
the initial condition; at the next level
Proposition~\ref{prop:rgfl-continuation} makes the
full step defined, and the fixed-point equations
identify its mass and restricted activity.
The induction reaches $k_*$ and proves control.
Solving the scalar vacuum recursion backwards gives
\[
 \varepsilon_j^{\mathrm v}
 =L^{-3(k_*-j)}\varepsilon_{k_*}^{\mathrm v}
     -\sum_{i=j}^{k_*-1}L^{-3(i-j)}
                     \big(\varepsilon_{\mathrm v}(E_i)
                                           +\delta\varepsilon_i^0\big).
\]
At $j=0$ its comparison with the chosen input forces
$\varepsilon_{k_*}^{\mathrm v}=0$, giving
\eqref{eq:rgfl-backward-vacuum}.
Conversely any controlled pair with the two terminal
zeros has its mass/activity sequence in $\mathcal B$
by the continuation proposition. It is the same fixed
point. Its initial vacuum is then the same scalar
sum. This proves uniqueness of the pair.

For the sharper mass bound put
$B_j=|m_j|/w_j^{\mathrm m}$.
The backward mass equation and the unit activity
bound imply $B_j\le B_{j+1}/4+1/16$, while
$B_{k_*}=0$.
Summing gives $B_j\le(1/16)\sum_{s\ge0}4^{-s}
=1/12\le1/8$.
The activity bound in \eqref{eq:rgfl-tuned-bounds}
is the continuation estimate.

It remains to bound the vacuum sum.
Extraction gives
$|\varepsilon_{\mathrm v}(E_i)|
\le4c_{\mathrm{fl}}e_i^{1/4-6\varepsilon}$.
For $i\ge j$ the associated summand obeys
\[
 L^{-3(i-j)}e_i^{1/4-6\varepsilon}
       =e_j^{1/4-6\varepsilon}
                         L^{-(i-j)(23/8+3\varepsilon)}
       \le e_j^{1/4-6\varepsilon}16^{-(i-j)}.
\]
Thus its sum is at most
$(64/15)c_{\mathrm{fl}}e_j^{1/4-6\varepsilon}$.
The box normalization satisfies
$0\le\delta\varepsilon_i^0\le e_i^7$ by the
step estimate.
Its remaining sum is bounded by
\[
 \sum_{i=j}^{k_*-1}L^{-3(i-j)}e_i^7
        =e_j^6\sum_{i=j}^{k_*-1}e_i
        \le e_j^6 e_*(\sqrt L-1)^{-1}
        \le\tfrac1{48}e_j^{1/4-6\varepsilon}.
\]
Here $\sqrt L\ge2$ and $e_*\le1/48$.
Since $c_{\mathrm{fl}}\ge1$,
$64/15+1/48=343/80<5$, giving the claimed bound.
The sums are empty at the stop.

At zero $E_0=0$. The backward mass equation reads
$m_{\mathrm{in}}=L^{-1}m_1$, and the forward activity is
$\mathcal LE_0^{\mathrm{det}}+E_0^*(m_{\mathrm{in}},0)$.
The determinant and correction estimates used above
include their full chamber trace and the first
weighted fluctuation integral. The initial vacuum
sum contains $\delta\varepsilon_0^0$.
No additional coupling is selected by the fixed point.
\end{proof}

\subsection{One pair on all tori of the construction}
\label{subsec:rgfl-volume-comparison}

For $N\ge N_{\min}(R_{\mathrm{top}},e)$ define
\[
 P_N(e)=\big(m_{0,N}(e),\varepsilon_{0,N}^{\mathrm v}(e)\big)
\]
to be the pair of Theorem~\ref{thm:rgfl-terminal-zero-pair}
on the prescribed side $R_{\mathrm{top}}(e)$.
The dependence on the fixed $\bar m$ and fixed
construction parameters is suppressed.
For lower cutoffs set $m_{0,N}(e)=0$.

\begin{theorem}\label{thm:rgfl-common-pair}
The pair $P_N(e)$ is controlled on every
$R\in\mathcal T(e)$ with $N\ge N_{\min}(R,e)$.
Up to its stop, its mass and vacuum sequence agree
with those on the top torus; its activity and
external scalar families agree there on flat polymers.
For every such level,
\begin{equation}\label{eq:rgfl-common-bounds}
 |m_j|\le\tfrac18e_j^{3/4-8\varepsilon},\qquad
 \|E_j\|_j\le\|\mathscr E_j\|_j^x
                             \le4e_j^{1/4-6\varepsilon},\qquad
 |\varepsilon_j^{\mathrm v}|
                         \le5c_{\mathrm{fl}}e_j^{1/4-6\varepsilon}.
\end{equation}
The real mass sequence is independent of the side
datum and of the twist. At all $N\ge0$,
\begin{equation}\label{eq:rgfl-bare-mass-bound}
 |m_{0,N}(e)|\le\tfrac18q^{3/4-8\varepsilon}<\tfrac18,
                       \qquad q=eL^{-N/2}.
\end{equation}
The total lattice mass is
$\boldsymbol\mu_N=\bar m\eta+m_{0,N}(e)$.
\end{theorem}

\begin{proof}
The pair is defined because
$N_{\min}(R_{\mathrm{top}},e)\le N_{\min}(R,e)$.
The two trajectories have the same initial
couplings and zero initial activities.
Suppose they agree through $j<k_*(R)$.
The top mass is bounded by $w_j^{\mathrm m}/8$;
the same bound on the smaller side permits continuation
to $j+1$ by Proposition~\ref{prop:rgfl-continuation}.
Both inputs are now in the unit balls, are real,
and have locally agreeing activity and scalar families.
The scale proposition verifies their primitive,
domain and next-side comparison hypotheses.
Theorem~\ref{thm:rgcov-step-comparison} gives the
same extracted mass and vacuum, including the same
$\delta\varepsilon_j^0$, and local agreement of the
outgoing activity and scalar family.
The scalar recursions therefore give the same
$m_{j+1},\varepsilon_{j+1}^{\mathrm v}$.
Restriction to zero external slot preserves agreement
of $E_{j+1}$.
This proves the induction up to $k_*(R)$.

The common mass bound follows from the top trajectory.
The local comparison alone is not an estimate of
the entire activity on the smaller torus; that
bound follows from its own continuation estimate.
The vacuum is a common scalar and satisfies the
top bound at this same level. Thus
\eqref{eq:rgfl-common-bounds} follows.
At a smaller-volume stop the couplings need not
vanish; the zeros defining $P_N(e)$ are imposed
at $k_{\mathrm{top}}$.

The fixed-point map on the top torus is determined
by $(N,e,\bar m)$, the fixed construction parameters
and photon coordinate basis, and the top primitive
families constructed from them.
The selected top side is itself a function of
$e,\bar m$ and these parameters.
The smaller side $R$ enters no part of this definition.
Nor does a twist: the twist is a harmonic part
of a connection inserted in the density identity,
whereas the coupling maps are maps of entire
local activity families and contain no such insertion.
The full forward activity recursion retains the external
slot as a variable, whereas the fixed-point map uses its
restriction to zero. No external connection or probe is
inserted when selecting the pair.
This proves the asserted independence within
the fixed construction.

For a cutoff on which the pair is defined,
\eqref{eq:rgfl-bare-mass-bound} is its initial mass
bound on the top torus. At the other finitely many
cutoffs it holds because the assigned mass is zero.
Finally $h\ge1$ implies $e\le e_1\le1/3$,
so $q\le1/3$ and its positive power is below $1$.
The sequence is real in both cases.
The sign of $\boldsymbol\mu_N$ is not fixed by this
estimate; the lattice construction admits a real mass.
\end{proof}

On the reference torus, for $N\ge N_{\mathrm{ref}}(e)$,
the outputs needed at the terminal scale are
\begin{equation}\label{eq:rgfl-reference-terminal-input}
 a_*\in\{h/L,h/L^2\},\qquad
 |m_*|\le\tfrac18e_*^{3/4-8\varepsilon},\qquad
 \|E_*\|_*\le\|\mathscr E_*\|_*^x
                                 \le4e_*^{1/4-6\varepsilon}.
\end{equation}
The primitive estimates along this flow have the
same fixed constants on every torus and at every
cutoff. The remaining terminal density and partition
estimates concern the integration of these fields.


\section{A positive continuum fourth coefficient}\label{sec:positive-probe}

The fourth-order response calculation produces the coefficient
$-6\tr[(S_{\bar m}i\gamma\cdot d^*\beta)^4]$. Its sign is not apparent
from this expression. We determine the leading term at large mass and
then dilate a compactly supported probe. The large-mass calculation is
local: a gauge choice removes the constant connection, and quartic
homogeneity makes the leading term depend only on the antisymmetric
first derivative. We compute that coefficient in a constant field.

For a real $a\in C_c^\infty(\R^3;\R^3)$ put
\[
 V_a=i\gamma\cdot a,\qquad S_M=(M+\gamma\cdot\partial)^{-1},\qquad
 q_{4,M}(a)=-\tfrac14\tr[(S_MV_a)^4],\quad M>0.
\]
The trace is on $L^2(\R^3;\C^4)$; $V_a$ denotes multiplication.
We use Euclidean norms in this section.

\begin{proposition}\label{prop:positive-probe}
The operator $(S_MV_a)^4$ is trace class and its trace is real. There
is $C_a<\infty$ such that, for $M\ge1$,
\begin{equation}\label{eq:large-mass-quartic}
 \left|q_{4,M}(a)-\frac1{240\pi M^5}
       \int_{\R^3}|\operatorname{curl}a|^4\,\dd x\right|
 \le C_aM^{-7}.
\end{equation}
In particular, for every $\bar m>0$ there is a real
$\beta\in C_c^\infty(\R^3;\Lambda^2)$ for which $f=dd^*\beta$ satisfies
$b(f)>0$ in \eqref{eq:positive-coefficient}.
\end{proposition}

\subsection{Kernels and differentiation in the mass}

The kernel and Fourier multiplier of $S_M$ are
\begin{align}
 S_M(z)&=(M-\gamma\cdot\partial)
              \frac{e^{-M|z|}}{4\pi|z|},\label{eq:probe-free-kernel}\\
 \widehat S_M(p)&=\frac{M-i\gamma\cdot p}{M^2+|p|^2},\qquad
 |\widehat S_M(p)|=(M^2+|p|^2)^{-1/2}.
 \label{eq:probe-free-symbol}
\end{align}
In particular,
\[
 |S_M(z)|\le\frac{e^{-M|z|}(2M|z|+1)}{4\pi|z|^2},\qquad
 S_M(z)=M^2S_1(Mz),\qquad \|S_M\|_{L^1}\le3/M.
\]
The operator $G=S_MV_aS_M$ is Hilbert--Schmidt. Indeed, its momentum
kernel, with measure $(2\pi)^{-3}\dd p$, is
$\widehat S_M(p)i\gamma\cdot\widehat a(p-q)\widehat S_M(q)$.
Writing $u(p)=(M^2+|p|^2)^{-1}$, the matrix norm bound and Young's
inequality give
\[
 \|G\|_{\mathrm{HS}}^2
 \le8(2\pi)^{-6}\int u(p)(|\widehat a|^2*u)(p)\,\dd p
 \le8(2\pi)^{-3}\|u\|_2^2\|a\|_2^2
 =\frac{\|a\|_2^2}{\pi M}.
\]
Here $\|u\|_2^2=\pi^2/M$. Since $(S_MV_a)^4=GV_aGV_a$,
\begin{equation}\label{eq:probe-trace-bound}
 |q_{4,M}(a)|\le\frac{\|a\|_2^2\|a\|_\infty^2}{4\pi M}.
\end{equation}
The matrix $\gamma_5=\left(\begin{smallmatrix}0&1_2\\1_2&0\end{smallmatrix}\right)$
conjugates $S_M$ to $S_M^*$ and $V_a$ to $V_a^*=-V_a$.
Conjugation of the fourth power and cyclicity of the trace therefore
give a real trace. Cyclicity here involves a bounded factor and the
trace-class product $GV_aG$, so it does not presume that $S_MV_a$
itself is trace class.

We will evaluate the mass derivative on the diagonal of a kernel with
five propagators. This requires continuity at the diagonal, including
for the linear connections used below. For either a compact smooth
connection or a linear connection, $|a(z)|\le c_a(1+|z|)$ and
$a$ is Lipschitz. Choose
$s(z)=K|z|^{-2}e^{-M|z|/2}$ so that
$|S_M(z)|(1+|z|)^{24}\le s(z)$. Its convolution powers obey
\begin{equation}\label{eq:probe-convolutions}
 s^{*2}(z)\le C|z|^{-1}e^{-M|z|/2},\quad
 s^{*2}(z)\le C|z|^{-3/2},\quad
 s^{*3}(z)\le C|z|^{-1/2}.
\end{equation}
For example, the bound for the second convolution follows by factoring
out $e^{-M|z|/2}$ and scaling
$\int|z-w|^{-2}|w|^{-2}\,\dd w=C|z|^{-1}$; the third follows by
convolving $|z|^{-2}$ and $|z|^{-3/2}$. Thus $s^{*2}\in L^2$,
$s^{*4}=s^{*2}*s^{*2}$ is bounded and continuous, and so is
$s^{*5}=s*s^{*4}$.

Let $R_a^{(j)}(x,y)$ be the integral kernel of
$(-1)^j(S_MV_a)^jS_M$, initially off the diagonal, for $0\le j\le4$.
Distributing the factors $1+|z_l|$ along the chain of propagators gives
\begin{equation}\label{eq:probe-kernel-majorant}
 |R_a^{(j)}(x,y)|\le[c_a(1+|x|)]^j s^{*(j+1)}(x-y).
\end{equation}
These integrals act as the indicated operators on Schwartz spinors;
absolute integration against such a spinor justifies every composition.
The fourth kernel extends continuously to all $(x,y)$. To see this,
use the four independent increments of the first three and last
propagators. Changing $(x,y)$ then changes four Lipschitz factors and
the middle propagator. On $|x|,|y|,|x'|,|y'|\le r$, the resulting
difference is bounded by
\begin{multline}\label{eq:probe-kernel-continuity}
 C_{a,r}\bigl((|x-x'|+|y-y'|)\|s^{*5}\|_\infty\\
 +\|s^{*4}\|_\infty
       \|S_M-S_M(\,\cdot+(x'-y')-(x-y))\|_{L^1}\bigr).
\end{multline}
This tends to zero. Define
$r_4(x;M,a)=\tr R_a^{(4)}(x,x)$ using this extension.

For compact $a$, differentiation in $M$ now gives
\begin{equation}\label{eq:probe-mass-derivative}
 \frac{\dd}{\dd M}q_{4,M}(a)=\int r_4(x;M,a)\,\dd x,
 \qquad q_{4,M}(a)\longrightarrow0\quad(M\longrightarrow\infty).
\end{equation}
To justify this differentiation, write the trace as the absolutely convergent
four-vertex integral
\[
 -\tfrac14\int\tr\bigl[S_M(z_4-z_1)V_a(z_1)
 S_M(z_1-z_2)V_a(z_2)S_M(z_2-z_3)V_a(z_3)
 S_M(z_3-z_4)V_a(z_4)\bigr]\,\dd z_1\cdots\dd z_4.
\]
Its absolute integral is at most
$4\|a\|_\infty^3\|a\|_{L^1}s^{*4}(0)$ before the factor $1/4$.
On any compact interval of positive masses, $S_M$ and
$\partial_MS_M$ have a common majorant of the same form. Since
$\partial_MS_M=-S_M*S_M$, the four differentiated terms have equal
integrals by cyclicity. The factor $-1/4$ and the minus sign in this
derivative cancel. Inserting the convolution yields the diagonal
five-propagator integral in \eqref{eq:probe-mass-derivative}; its absolute
integral is bounded by $4\|a\|_\infty^3\|a\|_{L^1}s^{*5}(0)$.
This also proves continuous differentiability. The limit at infinity
follows from \eqref{eq:probe-trace-bound}.

\subsection{The local large-mass expansion}

Write
\[
 V_j(x;w)=i\gamma\cdot\sum_{|\alpha|=j}
                  \frac{w^\alpha}{\alpha!}\partial^\alpha a(x).
\]
The substitution $z_l=x+M^{-1}w_l$ in the diagonal kernel gives
\begin{multline}\label{eq:probe-scaled-kernel}
 r_4(x;M,a)=M^{-2}\tr\int S_1(-w_1)V_a(x+M^{-1}w_1)
 S_1(w_1-w_2)V_a(x+M^{-1}w_2)\\
 {}\cdot S_1(w_2-w_3)V_a(x+M^{-1}w_3)
 S_1(w_3-w_4)V_a(x+M^{-1}w_4)S_1(w_4)\,\dd w.
\end{multline}
Let $\rho_j(x;a)$ be the sum of these integrals with the four
connection factors replaced by $V_{j_l}(x;w_l)$ and
$j_1+\cdots+j_4=j$. Taylor expansion through total order five gives
\begin{equation}\label{eq:probe-local-expansion}
 r_4(x;M,a)=\sum_{j=0}^5M^{-2-j}\rho_j(x;a)+O_a(M^{-8}),
 \qquad M\ge1,
\end{equation}
uniformly in $x$. The remainder contains four factors bounded by
$C\|a\|_{C^6}(1+|w_l|)^6$ and at least the factor $M^{-6}$.
With $w_0=0$,
$\prod_l(1+|w_l|)^6\le\prod_l(1+|w_l-w_{l-1}|)^{24}$.
Use the weighted majorant at unit mass for the five factors $S_1$.
Its fifth convolution at zero bounds the remainder integral by a
constant independent of $M$.
The coefficients vanish off $\supp a$. There is also a spatial tail
bound
\begin{equation}\label{eq:probe-spatial-tail}
 \int_{\operatorname{dist}(x,\supp a)\ge1}|r_4(x;M,a)|\,\dd x
 \le C_ae^{-M}.
\end{equation}
Indeed, the first and last propagators each travel at least
$d=\operatorname{dist}(x,\supp a)$; their bound is $Me^{-Md}$ for
$d\ge1$. Integrating the three intervening propagators gives
\[
 |r_4(x;M,a)|\le108\|a\|_\infty^4|\supp a|M^{-1}e^{-2Md}.
\]
For $M,d\ge1$, $2Md\ge M+d$, which proves
\eqref{eq:probe-spatial-tail}.

The reduction to field strength must be justified at the level of this
diagonal kernel. Let $\lambda\in C_c^\infty(\R^3)$ be real. On Schwartz
spinors, as a formal power series in $u$,
\[
 (M+\gamma\cdot\partial+ui\gamma\cdot(a+d\lambda))^{-1}
 =e^{-iu\lambda}(M+\gamma\cdot\partial+ui\gamma\cdot a)^{-1}
                       e^{iu\lambda}.
\]
All coefficients preserve Schwartz space. Comparing the coefficient
of $u^4$ gives, almost everywhere off the diagonal,
\begin{equation}\label{eq:probe-gauge-kernel}
 R_{a+d\lambda}^{(4)}(x,y)=\sum_{j=0}^4R_a^{(j)}(x,y)
       \frac{[-i(\lambda(x)-\lambda(y))]^{4-j}}{(4-j)!}.
\end{equation}
For $j=0,1,2,3$, the bounds in
\eqref{eq:probe-convolutions}--\eqref{eq:probe-kernel-majorant} show
that these terms are respectively $O(|x-y|^2)$,
$O(|x-y|^2)$, $O(|x-y|^{3/2})$, and $O(|x-y|)$ locally in $x$.
Continuity of the fourth kernel therefore implies
\[
 r_4(x;M,a+d\lambda)=r_4(x;M,a).
\]
Uniqueness of the finite asymptotic expansion in
\eqref{eq:probe-local-expansion} gives the same invariance for
$\rho_j$, $0\le j\le5$.

Fix $x$. If $p^{(k)}(z)$ is the homogeneous part of degree $k$ of
the fifth-order Taylor polynomial of $a(x+z)$, choose a compact
$\lambda$ equal near $x$ to
\[
 \lambda(x+z)=-\sum_{k=0}^5\frac{z\cdot p^{(k)}(z)}{k+1}.
\]
Euler's identity gives $z\cdot(a+d\lambda)^{[5]}(x+z)=0$ for
the Taylor polynomial. Thus the transformed constant term vanishes,
and its first derivative is the antisymmetric matrix
$F_{\nu\mu}(x)/2$, where
$F_{\nu\mu}=\partial_\nu a_\mu-\partial_\mu a_\nu$.
Each term defining $\rho_j$ has four connection factors. Hence it
vanishes for $j<4$, and at order four only $(j_1,j_2,j_3,j_4)=(1,1,1,1)$
remains. Put
\[
 A_{F,\mu}(z)=\tfrac12\sum_\nu z_\nu F_{\nu\mu},\qquad
 \Psi(F)=r_4(0;1,A_F).
\]
Then $\rho_4(x;a)=\Psi(F(x))$. The odd coefficients vanish as well:
substitution $w_l\mapsto-w_l$ and conjugation by $\gamma_5$ leave
the trace unchanged and multiply it by $(-1)^j$. In particular
$\rho_5=0$. Consequently
\begin{equation}\label{eq:probe-field-strength-reduction}
 r_4(x;M,a)=M^{-6}\Psi(F(x))+O_a(M^{-8}).
\end{equation}

\subsection{The constant-field coefficient}

Rotation of the variables and conjugation by the spin lift in $SU(2)$
show that $\Psi(F)$ depends on the vector
$B_l=\tfrac12\sum_{\nu,\mu}\varepsilon_{l\nu\mu}F_{\nu\mu}$ only
through $|B|$. Take $B=(0,0,b)$, $b\ge0$, and
$A=(-bz_2/2,bz_1/2,0)$. Set $D_u=\partial+iuA$ and
\[
 h_u=1-\Delta+uW_1+u^2W_2,\qquad
 W_1=-2iA\cdot\partial+b\sigma_3,\quad
 W_2=|A|^2.
\]
The Clifford relations give
\begin{equation}\label{eq:probe-squared-dirac}
 (1-\gamma\cdot D_u)(1+\gamma\cdot D_u)
   =\diag(h_u,h_u),\qquad
 (1+\gamma\cdot D_u)^{-1}
   =(1-\gamma\cdot D_u)\diag(h_u^{-1},h_u^{-1}).
\end{equation}
These are identities of formal series acting on Schwartz spinors;
no expansion of an operator with an unbounded connection is being
assumed in operator norm.

We compute the fourth coefficient of $h_u^{-1}$ from the kernel
\begin{multline}\label{eq:probe-Mehler}
 k_u(t;x,y)=(4\pi t)^{-3/2}e^{-|x-y|^2/(4t)}
 \frac{\tau}{\sinh\tau}\exp\left(-\frac{(\tau\coth\tau-1)|x_p-y_p|^2}{4t}
                           +\frac{iub}{2}x\wedge y\right)
 e^{-\tau\sigma_3},\\
 \tau=ubt,\quad x_p=(x_1,x_2),\quad x\wedge y=x_1y_2-x_2y_1.
\end{multline}
It is analytic in $u$ for $|u|bt<\pi$. The kernel satisfies
$(\partial_t+h_u-1)k_u=0$. The planar verification uses
$\alpha=(ub/4)\coth(ubt)$ and
$\phi=-\alpha|x_p-y_p|^2+iub(x\wedge y)/2$:
\[
 \frac{\check h e^\phi}{e^\phi}
   =4\alpha-(4\alpha^2-u^2b^2/4)|x_p-y_p|^2,
 \qquad \alpha'=-(4\alpha^2-u^2b^2/4),
\]
where $\check h=-\partial_1^2-\partial_2^2
-iub(x_1\partial_2-x_2\partial_1)+u^2b^2|x_p|^2/4$.
Differentiation of the prefactor supplies $-4\alpha$; differentiation
of $e^{-\tau\sigma_3}$ supplies the spin term. The longitudinal
Gaussian satisfies the one-dimensional heat equation.

Here only the coefficients of \eqref{eq:probe-Mehler} are needed.
Write $k_u=\sum_{n\ge0}u^nk^{(n)}$. Each coefficient is the Gaussian
times a finite polynomial in $t,x-y,y$, whose monomials
$t^j(x-y)^\alpha y^\beta$ have $2j+|\alpha|\ge n$.
The factor $e^{iub(x-y)\wedge y/2}$ has weight at least one per
power of $u$, and the other factors have weight at least two,
if $t$ has weight two and $x-y$ weight one.
Thus
\[
 K^{(n)}(x,y)=\int_0^\infty e^{-t}k^{(n)}(t;x,y)\,\dd t
\]
converges off the diagonal and, for $n\ge1$, on the diagonal.
Its terms are $y^\beta$ times
\[
 m_{j,\alpha}(x-y)=\int_0^\infty e^{-t}t^j(x-y)^\alpha
                 (4\pi t)^{-3/2}e^{-|x-y|^2/(4t)}\,\dd t.
\]
Their Fourier transforms are
$(i\partial_p)^\alpha[j!(1+|p|^2)^{-j-1}]$, so they define
operators preserving Schwartz space.

Coefficientwise integration of the heat equation identifies these
operators with the formal inverse of $h_u$. Indeed, for a Schwartz
spinor $\varphi$, the zeroth heat coefficient converges to $\varphi$
as $t\downarrow0$. The $n$th coefficient for $n\ge1$ tends to zero
uniformly, since every polynomial term is bounded after integration
against $\varphi$ by $C_\varphi t^{j+|\alpha|/2}$.
Integrating from $\epsilon$ to infinity and sending $\epsilon$ to
zero therefore gives
\[
 (1-\Delta)K^{(n)}+W_1K^{(n-1)}+W_2K^{(n-2)}=\delta_{n0}1.
\]
The convergence can be taken in tempered distributions; the
polynomial differential operators are continuous there. Since the
constant operator $1-\Delta$ is invertible on Schwartz space, the
right inverse is the unique formal inverse.

To evaluate the diagonal, the kernel of its fourth coefficient must
agree with the continuous kernel already constructed. The functions
$m_{j,\alpha}$ are continuous when $2j+|\alpha|\ge2$ and continuously
differentiable when $2j+|\alpha|\ge3$. For this, write
$m_{j,\alpha}(z)=z^\alpha\mu_j(|z|)$ with
\[
 \mu_j(s)=\int_0^\infty e^{-t}t^j(4\pi t)^{-3/2}e^{-s^2/(4t)}\,\dd t.
\]
For $j\ge1$, $\mu_j$ is bounded and continuous at zero and
$\mu_j'(s)=-s\mu_{j-1}(s)/2$ for $s>0$; for $j=0$ use
$\mu_0(s)=e^{-s}/(4\pi s)$. These formulas give the asserted
regularity at zero. Consequently $K^{(3)},K^{(4)}$ and
$\partial_xK^{(4)}$ are continuous. From
\eqref{eq:probe-squared-dirac} the fourth Dirac coefficient has kernel
\[
 (1-\gamma\cdot\partial_x)\diag(K^{(4)},K^{(4)})(x,y)
       -i\gamma\cdot A(x)\diag(K^{(3)},K^{(3)})(x,y).
\]
It equals $R_A^{(4)}$ because both continuous kernels give the same
operator on Schwartz space. At $x=y=0$, $A(0)=0$, and the derivative
term has zero trace: the two Pauli blocks have opposite signs.
It follows that
\begin{align}
 \Psi(F)&=2\int_0^\infty e^{-t}\tr k^{(4)}(t;0,0)\,\dd t\\
 &=-\frac4{45}(4\pi)^{-3/2}\Gamma(7/2)b^4
   =-\frac1{48\pi}|B|^4.\label{eq:probe-constant-field}
\end{align}
In the second line we used
$\tr k_u(t;0,0)=2(4\pi t)^{-3/2}\tau\coth\tau$ and
$\tau\coth\tau=1+\tau^2/3-\tau^4/45+O(\tau^6)$.

\subsection{Integration and dilation}

Integrate \eqref{eq:probe-field-strength-reduction} over the unit
neighbourhood of $\supp a$, and use \eqref{eq:probe-spatial-tail}
on its complement. Equations \eqref{eq:probe-mass-derivative} and
\eqref{eq:probe-constant-field} give
\[
 \frac{\dd}{\dd M}q_{4,M}(a)
 =-\frac1{48\pi M^6}\int|\operatorname{curl}a|^4\,\dd x
                                    +O_a(M^{-8}).
\]
Integration from $M$ to infinity proves
\eqref{eq:large-mass-quartic}.

Finally let $\varphi(x)=e^{-1/(1-|x|^2)}$ for $|x|<1$ and zero
otherwise, set $\beta_{12}=\varphi$, $\beta_{13}=\beta_{23}=0$, and
$a=d^*\beta=(\partial_2\varphi,-\partial_1\varphi,0)$.
The third component of $\operatorname{curl}a$ is
$-(\partial_1^2+\partial_2^2)\varphi$ and has value $4e^{-1}$ at zero.
Thus $I=\int|\operatorname{curl}a|^4\,\dd x>0$.
For $\ell\ge1$ put
$\beta^{(\ell)}(x)=\ell^2\beta(x/\ell)$,
$a_\ell(x)=\ell a(x/\ell)$, and $f_\ell=dd^*\beta^{(\ell)}$.
The unitary dilation $U_\ell\psi(x)=\ell^{-3/2}\psi(x/\ell)$ gives
\[
 U_\ell^{-1}S_{\bar m}V_{a_\ell}U_\ell
       =\ell^2S_{\bar m\ell}V_a,
 \qquad q_{4,\bar m}(a_\ell)=\ell^8q_{4,\bar m\ell}(a).
\]
For $\bar m\ell\ge1$, therefore,
\begin{equation}\label{eq:probe-positive-dilation}
 b(f_\ell)=24q_{4,\bar m}(a_\ell)
   =\frac{\ell^3I}{10\pi\bar m^5}
                      +O_a(\ell\bar m^{-7}).
\end{equation}
This is positive for sufficiently large $\ell$ and completes the proof
of Proposition~\ref{prop:positive-probe}.


\section{Free response coefficients}\label{sec:free-response}

The coefficient controlling the fourth cumulant is a derivative
of the free chamber determinant in a smooth external connection.
The lattice derivative has contact terms, because one bond phase
can be differentiated several times. We first show that these
terms vanish in the continuum limit, then compare the trace with
four first-order vertices to the continuum Dirac trace. The
second-order coefficient requires a different cancellation:
the two-vertex trace and its contact term must be combined
before their limit is taken.

\subsection{The probe and the finite determinant}\label{subsec:response-conventions}

Fix $0<\rho<1$, the standard chamber weights,
and $0<\bar m\le1$. Let $\beta$ be a real nonzero compactly
supported smooth two-form, and put $a=d^*\beta$, $f=dd^*\beta$.
Thus $a_\nu=-\sum_{\mu\ne\nu}\partial_\mu\beta_{\mu\nu}$.
Let $D_\beta$ be the least side of a closed coordinate cube
$Q_\beta$ containing its support. In this section $D_\beta$
measures the support itself, without a buffer. For $0<\eta\le1$
use the seminorm
\[
 |\beta|_{C^j}=\max_{\mu<\nu,\,|\alpha|\le j}
                         \|\partial^\alpha\beta_{\mu\nu}\|_\infty.
\]
Define
\[
 J_\eta v(y)=\eta^{-3}\int_{[-\eta/2,\eta/2]^3}v(y+z)\,\dd z,
 \qquad
 \delta_{\eta,\mu}v(y)=
      \frac{v(y+\eta e_\mu/2)-v(y-\eta e_\mu/2)}{\eta}.
\]
The interpolating functions
\begin{equation}\label{eq:response-lattice-probe}
 a_{\eta,\nu}^0(y)=-\sum_{\mu\ne\nu}
               \delta_{\eta,\mu}J_\eta\beta_{\mu\nu}(y)
\end{equation}
are defined without periodization. If $n=R/\eta\in\N$, $n\ge2$,
the lattice one-cochain on the physical torus is
\begin{equation}\label{eq:response-periodic-probe}
 a_{\eta,\nu}(\eta c)=\sum_{z\in\Z^3}
       a_{\eta,\nu}^0(\eta c+\eta e_\nu/2+Rz).
\end{equation}
It is the adjoint lattice coboundary of the cube-averaged probe
at plaquette centers. In lattice units the external connection
is $B_t(c,c+e_\nu)=t\eta a_{\eta,\nu}(\eta c)$, extended with
the orientation sign to reversed bonds. The free mass in
these units is $m=\bar m\eta$.
The cochains are periodic on the oriented nearest-neighbor
edges of the covering lattice. When $n=2$, distinct edge
orbits under period translations remain distinct even when
their quotient endpoints coincide.

On antiperiodic chamber spinors let
\begin{equation}\label{eq:response-free-ratio}
 N_0(t)=\frac{\det(D^{\mathrm{ch}}[B_t]+m)}
                  {\det(D^{\mathrm{ch}}[0]+m)},\qquad
 c_j=N_0^{(j)}(0),\qquad
 a_{\mathrm{ch}}=(\log N_0)''(0),\quad
 b_{\mathrm{ch}}=(\log N_0)^{(4)}(0).
\end{equation}
The logarithm is the holomorphic branch equal to zero at the
origin. Its denominator is nonzero: the real part of
$D^{\mathrm{ch}}[0]+m$ is at least $m>0$. For comparison
with the charge convention, $q=e\eta^{1/2}$ and the normalized
probe is $\widehat a=\eta^{1/2}a_\eta$, so
$B_t=q(t/e)\widehat a$. The free ratio and its coefficients
are independent of $e$.

\subsection{The singular part of the Bloch inverse}\label{subsec:response-bloch}

Use the fiber $V=\ell^2(O_{\mathrm{cub}};\C^4)$ with measure
$48^{-1}\sum_g$ and the constant-label isometry $J$. Its
dimension is $192$, and traces use this full dimension;
the label measure does not multiply the matrix trace by
$1/48$. Write $\mathcal D_m(p)$ for the free symbol of
\S\ref{subsec:schur-symbol}, and $R_m(p)=\mathcal D_m(p)^{-1}$.
For real $p$ it is invertible when $m>0$.

We need a bound on the difference between the full inverse
and its single light Dirac part, including mixed and heavy
blocks. There is $C_\rho<\infty$ such that
\begin{equation}\label{eq:response-bloch-remainder}
 R_m(p)=J(m+i\gamma\cdot p)^{-1}J^*+\mathcal E_m(p),\qquad
 \|\mathcal E_m(p)\|\le C_\rho,\qquad p\in[-\pi,\pi]^3,\quad0<m\le1.
\end{equation}
To prove this, use the block inversion formula
\eqref{eq:inverse-symbol-blocks}. The off-diagonal blocks
vanish at zero momentum, and are $O(|p|)$ because only
the external hop varies with $p$. The heavy inverse is
bounded by \eqref{eq:symbol-heavy}. The light Schur complement
therefore has the form
\[
 B_m(p)=m+i\gamma\cdot p+r_m(p),\qquad
                    \|r_m(p)\|\le C_\rho|p|^2.
\]
For $|p|$ below a fixed $\rho$-dependent radius, the product
$r_m(p)(m+i\gamma\cdot p)^{-1}$ has norm at most one half.
Its Neumann inverse gives
\[
 \|B_m(p)^{-1}-(m+i\gamma\cdot p)^{-1}\|
        \le2C_\rho |p|^2(m^2+|p|^2)^{-1}\le2C_\rho,
 \qquad |p|\|B_m(p)^{-1}\|\le2.
\]
In the other four terms of the block inverse, each mixed
block supplies a factor $|p|$ beside this light inverse.
The heavy terms are bounded. This proves
\eqref{eq:response-bloch-remainder} near zero. Away from
zero, the real-part and high-momentum ellipticity bounds
of \S\ref{subsec:schur-symbol} bound the full inverse uniformly,
and the continuum inverse is bounded by $1/|p|$. This
proves the estimate on the whole cube. In particular
\begin{equation}\label{eq:response-bloch-bounds}
 \|R_m(p)\|\le C_\rho+(m^2+|p|^2)^{-1/2}
                    \le\frac{C'_\rho}{m+|p|}.
\end{equation}

Define external-hop matrices
\[
 (T_\nu^\pm u)(g)=1_{\{ge_1=\pm e_\nu\}}
                          3(\rho\mp\gamma_\nu)u(gr_{e_1}).
\]
Then
\begin{align}\label{eq:response-symbol-vertices}
 \mathcal D_m(p)&=\mathcal D_{m,\mathrm{loc}}
       -\sum_\nu(e^{ip_\nu}T_\nu^++e^{-ip_\nu}T_\nu^-),\notag\\
 V_\nu(p):=\partial_\nu\mathcal D_m(p)
       &=-i(e^{ip_\nu}T_\nu^+-e^{-ip_\nu}T_\nu^-),\notag\\
 W_\nu(p):=\partial_\nu^2\mathcal D_m(p)
       &=e^{ip_\nu}T_\nu^++e^{-ip_\nu}T_\nu^-.
\end{align}
Mixed second derivatives vanish. The disjoint rows indexed
by $ge_1=\pm e_\nu$ imply
$\|\sum z_\nu^\pm T_\nu^\pm\|\le3(1+\rho)\max|z_\nu^\pm|$.
Since eight labels contribute to each direction,
\begin{equation}\label{eq:response-light-vertex}
 J^*T_\nu^\pm J=\tfrac12(\rho\mp\gamma_\nu),\qquad
 J^*V_\nu(0)J=i\gamma_\nu.
\end{equation}
All nonzero derivatives of the symbol are bounded by six
on real momenta. Differentiating $R_m\mathcal D_m=1$ gives
\[
 \partial^\alpha R_m=-\sum_{0<\delta\le\alpha}
   \binom\alpha\delta(\partial^{\alpha-\delta}R_m)
                    (\partial^\delta\mathcal D_m)R_m.
\]
Starting with \eqref{eq:response-bloch-bounds}, induction bounds
each summand by a constant times
$(\operatorname{dist}(p,2\pi\Z^3)+m)^{-1-|\alpha|}$;
the remaining factor of the distance has nonnegative exponent
and is bounded by a power of seven. Thus
\begin{equation}\label{eq:response-inverse-derivatives}
 \|\partial^\alpha R_m(p)\|
   \le C_{\alpha,\rho}
       (\operatorname{dist}(p,2\pi\Z^3)+m)^{-1-|\alpha|}.
\end{equation}
At $m=\bar m\eta$, \eqref{eq:response-bloch-remainder} and
the first two differentiated inverse identities give
\begin{equation}\label{eq:response-scaled-inverse}
 \eta^{1+|\alpha|}(\partial^\alpha R_{\bar m\eta})(\eta P_\eta)
       \longrightarrow J\partial^\alpha
                 (\bar m+i\gamma\cdot P)^{-1}J^*,
 \quad |\alpha|\le2,\quad P_\eta\to P.
\end{equation}
In the second derivative the contact term has two inverse
factors and one bounded $W_\nu$, so its rescaled norm has
an extra factor $\eta$ and vanishes.

The involution $(\mathcal Uu)(g)=\gamma_5u(-g)$ satisfies
\begin{equation}\label{eq:response-bloch-parity}
 \mathcal U\mathcal D_m(p)\mathcal U=\mathcal D_m(-p),\qquad
 \mathcal UV_\nu(p)\mathcal U=-V_\nu(-p).
\end{equation}
Indeed $-g$ reverses every bond direction, and
$\gamma_5\gamma_\nu\gamma_5=-\gamma_\nu$ restores the
hopping matrix. The external phase then changes $p$ to
$-p$. This parity will cancel the first momentum derivative
of the polarization kernel.

\subsection{Derivative and contact terms}\label{subsec:response-contact}

Let $\mathcal P_n=\frac{2\pi}{n}(\Z+\frac12)^3\cap(-\pi,\pi]^3$.
The normalized Bloch vectors are $n^{-3/2}e^{ip\cdot c}$,
and the lattice free propagator $G$ is the multiplier
$R_m(p)$ on these antiperiodic momenta. Write
$A_\nu(c)=a_{\eta,\nu}(\eta c)$ and
$A_\nu^-(c)=A_\nu(c-e_\nu)$, and let
$U_\nu^\pm\psi(c)=\psi(c\pm e_\nu)$. The bond values are
$b_\nu^+=\eta A_\nu$, $b_\nu^-=-\eta A_\nu^-$.
The exact connection derivatives are
\begin{equation}\label{eq:response-connection-derivatives}
 D_j=-\sum_{\nu,\pm}(ib_\nu^\pm)^jT_\nu^\pm U_\nu^\pm,
 \quad D^{\mathrm{ch}}[B_t]-D^{\mathrm{ch}}[0]
                  =\sum_{j\ge1}\frac{t^j}{j!}D_j,
 \quad Y_j=GD_j.
\end{equation}
The shifts act first. These $D_j$ are derivatives, with
the factorials outside the vertices.

For small complex $t$, $N_0(t)=\det(1+\sum_{j\ge1}t^jY_j/j!)$
and its zero-at-zero logarithm is the convergent finite-matrix
trace series $\sum_{l\ge1}(-1)^{l+1}\operatorname{Tr}(\sum_j
t^jY_j/j!)^l/l$. The tail $l\ge5$ is $O(t^5)$, so it
contributes no fourth derivative. Cyclicity of the trace
in the terms $l\le4$ gives
\begin{equation}\label{eq:response-fourth-identity}
 \begin{split}
 b_{\mathrm{ch}}={}&\operatorname{Tr}Y_4
       -4\operatorname{Tr}(Y_1Y_3)-3\operatorname{Tr}Y_2^2\\
       &+12\operatorname{Tr}(Y_1^2Y_2)-6\operatorname{Tr}Y_1^4.
 \end{split}
\end{equation}
For example, the cyclically equal orders $(1,1,2)$,
$(1,2,1)$ and $(2,1,1)$ contribute
$24\cdot\frac13\cdot\frac32=12$. The quadratic identity
is $a_{\mathrm{ch}}=\operatorname{Tr}Y_2-\operatorname{Tr}Y_1^2$.
The determinant series is used only near zero; the determinant
itself is entire in $t$.

We bound all four contact traces in
\eqref{eq:response-fourth-identity} uniformly in $R$ and the
position of the probe. The interpolating periodic field is
\[
 \widetilde A_\nu(y)=\sum_z
             a_{\eta,\nu}^0(y+\eta e_\nu/2+Rz),
 \qquad A_\nu(c)=\widetilde A_\nu(\eta c).
\]
When $R\ge D_\beta+1$, averaging translates of the derivatives
of $\beta$ shows
\begin{equation}\label{eq:response-probe-estimates}
 \begin{gathered}
 \|\widetilde A\|_\infty+\|\partial\widetilde A\|_\infty\le C_\beta^{\mathrm{free}},
 \qquad \|\widetilde A-a^R\|_\infty\le C_\beta^{\mathrm{free}}\eta,\\
 \eta^3|\{c:A_\nu(c)\ne0\text{ for some }\nu\}|\le(D_\beta+3)^3,
 \qquad
 \|\eta A_\nu\|_{\ell^4}\le C_\beta^{\mathrm{free}}\eta^{1/4},\\
 \|\eta(A_\nu^--A_\nu)\|_{\ell^4}\le C_\beta^{\mathrm{free}}\eta^{5/4},
 \qquad\|\widetilde A-a^R\|_{L^4(T_R)}\le C_\beta^{\mathrm{free}}\eta.
 \end{gathered}
\end{equation}
Every averaging displacement has size at most $\eta$ in
sup norm. The translates of $\supp\beta$ are disjoint,
so their derivative sup norms are unchanged. The supports
of the interpolants lie in $Q_\beta+[-\eta,\eta]^3$ modulo
the periods. Its volume and number of lattice points give
the last estimates, including the union with its one-step
translate for the difference.

Let $\widehat{A_\nu^j}(k)=n^{-3}\sum_cA_\nu(c)^je^{-ik\cdot c}$
on the finite momentum group $\mathcal G_n=\frac{2\pi}{n}\Z^3/
2\pi\Z^3$, and set $F_j(k)=\sum_\nu|\widehat{A_\nu^j}(k)|$.
There are constants $C_{\mathrm{Four},1},C_{\mathrm{Four},0}$ depending only on $\beta$ with
\begin{equation}\label{eq:response-probe-fourier-norms}
 \|F_j\|_{\ell^1(\mathcal G_n)}\le C_{\mathrm{Four},1}^j,
 \qquad R^3\|F_j\|_\infty\le C_{\mathrm{Four},1}^{j-1}C_{\mathrm{Four},0}.
\end{equation}
The transform of the centered difference and the cube average
in \eqref{eq:response-lattice-probe} is
\[
 -\sum_{\mu\ne\nu}iP_\mu\operatorname{sinc}(\eta P_\mu/2)
 \prod_l\operatorname{sinc}(\eta P_l/2)\widehat\beta_{\mu\nu}(P),
\]
where $\operatorname{sinc}(z)=\sin z/z$ and $\operatorname{sinc}(0)=1$.
The midpoint shift multiplies this transform by a phase of modulus one.
The shifted transform has modulus at most
$C_\beta^{\mathrm{free}}(1+|P|)^{-4}$. Poisson summation gives
\[
 \widehat A_\nu(k)=R^{-3}\sum_{z\in\Z^3}
      \widehat{a_{\eta,\nu}^0(\,\cdot+\eta e_\nu/2)}
                                      ((k+2\pi z)/\eta).
\]
Summing over representatives $k$ and $z$ gives a single
sum over $\frac{2\pi}{R}\Z^3$, bounded after multiplication
by $R^{-3}$ by its integrable radial majorant, uniformly
for $R\ge1$. For $|k|_\infty\le\pi$, nonzero $z$ have
$|(k+2\pi z)/\eta|\ge\pi|z|_\infty/\eta$; the convergent
sum $\sum_{z\ne0}|z|_\infty^{-4}$ gives the sup bound.
Powers of $A_\nu$ give convolutions on $\mathcal G_n$;
$\ell^1*\ell^1\to\ell^1$ and $\ell^1*\ell^\infty\to\ell^\infty$
give \eqref{eq:response-probe-fourier-norms} for all $j$.

The momentum blocks of a vertex are
\[
 D_j(p,p')=-\eta^j\sum_\nu\widehat{A_\nu^j}(p-p')
       [i^je^{ip'_\nu}T_\nu^++(-i)^je^{-ip_\nu}T_\nu^-],
 \qquad\|D_j(p,p')\|\le6\eta^jF_j(p-p').
\]
For $l\le3$ and $j_1+\cdots+j_l=4$, apply these bounds
and \eqref{eq:response-bloch-bounds} to the cyclic trace.
With differences $k_i=p_i-p_{i+1}$, the sum is parametrized
by $p_1,k_1,\ldots,k_{l-1}$. H\"older in $p_1$ bounds
the product of its $l$ translates of $(m+|p|)^{-1}$ by
$\sum_p(m+|p|)^{-l}$. Bound the final $F_{j_l}$ in sup
norm and the others in $\ell^1$. Thus
\begin{equation}\label{eq:response-cyclic-contact}
 |\operatorname{Tr}(Y_{j_1}\cdots Y_{j_l})|
   \le192(6C'_\rho)^l C_{\mathrm{Four},1}^3C_{\mathrm{Four},0}\,
          \eta\left(n^{-3}\sum_{p\in\mathcal P_n}(m+|p|)^{-l}\right).
\end{equation}
The factor $\eta$ is $\eta^4n^3/R^3$.

For antiperiodic momenta, comparison with their disjoint
cubes of side $2\pi/n$ gives, for a decreasing radial
function $h$, $n^{-3}\sum_ph(|p|)\le\pi^{-3}\int h(|y|)\,\dd y$.
Indeed in the cube of $p$ one has $|y|\le2|p|$, and a
rescaling by two gives the factor $\pi^{-3}$. On the
bounded momentum cube this yields constants for $l=1,2$
and $\log(1+5/m)$ for $l=3$. Inserting the multiplicities
$1,4,3,12$ from \eqref{eq:response-fourth-identity} proves
\begin{equation}\label{eq:response-contact-bound}
 \begin{split}
 &|\operatorname{Tr}Y_4|+4|\operatorname{Tr}(Y_1Y_3)|
          +3|\operatorname{Tr}Y_2^2|+12|\operatorname{Tr}(Y_1^2Y_2)|\\
 &\hspace{25mm}\le C_{\mathrm{ct}}(\rho,\beta)\,
               \eta\left(1+\log(1+5/(\bar m\eta))\right).
 \end{split}
\end{equation}
The constant does not depend on $\bar m$, $\eta$, $R$ or
the position of $\supp\beta$.

\subsection{The fourth-order continuum coefficient}\label{subsec:response-fourth}

Write $s_m(p)=(m+i\gamma\cdot p)^{-1}$. On the continuum
antiperiodic torus its momentum set is
$\mathcal P_R=\frac{2\pi}{R}(\Z+\frac12)^3$, and
$S^R$ is the multiplier $s_{\bar m}(P)$ there. For
$R\ge D_\beta+1$ let $a^R=\sum_za(\,\cdot+Rz)$, and put
\begin{equation}\label{eq:response-continuum-fourth}
 b_R(f)=-6\operatorname{Tr}(S^Ri\gamma\cdot a^R)^4,
 \qquad b(f)=-6\operatorname{Tr}(S_{\bar m}i\gamma\cdot a)^4.
\end{equation}
Both traces will be justified below. The compactly supported
potential $a=d^*\beta$ is determined by $f=da$: two such
potentials differ by a compactly supported smooth one-form
which is closed and coclosed, hence harmonic and zero.
Thus the notation $b(f)$ and $b_R(f)$ depends only on $f$.

\begin{proposition}\label{prop:response-fourth}
The operators $S^Ri\gamma\cdot a^R$ and
$S_{\bar m}i\gamma\cdot a$ belong to the Schatten class
of order four; their fourth-power traces are real.
There are constants
$C_1(\rho,\bar m,\beta)$ and $C_2(\bar m,\beta)$ such that
\begin{align}\label{eq:response-fourth-errors}
 |b_{\mathrm{ch}}(\eta,R)-b_R(f)|&\le C_1\eta^{1/4},\notag\\
 |b_R(f)-b(f)|&\le C_2e^{-\bar m(R-D_\beta)},
 \qquad R\ge D_\beta+1,\quad0<\eta\le1,\quad R/\eta\in\N.
\end{align}
The constants are unchanged by every real translate of
the probe. In particular $b_{\mathrm{ch}}(\eta,R)$ tends
to $b(f)$ whenever $\eta\to0$ and $R\to\infty$.
\end{proposition}

\begin{proof}
We first establish the trace bounds used in the comparisons.
For a lattice multiplier $M_h$ with fiber symbol $h(p)$
and a scalar multiplication operator $v$, one has
\begin{equation}\label{eq:response-four-ideal-lattice}
 \|M_hv\|_4\le
       \left(n^{-3}\sum_{p\in\mathcal P_n}\|h(p)\|_4^4\right)^{1/4}
                                         \|v\|_{\ell^4}.
\end{equation}
To prove it put $U=M_{h^*h}$ and $V=|v|^2$. Expanding the
nonnegative trace of $(UV-VU)^*(UV-VU)$ gives
$\operatorname{Tr}(UVUV)\le\operatorname{Tr}(U^2V^2)$.
The left side is $\|M_hv\|_4^4$, and each normalized
Bloch vector has constant squared modulus $n^{-3}$.
Evaluating the right side in that basis gives
\eqref{eq:response-four-ideal-lattice}. On the continuum
torus the same proof, first for a finite set of momenta,
gives
\begin{equation}\label{eq:response-four-ideal-torus}
 \|M_hv\|_4\le
       \left(R^{-3}\sum_{P\in\mathcal P_R}\|h(P)\|_4^4\right)^{1/4}
                                           \|v\|_{L^4(T_R)}.
\end{equation}
Truncation in momentum converges in the Schatten norm
when the sum is finite, and its operator-norm limit is
$M_hv$, proving the general form.

The cube comparison used above also gives
\begin{equation}\label{eq:response-four-momentum-sums}
 \begin{gathered}
 n^{-3}\sum_{p\in\mathcal P_n}(m^2+|p|^2)^{-2}\le\frac1{\pi m},
 \qquad
 R^{-3}\sum_{P\in\mathcal P_R}(\bar m^2+|P|^2)^{-2}
                                             \le\frac1{\pi\bar m},\\
 R^{-3}\sum_{|P|_\infty\le\pi/(2\eta)}|P|^{-2}\le C\eta^{-1},
 \qquad
 R^{-3}\sum_{|P|_\infty>\pi/(2\eta)}|P|^{-4}\le C\eta.
 \end{gathered}
\end{equation}
For the last bound use the decreasing radial majorant
$\min((2\eta/\pi)^4,|P|^{-4})$; its integral is $O(\eta)$.
The preceding bounds and
$\|s_{\bar m}(P)\|_4^4=4(\bar m^2+|P|^2)^{-2}$ show that
$S^Ri\gamma\cdot a^R$ has a bounded Schatten-four norm,
uniformly in $R\ge D_\beta+1$. The same estimate applies
to the interpolating field $\widetilde A$.

On $\R^3$, let
\[
 Y_{\bar m}(x)=\frac{e^{-\bar m|x|}}{4\pi|x|},\qquad
 S_{\bar m}(x)=(\bar m-\gamma\cdot\partial)Y_{\bar m}(x).
\]
The scalar Fourier transform is $(\bar m^2+|P|^2)^{-1}$;
the derivative has no point mass, since its boundary
term on a sphere of radius $r$ is $O(r)$. Its Fourier
transform therefore gives $s_{\bar m}(P)$. These kernels
are integrable, so their convolutions are the corresponding
bounded Fourier multipliers on $L^2$. Also
$\|Y_{\bar m}\|_{L^2}^2=(8\pi\bar m)^{-1}$.
For bounded compactly supported $v$, the kernel of
$v^*S_{\bar m}^*S_{\bar m}v$ is
$\overline{v(x)}Y_{\bar m}(x-y)v(y)1_4$. Tonelli and
Cauchy--Schwarz in the translated factors give
\[
 \|S_{\bar m}v\|_4^4
   \le4\|Y_{\bar m}\|_{L^2}^2\|v\|_{L^4}^4
   =\frac{\|v\|_{L^4}^4}{2\pi\bar m}.
\]
For the matrix multiplication $i\gamma\cdot a$, the
Clifford identity gives the same proof with $|a|$ in
place of $|v|$. Thus
\begin{equation}\label{eq:response-four-ideal-infinite}
 \|S_{\bar m}i\gamma\cdot a\|_4^4\le
                    (2\pi\bar m)^{-1}\int_{\R^3}|a|^4\,\dd x.
\end{equation}
The traces in
\eqref{eq:response-continuum-fourth} exist by Schatten
H\"older.

To check reality, treat $S=S_{\bar m}$ or $S^R$ and
$v=i\gamma\cdot a$ or $i\gamma\cdot a^R$ together.
Factor $v=|a|^{1/2}(i\gamma\cdot a/|a|)|a|^{1/2}$,
with the middle factor zero at its zeros. The preceding
estimates put $SvS$ in the Hilbert--Schmidt class.
One has $\gamma_5S\gamma_5=S^*$ and
$\gamma_5v\gamma_5=v^*$. Hence conjugating
$\operatorname{Tr}[(SvS)v(SvS)v]$ by $\gamma_5$ gives
its complex conjugate after cyclically moving the
bounded $v$ past a trace-class product of the two
Hilbert--Schmidt factors. This proves reality.

We now compare the lattice fourth-power trace with its
continuum torus value. The four-power inequality used
throughout is
\begin{equation}\label{eq:response-fourth-trace-difference}
 |\operatorname{Tr}X^4-\operatorname{Tr}Z^4|
 \le4\max(\|X\|_4,\|Z\|_4)^3\|X-Z\|_4.
\end{equation}
It follows by expanding $X^4-Z^4$ into four products
with one $X-Z$ factor and applying Schatten H\"older.
All the operators compared below have uniformly bounded
Schatten-four norms; each bound follows from
\eqref{eq:response-four-ideal-lattice} or
\eqref{eq:response-four-ideal-torus}. In particular
\eqref{eq:response-bloch-bounds} gives
$n^{-3}\sum_p\|R_m(p)\|_4^4\le
C_\rho(1+m^{-1})$, which, with
$\|\eta A_\nu\|_{\ell^4}\le C_\beta^{\mathrm{free}}\eta^{1/4}$,
proves $\|Y_1\|_4\le C_{\rho,\bar m,\beta}$.

\emph{Removing the bounded part of the inverse.}
Let $\chi(p)=1_{\{|p|_\infty\le\pi/2\}}$ and
$S_c$ be the lattice multiplier $\chi(p)s_m(p)$ on
four-component antiperiodic spinors. The complementary
inverse
$R_m(p)-J\chi(p)s_m(p)J^*$ is uniformly bounded:
\eqref{eq:response-bloch-remainder} controls the whole
cube, and $s_m$ is bounded on $1-\chi$. Thus
\[
 \|Y_1-JS_cJ^*D_1\|_4\le C_{\rho,\beta}\eta^{1/4}.
\]
The two norms in \eqref{eq:response-fourth-trace-difference}
are bounded by the same estimates. Since $J^*J=1$,
cyclicity reduces the fourth-power trace of the light
operator to that of $S_cw_1$, where
\begin{equation}\label{eq:response-compressed-vertex}
 w_1=J^*D_1J=-\frac{i\eta}{2}\sum_\nu
       [A_\nu(\rho-\gamma_\nu)U_\nu^+
             -A_\nu^-(\rho+\gamma_\nu)U_\nu^-].
\end{equation}
This comparison costs $C\eta^{1/4}$ in the trace. The
compression uses \eqref{eq:response-light-vertex}, and
therefore has one four-component light species.

\emph{Replacing the vertex by multiplication.}
Move shifts to the left by
$A_\nu U_\nu^+=U_\nu^+A_\nu^-$ and
$A_\nu^-U_\nu^-=U_\nu^-A_\nu$. Then
\begin{equation}\label{eq:response-vertex-error}
 \begin{split}
 w_1-i\eta\gamma\cdot A=-\frac{i\eta}{2}\sum_\nu\big[&
  (\rho-\gamma_\nu)(U_\nu^+-1)A_\nu^-
       -(\rho+\gamma_\nu)(U_\nu^--1)A_\nu\\
       &+(\rho-\gamma_\nu)(A_\nu^--A_\nu)\big].
 \end{split}
\end{equation}
The multiplier
$\chi(p)s_m(p)(\rho\mp\gamma_\nu)(e^{\pm ip_\nu}-1)$
has norm at most two, since $|e^{ip_\nu}-1|\le|p_\nu|$.
Its Schatten-four multiplier average is consequently
bounded independently of the mass. The first two terms
of \eqref{eq:response-vertex-error} then cost
$O(\eta^{1/4})$. The last term costs
$O(m^{-1/4}\eta^{5/4})=O(\eta)$ by
\eqref{eq:response-probe-estimates}. Therefore
$\|S_c(w_1-i\eta\gamma\cdot A)\|_4\le C\eta^{1/4}$,
and \eqref{eq:response-fourth-trace-difference} gives
the same trace error.

\emph{Identifying the truncated continuum momenta.}
The map $p\mapsto P=p/\eta$ identifies the retained
lattice momenta with
$\mathcal P_R^<=\{P\in\mathcal P_R:|P|_\infty\le\pi/(2\eta)\}$,
and $\eta s_m(p)=s_{\bar m}(P)$. The operator
$S_ci\eta\gamma\cdot A$ has range in this truncated
space. Its fourth trace is thus that of its compression
to the space. In the normalized momentum basis this
compression has blocks
\[
 T_{\mathrm{lat}}(P,P')=s_{\bar m}(P)i\sum_\nu\gamma_\nu
                 n^{-3}\sum_cA_\nu(c)e^{-i\eta(P-P')\cdot c}.
\]
Let $\widetilde A$ be the smooth interpolant and
$\widehat A_R(k)=R^{-3}\int_{T_R}\widetilde A(y)e^{-ik\cdot y}\,\dd y$.
Absolute Fourier convergence, using four derivatives
of the interpolant, gives the exact alias identity
\begin{equation}\label{eq:response-aliasing}
 n^{-3}\sum_cA_\nu(c)e^{-i\eta k\cdot c}
    =\widehat A_{R,\nu}(k)+
          \sum_{z\ne0}\widehat A_{R,\nu}(k+2\pi z/\eta).
\end{equation}
For $k=P-P'$, $|k|_\infty\le\pi/\eta$. Four integrations
by parts give
$|\widehat A_{R,\nu}(k')|\le C_\beta^{\mathrm{free}} R^{-3}|k'|_\infty^{-4}$;
the nonzero aliases in \eqref{eq:response-aliasing} are
therefore bounded by $C_\beta^{\mathrm{free}} R^{-3}\eta^4$.

Replacing the discrete matrix elements by $\widehat A_R$
gives the compression $T_{\mathrm{int}}$ of
$\chi S^Ri\gamma\cdot\widetilde A$. Its norm is bounded.
For the difference, sum the squared Hilbert--Schmidt
norms of its blocks. There are at most $R^3\eta^{-3}$
choices of $P'$, while
$\sum_{P\in\mathcal P_R^<}|P|^{-2}\le CR^3\eta^{-1}$.
Since $\|\gamma\cdot z\|_2^2=4|z|^2$, the alias bound
gives
\[
 \|T_{\mathrm{lat}}-T_{\mathrm{int}}\|_2^2
     \le C_\beta^{\mathrm{free}}(R^3\eta^{-1})(R^3\eta^{-3})
                                     (R^{-3}\eta^4)^2
     =C_\beta^{\mathrm{free}}\eta^4.
\]
All powers of $R$ cancel. The Schatten-four norm is
at most the Hilbert--Schmidt norm, so the trace error
in this comparison is $O(\eta^2)$.

\emph{Replacing the interpolant and removing the cutoff.}
The difference between the compressions of
$\chi S^Ri\gamma\cdot\widetilde A$ and
$\chi S^Ri\gamma\cdot a^R$ has Schatten-four norm
$O(\eta)$ by \eqref{eq:response-four-ideal-torus} and
\eqref{eq:response-probe-estimates}. For an operator
$X$ with $X=\chi X$, cyclicity gives
$\operatorname{Tr}X^4=\operatorname{Tr}(\chi X\chi)^4$;
the same holds for $X=X\chi$. Thus we can remove the
compression and compare $\chi S^Ri\gamma\cdot a^R$
directly with $S^Ri\gamma\cdot a^R$. The removed momentum
tail in \eqref{eq:response-four-momentum-sums} gives
\[
 \|(1-\chi)S^Ri\gamma\cdot a^R\|_4\le C_\beta^{\mathrm{free}}\eta^{1/4}.
\]
These are respectively the fourth and fifth comparisons;
their trace errors are $O(\eta)$ and $O(\eta^{1/4})$.
The preceding three comparisons, together with
\eqref{eq:response-contact-bound}, prove the first line
of \eqref{eq:response-fourth-errors}. Here
$\eta(1+\log(1+5/(\bar m\eta)))\le C_{\bar m}\eta^{1/4}$.

It remains to compare torus and infinite-space traces.
Because the multiplication by $a$ is supported in
$Q_\beta$, both traces can be computed on $L^2(Q_\beta;\C^4)$,
using extension by zero to $\R^3$ or antiperiodic extension
to $T_R$. The same compression identity proves this
localization. Convolution by the integrable kernel
$S_{\bar m}(x)$ on antiperiodic functions agrees with
$S^R$ on its momentum basis and hence on $L^2$.
After restricting its two variables to $Q_\beta$, its
difference from the infinite-space kernel is
\[
 W_R(x,y)=\sum_{z\in\Z^3\setminus\{0\}}
                  (-1)^{z_1+z_2+z_3}S_{\bar m}(x-y-Rz).
\]
For $|u|\ge1$ the Yukawa formula gives
$\|S_{\bar m}(u)\|\le(3/(4\pi))e^{-\bar m|u|}$.
Moreover $|x-y-Rz|\ge R|z|_\infty-D_\beta\ge
(R-D_\beta)+|z|_\infty-1\ge1$. Counting
$24j^2+2$ images at $|z|_\infty=j$ yields
\[
 \|W_R(x,y)\|\le C_{\bar m}e^{-\bar m(R-D_\beta)},
 \qquad x,y\in Q_\beta.
\]
The series converges absolutely and uniformly there.
Multiplication by $i\gamma\cdot a(y)$ then gives a
Hilbert--Schmidt difference with this same exponential
bound, since $Q_\beta$ has finite volume. The infinite
localized operator has bounded Schatten-four norm,
and the torus one is bounded by it and the difference.
Applying \eqref{eq:response-fourth-trace-difference}
proves the second line of \eqref{eq:response-fourth-errors}.

The constants above involve only $\rho$, $\bar m$,
the side of the support cube, derivative sup and
$L^1$ norms of $\beta$, and moduli of its Fourier
transform. All are unchanged by a real translation.
This gives the asserted uniformity and the limit.
\end{proof}

\subsection{Evenness and placement of the probe}\label{subsec:response-symmetry}

The coefficients in \eqref{eq:response-free-ratio} are real, and
\begin{equation}\label{eq:response-coefficient-relations}
 c_{2j+1}=0,\qquad c_2=a_{\mathrm{ch}},\qquad
                 c_4=b_{\mathrm{ch}}+3c_2^2.
\end{equation}
We verify the determinant identities for complex external connections:
the later comparison of coefficients uses them on a complex circle.

For a real connection, each chamber bond involution with its phase,
$R_i(B)$, is self-adjoint and unitary. The direction of a bond changes
sign at its other endpoint. Consequently
\[
 D^{\mathrm{ch}}[B]=\frac{\rho}{2}
       \sum_iw_i(1-R_i(B))+\frac12
       \sum_iw_i(\gamma\cdot\widehat d_i)R_i(B),
\]
where the first sum is nonnegative and the second is anti-self-adjoint.
The two spin sectors $\gamma_\mu=\diag(\sigma_\mu,-\sigma_\mu)$
give adjoint restrictions, as in \eqref{eq:reg-two-blocks}. Thus
$\det(D^{\mathrm{ch}}[B]+m)$ is a squared absolute value.
The real part of $D^{\mathrm{ch}}[B]+m$ is at least $m>0$,
so the determinant is strictly
positive. In particular $N_0(t)>0$ for real $t$.

The transpose identity extends charge symmetry to complex connections.
Let $C$ be the charge-conjugation matrix of
\S\ref{subsec:reg-symmetries}. Transposing a phased bond reverses its
direction and changes its phase from $B$ to $-B$.
The identity $C\gamma_v^\top C^{-1}=-\gamma_v$ changes the transposed
matrix at the other endpoint back to the matrix at the original one.
The diagonal part is scalar; hence
\[
 C(D^{\mathrm{ch}}[B]+m)^\top C^{-1}
                   =D^{\mathrm{ch}}[-B]+m.
\]
The transpose here is for the bilinear chamber pairing. In the site
basis this pairing is $1/48$ times the identity, so this is the
ordinary matrix transpose and preserves the determinant. Every matrix
entry is a finite sum of constants and exponentials of linear
functions of $t$. Therefore $N_0$ is entire, and the transpose
identity gives $N_0(-t)=N_0(t)$ for all $t\in\C$.
Expanding its logarithm at zero proves
\eqref{eq:response-coefficient-relations}.

We also record exact covariance under lattice translations and cubic
reflections. Let $g=(v,r)$, where $v\in\Z^3$ and $r$ is a signed
permutation. A reflection with normal $\widehat n$ has the spin lift
$S_{\widehat n}=(\gamma\cdot\widehat n)\gamma_5$, for which
\[
 S_{\widehat n}(\gamma\cdot z)S_{\widehat n}^{-1}
       =\gamma\cdot(z-2(z\cdot\widehat n)\widehat n),
                         \qquad z\in\C^3.
\]
Choose a word representing $r$ in the three linear chamber
reflections and let $S_g$ be the corresponding product. It follows
that $S_g(\gamma\cdot z)S_g^{-1}=\gamma\cdot rz$.
The site map $(c,h)\mapsto(rc+v,rh)$ commutes with each bond
involution and carries $\widehat d_i(c,h)$ to
$r\widehat d_i(c,h)$. Therefore, for
\[
 (U_g\psi)(x)=S_g\psi(g^{-1}x),\qquad
 (g_*B)(c,c')=B(g^{-1}c,g^{-1}c'),
\]
the hopping formula gives
\begin{equation}\label{eq:response-complex-covariance}
 D^{\mathrm{ch}}[g_*B]+m
       =U_g(D^{\mathrm{ch}}[B]+m)U_g^{-1},
                                             \qquad B\ \text{complex}.
\end{equation}
This identity descends to the antiperiodic torus. Indeed
$r^{-1}e_\mu=\pm e_\nu$, and either translation by $ne_\nu$
or by $-ne_\nu$ changes an antiperiodic spinor's sign.
Thus $U_g$ and its inverse preserve the spinor space.
Taking determinants in \eqref{eq:response-complex-covariance}
shows that $N_0$, $c_j$, $a_{\mathrm{ch}}$ and $b_{\mathrm{ch}}$
are unchanged when the external cochain is replaced by $g_*B_t$.

For a physical translation $v\in\eta\Z^3$, cube averaging and
centered differences commute with translation. Periodization in
\eqref{eq:response-periodic-probe} then gives
$a_\eta[\beta(\,\cdot-v)](y)=a_\eta[\beta](y-v)$.
This is exactly the cochain action with lattice translation $v/\eta$.
The same determinant coefficients can therefore be used at each
lattice placement of the probe and at each of its reflected images.

\subsection{The two-photon coefficient}\label{subsec:response-polarization}

The coefficient is
\[
 a_{\mathrm{ch}}=\operatorname{Tr}(GD_2)-
                     \operatorname{Tr}(GD_1GD_1).
\]
Estimating the two traces separately would lose a power of the mesh.
We keep their cancellation by expressing their difference
as a polarization kernel. Its integral vanishes at zero momentum;
parity removes the first momentum derivative.

Define the bosonic momentum representatives
\[
 \mathcal K_{\eta,R}=
 \left\{\frac{2\pi}{R}j:j\in\Z^3,\ -n/2<j_\nu\le n/2\right\},
                   \qquad n=R/\eta,
\]
and the midpoint transform of the unperiodized probe
\begin{equation}\label{eq:response-midpoint-transform}
 h_{\eta,\nu}(K)=\eta^3
       \sum_{y\in\eta\Z^3+\eta e_\nu/2}
                     a_{\eta,\nu}^0(y)e^{-iK\cdot y}.
\end{equation}
The sum is finite. For real $p,k$ put
\begin{align}
 L_{\mu\nu}(p,k)
  &=\delta_{\mu\nu}\tr[R_m(p)W_\nu(p)]\notag\\
  &\quad-\tr[R_m(p)V_\mu(p+k/2)
                     R_m(p+k)V_\nu(p+k/2)],\label{eq:response-polarization-integrand}\\
 \Pi^{\eta,R}_{\mu\nu}(K)
  &=\eta^{-1}n^{-3}\sum_{p\in\mathcal P_n}
                       L_{\mu\nu}(p,\eta K),\notag\\
 \Pi^\eta_{\mu\nu}(K)
  &=\eta^{-1}\int_{[-\pi,\pi]^3}
          L_{\mu\nu}(p,\eta K)\frac{\dd p}{(2\pi)^3}.
                                  \label{eq:response-polarization-kernels}
\end{align}
Only $\Pi^{\eta,R}$ is a finite-volume polarization; the integral
$\Pi^\eta$ retains the lattice symbol but has no volume parameter.
The exact coefficient is
\begin{equation}\label{eq:response-polarization-identity}
 a_{\mathrm{ch}}(\eta,R)
    =R^{-3}\sum_{K\in\mathcal K_{\eta,R}}\sum_{\mu,\nu}
       \overline{h_{\eta,\mu}(K)}h_{\eta,\nu}(K)
                                      \Pi^{\eta,R}_{\mu\nu}(K).
\end{equation}

To check the midpoint and normalization in this formula, write
$b_\nu(c)=\eta A_\nu(c)$ and
$\widehat b_\nu(k)=n^{-3}\sum_cb_\nu(c)e^{-ik\cdot c}$.
The derivative vertices of \S\ref{subsec:response-contact}
give the blocks
\begin{align*}
 D_1[p+k,p]&=\sum_\nu\widehat b_\nu(k)e^{-ik_\nu/2}
                                          V_\nu(p+k/2),\\
 D_1[p,p+k]&=\sum_\mu
         \overline{\widehat b_\mu(k)e^{-ik_\mu/2}}
                                          V_\mu(p+k/2),\\
 D_2[p,p]&=\sum_\nu\sum_k|\widehat b_\nu(k)|^2W_\nu(p).
\end{align*}
Here $k$ ranges over $\eta\mathcal K_{\eta,R}$; the last
identity uses Parseval for the real function $b_\nu$.
For example, the positive hop in $D_1[p+k,p]$ is
$-i\widehat b_\nu(k)e^{ip_\nu}T_\nu^+$, and the negative hop
is $i\widehat b_\nu(k)e^{-i(p_\nu+k_\nu)}T_\nu^-$.
Their sum is the stated midpoint vertex. Under $k_\nu\mapsto
k_\nu+2\pi$ both the phase and that vertex change sign, so
their product is well-defined on momentum representatives.

Unfolding the periodization gives, for $k=\eta K$,
\[
 \widehat b_\nu(k)e^{-ik_\nu/2}
     =n^{-3}\eta\sum_{c\in\Z^3}
          a_{\eta,\nu}^0(\eta c+\eta e_\nu/2)
                            e^{-iK\cdot(\eta c+\eta e_\nu/2)}
     =\eta R^{-3}h_{\eta,\nu}(K).
\]
Substitute these three block identities into the two traces.
The scalar normalization is
$\eta^2R^{-6}=R^{-3}\eta^{-1}n^{-3}$, which proves
\eqref{eq:response-polarization-identity}. This identity
requires no separation of the probe from its periodic images.

\begin{proposition}\label{prop:response-polarization-bounds}
There is $C_\rho<\infty$, independent of $0<\bar m\le1$,
$0<\eta\le1$, $R/\eta\in\N$, $R/\eta\ge2$, and $K\in\R^3$,
such that
\begin{align}
 |\Pi^{\eta,R}_{\mu\nu}(K)-\Pi^\eta_{\mu\nu}(K)|
       &\le C_\rho\bar m^{-3}R^{-4},
                            \label{eq:response-polarization-quadrature}\\
 |\Pi^\eta_{\mu\nu}(K)|&\le C_\rho\bar m^{-1}|K|^2,\qquad
 |\nabla_K\Pi^\eta_{\mu\nu}(K)|
                    \le C_\rho\bar m^{-1}|K|.
                            \label{eq:response-polarization-quadratic}
\end{align}
\end{proposition}

\begin{proof}
We first estimate the quadrature error. If $F$ is a $C^4$
periodic function on the momentum torus, its Fourier coefficients
satisfy
\[
 |\widehat F(z)|\le |z|_\infty^{-4}
       \max_j\int_{[-\pi,\pi]^3}|\partial_j^4F(p)|
                                      \frac{\dd p}{(2\pi)^3},
                                                \qquad z\ne0.
\]
This follows by four integrations by parts in a coordinate attaining
$|z|_\infty$. The coefficients are absolutely summable, since
$\sum_{z\ne0}|z|_\infty^{-4}<42$. Their Fourier series equals
$F$: its difference from $F$ has zero Fourier coefficients,
and density of trigonometric polynomials in continuous periodic
functions makes the difference zero. The antiperiodic mean
selects the coefficients at $nz$ with sign $(-1)^{z_1+z_2+z_3}$.
Consequently
\begin{equation}\label{eq:response-quadrature-basic}
 \left|n^{-3}\sum_{p\in\mathcal P_n}F(p)
                  -\int_{[-\pi,\pi]^3}F(p)\frac{\dd p}{(2\pi)^3}\right|
 \le42n^{-4}\max_j
        \int_{[-\pi,\pi]^3}|\partial_j^4F(p)|
                                      \frac{\dd p}{(2\pi)^3}.
\end{equation}

Set $x(p)=\operatorname{dist}(p,2\pi\Z^3)+m$. On a period
cube $x(p)=|p|+m$ and $x(p)<7$. The inverse derivative
estimate \eqref{eq:response-inverse-derivatives}, used up
to order four, and the bounded symbol derivatives give
\begin{equation}\label{eq:response-polarization-fourth-derivative}
 |\partial_j^4L_{\mu\nu}(p,k)|
                 \le C_\rho[x(p)^{-6}+x(p+k)^{-6}].
\end{equation}
For the contact term, each Leibniz summand has at most five
powers of $x(p)^{-1}$; the remaining power is supplied by
$x(p)<7$. For the bubble term, a summand with $a$ derivatives
on the first inverse and $b$ on the second has the factor
$x(p)^{-1-a}x(p+k)^{-1-b}$, with $a+b\le4$.
The other derivatives fall on bounded vertices. Since
$u^{-s}v^{-t}\le u^{-s-t}+v^{-s-t}$ for $s,t>0$,
and $x<7$, the factor is bounded by a constant times the
right side of \eqref{eq:response-polarization-fourth-derivative}.
This bound is uniform in the shift $k$.

For $j>3$, periodicity and the radial integral give
\begin{equation}\label{eq:response-translated-integral}
 \int_{[-\pi,\pi]^3}x(p+u)^{-j}\,\dd p
 \le4\pi\int_0^\infty\frac{r^2\,\dd r}{(r+m)^j}
       =\frac{8\pi m^{3-j}}{(j-1)(j-2)(j-3)},
                                             \qquad u\in\R^3.
\end{equation}
Indeed the integral over a translated period cube is unchanged,
and on the original cube it is bounded by the integral over $\R^3$.
The case $j=6$ bounds the integral in
\eqref{eq:response-quadrature-basic} by $C_\rho m^{-3}$.
The definitions of the two polarizations then give the first
estimate, because
\[
 \eta^{-1}n^{-4}m^{-3}
           =\eta^{-1}(\eta/R)^4(\bar m\eta)^{-3}
           =\bar m^{-3}R^{-4}.
\]

For the quadratic bound, use the symmetric bubble
\[
 F_{\mu\nu}(p,k)=
 \tr[R_m(p-k/2)V_\mu(p)R_m(p+k/2)V_\nu(p)].
\]
It is periodic in $p$. Translating the period cube by $-k/2$
shows that its integral is the integral of the bubble in
\eqref{eq:response-polarization-integrand}. At $k=0$ the
whole integrand is a total derivative:
\begin{equation}\label{eq:response-polarization-total-derivative}
 L_{\mu\nu}(p,0)=
    \partial_\mu\tr[R_m(p)V_\nu(p)].
\end{equation}
In fact $\partial_\mu R_m=-R_mV_\mu R_m$ and
$\partial_\mu V_\nu=\delta_{\mu\nu}W_\nu$. Its integral
therefore vanishes, including the contact term. We obtain
\begin{equation}\label{eq:response-polarization-subtracted}
 \Pi^\eta_{\mu\nu}(K)=
     -\eta^{-1}\int_{[-\pi,\pi]^3}
          [F_{\mu\nu}(p,\eta K)-F_{\mu\nu}(p,0)]
                                      \frac{\dd p}{(2\pi)^3}.
\end{equation}

Conjugation by the involution in
\eqref{eq:response-bloch-parity} changes $(p,k)$ to $(-p,-k)$
in $F_{\mu\nu}$; the signs of its two vertices cancel.
Thus $F_{\mu\nu}(-p,-k)=F_{\mu\nu}(p,k)$.
Its first $k$ derivative at zero is odd in $p$, and hence
\begin{equation}\label{eq:response-polarization-parity-cancellation}
 \int_{[-\pi,\pi]^3}
              (\partial_{k_a}F_{\mu\nu})(p,0)\,\dd p=0.
\end{equation}
This removes the linear term as well as the constant term
removed by \eqref{eq:response-polarization-total-derivative}.

The Hessian has two derivatives in total on the inverse factors;
the midpoint vertices no longer depend on $k$. With $p_\pm=p\pm k/2$
and the vertices evaluated at $p$, it is
\begin{align*}
 \partial_{k_a}\partial_{k_b}F_{\mu\nu}
 =\tfrac14\big(&
  \tr[(\partial_a\partial_bR_m)(p_-)V_\mu R_m(p_+)V_\nu]\\
 &-\tr[(\partial_aR_m)(p_-)V_\mu(\partial_bR_m)(p_+)V_\nu]\\
 &-\tr[(\partial_bR_m)(p_-)V_\mu(\partial_aR_m)(p_+)V_\nu]\\
 &+\tr[R_m(p_-)V_\mu(\partial_a\partial_bR_m)(p_+)V_\nu]\big).
\end{align*}
Apply \eqref{eq:response-inverse-derivatives} through order two
and the same product inequality used above. It follows that
\[
 |\partial_{k_a}\partial_{k_b}F_{\mu\nu}(p,k)|
               \le C_\rho[x(p-k/2)^{-4}+x(p+k/2)^{-4}],
\]
and \eqref{eq:response-translated-integral} with $j=4$ gives
\begin{equation}\label{eq:response-polarization-hessian}
 \int_{[-\pi,\pi]^3}
     |\partial_{k_a}\partial_{k_b}F_{\mu\nu}(p,k)|
                             \frac{\dd p}{(2\pi)^3}
                  \le C_\rho m^{-1},\qquad k\in\R^3.
\end{equation}

Taylor's formula with integral remainder, followed by
\eqref{eq:response-polarization-parity-cancellation}, now makes
the two momentum powers explicit:
\begin{equation}\label{eq:response-polarization-taylor}
 \Pi^\eta_{\mu\nu}(K)=
   -\eta\int_{[-\pi,\pi]^3}\int_0^1(1-s)
      \sum_{a,b}K_aK_b
       (\partial_{k_a}\partial_{k_b}F_{\mu\nu})(p,s\eta K)
                           \,\dd s\,\frac{\dd p}{(2\pi)^3}.
\end{equation}
This gives the bound $C_\rho(\eta/m)|K|^2$.
For the gradient, differentiate
\eqref{eq:response-polarization-subtracted}, subtract its
zero-momentum derivative using parity, and apply the
first-order integral remainder:
\[
 \partial_{K_a}\Pi^\eta_{\mu\nu}(K)=
  -\eta\int_{[-\pi,\pi]^3}\int_0^1
       \sum_bK_b(\partial_{k_a}\partial_{k_b}F_{\mu\nu})(p,s\eta K)
                              \,\dd s\,\frac{\dd p}{(2\pi)^3}.
\]
All differentiations are under an integral over a compact cube
of a smooth function, since $m>0$. Using
\eqref{eq:response-polarization-hessian} and $\eta/m=\bar m^{-1}$
proves \eqref{eq:response-polarization-quadratic}.
\end{proof}

\subsection{Uniform bounds on the determinant coefficients}
\label{subsec:response-uniform-coefficients}

We finish by pairing the polarization bounds with decay of the
probe transform. This step needs no continuum evaluation of
the two-photon kernel.

For each integer $j\ge0$ there is $C_{\beta,j}<\infty$ such that
\begin{equation}\label{eq:response-midpoint-decay}
 |h_{\eta,\nu}(K)|\le C_{\beta,j}(1+|K|)^{-j},
                \qquad |K|_\infty\le\pi/\eta,\quad0<\eta\le1.
\end{equation}
The constants are unchanged by any real translation of $\beta$.
Here is the discrete integration by parts that makes this
uniform in the mesh. Suppose a smooth function $v$ is supported
in a cube of side $\ell$, and sum it over a translate of $\eta\Z^3$.
The backward difference
$\Delta_{\eta,\mu}v(y)=\eta^{-1}(v(y)-v(y-\eta e_\mu))$
satisfies
$\|\Delta_{\eta,\mu}^jv\|_\infty
\le\|\partial_\mu^jv\|_\infty$, by its integral representation.
Its support is enlarged by at most $j\eta$ in one direction,
and the number of sampling points, multiplied by $\eta^3$,
is at most $(\ell+j+1)^3$. Summation by parts gives
\[
 [\eta^{-1}(1-e^{-i\eta K_\mu})]^j
       \eta^3\sum_yv(y)e^{-iK\cdot y}
          =\eta^3\sum_y\Delta_{\eta,\mu}^jv(y)e^{-iK\cdot y}.
\]
On $|\eta K_\mu|\le3\pi/2$ the modulus of the bracket
is at least $0.3|K_\mu|$. Choose a coordinate attaining
$|K|_\infty$. For $|K|\ge1$ this is at least
$(1+|K|)/12$; for $|K|<1$ use the point-count estimate directly.
We thus obtain, on that momentum cube,
\[
 \left|\eta^3\sum_yv(y)e^{-iK\cdot y}\right|
 \le12^j(\ell+j+1)^3
       \max(\|v\|_\infty,\max_\mu\|\partial_\mu^jv\|_\infty)
                                               (1+|K|)^{-j}.
\]

Apply this to $v=a_{\eta,\nu}^0$ on
$\eta\Z^3+\eta e_\nu/2$. Cube averaging has norm one on
derivative sup norms, and a centered difference is an average
of the corresponding first derivative. Consequently
$\|\partial_\mu^ja_{\eta,\nu}^0\|_\infty
\le2|\beta|_{C^{j+1}}$, while its support has side at most
$D_\beta+2\eta$. One may take
\[
 C_{\beta,j}=2\cdot12^j(D_\beta+j+3)^3|\beta|_{C^{j+1}}.
\]
This proves \eqref{eq:response-midpoint-decay}, with position-independent
constants and no assumption on $R$.

\begin{proposition}\label{prop:response-uniform-coefficients}
For each fixed $\rho$, $\bar m$ and $\beta$ as above,
\[
 |c_2|\le C_{\rho,\bar m,\beta},
              \qquad 0<\eta\le1,\quad R\ge1,\quad
                         R/\eta\in\N,\quad R/\eta\ge2.
\]
If in addition $R\ge D_\beta+1$, then
$|c_4|\le C_{\rho,\bar m,\beta}$.
The bounds are uniform in all real translations of $\beta$,
and the coefficients have no dependence on $e$.
\end{proposition}

\begin{proof}
Proposition~\ref{prop:response-polarization-bounds} gives
\[
 |\Pi^{\eta,R}_{\mu\nu}(K)|
   \le C_\rho(\bar m^{-3}R^{-4}+\bar m^{-1}|K|^2).
\]
In \eqref{eq:response-polarization-identity} use the decay
estimate with $j=2$ for the quadrature term and $j=3$ for
the quadratic term. Since $R\ge1$, both contributions are
bounded by a constant times $(1+|K|)^{-4}$; hence
\[
 |c_2|=|a_{\mathrm{ch}}|
 \le9C_\rho
       (\bar m^{-3}C_{\beta,2}^2+\bar m^{-1}C_{\beta,3}^2)
       R^{-3}\sum_{K\in\mathcal K_{\eta,R}}(1+|K|)^{-4}.
\]
The last sum is bounded uniformly for $R\ge1$. To see this,
tile $\R^3$ with cubes of side $2\pi/R$ centered at
$K\in(2\pi/R)\Z^3$. On each cube,
$1+|P|\le(1+\sqrt3\pi)(1+|K|)$, so its volume times
$(1+|K|)^{-4}$ is bounded by a constant times the integral
of $(1+|P|)^{-4}$ there. Sum over the cubes; that integral
over $\R^3$ is finite. Restricting the sum to
$\mathcal K_{\eta,R}$ only decreases it. This proves the
uniform second-order bound.

For the fourth coefficient, use
\eqref{eq:response-coefficient-relations} and
Proposition~\ref{prop:response-fourth}. The infinite-space
trace $b(f)$ is finite by \eqref{eq:response-four-ideal-infinite},
and the two errors in \eqref{eq:response-fourth-errors} are
bounded for $R\ge D_\beta+1$ and $0<\eta\le1$.
Thus $|b_{\mathrm{ch}}|$ is bounded, and adding $3c_2^2$
gives the asserted bound for $c_4$. All constants in the
proof depend on support sizes and derivative norms of $\beta$,
not on its position. Independence of $e$ follows directly
from the external connection convention.
\end{proof}


\section{Probes on the reference cells}\label{sec:reference-cells}

The external-source identity uses a finitely supported bond cochain.
The continuum cumulant uses the test $f=dd^*\beta$. We identify the
discrete source with a nearby smooth test and place it strictly inside
one reference cell. Copies in distinct cells then have orthogonal
curvatures, which makes the normal-ordering factor separate over cells.

\subsection{The discrete source and the smooth test}

Fix a real nonzero $\beta\in C_c^\infty(\R^3;\Lambda^2)$.
Let $D_\beta$ retain the unbuffered support-cube side of
Section~\ref{subsec:response-conventions}, and write
$D_\beta^{\mathrm{buf}}=D_\beta+4$. Choose a cube
$Q_\beta=z_\beta+[-D_\beta^{\mathrm{buf}}/2,
D_\beta^{\mathrm{buf}}/2]^3$ containing
$\supp\beta+[-2,2]^3$. Use the centred differences
$\delta_{\eta,\mu}$ and cube average $J_\eta$ of that section.
Put
\begin{align}
 a'_{\eta,\nu}&=-\sum_{\mu\ne\nu}
                          \delta_{\eta,\mu}\beta_{\mu\nu},\notag\\
 f^\eta_{\mu\nu}&=\delta_{\eta,\mu}a'_{\eta,\nu}
                         -\delta_{\eta,\nu}a'_{\eta,\mu}.
                                      \label{eq:cells-smooth-test}
\end{align}
All these functions are smooth and real. Since each centred
difference enlarges support by at most $\eta/2$ in its direction,
\[
 \supp f^\eta,\ \supp J_\eta a'_\eta
                   \subset\supp\beta+[-\eta,\eta]^3.
\]
In particular their supports lie in $Q_\beta$ for $0<\eta\le1$.

On the physical lattice $\eta\Z^3$ set
\[
 b_{\eta,\mu\nu}(c)
       =(J_\eta\beta_{\mu\nu})
                       (c+\eta(e_\mu+e_\nu)/2),\qquad
 a_\eta^o=d_\eta^*b_\eta.
\]
Here $d_\eta$ is the mesh-normalized forward cochain differential
and $d_\eta^*$ its adjoint for the pairing $\eta^3\sum$.
The adjoint formula gives
\begin{equation}\label{eq:cells-midpoint-source}
 a_{\eta,\nu}^o(c)=(J_\eta a'_{\eta,\nu})(c+\eta e_\nu/2).
\end{equation}
Thus $J_\eta a'_\eta$ is the interpolating vector field
$a_\eta^0$ already used in the free response.
Transport to the unit lattice and define
\[
 \hat{\mathfrak a}_\eta^o(x,x+e_\nu)
                    =\eta^{1/2}a_{\eta,\nu}^o(\eta x).
\]
Its coboundary uses unscaled unit-lattice differences. The
midpoint identities yield
\begin{equation}\label{eq:cells-source-curvature}
 d\hat{\mathfrak a}_\eta^o(\sigma)
                        =f^\eta_{[\eta]}(\sigma).
\end{equation}
Indeed, the bond values in the curl are evaluated at the two
points displaced by $\pm\eta e_\mu/2$ from the plaquette centre.
Their difference is $\eta\delta_{\eta,\mu}J_\eta a'_{\eta,\nu}$.
Commutation of the average and differences then gives
$\eta^{3/2}J_\eta f^\eta$ at that centre, which is the right
side of \eqref{eq:cells-source-curvature}.

Periodize the cochain on a torus $C(n)$ of physical side
$R=n\eta$. Formula \eqref{eq:limit-step-pairing} gives the
exact observable identity
\begin{equation}\label{eq:cells-observable}
 \langle d\hat{\mathfrak a}_\eta,F^{\mathrm u}\rangle
                            =F^\sharp(f^\eta).
\end{equation}
The source is transverse: $a_\eta^o=d_\eta^*b_\eta$ is
orthogonal to every closed bond cochain, by summation by
parts. This applies to both the covering cochain and its
periodization.

The continuum approximation has the uniform bound
\begin{equation}\label{eq:cells-test-error}
 \|f^\eta-f\|_{H^s}\le\eta^2\|\beta\|_{H^{s+4}}
                         \qquad(s\in\R,\ 0<\eta\le1).
\end{equation}
To see this, the difference symbol is
$is_\mu(p)$ with $s_\mu(p)=2\eta^{-1}\sin(\eta p_\mu/2)$.
One has $|s_\mu|\le|p_\mu|$ and
$|s_\mu-p_\mu|\le\eta^2|p_\mu|^3/24$, so
\[
 |s_\mu s_\nu-p_\mu p_\nu|
                               \le\eta^2|p|^4/12.
\]
Each component of $f^\eta-f$ contains four such terms.
Cauchy--Schwarz on these terms, followed by summation over
the three two-form components, gives
$|\widehat{f^\eta-f}|^2\le(\eta^2/3)^2|p|^8|\hat\beta|^2$.
Plancherel's theorem proves \eqref{eq:cells-test-error}.

\subsection{Reference geometry and placement}

Let $T^*\ge1$ be the reference exponent furnished by the
construction, and put
\[
 R^*=hL^{T^*},\qquad n_*=R^*/\eta,\qquad
 p=n_*/L,\qquad w=p/2,\qquad \ell_c=\eta w=R^*/(2L).
\]
Since $h$ and $\eta^{-1}$ are powers of the odd integer $L$,
$p$ is an odd integer, $w\in\Z+1/2$, and the open unit-lattice
reference cells are $Q_c^\circ=w(c+(0,1)^3)$.
The reference torus $C(n_*)$ has $2L$ cells per side and
$N_*=8L^3$ in all.
The second reference torus has side $Ln_*$.
Every cutoff torus $n=8hL^{T+N}$ with $T\ge T^*+1$
is an integer multiple of $Ln_*$. These are in the range
of Proposition~\ref{prop:transfer-general}, with
$s_0=2L$ and $j=L$.

The probe fits the reference cells if
\begin{equation}\label{eq:cells-fit}
 D_\beta^{\mathrm{buf}}+3\eta\le\ell_c.
\end{equation}
A placement is a vector $v\in\eta\Z^3$ satisfying
\[
 |v+z_\beta-(\ell_c/2)(1,1,1)|_\infty\le\eta/2.
\]
Such a vector always exists by coordinatewise rounding.
It translates $Q_\beta$ into $[\eta,\ell_c-\eta]^3$:
the displacement of its centre is at most $\eta/2$, and
the half-side is at most $(\ell_c-3\eta)/2$.
Hence
\begin{equation}\label{eq:cells-probe-margin}
 \supp(\tau_v\beta)
                     \subset[2+\eta,\ell_c-2-\eta]^3.
\end{equation}
Form the source $\hat{\mathfrak a}^o$ with this translated
probe.

Every endpoint of a bond in the source support is within
$3\eta/2$ in sup distance of $\supp(\tau_v\beta)$.
By \eqref{eq:cells-probe-margin}, its lattice coordinates
therefore lie in
$[2/\eta-1/2,w-2/\eta+1/2]$.
If a plaquette has nonzero $d\hat{\mathfrak a}^o$, at least
one of its boundary bonds belongs to that support. Its
other two corners are one unit away. Thus the whole
plaquette lies in the box
\begin{equation}\label{eq:cells-support-box}
 B_0=(w/2)(1,1,1)+[-r_0,r_0]^3,\qquad
 r_0=w/2-2/\eta+3/2.
\end{equation}
Condition \eqref{eq:cells-fit} gives $r_0\ge3$.
The distance of this box from each bounding wall is
$2/\eta-3/2\ge1/2$, so $B_0\subset Q_0^\circ$.
Its reflected copies lie strictly inside the corresponding
cells. This is the support margin needed by the local
source identity.

\subsection{Orthogonal sources and translated observables}

Let $g_c$ be the wall-group element carrying $Q_0^\circ$ to
$Q_c^\circ$. Put
$\hat{\mathfrak a}_c^o=(g_c)_*\hat{\mathfrak a}^o$ on the
covering lattice and periodize it on $C(n)$, where $p$ divides
$n$ and $n\ge n_*$, to obtain $\hat{\mathfrak a}_c$.
The period lattice identifies cell indices modulo
$(2n/p)\Z^3$. The coboundary commutes with these actions
and periodization, and its support lies in $B_c=g_cB_0$.
For cell representatives on this torus,
\begin{equation}\label{eq:cells-orthogonality}
 \langle d\hat{\mathfrak a}_c,d\hat{\mathfrak a}_{c'}\rangle_n
             =v_\eta\,\mathbf1_{\{c=c'\}},\qquad
 v_\eta=\|d\hat{\mathfrak a}^o\|^2.
\end{equation}
The covering norm is the unweighted $\ell^2$ norm, and the
torus pairing is the counting pairing.

To prove the identity, unfold one factor:
\[
 \langle d\hat{\mathfrak a}_c,d\hat{\mathfrak a}_{c'}\rangle_n
  =\sum_{\sigma}\sum_{z\in\Z^3}
       d\hat{\mathfrak a}_c^o(\sigma)
       d\hat{\mathfrak a}_{c'}^o(\sigma+nz).
\]
A nonzero summand has $\sigma\subset Q_c^\circ$ and
$\sigma+nz\subset Q_{c'}^\circ$. The open cells are disjoint,
so necessarily $c'=c+(2n/p)z$. For distinct torus cells
there is no such $z$. For the same cell exactly one translate
can contribute; its summands are the squares of
$d\hat{\mathfrak a}_c^o$. The wall action preserves the
counting norm, which proves \eqref{eq:cells-orthogonality}.
Changing the placement by a vector in $\eta\Z^3$ translates
the covering source by an integer vector. Thus $v_\eta$
depends only on $\beta$ and $\eta$, and is independent of
the torus and the placement.

For a closed cochain $u$ define
$X_c(u)=\langle d\hat{\mathfrak a}_c,u\rangle_n$.
Every polynomial in this observable is read by one and the
same local integrand on all these tori. On the cover its
local expression is
\[
 \widehat A\longmapsto
 P\left(\sum_{\sigma\subset B_c}
             d\hat{\mathfrak a}_c^o(\sigma)(d\widehat A)(\sigma)\right).
\]
This is a scalar gauge-invariant polynomial. Cauchy--Schwarz
bounds its argument by
\[
 \sqrt{v_\eta}\left(\sum_{\sigma\subset B_c}
                      |d\widehat A(\sigma)|^2\right)^{1/2}.
\]
The polynomial is therefore bounded by
\[
 C_{\mathrm{pol}}\exp\left(\frac14\sum_{\sigma\subset B_c}
                                      |d\widehat A(\sigma)|^2\right)
\]
for a finite constant $C_{\mathrm{pol}}$ depending on the polynomial
and the covering source. This is the growth condition for
$\mathcal A(Q_c^\circ)$.

The local boxes embed in the torus. Indeed, writing
$p=2m_p+1$, condition \eqref{eq:cells-fit} gives $m_p\ge7$,
and
\[
 2r_0+2=w-4/\eta+5\le w+1<p\le n.
\]
For each $c$, the translates of $B_c+[-1,1]^3$ by
$nz$, $z\in\Z^3$, are consequently disjoint.
On a local lift, $d\widehat A=F^{\mathrm u}$; unfolding
the finite cochain sum reads the polynomial as $P(X_c)$.
For a real polynomial its reflected copy from the base
cell is exactly $P(X_c)$, since conjugation fixes its
coefficients. This applies to the centered Wick polynomials.

The variance in \eqref{eq:cells-orthogonality} approaches
the continuum normal-ordering variance:
\begin{equation}\label{eq:cells-variance-limit}
 \left|\sqrt{v_\eta}-\|f\|_2\right|
 \le21\eta|\beta|_{C^3}(D_\beta^{\mathrm{buf}})^{3/2}
                      +\eta^2\|\beta\|_{H^4}.
\end{equation}
To see this, first \eqref{eq:cells-source-curvature} gives
$v_\eta=|(f^\eta)_\eta|_\eta^2$. Let $\Pi_\eta f^\eta$
be its componentwise cell-average step function on a
period cube. Its $L^2$ norm is $\sqrt{v_\eta}$.
Each first derivative of $f^\eta_{\mu\nu}$ is bounded
by $4|\beta|_{C^3}$, from the four centred-difference
terms in \eqref{eq:cells-smooth-test} and the identity
$\delta_{\eta,\mu}g=\int_{-1/2}^{1/2}
(\partial_\mu g)(\,\cdot+t\eta e_\mu)\,\dd t$.
The pointwise error of the cell average is at most
$12\eta|\beta|_{C^3}$. It vanishes outside the periodic
images of the buffered cube, whose measure in one period
is $(D_\beta^{\mathrm{buf}})^3$.
Summing the three component squares bounds the $L^2$
error by $21\eta|\beta|_{C^3}(D_\beta^{\mathrm{buf}})^{3/2}$.
The norm triangle inequality and
\eqref{eq:cells-test-error} prove
\eqref{eq:cells-variance-limit}.

Finally the even cell indices give pure translations.
For $y\in(R^*/L)\Z^3=p\eta\Z^3$ set
$c(y)=2y/(p\eta)\in2\Z^3$. Then $g_{c(y)}$ is the unit-lattice
translation by $y/\eta$, and
\begin{equation}\label{eq:cells-translation-observable}
 X_{c(y)}=F^\sharp(\tau_{v+y}f^\eta).
\end{equation}
Two such cells agree on $C(n)$ exactly when their physical
translations agree modulo $R\Z^3$. Any fixed finite set
of distinct translations therefore occupies distinct
cells once the volume is sufficiently large.
Lattice translation invariance removes the common placement
$v$ from the expectation of any polynomial in these
observables. Its integrability follows from Gaussian
domination. Thus the reference-cell product estimates
give estimates for the translates $\tau_yf^\eta$ used
in Section~\ref{sec:ultraviolet-limits}.


\section{Exact external-response identities}\label{sec:external-response}

We express the generating function of a curvature observable as a
determinant in an external connection. The identity is first proved
in each integer flux sector. In the zero-flux sector an unfolding
of the integer cochains replaces the compact photon integral by a
Gaussian integral over transverse cochains and an average over the
three twists. We compute its normalization before applying the
source shift.

\subsection{The determinant at complex connections}
\label{subsec:external-complex-determinant}

Use lattice units and fix $q>0$, $r_q=2\pi/q$ and a real lattice
mass $\mu_{\mathrm{lat}}$. Put
$M_{\mathrm{ch}}=\mu_{\mathrm{lat}}+\frac{\rho}{2}\sum_iw_i$.
For the cutoff laws of Theorem~\ref{thm:main},
$\mu_{\mathrm{lat}}=\boldsymbol\mu_N$.
The weights are the standard chamber weights, and $0<\rho<1$.
For an integer triple $n=(n_1,n_2,n_3)$ let
\[
 V=|n|=n_1n_2n_3,\qquad
 B_n=\prod_{\mu=1}^3\{0,\ldots,n_\mu-1\},\qquad
 \mathcal L_n=B_n\times\{1,2,3\}.
\]
In this subsection each $n_\mu$ may be as small as one.
Cochains are periodic functions on the oriented covering cells,
so the $3V$ link orbits remain distinct even at those small sides.
The chamber spinors are antiperiodic in each direction and have
dimension $192V$. Write $D_B$ for the massless chamber operator
with phases $\exp(iqB)$, and $\det_n$ for its finite-dimensional
determinant. In the phase convention of
\S\ref{subsec:response-conventions}, $D_B=D^{\mathrm{ch}}[qB]$.
The weight $1/48$ in the spinor pairing does not rescale this
determinant.

The determinant is a polynomial in the $6V$ variables
$\exp(iqB(b))$ and $\exp(-iqB(b))$, of degree at most $32$
in each. In particular it is entire in the complex link values.
To see the degree bound, choose the site basis on
$B_n\times O_{\mathrm{cub}}$ and include the antiperiodic sign
whenever a neighbor is reduced into $B_n$. Each of the four
rows at a site $(c,g)$ has constant entries except for the
one external-hop block. If $ge_1=e_\mu$, this block contains
$\exp(iqB_\mu(c))$; if $ge_1=-e_\mu$, it contains
$\exp(-iqB_\mu(c-e_\mu))$. For a fixed oriented photon link
there are eight labels at each end. Thus each phase variable
occurs in exactly $8\cdot4=32$ rows. The Leibniz determinant
chooses one entry in each row, proving the bound. The five
site labels of the diagonal and the four hops are distinct,
even if a side is one, so no other phase is hidden in one
of these blocks.

The two invariances needed below hold for complex cochains:
\begin{equation}\label{eq:external-determinant-invariances}
 \det_n(D_{B+r_qk}+\mu_{\mathrm{lat}})
       =\det_n(D_B+\mu_{\mathrm{lat}}),\qquad
 \det_n(D_{B+d\lambda}+\mu_{\mathrm{lat}})
       =\det_n(D_B+\mu_{\mathrm{lat}}),
\end{equation}
where $k$ is integer valued and $\lambda$ is a periodic complex
zero-cochain. The first identity follows directly from the phases.
For the second, multiplication by $\exp(iq\lambda(c))$ preserves
antiperiodicity and has inverse multiplication by
$\exp(-iq\lambda(c))$. On an external hop the factor
\[
 e^{iq(B-d\lambda)(c,c')}e^{iq\lambda(c')}
                                  =e^{iq\lambda(c)}e^{iqB(c,c')}
\]
gives $D_{B-d\lambda}=e^{iq\lambda}D_Be^{-iq\lambda}$.
Internal hops lie over one photon site and give the same
conjugation. This proves the determinant identity.

For $s\ge0$ define
\begin{align}
 K_n(s)&=
 \left(|M_{\mathrm{ch}}|+\frac{1+\rho}{2}
           \left(\sum_iw_i-6+6e^{qs}\right)\right)^{192V},
           \notag\\
 |\det_n(D_B+\mu_{\mathrm{lat}})|
                   &\le K_n(\|\operatorname{Im}B\|_\infty).
                   \label{eq:external-determinant-strip}
\end{align}
Indeed each chamber bond involution permutes the site representatives.
For an internal hop its phase has modulus one; for an external
hop the modulus is at most $e^{q\|\operatorname{Im}B\|_\infty}$.
The hopping matrix has norm $w_i(1+\rho)/2$. The operator norm
is therefore bounded by the base of the power in
\eqref{eq:external-determinant-strip}, and a determinant in dimension
$192V$ is bounded by that norm to the same power.
This estimate is uniform in the real part of the connection.

The twist connection is
\begin{equation}\label{eq:external-twist}
 H_\theta=\sum_\mu\frac{\theta_\mu}{qn_\mu}\mathbf1_\mu,
 \qquad
 \mathbf1_\mu=\sum_{c\in B_n}\mathbf1_{(c,\mu)}.
\end{equation}
Its determinant has period $2\pi$ in each complex twist variable.
For later use we need the cochain implementing this period. Set
\[
 k^{(\mu)}_\nu(c)=\delta_{\mu\nu}1_{\{c_\mu=n_\mu-1\}},
 \qquad
 \lambda^{(\mu)}(c)=-\frac{r_q}{n_\mu}c_\mu,\qquad c\in B_n,
\]
with periodic extension. On links not crossing the $\mu$ seam,
$d\lambda^{(\mu)}=-r_q/n_\mu$ in direction $\mu$ and zero
in the other directions. On a seam link its value is
$r_q-r_q/n_\mu$. Thus
\begin{equation}\label{eq:external-twist-period}
 H_{2\pi e_\mu}=r_qk^{(\mu)}-d\lambda^{(\mu)}.
\end{equation}
The two invariances in
\eqref{eq:external-determinant-invariances} prove the asserted
twist period. For $n_\mu=1$, $\lambda^{(\mu)}=0$ and
$k^{(\mu)}=\mathbf1_\mu$, so the same formula holds.

\subsection{Real and integer cochains on the torus}
\label{subsec:external-cochains}

From now on $n_\mu\ge4$. The counting pairings on cochains
are extended bilinearly when their arguments are complex.
Use
\[
 V_{\mathrm g}=dC^0(n),\qquad
 V_{\mathrm h}=\operatorname{span}_{\R}\{\mathbf1_1,\mathbf1_2,\mathbf1_3\},
 \qquad
 V_{\mathrm t}=(V_{\mathrm g}+V_{\mathrm h})^\perp.
\]
We recall the cochain arguments needed for the unfolding and
for its Jacobian. In particular the integer condition in the
flux sum requires an integer primitive, not merely a real one.

Let
\[
 \mathcal P_n^{\mathrm{per}}=
   \left\{p:p_\mu=2\pi j_\mu/n_\mu,\ 0\le j_\mu<n_\mu\right\},
 \qquad e_p(c)=V^{-1/2}e^{ip\cdot c},\qquad
 \xi_\mu(p)=e^{ip_\mu}-1.
\]
The characters $e_p$ are orthonormal. At each momentum $p$ the
coboundary on exterior-valued cochains is exterior multiplication
by $\xi(p)$. If $\iota_{\bar\xi}$ denotes contraction by $\bar\xi$,
then
\[
 \iota_{\bar\xi}(\xi\wedge u)+
                  \xi\wedge\iota_{\bar\xi}u=|\xi|^2u.
\]
For $p\ne0$ a closed coefficient $u$ is consequently
$\xi\wedge(|\xi|^{-2}\iota_{\bar\xi}u)$, and is exact.
At zero momentum the cochains are constant and the
coboundary vanishes. Nonzero and zero modes are orthogonal,
and taking real parts of a complex primitive gives a real
primitive for a real cochain. Therefore
\begin{equation}\label{eq:external-Hodge}
 \ker(d:C^1\to C^2)=V_{\mathrm g}\oplus V_{\mathrm h},
 \qquad C^1=V_{\mathrm g}\oplus V_{\mathrm h}\oplus V_{\mathrm t},
 \qquad
 \ker(d:C^2\to C^3)=dC^1\oplus\operatorname{Harm}^2.
\end{equation}
Here $\operatorname{Harm}^2$ is the space of constant two-cochains.
The same Fourier formula gives $d^2=0$.

For a closed integer two-cochain $z$, define
$w_{\mu\nu}(z)$ as its sum over one coordinate $(\mu,\nu)$
two-torus. The sum is independent of the transverse coordinate:
sum $dz=0$ over the layer of cubes between two consecutive
planes. Differences in the two tangential directions cancel
by periodicity; the remaining difference is that of the two
plane sums. The fluxes give
\begin{equation}\label{eq:external-integer-two-exactness}
 z\in dC^1(n;\Z)\quad\Longleftrightarrow\quad
                         w_{\mu\nu}(z)=0\ \text{for all }\mu<\nu.
\end{equation}
The integer coefficients follow from the
cellular cohomology of the product of the three subdivided
circles. Its cochain complex is precisely the integer complex
used here. The three coordinate two-tori give a basis of
$H_2(T^3;\Z)$, and $H_1(T^3;\Z)$ is free. The universal
coefficient theorem identifies $H^2(T^3;\Z)$ with
$\operatorname{Hom}(H_2(T^3;\Z),\Z)$ by evaluation.
The class of $z$ evaluates to $w_{\mu\nu}(z)$ on each basis
two-torus, so it vanishes exactly when these sums vanish.
Vanishing of this cellular class is membership in
$dC^1(n;\Z)$.
\cite[Theorems~3.2, 3.5 and~3B.6]{Hatcher2002}

We next describe the closed integer one-cochains. If $d\lambda$
is integer valued, its sum along a path from $0$ to $c$ shows
that $\lambda(c)-\lambda(0)$ is integer valued. If $d\lambda=0$,
the same path argument shows that $\lambda$ is constant.
It follows that
\[
 C^1(n;\Z)\cap V_{\mathrm g}=dC^0(n;\Z).
\]
Each seam cochain $k^{(\mu)}$ in
\eqref{eq:external-twist-period} is closed, because its
multiple by $r_q$ lies in $V_{\mathrm g}+V_{\mathrm h}$.
For an arbitrary closed integer one-cochain $k$,
\eqref{eq:external-Hodge} gives
$k=d\lambda+\sum_\mu t_\mu\mathbf1_\mu$.
Its sum along a periodic $\mu$ loop is $n_\mu t_\mu$,
an integer $j_\mu$. Subtracting
$\sum_\mu j_\mu k^{(\mu)}$ leaves an integer exact cochain,
so
\begin{equation}\label{eq:external-closed-integer-one}
 k=d\lambda'+\sum_\mu j_\mu k^{(\mu)},\qquad
 \lambda'\in C^0(n;\Z),\quad j\in\Z^3.
\end{equation}
The harmonic projection of $k$ determines $j$, and then
determines $d\lambda'$. Moreover
\begin{equation}\label{eq:external-integer-Hodge}
 r_qk=
   \left(r_qd\lambda'+\sum_\mu j_\mu d\lambda^{(\mu)}\right)
                         +H_{2\pi j}.
\end{equation}
The first term is exact and the second is harmonic. Their
coupling through $j$ will matter when choosing a fundamental
domain.

The scalar Laplacian $\Delta_0=d^*d$ has eigenvalues
$|\xi(p)|^2=\sum_\mu(2-2\cos p_\mu)$ and a one-dimensional
kernel. Write
\begin{equation}\label{eq:external-scalar-determinant}
 \det{}'\Delta_0=\prod_{p\in\mathcal P_n^{\mathrm{per}}\setminus\{0\}}
                       |\xi(p)|^2,\qquad
 \mathcal L(n)=\log\det{}'\Delta_0.
\end{equation}
At each nonzero momentum the transverse one-cochains
are the two-dimensional space $\xi(p)^\perp$.
For $v\in\C^3$,
\[
 \sum_{\mu<\nu}|\xi_\mu v_\nu-\xi_\nu v_\mu|^2
           =|\xi|^2|v|^2-|\bar\xi\cdot v|^2.
\]
Thus $d^*d$ acts on that transverse space as multiplication
by $|\xi(p)|^2$, and
\begin{equation}\label{eq:external-transverse-determinant}
 \dim V_{\mathrm t}=2V-2,\qquad
 \det(d^*d|_{V_{\mathrm t}})=(\det{}'\Delta_0)^2.
\end{equation}
Diagonalizing this positive real form gives its Gaussian integral
\begin{equation}\label{eq:external-transverse-Gaussian}
 \mathcal G_n:=\int_{V_{\mathrm t}}e^{-\|dA\|^2/2}\,\dd\sigma(A)
                    =\frac{(2\pi)^{V-1}}{\det{}'\Delta_0}.
\end{equation}
The measure $\sigma$ is Euclidean Lebesgue measure for the
counting pairing, not a coordinate measure chosen along a
nonorthogonal basis.

There is one more determinant needed for the exact directions.
The cochains $d\mathbf1_c$, $c\in B_n\setminus\{0\}$,
form both a real basis of $V_{\mathrm g}$ and an integer
basis of $dC^0(n;\Z)$. Indeed one subtracts the value at
zero from a zero-cochain; independence follows because
the resulting zero-cochain has value zero there and
cannot be a nonzero constant. Their Gram determinant is
\begin{equation}\label{eq:external-exact-Gram}
 \det\big(\langle d\mathbf1_c,d\mathbf1_{c'}\rangle\big)_{c,c'\ne0}
                           =\frac{\det{}'\Delta_0}{V}.
\end{equation}
To prove this formula, let $S$ be the scalar Laplacian matrix
in the site basis. The coefficient of $t$ in $\det(t+S)$
is $\det{}'\Delta_0$, by its eigenvalues. Multilinearity
of the determinant in columns identifies that coefficient
with the sum of the $V$ principal minors obtained by
deleting one row and column. All these minors have the same
determinant, since site translations conjugate $S$ by
permutations. The minor at zero is the Gram matrix above.

\subsection{Unfolding the zero-flux sector}
\label{subsec:external-zero-flux}

Let $\Phi$ be a Borel function on $C^2(n)$ with
$|\Phi(u)|\le C_\Phi e^{C_\Phi\|u\|}$.
Let $W$ be a bounded continuous complex function on
$C^1(n)$, invariant under addition of real exact cochains
and of $r_q$ times integer one-cochains. Then
\begin{align}
 &\sum_{\substack{z\in C^2(n;\Z),\,dz=0\\w(z)=0}}
       \int_{[0,r_q)^{\mathcal L_n}}
            \Phi(dA+r_qz)e^{-\|dA+r_qz\|^2/2}W(A)\,\dd A
       \notag\\
 &\quad=\mathcal N_n
       \int_{[-\pi,\pi]^3}\frac{\dd\theta}{(2\pi)^3}
       \int_{V_{\mathrm t}}
              \Phi(dA)e^{-\|dA\|^2/2}W(A+H_\theta)\,\dd\sigma(A),
                                     \label{eq:external-zero-unfolding}\\
 &\mathcal N_n=r_q^{V+2}(\det{}'\Delta_0)^{1/2}.
                                     \label{eq:external-zero-normalization}
\end{align}
Both sides converge absolutely. The coefficient comes from
the lattice of exact gauge transformations and the three
periods of the harmonic directions; we give the unfolding
with both contributions.

First, the Gaussian sum \eqref{eq:reg-gaussian-sum} gives
\[
 \sum_{z\in C^2(n;\Z)}\int_{[0,r_q)^{\mathcal L_n}}
              e^{-\epsilon\|dA+r_qz\|^2}\,\dd A
                     \le(r_qC_q(\epsilon))^{3V}.
\]
Here the integer closedness and flux restrictions have
been dropped. Completing the scalar square bounds
$|\Phi(u)|e^{-\|u\|^2/2}$ by
$C_\Phi e^{C_\Phi^2}e^{-\|u\|^2/4}$.
Together with boundedness of $W$, this proves absolute
convergence of the left side.

Let $\mathcal Z^1_{\Z}=\ker d\cap C^1(n;\Z)$ and
choose representatives $\mathcal R$ of
$C^1(n;\Z)/\mathcal Z^1_{\Z}$. By
\eqref{eq:external-integer-two-exactness}, $k\mapsto dk$
is a bijection from $\mathcal R$ onto the closed
zero-flux integer two-cochains. Put
\[
 P(A)=\Phi(dA)e^{-\|dA\|^2/2}W(A),\qquad
 \mathcal D_1=\bigcup_{k\in\mathcal R}
                     ([0,r_q)^{\mathcal L_n}+r_qk).
\]
The two invariances of $W$ turn the left side of
\eqref{eq:external-zero-unfolding} into
$\int_{\mathcal D_1}P(A)\,\dd A$.
The boxes are disjoint. Every $A$ has a unique decomposition
into its coordinatewise remainder in $[0,r_q)$ and
$r_q$ times an integer cochain. Decomposing that cochain
into its chosen representative and a closed one shows
that $\mathcal D_1$ is a Borel fundamental domain for
$r_q\mathcal Z^1_{\Z}$. The function $P$ is invariant
under that group.

We replace this domain by
\[
 \mathcal F_{\mathrm g}
 =\left\{r_q\sum_{c\ne0}t_c\,d\mathbf1_c:0\le t_c<1\right\},
 \qquad
 \mathcal F_{\mathrm h}=\{H_\theta:\theta\in[-\pi,\pi)^3\},
 \qquad
 \mathcal D_2=\mathcal F_{\mathrm g}+\mathcal F_{\mathrm h}
                                      +V_{\mathrm t}.
\]
To check that it is a fundamental domain for the same
group, write $A=A_{\mathrm g}+A_{\mathrm h}+A_{\mathrm t}$.
Formula \eqref{eq:external-integer-Hodge} shows that
subtracting $r_qk$ changes the harmonic part by
$H_{2\pi j}$ and the exact part by
$\sum_\mu j_\mu d\lambda^{(\mu)}+r_qd\lambda'$.
There is exactly one $j$ placing the harmonic part
in $\mathcal F_{\mathrm h}$. For this $j$ there is
exactly one $d\lambda'$ placing the exact part
in $\mathcal F_{\mathrm g}$, by the integer basis
used in \eqref{eq:external-exact-Gram}. The transverse
part is unchanged. Uniqueness in
\eqref{eq:external-closed-integer-one} gives uniqueness
of $k$. The orthogonal decomposition identifies
$\mathcal D_2$ with a product of Borel sets, so it
is Borel.

For an integrable periodic Borel function, the integral
over two such fundamental domains is the same.
Indeed partition $\mathcal D_1$ into its intersections
with $\mathcal D_2+h$, $h\in r_q\mathcal Z^1_{\Z}$,
and translate each intersection by $-h$. They partition
$\mathcal D_2$, and invariance of both the measure and
the function gives equality. Tonelli proves it first
for nonnegative functions; applying it to the modulus
then justifies the complex case by absolute integrability.
Thus $\int_{\mathcal D_1}P=\int_{\mathcal D_2}P$.

The two finite-dimensional volumes are
\begin{align}
 \operatorname{vol}(\mathcal F_{\mathrm g})
    &=r_q^{V-1}(\det{}'\Delta_0/V)^{1/2},
                                     \label{eq:external-gauge-volume}\\
 \dd A_{\mathrm h}
    &=V^{1/2}q^{-3}\,\dd\theta
      =r_q^3V^{1/2}\frac{\dd\theta}{(2\pi)^3}.
                                     \label{eq:external-twist-Jacobian}
\end{align}
The first follows from the square root of the Gram
determinant. For the second, the constant cochains
$V^{-1/2}\mathbf1_\mu$ are orthonormal, and the
matrix of $\theta\mapsto H_\theta$ in this basis
has diagonal entries $V^{1/2}/(qn_\mu)$.
Their product is $V^{1/2}q^{-3}$.
Their combined factor is precisely $\mathcal N_n$.
Since $P(A_{\mathrm g}+A_{\mathrm h}+A_{\mathrm t})
=\Phi(dA_{\mathrm t})e^{-\|dA_{\mathrm t}\|^2/2}
W(A_{\mathrm t}+A_{\mathrm h})$, it is independent
of $A_{\mathrm g}$. Fubini over the orthogonal
decomposition and the two displayed volume formulas
prove \eqref{eq:external-zero-unfolding}, including
absolute convergence of its right side.

Apply this identity to
$W(A)=\det_n(D_A+\mu_{\mathrm{lat}})$. Its hypotheses
follow from \eqref{eq:external-determinant-invariances}
and \eqref{eq:external-determinant-strip}. If
$Z_n^{(0)}[\Phi]$ is the zero-flux part of the
configuration integral, we have proved
\begin{equation}\label{eq:external-zero-sector}
 Z_n^{(0)}[\Phi]=\mathcal N_n
       \int\frac{\dd\theta}{(2\pi)^3}
       \int_{V_{\mathrm t}}\Phi(dA)e^{-\|dA\|^2/2}
          \det_n(D_{A+H_\theta}+\mu_{\mathrm{lat}})\,\dd\sigma(A).
\end{equation}
The twist integral always runs over $[-\pi,\pi]^3$.
For real $A$ and $\theta$ the determinant is
nonnegative. Also $Z_n^{(0)}[1]>0$ at every real
$\mu_{\mathrm{lat}}$: the lower bound in
Proposition~\ref{prop:reg-partition} was obtained
from the single term $z=0$, which belongs to this
sector. Dividing \eqref{eq:external-zero-sector}
by its value at $\Phi=1$ identifies the zero-flux
curvature law with the law of $dA$ under the
harmonic-averaged transverse measure.

\subsection{The Gaussian normalization and matter density}
\label{subsec:external-Gaussian-law}

Let $F_n=dC^1(n)=dV_{\mathrm t}$ and let $\gamma_n$
be the standard Gaussian probability measure on this
Euclidean space. The map $d_{\mathrm t}=d|_{V_{\mathrm t}}$
is a linear bijection onto $F_n$, by
\eqref{eq:external-Hodge}; write $A(u)=d_{\mathrm t}^{-1}u$.
Changing variables along this map gives
\begin{equation}\label{eq:external-Gaussian-change}
 \int_{V_{\mathrm t}}G(dA)e^{-\|dA\|^2/2}\,\dd\sigma(A)
                     =\mathcal G_n\int G(u)\,\dd\gamma_n(u).
\end{equation}
Indeed the Jacobian is $|\det d_{\mathrm t}|^{-1}$,
and $\dim F_n=2V-2$. The Gaussian density contributes
$(2\pi)^{V-1}$; setting $G=1$ identifies the constant
with \eqref{eq:external-transverse-Gaussian}.
The formula holds for nonnegative Borel functions,
and for the exponentially bounded complex functions
used above by their Gaussian integrability.

Define the photon factor and normalized matter factor by
\begin{align}
 Z_n^{\mathrm{ph}}&=\mathcal N_n\mathcal G_n
       =r_q^{V+2}(2\pi)^{V-1}
                            e^{-\mathcal L(n)/2},
                                  \label{eq:external-photon-factor}\\
 z_{r,n}(q,\mu_{\mathrm{lat}})
   &=\int\frac{\dd\theta}{(2\pi)^3}
       \int\det_n(D_{A(u)+H_\theta}+\mu_{\mathrm{lat}})
                                             \,\dd\gamma_n(u).
                                  \label{eq:external-matter-factor}
\end{align}
Then \eqref{eq:external-zero-sector} and
\eqref{eq:external-Gaussian-change} imply
\begin{equation}\label{eq:external-zero-factorization}
 Z_n^{(0)}[1]=Z_n^{\mathrm{ph}}z_{r,n},\qquad
                         0<z_{r,n}\le K_n(0).
\end{equation}
Positivity follows from that of $Z_n^{(0)}[1]$;
the upper bound follows because both integration
measures in \eqref{eq:external-matter-factor} are
probability measures.

On a cubic torus $V=n^3$, fix a positive reference mass
$\bar m$, mesh $\eta$ and a real pair
$(m_{\mathrm{in}},\varepsilon_{\mathrm{in}}^{\mathrm v})$,
with $\mu_{\mathrm{lat}}=\bar m\eta+m_{\mathrm{in}}$.
The density with external cochain $a$ is
\begin{equation}\label{eq:external-matter-density}
 \mathrm{dens}^{\theta,a}(u)=
 e^{-\varepsilon_{\mathrm{in}}^{\mathrm v}n^3}
 \frac{\det_n(D_{A(u)+H_\theta+a}+\mu_{\mathrm{lat}})}
                       {\det_n(D^{\mathrm{ch}}[0]+\bar m\eta)}.
\end{equation}
The denominator is positive, since the real part
of $D^{\mathrm{ch}}[0]+\bar m\eta$ is at least $\bar m\eta>0$.
At $a=0$ this is a nonnegative bounded continuous
density. Let $\nu^{\theta,0}$ be its measure against
$\gamma_n$ and let $\bar\nu$ be the twist average.
Equations \eqref{eq:external-zero-sector} and
\eqref{eq:external-Gaussian-change} give
\begin{align}
 \bar\nu(\Phi)
   &=\frac{e^{-\varepsilon_{\mathrm{in}}^{\mathrm v}n^3}}
        {\det_n(D^{\mathrm{ch}}[0]+\bar m\eta)Z_n^{\mathrm{ph}}}
                                      Z_n^{(0)}[\Phi],
                                  \label{eq:external-density-zero-sector}\\
 \bar\nu(1)
   &=\frac{e^{-\varepsilon_{\mathrm{in}}^{\mathrm v}n^3}
                                   z_{r,n}}
                       {\det_n(D^{\mathrm{ch}}[0]+\bar m\eta)}
                                  \in(0,\infty).
                                  \label{eq:external-density-normalization}
\end{align}
Thus $\bar\nu/\bar\nu(1)$ is exactly the zero-flux
curvature law. The vacuum normalization cancels in
this quotient. The quantities $Z_n^{\mathrm{ph}}$
and $z_{r,n}$ depend only on $(n,q,\mu_{\mathrm{lat}})$
and the fixed chamber parameters; the auxiliary
split $\mu_{\mathrm{lat}}=\bar m\eta+m_{\mathrm{in}}$
does not change them.

\subsection{Holomorphic external partition functions}
\label{subsec:external-holomorphy}

Write $\mathfrak F(n)$ for the exponentially bounded
Borel functions $\Phi$ used above. For $w\in\Z^3$ let
$\mathcal M_w$ be the closed integer two-cochains of
flux $w$. For a complex one-cochain $a$ define
\begin{align}
 Z_n^{\mathrm{x},(w)}(a;\Phi)
  &=\sum_{z\in\mathcal M_w}\int_{[0,r_q)^{\mathcal L_n}}
       \Phi(F)e^{-\|F\|^2/2}
                \det_n(D_{A+a}+\mu_{\mathrm{lat}})\,\dd A,
       \qquad F=dA+r_qz,\label{eq:external-sector-source}\\
 \mathbf Z_n^{\mathrm{x}}(a;\Phi)
  &=\int\frac{\dd\theta}{(2\pi)^3}
       \int_{V_{\mathrm t}}\Phi(dA)e^{-\|dA\|^2/2}
       \det_n(D_{A+H_\theta+a}+\mu_{\mathrm{lat}})\,\dd\sigma(A).
                                  \label{eq:external-averaged-source}
\end{align}
When $\Phi=1$ we omit it from the notation. These are
partition functions in an external connection; the
curvature $F$ in their photon weight is not shifted.

\begin{proposition}\label{prop:external-entire}
For every $\Phi\in\mathfrak F(n)$ both expressions
converge absolutely, locally uniformly in $a$, and are
entire in its $3V$ complex link values. The sum over
all sectors converges absolutely, locally uniformly
as well. Moreover
\begin{equation}\label{eq:external-source-zero-sector}
 Z_n^{\mathrm{x},(0)}(a;\Phi)
             =\mathcal N_n\mathbf Z_n^{\mathrm{x}}(a;\Phi),
 \qquad
 Z_n^{\mathrm{x},(w)}(0;\Phi)=Z_n^{(w)}[\Phi].
\end{equation}
In particular $\mathbf Z_n^{\mathrm{x}}(0)>0$.
\end{proposition}

\begin{proof}
For real $A,\theta$ the imaginary parts of $A+a$ and
$A+H_\theta+a$ are both $\operatorname{Im}a$. On a
polydisc $\max_b|a(b)|\le s$ the determinant is
bounded by $K_n(s)$. If
$|\Phi(u)|\le C_\Phi e^{C_\Phi\|u\|}$, the absolute
value of each integrand is consequently bounded by
\[
 K_n(s)C_\Phi e^{C_\Phi^2}e^{-\|F\|^2/4}
 \quad\text{or}\quad
 K_n(s)C_\Phi e^{C_\Phi^2}e^{-\|dA\|^2/4}.
\]
The first majorant can be summed over every integer
two-cochain and integrated over the compact box,
giving
$K_n(s)C_\Phi e^{C_\Phi^2}(r_qC_q(1/4))^{3V}$.
The second is integrable by
\eqref{eq:external-transverse-Gaussian}: rescaling
$A$ by $\sqrt2$ gives
\[
 \int_{V_{\mathrm t}}e^{-\|dA\|^2/4}\,\dd\sigma(A)
                          =2^{V-1}\mathcal G_n.
\]
These bounds also show local uniform absolute
convergence of the sector sum.

We recall the parameter-integral argument to justify
entireness and all later differentiations. Let
$(Y,\nu)$ be a $\sigma$-finite measure space. Suppose
$v(y,a)$ is measurable in $(y,a)$, entire in
$a\in\C^d$ for each $y$, and bounded on each
polydisc by a $\nu$-integrable $M_s(y)$. If $v_\alpha(y)$
are its Taylor coefficients at zero, integration
over the product of circles of radius $s$ gives
\[
 v_\alpha(y)=s^{-|\alpha|}
       \int_{[0,2\pi]^d}
       v(y,se^{i\theta})e^{-i\alpha\cdot\theta}
                         \frac{\dd\theta}{(2\pi)^d},
 \qquad |v_\alpha(y)|\le s^{-|\alpha|}M_s(y).
\]
The integral formula makes the coefficients measurable.
For $\max_j|a_j|\le s$, applying the bound at radius
$2s$ yields
\[
 \sum_\alpha\int|v_\alpha(y)a^\alpha|\,\dd\nu(y)
       \le\sum_\alpha2^{-|\alpha|}\int M_{2s}(y)\,\dd\nu(y)
       =2^d\int M_{2s}(y)\,\dd\nu(y)<\infty.
\]
Thus the Taylor series may be integrated term by term,
and its integral is an entire function with coefficients
$\int v_\alpha$. This also proves differentiation under
the integral at zero. Translating the parameter to an
arbitrary $a^0$ proves it there, using the majorant
$M_{s+\|a^0\|_\infty}$. This argument applies to a
counting measure times Lebesgue measure as well as to
the transverse integral.

Apply this argument to the two displayed majorants.
The dependence on $a$ in each integrand is a finite
sum of exponentials of linear forms, by the determinant
polynomial. The other factors are Borel and independent
of $a$, so the hypotheses hold also when $\Phi$ is
merely Borel. This proves the asserted holomorphy.
For fixed complex $a$, the function
$W_a(A)=\det_n(D_{A+a}+\mu_{\mathrm{lat}})$ on real
cochains is bounded by $K_n(\|\operatorname{Im}a\|_\infty)$,
continuous, and has the two invariances required in
\eqref{eq:external-zero-unfolding}. That formula gives
\eqref{eq:external-source-zero-sector}; the value at zero
follows from the definitions. Its positive zero value
follows from $Z_n^{(0)}[1]>0$ and $\mathcal N_n>0$.
\end{proof}

The complex densities in
\eqref{eq:external-matter-density} define finite complex
measures $\nu^{\theta,a}$ and their twist average
$\bar\nu^a$ against $\gamma_n$. Their absolute values
are bounded by
\begin{equation}\label{eq:external-complex-density-bound}
 |\mathrm{dens}^{\theta,a}(u)|
 \le
 \frac{e^{-\varepsilon_{\mathrm{in}}^{\mathrm v}n^3}
                   K_n(\|\operatorname{Im}a\|_\infty)}
                   {\det_n(D^{\mathrm{ch}}[0]+\bar m\eta)}.
\end{equation}
This bound is independent of $(\theta,u)$.
Fubini therefore writes the averaged measure as a
bounded Borel density against $\gamma_n$. Applying
\eqref{eq:external-Gaussian-change} at each twist,
and then \eqref{eq:external-source-zero-sector}, gives
\begin{equation}\label{eq:external-complex-measure}
 \bar\nu^a(\Phi)=
 \frac{e^{-\varepsilon_{\mathrm{in}}^{\mathrm v}n^3}}
      {\det_n(D^{\mathrm{ch}}[0]+\bar m\eta)\mathcal G_n}
                     \mathbf Z_n^{\mathrm{x}}(a;\Phi)
 =
 \frac{e^{-\varepsilon_{\mathrm{in}}^{\mathrm v}n^3}}
      {\det_n(D^{\mathrm{ch}}[0]+\bar m\eta)Z_n^{\mathrm{ph}}}
                     Z_n^{\mathrm{x},(0)}(a;\Phi).
\end{equation}
The exponential majorants in the proof above justify
both changes of integration for $\Phi\in\mathfrak F(n)$.
In particular, for every Borel set $E\subset F_n$,
$a\mapsto\bar\nu^a(E)$ is entire and
\[
 |\bar\nu^a(E)|\le
 \frac{e^{-\varepsilon_{\mathrm{in}}^{\mathrm v}n^3}
                   K_n(\|\operatorname{Im}a\|_\infty)}
                   {\det_n(D^{\mathrm{ch}}[0]+\bar m\eta)}.
\]
Indeed $F_n$ is a closed linear subspace, so its
Borel set indicator extends to a bounded Borel
function on $C^2(n)$. Composition with any linear
parameter map $t\mapsto a(t)$ remains entire.
These statements concern the unnormalized complex
measures; the positive normalization at $a=0$ is
\eqref{eq:external-density-normalization}.

\subsection{The source shift and Wick derivatives}
\label{subsec:external-source-shift}

\begin{proposition}\label{prop:external-shift}
For every flux $w$ and complex one-cochain $a$,
\begin{equation}\label{eq:external-shift-sector}
 Z_n^{(w)}[e^{\langle da,F\rangle}]
      =e^{\langle da,da\rangle/2}Z_n^{\mathrm{x},(w)}(a).
\end{equation}
Both sides, and their sums over $w$, converge
absolutely. Consequently
\begin{equation}\label{eq:external-shift-all-flux}
 Z_n[e^{\langle da,F\rangle}]
      =e^{\langle da,da\rangle/2}
                    \sum_w Z_n^{\mathrm{x},(w)}(a).
\end{equation}
The quadratic form in the prefactor is bilinear;
there is no complex conjugation of $da$.
\end{proposition}

\begin{proof}
First let $a$ be real. On real cochains define
\[
 P_a(A)=\sum_{z\in\mathcal M_w}
  e^{-\|dA+r_qz\|^2/2+\langle da,dA+r_qz\rangle}
                            \det_n(D_A+\mu_{\mathrm{lat}}).
\]
This is a nonnegative Borel function. Adding
$r_q\mathbf1_b$ to $A$ adds $d\mathbf1_b$ to the
integer two-cochain in its curvature. This map
permutes $\mathcal M_w$: its image is closed
and has the same flux, since each exact plaquette
cochain has zero plane sums. The determinant
is unchanged. Hence $P_a$ is invariant under
$r_qC^1(n;\Z)$. Both the compact box and its
translate by $a$ are fundamental domains for
this group. The domain argument in
\S\ref{subsec:external-zero-flux} gives
\[
 \int_{[0,r_q)^{\mathcal L_n}}P_a(A)\,\dd A
        =\int_{[0,r_q)^{\mathcal L_n}}P_a(A+a)\,\dd A.
\]
On the right the exponent completes to
\[
 -\tfrac12\|F+da\|^2+\langle da,F+da\rangle
                  =-\tfrac12\|F\|^2+\tfrac12\|da\|^2.
\]
Tonelli now identifies the two sides with
\eqref{eq:external-shift-sector} for real $a$.
Their finiteness also follows from the Gaussian
majorant used below.

For complex $a$, the right side of
\eqref{eq:external-shift-sector} is entire by
Proposition~\ref{prop:external-entire}.
The left side is entire by the same
parameter-integral argument. In fact, on
$\max_b|a(b)|\le s$,
\[
 \|d\operatorname{Re}a\|\le4s\sqrt{3V}=:c_s,
 \qquad
 \left|e^{-\|F\|^2/2+\langle da,F\rangle}
                            \det_n(D_A+\mu_{\mathrm{lat}})\right|
                  \le K_n(0)e^{c_s^2}e^{-\|F\|^2/4}.
\]
This can be summed and integrated over all integer
two-cochains, so it supplies an integrable polydisc
majorant. Both sides agree on the real cochain
space. An entire function vanishing there has
all complex partial derivatives zero there:
at a real point these derivatives can be taken
along real increments. Repeating this for all
derivatives shows that every Taylor coefficient
at zero vanishes. Thus the two entire functions
agree for all complex $a$.
Finally the same Gaussian majorant permits
summation over sectors; the sector sum on the
right converges absolutely by
Proposition~\ref{prop:external-entire}.
\end{proof}

Let $a^{(1)},\ldots,a^{(r)}$ be real cochains and
$X_j=\langle da^{(j)},F\rangle$. Dividing
\eqref{eq:external-shift-sector} at $w=0$ by its
positive zero-source value yields the normalized
generating identity
\begin{equation}\label{eq:external-generating-general}
 \E_n^{\mathrm{av}}\exp\left(
       \sum_jz_jX_j-\tfrac12\sum_{i,j}
                  z_iz_j\langle da^{(i)},da^{(j)}\rangle\right)
 =
 \frac{\mathbf Z_n^{\mathrm{x}}(\sum_jz_ja^{(j)})}
                         {\mathbf Z_n^{\mathrm{x}}(0)},
                                                   \qquad z\in\C^r.
\end{equation}
The expectation $\E_n^{\mathrm{av}}$ is the zero-flux
law identified in \eqref{eq:external-zero-sector}.
The Gaussian factor on the left is determined by
the counting pairing of the curvatures of the
source cochains.

Use the reference cells of Section~\ref{sec:reference-cells}
on the torus $n_*$. Their index set is
$\mathcal C_*=\{0,\ldots,2L-1\}^3$, and
$\widehat a_c$ is the periodization of
$\hat{\mathfrak a}_c^o$. These are real periodic cochains.
Equation \eqref{eq:cells-orthogonality} gives
\begin{equation}\label{eq:external-cell-pairing}
 \langle d\widehat a_c,d\widehat a_{c'}\rangle
                         =v_\eta\,1_{\{c=c'\}}.
\end{equation}
The variance $v_\eta=\|d\hat{\mathfrak a}_\eta^o\|^2$
is independent of the cell and the placement.
This is the cell-geometry input needed here:
the external-response proof uses neither the
cell width nor its offsets otherwise. Put
\[
 a(t)=e^{-1}\sum_{c\in\mathcal C_*}t_c\widehat a_c,\qquad
 \bar Z^{\mathrm{ext}}(t)=
     \int\frac{\dd\theta}{(2\pi)^3}
        \int\det_{n_*}(D_{A(u)+H_\theta+a(t)}
                              +\mu_{\mathrm{lat}})\,\dd\gamma_{n_*}(u).
\]
The reference torus has $n_*\ge4$, and $e>0$ is
the physical charge in this convention.
Equation \eqref{eq:external-Gaussian-change} gives
\begin{equation}\label{eq:external-reference-partition}
 \bar Z^{\mathrm{ext}}(t)=
          \mathcal G_{n_*}^{-1}\mathbf Z_{n_*}^{\mathrm{x}}(a(t)),
 \qquad
 \mathring\Phi(t):=\frac{\bar Z^{\mathrm{ext}}(t)}
                           {\bar Z^{\mathrm{ext}}(0)}
        =\frac{\mathbf Z_{n_*}^{\mathrm{x}}(a(t))}
                           {\mathbf Z_{n_*}^{\mathrm{x}}(0)}.
\end{equation}
Its zero value in the denominator is positive,
and $\mathring\Phi$ is entire. With
$X_c=\langle d\widehat a_c,F\rangle$,
\eqref{eq:external-generating-general} becomes
\begin{equation}\label{eq:external-generating-cells}
 \E_{n_*}^{\mathrm{av}}
      \exp\sum_c(z_cX_c-\tfrac12v_\eta z_c^2)
                                      =\mathring\Phi(ez).
\end{equation}

Define the Hermite polynomials with this variance by
$H_j(y)=\partial_z^j e^{zy-v_\eta z^2/2}|_{z=0}$.
In particular
\[
 H_2(y)=y^2-v_\eta,\qquad
 H_4(y)=y^4-6v_\eta y^2+3v_\eta^2.
\]
To differentiate \eqref{eq:external-generating-cells}
under the expectation, on a polydisc
$\max_c|z_c|\le s$ use
\[
 \left|\exp\sum_c(z_cX_c-\tfrac12v_\eta z_c^2)\right|
 \le
 \exp\left(s|\mathcal C_*|\sqrt{v_\eta}\|F\|
                       +\tfrac12|\mathcal C_*|s^2v_\eta\right).
\]
This is exponentially bounded in the real curvature
and is integrable under the zero-flux law, by the
Gaussian majorant. The parameter-integral argument
therefore applies to this probability measure.
The resulting products of $H_2$ or $H_4$ are also
integrable: $|X_c|\le\sqrt{v_\eta}\|F\|$ bounds them by a
finite polynomial in $\|F\|$, which is bounded by
a constant times $e^{\|F\|}$.

For $m=2$ or $4$ and a subset $S\subset\mathcal C_*$,
differentiate in precisely those variables. The
factor $ez$ on the right supplies $e^{m|S|}$, giving
\[
 \E_{n_*}^{\mathrm{av}}\prod_{c\in S}e^{-m}H_m(X_c)
         =\left(\prod_{c\in S}\partial_{t_c}^m
                                      \mathring\Phi\right)(0).
\]
Expanding the product over subsets proves, for
every real $c_m$,
\begin{equation}\label{eq:external-Wick-product}
 \E_{n_*}^{\mathrm{av}}
        \prod_{c\in\mathcal C_*}(e^{-m}H_m(X_c)-c_m)
     =\left(\prod_{c\in\mathcal C_*}
                (\partial_{t_c}^m-c_m)\mathring\Phi\right)(0).
\end{equation}
The derivative operators commute because
$\mathring\Phi$ is entire. This identity is the
finite-volume quantity to which the terminal
interacting estimates will be applied; all its
normalization factors have already been included
in \eqref{eq:external-reference-partition}.


\section{Terminal operators and Gaussian bounds}
\label{sec:terminal-foundations}

The terminal estimates use the fine and coarse field maps, their
weighted Grassmann norms, and Gaussian domination in the zero-flux
sector. We establish these inputs before the partition comparison
and the response estimates.

\subsection{Reference scales and the probe}
\label{subsec:interacting-stop-probe}

We use the reference torus of the tuned flow.
Fix $\bar m\in(0,1]$, put
$h=L^{\lfloor\log_L(1/\bar m)\rfloor}$, and assume
$0<e\sqrt h\le e_1$ and $N\ge N_{\rm ref}(e)$.
By \eqref{eq:rgfl-physical-scales} and
\eqref{eq:rgfl-reference-threshold},
\begin{equation}\label{eq:interacting-physical-scales}
 \begin{aligned}
 \eta=L^{-N},\qquad K=N+\log_Lh,\qquad
 &a_j=L^{j-N},\\
 e_j=e\sqrt{a_j},\qquad
 &\eta_j=L^{-j}=\eta/a_j,\qquad \ell_j=R^*/a_j .
 \end{aligned}
\end{equation}
Thus a point with level-$j$ coordinates $c$
has physical coordinates $a_jc$, and
$T_j=\eta_j\Z^3/\ell_j\Z^3$.
At $j=0$ the charge is $q=e\sqrt\eta$.
The symbols $e_0,e_1$ without a running
index retain their meanings as charge ceilings.
For a bond cochain $a$ on the fine unit torus,
its rescaling is
\begin{equation}\label{eq:interacting-connection-rescaling}
                    a^{[j]}(b)=L^{j/2}a(L^jb).
\end{equation}
The physical scale $a_j$ is a number, while
$a$ denotes a connection.
On the reference torus,
\eqref{eq:rgfl-reference-stop} and
\eqref{eq:rgfl-scale-ranges} give
\[
 k_*\in\{K-1,K-2\},\qquad
 a_*\in\{h/L,h/L^2\},\qquad
 \ell_*=Mr_*,\qquad r_*>31L+16.
\]
These stop-scale quantities are independent
of $N\ge N_{\rm ref}(e)$.
Write $r_{\rm loc}$ for the localization
power in Section~\ref{sec:rg-flow}.

Let $\beta\in C_c^\infty(\R^3;\Lambda^2)$ be real
and nonzero. Use the cube
$Q_\beta=z_\beta+[-D_\beta^{\rm buf}/2,D_\beta^{\rm buf}/2]^3$
of Section~\ref{sec:reference-cells},
where $D_\beta^{\rm buf}=D_\beta+4$.
In this section a placement $v\in\eta\Z^3$
is admitted if, in physical coordinates,
\begin{equation}\label{eq:interacting-admitted-placement}
              v+Q_\beta\subset[\ell_c/8,7\ell_c/8]^3,
              \qquad \ell_c=R^*/(2L).
\end{equation}
A centered placement has
$|v+z_\beta-(\ell_c/2)(1,1,1)|_\infty\le\eta/2$.
It is admitted when
$D_\beta^{\rm buf}/2+\eta/2\le3\ell_c/8$.
This follows by adding the displacement
of the cube center to its half-side.

Form $\widehat{\mathfrak a}$ with $\tau_v\beta$
as in \eqref{eq:cells-midpoint-source}, and
periodize it on $C(n_*)$, $n_*=R^*/\eta$.
For $c\in\mathcal C_*=\{0,\ldots,2L-1\}^3$,
write $\widehat{\mathfrak a}_c$ for its
cell copy. In physical coordinates the
wall map carrying the first cell to cell $c$
acts in coordinate $\mu$ by translation through
$2i\ell_c$ if $c_\mu=2i$, and by reflection
$y_\mu\mapsto(2i+2)\ell_c-y_\mu$ if $c_\mu=2i+1$.
Define the marked connections by
\begin{equation}\label{eq:interacting-marked-connections}
 a(t)=e^{-1}\sum_ct_c\widehat{\mathfrak a}_c,\qquad
 \widehat c_c=e^{-1}\widehat{\mathfrak a}_c^{[k_*]},\qquad
 \mathscr C(t)=a(t)^{[k_*]}=\sum_ct_c\widehat c_c .
\end{equation}
For a level-$j$ bond function $A$, $\partial A$
uses forward differences divided by $\eta_j$,
and
\[
 |\delta_\alpha\partial A|_\infty
  =\sup_{\mu,\nu}\sup_{0<d(x,x')\le1}
     \frac{|\partial_\nu A_\mu(x')-\partial_\nu A_\mu(x)|}
          {d(x,x')^\alpha}.
\]
The external norm of a fine connection is
$|a|_{\rm ext}=|a^{[k_*]}|_\infty+
|\partial a^{[k_*]}|_\infty+
|\delta_\alpha\partial a^{[k_*]}|_\infty$.

We define the uniform probe constants over both
possible stop scales. For $\eta'=L^{-N'}$ put
\[
 g_\nu^{\eta'}[\beta']
      =-\sum_{\mu\ne\nu}\delta_{\eta',\mu}J_{\eta'}\beta'_{\mu\nu}.
\]
This is the smooth vector field
$J_{\eta'}a'_{\eta'}$ in the midpoint formula.
Let $\mathcal I_\beta$ consist of triples
$(N',a_{\rm s},r')$ with
$N'\ge0$, $a_{\rm s}\in\{h/L,h/L^2\}$,
$N'+\log_La_{\rm s}\ge1$, and
$r'\in\{\pm1\}^3$.
$r'$ is identified with its diagonal orthogonal
matrix, and $U_{r'}\beta(x)=(\wedge^2r')\beta(r'^{-1}x)$.
At such a triple set $\eta''=\eta'/a_{\rm s}$,
$g^0=g^{\eta'}[U_{r'}\beta]$, and
\[
 G_\nu(c)=\sqrt{a_{\rm s}}\,
                  g_\nu^0(a_{\rm s}c+\eta'e_\nu/2)
                       \quad(c\in\eta''\Z^3),\qquad
 W_\nu(x)=\sqrt{\eta'}\,g_\nu^0(\eta'x+\eta'e_\nu/2)
                       \quad(x\in\Z^3).
\]
The three norms of $G$ are taken at mesh $\eta''$.
Set
\begin{equation}\label{eq:interacting-uniform-probe-constants}
 C_\beta=\sup_{\mathcal I_\beta}
              (|G|_\infty+|\partial G|_\infty+
                                    |\delta_\alpha\partial G|_\infty),
 \qquad
 \mathbf D=\sup_{\mathcal I_\beta}\|dW\|_{\ell^2(\text{plaquettes})}^2.
\end{equation}
These suprema depend on $\bar m,L,\alpha,\beta$,
and contain no charge, torus, or placement.
They have the finite bounds
\begin{equation}\label{eq:interacting-uniform-probe-finite}
 C_\beta\le2\sqrt h(1+h+3h^2)|\beta|_{C^3},
 \qquad
 \mathbf D\le48|\beta|_{C^2}^2(D_\beta^{\rm buf}+3)^3 .
\end{equation}

Indeed the centered difference is the average
of the corresponding derivative along its
segment, and $J_{\eta'}$ has norm one on
each derivative sup norm. Each component
of $g^0$ has two summands, hence
$|D^ig_\nu^0|_\infty\le2|\beta|_{C^{i+1}}$ for
$i=0,1,2$. This gives
$|G|_\infty\le2a_{\rm s}^{1/2}|\beta|_{C^1}$.
A first lattice difference of $G$ is
$a_{\rm s}^{3/2}$ times an average of
$\partial_\mu g_\nu^0$, so
$|\partial G|_\infty\le2a_{\rm s}^{3/2}|\beta|_{C^2}$.
The difference of these averages at $c,c'$
is bounded by
$6a_{\rm s}|c-c'|\,|\beta|_{C^3}$.
Since $|c-c'|\le1$ and $\alpha<1$,
the third norm is at most
$6a_{\rm s}^{5/2}|\beta|_{C^3}$.
Sum these bounds and use $a_{\rm s}\le h$.

For the energy, each curl component of $W$
is a difference of two first differences,
and has modulus at most
$4(\eta')^{3/2}|\beta|_{C^2}$.
A supporting plaquette has a supporting
bond, whose endpoints in physical coordinates
lie in $r'Q_\beta$.
Its basepoint is therefore in that cube
enlarged by $\eta'$ in each negative
coordinate direction. There are at most
$(D_\beta^{\rm buf}/\eta'+2)^3$ basepoints
and three plaquette directions. Squaring
the pointwise curl bound gives
$48|\beta|_{C^2}^2(D_\beta^{\rm buf}+2\eta')^3$,
which implies \eqref{eq:interacting-uniform-probe-finite}.

At the reference stop, $(N,a_*,1)\in\mathcal I_\beta$.
The midpoint formula and
\eqref{eq:interacting-connection-rescaling} give
\[
 \widehat{\mathfrak a}^{[k_*]}_\nu(c)
       =\sum_{z\in\Z^3}G_\nu(c-v/a_*+\ell_*z).
\]
Supporting bond endpoints lie in
$a_*^{-1}(v+Q_\beta)$.
Distinct periodic copies of that cube have
distance at least
$\ell_*-3\ell_*/(8L)>\ell_*/2>1+2\eta_*$.
Each functional defining one of the three
norms reads only bonds with endpoint distances
at most $1+2\eta_*$. It therefore reads
at most one periodic copy. Taking its
supremum identifies each torus norm with
the corresponding norm of $G$.
Translations and coordinate reflections
preserve these norms, including the one-bond
shift that reverses an oriented component.
The same argument applies to the cell copies.

The image of $v+Q_\beta$ under each wall map
lies in
$\prod_\mu[c_\mu\ell_c+\ell_c/8,(c_\mu+1)\ell_c-\ell_c/8]$.
For distinct cells, a separating coordinate
on the periodic circle gives distance
at least $\ell_c/4$: any joining arc
contains the two margins of size $\ell_c/8$.
Thus supporting endpoints of distinct
$\widehat c_c$ have stop-coordinate distance
at least $\ell_*/(8L)$, which is larger
than $1+2\eta_*$. Again each external-norm
functional reads at most one copy, so
\begin{equation}\label{eq:interacting-external-probe-bound}
 |a(t)|_{\rm ext}\le C_\beta|t|_\infty/e,\qquad
                 v_\eta=\|d\widehat{\mathfrak a}\|^2\le\mathbf D.
\end{equation}
For the second bound, the covering curvature
$dW$ is supported in the cube of side
$D_\beta^{\rm buf}/\eta+2<n_*$
obtained by enlarging $\eta^{-1}(v+Q_\beta)$.
Its periodic translates do not overlap,
so periodization preserves its counting
norm. Distinct cell curvatures also have
disjoint plaquette supports: a common
plaquette would have vertices in both
endpoint sets at physical distance at
most $\eta\le1$, whereas $\ell_c/4>1$.
This gives \eqref{eq:cells-orthogonality}
for every admitted placement.

Let $P_c$ be the unit-block centers whose
blocks contain an endpoint of a supporting
bond of $\widehat c_c$, and let $P_c^+$
be their $(M+1)$-neighborhood.
For $y\in P_c$, $y'\in P_{c'}$ choose the
corresponding endpoints in $\Delta_y,\Delta_{y'}$.
Each endpoint is at distance less than $1/2$
from its center. Consequently
\begin{align}
 d_{\rm sep}:=\min_{c\ne c'}d(P_c,P_{c'})
    &\ge\ell_*/(8L)-1=Mr_*/(8L)-1,
                              \label{eq:interacting-probe-separation}\\
 |P_c^+|&\le p_\beta:=
                  (D_\beta^{\rm buf}/a_*+2M+4)^3 .
                              \label{eq:interacting-support-number}
\end{align}
For the count, the endpoints have a lift in
a cube of side $D_\beta^{\rm buf}/a_*$.
A center adds less than $1/2$ to each
coordinate bound, and its neighborhood adds
$M+1$. The resulting open interval has
length $D_\beta^{\rm buf}/a_*+2M+3$
and contains fewer than one plus that
many integers. Multiply the three counts.
If the lattice probe vanishes, all these
sets are empty and $d_{\rm sep}=\infty$.

\subsection{Terminal operators and local insertions}
\label{subsec:interacting-terminal-operators}

At the stop level write $\eta_*=L^{-k_*}$,
$T_*=\eta_*\Z^3/\ell_*\Z^3$ and
$T_*^0=\Z^3/\ell_*\Z^3$. The fine chamber spinors
have counting measure $\eta_*^3\,48^{-1}\sum_{c,g}$,
and the coarse spinors have four components per unit
site and weight one. We use
\[
 \begin{gathered}
 k_*\ge1,\qquad \ell_*>2M+2,\qquad \gamma^FM\ge16,\\
 1-L^{-1}=:b_-\le b_*\le1,\qquad
 L^{-3}=:\mu_-<\mu_*\le L^{-1}.
 \end{gathered}
\]
The running mass $\mu_*$ is positive; the bare mass
coupling is treated separately. Let
\begin{align}
 F_*(A)&=D_A+\mu_*+b_*JP_*(A)J^*,&
 S_*(A)&=F_*(A)^{-1},\notag\\
 H_*(A)&=b_*S_*(A)JQ_*^T(-A),&
 D_*(A)&=b_*-b_*Q_*(A)J^*H_*(A).
                                  \label{eq:interacting-terminal-operators}
\end{align}
These are the blocked operators already constructed
in \eqref{eq:block-algebra} and
\eqref{eq:coarse-dirac-definition}, at $k=k_*$.
The symbol $F_*$ here is an operator on spinors,
distinct from the curvature variable in the external
partition function.

For a unit site $y$ let $\Delta_y$ be its fine block.
The averaging phases have the form
$\varphi_A(y,c)=e^{i\phi_A(y,c)}$, where
$\phi_A=e_*\eta_*\tau_{k_*}A$ is linear in $A$.
The nested block centers give
\begin{equation}\label{eq:interacting-phase-bounds}
 |\phi_B(y,c)|\le\tfrac32e_*|B|_\infty,\qquad
 |\upsilon_B(c,c')|\le3|c-c'|e_*|B|_\infty.
\end{equation}
The first phase uses only bonds inside $\Delta_y$;
the second, which defines $U_A(c,c')=e^{i\upsilon_A(c,c')}$,
uses only bonds in the coordinate box between its endpoints.
For the first bound, the path lengths at block scales
$L^j\eta_*$ add up to
$\frac32(L-1)\eta_*\sum_{j<k_*}L^j\le\frac32$.
Each of the six coordinate-order paths between
$c$ and $c'$ has length at most $3|c-c'|$,
which gives the second bound.

For a complex direction $B$, let $P_B$ be the unit sites
$y$ whose block $\Delta_y$ contains an endpoint of a bond
in $\supp B$. Define
\[
 \epsilon_B=e_*|B|_\infty,\qquad
 P_B^+=\{y:d(y,P_B)\le M+1\}.
\]
The distance is the periodic sup distance on the unit
lattice. For an operator $O$ between fine, coarse or
pair spaces, localize each fine spinor in its unit
block and each pair at its first point, and put
\begin{align}
 |O|(y,y')&=\sup_{\|v\|_\infty\le1}
                         \|\mathbf1_yO\mathbf1_{y'}v\|_\infty,\notag\\
 \|O\|_\gamma&=\max\left\{
       \sup_y\sum_{y'}|O|(y,y')e^{\gamma d(y,y')},
       \sup_{y'}\sum_y|O|(y,y')e^{\gamma d(y,y')}\right\}.
                                  \label{eq:interacting-operator-majorants}
\end{align}
The spinor sup norm uses the maximum of component
moduli. Inserting the unit-block partition between
two factors gives
$|O_1O_2|\le|O_1||O_2|$ and
$\|O_1O_2\|_\gamma\le\|O_1\|_\gamma\|O_2\|_\gamma$.
For coarse spinor operators,
\begin{equation}\label{eq:interacting-matrix-norms}
 \|O\|_{\ell^2\to\ell^2}\le2\|O\|_0,\qquad
 \|O\|_1\le4\sum_{y,y'}|O|(y,y').
\end{equation}
Indeed the matrix row sums are bounded by $\|O\|_0$
and the column sums by $4\|O\|_0$, so Schur's test
gives the first bound. Polar decomposition bounds
the trace norm by the sum of all entry moduli, and
the entries in one block have total modulus at most
four times its maximum row sum.

We use the fixed constants $C^F,\gamma^F$
of the fermion construction, with its calibration
$a'\ge1$ and
\[
 C^F\ge608\exp(1)\max(1,b)a',
 \qquad C^F\ge64\exp(1)b^2a',\qquad b=1.
\]
The first inequality also dominates the component
norm bounds $64a'$ and $608a'$ for the fine
inverse and its short covariant quotients.
The charge ceiling satisfies
$6Le_*^\varepsilon\le1$ and
$3Me_*^\varepsilon\le1$; on $\widetilde G_*$,
$\mathfrak i(A)<e_*^\varepsilon$. The phase
estimates therefore give
$|\varphi_A^{\pm1}|,|U_A^{\pm1}|\le\exp(1)$,
including every allowed long pair.
The unweakened estimate
\eqref{eq:global-propagator-bounds}, at its
spectral parameter $y=\infty$, and the block
phases give the following terminal bounds.
On the outer field domain $\widetilde G_*$ the inverse
$S_*$ is analytic, and, with
$E(y,y')=e^{-\gamma^Fd(y,y')}$,
\begin{equation}\label{eq:interacting-terminal-input}
 |S_*|,\quad|\delta_{\alpha,A}S_*|,\quad
 |H_*|,\quad|\delta_{\alpha,A}H_*|
                                  \le C^FE.
\end{equation}
The two pair majorants are restricted to $d_X<1$,
including pairs inside a chamber cluster.
The field domain
\[
 \mathcal W=\{A_r+\mathscr C:
          A_r\in\widetilde R_*\text{ real},\
          \mathscr C\in\mathfrak X_*\}
\]
is contained in $\widetilde G_*$, is invariant under
$A\mapsto-A$, and satisfies
$\mathfrak i(A):=e_*|\operatorname{Im}A|_\infty<1/16$
and $3M\mathfrak i(A)<1/4$. For long pairs,
$|\delta_{\alpha,A}v(x,x')|_c
\le\mathrm e|v(x')|_c+|v(x)|_c$.
The constants $C^F\ge1$ and $0<\gamma^F\le1$ are
uniform in the mesh and reference torus. These
are the stop-level forms of the preceding walk
and local inverse estimates. To obtain the bounds
for $H_*$, write
$H_*=b_*S_*JQ_*^T(-A)$. The last factor has
Euclidean spinor sup norm at most $\exp(1)$,
preserves each unit block, and has no chamber
factor: it repeats the same coarse spinor over
the $48$ labels, multiplied by the block phase.
Composition with the localized inverse bounds
gives constants $32\exp(1)ba'$ and
$304\exp(1)ba'$. Since
$|v|_c\le|v|_2\le2|v|_c$ on $\C^4$,
their component sup norms are at most twice
these constants, hence at most $C^F$.

We shall also need the global component row
bound, without an exponential weight:
\begin{equation}\label{eq:interacting-global-field-rows}
 |H_*^\#(A)|_{1,\infty}\le C^F,\qquad
 |\bar H_*^\#(A)|_{1,\infty}\le C^F .
\end{equation}
Here the rows include every allowed pair, and
the columns are the four coarse components
at the unit sites. Indeed, the unweighted
consequences of \eqref{eq:global-propagator-bounds}
give
\[
 \|H_*(A)v\|_{\mathrm{su}}
       \le32\exp(1)ba'\|v\|_{\mathrm{su}},\qquad
 \|\delta_{\alpha,A}H_*(A)v\|_{\mathrm{su}}
       \le304\exp(1)ba'\|v\|_{\mathrm{su}}
\]
for short pairs. These include distinct
labels at the same photon site.
For a long pair, $1\le d_X(x,x')\le M$,
so $d_X^{-\alpha}\le1$ and $|U_A|\le\exp(1)$.
Consequently
\[
 |(\delta_{\alpha,A}H_*(A)v)(x,x')|_2
 \le(1+\exp(1))\|H_*(A)v\|_{\mathrm{su}}
 \le128\exp(1)ba'\|v\|_{\mathrm{su}} .
\]
Thus $304\exp(1)ba'$ bounds the Euclidean
sup norm for all pair rows. Conversion to
the component sup norm costs a factor of
two, and the calibration of $C^F$ gives
\eqref{eq:interacting-global-field-rows}.
To identify this operator norm with the row
sum, fix one output component and prescribe
each input component to have the conjugate
phase of its matrix entry. The coarse
spinor values on a fundamental domain can
be prescribed independently; their
antiperiodic extension does not affect
the entry moduli. All unit sites have
weight one.

The barred bound follows from
$\bar H_*(A)=\mathsf C^{-1}H_*(-A)\mathsf C$.
The charge-conjugation matrix $\mathsf C$
is a signed permutation of the four
components, and commutes with the scalar
phases in the covariant quotient. It
preserves both spinor norms. Applying
the preceding argument at $-A$ proves
the second inequality in
\eqref{eq:interacting-global-field-rows}.
Finally, localization of the long-pair
argument gives the bound
$C^F\exp(-\gamma^F\min\{d([x],y'),d([x'],y')\})$
for the coefficient at a coarse site $y'$.
The intermediate Euclidean constant is
$128\exp(1)ba'$, and its component constant
$256\exp(1)ba'$ is again at most $C^F$.

Put
\[
 \mathcal S=C^F(24/\gamma^F)^3,\qquad
 h_{\mathrm L}=6\mathcal S,\qquad
 C_H=203M(C^F)^2,\qquad C_K=226\mathcal S^2.
\]
The lattice sum
$\sum_{y'}e^{-\gamma d(y,y')}\le(12/\gamma)^3$,
$0<\gamma\le6$, follows by choosing nearest periodic
representatives and summing
$e^{-\gamma|z|_\infty}\le e^{-\gamma|z|_1/3}$ over
$\Z^3$. The one-dimensional sum is
$\coth(\gamma/6)\le12/\gamma$.

For a fine spinor or Hölder-pair index $u$, let
$\operatorname{Leg}_A(u;y',\beta')$ be the coefficient
of $\Psi_{\beta'}(y')$ after substituting
$H_*^\#(A)\Psi$ in that generator. Barred generators
use the corresponding barred coefficient.
The set $\operatorname{pos}(u)$ consists of the unit
block containing the first endpoint, and, for a
pair with $d_X(x,x')\ge1$, also the block
containing its second endpoint. Thus it has at
most two elements.

\begin{proposition}\label{prop:interacting-local-insertions}
For $A\in\mathcal W$, the blocked field maps and
their Hölder maps are analytic in $A$. The first
coarse-operator derivative satisfies
\begin{align}
 |\partial_BD_*(A)|(y,y')
  &\le226(C^F)^2\epsilon_B
       \sum_{z\in P_B}E(y,z)E(z,y'),\notag\\
 \|\partial_BD_*(A)\|_{\gamma^F/2}
  &\le C_K\epsilon_B,\qquad
 \sum_{y,y'}|\partial_BD_*(A)|(y,y')
                           \le C_K\epsilon_B|P_B|.
                                  \label{eq:interacting-coarse-first}
\end{align}
If $d(P_B,P_{B'})\ge M+2$, then
\begin{align}
 |\partial_B\partial_{B'}D_*(A)|(y,y')
 \le{}& C_H^2\epsilon_B\epsilon_{B'}
       \sum_{z\in P_B,z'\in P_{B'}}
       \big(E(y,z)E(z,z')E(z',y')\notag\\
       &\hspace{38mm}+E(y,z')E(z',z)E(z,y')\big).
                                  \label{eq:interacting-coarse-mixed}
\end{align}
The corresponding fine-field legs satisfy
\begin{align}
 |\operatorname{Leg}_A|(u;y')
   &\le3C^F\sum_{w\in\operatorname{pos}(u)}E(w,y'),\notag\\
 |\partial_B\operatorname{Leg}_A|(u;y')
   &\le C_H\epsilon_B
       \sum_{w\in\operatorname{pos}(u),z\in P_B^+}E(w,z)E(z,y'),
                                  \label{eq:interacting-leg-first}\\
 |\partial_B\partial_{B'}\operatorname{Leg}_A|(u;y')
   &\le C_H^2\epsilon_B\epsilon_{B'}
       \sum_{\substack{w\in\operatorname{pos}(u)\\z\in P_B^+,z'\in P_{B'}^+}}
       \big(E(w,z)E(z,z')E(z',y')\notag\\
   &\hspace{37mm}+E(w,z')E(z',z)E(z,y')\big).
                                  \label{eq:interacting-leg-mixed}
\end{align}
The last bound has the same separation hypothesis.
These estimates hold for barred and unbarred legs.
\end{proposition}

\begin{proof}
The block averages and their transposes are diagonal
in unit blocks. From
\eqref{eq:interacting-phase-bounds} and
$\mathfrak i(A)\le1/4$ their majorants are at most
$3/2$. Their directional derivatives are bounded by
\[
                 (9/4)\epsilon_B1_{\{y\in P_B\}}.
\]
Differentiating the only phased Dirac hop gives
\[
 (\partial_BD_Av)(x)=
        -ie_*B(c,c_3)M_3(x)u_A(c,c_3)v(x^{(3)}).
\]
Its majorant is at most $8\epsilon_B$, vanishes
unless the unit blocks coincide or are neighbors,
and vanishes unless both belong to $P_B$.
The derivative of $JP_*(A)J^*$ has diagonal
majorant at most $7\epsilon_B$. Thus for
$F_B=\partial_BF_*(A)$,
\[
 |F_B|(z,z')\le\epsilon_B
       (15\,1_{\{z=z'\}}+8\,1_{\{z,z'\text{ neighbors}\}})
                                      1_{\{z,z'\in P_B\}}.
\]
Summing the six neighboring blocks gives
$\|F_B\|_\gamma\le146\epsilon_B$ for $0\le\gamma\le1$.
Also, since $E(z',y')\le\mathrm eE(z,y')$ for
neighboring $z,z'$,
\begin{equation}\label{eq:interacting-insertion-chain}
 \sum_{z,z'}E(y,z)|F_B|(z,z')E(z',y')
        \le146\epsilon_B\sum_{z\in P_B}E(y,z)E(z,y').
\end{equation}

If the supports are separated by $M+2$, the mixed
derivatives of $D_A$, the averages, $JP_*J^*$,
$F_*$ and $\delta_{\alpha,A}$ vanish. Each entry
is an exponential of a linear phase supported
in a unit block, a single external bond or a
coordinate box of diameter at most $M$.
Two separated directions cannot both occur in
one phase. For $JP_*J^*$ the remaining product
terms have disjoint unit-block supports.

The inverse derivative and the blocked-field
definition give
\[
 \partial_BH_*=-S_*F_BH_*+
               b_*S_*\partial_B(JQ_*^T(-A)).
\]
Use \eqref{eq:interacting-terminal-input} and
\eqref{eq:interacting-insertion-chain} in the first
term. The diagonal derivative in the second
is absorbed as the summand $z=y'$. The combined
coefficient is at most $149(C^F)^2$:
\begin{equation}\label{eq:interacting-H-first}
 |\partial_BH_*|(y,y')
 \le149(C^F)^2\epsilon_B
                    \sum_{z\in P_B}E(y,z)E(z,y').
\end{equation}
The same estimate holds for
$\delta_{\alpha,A}\partial_BH_*$ on short pairs,
since the first factor $S_*$ can be replaced
by $\delta_{\alpha,A}S_*$.

For separated directions differentiating once more
leaves four terms:
\begin{align*}
 \partial_B\partial_{B'}H_*={}&
 S_*F_BS_*F_{B'}H_*+S_*F_{B'}S_*F_BH_*\\
 &-b_*S_*F_BS_*\partial_{B'}(JQ_*^T(-A))
  -b_*S_*F_{B'}S_*\partial_B(JQ_*^T(-A)).
\end{align*}
Apply \eqref{eq:interacting-insertion-chain} twice
to the first two terms. Each gives an ordered
three-link exponential chain with coefficient
$146^2(C^F)^3$. The diagonal last factors in
the other two terms give the same chains, with
coefficient $(657/2)(C^F)^2$. Each ordered
chain therefore has coefficient at most
$21645(C^F)^3$. Again the first factor may
be replaced by the short-pair Hölder inverse.

Differentiate $D_*=b_*-b_*Q_*J^*H_*$.
The derivative of the average contributes
$(9/4)C^F$ to the first exponential chain and
the derivative of $H_*$ contributes
$(3/2)149(C^F)^2$. Their sum is at most
$226(C^F)^2$, proving the first line of
\eqref{eq:interacting-coarse-first}.
Weighted row and column sums give two factors
$(24/\gamma^F)^3$, hence $C_K$; without weights
the same sum retains one factor $|P_B|$.
For the mixed derivative the three product
terms contain one derivative on each factor,
or the mixed derivative of $H_*$. Each chain
has coefficient at most
$(1341/4)(C^F)^2+(64935/2)(C^F)^3
\le32803(C^F)^3\le C_H^2$.
This proves \eqref{eq:interacting-coarse-mixed}.

A leg is a component row of $H_*$ or of
$\delta_{\alpha,A}H_*$. For a short pair its
position is its first unit block; for a long
pair its position set contains both endpoint
blocks. The sum over spin components of the
row is bounded by the relevant majorant.
For long pairs use
$|\delta_{\alpha,A}v|\le\mathrm e|v(x')|+|v(x)|$;
for short pairs use the Hölder input directly.
This gives the first line of
\eqref{eq:interacting-leg-first}, and its
weighted row sum is at most $h_{\mathrm L}$.

The derivative of a Hölder leg has two terms:
$(\partial_B\delta_{\alpha,A})H_*+
\delta_{\alpha,A}\partial_BH_*$.
The first is bounded by
$4d_X^{1-\alpha}\epsilon_B|H_*v(x')|$,
and vanishes unless the pair is within $M$
of $P_B$. It vanishes for pairs inside a
cluster, since their phase is one.
For short pairs this adds at most
$11C^F$ to $149(C^F)^2$; for long pairs
it adds $4MC^F$ to $(4/3)149(C^F)^2$.
Both are at most $C_H$, after enlarging
$P_B$ to $P_B^+$. This proves the derivative
bound in \eqref{eq:interacting-leg-first}.

For its mixed derivative the second derivative
of the Hölder phase vanishes by separation.
There remain
$(\partial_B\delta_{\alpha,A})\partial_{B'}H_*$
and its interchange, and
$\delta_{\alpha,A}\partial_B\partial_{B'}H_*$.
On short pairs each ordered chain has
coefficient at most
$1639(C^F)^2+21645(C^F)^3$.
On long pairs the coefficient at the second
endpoint is at most
$596M(C^F)^2+28860(C^F)^3$;
the first endpoint has the smaller coefficient
$21645(C^F)^3$. All are at most $C_H^2$,
proving \eqref{eq:interacting-leg-mixed}.
Barred legs follow by the monomial charge
conjugation matrix and the replacement $A\mapsto-A$,
which preserve component row sums; the domain
is invariant under that replacement.
\end{proof}

\subsection{Complex neighborhoods and the real-field covariance}
\label{subsec:interacting-covariance}

Choose
\[
 c_\tau=\min\{1/8,\mu_-/48,(292\mathcal S)^{-1}\},\qquad
 C_K'=1328\mathcal S^2,\qquad
 \mathcal W^+=\{A+A_1:A\in\mathcal W,\ e_*|A_1|_\infty\le c_\tau\}.
\]
Every point of $\mathcal W^+$ has imaginary part
$\mathfrak i<1/4$. The elementary averages and
$F_*$ therefore satisfy the same local derivative
bounds as above, although the exponential propagator
input is initially given only on $\mathcal W$.
The enlargement preserves invertibility:
\begin{equation}\label{eq:interacting-enlarged-inverse}
 \|S_*(A+A_1)\|_{\gamma^F/2}\le2\mathcal S.
\end{equation}
To prove this, integrate $\partial_{A_1}F_*$
along the segment to get
$\|F_*(A+A_1)-F_*(A)\|_{\gamma^F/2}
\le146c_\tau$. Multiplying by $S_*(A)$ gives
norm at most $146c_\tau\mathcal S\le1/2$.
The Neumann inverse in the submultiplicative
majorant norm proves \eqref{eq:interacting-enlarged-inverse}.
Since its finite-dimensional determinant is
nonzero, $S_*$ is analytic in a neighborhood
of each point of $\mathcal W^+$, as is $D_*$.

Differentiate
$D_*=b_*-b_*^2Q_*J^*S_*JQ_*^T(-A)$.
The three terms have majorant-norm coefficients
\[
 \tfrac94(2\mathcal S)\tfrac32,\qquad
 \tfrac32(2\mathcal S)146(2\mathcal S)\tfrac32,\qquad
 \tfrac32(2\mathcal S)\tfrac94.
\]
Their sum is
$\frac{27}{2}\mathcal S+1314\mathcal S^2\le C_K'$.
Keeping the support indicator of either the
average derivative or $F_B$ gives, respectively,
\begin{equation}\label{eq:interacting-enlarged-derivative}
 \|\partial_BD_*(A)\|_{\gamma^F/2}\le C_K'\epsilon_B,\qquad
 \sum_{y,y'}|\partial_BD_*(A)|(y,y')
                                  \le C_K'\epsilon_B|P_B|,
                                      \qquad A\in\mathcal W^+.
\end{equation}
In the middle term, sum first the two outer
$S_*$ factors using their row and column bounds
$2\mathcal S$, and then sum the local derivative
over its first endpoint in $P_B$. Integrating
\eqref{eq:interacting-enlarged-derivative} along
any segment inside $\mathcal W^+$ gives the
same two bounds for $D_*(A+B)-D_*(A)$.

The covariance at a real field requires no
small-field hypothesis. If $A_r$ is real,
$R=D_{A_r}+\mu_*$ has real part at least $\mu_*$.
The averaging map $Q=Q_*(A_r)J^*$ satisfies
$Q^*=JQ_*^T(-A_r)$ and $QQ^*=1$, because its
phases have modulus one and every block has
fine measure one. Thus $F_*(A_r)=R+b_*Q^*Q$
also has real part at least $\mu_*$.
Set $C=QR^{-1}Q^*$ and $Y=(b_*^{-1}+C)^{-1}$.
The inverse $R^{-1}$ is accretive, so
$b_*^{-1}+C$ is invertible. Multiplication
verifies
\[
 S_*=R^{-1}-R^{-1}Q^*YQR^{-1},\qquad
 Q S_*Q^*=b_*^{-1}(1-b_*^{-1}Y).
\]
The definition of $D_*$ therefore gives
\begin{equation}\label{eq:interacting-coarse-covariance}
 D_*(A_r)^{-1}=b_*^{-1}+QR^{-1}Q^*,\qquad
 \|D_*(A_r)^{-1}\|_2\le b_*^{-1}+\mu_*^{-1}.
\end{equation}

Its block decay follows by conjugating the fine
inverse with a weight. Fix $y'$ and put
$\psi(c)=\vartheta\,d(c,\Delta_{y'})$ on the
periodic fine lattice, with
$\vartheta=\mu_*/(12\mathrm e)\le1$.
This weight is zero on $\Delta_{y'}$, is
$\vartheta$-Lipschitz, and on $\Delta_y$ is
at least $\vartheta(d(y,y')-1)$.
Only the external Dirac hop changes under
conjugation by $e^\psi$, since the internal
hops stay over one photon site. Consequently
\[
 \|e^\psi D_{A_r}e^{-\psi}-D_{A_r}\|_2
 \le6\eta_*^{-1}(e^{\vartheta\eta_*}-1)
 \le6\mathrm e\vartheta=\mu_*/2.
\]
The conjugated $R$ has real part at least
$\mu_*/2$, so its inverse has norm at most
$2/\mu_*$. Apply it between $Q^*$ times
unit spinors at $y$ and $y'$; these have
norm one and supports in their respective
blocks. Moving the weight across the
scalar product gives
\begin{equation}\label{eq:interacting-real-covariance-decay}
 \|D_*(A_r)^{-1}(y,y')\|_2
  \le b_*^{-1}1_{\{y=y'\}}+
       \frac{2\mathrm e}{\mu_*}
                       e^{-\mu_*d(y,y')/(12\mathrm e)}.
\end{equation}
The factor $\mathrm e$ accounts for the unit
block diameter in the lower bound for $\psi$.

Set
\[
 g=(b_-^{-1}+\mu_-^{-1})^{1/2},\qquad
 h_G=1+g,\qquad g_1=b_-^{-1}+4/\mu_-,
 \qquad \theta_{\mathrm{CT}}=\mu_-/(12\mathrm e),
\]
\[
 \vartheta_0=\min\{\theta_{\mathrm{CT}}/2,\gamma^F/8,1/M\},
 \qquad
 \alpha_G=2b_-^{-1}
        +4\mathrm e\mu_-^{-1}(24/\theta_{\mathrm{CT}})^3.
\]
A component row or column of a $4\times4$ matrix
has sum at most twice its Hilbert-space norm.
Using \eqref{eq:interacting-real-covariance-decay}
and the lattice exponential sum, every row and
column of $G=D_*(A_r)^{-1}$ therefore satisfies
\begin{equation}\label{eq:interacting-Gaussian-covariance}
 \sum_{u'}|G(u,u')|e^{\vartheta_0d(u,u')}
       \le\alpha_G,\qquad
 \sum_u|G(u,u')|e^{\vartheta_0d(u,u')}
       \le\alpha_G,\qquad \|G\|_2\le g^2.
\end{equation}
These constants depend on $L,C^F,\gamma^F,M$,
not on the size of the real connection.

For later complex twists we also need the
coarse inverse away from real fields. Suppose
$A=A_r+A_1$, with $A_r$ real,
$e_*|A_1|_\infty\le\mu_*/24$, and $F_*(A)$
invertible. The external-hop phase difference
gives $\|D_A-D_{A_r}\|_2\le12e_*|A_1|_\infty
\le\mu_*/2$. Hence $D_A+\mu_*$ is accretive
with inverse norm at most $2/\mu_*$. Its two
average maps have norm at most $\sqrt2$,
because $|\varphi_A^{\pm1}|\le e^{1/16}\le\sqrt2$.
The same finite-rank inversion calculation
then gives
\begin{equation}\label{eq:interacting-complex-covariance}
 D_*(A)^{-1}=b_*^{-1}+
 Q_*(A)J^*(D_A+\mu_*)^{-1}JQ_*^T(-A),
 \qquad
 \|D_*(A)^{-1}\|_2\le g_1.
\end{equation}
Here invertibility of $b_*^{-1}+QR^{-1}Q'$
follows from that of $F_*(A)$ and $R$:
$\det(1+b_*R^{-1}Q'Q)=\det(1+b_*QR^{-1}Q')$.
This determinant identity follows by the two
block-triangular factorizations of
$\left(\begin{smallmatrix}1&-b_*R^{-1}Q'\\Q&1\end{smallmatrix}\right)$.
The inverse formula is then checked by
multiplication exactly as above.
Each covariance entry is bounded by $g_1$,
so each row has sum at most $4\ell_*^3g_1$.

\subsection{The terminal Grassmann norms}
\label{subsec:interacting-terminal-norms}

The terminal integral is over four coarse
components per unit site, with Grassmann
weight one. Set
\[
 a_1=e_*^{-1/4},\qquad
 a_2=C_0^{-1}e_*^{-1/4+\varepsilon},\qquad
 h_*=(e_*^{-1/4},e_*^{-1/4+\varepsilon}).
\]
By \eqref{eq:interacting-global-field-rows},
the substitution $H_*^\#(A)\Psi$ has row
norm at most $C^F\le C_0$, for both spinor
and H\"older legs, including all long
and cluster pairs. The coarse kernel
$b_*-D_*(A)=b_*^2Q_*J^*S_*JQ_*^T(-A)$
is bounded by $C^FE(y,y')$: block
compression gives Euclidean constant
$32\exp(1)b^2a'$ by
\eqref{eq:block-compression-bound}, and
conversion to the component norm doubles
it. The calibration of $C^F$ applies.

The weighted Grassmann substitution estimate
applies with $C_0a_2=h_{*,2}\le h_{*,1}$.
Therefore
\begin{equation}\label{eq:interacting-terminal-substitution}
 |F(H_*^\#(A)\Psi)|_{a_2}\le\|F\|_{h_*}.
\end{equation}
For the mass form before substitution, the
fine measure sums to $\ell_*^3$ and there
are four spin components. Its norm at
weight $C_0a_2$ is exactly
$4\ell_*^3(C_0a_2)^2$. In particular
\begin{equation}\label{eq:interacting-terminal-mass-form}
 |\langle\bar\psi,\psi\rangle|_{a_2}
                   \le4e_*^{-1/2+2\varepsilon}\ell_*^3.
\end{equation}
The $48$ labels are canceled by the fine
measure $1/48$; there is no factor $192$
in this estimate. Similarly the coarse
bilinear form has norm $a_1^2$ times the
sum of its matrix entry moduli. Its diagonal
term contributes $4b_*$ per unit site,
and its decaying kernel contributes at
most $4C^F\sum_{y'}E(y,y')$.
For $c(\gamma)=\sum_{z\in\Z^3}e^{-\gamma|z|_\infty}$,
the earlier lattice-sum estimate gives
$c(\gamma)\le(12/\gamma)^3$ when $0<\gamma\le6$.
The chosen $C_0\ge4+16C^Fc(\gamma^F)$ therefore gives
\begin{equation}\label{eq:interacting-terminal-Dirac-form}
 |\langle\bar\Psi,D_*(A)\Psi\rangle|_{a_1}
                                  \le C_0e_*^{-1/2}\ell_*^3.
\end{equation}
These estimates hold on $\widetilde G_*$.
For the external activity evaluated at
$\mathscr T=\mathscr A+\mathscr C$, the
same substitution applies because
$\widetilde R_*+\mathfrak X_*\subset\widetilde G_*$.

\subsection{Gaussian bounds in the zero-flux sector}
\label{subsec:interacting-zero-flux-Gaussian}

The complement estimate below concerns the normalized
zero-flux law, denoted $\bar\mu$.
The all-flux Gaussian bound does not directly
give a bound for this conditional law.
We therefore retain the sector projection in
the coordinate-slice trace.

On $C(n_1,n_2,n_3)$ use the spaces $H^s,H^l$,
$P_G$, $E$, and $M_0$ of
\eqref{eq:slab-electric-operators} and
\eqref{eq:slab-magnetic-multiplier}.
Let $P_k$ project onto the magnetic-flux
summand $k$. Choose coordinate slice cycles
$S_2,S_3$, through a fixed slice site, and
for $\xi\in[0,1]^2$ set
$b_\xi=q(\xi_2\mathbf1_{S_2}+\xi_3\mathbf1_{S_3})$.
For bounded magnetic factors $u_j$ and
electric factors $v_j$ with sub-Gaussian growth,
let $K_j^\xi$ have the periodized convolution kernel
\[
 v_j(A'-A)e^{i\langle b_\xi,A'-A\rangle}
                         \prod_bg_1(A'(b)-A(b)).
\]
The zero-flux trace formula is
\begin{equation}\label{eq:interacting-zero-flux-Bloch}
 Z^{(0)}\left[\prod_ju_j(F_j^{\rm mag})v_j(F_j^{\rm el})\right]
 =(2\pi)^{n_1n_2n_3}\int_{[0,1]^2}
 \tr\left(P_0\prod_{j=n_1-1,\ldots,0}
                  r_q^{n_2n_3}K_j^\xi P_GM_0u_j(F^{\rm mag})M_0^*\right)
                                                        \,\dd\xi .
\end{equation}
To derive it, sum the electric curvature along
$S_2,S_3$ and over time. The sums are
$r_qk_{12}$ and $r_qk_{13}$.
The added electric phases are therefore
$e^{2\pi i(\xi_2k_{12}+\xi_3k_{13})}$.
Restricting the first magnetic slice to
flux $k_{23}$ inserts $P_{k_{23}}$ in the
word trace, since the sum of $d_sA$ on
the slice vanishes. All operators preserve
that summand. The flux sum is absolutely
convergent by the curvature Gaussian bound
on the modulus of the integrand.
Set $k_{23}=0$ and integrate over $\xi$;
Fourier orthogonality keeps precisely
$k_{12}=k_{13}=0$. This proves
\eqref{eq:interacting-zero-flux-Bloch}.

\begin{proposition}\label{prop:interacting-zero-flux-domination}
For the harmonic-average zero-flux law at every
real lattice mass and every coordinate side
$n_\mu\ge4$,
\begin{equation}\label{eq:interacting-zero-flux-Gaussian}
 \bar\mu(e^{\langle g,F^{\rm u}\rangle})
                  \le e^{\|\operatorname{Pr}_{F_N}g\|^2/2}.
\end{equation}
In particular, for $\phi\in F_N$ with
$\|\phi\|\le1$,
$\bar\mu\{|\langle\phi,F^{\rm u}\rangle|>a\}\le2e^{-a^2/2}$.
\end{proposition}

\begin{proof}
Put $\mathcal Z(\beta)=
Z^{(0)}[e^{-\langle\beta,F\rangle-\|\beta\|^2/2}]$.
It is finite by the all-flux bound and is
positive because the sector contains the
$z=0$ term, which has positive integral;
the shifted photon weight is strictly positive.
For time direction $1$, split $\beta$
into electric parts $\beta_j^{\rm el}$ and
magnetic parts $\beta_j^{\rm mag}$.
Take the corresponding exponential $u_j$
and $v_j$ in
\eqref{eq:interacting-zero-flux-Bloch}, replacing
$u_j$ first by $u_j^\delta=\min(u_j,\delta^{-1})$.
The bounded magnetic factors make the
word trace applicable.

Completing the electric Gaussian gives
\[
 K_j^\xi=z_jU_j e^{-|E-b_\xi|^2/2},\qquad
 U_j=e^{i\langle\beta_j^{\rm el},E\rangle},\qquad
 |z_j|=1 .
\]
Let $K_\xi=r_q^{n_2n_3}e^{-|E-b_\xi|^2/2}P_G$.
It is positive, and its square root is
$r_q^{n_2n_3/2}e^{-|E-b_\xi|^2/4}P_G$.
The operator $K_\xi^{1/2}M_0$ is Hilbert--Schmidt.
The periodized square-root Gaussian has the
additional phase $e^{i\langle b_\xi,A'-A\rangle}$;
its modulus is bounded by the kernel majorant
used in \eqref{eq:slab-shifted-hs}.
The same real-slice integral and flux sum
therefore proves this assertion uniformly
in $\xi$. Multiplication by bounded
$(u_j^\delta)^{1/2}$ preserves it.

Put $Y_j=K_\xi^{1/2}M_0(u_j^\delta)^{1/2}$.
Cyclicity turns the trace into
$\prod_jz_j$ times the closed chain
\[
 \tr\prod_j U_jP_0Y_j(P_0Y_j)^*.
\]
All factors preserve the magnetic-flux
summand; $U_j$ is unitary. Schatten Hölder
with $2n_1$ factors of exponent $2n_1$
bounds this trace by
\[
 \prod_j\left(\tr_{H^s}P_0(Y_j^*Y_j)^{n_1}\right)^{1/n_1}.
\]
Hölder in $\xi$ and the Bloch trace formula
identify its $j$th integrated factor with
the zero-flux functional having the same
magnetic factor at every slice and no
electric factor. Let $\beta^{(j)}$ denote
that purely magnetic two-cochain. Monotone
convergence as $\delta\downarrow0$ gives
\[
 \mathcal Z(\beta)\le
                \prod_j\mathcal Z(\beta^{(j)})^{1/n_1}.
\]
Rotate the torus by an orientation-preserving
cyclic coordinate permutation. Zero flux
and the pairing are preserved, and
$\beta^{(j)}$ becomes purely electric for
the new time direction. The same inequality
then bounds $\mathcal Z(\beta^{(j)})$
by $\mathcal Z(0)$. Hence
$\mathcal Z(\beta)\le Z^{(0)}[1]$.

The zero-sector identity
\eqref{eq:external-zero-sector} identifies its
normalized curvature law with $\bar\mu$.
Since this law is supported on $F_N$,
the pairing with $g$ equals that with
$\operatorname{Pr}_{F_N}g$. Apply the
preceding inequality with the negative of
that projection and divide by the positive
zero-sector partition function. This proves
\eqref{eq:interacting-zero-flux-Gaussian}.
Exponential Markov applied to the two
signs and optimized at parameter $a$
gives the stated tail bound.
\end{proof}


\section{The terminal integral and partition bounds}
\label{sec:terminal-partition}

We first put the curvature Gaussian in coordinates
adapted to all block steps. The energy splits into
the terminal photon energy and independent unit
Gaussian fluctuation coordinates. Iterating the exact
density identity in these coordinates then gives
the terminal fermion integral on the small-field set.

\subsection{The downward photon coordinates}
\label{subsec:term-photon-tower}

Fix a flow side and cutoff as in
Proposition~\ref{prop:rgfl-scales}. Put $k=k_*(R)$ and
$n=R/\eta=L^k\ell_*$.
On the unit torus of side $\ell_j=n/L^j$ use the
fluctuation space $V_j$ and its fixed orthonormal
basis, and the maps $A_j^{\min},\mathrm{Sq}_j$ of
the photon construction.
For a real terminal unit bond field $A$ and
$\zeta=(\zeta_j)_{j<k}$ define
\begin{equation}\label{eq:term-downward-recursion}
 A_k^\downarrow=A,\qquad
 A_j^\downarrow=A_j^{\min}(A_{j+1}^\downarrow)
                         +\mathrm{Sq}_j\zeta_j,\qquad j=k-1,\ldots,0 .
\end{equation}
The last field is denoted $A_0^\downarrow(A,\zeta)$.
These fields are linear in their data.

For any fine unit field $A_0$ define its block fields by
$A_j=(\mathbf Q^{(j)}A_0)^{[j]}$, where $\mathbf Q^{(j)}$
is the $j$-fold photon average on the original unit torus and
$B^{[j]}(b)=L^{j/2}B(L^jb)$.
The block-averaging/scaling identities give
$(A_{j+1})_L=\mathbf QA_j$.
Let $\mathrm{Ax}_k$ be the fields $A_0$ for which
$\tau A_j=0$ for $j<k$.
For each unit torus put
$V_{\mathrm t}=\delta C^2$, the real coclosed bond
fields with mean zero.
On the fine torus let $F_N=dC^1(n)$ with its
counting inner product and standard Gaussian
$\gamma_N$. The slice is
\[
 \mathrm{Sl}_k=\{A_0\in\mathrm{Ax}_k:A_k\in V_{\mathrm t}(T_k^0)\}.
\]

We will use the elementary cochain identities
\begin{equation}\label{eq:term-unit-Hodge}
 C^1=dC^0\oplus\{\text{constant bond fields}\}\oplus\delta C^2,
       \qquad \ker d=dC^0\oplus\{\text{constant bond fields}\}.
\end{equation}
The sum is orthogonal. Indeed $dd=0$, so $dC^0$
and $\delta C^2$ are orthogonal, and their common
orthogonal complement is $\ker(d\delta+\delta d)$.
Commuting lattice differences make this operator
the scalar Laplacian in each component.
Its kernel consists of constants: the energy identity
forces every nearest-neighbor difference to vanish.
This proves \eqref{eq:term-unit-Hodge}.
Consequently $d$ is injective on $V_{\mathrm t}$ and,
on a unit torus with $\ell^3$ sites,
$\dim V_{\mathrm t}=\dim dC^1=2\ell^3-2$.

\begin{proposition}\label{prop:term-photon-tower}
The map $(A,\zeta)\mapsto A_0^\downarrow(A,\zeta)$
is a linear bijection from
$C^1(T_k^0)\times\prod_{j<k}V_j$ to $\mathrm{Ax}_k$.
The block fields of its image are exactly
$A_j^\downarrow$.
The coboundary is a linear bijection
$d:\mathrm{Sl}_k\to F_N$.
Thus curvature has coordinates
\begin{equation}\label{eq:term-curvature-coordinates}
 \mathcal J_{\mathrm{ph}}:
 F\longmapsto(A_*;\zeta_0,\ldots,\zeta_{k-1}),\qquad
 F=dA_0^\downarrow(A_*,\zeta),
\end{equation}
where $A_*\in V_{\mathrm t}(T_k^0)$.
They satisfy the exact energy identity
\begin{equation}\label{eq:term-tower-energy}
 \|F\|^2=\|dH_k^{\mathrm{ph}}A_*\|_2^2
                                  +\sum_{j<k}\|\zeta_j\|^2 .
\end{equation}
The image of $\gamma_N$ under
$\mathcal J_{\mathrm{ph}}$ is
\begin{equation}\label{eq:term-tower-Gaussian}
                  \dd\mathsf m_*(A_*)\prod_{j<k}\dd\mu_I(\zeta_j),
\end{equation}
where the terminal centered Gaussian on
$V_{\mathrm t}(T_k^0)$ is nondegenerate and has density
proportional to
\[
                  \exp(-\|dH_k^{\mathrm{ph}}A\|_2^2/2).
\]

There are orthonormal vectors
$\phi_{j,\Upsilon}\in F_N$ such that
$\zeta_j(\Upsilon)=\langle\phi_{j,\Upsilon},F\rangle$.
For each terminal unit plaquette $p$ there is
$\phi_p\in F_N$ with
\begin{equation}\label{eq:term-tower-row-norms}
 (dA_*)(p)=\langle\phi_p,F\rangle,\qquad \|\phi_p\|\le1.
\end{equation}
There are $3\ell_*^3$ such plaquettes, and
\begin{equation}\label{eq:term-tower-coordinate-count}
 \dim V_j=(2-2L^{-3})\ell_j^3
                       \le2\ell_*^3L^{3(k-j)} .
\end{equation}
\end{proposition}

\begin{proof}
The one-step minimizer and fluctuation maps obey
$\mathbf QA_j^{\min}(A')=A'_L$, $\tau A_j^{\min}(A')=0$,
$\mathbf Q\mathrm{Sq}_j=0$ and $\tau\mathrm{Sq}_j=0$.
Apply $\mathbf Q$ to \eqref{eq:term-downward-recursion};
the injective scaling map then identifies the next
block field with $A_{j+1}^\downarrow$.
Starting from $A_0=A_0^\downarrow$ proves this at all
levels, including $\tau A_j=0$.
Conversely, for $A_0\in\mathrm{Ax}_k$ its difference
$A_j-A_j^{\min}(A_{j+1})$ is annihilated by $\mathbf Q$ and
$\tau$. The coordinate bijection
\eqref{eq:photon-coordinate-bijection} and the invertible
square root give exactly one $\zeta_j$ with this
difference equal to $\mathrm{Sq}_j\zeta_j$.
Downward recursion recovers $A_0$.
This proves the first bijection and its inverse.

The orthogonality
\eqref{eq:photon-block-minimizer-orthogonality}
and the Gaussian fluctuation decomposition give
at each level
\[
 \|dH_j^{\mathrm{ph}}A_j\|_2^2
       =\|dH_{j+1}^{\mathrm{ph}}A_{j+1}\|_2^2+\|\zeta_j\|^2 .
\]
Summing, and using $H_0^{\mathrm{ph}}=I$, proves
\eqref{eq:term-tower-energy}.
The dimension of the slice is
\[
 2\ell_k^3-2+\sum_{j<k}(2\ell_j^3-2\ell_{j+1}^3)
                                                    =2n^3-2=\dim F_N.
\]
If $dA_0=0$ on the slice, the energy identity forces
every $\zeta_j=0$ and $dH_k^{\mathrm{ph}}A_k=0$.
Since $\mathbf Q_kH_k^{\mathrm{ph}}=I$ and
$d\mathbf Q_k=\mathbf Q_k^{(2)}d$,
this implies $dA_k=0$.
The injectivity of $d$ on $V_{\mathrm t}$ gives
$A_k=0$, whence $A_0=0$ by the downward bijection.
Equality of dimensions therefore proves that
$d:\mathrm{Sl}_k\to F_N$ is bijective and gives
\eqref{eq:term-curvature-coordinates}.

Equip the coordinate space with the scalar product
\[
 \langle(A,\zeta),(A',\zeta')\rangle_{\mathrm{ph}}
       =\langle A,\Delta_kA'\rangle
                                      +\sum_{j<k}\langle\zeta_j,\zeta_j'\rangle.
\]
Its terminal part is positive definite:
$\langle A,\Delta_kA\rangle
=\|dH_k^{\mathrm{ph}}A\|_2^2\ge\|dA\|_2^2>0$
for nonzero $A\in V_{\mathrm t}$.
Here the block curvature contraction follows from
\eqref{eq:photon-averaging-identities}.
The energy identity and polarization make
$\mathcal J_{\mathrm{ph}}$ a linear isometry for
this scalar product.
It maps orthonormal bases to orthonormal bases;
the Jacobian is therefore $1$ for their Euclidean
volume measures. These are the measures in the
normalized standard Gaussians.
In an orthogonal direct sum, that volume measure
and the Gaussian density factor into those of
the summands. This proves
\eqref{eq:term-tower-Gaussian}, including its normalization.
The volume measure for the terminal scalar product
differs from the ordinary Lebesgue measure on
$V_{\mathrm t}$ by one constant factor.
Thus its Gaussian has exactly the stated density.
Its quadratic form is positive definite, so it is
nondegenerate.

For a basis vector $\Upsilon$ in $V_j$, let
$\mathbf e_{j,\Upsilon}$ be the coordinate vector
with that single nonzero entry.
Its inverse image under the isometry is
$\phi_{j,\Upsilon}$.
The coordinate vectors are orthonormal, proving
orthonormality and the required coordinate formula.
The terminal curvature functional satisfies
\[
 |(dA_*)(p)|
 \le\|dA_*\|_2
   =\|\mathbf Q_k^{(2)}dH_k^{\mathrm{ph}}A_*\|_2
 \le\|dH_k^{\mathrm{ph}}A_*\|_2\le\|F\|.
\]
It is consequently represented by a vector of norm
at most $1$ in the finite Euclidean space $F_N$,
giving \eqref{eq:term-tower-row-norms}.
Finally \eqref{eq:photon-fluctuation-dimension} and
$\ell_j=L^{k-j}\ell_*$ give
\eqref{eq:term-tower-coordinate-count}.
\end{proof}


If a terminal field $A^c$ is closed, linearity gives
the difference of downward fields as
$A_0^\downarrow(A^c,0)$.
For a closed input the background identity
\eqref{eq:photon-background-one-step} gives zero
curvature energy at the preceding level. Applying
the average shows that its block minimizer is closed.
The identity $\mathbf QA_j^{\min}(A')=A'_L$ and
preservation of means by the block average give
$\operatorname{mean}(A_j^{\min}(A'))=
L^{-1/2}\operatorname{mean}(A')$.
Its final mean is $L^{-k/2}\operatorname{mean}(A^c)$.
Equation~\eqref{eq:term-unit-Hodge} therefore gives
\begin{equation}\label{eq:term-closed-terminal-shift}
 A_0^\downarrow(A+A^c,\zeta)
    =A_0^\downarrow(A,\zeta)
        +L^{-k/2}\operatorname{mean}(A^c)+d\lambda .
\end{equation}
In particular adding a coarse gradient adds a fine
gradient. For the harmonic twist
$c_{\theta,\mu}=\theta_\mu/(e_*\ell_*)$,
$L^{-k/2}c_\theta=H_\theta$, where
$H_{\theta,\mu}=\theta_\mu/(qn)$.
Thus its downward shift is $H_\theta+d\lambda$.

\subsection{The terminal density with an external connection}
\label{subsec:term-density}

Let $(m_{\mathrm{in}},\varepsilon_{\mathrm{in}}^{\mathrm v})$
be a controlled real pair on $R$, with the trajectory
of Section~\ref{sec:rg-flow}.
An external connection is a complex fine unit bond
field $a$ satisfying
$a^{[j]}\in\mathfrak X_j$ for every $j\le k$;
here $a^{[j]}(b)=L^{j/2}a(L^jb)$.
Write $\mathscr C_*=a^{[k]}$.
For $A\in\mathrm{Dom}_k$ put
$\mathscr A=H_k^{\mathrm{ph}}A$,
$B=\mathscr A+\mathscr C_*$, and insert the fermion
fields $H_k^\#(B)\Psi$ and their barred partners.
The terminal function is
\begin{align}
 Y_*^x(A;\mathscr C_*)
   ={}&\det D_k(0)^{-1}
        \exp\left(-\varepsilon_k^{\mathrm v}\ell_*^3
                       +\sum_X\mathfrak P_k(X,\mathscr C_*)\right)
       \nonumber\\
   &\times\int\mathrm D\Psi\,
      \exp\left(-\langle\bar\Psi,D_k(B)\Psi\rangle
                -m_k\langle\bar\psi,\psi\rangle
                       +\sum_X\mathscr E_k(X;\mathscr A,\mathscr C_*,
                                                        \psi,\chi)\right).
                                      \label{eq:term-Y-density}
\end{align}
Since $k\ge1$, this Berezin integral has $4\ell_*^3$
pairs with weight $1$. The inserted chamber fields
and kernels retain their fine weight $\eta_k^3/48$.
The inverse and activity domains in the preceding
sections make every term well defined, and
$\det D_k(0)\ne0$.
At zero external slot,
$\mathfrak P_k(X,0)=0$ and $\mathscr E_k|_0=E_k$;
we write $Y_*(A)=Y_*^x(A;0)$.

For the original fine unit field let
\begin{equation}\label{eq:term-fine-weight}
 w_N(B)=
 \exp(-\varepsilon_{\mathrm{in}}^{\mathrm v}n^3)
       \frac{\det(D^{\mathrm{ch}}[qB]+\bar m\eta+m_{\mathrm{in}})}
                              {\det(D^{\mathrm{ch}}[0]+\bar m\eta)} .
\end{equation}
The matrices have size $192n^3$.
The bracket argument of $D^{\mathrm{ch}}$ is the bond phase,
as in \eqref{eq:response-free-ratio}; an unabsorbed unit
connection $B$ therefore has argument $qB$.
The denominator is nonzero by the free positive-mass
coercivity estimate; no invertibility of the numerator
is required.
The weighted Gaussian Berezin identity gives
the initial integrated density exactly as $w_N(A_0+a)$.
In particular all initial measure and vacuum factors
are included in \eqref{eq:term-fine-weight}.

This fine weight is entire in $B$ and is invariant
under real gauge shifts.
Its entries involve only the four chamber neighbors,
with link phases in the fourth neighbor.
Multiplication by $\exp(-\mathrm i q\lambda(c))$ on
all chamber components over $c$ conjugates the operator
at $B$ to the operator at $B+d\lambda$: the endpoint
phases cancel the gradient in the fourth hopping phase,
and the other hops stay over the same $c$.
The determinants therefore agree.
The same stencil bounds every row sum at $B_0+a$,
with $B_0$ real, by
$|\boldsymbol\mu|+41+6\exp(q\|\operatorname{Im}a\|_\infty)$,
where $\boldsymbol\mu=\bar m\eta+m_{\mathrm{in}}$.
Indeed the three internal hopping block norms sum
with the diagonal to less than $|\boldsymbol\mu|+41$,
and the fourth has row norm at most
$6\exp(q\|\operatorname{Im}a\|_\infty)$.
The determinant expansion bounds its modulus by
the product of the absolute row sums: its permutation
sum is a subsum of the sum over all column choices.
Consequently
\begin{equation}\label{eq:term-finite-density-bound}
 |w_N(B_0+a)|\le B_N(a):=
 \frac{\exp(-\varepsilon_{\mathrm{in}}^{\mathrm v}n^3)}
            {|\det(D^{\mathrm{ch}}[0]+\bar m\eta)|}
       \big(|\boldsymbol\mu|+41
                     +6\exp(q\|\operatorname{Im}a\|_\infty)\big)^{192n^3}
                                                        <\infty .
\end{equation}
This bound is used for integrability at a fixed cutoff.

Let $A(F)\in V_{\mathrm t}$ be the transverse fine
potential of $F\in F_N$, and define the finite complex
measure
\[
 \dd\nu^{\theta,a}(F)=
               w_N(A(F)+H_\theta+a)\dd\gamma_N(F).
\]
The small-field event and the terminal indicator are
\begin{equation}\label{eq:term-small-field-event}
 \mathfrak G=\{F:|\zeta_j(\Upsilon)|\le p_{0,j}\ (j<k),
                                 \ \|dA_*\|_\infty\le p_*\},\qquad
 \chi_*(A)=\mathbf1[\|dA\|_\infty\le p_*].
\end{equation}
Here $A_*,\zeta$ are the actual coordinates
\eqref{eq:term-curvature-coordinates}.
Products with $\chi_*$ are extended by zero off
its support.

\begin{theorem}\label{thm:term-density}
For every $A\in\mathrm{Dom}_k$,
\begin{equation}\label{eq:term-iterated-density}
 Y_*^x(A;\mathscr C_*)
      =\int\prod_{j<k}
                    \big(\mathbf1_{\mathrm{Box}_j}(\zeta_j)
                                             \dd\mu_I(\zeta_j)\big)
                       w_N(A_0^\downarrow(A,\zeta)+a).
\end{equation}
The integrand is continuous on the compact product
of boxes.
If $A,A+d\lambda\in\mathrm{Dom}_k$, with $\lambda$
real, their terminal functions agree.
For every bounded Borel function $\Phi$ on
$V_{\mathrm t}(T_k^0)$ and every real twist $\theta$,
\begin{equation}\label{eq:term-Borel-observable}
 \nu^{\theta,a}(\mathbf1_{\mathfrak G}\Phi(A_*))
    =\int\dd\mathsf m_*(A)\,
                        \Phi(A)\chi_*(A)Y_*^x(A+c_\theta;\mathscr C_*).
\end{equation}
Its right-hand density is bounded and Borel, and
$Y_*^x(A+c_\theta;\mathscr C_*)$ is jointly continuous
in $(\theta,A)$ on $\{\chi_*=1\}$.
\end{theorem}

\begin{proof}
The photon domain conditions
\eqref{eq:photon-domain-small-charge} hold at each
running scale by the fixed charge choice.
Proposition~\ref{prop:photon-domain-preservation} gives
two consequences used here: a real unit field with
curvature bounded by $p_j$ lies in $\mathrm{Dom}_j$;
and a field in $\mathrm{Dom}_{j+1}$ descends to
$\mathrm{Dom}_j$ after addition of any box fluctuation.
Downward induction therefore puts every field of
\eqref{eq:term-downward-recursion} in its respective
domain whenever its top field lies in $\mathrm{Dom}_k$ and
all fluctuations lie in their boxes.

Let $y_j^x(A;a^{[j]})$ be the integrated density
at level $j$ along this trajectory.
The exact integrated step gives
\[
 y_{j+1}^x(A';a^{[j+1]})
     =\int_{\mathrm{Box}_j}\dd\mu_I(\zeta_j)\,
                   y_j^x(A_j^{\min}(A')+\mathrm{Sq}_j\zeta_j;
                                                               a^{[j]}).
\]
Its external slot is the required one because
$(a^{[j+1]})_L=a^{[j]}$, directly from the two
scaling factors.
Its coupling maps are the actual trajectory maps,
and its vacuum includes the box normalization.
The first such identity integrates the full
$192n^3$ weighted chamber pairs into
$4\ell_1^3$ coarse pairs.

For a top field in $\mathrm{Dom}_j$, the function
$w_N(A_0^\downarrow(A,\zeta)+a)$ is continuous in
the finitely many fluctuation coordinates.
Every finite product of boxes is compact,
and \eqref{eq:term-finite-density-bound} bounds it.
Thus all repeated integrals are absolutely integrable.
Starting with the initial identity
$y_0^x(A_0;a)=w_N(A_0+a)$, use the one-step formula,
the downward recursion and Fubini at each level.
This proves \eqref{eq:term-iterated-density} at $j=k$,
with its unnormalized factors
$\mathbf1_{\mathrm{Box}_j}\dd\mu_I$.
The continuity assertion follows from the same
finite-dimensional composition.
Equation~\eqref{eq:term-closed-terminal-shift}
turns a coarse real gradient into a fine gradient;
gauge invariance of $w_N$ inside this integral
then proves gauge invariance of $Y_*^x$.

To integrate against an arbitrary bounded Borel
observable, first express the fine density in the
curvature coordinates.
The slice potential $A_0^\downarrow(A_*,\zeta)$
has mean zero: its block average has the mean of
the fine field, and its top field has mean zero.
Its difference from $A(F)$ is closed and has mean
zero, hence is a gradient by
\eqref{eq:term-unit-Hodge}.
The harmonic shift identity gives
\[
 w_N(A(F)+H_\theta+a)
             =w_N(A_0^\downarrow(A_*+c_\theta,\zeta)+a).
\]
Under the tower isometry the Gaussian becomes
\eqref{eq:term-tower-Gaussian}, and the event
$\mathfrak G$ has indicator
\[
              \chi_*(A_*)\prod_{j<k}\mathbf1_{\mathrm{Box}_j}(\zeta_j).
\]
The integrand is a bounded Borel function with bound
$\|\Phi\|_\infty B_N(a)$.
The change-of-variables identity for its image
measure and Fubini therefore give
\begin{align*}
 \nu^{\theta,a}(\mathbf1_{\mathfrak G}\Phi(A_*))
   =\int\dd\mathsf m_*(A)\,\Phi(A)\chi_*(A)
      \int\prod_{j<k}
                      \big(\mathbf1_{\mathrm{Box}_j}\dd\mu_I(\zeta_j)\big)
                         w_N(A_0^\downarrow(A+c_\theta,\zeta)+a).
\end{align*}
Where $\chi_*(A)=1$, the field $A+c_\theta$ has
the same curvature bound and belongs to
$\mathrm{Dom}_k$. The inner integral is
$Y_*^x(A+c_\theta;\mathscr C_*)$ by
\eqref{eq:term-iterated-density}, proving
\eqref{eq:term-Borel-observable}.

The inner integral defines a function on all real
$(\theta,A)$, whether or not $A+c_\theta\in\mathrm{Dom}_k$.
Its modulus is at most $B_N(a)$.
Its integrand is continuous and uniformly continuous
on a compact neighborhood of any $(\theta,A)$
times the compact box product. Integrating this
uniform bound proves joint continuity.
Multiplying by $\chi_*$ gives a bounded Borel
function, since $\{\chi_*=1\}$ is closed.
On that set it is the asserted terminal density.
\end{proof}

\subsection{The large-field complement}
\label{subsec:term-complement}

At zero external connection write $\nu^\theta=\nu^{\theta,0}$,
and let $\bar\nu$ be its average over $[-\pi,\pi]^3$
with measure $\dd\theta/(2\pi)^3$.
The real chamber operator has two adjoint spin blocks;
thus its determinant at every real lattice mass is
the squared modulus of the determinant of one block.
Consequently $\nu^\theta$ is positive.
The zero-sector identity
\eqref{eq:external-density-zero-sector} gives
$0<\bar\nu(1)<\infty$ and identifies
$\bar\mu=\bar\nu/\bar\nu(1)$ with the zero-flux
curvature probability law.

Fix $C_0\ge1$, which will also dominate the terminal
field-map and Dirac-form constants in
Section~\ref{subsec:interacting-terminal-norms}. Set
\begin{equation}\label{eq:term-complement-errors}
 \varepsilon_3
 =6\ell_*^3\left(e^{-p_{0,*}/2}
                    +e^{-p_*^2/(2C_0^2)}\right),\qquad
 \varepsilon_4=6\ell_*^3e^{-p_*^2/(2C_0^2)}.
\end{equation}

\begin{proposition}\label{prop:term-complement}
Suppose $v_*\ge8$, the integer $p_0\ge1$ satisfies
$p_0v_*^{p_0-1}\ge28$, and $L\ge5$. Then, at every
real initial pair,
\begin{equation}\label{eq:term-complement-bound}
 \bar\mu(\mathfrak G^c)\le\varepsilon_3,\qquad
             \mathsf m_*(\chi_*=0)\le\varepsilon_4.
\end{equation}
For fixed $(\bar m,e,R)$, these bounds are independent
of the ultraviolet cutoff and of the initial pair.
\end{proposition}

\begin{proof}
The zero-flux domination estimate in
Proposition~\ref{prop:interacting-zero-flux-domination}
and the rows of \eqref{eq:term-tower-row-norms} give
\[
 \bar\mu\{|\zeta_j(\Upsilon)|>p_{0,j}\}
                   \le2e^{-p_{0,j}^2/2},\qquad
 \bar\mu\{|(dA_*)(p)|>p_*\}\le2e^{-p_*^2/2}.
\]
Using the coordinate counts in
\eqref{eq:term-tower-coordinate-count}, the union bound is
\begin{equation}\label{eq:term-complement-union}
 \bar\mu(\mathfrak G^c)
 \le4\ell_*^3\sum_{i=1}^{k}L^{3i}
                  e^{-(v_*+i\log L/2)^{2p_0}/2}
                +6\ell_*^3e^{-p_*^2/2}.
\end{equation}
Here $i=k-j$ and $v_j=v_*+i\log L/2$.
The first two terms of the binomial expansion give
\[
 p_{0,j}=(v_*+i\log L/2)^{p_0}
       \ge p_{0,*}+28i\log L/2.
\]
Since $p_{0,j}\ge1$, its square is at least $p_{0,j}$.
It follows that
\[
 L^{3i}e^{-p_{0,j}^2/2}
                  \le e^{-p_{0,*}/2}L^{-4i}.
\]
Summing the geometric series, whose sum is
$(L^4-1)^{-1}\le1$, bounds
\eqref{eq:term-complement-union} by
$4\ell_*^3e^{-p_{0,*}/2}+6\ell_*^3e^{-p_*^2/2}$.
This is at most $\varepsilon_3$, since $C_0\ge1$.

The image of $\gamma_N$ under $F\mapsto A_*$ is
$\mathsf m_*$ by \eqref{eq:term-tower-Gaussian}.
For each plaquette the variable
$\langle\phi_p,F\rangle$ is Gaussian with variance
$\|\phi_p\|^2\le1$. Its two-sided exponential Markov
bound is $2e^{-p_*^2/2}$.
Summing over the $3\ell_*^3$ plaquettes proves
$\mathsf m_*(\chi_*=0)\le6\ell_*^3e^{-p_*^2/2}
\le\varepsilon_4$.
Finally, $\ell_*$ and $e_*$ are independent of $N$
after the mesh threshold, by
Proposition~\ref{prop:rgfl-scales}; hence the displayed
errors have the asserted dependence.
\end{proof}

\subsection{Free pressure and twist dependence}
\label{subsec:term-free-determinants}

The partition comparison requires an error in the free pressure
which depends on the physical size, rather than on the number of
lattice sites. It also requires a twist estimate for the operator
with a block penalty. The latter penalty removes the small-mass
singularity before the fine mesh is taken to zero.

For $0<m\le1$ let
$\mathfrak d_m(p)=\det\mathcal D_m(p)$, where
$\mathcal D_m$ is the $192$-dimensional chamber symbol in
Section~\ref{subsec:schur-symbol}, and set
\begin{equation}\label{eq:term-free-pressure}
 p_{\rm ch}^0(m)=
       \int_{[-\pi,\pi]^3}\log\mathfrak d_m(p)
                                      \frac{\dd p}{(2\pi)^3}.
\end{equation}
The spin-block identity and positive real part imply
$\mathfrak d_m(p)>0$ for real $p$.
In this subsection $|z|=|z|_\infty$ for $z\in\Z^3$.
Write $\kappa_{(0)}=\kappa_{\mathrm{strip}}$ from
Proposition~\ref{prop:schur-symbol}.

\begin{proposition}\label{prop:term-free-pressure}
There is a finite $C_6$, depending only on $\rho$, such that,
with
\begin{equation}\label{eq:term-epsilon-five}
 \varepsilon_5(x)=
 C_6\log\frac1{1-e^{-\kappa_{(0)}x}}\qquad(x>0),
\end{equation}
the free twisted determinant on a unit torus of side $n\ge1$
satisfies
\begin{equation}\label{eq:term-free-pressure-error}
 \left|\log\det(D^{\rm ch}[qH_\theta]+m)
                     -n^3p_{\rm ch}^0(m)\right|
                                 \le\varepsilon_5(mn).
\end{equation}
Here $H_{\theta,\mu}=\theta_\mu/(qn)$ and $\theta\in\R^3$.
For $m=\bar m\eta$ and $n=R/\eta$ define
\[
 z_{r,n}^{\rm free}
       =\int_{[-\pi,\pi]^3}
            \det(D^{\rm ch}[qH_\theta]+\bar m\eta)
                                         \frac{\dd\theta}{(2\pi)^3}.
\]
Then $z_{r,n}^{\rm free}>0$ and
\begin{equation}\label{eq:term-free-average-error}
 \left|\log z_{r,n}^{\rm free}
                 -n^3p_{\rm ch}^0(\bar m\eta)\right|
                              \le\varepsilon_5(\bar mR).
\end{equation}
\end{proposition}

\begin{proof}
The plane waves with momenta
$P_{n,\theta}=\{(2\pi j+\pi\mathbf1+\theta)/n:
j\in\{0,\ldots,n-1\}^3\}$ diagonalize the translation
part of the chamber operator. Inserting a plane wave
in the fourth hop shifts its momentum by $\theta/n$,
whereas the internal hops are unchanged. Thus
\begin{equation}\label{eq:term-free-product}
 \det(D^{\rm ch}[qH_\theta]+m)
                   =\prod_{p\in P_{n,\theta}}\mathfrak d_m(p)>0.
\end{equation}
The phases contain $qH_\theta=\theta/n$, so the determinant
is independent of $q$.

Factor the determinant by the light--heavy decomposition:
\[
 \log\mathfrak d_m(p)=f_H(p)+f_S(p),\qquad
 f_H(p)=\log\det D_{\rm hh}(p),\quad
 f_S(p)=\log\det B_m(p).
\]
Both logarithms are real and smooth on real momenta.
Their derivatives, which suffice for contour shifting, are
\[
 \partial_\mu f_H=\tr(D_{\rm hh}^{-1}\partial_\mu D_{\rm hh}),
 \qquad
 \partial_\mu f_S=\tr(B_m^{-1}\partial_\mu B_m).
\]
The heavy block has dimension $188$. On the strip of
width $\theta_0$ its inverse has norm at most $2/g_{\rm h}$,
and the fourth-hop derivative has norm at most
$3(1+\rho)e^{\theta_0}=3(1+\rho)+g_{\rm h}/2$.
Hence
\[
 |\partial_\mu f_H|\le
 C_{\rm lH}:=188\left(1+\frac{6(1+\rho)}{g_{\rm h}}\right).
\]
For $z\ne0$ choose $|z_\mu|=|z|$. One integration by parts
and a shift of that contour by
$-\mathrm i\theta_0\operatorname{sgn}(z_\mu)$ give
\begin{equation}\label{eq:term-heavy-log-Fourier}
 |\widehat f_H(z)|\le C_{\rm lH}|z|^{-1}e^{-\theta_0|z|}.
\end{equation}
The derivative trace is holomorphic in a neighborhood of
the closed strip. Its periodicity cancels the vertical sides
of the contour rectangle; no choice of a complex logarithm
is needed.

For the light factor put $r_m(p')=(|p'|_2^2+m^2)^{1/2}$.
The inverse derivative bounds
\eqref{eq:schur-inverse-derivatives} hold on
$|\operatorname{Im}p_\nu|\le\kappa_{(0)}m$.
The bound $K_s$ for $B_m$ on the strip of width $\theta_0$
and Cauchy's formula on polydiscs of radius $\theta_0/2$
give
$\|\partial^\beta B_m\|
\le24K_s(2/\theta_0)^{|\beta|}$ for $1\le|\beta|\le4$.
Leibniz's rule in the derivative trace therefore gives
\begin{equation}\label{eq:term-light-log-derivatives}
 |\partial^\alpha\partial_\mu f_S(p'+\mathrm ip'')|
 \le C_{\rm lS}\,r_m(p')^{-1-|\alpha|}
 \le C_{\rm lS}|p'|_2^{-1-|\alpha|}
       \quad(|\alpha|\le3,\ p'\ne0),
\end{equation}
where one may take
\[
 C_{\rm lS}
 =\max_{|\alpha|\le3}96K_s
   \sum_{\beta\le\alpha}
    \binom{\alpha}{\beta}C_\beta(\rho)
    (2/\theta_0)^{|\alpha|-|\beta|+1}
    (3\pi^2+1)^{(|\alpha|-|\beta|)/2}.
\]
This constant is independent of $m$.
Shift the differentiated Fourier integral by
$-\mathrm i\kappa_{(0)}m\operatorname{sgn}(z_\mu)e_\mu$.
The periodic Fourier estimate proved in
Proposition~\ref{prop:free-chamber-kernel}, with singular
order $1$ and three derivatives, bounds this integral by
$c_{\rm F}C_{\rm lS}(1+|z|)^{-2}$, where $c_{\rm F}$
is its constant at those orders. Thus
\begin{equation}\label{eq:term-light-log-Fourier}
 |\widehat f_S(z)|
    \le c_{\rm F}C_{\rm lS}|z|^{-1}(1+|z|)^{-2}
                          e^{-\kappa_{(0)}m|z|}.
\end{equation}
The strip is contained in the heavy strip since
$\kappa_{(0)}m\le\theta_0/2$.
The inverse is holomorphic near the closed strip, so
the contour shift is legitimate.

The Fourier coefficients of $f=f_H+f_S$ are summable.
For a summable periodic Fourier series, summing over
$P_{n,\theta}$ keeps precisely the indices $z=nw$:
\begin{equation}\label{eq:term-Fourier-sampling}
 \sum_{p\in P_{n,\theta}}f(p)
 =n^3\sum_{w\in\Z^3}\widehat f(nw)
                       (-1)^{w_1+w_2+w_3}e^{\mathrm i\theta\cdot w}.
\end{equation}
This follows by summing the geometric progression in each
coordinate; absolute convergence permits the interchange.
Its $w=0$ term is $n^3p_{\rm ch}^0(m)$.
There are $24t^2+2\le26t^2$ indices with $|w|=t$.
The contribution of \eqref{eq:term-light-log-Fourier} is at most
\[
 26c_{\rm F}C_{\rm lS}
                  \sum_{t\ge1}t^{-1}e^{-\kappa_{(0)}mnt}.
\]
For the heavy term use
$\kappa_{(0)}m\le\theta_0/2$ and
$x^2e^{-\theta_0x/2}\le16/(\exp(1)\theta_0)^2$.
Its contribution is at most
\[
 \frac{416C_{\rm lH}}{(\exp(1)\theta_0)^2}
                     \sum_{t\ge1}t^{-1}e^{-\kappa_{(0)}mnt}.
\]
Since $\sum_{t\ge1}e^{-xt}/t=\log(1-e^{-x})^{-1}$,
\eqref{eq:term-free-pressure-error} follows with
$C_6=26[c_{\rm F}C_{\rm lS}
          +16C_{\rm lH}/(\exp(1)\theta_0)^2]$.
Finally every twisted determinant lies between
$\exp(n^3p_{\rm ch}^0(m)\pm\varepsilon_5(mn))$.
Its harmonic average lies in the same positive interval.
Taking logarithms proves \eqref{eq:term-free-average-error}.
\end{proof}

We next retain the block penalty. Let $\mathfrak N\ge3$
be odd and put $\tau=\mathfrak N^{-1}$.
On the fine torus of side $\ell\in\N$ let
$\mathcal H_\theta$ be the chamber spinors with mesh $\tau$
and boundary multipliers $-e^{\mathrm i\theta_\mu}$.
The scalar product uses $\tau^3/48$.
The coarse space has four components per unit site and
counting scalar product. The average $Q$ takes the mean
over each unit block and its $48$ labels, and $Q^*$
extends a coarse field constantly on that block.
Thus $QQ^*=1$, $\|Q\|=\|Q^*\|=1$, and
$\Pi_{\rm bl}=Q^*Q$ is an orthogonal projection.
For $0<\mu\le1$ and $b_-\le b\le b_+$ set
\begin{equation}\label{eq:term-free-block-operators}
 T_\theta=D+\mu+b\Pi_{\rm bl},\qquad
 D_\theta^{\rm bl}=b-b^2QT_\theta^{-1}Q^*.
\end{equation}
All determinants below are determinants of linear maps.
The weighting of a scalar product changes their matrix
representations, but does not add a determinant factor.

\begin{proposition}\label{prop:term-block-twist}
The determinants of $T_\theta$, $D+\mu$, and
$D_\theta^{\rm bl}$ are positive, and
\begin{align}
 (D_\theta^{\rm bl})^{-1}
     &=b^{-1}+Q(D+\mu)^{-1}Q^*,
                                    \label{eq:term-free-block-inverse}\\
 \det D_\theta^{\rm bl}
     &=b^{4\ell^3}\frac{\det(D+\mu)}{\det T_\theta}.
                                    \label{eq:term-free-block-determinant}
\end{align}
There are $C_\Sigma<\infty$ and $0<\gamma_\Sigma\le1$,
depending only on $(b_-,b_+,\rho)$, such that
\begin{equation}\label{eq:term-penalty-twist-bound}
 \left|\log\frac{\det T_\theta}{\det T_0}\right|
                  \le C_\Sigma\ell^{-1}e^{-\gamma_\Sigma\ell}.
\end{equation}
For the level $k\ge1$ data
\[
             (\mathfrak N,\ell,\mu,b)=(L^k,\ell_*,\mu_*,b_k),
             \qquad c_{\theta,\mu}=\theta_\mu/(e_*\ell_*),
\]
put
$R_\theta^{\rm tw}=\det D_k(c_\theta)/\det D_k(0)$.
Then
\begin{equation}\label{eq:term-twist-determinant-identity}
 R_\theta^{\rm tw}
 =\frac{\det(D^{\rm ch}[qH_\theta]+\bar m\eta)}
              {\det(D^{\rm ch}[0]+\bar m\eta)}
                  \frac{\det T_0}{\det T_\theta}>0.
\end{equation}
\end{proposition}

\begin{proof}
The real part of $D+\mu$ is at least $\mu$.
The same holds for $T_\theta$, since its penalty is positive.
Both operators have two adjoint spin blocks; their determinants
are therefore positive. Put $G=(D+\mu)^{-1}$,
$S=QGQ^*$ and $Z=(b^{-1}+S)^{-1}$.
Their inverses exist since $\operatorname{Re}G\ge0$ and
$\operatorname{Re}(b^{-1}+S)\ge b^{-1}$.
Multiplication gives
$T_\theta^{-1}=G-GQ^*ZQG$ and
$D_\theta^{\rm bl}=Z$, proving
\eqref{eq:term-free-block-inverse}.
Sylvester's determinant identity, applied to
$-\!bQ$ and $T_\theta^{-1}Q^*$, gives
\eqref{eq:term-free-block-determinant}. Its right side is
positive.

We prove \eqref{eq:term-penalty-twist-bound} with constants
independent of $\tau$ and $\mu$.
Let
\[
 c_{\rm nz}=\min\left\{\frac\rho\pi,
             \frac{g_{\rm h}\rho^2}{49\sqrt3\,\pi^3}\right\}>0.
\]
For real $q'\in[-\pi,\pi]^3$,
\begin{equation}\label{eq:term-symbol-nozero}
             \|\mathcal D_0(q')u\|
                         \ge c_{\rm nz}|q'|_2\|u\|.
\end{equation}
Normalize $\|u\|=1$ and
write $u=Pu+\Pi_{\rm h}u$ for its light and heavy parts,
where $\Pi_{\rm h}=1-P$ is the cluster-heavy projection.
If $d=\|\mathcal D_0(q')u\|$, then
$g_{\rm h}\|\Pi_{\rm h}u\|^2\le d$.
The light compression is
$W(q')=\rho\sum_\mu(1-\cos q'_\mu)+
\mathrm i\gamma\cdot\sin q'$; its squared modulus is
at least $(2\rho/\pi)^2|q'|_2^2$.
Also $J^*\mathcal D_0(0)=0$, and the fourth-hop
difference has norm at most $6|q'|_2$.
Projecting $\mathcal D_0(q')u$ onto the light space gives
\[
 d\ge |q'|_2\left(\frac{2\rho}{\pi}-7\|\Pi_{\rm h}u\|\right).
\]
If $\|\Pi_{\rm h}u\|\le\rho/(7\pi)$ this is at least
$(\rho/\pi)|q'|_2$. Otherwise
$d\ge g_{\rm h}\rho^2/(49\pi^2)$; since
$|q'|_2\le\sqrt3\pi$, this proves
\eqref{eq:term-symbol-nozero}.

On one unit cell use the Bloch space
$\mathcal H(p)=\{f(c+e_\mu,g)=e^{\mathrm ip_\mu}f(c,g)\}$.
Its modes are $p_j=p+2\pi j$, with
$j\in I_{\mathfrak N}^3$,
$I_{\mathfrak N}=\{-(\mathfrak N-1)/2,\ldots,
(\mathfrak N-1)/2\}$.
The map
$\{u_j\}\mapsto\sum_j e^{\mathrm ip_j\cdot c}u_j(g)$
is unitary for the one-cell scalar product $\tau^3/48$.
It diagonalizes $D$ with blocks
$X_j=\tau^{-1}\mathcal D_0(\tau p_j)$.
Define
\[
 u_1(t)=\mathfrak N^{-1}\sum_{s\in I_{\mathfrak N}}
                                      e^{-\mathrm its/\mathfrak N},
 \qquad u(p)=\prod_\mu u_1(p_\mu).
\]
The block projection in this representation has kernel
$u(p_j)u(p_{j'})P$. This follows by averaging the plane
waves on the centered unit block; discrete Fourier
orthogonality expands the constant block extension.
For $p\in[-\pi,\pi]^3$ one has $\tau p_j\in[-\pi,\pi]^3$
and $|p_j|_2\ge\pi|j|$ when $j\ne0$.
Hence \eqref{eq:term-symbol-nozero} gives
$\|X_ju\|\ge c_{\rm nz}|p_j|_2\|u\|$.

The inverse on $\mathcal H(p)$ is uniformly bounded.
To see this, let $\|f\|=1$ and
$d=\|(D+\mu+b\Pi_{\rm bl})f\|\le1$.
The real-part estimate gives
\[
 \mu\le d,\qquad
 \|\Pi_{\rm bl}f\|\le(d/b_-)^{1/2},\qquad
 \|\Pi_{\rm h}f\|\le(d/(3g_{\rm h}))^{1/2},
 \qquad \|Df\|\le(2+\sqrt{b_+})\sqrt d.
\]
Separate the mode $j=0$ from $f^{\rm r}$, the remaining
modes. The bound on $X_j$ gives
$\|f^{\rm r}\|\le\|Df\|/(\pi c_{\rm nz})$.
Moreover
$u_1(t)=\sin(t/2)/[\mathfrak N\sin(t/(2\mathfrak N))]
\ge2/\pi$ for $|t|\le\pi$, with value $1$ at zero.
Thus the block projection of the zero mode controls
its light part by the factor $8/\pi^3$.
Using $\|\Pi_{\rm bl}\|=1$ and subtracting the other modes,
\[
 \sqrt{d/b_-}\ge
 \frac8{\pi^3}(1-\|\Pi_{\rm h}f\|)
                     -(1+8/\pi^3)\|f^{\rm r}\|.
\]
It follows that $8/\pi^3\le K_{\rm iv}\sqrt d$, where
\[
 K_{\rm iv}=b_-^{-1/2}+\frac8{\pi^3}(3g_{\rm h})^{-1/2}
       +(1+8/\pi^3)\frac{2+\sqrt{b_+}}{\pi c_{\rm nz}}.
\]
Consequently the inverse norm is at most
$C_{\rm iv}=\max\{1,(\pi^3K_{\rm iv}/8)^2\}$.

For complex momenta use the periodic gauge
$f(c,g)\mapsto e^{\mathrm ip\cdot c}f(c,g)$ to work
on the fixed periodic one-cell space. Let $D(p)$
and $\Pi(p)$ be the resulting operators, so
$\mathcal T(p)=D(p)+\mu+b\Pi(p)$.
The fourth hopping phase is
$e^{\mathrm i\tau p\cdot ge_1}$, and the block kernel is
$e^{-\mathrm ip\cdot(c-c')}$. A shift of $p_\nu$
by $2\pi$ conjugates these operators by multiplication
by $e^{2\pi\mathrm ic_\nu}$.
Thus their determinant and derivative traces are periodic.
For $p=p'+\mathrm it e_\nu$, $|t|\le1$, and $1\le j\le4$,
direct differentiation gives
\begin{equation}\label{eq:term-Bloch-derivatives}
 \|\partial_\nu^jD(p)\|\le9\tau^{j-1},\qquad
 \|\partial_\nu^j\Pi(p)\|\le\exp(1),\qquad
 \operatorname{rank}\partial_\nu^j\Pi(p)\le4(j+1).
\end{equation}
The last bound follows by expanding $(c_\nu-c'_\nu)^j$
into $j+1$ products of a function of $c$ and a function
of $c'$, each with four spin components.
The norm estimates use $|c_\nu-c'_\nu|<1$ in one block,
$\tau\le1/3$, and $\|M_3\|\le6$.
The difference from the real momentum has norm at most
$(6+b_+)|t|e^{|t|}$. With
\[
 \gamma_\Sigma=
 \min\{1,[2\exp(1)(6+b_+)C_{\rm iv}]^{-1}\},
\]
a Neumann series gives
$\|\mathcal T(p)^{-1}\|\le2C_{\rm iv}$ on this strip.
Its singular values are at most twice those at $p'$.

The growing dimension of the fine cell cannot be used
in a trace bound. Instead write $G(p)=\mathcal T(p)^{-1}$
and compare at real $p'$ with $(D(p')+1)^{-1}$.
The resolvent identity, the bound just proved, and the
singular-value multiplication inequality
\cite[slides~3--9]{PongeSchatten2025} imply
\[
 s_j(G(p))\le C_{\rm sv}s_j((D(p')+1)^{-1}),\qquad
 C_{\rm sv}=2[1+(1+b_+)C_{\rm iv}].
\]
The comparison operator is the direct sum of
$(X_j+1)^{-1}$. Since $\operatorname{Re}X_j\ge0$,
each of its $192$ singular values is at most
$(1+c_{\rm nz}^2|p'_j|_2^2)^{-1/2}$.
Set
\[
 S_4=\sum_{j\in\Z^3\setminus\{0\}}|j|^{-4}
        =4\pi^2+\pi^4/45,\quad
 C_{\rm uv}=(1+3\pi^2c_{\rm nz}^2)^{1/2},\quad
 S_{\rm jf}=1+S_4/(\pi c_{\rm nz})^4.
\]
For $1\le s\le4$ we obtain
\begin{equation}\label{eq:term-Bloch-Schatten}
 \tau^{4-s}\|G(p)\|_{\mathfrak S_s}^{\,s}
           \le192C_{\rm sv}^{\,s}C_{\rm uv}^{\,4-s}S_{\rm jf}.
\end{equation}
Indeed
$\tau^2(1+c_{\rm nz}^2|p'_j|_2^2)\le C_{\rm uv}^2$,
so multiplication by $\tau^{4-s}$ makes each summand
a multiple of $(1+c_{\rm nz}^2|p'_j|_2^2)^{-2}$.
The nonzero indices are bounded by
$(\pi c_{\rm nz}|j|)^{-4}$ and have the summable
majorant displayed above.

Let $F(p)=\log\det\mathcal T(p)$ on real momenta.
The fourth derivative is the trace expression
\begin{align}
 F_4(p)={}&\tr(GT_4)-4\tr(GT_1GT_3)
                  -3\tr(GT_2GT_2)\nonumber\\
          &+12\tr(GT_1GT_1GT_2)-6\tr((GT_1)^4),
                                      \label{eq:term-Bloch-fourth}
\end{align}
where $T_j=\partial_\nu^j\mathcal T$.
This follows by differentiating
$\partial_\nu F=\tr(GT_1)$ and
$\partial_\nu G=-GT_1G$, using trace cyclicity.
The expression is holomorphic near the closed strip
in the $\nu$ coordinate.
Each trace in \eqref{eq:term-Bloch-fourth} has $s\le4$
inverse factors and derivatives of total order four.
Expand every $T_j$ as $D_j+b\Pi_j$.
For the term with only $D_j$, Schatten Hölder and
\eqref{eq:term-Bloch-derivatives} give the bound
$9^s\tau^{4-s}\|G\|_{\mathfrak S_s}^{\,s}$.
Equation~\eqref{eq:term-Bloch-Schatten} bounds this by
$9^4\,192C_{\rm sv}^4C_{\rm uv}^3S_{\rm jf}$.
In each other term estimate one $b\Pi_j$ in trace norm:
its rank bound gives $\|b\Pi_j\|_{\mathfrak S_1}
\le20\exp(1)b_+$. The other factors have operator
norms at most $2C_{\rm iv}$ and
$\max\{9,\exp(1)b_+\}$.
There are at most $15$ such terms in any trace.
Since the five numerical coefficients have absolute
sum $26$, this proves $|F_4(p)|\le C_{\rm Flg}$, where
\[
 C_{\rm Flg}=26\left[
       9^4\,192C_{\rm sv}^4C_{\rm uv}^3S_{\rm jf}
       +300\exp(1)b_+(2C_{\rm iv})^4
                           \max\{9,\exp(1)b_+\}^3\right].
\]
All Schatten applications are on a finite-dimensional
space, with $s$ factors of exponent $s$ and reciprocal
exponents summing to one
\cite[slides~54--56]{PongeSchatten2025}.

For $z\ne0$ choose $|z_\nu|=|z|$.
Four integrations by parts and a contour shift by
$-\mathrm i\gamma_\Sigma\operatorname{sgn}(z_\nu)e_\nu$
now give
$|\widehat F(z)|\le C_{\rm Flg}|z|^{-4}
e^{-\gamma_\Sigma|z|}$.
The trace expression is periodic by the preceding
conjugacy and is holomorphic near the closed strip.
Decomposition into unit-cell Bloch spaces gives
$\log\det T_\theta=\sum_{p\in P_{\ell,\theta}}F(p)$.
Using \eqref{eq:term-Fourier-sampling} at $\theta$ and
at zero, and subtracting, proves
\[
 \left|\log\frac{\det T_\theta}{\det T_0}\right|
 \le2\ell^3\sum_{w\ne0}C_{\rm Flg}(\ell|w|)^{-4}
                                    e^{-\gamma_\Sigma\ell|w|}
 \le2S_4C_{\rm Flg}\ell^{-1}e^{-\gamma_\Sigma\ell}.
\]
Thus $C_\Sigma=2S_4C_{\rm Flg}$.

Finally constant connections are twists by multiplication
by $e^{\mathrm i\theta\cdot c/\ell_*}$.
Every rectilinear path contributing to the average has
the same constant-connection integral, so its phase is
$e^{\mathrm i\theta\cdot(c-y)/\ell_*}$.
Consequently the covariant chamber averages and fine
inverse are conjugate to $Q,Q^*$ and $T_\theta^{-1}$,
and $D_k(c_\theta)$ is conjugate to $D_\theta^{\rm bl}$.
Dilating the mesh $\tau$ to one gives
\[
 \det(D+\mu\text{ on }\mathcal H_\theta)
          =\tau^{-192n^3}\det(D^{\rm ch}[qH_\theta]+\tau\mu),
                \qquad n=\ell_*\mathfrak N .
\]
Here $\tau\mu_*=\bar m\eta$.
Divide \eqref{eq:term-free-block-determinant} at the two
twists. The factors $b^{4\ell_*^3}$ and
$\tau^{-192n^3}$ cancel, proving
\eqref{eq:term-twist-determinant-identity}.
\end{proof}

\subsection{The terminal remainder}
\label{subsec:term-remainder}

Choose $C_0$ to bound the photon row sums, the field maps
and the Dirac form in
Section~\ref{subsec:interacting-terminal-norms}.
Set
\[
 c_\kappa=\sup_\square\sum_{X\supset\square}
                                      e^{-\kappa d_M(X)}\le K_0.
\]
In addition to the primitive and flow hypotheses, impose
the following inequalities at the stop:
\begin{align}
 8C_0^2\ell_*^6(b_-^{-1}+\mu_*^{-1})
                   &\le e_*^{-\varepsilon/2},\nonumber\\
 4(2C_0+c_\kappa+5c_{\rm fl}+6)\ell_*^3
                   &\le e_*^{-\varepsilon},\nonumber\\
 e_*^{1/4-9\varepsilon}&\le\tfrac18,\qquad
 4C_0\ell_*^2p_*<e_*^{-\varepsilon},\nonumber\\
 12\ell_*^3\left(e^{-p_{0,*}/2}
                         +e^{-p_*^2/(2C_0^2)}\right)&\le e_*,
 \qquad
 C_\Sigma e^{-\gamma_\Sigma\ell_*}\le e_*.
                                      \label{eq:term-stop-conditions}
\end{align}
These are the terminal conditions used in the proof.
The choice in Proposition~\ref{prop:ordered-constants}
is required to satisfy them: the scale estimates give
$Mv_*^r\le\ell_*<LMv_*^r$ and
$\mu_*^{-1}<L\ell_*$. Substitute the latter two bounds
in the first four inequalities; use the lower bound
in the last exponential inequality.
Thus the conditions are uniform in the cutoff.

Put
\begin{equation}\label{eq:term-remainder-errors}
 \varepsilon_1=\tfrac12e_*^{1/4-9\varepsilon},\qquad
 \varepsilon_2=C_\Sigma\ell_*^{-1}e^{-\gamma_\Sigma\ell_*}.
\end{equation}
Together with \eqref{eq:term-complement-errors},
\eqref{eq:term-stop-conditions} gives
$\varepsilon_1\le1/16$, $\varepsilon_2\le e_*$,
and $\varepsilon_4\le\varepsilon_3\le e_*/2$.

\begin{proposition}\label{prop:term-remainder}
For a controlled real pair, every mean-zero real coarse
field $A$ with $\|dA\|_\infty\le p_*$ and every
$\theta\in\R^3$ satisfy
\begin{equation}\label{eq:term-density-remainder}
 Y_*(A+c_\theta)
     =e^{-\varepsilon_k^{\rm v}\ell_*^3}
                R_\theta^{\rm tw}(1+r_*(\theta,A)),
                         \qquad |r_*(\theta,A)|\le\varepsilon_1.
\end{equation}
The remainder is continuous in $(\theta,A)$ on this set.
\end{proposition}

\begin{proof}
The gauge construction in
\eqref{eq:photon-inner-basepoint-bound}--\eqref{eq:photon-seam-basepoint-bound}
gives a real $\lambda$ such that
$A^g=A-d\lambda$ satisfies $\|A^g\|_\infty\le2\ell_*^2p_*$.
The curvature bound puts both $A+c_\theta$ and
$A^g+c_\theta$ in $\mathrm{Dom}_k$, by
Proposition~\ref{prop:photon-domain-preservation};
Theorem~\ref{thm:term-density} gives equality of their
terminal densities.
Set $\mathscr A^0=H_k^{\rm ph}A^g$.
Summing each of the three photon kernel bounds gives
a norm at most $C_0\|A^g\|_\infty$; hence
\[
 \|\mathscr A^0\|_{\rm each}
                    \le2C_0\ell_*^2p_*<\tfrac12e_*^{-\varepsilon}.
\]
The photon minimizer preserves constants, so the
background is $\mathscr A=\mathscr A^0+c_\theta$.
For $\rho_*=e_*^{-3/4+4\varepsilon}\ge8$ and
$|t|<2\rho_*$, each norm of $t\mathscr A^0$ is
less than $e_*^{-3/4+3\varepsilon}$, the smallest
analytic-domain threshold. Invariance under real
constant shifts gives
\begin{equation}\label{eq:term-terminal-disc}
 \mathscr A^{(t)}=t\mathscr A^0+c_\theta
                  \in\widetilde R_k\subset\widetilde G_k
                         \qquad(|t|<2\rho_*).
\end{equation}

Let $D^{\rm tw}=D_k(c_\theta)$ and
$C_\theta=(D^{\rm tw})^{-1}$.
The free block formula and its twist conjugacy give
$\|C_\theta\|_{\rm op}\le b_-^{-1}+\mu_*^{-1}$.
At the connection $\mathscr A$, after the fermion
field substitution, put
\begin{align}
 V_*={}&-\langle\bar\Psi,
          [D_k(\mathscr A)-D_k(c_\theta)]\Psi\rangle
               -m_k\langle\bar\psi,\psi\rangle\nonumber\\
       &+\sum_XE_k(X;\mathscr A,\psi,\chi).
                                      \label{eq:term-terminal-interaction}
\end{align}
All terms are even. Completing the Grassmann Gaussian
in \eqref{eq:term-Y-density}, at zero external slot, gives
\begin{equation}\label{eq:term-terminal-Gaussian}
 Y_*(A+c_\theta)=e^{-\varepsilon_k^{\rm v}\ell_*^3}
             R_\theta^{\rm tw}\int\dd\mu_{C_\theta}(\Psi)e^{V_*}.
\end{equation}
This is the weight-one identity of the terminal coarse
space; the fine fields retain their $\eta_k^3/48$
pairing in $V_*$.

We estimate the three terms before integrating.
Let $a_1=e_*^{-1/4}$ and
$a_2=C_0^{-1}e_*^{-1/4+\varepsilon}\le a_1$.
The analytic quadratic form
$f(t)=\langle\bar\Psi,D_k(\mathscr A^{(t)})\Psi\rangle$
has $|f(t)|_{a_1}\le C_0e_*^{-1/2}\ell_*^3$ on
\eqref{eq:term-terminal-disc}, by
\eqref{eq:interacting-terminal-Dirac-form}.
Coefficientwise Cauchy's formula with kernel
$[t(t-1)]^{-1}$ on $|t|=\rho_*$ gives
\[
 |f(1)-f(0)|_{a_2}
 \le |f(1)-f(0)|_{a_1}
 \le\frac{C_0e_*^{-1/2}\ell_*^3}{\rho_*-1}
                  \le2C_0e_*^{1/4-4\varepsilon}\ell_*^3 .
\]
The mass-form and activity-substitution bounds
\eqref{eq:interacting-terminal-mass-form} and
\eqref{eq:interacting-terminal-substitution} give
\[
 |m_k\langle\bar\psi,\psi\rangle|_{a_2}
                \le4e_*^{1/4-8\varepsilon}\ell_*^3,\qquad
 |E_k(X;\mathscr A,\psi,\chi)|_{a_2}
                 \le e^{-\kappa d_M(X)}w_k^{\rm E}.
\]
Here $|m_k|\le e_*^{3/4-8\varepsilon}$ and
$w_k^{\rm E}=e_*^{1/4-7\varepsilon}$.
These norm estimates include the ordered pairs within
one chamber cluster. The mass pairing uses the fine site measure
$\eta_k^3/48$;
the four spin components in the mass form therefore
contribute $4\ell_*^3$, rather than $192\ell_*^3$.
Every polymer contains an $M$-cube, so
$\sum_Xe^{-\kappa d_M(X)}
\le(\ell_*/M)^3c_\kappa\le c_\kappa\ell_*^3$.
Consequently
\begin{equation}\label{eq:term-terminal-interaction-bound}
 |V_*|_{a_2}\le
            (2C_0+c_\kappa+4)e_*^{1/4-8\varepsilon}\ell_*^3.
\end{equation}

The covariance integral is the unit-covariance integral
after the one-sided substitution
$\Psi\mapsto C_\theta\Psi$, $\bar\Psi\mapsto\bar\Psi$.
Each covariance row has $4\ell_*^3$ entries of modulus
at most $b_-^{-1}+\mu_*^{-1}$.
Its row norm, together with the identity on barred
variables, is bounded by
\[
 4\ell_*^3(b_-^{-1}+\mu_*^{-1})
 \le\frac{e_*^{-\varepsilon/2}}{2C_0^2\ell_*^3}
 =a_2\,\frac{e_*^{1/4-3\varepsilon/2}}{2C_0\ell_*^3}
 \le a_2.
\]
The weighted substitution estimate therefore sends
$V_*$ to $V_*'$ with $|V_*'|_1\le|V_*|_{a_2}$.
The second inequality of \eqref{eq:term-stop-conditions}
and \eqref{eq:term-terminal-interaction-bound} imply
$|V_*'|_1\le e_*^{1/4-9\varepsilon}/4\le1/32$.
The unit-covariance Wick integral is bounded by the
weight-one Grassmann norm and integrates $1$ to $1$.
Thus
\[
 \left|\int\dd\mu_{C_\theta}e^{V_*}-1\right|
 \le |e^{V_*'}-1|_1
 \le e^{|V_*'|_1}-1
 \le2|V_*'|_1\le\varepsilon_1.
\]
Together with \eqref{eq:term-terminal-Gaussian}, this
proves \eqref{eq:term-density-remainder}.
The density is jointly continuous by
Theorem~\ref{thm:term-density}. The twist determinant
is a continuous positive function, so its quotient
in \eqref{eq:term-density-remainder} is continuous
on the asserted set.
\end{proof}

\subsection{Partition bounds and comparison of torus sizes}
\label{subsec:term-partition-bounds}

Let $z_{r,n}$ be the matter factor in
\eqref{eq:external-matter-factor}, formed at the real mass
$\bar m\eta+m_{\mathrm{in}}$. The fine normalization is
\begin{equation}\label{eq:term-matter-normalization}
 z_{r,n}
 =e^{\varepsilon_{\mathrm{in}}^{\rm v}n^3}
       \det(D^{\rm ch}[0]+\bar m\eta)\,\bar\nu(1)>0.
\end{equation}
Both signs of the total lattice mass are allowed here.
The positive free reference mass $\bar m\eta$ is used
only in the determinant normalization.

\begin{proposition}\label{prop:term-partition-bounds}
For every controlled real pair on $R$, the following
two-sided estimates hold:
\begin{align}
 e^{-\varepsilon_k^{\rm v}\ell_*^3}
 R_\theta^{\rm tw}(1-\varepsilon_1)(1-\varepsilon_4)
 &\le\nu^\theta(\mathfrak G)
 \le e^{-\varepsilon_k^{\rm v}\ell_*^3}
 R_\theta^{\rm tw}(1+\varepsilon_1),
                                   \label{eq:term-small-field-mass}\\
 (1-\varepsilon_3)\bar\nu(1)
 &\le\bar\nu(\mathfrak G)\le\bar\nu(1).
                                   \label{eq:term-full-mass}
\end{align}
In particular
\begin{align}
 \left|\log z_{r,n}-\log z_{r,n}^{\rm free}
       -\varepsilon_{\mathrm{in}}^{\rm v}n^3
                   +\varepsilon_k^{\rm v}\ell_*^3\right|
 &\le2\varepsilon_1+\varepsilon_2
                              +2\varepsilon_3+2\varepsilon_4
 \le2e_*^{1/4-9\varepsilon},\label{eq:term-partition-free-bound}\\
 \left|\log z_{r,n}
       -[p_{\rm ch}^0(\bar m\eta)
                     +\varepsilon_{\mathrm{in}}^{\rm v}]n^3
                    +\varepsilon_k^{\rm v}\ell_*^3\right|
 &\le\varepsilon_5(\bar mR)+2e_*^{1/4-9\varepsilon}.
                                   \label{eq:term-partition-pressure-bound}
\end{align}
\end{proposition}

\begin{proof}
Apply \eqref{eq:term-Borel-observable} at $a=0$ and
$\Phi=1$, then insert \eqref{eq:term-density-remainder}.
The result is
\[
 \nu^\theta(\mathfrak G)=
 e^{-\varepsilon_k^{\rm v}\ell_*^3}R_\theta^{\rm tw}
       \int\chi_*(A)(1+r_*(\theta,A))\,\dd\mathsf m_*(A).
\]
The left side and the prefactor are real, and the
prefactor is positive. The last integral therefore
equals its real part.
The remainder bound and
$1-\varepsilon_4\le\mathsf m_*(\chi_*)\le1$
give \eqref{eq:term-small-field-mass}.
The first inequality of \eqref{eq:term-complement-bound}
gives
$\bar\nu(\mathfrak G^c)\le\varepsilon_3\bar\nu(1)$.
This and positivity prove \eqref{eq:term-full-mass}.

Multiply \eqref{eq:term-small-field-mass} by
$\det(D^{\rm ch}[0]+\bar m\eta)>0$.
The determinant identity
\eqref{eq:term-twist-determinant-identity} and the
penalty twist bound replace this product by the fine
twisted free determinant, with factors between
$e^{-\varepsilon_2}$ and $e^{\varepsilon_2}$.
Integrating in $\theta$ gives
\begin{align*}
 &e^{-\varepsilon_k^{\rm v}\ell_*^3}z_{r,n}^{\rm free}
       e^{-\varepsilon_2}(1-\varepsilon_1)(1-\varepsilon_4)
 \\
 &\hspace{15mm}\le
     \det(D^{\rm ch}[0]+\bar m\eta)\bar\nu(\mathfrak G)
 \le e^{-\varepsilon_k^{\rm v}\ell_*^3}z_{r,n}^{\rm free}
                              e^{\varepsilon_2}(1+\varepsilon_1).
\end{align*}
All integrands are continuous in $\theta$, and the
factors being compared are positive.
Use \eqref{eq:term-full-mass} and
\eqref{eq:term-matter-normalization}, then take logarithms.
Writing the expression on the left of
\eqref{eq:term-partition-free-bound} without absolute
values as $Q_n$, we obtain
\[
 -\varepsilon_2+\log(1-\varepsilon_1)+\log(1-\varepsilon_4)
 \le Q_n\le
 \varepsilon_2+\log(1+\varepsilon_1)-\log(1-\varepsilon_3).
\]
The errors lie in $[0,1/2]$. The inequalities
$-\log(1-x)\le2x$ and $\log(1+x)\le x$ therefore give
the first bound in \eqref{eq:term-partition-free-bound}.
Its right side is at most
$e_*^{1/4-9\varepsilon}+3e_*$.
Since $e_*^{1/4}\le1/8$, we have
$3e_*\le e_*^{1/4}\le e_*^{1/4-9\varepsilon}$,
which gives the second bound.
Adding \eqref{eq:term-free-average-error} proves
\eqref{eq:term-partition-pressure-bound}.
\end{proof}

We now use the pair $P_N(e)$ tuned on the top torus
in Theorem~\ref{thm:rgfl-common-pair}.
Define the common pressure
\begin{equation}\label{eq:term-common-pressure}
 p_F(N,e,\bar m)=
       p_{\rm ch}^0(\bar m\eta)+\varepsilon_{0,N}^{\rm v}(e).
\end{equation}
The pair and this pressure are independent of $R$
within the fixed torus scheme.

\begin{theorem}\label{thm:term-common-partition}
For every $R\in\mathcal T(e)$ and
$N\ge N_{\min}(R,e)$, the tuned matter factor satisfies
\begin{align}
 \left|\log z_{r,n}-\log z_{r,n}^{\rm free}
                      -\varepsilon_{0,N}^{\rm v}(e)n^3\right|
          &\le3e_*^{1/4-9\varepsilon},\nonumber\\
 \left|\log z_{r,n}-p_Fn^3\right|
          &\le\varepsilon_P(R)
             :=\varepsilon_5(\bar mR)+3e_*^{1/4-9\varepsilon}
          \le\varepsilon_5(1/L)+3e_K^{1/4-9\varepsilon}.
                                     \label{eq:term-common-pressure-bound}
\end{align}
In particular let $R_*=R_{\mathrm{top}}/L$,
$n_*=R_*/\eta$, and $N\ge N_{\rm ref}(e)$.
The two tori $R_*$ and $LR_*$ use the same mass and
the same pressure. Their matter comparison satisfies
\begin{align}
 |D_{\rm mat}|
 &:=\left|L^3\log z_{r,n_*}-\log z_{r,Ln_*}\right|
       \le L^3\varepsilon_P(R_*)+\varepsilon_P(LR_*)
       \nonumber\\
 &\le (L^3+1)
          \left[\varepsilon_5(\bar mR_*)
                                      +3e_K^{1/4-9\varepsilon}\right]
       \le C_{\rm mat},\nonumber\\
 C_{\rm mat}&=(L^3+1)\bigl(\varepsilon_5(1/L)+1\bigr).
                                     \label{eq:term-matter-comparison}
\end{align}
The displayed bounds are uniform in $N$; $C_{\rm mat}$
depends only on $L,\rho$ and the fixed chamber weights.
\end{theorem}

\begin{proof}
The common pair is controlled on every torus in the
asserted range, by Theorem~\ref{thm:rgfl-common-pair}.
Its terminal vacuum obeys
$|\varepsilon_k^{\rm v}|
\le5c_{\rm fl}e_*^{1/4-6\varepsilon}$.
The second inequality of \eqref{eq:term-stop-conditions}
therefore gives
\[
 |\varepsilon_k^{\rm v}|\ell_*^3
          \le\tfrac14e_*^{1/4-7\varepsilon}
          \le e_*^{1/4-9\varepsilon}.
\]
Remove this term from
\eqref{eq:term-partition-free-bound} and
\eqref{eq:term-partition-pressure-bound}. This proves
the first two estimates in
\eqref{eq:term-common-pressure-bound}.
The function $\varepsilon_5$ is decreasing.
The scheme's scale bounds give
$\bar mR\ge\mu_K>1/L$ and $e_*\le e_K$.
Since $1/4-9\varepsilon>0$, they prove its last estimate.

Both $R_*$ and $LR_*=R_{\mathrm{top}}$ belong to
$\mathcal T(e)$. The mesh threshold is monotone:
$N_{\min}(LR_*,e)\le N_{\min}(R_*,e)=N_{\rm ref}(e)$.
Thus the common pair and pressure are available on
both tori at every cutoff in the theorem.
The bulk pressure cancels exactly:
\[
 D_{\rm mat}=
 L^3(\log z_{r,n_*}-p_Fn_*^3)
             -(\log z_{r,Ln_*}-p_F(Ln_*)^3).
\]
The pressure bounds give the first inequality in
\eqref{eq:term-matter-comparison}.
Monotonicity of $\varepsilon_5$ and $e_*\le e_K$
give the second. Finally
$\bar mR_*>1/L$ and
$e_K^{1/4-9\varepsilon}\le1/8$ show that the bracket
is at most $\varepsilon_5(1/L)+1$.
The terminal scales, $R_*$ and $e_K$, are independent
of $N$; this proves the stated uniformity.
\end{proof}


\section{Interacting response on the reference torus}
\label{sec:interacting-response}

The free coefficients do not by themselves control the
reference Wick products. We must compare the normalized
interacting generating function with its one-cell free
factors, including mixed derivatives in different cells.
The quantity to be estimated is
\begin{equation}\label{eq:interacting-reference-error}
 \Theta_m^{(0)}=
 \left|\left(\prod_{c\in\mathcal C_*}
       (\partial_{t_c}^m-c_m)\mathring\Phi\right)(0)\right|^{1/N_*},
 \qquad m=2,4,\qquad N_*=|\mathcal C_*|=8L^3,
\end{equation}
where $c_m=N_0^{(m)}(0)$. The exact relation to Wick
products is \eqref{eq:external-Wick-product}.
The proof uses decay between separate insertions,
an analytic logarithm for the terminal fermion integral,
and a separate estimate of the complement of the
small-field set.

\subsection{The terminal determinant ratio}
\label{subsec:interacting-determinant-ratio}

For the marked connections of
\eqref{eq:interacting-marked-connections}, set
\begin{equation}\label{eq:interacting-mark-radii}
 C_b=1+\sqrt h\,C_\beta,\qquad
 t_x=\frac{e_*^{3\varepsilon}}{C_b},\qquad
 t_x^+=\frac{e_*^{3\varepsilon}}{\sqrt h\,C_\beta},
 \qquad r=t_x/4.
\end{equation}
If $C_\beta=0$, interpret $t_x^+=\infty$.
For $t$ in the polydisc
$\mathbb D(t_x^+)=\{|t_c|<t_x^+\text{ for every }c\}$,
all the rescaled connections $a(t)^{[j]}$
belong to $\mathfrak X_j$ and
$e_*|\mathscr C(t)|_\infty<e_*^{3\varepsilon}$.
Indeed, \eqref{eq:interacting-external-probe-bound}
and $e_*\le e\sqrt h$ give
$|a(t)|_{\rm ext}<e_*^{-1+3\varepsilon}$.
This bounds the stop sup and first-difference
norms by the defining thresholds of $\mathfrak X_*$.
For $i=k_*-j$ the rescaling multiplies
the sup norm by $L^{-i/2}$ and the first
differences by $L^{-3i/2}$, while
$e_j=L^{-i/2}e_*$. Since
$\varepsilon<1/3$, the rescaled bounds are
strictly below $e_j^{-1+2\varepsilon}$ and
$e_j^{-1+3\varepsilon}$, respectively.

We first control the determinant activities accumulated
along the flow. Recall the local determinant
function $E_j^d(X,A)$ of
Proposition~\ref{prop:local-determinant-expansion} and set
$E_j^{\rm det}(X,A)=E_j^d(X,A)-E_j^d(X,0)$.
Its analyticity, locality, real gauge invariance,
and evenness will all be used below; the relevant
bound is
\begin{equation}\label{eq:interacting-determinant-input}
 |E_j^d(X,A)|\le C_dM^3e^{-(\kappa+1)d_M(X)}
                       \qquad(A\in\widetilde G_j).
\end{equation}
Let $\mathfrak P_0=0$ and
\begin{equation}\label{eq:interacting-P-recursion}
 \mathfrak P_{j+1}(Y,C)=
 \sum_{\overline X^{\,L}=LY}
       \bigl(\mathfrak P_j(X,C_L)+E_j^{\rm det}(X,C_L)\bigr).
\end{equation}
Then $\mathfrak P_j(X,\cdot)$ is analytic on
$\mathfrak X_j$, vanishes at zero, and depends only
on bonds meeting $X$. To see locality under
reblocking, a bond meeting $X$ has, after division
by $L$, an endpoint in $Y$ whenever
$\overline X^{\,L}=LY$. Analyticity follows from
$C_L\in\mathfrak X_j$.

For $d_M(Y)\ge1$, the reblocking estimate gives
\[
 \sum_{\overline X^{\,L}=LY}e^{-\kappa d_M(X)}
 \le c_{\rm rb}L^3e^{-L(\kappa-\kappa_0-1)d_M(Y)}.
\]
With
$L(\kappa-\kappa_0-1)\ge\kappa$ and
$c_{\rm rb}L^3e^{-[L(\kappa-\kappa_0-1)-\kappa]}\le1/2$,
the weighted supremum $q_j$ of $\mathfrak P_j$
over $d_M(X)\ge1$ satisfies
$q_{j+1}\le(q_j+2C_dM^3)/2$. Since $q_0=0$,
\begin{equation}\label{eq:interacting-P-bound}
 |\mathfrak P_j(X,C)|\le2C_dM^3e^{-\kappa d_M(X)}
                           \qquad(d_M(X)\ge1).
\end{equation}
At $j=0$ the determinant activity includes all
$192$ chamber components. The primitive bound and
reblocking estimate have the same constants at
that first step.

The determinant identity
\eqref{eq:rg-step-determinant-normalization} gives
\begin{align}
 \frac{\det(D^{\rm ch}[qa]+\bar m\eta)}{\det(D^{\rm ch}[0]+\bar m\eta)}
 &=\exp\left(\sum_X\mathfrak P_*(X,a^{[k_*]})\right)
                    \frac{\det D_*(a^{[k_*]})}{\det D_*(0)},
                           \label{eq:interacting-determinant-telescope}\\
 \sum_X\mathfrak P_*(X,a^{[k_*]})
 &=\sum_{j<k_*}\sum_XE_j^{\rm det}(X,a^{[j]}).
                           \label{eq:interacting-P-telescope}
\end{align}
Here $a^{[j]}\in\mathfrak X_j$ for every $j$.
The fine determinant uses the phase-argument convention
of \eqref{eq:response-free-ratio}: its connection argument
is $qa(t)=\sqrt\eta\sum_ct_c\widehat{\mathfrak a}_c$.
For a single cell this is precisely the phase connection
$B_t$ in \eqref{eq:response-free-ratio}.
The one-step identity reads
\[
 \det D_{j+1}(a^{[j+1]})
       =b^{4n_{j+1}}\det D_j(a^{[j]})/\delta Z_j(a^{[j]}).
\]
Divide it by its zero-field version and use
\eqref{eq:local-determinant-quotient}. The resulting
ratio acquires the factor
\[
 \exp\left(-\sum_XE_j^{\rm det}(X,a^{[j]})\right)
\]
at each
step, which proves the first identity by iteration.
Summing \eqref{eq:interacting-P-recursion} over $Y$
proves the second: each fine polymer has exactly
one reblocked image. The zero-field determinants
are nonzero by the real covariance formula.
At $j=0$ the incoming determinant has dimension
$192n_0$, the outgoing one dimension $4n_1$;
the factor $b^{4n_1}$ cancels in the ratio.

Fix a coarse real photon field $A$ with
$\chi_*(A)=1$, and a real twist $\theta$.
The real gauge construction in
\eqref{eq:photon-inner-basepoint-bound}--%
\eqref{eq:photon-seam-basepoint-bound}
gives a representative $A^g$ with
$|A^g|_\infty\le2\ell_*^2p_*$.
The curvature bound and
Proposition~\ref{prop:photon-domain-preservation}
put both $A+c_\theta$ and $A^g+c_\theta$
in $\mathrm{Dom}_{k_*}$.
Write
\[
 \mathscr A^0=H_*^{\rm ph}A^g,\qquad
 c_{\theta,\mu}=\frac{\theta_\mu}{e_*\ell_*},
 \qquad \mathscr A=\mathscr A^0+c_\theta.
\]
Summing the three photon kernel estimates in
\eqref{eq:photon-kernel-regularity} bounds each
field norm of $\mathscr A^0$ by
$2C_0\ell_*^2p_*<e_*^{-\varepsilon}/2$;
the strict inequality is the corresponding
condition in \eqref{eq:term-stop-conditions}.
The local photon domains and their invariance
under real constants then give
$s\mathscr A^0+c_\theta\in\frac14\widetilde R_*$
for $0\le s\le1$. In particular all
connections needed below belong to $\mathcal W$.

The flow bound \eqref{eq:rgfl-reference-terminal-input}
gives an activity with even Grassmann kernels,
with equal barred and unbarred degrees, and
\begin{equation}\label{eq:interacting-activity-input}
 \|\mathscr E_*\|_*^x
       \le4e_*^{1/4-6\varepsilon},\qquad
 |m_*|\le e_*^{3/4-8\varepsilon}/8.
\end{equation}
The coefficient norm includes the weights
$\eta_*^3/48$ for a fine spinor index and
$(\eta_*^3/48)^2$ for a pair index, with all
admissible pairs, including pairs inside a cluster.
Its polymer factor is $e^{\kappa d_M(X)}$.
For $(A,C)\in\widetilde R_*\times\mathfrak X_*$ define
\begin{equation}\label{eq:interacting-two-field-integral}
 Z_{\rm in}(A,C)=\int D\Psi\,
 \exp\left\{-\langle\bar\Psi,D_*(A+C)\Psi\rangle
       -m_*\langle\bar\psi,\psi\rangle
       +\sum_X\mathscr E_*(X;A,C,\psi,\chi)\right\},
\end{equation}
where $\psi,\bar\psi$ are substituted using
$H_*^\#(A+C)\Psi$ and their barred partners,
and $\chi,\bar\chi$ using the corresponding
covariant Hölder quotients.
The inclusion
$\widetilde R_*+\mathfrak X_*\subset\widetilde G_*$
makes all these maps analytic; the activity
kernels are analytic in $(A,C)$.
Berezin integration is coefficient extraction,
and the exponential is its scalar exponential
times a finite polynomial in the positive-degree
coefficients. Thus $Z_{\rm in}$ is analytic
on this two-field domain.
Set
\begin{align}
 Z_{\rm tm}(t)&=Z_{\rm in}(\mathscr A,\mathscr C(t)),
                               \label{eq:interacting-terminal-integral}\\
 f_{\rm int}(t)&=
   \frac{Z_{\rm tm}(t)}{\det D_*(c_\theta+\mathscr C(t))}.
                               \label{eq:interacting-factor-definition}
\end{align}
The Berezin measure has one pair for each
of the $4\ell_*^3$ coarse spinor coordinates.
The two-field analyticity gives analyticity of
$Z_{\rm tm}$ on
$\mathbb D(t_x^+)$.

\begin{proposition}\label{prop:interacting-factor}
Suppose
\begin{gather}
 4C_0\ell_*^3g_1\le e_*^{-1/4+\varepsilon},\qquad
 4\ell_*^3(2C_K+K_0+4)\le e_*^{-2\varepsilon},\qquad
 e_*^{1/4-9\varepsilon}\le1/8,
                                      \label{eq:interacting-factor-conditions}\\
 48e_*^{2\varepsilon}\le\mu_- .
                                      \label{eq:interacting-small-complex-mass}
\end{gather}
Then $f_{\rm int}$ is analytic and
\begin{equation}\label{eq:interacting-factor-bound}
 |f_{\rm int}(t)-1|\le\frac12e_*^{1/4-9\varepsilon}
                              \le\frac1{16}.
\end{equation}
In particular $Z_{\rm tm}$ has no zero on the
polydisc. The normalized terminal ratio
\begin{equation}\label{eq:interacting-terminal-ratio}
 R_{\rm tm}(t;\theta,A)=
 \frac{Y_*^{\rm x}(A+c_\theta;\mathscr C(t))}
      {Y_*^{\rm x}(A+c_\theta;0)}
 =\exp\left(\sum_X\mathfrak P_*(X,\mathscr C(t))\right)
       \frac{Z_{\rm tm}(t)}{Z_{\rm tm}(0)}
\end{equation}
is nonzero, continuous in $(t,\theta,A)$ on
$\mathbb D(t_x^+)\times\R^3\times\{\chi_*=1\}$,
and analytic in $t$. It has a unique continuous
logarithm $\mathcal G(t;\theta,A)$ that vanishes
at zero, and this logarithm is analytic in $t$.
\end{proposition}

\begin{proof}
Use the covariance
\[
 C_t=D_*(c_\theta+\mathscr C(t))^{-1}.
\]
The inverse exists by
\eqref{eq:interacting-complex-covariance}, because
$e_*|\mathscr C(t)|_\infty<e_*^{3\varepsilon}
<\mu_*/24$. Each covariance row has sum at
most $4\ell_*^3g_1$. Put
$a_w=C_0^{-1}e_*^{-1/4+\varepsilon}$.
The first condition in
\eqref{eq:interacting-factor-conditions} gives
$a_w\ge4\ell_*^3g_1\ge1$.

Separate the quadratic form with covariance $C_t$
from the remaining even interaction $V_t$.
Normalized Gaussian Grassmann integration gives
\begin{equation}\label{eq:interacting-crude-Gaussian}
 |f_{\rm int}(t)-1|\le e^{|V_t|_{a_w}}-1.
\end{equation}
To justify this particular norm bound, replace
the unbarred generators by $C_t\Psi$ and leave
the barred generators fixed. Their row sums
are at most $a_w$, so this substitution maps
the $a_w$ coefficient norm into the unit
coefficient norm. Integration against the identity
covariance has norm one in that latter norm.
Apply this to $e^{V_t}-1$, and use the
submultiplicativity of the coefficient norm.

There are three terms in $V_t$. The difference
between the two coarse quadratic forms has
entry sum at most
$C_Ke_*^{1-\varepsilon}\ell_*^3/2$, by
\eqref{eq:interacting-coarse-first} integrated along
$c_\theta+\mathscr C(t)+s\mathscr A^0$.
A quadratic Grassmann form has coefficient norm
at most $4a_w^2$ times this block sum, so
its contribution is bounded by
$2C_K\ell_*^3e_*^{1/2+\varepsilon}$.
The terminal mass estimate and
\eqref{eq:interacting-activity-input} give respectively
\[
 \frac12\ell_*^3e_*^{1/4-6\varepsilon},
 \qquad K_0\ell_*^3e_*^{1/4-7\varepsilon}.
\]
For the activity sum, substitute the terminal
field maps in the coefficient norm, then sum
$e^{-\kappa d_M(X)}$ by assigning each polymer
one of its $M$-cubes and using the polymer
sum bound $K_0$. We have used
$4e_*^\varepsilon\le1$.
Consequently
\[
 |V_t|_{a_w}\le
 \ell_*^3(2C_K+K_0+4)e_*^{1/4-7\varepsilon}
 \le\frac14e_*^{1/4-9\varepsilon}.
\]
The inequality $e^x-1\le2x$ proves
\eqref{eq:interacting-factor-bound}.
The denominator in \eqref{eq:interacting-factor-definition}
is analytic and nonzero, so $f_{\rm int}$ is analytic.

By \eqref{eq:term-Y-density} the terminal density has the form
\[
 Y_*^{\rm x}(A';C)=
 (\det D_*(0))^{-1}
 \exp\left(-\varepsilon_*^{\rm v}\ell_*^3+
                         \sum_X\mathfrak P_*(X,C)\right)
 Z_{\rm in}(H_*^{\rm ph}A',C).
\]
Its gauge invariance in Theorem~\ref{thm:term-density}
permits replacing $A$ by
$A^g$. The photon minimizer preserves real
constants. Dividing the density at $C=\mathscr C(t)$
by its value at zero proves
\eqref{eq:interacting-terminal-ratio}; the quotient
also shows that the result is independent of
the chosen gauge representative. For continuity,
use the ungauged linear map
$A+c_\theta\mapsto H_*^{\rm ph}(A+c_\theta)$
in the displayed formula. The terminal Grassmann
integral is analytic in its two field arguments,
and its zero-mark denominator is nonzero.

The required logarithm is
\begin{equation}\label{eq:interacting-logarithm}
 \mathcal G(t;\theta,A)=
 \int_0^1\sum_ct_c
       \frac{\partial_{t_c}R_{\rm tm}(st;\theta,A)}
            {R_{\rm tm}(st;\theta,A)}\,\dd s .
\end{equation}
Local Cauchy formulas show that the derivatives
in this integrand are jointly continuous;
integration over $[0,1]$ preserves continuity
and analyticity. Along a ray the integrand is
the logarithmic derivative of $R_{\rm tm}(st)$.
Differentiating the product of this ratio with
the exponential of minus the integrated
logarithmic derivative gives zero, and its
value at $s=0$ is one. Hence
$e^{\mathcal G}=R_{\rm tm}$.
Any two continuous logarithms differ by
$2\pi i$ times an integer-valued continuous
function on the connected polydisc. Their
common value at zero proves uniqueness.
\end{proof}

\subsection{The one-cell free logarithm and twists}
\label{subsec:interacting-one-cell}

For a fixed cell $c$ and real twist $\theta$, let
\begin{equation}\label{eq:interacting-free-coarse-log}
 g_D^\theta(t)=\tr\log\left(1+D_*(c_\theta)^{-1}
                   [D_*(c_\theta+t\widehat c_c)-D_*(c_\theta)]\right).
\end{equation}
The logarithm here is the matrix power series.
Set
\begin{align}
 C_{\rm fr}&=2592C_dK_0C_\beta^2
       \left(M^2(D_\beta^{\rm buf})^3L^4h^{-2}+8M^5h\right),
                                                   \notag\\
 B_{\rm wr}&=16\exp(6)C_dM^3K_0p_\beta e^{-r_*},
 \qquad B_1=16g_1C_K'p_\beta e_*^{2\varepsilon},       \notag\\
 B_0&=C_{\rm fr}e_*^{4\varepsilon}+B_{\rm wr}+B_1 .
                                \label{eq:interacting-free-log-constants}
\end{align}

\begin{proposition}\label{prop:interacting-one-cell}
Suppose \eqref{eq:interacting-factor-conditions}--%
\eqref{eq:interacting-small-complex-mass} hold,
$r_*>31L+16$, and
\begin{equation}\label{eq:interacting-one-cell-conditions}
 2g_1C_K'e_*^{2\varepsilon}\le1/2,\quad
 B_0\le1/4,\quad
 \tau_s:=c_\tau\ell_*\ge1,\quad
 52B_1e^{-\tau_s}\le e_*^{1/4-9\varepsilon}.
\end{equation}
The free ratio $N_0$ has the logarithm
\begin{equation}\label{eq:interacting-one-cell-free-log}
 \ell_0(t)=\sum_X\mathfrak P_*(X,t\widehat c_c)+g_D^0(t),
 \qquad N_0(t)=e^{\ell_0(t)} ,
\end{equation}
analytic on $|t|<t_x^+$. It is independent of
$c$ and of the admitted placement, is even
and real on the real axis, and satisfies
\begin{align}
 \left|\sum_X\mathfrak P_*(X,t\widehat c_c)\right|
        &\le C_{\rm fr}e_*^{-2\varepsilon}|t|^2+B_{\rm wr},
                                  \label{eq:interacting-P-probe-gain}\\
 |\ell_0(t)|&\le B_0\qquad(|t|\le t_x).
                                  \label{eq:interacting-free-log-bound}
\end{align}
The one-cell interacting error
$w_c(t)=\mathcal G(te_c;\theta,A)-\ell_0(t)$
is analytic, vanishes at zero, and obeys
\begin{equation}\label{eq:interacting-one-cell-error}
 |w_c(t)|\le3e_*^{1/4-9\varepsilon}
                         \qquad(|t|<t_x^+).
\end{equation}
\end{proposition}

\begin{proof}
Write $K=D_*(c_\theta+t\widehat c_c)-D_*(c_\theta)$.
The real covariance and the integrated insertion
bounds give
\[
 \|D_*(c_\theta)^{-1}K\|_2\le
           2g_1C_Ke_*^{2\varepsilon}\le1/2,\qquad
 \|D_*(c_\theta)^{-1}K\|_1
           \le4g_1C_Kp_\beta e_*^{2\varepsilon}.
\]
Thus the series in \eqref{eq:interacting-free-coarse-log}
converges, is analytic, vanishes at zero, and
has modulus at most $B_1$. Triangularizing
the matrix identifies its exponential with
$\det D_*(c_\theta+t\widehat c_c)/\det D_*(c_\theta)$.
Combine this with
\eqref{eq:interacting-determinant-telescope}.
The lattice symmetries carrying one cell
probe to another preserve the fine determinant,
so the resulting fine ratio is $N_0$.

The bound \eqref{eq:interacting-P-probe-gain}
requires a cancellation before summing the levels.
Put $i=k_*-j$ and $B_j=e^{-1}\widehat{\mathfrak a}_c^{[j]}$.
Rescaling gives
\begin{equation}\label{eq:interacting-probe-derivative}
 e_j|\partial B_j|_\infty\le
                         a_j^2a_*^{-3/2}C_\beta .
\end{equation}
Let $\mathcal Q_j$ be the $M$-cubes containing
endpoints of the supporting bonds of $B_j$.
The buffered support cube yields
\begin{equation}\label{eq:interacting-probe-cube-count}
 |\mathcal Q_j|\le
 4\left((D_\beta^{\rm buf})^3M^{-3}a_j^{-3}+8\right),
 \qquad |\mathcal Q_j|\le L^{3i}p_\beta .
\end{equation}
By locality only polymers containing a cube
of $\mathcal Q_j$ contribute to
\eqref{eq:interacting-P-telescope}.

For such a polymer $X$, put $d=d_M(X)$ and
$s_X=M(d+3)$. Suppose first that $2s_X<\ell_j$.
All bonds meeting $X$ lie in a coordinate
cube of radius $M(d+2)+1$ around a chosen
site $\tilde c$ of $X$. The strict inequality
$2[M(d+2)+1]+\eta_j<2s_X<\ell_j$
makes this cube a lift to the lattice.
Define a real function there by
$\lambda(c)=\sum_\nu B_{j,\nu}(\tilde c)(c-\tilde c)_\nu$,
and extend it by zero outside. On bonds meeting
$X$, its gradient is the constant value
$B_j(\tilde c)$. If $[\,\cdot\,]_X$ restricts
a connection to those bonds, then
\[
 B_X=[B_j-d\lambda]_X,\qquad
 |B_X|_\infty\le3s_X|\partial B_j|_\infty.
\]
The last inequality follows by integrating
the first differences along coordinate paths
inside the lifted cube.

Real gauge invariance must be used before
continuing to the complex gauge parameter.
On the convex open set where
$|[tB_j-zd\lambda]_X|_\infty<e_j^{-1+\varepsilon}$,
the analytic function
$z\mapsto E_j^d(X,[tB_j-zd\lambda]_X)$
equals $E_j^d(X,tB_j)$ for real $z$ near zero:
the unrestricted connection is related to
$tB_j$ by a real gauge transformation, and
locality permits the restriction. The identity
theorem gives this equality throughout that
set. In particular it holds at $z=t$
whenever $|t|<\rho_X$, where
$\rho_X=e_j^{-1+\varepsilon}/(3s_X|\partial B_j|_\infty)$.

The function
$z\mapsto E_j^d(X,zB_X)-E_j^d(X,0)$ is even,
bounded by $2C_dM^3e^{-(\kappa+1)d}$ on
$|z|<\rho_X$, and has a double zero at zero.
Apply the maximum principle after dividing
by $z^2$. For $|t|\ge\rho_X$ the original
primitive bound gives the same inequality
with the factor $(|t|/\rho_X)^2\ge1$.
Thus, in both cases,
\begin{equation}\label{eq:interacting-local-constant-removal}
 |E_j^{\rm det}(X,tB_j)|\le
 18C_dM^5C_\beta^2(d+3)^2e^{-(\kappa+1)d}
          |t|^2e_j^{-2\varepsilon}a_j^4a_*^{-3}.
\end{equation}
When $\partial B_j=0$, $B_X=0$ and this contribution
vanishes. This proves the cancellation also
at the chamber first step.

Since $(d+3)^2e^{-(\kappa+1)d}\le9e^{-\kappa_0d}$,
summing \eqref{eq:interacting-local-constant-removal}
over cubes and polymers gives at level $j$
the bound
\[
 648C_dK_0C_\beta^2|t|^2e_j^{-2\varepsilon}a_*^{-3}
           \left(M^2(D_\beta^{\rm buf})^3a_j+8M^5a_j^4\right).
\]
For the remaining polymers $2s_X\ge\ell_j$,
$d\ge r_*L^i/2-3$. The extra exponential in
the primitive bound then gives at level $j$
\[
 2\exp(6)C_dM^3K_0p_\beta L^{3i}e^{-r_*L^i}.
\]
Sum over $i\ge1$, using $a_j=L^{-i}a_*$,
$e_j^{-2\varepsilon}=e_*^{-2\varepsilon}L^{i\varepsilon}$,
$\sum_iL^{-i(1-\varepsilon)}\le1$, and
$a_*^{-2}\le L^4h^{-2}$, $a_*\le h$.
For $r_*>31L+16$, the second sum is bounded
by a constant times $e^{-r_*}$.
The stated definitions of $C_{\rm fr}$ and
$B_{\rm wr}$ give \eqref{eq:interacting-P-probe-gain}.
Together with $|g_D^0|\le B_1$ this proves
\eqref{eq:interacting-free-log-bound}.

Evenness and reality follow from the corresponding
fine determinant symmetries. Both $\ell_0(t)$
and $\ell_0(-t)$ vanish at zero and have
exponential $N_0(t)$, so they coincide.
For real $t$, the nonzero real ratio $N_0(t)$
is positive by continuity from one at zero;
the normalized logarithm is therefore real.
Cell and placement independence follows in
the same way from the determinant translation
and reflection symmetries.

It remains to control the twist. For
$|\operatorname{Im}\theta_\mu|\le\tau_s$
the connections $c_\theta+st\widehat c_c$
lie in $\mathcal W^+$, and
$\|D_*(c_\theta)^{-1}\|_2\le g_1$.
There is a positive strip margin
$\ell_*(e_*^{2\varepsilon}-e_*^{3\varepsilon})$:
split the imaginary constant connection
into a part bounded by $c_\tau/e_*$ and a
remaining part absorbed into $\mathfrak X_*$.
The latter has vanishing first differences.
The mass condition bounds its total imaginary
part by $\mu_*/24$. On this larger open strip,
\eqref{eq:interacting-enlarged-derivative} gives
the same matrix power-series bound $|g_D^\theta|\le B_1$.
The function is holomorphic there and
$2\pi$-periodic in each twist variable.
For real twists, periodicity is conjugation by
the periodic gauge factor
$\exp(-2\pi i c_\mu/\ell_*)$; the identity
theorem extends it to the strip.

We recall the short Fourier estimate to make the
size of this suppression explicit. A bounded
periodic holomorphic function $g$ in a closed
strip of width $\tau\ge1$ has
$|\widehat g(m)|\le B e^{-\tau|m|_1}$:
shift each Fourier contour towards
$-i\tau\operatorname{sgn}(m_\mu)$, first using
an interior width and then taking its limit.
The vertical sides cancel by periodicity.
The nonconstant Fourier sum is bounded by
\[
 B\left[\left(\frac{1+e^{-\tau}}{1-e^{-\tau}}\right)^3-1\right]
                                      \le25Be^{-\tau}.
\]
The last inequality follows by expanding the
bracket divided by $e^{-\tau}$ as a power
series with positive coefficients and evaluating
at $e^{-\tau}=e^{-1}$. Thus the difference
between two real values is at most $52Be^{-\tau}$.
Apply this to $g_D^\theta(t)$ to obtain
\begin{equation}\label{eq:interacting-twist-suppression}
 |g_D^\theta(t)-g_D^0(t)|\le52B_1e^{-\tau_s}
                                  \le e_*^{1/4-9\varepsilon}.
\end{equation}

Finally the normalized logarithm identity and
\eqref{eq:interacting-factor-definition} give
\[
 w_c(t)=g_D^\theta(t)-g_D^0(t)
                +\log f_{\rm int}(te_c)-\log f_{\rm int}(0).
\]
The principal logarithm of $f_{\rm int}$ is
analytic and has modulus at most
$e_*^{1/4-9\varepsilon}$. Both sides vanish
at zero and exponentiate to the same ratio;
uniqueness of the logarithm justifies the
identity. Equation
\eqref{eq:interacting-twist-suppression} proves
\eqref{eq:interacting-one-cell-error}.
\end{proof}

Let $\phi_0(t)=\sum_{i=0}^4N_0^{(i)}(0)t^i/i!$.
Cauchy's estimate applied to
$|N_0(t)-1|\le2B_0$ on $|t|=t_x=4r$ gives
\begin{equation}\label{eq:interacting-free-jet}
 |c_2|\le r^{-2}/16,\qquad |c_4|\le r^{-4}/16,\qquad
 \|\phi_0-1\|_r\le1/6,\qquad \|\phi_0\|_r\le2.
\end{equation}
The odd coefficients vanish. Expanding $\log N_0$
through degree four gives
$\ell_0^{(4)}(0)=c_4-3c_2^2$; hence this
is the chamber fourth coefficient estimated
in the free-response section.

\subsection{Derivatives in separated cells}
\label{subsec:interacting-separated}

To estimate mixed logarithmic derivatives, we use the
real covariance $G=D_*(\mathscr A)^{-1}$. On the
coarse Grassmann algebra, write an even element
without constant term as
$F=\sum_{p\ge2}(p!)^{-1}\sum F_p(u_1,\ldots,u_p)
\Psi(u_1)\cdots\Psi(u_p)$ with antisymmetric kernels.
For a symmetric weight $w\ge1$, define
\begin{align}
 \rho_G(F;w)&=\sum_{p\ge2}\frac{h_G^p}{(p-1)!}
       \sup_{u_1}\sum_{u_2,\ldots,u_p}|F_p(u_1,\ldots,u_p)|w,
                                                     \notag\\
 N_G(F;w)&=\sum_{p\ge2}\frac{h_G^p}{(p-1)!}
       \sum_{u_1,\ldots,u_p}|F_p(u_1,\ldots,u_p)|w .
                                     \label{eq:interacting-tree-norms}
\end{align}
The supremum in $\rho_G$ is taken separately
at each degree. The indices include site, spin,
and barred or unbarred status; their distance
is the distance between the sites. We use
\[
 w_\delta(u_1,\ldots,u_p)=
          e^{\vartheta_0\max_{i,j}d(u_i,u_j)},\qquad
 w_P(u_1,\ldots,u_p)=e^{\vartheta_0\max_i d(u_i,P)} .
\]
The constants $h_G,\vartheta_0,\alpha_G$ are
those of \eqref{eq:interacting-Gaussian-covariance}.
Antisymmetrizing a kernel that is not antisymmetric
bounds $N_G$ by the same expression for that
kernel. For $\rho_G$, replace the supremum at
the first argument by the largest supremum
over all argument positions.

In the fine algebra absorb the terminal mass
into each single-cube activity:
\[
 E^a(X;t)=\mathscr E_*(X;\mathscr A,\mathscr C(t))
              -\mathbf1_{\{X\text{ is an }M\text{-cube}\}}
                         m_*\langle\bar\psi,\psi\rangle_X.
\]
Let $E^a_{nm}(X;t)$ denote its fine spinor/pair
kernels. The antisymmetric mass kernel has
norm $8M^3$: its two fine-index weights multiply
$48\eta_*^{-3}$, its diagonal Kronecker factors
leave $48(M/\eta_*)^3$ chamber sites, and there
are four spin components and two orders.
It follows from \eqref{eq:interacting-activity-input}
that, when $a_w\le e_*^{-1/4+\varepsilon}$ and
$(4+M^3/2)e_*^\varepsilon\le1$,
\begin{equation}\label{eq:interacting-absorbed-activity}
 \sum_{n,m}\frac{a_w^{n+m}}{n!m!}\|E^a_{nm}(X;t)\|
          \le w_*^{\rm E}e^{-\kappa d_M(X)},\qquad
 w_*^{\rm E}=e_*^{1/4-7\varepsilon}
                                      \quad(|t|_\infty\le t_x).
\end{equation}
Cauchy's formula on the circles of radius $t_x/2$
in one or two distinct mark variables gives,
on $|t|_\infty\le t_x/2$, the same bound
for $\partial_cE^a_{nm}$ and
$\partial_c\partial_{c'}E^a_{nm}$ multiplied
respectively by $q_x=2/t_x$ and $q_x^2$.
The circle integrals precede the finite weighted
kernel sums, so the coefficient norm is bounded
without replacing it by a supremum of individual
coefficients. Locality gives
$\partial_cE^a(X;t)=0$ unless $X\cap P_c\ne\varnothing$.

After substituting the terminal field maps, put
\begin{align}
 V(t)&=-\langle\bar\Psi,K(t)\Psi\rangle
           +\sum_X E^a(X;t)[H_*^\#(\mathscr A+\mathscr C(t))\Psi],
                                                        \notag\\
 K(t)&=D_*(\mathscr A+\mathscr C(t))-D_*(\mathscr A),\qquad
 V_0(t)=\sum_XE^a_{00}(X;t),\qquad V_+(t)=V(t)-V_0(t).
                                    \label{eq:interacting-tree-interaction}
\end{align}
It is even and analytic, and $V_+$ has no
constant term. Gaussian integration gives
\begin{equation}\label{eq:interacting-real-Gaussian-factor}
 Z_{\rm tm}(t)=\det D_*(\mathscr A)e^{V_0(t)}
                                  \int\dd\mu_G\,e^{V_+(t)} .
\end{equation}
The choice of the real covariance will allow
the exponential weights in the tree expansion.

\begin{proposition}\label{prop:interacting-differentiated-norms}
Suppose
\begin{equation}\label{eq:interacting-tree-smallness}
 5h_{\rm L}h_G\le e_*^{-1/4+\varepsilon},\qquad
 (4+M^3/2)e_*^\varepsilon\le1,\qquad
 \alpha_G\bar r\le1/2,
 \quad
 \bar r=4h_G^2C_Ke_*^{2\varepsilon}
                            +650K_0e_*^{1/4-7\varepsilon}.
\end{equation}
Let $d_{\rm sep}=\min_{c\ne c'}d(P_c,P_{c'})$ and
suppose $d_{\rm sep}\ge6M+8$.
With
$s_F=(24/\gamma^F)^3$ and
$C_V^{\mathrm{ins}}=10^3C_H^2s_F^4K_0$, we have
\begin{align}
 \rho_G(V_+(t);w_\delta)&\le\bar r
                                      &&(|t|_\infty\le t_x),\notag\\
 N_G(\partial_cV_+(t);w_{P_c})&\le q_xC_V^{\mathrm{ins}}h_G^2p_\beta,
 \quad
 \rho_G(\partial_cV_+(t);w_{P_c})\le q_xC_V^{\mathrm{ins}}h_G^2,\notag\\
 N_G(\partial_c\partial_{c'}V_+(t);1)
          +|\partial_c\partial_{c'}V_0(t)|
 &\le q_x^2C_V^{\mathrm{ins}}h_G^2p_\beta
                    e^{-\vartheta_0(d_{\rm sep}-4M-6)} .
                             \label{eq:interacting-differentiated-norms}
\end{align}
The last two lines hold on $|t|_\infty\le t_x/2$,
with $c\ne c'$ in the last line. If the lattice
probe vanishes, all marked derivatives vanish.
\end{proposition}

\begin{proof}
We give the kernel summation which preserves
the marked supports. For a polymer $X$, let
$X^+$ be its unit-site one-neighborhood.
Every fine index in an activity supported on
$X$ has its position set in $X^+$. This
includes both chamber sites in an internal
pair and the first site of an external
nearest-neighbor pair whose other endpoint
lies over $X$. The polymer diameter estimate
gives
\begin{equation}\label{eq:interacting-X-diameter}
 \vartheta_0\operatorname{diam}(X^+)\le d_M(X)+3 .
\end{equation}
The weighted kernel after substituting the
field maps is, at degree $p=n+m$,
\[
 K_{nm}^X(u_1,\ldots,u_p)=
 \int E^a_{nm}(X;t;v_1,\ldots,v_p)
           \prod_{i=1}^p\operatorname{Leg}(v_i;u_i)\,\dd v .
\]
Here $\dd v$ includes the fine weights already
specified. Passing from the bidegree expression
to a degree-$p$ kernel produces the factor
$p!/(n!m!)$; in the tree norms its remaining
factor is $p h_G^p/(n!m!)$.

Suppose a weight is bounded by
$c_w\prod_i e^{\vartheta_0d(u_i,\operatorname{pos}(v_i))}$.
Each summed leg contributes at most $h_{\rm L}$.
If one argument $u$ is fixed, its leg contributes
at most
$h_{\rm L}e^{-(\gamma^F/4)d(u,X^+)}$.
Thus the full kernel sum is at most
$c_wh_{\rm L}^p\|E^a_{nm}(X;t)\|$, and the
sum with one fixed argument retains the
extra exponential in $d(u,X^+)$.
For $w_\delta$, the triangle inequality between
the two farthest arguments and their position
sets, followed by \eqref{eq:interacting-X-diameter},
permits $c_w=e^{d_M(X)+3}$.

The needed rooted polymer sum is
\begin{equation}\label{eq:interacting-rooted-polymer-sum}
 \sum_X e^{-(\gamma^F/4)d(u,X^+)}
                       e^{-(\kappa-1)d_M(X)}\le32K_0.
\end{equation}
Indeed the sum of
$e^{-(\gamma^F/4)d(u,\square^+)}$ over $M$-cubes
is at most $32$. The neighboring $27$ cubes
contribute at most $27$; a cube $j$ layers
away has distance at least $(j-1)M$ and there
are at most $(2j+1)^3$ such cubes.
Use $\gamma^FM/4\ge4$ to sum the remaining
geometric tail. Assigning a cube of $X$
and applying the anchored polymer sum $K_0$
proves \eqref{eq:interacting-rooted-polymer-sum}.

Apply \eqref{eq:interacting-absorbed-activity}
with $a_w=5h_{\rm L}h_G$. The factor $5^{-p}$
absorbs the polynomial degree factors.
The preceding full and rooted sums give
\[
 \rho_G\left(\sum_X E^a(X;t)[H_*^\#\Psi]_+;w_\delta\right)
                                       \le650K_0w_*^{\rm E}.
\]
More generally, for kernels bounded by
$\mathcal B e^{-\kappa d_M(X)}$ at that
weight, these sums give
$650K_0\mathcal B$ in $\rho_G$ and
$\exp(3)\mathcal B\sum_Xe^{-(\kappa-1)d_M(X)}$
in $N_G$. This follows by summing
$\sum_{p\ge2}p5^{-p}=9/80$ in the rooted
case and using $p\le5^p$ in the full case.
The quadratic term contributes at most
$4h_G^2C_Ke_*^{2\varepsilon}$, since
$\vartheta_0\le\gamma^F/2$. This proves the
first line of \eqref{eq:interacting-differentiated-norms}.

For the first derivative, there are three
types of term: the coarse quadratic derivative,
the activity-kernel derivative, and a derivative
of one leg. The first has respective bounds
$8h_G^2C_K\epsilon_cp_\beta$ and
$4h_G^2C_K\epsilon_c$ in $N_G$ and $\rho_G$,
with $\epsilon_c=e_*|\widehat c_c|_\infty\le q_x$.
In its chain through $z\in P_c$, the weight
$w_{P_c}$ is absorbed by the exponential
distances from the two coarse endpoints to $z$.

For the kernel derivative, $X$ must meet $P_c$.
The same weight comparison uses that meeting
point and \eqref{eq:interacting-X-diameter}.
The Cauchy bound multiplies
$\mathcal B=w_*^{\rm E}$ by $q_x$, and
the sum of polymers meeting $P_c$ is at most
$K_0p_\beta$. Its $N_G$ bound is
$\exp(3)K_0p_\beta q_xw_*^{\rm E}$ and its
$\rho_G$ bound is $650K_0q_xw_*^{\rm E}$.

For a differentiated leg at position $j$, the
insertion estimate has a chain
$w\in\operatorname{pos}(v_j)\to z\in P_c^+\to u_j$.
The weight obeys
\[
 w_{P_c}\le e^{d_M(X)+5}
 e^{\vartheta_0[d(w,z)+d(z,u_j)]}
            \prod_{i\ne j}e^{\vartheta_0d(u_i,\operatorname{pos}(v_i))}.
\]
The extra two in the exponent accounts for
$d(P_c^+,P_c)\le M+1$.
Summing the differentiated leg and the other
legs leaves, before the degree sum,
\[
 2e^{d_M(X)+5}h_{\rm L}^{p-1}
       C_H\epsilon_cs_F
       \sum_{z\in P_c^+}e^{-(\gamma^F/2)d(z,X^+)}
                                      \|E^a_{nm}(X;t)\|.
\]
There are $p$ choices of the differentiated
leg. The kernel norm contributes one further
factor $p$, and $p^2\le5^p$ absorbs both.
Use \eqref{eq:interacting-rooted-polymer-sum}
with $\gamma^F/2$ in place of $\gamma^F/4$.
The result is bounded by
$9500C_Hs_FK_0w_*^{\rm E}p_\beta\epsilon_c$
in $N_G$.

For $\rho_G$, if the fixed argument lies on
the differentiated leg, split its two-edge
exponent as $7\gamma^F/8=3\gamma^F/8+\gamma^F/2$.
Summing the intermediate site leaves
$s_Fe^{-(\gamma^F/4)d(u,X^+)}$.
If the fixed argument lies on another leg,
that leg supplies the same rooted exponential,
and the differentiated leg contributes
$2C_H\epsilon_cs_F^2$.
Both cases therefore have the bound
$9500C_Hs_F^2K_0w_*^{\rm E}\epsilon_c$.
Here the degree sum is
$\sum_{p\ge2}p^25^{-p}=43/160$.
Adding the three types of term, using
$w_*^{\rm E}\le1/8$, $C_K\le C_Hs_F^2$,
and $C_H\ge2$, proves the second line with
$C_V^{\mathrm{ins}}=10^3C_H^2s_F^4K_0$.

For the mixed derivative, put
\[
 d_0=d(P_c^+,P_{c'}^+),\quad
 d_1=d(X^+,P_c^+),\quad d_2=d(X^+,P_{c'}^+).
\]
Then $d_0\ge d_{\rm sep}-2M-2$.
The differentiated coarse quadratic term is
bounded in $N_G$ by
\[
 16h_G^2C_H^2s_F^3p_\beta q_x^2
                                      e^{-(\gamma^F/2)d_{\rm sep}}.
\]
For an activity kernel with $p$ legs, there
are $(p+1)^2$ placements of the two derivatives.
They may act on two legs, the same leg, a
kernel and a leg, or both on the kernel.
The summed first differentiated legs are
bounded by $2C_H\epsilon_cs_F^2e^{-(\gamma^F/2)d_1}$
and its $c'$ version. The mixed leg is
bounded by
\[
 2C_H^2\epsilon_c\epsilon_{c'}s_F^3e^{-(\gamma^F/2)d_0}
       \bigl(e^{-(\gamma^F/2)d_1}+e^{-(\gamma^F/2)d_2}\bigr).
\]
If a derivative acts on the kernel, its
corresponding distance $d_i$ is zero by
locality. Normalize the first kernel
derivatives by $q_x$ and the mixed one by
$q_x^2$. Each of the four resulting kernel
types satisfies \eqref{eq:interacting-absorbed-activity}.
Every product-rule term is consequently at most
$4q_x^2C_H^2s_F^4h_{\rm L}^p\|E_{\rm type}\|\mathcal D(X)$,
where
\[
 \mathcal D(X)=e^{-(\gamma^F/2)(d_1+d_2)}
       +e^{-(\gamma^F/2)d_0}
                  \bigl(e^{-(\gamma^F/2)d_1}+e^{-(\gamma^F/2)d_2}\bigr).
\]

The triangle inequality through $X^+$ gives
$d_0\le d_1+d_2+M d_M(X)+2M+2$.
Use $\vartheta_0\le\gamma^F/8$ and
$\vartheta_0M\le1$ to obtain
\[
 \mathcal D(X)e^{-\kappa d_M(X)}
 \le2e^{-\vartheta_0(d_0-2M-2)}
       \bigl(e^{-(\gamma^F/4)d_1}+e^{-(\gamma^F/4)d_2}\bigr)
                                e^{-(\kappa-1)d_M(X)}.
\]
Sum the two distance factors over the sites
of $P_c^+$ and $P_{c'}^+$, then use
\eqref{eq:interacting-rooted-polymer-sum}.
The degree factor $p(p+1)^2$ is at most
$5^p$ for $p\ge2$. At $p=0$ there is
one kernel term with no legs, and the same
bound controls $\partial_c\partial_{c'}V_0$.
The four kernel types contribute in total
at most
$256C_H^2s_F^4K_0p_\beta q_x^2
e^{-\vartheta_0(d_{\rm sep}-4M-4)}$.
Adding the quadratic term proves the last
line of \eqref{eq:interacting-differentiated-norms}.
If the lattice probe is zero, $\mathscr C(t)=0$
and there is no mark dependence.
\end{proof}

We next turn these kernel estimates into a bound
for the logarithm. The finite-dimensional tree
formula is recalled with its proof, since the
distance between its two marked vertices is
the source of the mixed-cell decay.

For $n$ copies of the generators, let
$\Delta^{ij}=\sum_{x,y}G(x,y)
\partial_{\bar\Psi^{(j)}(y)}\partial_{\Psi^{(i)}(x)}$
and $\Delta_{\{i,j\}}=\Delta^{ij}+\Delta^{ji}$.
The derivatives are left Grassmann derivatives.
For a forest $T$ and edge parameters
$\tau\in[0,1]^T$, let $s_{ij}^T(\tau)$ be the
minimum parameter on the path from $i$ to $j$,
zero if there is no path, and $s_{ii}^T=1$.
Write $\Delta[s]=\sum_{i,j}s_{ij}\Delta^{ij}$
and let $\epsilon_0$ extract the constant term.
For even elements $F_1,\ldots,F_n$ without
constant term, define
\begin{equation}\label{eq:interacting-connected-tree}
 T_n(F_1,\ldots,F_n)=
 \sum_{\text{trees }T}\int_{[0,1]^T}
 \epsilon_0\left\{\prod_{\ell\in T}\Delta_\ell
          e^{\Delta[s^T(\tau)]}\prod_iF_i(\Psi^{(i)})\right\}\,\dd\tau .
\end{equation}
For $n=1$ this is the Gaussian expectation.

\begin{proposition}\label{prop:interacting-tree-log}
If $V_+(t)$ satisfies
\eqref{eq:interacting-tree-smallness} and is analytic,
then
\[
 W(t)=V_0(t)+\sum_{n\ge1}\frac1{n!}
                       T_n(V_+(t),\ldots,V_+(t))
\]
is analytic on $\mathbb D(t_x)$ and
$Z_{\rm tm}(t)=\det D_*(\mathscr A)e^{W(t)}$.
For $c\ne c'$ and nonempty unit-site sets $P,P'$,
\begin{align}
 |\partial_c\partial_{c'}(W-V_0)|
 \le{}&2N_G(\partial_c\partial_{c'}V_+;1)\notag\\
 &+2\alpha_GN_G(\partial_cV_+;w_P)
       \rho_G(\partial_{c'}V_+;w_{P'})
                                  e^{-\vartheta_0d(P,P')}.
                                  \label{eq:interacting-tree-mixed}
\end{align}
\end{proposition}

\begin{proof}
The forest interpolation formula for a polynomial $H$ is
\[
 H(\mathbf1)=\sum_F\int
       \left(\prod_{\ell\in F}\partial_{s_\ell}H\right)
                                 (s^F(\tau))\,\dd\tau.
\]
For a partition $\pi$ of the vertices, keep fixed
entries $s_{\mathrm{in}}(i,j)$ inside its blocks and
put $X_\pi(u)_{ij}=u$ on pairs joining distinct blocks.
The chain rule and the fundamental theorem of calculus give
\begin{equation}\label{eq:interacting-forest-first-step}
 G(X_\pi(u))=G(X_\pi(0))
       +\sum_{\substack{\ell=(i,j)\\
                        i,j\text{ in distinct blocks of }\pi}}
          \int_0^u(\partial_{s_\ell}G)(X_\pi(w))\,\dd w .
\end{equation}
When $\ell$ joins two blocks, merge them and assign
the value $w$ to every pair between those two blocks.
The interpolation point is unchanged. Begin with singleton
blocks, $G=H$ and $u=1$, and iterate this identity
on each integrand. Each derivative
merges two components, so after at most $n-1$
steps the procedure ends. The successive
parameters are ordered decreasingly; when two
vertices merge, their common value is the
minimum parameter on their eventual path.
For a fixed forest, its edge orderings divide
the parameter cube into these ordered simplices.
Summing them gives the formula.

The matrix $s^F$ is a Gram matrix of unit
vectors: at $u\in[0,1]$, assign to vertex $i$
the basis vector of its component among the
edges with $\tau_\ell>u$, and integrate their
scalar products in $u$. This scalar product
is the length of the interval on which the
two vertices remain connected, namely
$s_{ij}^F$. Wick's formula and the Gram
inequality therefore bound a residual monomial
of degree $p$ under $e^{\Delta[s^F]}$
by $g^p$, where $\|G\|_2\le g^2$.
Indeed, its determinant entries are
$\langle\omega_i\otimes\delta_x,\omega_j\otimes G\delta_y\rangle$.
These facts apply whether or not the
interpolated covariance is invertible.

Let vertex $i$ of a tree have degree $k_i$
and let the chosen homogeneous part of $F_i$
have degree $p_i$. The edge derivatives leave
the factorial $(p_i-k_i)!^{-1}$ and $p_i-k_i$
uncontracted fields at that vertex; terms with
$p_i<k_i$ vanish. Wick's Gram bound supplies
$g^{p_i-k_i}$. Sum contracted arguments from
the leaves to vertex $1$, using the weighted
row and column bound $\alpha_G$ for every
edge. If nonnegative symmetric weights $w_i$
and a distance $D$ satisfy
\[
 e^{\vartheta_0D}\le
       \prod_iw_i(\text{arguments at }i)
       \prod_{\ell=(i,j)}
          e^{\vartheta_0d(\text{edge arguments})},
\]
this yields $e^{-\vartheta_0D}\alpha_G^{n-1}$
times the full kernel sum at vertex $1$
and the rooted sums at the other vertices.
The number of trees with degrees $k_1,\ldots,k_n$
is $(n-2)!/\prod_i(k_i-1)!$.
Enlarge the sum over degrees and use
\[
 \sum_{k=1}^{p}
       \frac{g^{p-k}}{(k-1)!(p-k)!}
       =\frac{(1+g)^{p-1}}{(p-1)!}
       \le\frac{h_G^p}{(p-1)!}.
\]
We obtain, for $n\ge2$,
\begin{equation}\label{eq:interacting-tree-bound}
 |T_n(F_1,\ldots,F_n)|
 \le e^{-\vartheta_0D}(n-2)!\alpha_G^{n-1}
                N_G(F_1;w_1)\prod_{i=2}^n\rho_G(F_i;w_i).
\end{equation}
For $n=1$ the bound is $N_G(F_1;1)$.
This also explains why the tree norms use
$(p-1)!$ rather than $p!$.

Apply the forest formula to the replicated
Gaussian expectation of $V_+^n$.
Different forest components contain disjoint
generator sets and their even factors commute,
so their expectations factor. Grouping by
component partitions gives
\[
 \int\dd\mu_G\,V_+^n
       =\sum_{\{B_1,\ldots,B_q\}}
                     \prod_{a=1}^qT_{|B_a|}(V_+,\ldots,V_+).
\]
The exponential formula for partitions
identifies the Gaussian integral with the
exponential of the connected tree sum.
The identity first holds as a formal series
in a scalar multiplying $V_+$. Bound
\eqref{eq:interacting-tree-bound} and
$\alpha_G\rho_G(V_+;1)\le1/2$ make the
connected series absolutely convergent up
to that scalar equal to one, so the formal
identity is an analytic identity there.
The tree terms are polynomial in the
activity coefficients, and their series
converges locally uniformly in $t$. Thus $W$
is analytic and its derivatives are obtained
term by term.

Two derivatives of the symmetric multilinear
tree expression either act on the same vertex
or on two different vertices. The first
terms, including $n=1$, sum to at most
$2N_G(\partial_c\partial_{c'}V_+;1)$.
For the second terms put the two derivatives
at vertices $1$ and $2$. The weights there
are $w_P,w_{P'}$, and those at all other
vertices are $w_\delta$.
Along the unique path from $1$ to $2$,
the distance $d(P,P')$ is at most the
sum of the two endpoint distances, all
edge distances, and the argument diameters
at intermediate vertices. Hence the weight
condition for \eqref{eq:interacting-tree-bound}
holds with $D=d(P,P')$. The tree sum leaves
$\alpha_GN_G(\partial_cV_+;w_P)
\rho_G(\partial_{c'}V_+;w_{P'})
e^{-\vartheta_0d(P,P')}$ times
$\sum_{n\ge2}(\alpha_G\rho_G(V_+;w_\delta))^{n-2}\le2$.
This proves \eqref{eq:interacting-tree-mixed}.
\end{proof}

For a function $f$ of the marks, let
$t^T$ keep the coordinates indexed by $T$
and set all others to zero, and define
its component with variable set $S$ by
\[
 f_S(t)=\sum_{T\subset S}(-1)^{|S\setminus T|}f(t^T).
\]
Then $f=\sum_Sf_S$, $f_\varnothing=f(0)$,
and $f_S$ vanishes if any variable of $S$
vanishes. The latter statement follows by
pairing the terms with and without that
variable. In particular a variable absent
from $f$ makes every component containing
it zero. Two coordinate differences can
be represented by two integrals of the
mixed derivative. Consequently
\begin{equation}\label{eq:interacting-Mobius-bound}
 \sup_{|t|_\infty\le R}|f_S(t)|
 \le2^{|S|-2}R^2
          \sup_{|t|_\infty\le R}|\partial_c\partial_{c'}f(t)|
 \qquad(c,c'\in S,\ c\ne c').
\end{equation}

\begin{proposition}\label{prop:interacting-mixed-logarithm}
For $|S|\ge2$, the normalized logarithm satisfies
\begin{equation}\label{eq:interacting-mixed-logarithm}
 \sup_{|t|_\infty\le t_x/2}|\mathcal G_S(t;\theta,A)|
                     \le\varpi_{\rm sep}\le\varpi_{\rm sep}^+,
\end{equation}
where
\begin{align}
 C_W^{\mathrm{mix}}&=3C_V^{\mathrm{ins}}h_G^2+2\alpha_G(C_V^{\mathrm{ins}})^2h_G^4
                             +4\exp(6)C_dM^3K_0,\notag\\
 \varpi_{\rm sep}&=2^{N_*}C_W^{\mathrm{mix}}p_\beta
                    e^{-\vartheta_0(d_{\rm sep}-4M-6)},\notag\\
 \varpi_{\rm sep}^+&=2^{N_*}C_W^{\mathrm{mix}}p_\beta
                    e^{-\vartheta_0(Mr_*/(8L)-4M-7)} .
                                  \label{eq:interacting-separation-errors}
\end{align}
\end{proposition}

\begin{proof}
The tree representation and uniqueness of the
normalized logarithm give
\[
 \mathcal G(t)=\sum_X\mathfrak P_*(X,\mathscr C(t))+W(t)-W(0).
\]
The component of a determinant activity is zero
unless $X$ meets every $P_c$, $c\in S$.
Such a polymer has
$d_M(X)\ge d_{\rm sep}/M-2\ge1$.
Use \eqref{eq:interacting-P-bound}, sum the
polymers anchored at one of at most $p_\beta$
cubes, and keep $e^{-2d_M(X)}$ from
$\kappa-\kappa_0\ge2$.
Its total contribution is at most
$2^{|S|}2\exp(4)C_dM^3K_0p_\beta
e^{-\vartheta_0d_{\rm sep}}$.

Apply \eqref{eq:interacting-tree-mixed}
with $P=P_c$, $P'=P_{c'}$ and
\eqref{eq:interacting-differentiated-norms}.
The bound for $\partial_c\partial_{c'}W$ is
\[
 q_x^2p_\beta\left[
       2C_V^{\mathrm{ins}}h_G^2e^{-\vartheta_0(d_{\rm sep}-4M-6)}
           +2\alpha_G(C_V^{\mathrm{ins}})^2h_G^4e^{-\vartheta_0d_{\rm sep}}\right].
\]
In \eqref{eq:interacting-Mobius-bound}, take
$R=t_x/2$ and use $q_xR=1$.
Add the determinant contribution and
$|S|\le N_*$ to obtain the first inequality.
The placement estimate
$d_{\rm sep}\ge Mr_*/(8L)-1$
proves the second.
When the lattice probe is zero,
$R_{\rm tm}=1$, $\mathcal G=0$ and the
first separation error is interpreted as zero.
\end{proof}

\subsection{The small-field representation and its complement}
\label{subsec:interacting-complement}

Proposition~\ref{prop:term-photon-tower} identifies $F_N$
with $(A_*;\zeta_j)_{j<k_*}$. Their scalar
coordinates have the form
$\zeta_j(\Upsilon)=\langle\phi_{j,\Upsilon},F^{\rm u}\rangle$
with orthonormal $\phi_{j,\Upsilon}$; each
coarse plaquette coordinate $(dA_*)(p)$ has
the form $\langle\phi_p,F^{\rm u}\rangle$ with
$\|\phi_p\|\le1$.
There are at most $2\ell_*^3L^{3(k_*-j)}$
coordinates at level $j$ and $3\ell_*^3$
coarse plaquettes. The stop Gaussian
$\dd\mathsf m_*$ is nondegenerate on the
real mean-zero coclosed coarse fields.
The row and count statements are
\eqref{eq:term-tower-row-norms} and
\eqref{eq:term-tower-coordinate-count}; the
Gaussian assertion is \eqref{eq:term-tower-Gaussian}.

Let $\mathfrak G$ require
$|\zeta_j(\Upsilon)|\le p_{0,j}$
and $|dA_*|_\infty\le p_*$, and let
$\mathfrak G^-$ decrease each threshold by one.
Put
\begin{align}
 \varpi_0&=18\ell_*^3e^{-p_{0*}/2}
               +6\ell_*^3e^{-(p_*-1)^2/(2C_0^2)},\notag\\
 t_2&=e\,\min\left\{N_*^{-1/2},\,C_\beta^{-1},
                                     (2\sqrt{N_*\mathbf D})^{-1}\right\},\notag\\
 \varpi_L&=\sqrt{\varpi_0}\,e^{\mathbf D/2},
 \qquad M_L=2+e^{\mathbf D/2},\notag\\
 \lambda_L&=(\varpi_L^{1/3}M_L^{2/3})^{1/N_*}
                              5(2N_*r/t_2)^4 .
                                     \label{eq:interacting-complement-constants}
\end{align}
A reciprocal with zero denominator is interpreted
as $+\infty$. The number $\mathbf D$ is the
uniform probe energy bound, so $v_\eta\le\mathbf D$.

\begin{proposition}\label{prop:interacting-complement}
Suppose the preceding conditions hold, and, for
a fixed $0<s<1$, set
$\bar\kappa_1=6e_*^{1/4-9\varepsilon}$.
Assume
\begin{equation}\label{eq:interacting-jet-conditions}
 2\varpi_{\rm sep}^+\le\bar\kappa_1^{N_*},\quad
 N_*\bar\kappa_1^{1-s}\le1/2,\quad
 \varpi_0\le1/4,\quad 8N_*t_2\le t_x .
\end{equation}
There is a finite positive measure $v_{\rm sf}$
on $[-\pi,\pi]^3\times\{\chi_*=1\}$, of
mass in $[1-\varpi_0,1]$, such that
\begin{equation}\label{eq:interacting-small-field-representation}
 \frac{\bar\nu^{a(t)}(\mathfrak G)}{\bar\nu(F_N)}
       =\int e^{\mathcal G(t;\theta,A)}\,v_{\rm sf}(\dd\theta\,\dd A).
\end{equation}
The normalized complement
$L(t)=\bar\nu^{a(t)}(F_N\setminus\mathfrak G)/\bar\nu(F_N)$
is entire. Its fourth jet in each variable has
coefficient norm at radius $r$ at most
$\lambda_L^{N_*}$.
\end{proposition}

\begin{proof}
First, the union bound and
\eqref{eq:interacting-zero-flux-Gaussian} give
\begin{equation}\label{eq:interacting-small-field-tail}
 \bar\mu(F_N\setminus\mathfrak G^-)\le\varpi_0 .
\end{equation}
For $i=k_*-j\ge1$, write
$v_j=v_*+i(\log L)/2$.
The scale conditions
$v_*\ge8$, $p_0v_*^{p_0-1}\ge28$ give
$p_{0,j}\ge p_{0*}+14i\log L$.
Since $(p_{0,j}-1)^2\ge p_{0,j}$,
the individual tail is at most
$2e^{-p_{0*}/2}L^{-7i}$.
Multiply by the coordinate count
$2\ell_*^3L^{3i}$ and sum $L^{-4i}$.
The $3\ell_*^3$ coarse plaquettes give the
second term in \eqref{eq:interacting-complement-constants}.
These estimates yield \eqref{eq:interacting-small-field-tail},
with room in the coefficient $18$.

The exact density identity defines
\begin{equation}\label{eq:interacting-small-field-measure}
 v_{\rm sf}(\dd\theta\,\dd A)=
   \frac{\chi_*(A)Y_*^{\rm x}(A+c_\theta;0)}
             {\bar\nu(F_N)}\,\frac{\dd\theta}{(2\pi)^3}
                                      \dd\mathsf m_*(A).
\end{equation}
We check positivity of this density, which
does not follow from the formal Grassmann
representation. For real marks, the original
fine law $\nu^{\theta,a(t)}$ is nonnegative.
The identity \eqref{eq:term-Borel-observable} says, for every
nonnegative bounded Borel $\Phi$,
\[
 \int\Phi(A)\chi_*(A)Y_*^{\rm x}(A+c_\theta;\mathscr C(t))
                \,\dd\mathsf m_*(A)
       =\nu^{\theta,a(t)}(\mathbf1_{\mathfrak G}\Phi(A_*))\ge0.
\]
The terminal density is continuous. On the
interior of the compact convex set
$\{\chi_*=1\}=\{|dA|_\infty\le p_*\}$, a
nonzero imaginary part or negative real
part would persist on a ball of positive
Gaussian measure, contradicting this identity.
At the boundary, approach along $uA$,
$u\uparrow1$; zero is an interior point.
Thus the density is nonnegative everywhere
on this set. Proposition~\ref{prop:interacting-factor}
makes it nonzero there, hence strictly positive.
This reasoning also proves
$\mathcal G(t;\theta,A)\in\R$ for real $t$:
its exponential is a positive ratio and its
continuous value at zero is zero.

Integration of the exact terminal identity over
$\theta$, followed by
\eqref{eq:interacting-terminal-ratio}, proves
\eqref{eq:interacting-small-field-representation}.
At $t=0$ the mass is $\bar\mu(\mathfrak G)$,
which lies in $[1-\varpi_0,1]$.
All integrals here are well-defined Borel
integrals: the density is continuous on
the compact support and zero elsewhere.

For real $t=ez$ with $|t|_\infty\le t_2$,
put $G_z=\sum_cz_cd\widehat{\mathfrak a}_c$.
Disjoint supports imply
$\|G_z\|^2=v_\eta|z|_2^2\le1/4$.
The real source shift gives
\[
 L(ez)=e^{-v_\eta|z|_2^2/2}
 \bar\mu\left(\mathbf1_{(F_N\setminus\mathfrak G)+G_z}
                              e^{\langle G_z,F^{\rm u}\rangle}\right).
\]
Every photon coordinate changes by at most
$\|G_z\|\le1$ under the shift. Thus the
shifted complement is contained in
$F_N\setminus\mathfrak G^-$.
Cauchy--Schwarz, \eqref{eq:interacting-small-field-tail},
and Gaussian domination at $2G_z$ give
\begin{equation}\label{eq:interacting-real-complement}
 0\le L(t)\le\sqrt{\varpi_0}\,
          e^{v_\eta|z|_2^2/2}\le\varpi_L .
\end{equation}
This is the only place where the lowered
thresholds of $\mathfrak G^-$ are needed.

The source-shift identity also gives, for complex marks,
\[
 |\mathring\Phi(t)|\le
       \exp\left(\frac{v_\eta}{2e^2}
                                      \sum_c(\operatorname{Im}t_c)^2\right).
\]
It follows directly by bounding the expectation
of the real exponential and canceling its
real quadratic normalization.
On $|t|_\infty\le t_2$ this is at most
$e^{\mathbf D/2}$.
The small-field integral has modulus at most
$2$. Indeed, evenness and the double zero of
$\ell_0$ give $|\ell_0(t)|\le B_0|t|^2/t_x^2$,
while the zero of $w_c$ gives
$|w_c(t)|\le(3/8)|t|/t_x$.
Their sum over cells at
$t_2\le t_x/(8N_*)$ is at most
$1/(256N_*)+3/64$. The remaining logarithm
components sum to at most
$2^{N_*}\varpi_{\rm sep}^+
\le(2\bar\kappa_1)^{N_*}/2\le N_*^{-N_*}/2$,
by \eqref{eq:interacting-jet-conditions}.
Thus $|\mathcal G|\le1/2$ and
$|e^{\mathcal G}|\le2$.
Since $L$ is the difference between the
full generating function and the small-field
integral, $|L|\le M_L$ on that polydisc.
It is entire because the finite determinant
density is entire and is bounded on compact
mark sets uniformly in the real photon
field and real twist. The modulus bound
in the external-response section supplies
the required Gaussian integrable majorant
also for either Borel event.

We use the following two-constants estimate.
Let $n\ge1$ be an integer. If a holomorphic function
on a polydisc in $\C^n$ of radius $R$ has modulus at most $M$, and
at most $\delta\le M$ on its real points,
then its modulus is at most $\delta^{1/3}M^{2/3}$
on the polydisc of radius $R/(2n)$.
On the upper
half of the unit disc the harmonic function
$\omega(z)=\frac2\pi\arg[(1+z)/(1-z)]$
is zero on the diameter and one on the
semicircle. Comparing $\log|f(R'z)|$
with $(1-\omega)\log\delta+\omega\log M$,
for $R'<R$, the maximum principle gives
the same inequality in the half disc.
To handle its endpoints, subtract
$\epsilon\log[4/|1-z^2|]$ before applying
the principle, then let $\epsilon\downarrow0$.
Its harmonic conjugate gives a holomorphic
function to which the modulus principle
applies. On $|z|\le q$, $\omega(z)\le
(4/\pi)\arctan q$; reflection treats the
lower half disc. Let $R'\uparrow R$.
Apply this one-variable estimate successively
with $q=1/(2n)$. The remaining exponent of
$\delta$ is at least
$(1-\frac4\pi\arctan(1/(2n)))^n
\ge1-2/\pi>1/3$, proving the assertion.
The case $\delta=0$ follows by the identity
theorem, and $M=0$ is immediate.

Apply it to $L$, with $n=N_*$ and $R=t_2$.
Cauchy's coefficient bound gives
\[
 |L_\alpha|\le\varpi_L^{1/3}M_L^{2/3}(2N_*/t_2)^{|\alpha|}.
\]
In the fourth jet, every exponent is at most
four. Since $y=2N_*r/t_2\ge1$,
\[
 \sum_{\alpha_c\le4}|L_\alpha|r^{|\alpha|}
   \le\varpi_L^{1/3}M_L^{2/3}
                 (1+y+\cdots+y^4)^{N_*}
   \le\lambda_L^{N_*}.
\]
Reality on real marks gives real jet coefficients.
\end{proof}

\subsection{The fourth-jet error}
\label{subsec:interacting-fourth-jet}

Work in the finite ring
$\mathcal J=\C[t_c:c\in\mathcal C_*]/(t_c^5:c\in\mathcal C_*)$,
with coefficient norm
$\|P\|_r=\sum_{\alpha_c\le4}|P_\alpha|r^{|\alpha|}$.
This norm is submultiplicative, since discarding
terms with an exponent at least five can only
decrease the coefficient sum.
The jet map is a ring homomorphism; it
keeps degree at most four in each variable.
For $S=\{c\}$ set $k_S=e^{w_c}-1$ in this
ring, and for $|S|\ge2$ set
$k_S=e^{\mathcal G_S}-1$.
Every monomial of $k_S$ uses every variable
of $S$ and no other variable.

For a component $f_S$ bounded on a polydisc
of radius $R$, Cauchy's estimate at a smaller
radius $R'\le R/2$ gives
\[
 \|f_S\|_{R'}\le\sup|f_S|
             \left(\sum_{j=1}^4(R'/R)^j\right)^{|S|}
                         \le\sup|f_S|.
\]
Its monomials contain each variable of $S$
because it vanishes on every corresponding
coordinate hyperplane. Use this with
$R=t_x/2$, $R'=r$, and use the one-cell
estimate on radius $t_x$, to obtain
\begin{equation}\label{eq:interacting-error-jets}
 \|k_{\{c\}}\|_r\le\bar\kappa_1,\qquad
 \|k_S\|_r\le2\varpi_{\rm sep}^+
                       \le\bar\kappa_1^{N_*}\quad(|S|\ge2).
\end{equation}
We used $\|w_c\|_r\le3e_*^{1/4-9\varepsilon}\le3/8$
and $\|e^P-1\|_r\le2\|P\|_r$ for $\|P\|_r\le1/2$.
The coefficients are continuous in $(\theta,A)$:
express the logarithm coefficients by Cauchy
integrals on a fixed small torus, and then
take the finite polynomial expressions in
those coefficients. They are real because
the logarithm is real at real marks.

In $\mathcal J$ the normalized generating
function has the exact representation
\begin{equation}\label{eq:interacting-jet-representation}
 \mathring\Phi=
   \prod_c\phi_0(t_c)
       \int\prod_{\varnothing\ne S\subset\mathcal C_*}
                         (1+k_S)\,v_{\rm sf}(\dd\theta\,\dd A)
       +L^{[4]}.
\end{equation}
To justify taking the jet inside the integral,
use joint continuity, holomorphy, and
$|e^{\mathcal G}|\le2$ on a fixed
polydisc of radius at most $t_2$.
Cauchy's coefficient formulas and Fubini
then commute every required coefficient
with the finite measure integral.
Decompose $\mathcal G$ into its variable-set
components; its singleton components are
$\ell_0(t_c)+w_c(t_c)$.
Exponentiation, the jet homomorphism,
and $e^{\ell_0}=N_0$ prove the representation.

\begin{theorem}\label{thm:interacting-reference-error}
Assume the terminal inputs and conditions above,
and $\lambda_L\le\bar\kappa_1^s$.
For $m=2,4$,
\begin{equation}\label{eq:interacting-Theta-bound}
 \Theta_m^{(0)}\le
        m!r^{-m}\bigl(2\exp(1)\bar\kappa_1^s+\lambda_L\bigr)
                       \le8m!r^{-m}\bar\kappa_1^s.
\end{equation}
With
$\mathfrak e=s(1/4-9\varepsilon)-12\varepsilon>0$,
\begin{align}
 \Theta_4^{(0)}
    &\le3\cdot10^5 C_b^4e_*^{\mathfrak e}
           \le3\cdot10^5C_b^4h^{1/8}e^{\mathfrak e},\notag\\
 (1+|c_2|)\Theta_2^{(0)}
    &\le3\cdot10^4C_b^4e_*^{\mathfrak e}
           \le3\cdot10^4C_b^4h^{1/8}e^{\mathfrak e}.
                                    \label{eq:interacting-final-terminal-error}
\end{align}
The constants are independent of the ultraviolet
cutoff, the twist, and the admitted placement.
\end{theorem}

\begin{proof}
On the one-variable jet space define
$\ell_c(t_c^j)=m!\mathbf1_{\{j=m\}}-c_m\mathbf1_{\{j=0\}}$
and let $\ell=\bigotimes_c\ell_c$.
Then $(\Theta_m^{(0)})^{N_*}=|\ell(\mathring\Phi)|$.
The free jet satisfies $\ell_c(\phi_0)=0$.
From \eqref{eq:interacting-free-jet},
$|\ell(P)|\le(m!r^{-m})^{N_*}\|P\|_r$.
These statements follow by evaluating the
functionals on monomials with every exponent
in $\{0,m\}$.

Expand the product in
\eqref{eq:interacting-jet-representation} over
families $\mathcal Y$ of nonempty cell sets.
If $\mathcal Y$ fails to cover a cell $c$,
its error factors do not contain $t_c$;
the factor $\phi_0(t_c)$ is therefore
annihilated by $\ell_c$. Only covering
families contribute.
Assign cost one to each singleton and
cost $N_*$ to each larger set, and let
$a(\mathcal Y)$ be the sum of these costs.
A covering family has $a(\mathcal Y)\ge N_*$:
either it contains a larger set or it
contains every singleton.
By \eqref{eq:interacting-error-jets} its
coefficient norm is at most
$\bar\kappa_1^{a(\mathcal Y)}$.
The free factors have norm at most $2^{N_*}$.
Keep $\bar\kappa_1^{sN_*}$ from the cost
and sum the rest over all families. This
sum is
\[
 (1+\bar\kappa_1^{1-s})^{N_*}
  (1+\bar\kappa_1^{(1-s)N_*})^{2^{N_*}-1-N_*}
 \le\exp\left(N_*\bar\kappa_1^{1-s}
                  +2^{N_*}\bar\kappa_1^{(1-s)N_*}\right)
 \le\exp(1).
\]
The last step uses $N_*\bar\kappa_1^{1-s}\le1/2$.
The positive small-field measure has mass at
most one. Its total contribution is thus
bounded by
$(m!r^{-m})^{N_*}\exp(1)(2\bar\kappa_1^s)^{N_*}$.
The complement contributes at most
$(m!r^{-m}\lambda_L)^{N_*}$.
Taking the $N_*$th root proves
\eqref{eq:interacting-Theta-bound}.

Since $r^{-4}=256C_b^4e_*^{-12\varepsilon}$
and $\bar\kappa_1^s\le6e_*^{s(1/4-9\varepsilon)}$,
the $m=4$ bound has coefficient at most
$8\cdot24\cdot1536<3\cdot10^5$.
For $m=2$, $1+|c_2|\le(17/16)r^{-2}$,
so its coefficient is at most
$17\cdot1536<3\cdot10^4$.
Finally $0<\mathfrak e<1/4$ and
$e_*\le e\sqrt h$ imply
$e_*^{\mathfrak e}\le h^{1/8}e^{\mathfrak e}$.
The scale and probe bounds entering this
proof are uniform in the stated choices.
\end{proof}

\subsection{A positive charge threshold}
\label{subsec:interacting-positive-threshold}

We now verify that the response conditions admit
a common interval of charges, independent of the
ultraviolet cutoff. In addition to
\eqref{eq:interacting-factor-conditions},
\eqref{eq:interacting-small-complex-mass},
\eqref{eq:interacting-one-cell-conditions},
\eqref{eq:interacting-tree-smallness}, and
\eqref{eq:interacting-jet-conditions}, require
$r_*>31L+16$ and
\begin{gather}
 Mr_*/(8L)-1\ge6M+8,\qquad
 \lambda_L\le\bar\kappa_1^s,
                                  \label{eq:interacting-last-response-conditions}\\
 D_\beta^{\rm buf}/2+1/2\le3\ell_c/8,\qquad
 D_\beta^{\rm buf}+3\le\ell_c .
                                  \label{eq:interacting-placement-conditions}
\end{gather}
The construction of Section~\ref{sec:rg-flow}
fixes an odd $L\ge13$, a power
$M$ of $L$ with $M\ge L^2$, a localization
power $r_{\rm loc}\ge2$, and cutoff powers
$p_0\ge2$, $p>2p_0$. It also fixes
$0<\varepsilon\le1/200$ and $0<s<1$ such that
\begin{equation}\label{eq:interacting-response-exponents}
 \mathfrak e=s(1/4-9\varepsilon)-12\varepsilon,\qquad
 \sigma=\mathfrak e-\frac{2a_L}{N_*}>0,\qquad
 a_L=\frac{L^3-1}{(L-1)^3}.
\end{equation}
Only $\mathfrak e$ enters the terminal estimate;
the loss $2a_L/N_*$ belongs to the subsequent
finite-volume transfer.

\begin{proposition}\label{prop:interacting-positive-threshold}
For the fixed construction constants, $\bar m$,
and probe $\beta$, there is $e_\sharp(\beta)>0$
such that all response conditions hold for
$0<e\le e_\sharp(\beta)$ and every
$N\ge N_{\rm ref}(e)$. The threshold depends
on the probe only through
$C_\beta,\mathbf D,D_\beta^{\rm buf}$.
It is independent of $N$, the twist, and the
admitted placement. All centered placements
are admitted in this range.
\end{proposition}

\begin{proof}
Write
$v_K=-\log(e\sqrt h)$ and
$r_K=\min\{L^i:L^i\ge v_K^{r_{\rm loc}}\}$.
Then $R^*(e)=hMr_K/L$ and
$\ell_c=R^*(e)/(2L)$.
Proposition~\ref{prop:rgfl-scales} and
\eqref{eq:rgfl-reference-threshold} give
a finite $N_{\rm ref}(e)$.

For $N\ge N_{\rm ref}(e)$ the stop level
is $K-1$ or $K-2$ by
\eqref{eq:rgfl-reference-stop}. In particular its
scales are independent of $N$ and satisfy
\begin{gather}
 v_K\le v_*,\quad e_*=e^{-v_*},\quad
 a_*\in\{h/L,h/L^2\},\notag\\
 v_*^{r_{\rm loc}}\le r_*,\quad
 Mv_*^{r_{\rm loc}}\le\ell_*<LMv_*^{r_{\rm loc}},
 \quad p_{0*}=v_*^{p_0},\quad p_*=v_*^p,\notag\\
 p_\beta\le p_\beta^+:=
                (D_\beta^{\rm buf}L^2/h+2M+4)^3,\qquad
 \ell_c\ge hMv_K^{r_{\rm loc}}/(2L^2).
                                  \label{eq:interacting-threshold-scales}
\end{gather}
The strict wrapping condition $r_*>31L+16$
already follows from \eqref{eq:rgfl-scale-ranges}
whenever $e\sqrt h\le e_1$.
Thus every response condition is a condition
on $e$ alone once the construction and probe
are fixed.

A bound $A v^d e^{-\gamma v}$
with $A,d\ge0$ and $\gamma>0$ tends to zero,
uniformly on a sufficiently large half-line.
Split the exponential in half and use
$v^de^{-\gamma v/2}\le(2d/(\gamma\exp(1)))^d$.
Likewise, $cv^2$ exceeds $bv+\log^+A$
on a sufficiently large half-line for
$c>0$. We apply these two inequalities to
the following groups of conditions.

The local tree smallness and terminal factor
conditions follow from exponential decay
with positive rates
$\varepsilon$, $2\varepsilon$,
$1/4-\varepsilon$, and $1/4-9\varepsilon$;
the volume factors are bounded by
$(LM)^3v_*^{3r_{\rm loc}}$.
For $\bar r$, use
$1/4-7\varepsilon>2\varepsilon$.
The complex mass and coarse determinant
conditions similarly have rate $2\varepsilon$.
The support count is bounded by $p_\beta^+$.
All three summands of $B_0$ tend to zero:
their rates are $4\varepsilon$, $2\varepsilon$,
and the quadratic decay $e^{-r_*}\le e^{-v_*^2}$.
Also $\tau_s\ge c_\tau Mv_*^2$, so both
$\tau_s\ge1$ and
$52B_1e^{-\tau_s}\le e_*^{1/4-9\varepsilon}$
hold once $v_*$ is large.

The separation condition follows from
$r_*\ge v_*^2$. Its error obeys
\[
 2\varpi_{\rm sep}^+
 \le2^{N_*+1}C_W^{\mathrm{mix}}p_\beta^+
          e^{\vartheta_0(4M+7)}
          e^{-\vartheta_0Mv_*^2/(8L)} .
\]
Since $\bar\kappa_1\ge e^{-v_*/4}$,
quadratic decay makes this at most
$\bar\kappa_1^{N_*}$ for all sufficiently
large $v_*$. The family summation condition
$N_*\bar\kappa_1^{1-s}\le1/2$
uses the positive rate
$(1-s)(1/4-9\varepsilon)$.

For the photon complement, $p>2p_0$ implies
$(p_*-1)^2\ge C_0^2p_{0*}$ once $v_*$
is large. Therefore
\begin{equation}\label{eq:interacting-complement-threshold}
 \varpi_0\le24(LM)^3v_*^{3r_{\rm loc}}e^{-v_*^2/2}.
\end{equation}
This is eventually at most $1/4$.
Moreover $t_2/e$ is fixed and positive, while
$e$ is between $(L/h)^{1/2}e_*$
and $Lh^{-1/2}e_*$. Hence
$8N_*t_2/t_x$ is bounded by a fixed
constant times $e_*^{1-3\varepsilon}$,
which tends to zero.

The potentially largest loss is the fourth
derivative radius in $\lambda_L$.
From \eqref{eq:interacting-complement-threshold},
the definition of $\lambda_L$, and
$r/t_2\le C e^{v_*}$, we obtain
\begin{equation}\label{eq:interacting-lambda-threshold}
 \lambda_L\le A_8
   \exp\left(-\frac{v_*^2}{12N_*}
                     +\left(\frac{r_{\rm loc}}{2N_*}+4\right)v_*\right)
\end{equation}
for a finite $A_8>0$ depending only on the
fixed constants and probe bounds.
For clarity one can take
\[
 A_8=[24(LM)^3]^{1/(6N_*)}
       e^{\mathbf D/(6N_*)}M_L^{2/(3N_*)}
       5\left[
        \frac{N_*\sqrt h}{2C_b\sqrt L}(t_2/e)^{-1}\right]^4 .
\]
If
$v_*^2/(12N_*)\ge
(r_{\rm loc}/(2N_*)+17/4)v_*+\log^+A_8$,
then $\lambda_L\le e^{-v_*/4}\le\bar\kappa_1^s$.
Thus the complex continuation and the
fourth derivatives preserve a positive
small-charge interval.
Finally the two placement inequalities
hold when
$\ell_c\ge(4/3)D_\beta^{\rm buf}+3$;
this follows from
\eqref{eq:interacting-threshold-scales} as $v_K\to\infty$.

Choose a finite $V$ beyond all the half-lines
just specified and put
$e_V=h^{-1/2}\min(e_1,e^{-V})>0$.
Every $0<e\le e_V$ has $v_*\ge v_K\ge V$
and satisfies every response condition.
Define $e_\sharp$ as one half of the
supremum of the charges $e'\le e_1h^{-1/2}$
for which all conditions hold for every
$0<e\le e'$. Its defining set contains
$e_V$, so $e_\sharp\ge e_V/2>0$.
For $e\le e_\sharp$, a member of the
defining set exceeds $e_\sharp$, which
ensures that all inequalities hold at $e$
even if the scale functions have jumps.
This avoids assuming monotonicity of the
rounded stop scales.

The scale identities depend on $e,h$ and
the fixed construction constants; the probe
enters through
$C_\beta,\mathbf D,D_\beta^{\rm buf}$.
They do not involve the twist or placement,
and they are unchanged by the coordinate
reflections used to form the cell probes.
Centered placements satisfy the admission
inequality by
\eqref{eq:interacting-placement-conditions},
since their centers differ from the cell
center by at most $\eta/2\le1/2$.
The same conditions give
$D_\beta^{\rm buf}+3\eta\le\ell_c$ and
$R^*(e)\ge D_\beta+1$.
\end{proof}


\section{Choice of the construction constants}\label{sec:ordered-constants}

The estimates impose conditions on fixed constants and on the
running charge. The constants are chosen before the charge, mesh,
volume and twist. We first choose the polymer decay and the two
geometric scales; the charge bounds are then imposed for every
smaller charge. The local inverse constants are fixed before $M$,
whereas the weakened kernel constants may depend on $M$.
Neither group depends on the running charge. We verify this order
before imposing the small-charge conditions.

\begin{proposition}\label{prop:ordered-constants}
Fix $0<\rho<1$ and the standard chamber weights.
The constants of the photon and fermion primitives,
small-field step, cover comparison and tuned flow can
be chosen in the order described below, with positive
charge ceilings $e_0,e_1$ independent of the mesh,
volume, twist and reference mass.
For every $0<\bar m\le1$, the physical ceiling $e_G$
in \eqref{eq:constants-physical-charge} is positive.
The terminal conditions \eqref{eq:term-stop-conditions}
hold on the top and reference tori throughout
$0<e\le e_G$.
For a fixed compact real probe with $b(dd^*\beta)>0$,
there is a positive $e_\beta\le e_G$ for which
\eqref{eq:transfer-small-Wick} holds with
$\sigma>0$ and $C_fe^\sigma\le b(dd^*\beta)/2$.
\end{proposition}

\subsection{The geometric scales and the positive exponent}

\noindent\emph{Proof.} Set $\bar\varepsilon=1/200$, $c_X=54$,
$\kappa_0=54\log72$, $K_0=72^{54}/36$,
$c_{\mathrm{rb}}=54K_0$ and $c_{\mathrm{fl}}=61c_XK_0$.
Choose
\[
 1-\bar\varepsilon<\alpha<1,\qquad
 \kappa\ge7\kappa_0+7.
\]
An odd integer $L\ge5$ is required to satisfy
\begin{align}
 L^{1/8-3\bar\varepsilon/2}&\ge2400,\qquad
 c_{\mathrm{fl}}L^{-3/8+6\bar\varepsilon}\le1/16,\notag\\
 L(\kappa-6\kappa_0-4)-1&\ge\kappa,\notag\\
 c_{\mathrm{rb}}L^3
       \exp\bigl(-L(\kappa-\kappa_0-1)+\kappa\bigr)
                                                   &\le1/2.
                                      \label{eq:constants-L}
\end{align}
These are simultaneous lower-bound conditions on $L$.
The powers $1/8-3\bar\varepsilon/2$ and
$3/8-6\bar\varepsilon$ are positive, as are
$\kappa-6\kappa_0-4$ and $\kappa-\kappa_0-1$.
The last expression tends to zero as $L\to\infty$.
Thus every sufficiently large odd $L$ satisfies
\eqref{eq:constants-L}.

After $L$ is fixed, the primitive construction fixes
$a'$, $C_V$, $\gamma_w$ and $\gamma_P$, with $a',C_V>0$
and $\gamma_w,\gamma_P>0$. Put
$N_c=(2L+1)^3$ and
$\widehat\kappa=N_c\kappa+\log(6c_X)>0$.
The photon and fermion
constants have no dependence on $M$ at this stage.
Put $\alpha_0=1/54$, $\kappa'=\kappa_w+\kappa_0$ and
\[
 M_F=\max\{4L+1,\ (108\pi C_Va')^2,\
                         108\pi C_Va'\exp(27\kappa')\}.
\]
Choose $M$ to be a power of $L$ satisfying
\begin{equation}\label{eq:constants-M}
 \begin{split}
 M&\ge L^2,\qquad M>4L,\qquad \gamma_wM\ge16,\\
 M^{\alpha_0}&\ge1+\exp(\widehat\kappa),\qquad
 M\ge M_F,\qquad 804\alpha_0\log M\le\gamma_PM.
 \end{split}
\end{equation}
All these conditions hold for every sufficiently large
power of $L$. For the last one, $\log x\le\sqrt x$
on $x\ge1$ gives the sufficient lower bound
$M\ge(804\alpha_0/\gamma_P)^2$.
The definition of $M_F$ gives the two fermion
conditions in \eqref{eq:weakened-walk-contraction},
because $27\alpha_0=1/2$.
The last inequality gives \eqref{eq:photon-weakening-scale}.
The first four give the geometric separation, decay
and parameter-radius conditions used in localization.

The exponents can be fixed independently of the charge.
Take $r=2$, $p_0=9$, $p=19$ and
$1-\alpha<\varepsilon\le\bar\varepsilon$, and take $s=1/2$.
Then $p_0>3r+2$ and $p>2p_0$, while
\[
 \mathfrak e=s(1/4-9\varepsilon)-12\varepsilon
                      \ge17/400.
\]
For $N_*=8L^3$ and $a_L=(L^3-1)/(L-1)^3$,
\[
 \frac{2a_L}{N_*}
     <\frac1{4(L-1)^3}\le\frac1{256}<\frac{17}{400}.
\]
It follows that $\sigma=\mathfrak e-2a_L/N_*>0$.
In particular the positive exponent is fixed before
the charge is chosen.

\subsection{Calibration of the primitive and terminal bounds}

The chamber coercivity, heavy-sector gap and local inverse
construction fix $a'$ and $\gamma^F>0$ before $M$ is chosen;
see \eqref{eq:local-inverse-bounds} and
\eqref{eq:global-propagator-bounds}. The photon construction
likewise fixes $C_P^{\mathrm{ph}}$ and $\gamma_P>0$ before $M$;
its weakened families are considered only after
\eqref{eq:photon-weakening-scale} holds.
The sums $c(\gamma)=\sum_{z\in\Z^3}
\exp(-\gamma|z|_\infty)$ are finite for $\gamma>0$.
Thus the constants used to choose $M$ in
\eqref{eq:constants-M} have already been fixed.

Keep $b=1$, $b_-=1-L^{-1}$ and $b_+=1$ from the inverse
construction; these are already fixed before $M$.
The coarse resolvent and its auxiliary block matrix give
\[
 C_{\mathrm{By}}=(1+2\exp(1))\sqrt2\max(1,b_-^{-1}),
 \qquad
 C_\Gamma=\frac{2(1+2\exp(1))}{b_-}
                  +128\exp(1)(1+2\exp(1))^2a'.
\]
The latter is the constant in
\eqref{eq:weakened-Gamma-bound}. Define
\begin{equation}\label{eq:constants-fermion-calibration}
 \begin{split}
 C_1=\max\{&
 304\exp(1)\max(1,b)a'M^{27\alpha_0},
 C_\Gamma M^{27\alpha_0},\\
 &32\exp(1)\max(1,b^2)a'\exp(27\kappa')\},\\
 C^F=2\max\{&
 C_1,C_{\mathrm{By}},304\exp(1)\max(1,b)a',
 C_\Gamma,\exp(1)b^2(32a')^2\}.
 \end{split}
\end{equation}
The factor two converts Euclidean spinor bounds to
component bounds. The walk and coarse-family estimates
\eqref{eq:weakened-propagator-bound}--\eqref{eq:polymer-propagator-bound}
and \eqref{eq:weakened-H-bound}--\eqref{eq:weakened-Gamma-bound}
are bounded by $C^F$ in this convention. The last entry
also bounds the compressed quadratic kernel.
At level zero the normalized kernel has the same bound,
and $M^{8\alpha_0}\le M^{27\alpha_0}$ bounds its diagonal
field multiplier.
The unweakened all-pair field bound
\eqref{eq:interacting-global-field-rows} is at most
$608\exp(1)\max(1,b)a'\le C^F$; at zero its constant is
at most four, also below $C^F$.

Take the common primitive constants
\[
 C^P=116C_P^{\mathrm{ph}}M^{402\alpha_0},\qquad
 C_w^{(1)}=\max\{C^F,C^P,4\exp(2\gamma_w)M^{27\alpha_0}\},
 \qquad C_w=64K_0C_w^{(1)}.
\]
These are precisely \eqref{eq:rgcov-top-constants}.
The top systems and their projections therefore satisfy
the same bounds at every required level and side.
The determinant, photon-domain and localization constants
are given in Sections~\ref{subsec:local-determinant-terms},
\ref{subsec:photon-domain-bounds} and
\ref{subsec:rgnl-domains}. In particular,
\[
 \begin{aligned}
 c_C&=3L^{9/2},\qquad C_W=28C_w^2c_\gamma,\\
 C_d=16c_XC_w
       &\max\{\tfrac\pi2b\exp(1)(1+2\exp(1)),
                                      2+58\exp(1)+b\exp(1)\}.
 \end{aligned}
\]
The kernel-sum constant $c_\gamma$ is bounded by
$2L^3\exp((2L+1)\gamma_w)c(\gamma_w)$;
the local-gauge constant $C_B$ depends only on
$C_P^{\mathrm{ph}},\gamma_P,L,M$.
All these constants are finite after $M$ is fixed.

For the terminal norms fix
\begin{equation}\label{eq:constants-terminal-calibration}
 C_0=\max\{C_P^{\mathrm{ph}}c_{\mathrm{bd}}(\gamma_P),
                 C^F,4b_++16C^Fc(\gamma^F)\},
 \qquad c_{\mathrm{bd}}(\gamma)=3\exp(3\gamma)c(\gamma).
\end{equation}
The photon row sums, both field-map row norms and the
terminal Dirac form in
Section~\ref{subsec:interacting-terminal-norms} are then
bounded by this $C_0$. In particular $C_0\ge1$.
The polymer sum has $1\le c_\kappa\le K_0$.
The free Bloch construction fixes the twist constants
$C_\Sigma<\infty$ and $\gamma_\Sigma>0$ independently of
the running charge, mass in its prescribed range, mesh
and side; see \eqref{eq:term-penalty-twist-bound}.

The differentiated-response constants are now fixed by
the terminal inverse and field-map estimates:
$\mathcal S=C^F(24/\gamma^F)^3$,
$C_K=226\mathcal S^2$, $C_K'=1328\mathcal S^2$ and
$C_H=203M(C^F)^2$. The real covariance coefficients
$g,h_G,g_1,\alpha_G$ and its rate $\vartheta_0$ are
defined with \eqref{eq:interacting-Gaussian-covariance}.
The insertion and mixed-logarithm coefficients are
$C_V^{\mathrm{ins}}=10^3C_H^2(24/\gamma^F)^{12}K_0$ and
$C_W^{\mathrm{mix}}$ in \eqref{eq:interacting-separation-errors}.
Their inverse-decay rates are positive and their
coefficients finite because $b_->0$,
$\mu_-=L^{-3}>0$ and $\gamma^F>0$.
They depend on $L,M$ and the fixed chamber data;
they have no dependence on the later charge or probe.
The probe-dependent constants will be chosen last.

\subsection{Uniform running-charge bounds}

All coefficients of the small-charge conditions are
fixed after the preceding choices. The elementary
estimate used to meet them simultaneously is as follows.
For finitely many $A_i,d_i,\lambda_i$ with
$A_i,d_i\ge0$ and $\lambda_i>0$, and for $b>0$, put
\begin{equation}\label{eq:constants-tail-threshold}
 V(A,d,\lambda;b)
   =\max_i\frac2{\lambda_i}
       \log^+\left[
         \frac{nA_i}{b}
         \left(\frac{2d_i}{\lambda_i\exp(1)}\right)^{d_i}
       \right].
\end{equation}
Here $n$ is the number of terms, $\log^+0=0$
and a factor with $d_i=0$ is interpreted as one.
For all $v\ge V(A,d,\lambda;b)$ with $v>0$,
\begin{equation}\label{eq:constants-tail-bound}
 \sum_i A_iv^{d_i}\exp(-\lambda_iv)\le b.
\end{equation}
Indeed $v^d\exp(-\lambda v/2)$ has maximum
$(2d/(\lambda\exp(1)))^d$ for $d>0$.
For $d=0$ use $\exp(-\lambda v)\le\exp(-\lambda v/2)$.
The threshold makes each resulting term at most $b/n$.

The primitive domains, determinant, adjustment and nonlinear
step require \eqref{eq:photon-domain-small-charge},
\eqref{eq:rg-domain-charge-conditions},
\eqref{eq:rg-adjustment-charge-condition},
\eqref{eq:rg-contraction-parameters},
\eqref{eq:determinant-gain-conditions} and
\eqref{eq:rgnl-fixed-smallness}.
Their charge conditions are obtained from the following
finite families for $v=-\log e_k$.
Write $T(A,d,\lambda)=Av^d\exp(-\lambda v)$.
The polynomial-exponential conditions are
\begin{align}
 T(1,p_0,\varepsilon/4)&\le1,&
 T(1,0,5\varepsilon/2)&\le\tfrac12,\notag\\
 T(c_{\mathrm{fl}},0,\varepsilon)&\le1,&
 T(3M,0,\varepsilon)&\le1,\notag\\
 T(2C_w^2c_\gamma^2c_C,0,\varepsilon/4)&\le1,&
 T(16c_Xc_\gamma C_wC_T^2M^3,0,\varepsilon/2)&\le1,\notag\\
 T(96M^3,0,1/4)&\le1,&
 T(32C_wC_T,0,3\varepsilon/2)&\le1,\notag\\
 T(88M^3,0,1/4-6\varepsilon)&\le1,&
 T(112C_w^2c_\gamma,0,2\varepsilon)&\le3,\notag\\
 T(c_S,0,1/4-5\varepsilon)&\le1,&
 T(c_SL^3,0,\varepsilon)&\le1,\notag\\
 T(36c_X^2K_0^2,0,1/4-5\varepsilon)&\le1,&
 T(1,0,\varepsilon)&\le\tfrac14,\notag\\
 T(3M,0,1/4+\varepsilon)+T(3M,0,2\varepsilon)&\le1,\notag\\
 T(1,0,1/4)+T(1,0,\varepsilon)&\le1,\notag\\
 T(160L^{1+3\varepsilon/2}M,0,\varepsilon)&\le1,&
 T(90c_XLM,0,\varepsilon)&\le1,\notag\\
 T(9ML^\varepsilon,0,2\varepsilon)&\le1,\notag\\
 T(2M,1,\varepsilon)+T(32M,0,\varepsilon)&\le1,\notag\\
 T(8C_B,p,3/4-3\varepsilon)&\le1,&
 T(8,0,3/4-7\varepsilon/2)&\le1,\notag\\
 T(6L,0,\varepsilon)&\le1,&
 T(4C_Va',0,\varepsilon)&\le1,\notag\\
 T(4,0,1/4+2\varepsilon)+T(2,0,3\varepsilon)&\le1,\notag\\
 T(8c_{\mathrm{rb}}L^3C_dM^3,0,2\varepsilon)&\le1.
                                      \label{eq:constants-running-inventory}
\end{align}
Here $c_S=10^5c_X^5K_0^5\exp(6\kappa)$.
The inequality with $96M^3$ controls the first step;
the one with $C_T^2M^3$ controls the positive-level
fluctuation substitution. The conditions with $c_S$
and $36c_X^2K_0^2$ control the connected logarithm and its scaling.
The domain conditions occur with target $1/4$,
with rates $(1/4,\varepsilon)$, with coefficients
$8C_B,8$, and with coefficients $4,2$.
The remaining displayed charge hypotheses above are
individual inequalities of this inventory;
\eqref{eq:rgnl-cluster-smallness} follows from the
condition containing $c_S$.

Every decay rate in the inventory is positive for
$0<\varepsilon\le1/200$. The only additional
running-charge restrictions in this group are
\[
 v^{2p_0}/2-7v\ge\log12,\qquad v^2\ge10L,\qquad v\ge9.
\]
They hold respectively for $v\ge15$,
$v\ge\sqrt{10L}$ and $v\ge9$.
Taking the maximum of these thresholds and the finitely
many values in \eqref{eq:constants-tail-threshold}
gives a finite $V_0$. Choose a running-charge ceiling
$0<e_0<\exp(-V_0)$ so that every such condition
holds for every $v\ge-\log e_0$, equivalently for
every running charge in $(0,e_0]$.
The set of permitted ceilings is downward closed.

The terminal-scale conditions require a second ceiling
$e_1\le\min(e_0,1/3)$. Write $\ell(v)=LMv^r$.
They follow from
\begin{align}
 (c_{\mathrm{fl}}+2)\exp(-\varepsilon v)
                      &\le\tfrac1{16}(1-L^{-\varepsilon/2}),\notag\\
 8C_0^2\ell(v)^6(b_-^{-1}+L\ell(v))
                      &\le\exp(\varepsilon v/2),\notag\\
 4(2C_0+c_\kappa+5c_{\mathrm{fl}}+6)\ell(v)^3
                      &\le\exp(\varepsilon v),\notag\\
 \exp(-(1/4-9\varepsilon)v)&\le\tfrac18,\notag\\
 4C_0\ell(v)^2v^p&<\exp(\varepsilon v),\notag\\
 12\ell(v)^3\left(\exp(-v^{p_0}/2)
                    +\exp(-v^{2p}/(2C_0^2))\right)
                      &\le\exp(-v),\notag\\
 C_\Sigma\exp(-\gamma_\Sigma Mv^r)&\le\exp(-v),\notag\\
 L^2\exp(-2v)\ell(v)&\le1,\notag\\
 6N_*Mv^r\exp(-v)(v\exp(-v)+\tfrac12)&\le\pi^2.
                                      \label{eq:constants-terminal-inventory}
\end{align}
These are the scalar hypotheses of
\eqref{eq:rgfl-scale-smallness},
\eqref{eq:term-stop-conditions} and
\eqref{eq:transfer-tail-condition}, after substituting
\eqref{eq:rgfl-terminal-scales}.
Every polynomial-exponential rate here is positive for
$0<\varepsilon\le1/200$. The reference flux condition expands as
\[
          6N_*Mv^{r+1}\exp(-2v)+3N_*Mv^r\exp(-v)\le\pi^2.
\]
The superlinear tails reduce to these bounds:
for $v\ge\max(4,4C_0^2)$,
\[
 v-\tfrac12v^{p_0}\le-v,\qquad
 v-\frac{v^{2p}}{2C_0^2}\le-v.
\]
Hence
\[
 12\ell(v)^3
 \left[\exp\left(v-\tfrac12v^{p_0}\right)
       +\exp\left(v-\frac{v^{2p}}{2C_0^2}\right)\right]
                \le24(LM)^3v^{3r}\exp(-v).
\]
For the twist tail, if
$v\ge\max(1,2/(\gamma_\Sigma M),\log^+C_\Sigma)$,
then $C_\Sigma\exp(v-\gamma_\Sigma Mv^r)\le1$.
Put $\mathrm{lg}=\tfrac12\log L$.
The scale-ratio condition follows for $v\ge2r$:
\begin{equation}\label{eq:constants-scale-ratio}
 \left(1+\frac{3\mathrm{lg}}{v}\right)^r
        \le\exp\left(\frac{3r\mathrm{lg}}{v}\right)
        \le L^{3/4}<L.
\end{equation}
Together with $v\ge8$, $v\ge31L+17$ and
$p_0v^{p_0-1}\ge28$, these bounds give a finite
threshold $V_1$. Its maximum includes $2r$.
Let $\mathcal E_1$ consist of the numbers
$e'\in(0,\min(e_0,1/3)]$ for which the geometric
conditions and \eqref{eq:constants-terminal-inventory}
hold for every $v\ge-\log e'$. It contains
$\min(e_0,1/3,\exp(-V_1))>0$ and is downward closed.
Fix
\[
 e_1=\tfrac12\sup\mathcal E_1>0.
\]
There is an element of $\mathcal E_1$ larger than $e_1$;
downward closure therefore gives $e_1\in\mathcal E_1$.
Thus all terminal-scale conditions hold for every
$v\ge-\log e_1$.

\subsection{Physical and probe-dependent charge intervals}

Fix the reference mass and hence $h\ge1$. Set
\begin{equation}\label{eq:constants-physical-charge}
 e_G=
 \begin{cases}
 e_1,&h=1,\\
 \min\{e_1/\sqrt h,\,1/(\sqrt h\log h)\},&h>1.
 \end{cases}
\end{equation}
This is positive. For $0<e\le e_G$, every running charge
up to level $K$ is at most $e_1$, and
$e\sqrt h\log h\le1$.
The scale proposition therefore applies on the top and
reference tori. The terminal inventory gives
\eqref{eq:term-stop-conditions} at their respective stops,
and \eqref{eq:transfer-W3} holds on the reference torus.
The common pair of Theorem~\ref{thm:rgfl-common-pair}
is defined for $N\ge N_{\min}(R_{\mathrm{top}},e)$.
At lower cutoffs only its mass shift is assigned zero.
For $N\ge N_{\mathrm{ref}}(e)$ the same pair is
controlled on the reference torus. Its mass shift is real,
independent of the cutoff volume, and has absolute value
less than $1/8$.

Fix a compact real probe $\beta$. Its buffered support
diameter, energy and uniform rescaled norms are finite by
\eqref{eq:interacting-uniform-probe-finite}.
The differentiated inverse and interaction constants have
already been fixed; these probe constants are chosen last.
Proposition~\ref{prop:interacting-positive-threshold}
gives $e_\sharp(\beta)>0$ for which all its factor,
one-cell, tree, jet and placement conditions hold for
every $0<e\le e_\sharp(\beta)$, uniformly in
$N\ge N_{\mathrm{ref}}(e)$ and in the allowed placements.
Its proof uses $a_*\in\{h/L,h/L^2\}$ and
$v_*=v_K+O(1)$: polynomial factors have fixed degrees,
the separated-cell errors decay exponentially in $v_*^r$,
and the complement decays faster than the Cauchy-radius
loss. This gives one interval of charges for all cutoffs.

For $0<e\le\min(e_G,e_\sharp(\beta))$,
\eqref{eq:term-matter-comparison} and
\eqref{eq:interacting-final-terminal-error} supply the
two inputs in \eqref{eq:transfer-terminal-input}.
The probe placement inequalities imply
\eqref{eq:cells-fit} for $\eta\le1$.
The bounds for the free coefficients and the constant $C_U$ in
\eqref{eq:transfer-CU} are finite and independent of
$e,N,T$ and placement. They depend only on the fixed
construction and probe.

Suppose now that $b(dd^*\beta)>0$ and put
$C_f=8C_U+3C_U^2$.
Choose
\begin{equation}\label{eq:constants-probe-charge}
 0<e_\beta\le
 \min\{e_G,e_\sharp(\beta),1,
                         (b(dd^*\beta)/(2C_f))^{1/\sigma}\}.
\end{equation}
The minimum is positive, since $\sigma>0$ and $C_f>0$.
Proposition~\ref{prop:transfer-UV-inputs} holds throughout
$(0,e_\beta]$ and gives the cutoff hypotheses of
Proposition~\ref{prop:UV-component-cumulant}.
In particular $C_fe^\sigma\le b(dd^*\beta)/2$.
The construction constants were chosen before $h$ and
the probe; the physical charge interval precedes the
probe-dependent restriction.
\hfill\(\square\)


\section{Transfer from the reference volumes}\label{sec:uniform-transfer}

The renormalization estimates are obtained on two reference tori.
To use them for volume-first limits, we need bounds on products
of local observables in arbitrarily large cutoff tori. The
chessboard estimate reduces each such product to a folded
one-cell observable. A separate convexity estimate controls
the normalization by the partition functions of the two
reference volumes.

\subsection{Folded observables and the partition defect}

Use unit-lattice lengths in this subsection. Fix the charge $q>0$
and a real lattice mass. Choose coordinate cells of width
$w\in\tfrac12\Z_{>0}$ and offsets $o_\mu\in\tfrac12\Z$:
\[
 Q_c^\circ=\{y:o_\mu+wc_\mu<y_\mu<o_\mu+w(c_\mu+1),\
                           \mu=1,2,3\}.
\]
The walls are admissible coordinate mirrors. For even positive
integers $s_\mu$ with $ws_\mu\ge4$, let $C(ws)$ be the rectangular
torus with these side lengths and antiperiodic fermions.
It has $s_1s_2s_3$ cells. The wall-reflection group acts simply
transitively on the covering cells. Write $O^{[c]}$ for the
reflected copy of $O\in\mathcal A(Q_0^\circ)$ in $Q_c^\circ$,
and
\[
 O^\sharp_{ws}=\prod_{c\bmod s}O^{[c]},\qquad
 \mathcal Z_s(O)=Z_{C(ws)}[O^\sharp_{ws}].
\]
The copies use the antilinear reflection action, including
conjugation of scalar coefficients.

These arrays satisfy the rectangular chessboard estimate
\eqref{eq:slab-chessboard-rectangular}. The half-integer width
causes no change in its proof: all walls remain admissible,
and the torus sides are integers at least four.
For a support box $B=y+[-r,r]^3\subset Q_c^\circ$ one has
$2r<w$. A nonzero period has sup norm at least $2w>4r$
and at least four. Thus its enlarged boxes $B+[-1,1]^3$
are disjoint, by the same embedding estimate used in the
proof of Theorem~\ref{thm:slab-chessboard}.
The products therefore belong to the required slab classes.

Hold two entries of $s$ fixed and vary $s_\mu$.
The product of the copies in the layer $c_\mu=0$ gives
a Hilbert--Schmidt slab operator $Y$ independent of $s_\mu$.
Its reflected layer pairs fill the entire folded array.
Equation \eqref{eq:slab-fold-holder} gives
\begin{equation}\label{eq:transfer-fold-trace}
 \mathcal Z_s(O)=\operatorname{Tr}(Y^*Y)^{s_\mu/2}.
\end{equation}
In particular $\mathcal Z_s(O)\ge0$, and
$s_\mu\mapsto\mathcal Z_s(O)^{1/s_\mu}$ is nonincreasing.
If $Y\ne0$, the logarithm of this folded partition function
is convex in $s_\mu$. The first assertion is the monotonicity
of Schatten norms, and the second follows from H\"older's
inequality for the trace of positive powers.
The normalization of $Y$ here is that of the slab operators
in Section~\ref{subsec:slab-folds}.

Choose an even reference cell count $s_0\ge2$ with
$n_r=ws_0\ge4$, and an integer $j\ge2$. Put
\begin{equation}\label{eq:transfer-reference-data}
 N_r=s_0^3,\qquad
 D_j=j^3\log Z_{C(n_r)}[1]-\log Z_{C(jn_r)}[1],
 \qquad \delta_j=\frac{D_j}{(j-1)^3N_r}.
\end{equation}
For an observable in the base cell its normalized reference
seminorm is
\begin{equation}\label{eq:transfer-reference-seminorm}
 \|O\|_*=
 \left(\frac{\mathcal Z_{(s_0,s_0,s_0)}(O)}
                    {Z_{C(n_r)}[1]}\right)^{1/N_r}.
\end{equation}
The partition function is positive, and the numerator is
nonnegative, so this is well-defined also when it vanishes.

\begin{proposition}\label{prop:transfer-general}
One has $D_j\ge0$. Let $s$ be even, with either $s=s_0$
or $s\ge js_0$, and put $r=s/s_0$, $n=ws$.
For every set $\mathcal C$ of distinct cells of $C(n)$ and
$O_c\in\mathcal A(Q_0^\circ)$,
\begin{equation}\label{eq:transfer-general}
 \left|\E_{C(n)}\prod_{c\in\mathcal C}O_c^{[c]}\right|
                    \le\prod_{c\in\mathcal C}e^{\delta_j}\|O_c\|_*.
\end{equation}
The right side uses only the two reference partition functions
and the one-cell observables folded on $C(n_r)$.
\end{proposition}

\begin{proof}
Successively decreasing the three side counts in
\eqref{eq:transfer-fold-trace} gives, for $s\ge s_0$,
\begin{equation}\label{eq:transfer-fold-upper}
 \mathcal Z_{(s,s,s)}(O)
            \le\mathcal Z_{(s_0,s_0,s_0)}(O)^{r^3}.
\end{equation}
Taking $O=1$ and $r=j$ proves $D_j\ge0$.

We next obtain the lower bound for the normalizing partition
function. Define
$p(r_1,r_2,r_3)=\log Z_{C(r_1n_r,r_2n_r,r_3n_r)}[1]$
on ratios with integer side lengths at least four.
The pressure trace \eqref{eq:slab-pressure-trace} shows that
$p$ is convex in each argument and that $p/r_\mu$ is
nonincreasing in that argument. For $r\ge j$ put
$a=(r-1)/(j-1)$. Convexity gives
$p(r,\cdot,\cdot)\ge a p(j,\cdot,\cdot)
-(a-1)p(1,\cdot,\cdot)$.
Apply this in each argument, retaining the positive coefficient
$a$ at the next step. Bound the three subtracted terms by the
side monotonicity. The result is
\begin{align*}
 p(r,r,r)\ge{}&
 a^3p(j,j,j)\\
 &-(a-1)(a^2j^2+ajr+r^2)p(1,1,1)\\
 ={}&r^3p(1,1,1)-a^3D_j.
\end{align*}
The last equality uses $aj-(a-1)=r$ three times.
The negative coefficients are why upper bounds are used on
the three subtracted terms; no sign condition on $p$ is needed.
Thus, including the equality case $r=1$,
\begin{equation}\label{eq:transfer-partition-lower}
 \log Z_{C(n)}[1]\ge
 r^3\left(\log Z_{C(n_r)}[1]-\frac{D_j}{(j-1)^3}\right).
\end{equation}

Fill the cells outside $\mathcal C$ with the constant integrand.
The normalized rectangular chessboard estimate then gives
\[
 \left|\E_{C(n)}\prod_{c\in\mathcal C}O_c^{[c]}\right|
 \le\prod_{c\in\mathcal C}
       \left(\frac{\mathcal Z_{(s,s,s)}(O_c)}
                         {Z_{C(n)}[1]}\right)^{1/(r^3N_r)}.
\]
Use \eqref{eq:transfer-fold-upper} in the numerator and
\eqref{eq:transfer-partition-lower} in the denominator.
Each factor is at most $e^{\delta_j}\|O_c\|_*$, proving
\eqref{eq:transfer-general}.
\end{proof}

\subsection{The photon part of the defect}

The zero-flux Gaussian factor involves the determinant of the
unit-lattice scalar Laplacian with its constant mode removed.
Write its logarithm as
\[
 \mathcal L(n_1,n_2,n_3)
  =\sum_{\substack{k\in\prod_\mu\Z/n_\mu\Z\\k\ne0}}
      \log\left(\sum_{\mu=1}^3
                       \left(2-2\cos\frac{2\pi k_\mu}{n_\mu}\right)\right),
 \qquad \mathcal L(n)=\mathcal L(n,n,n).
\]
For $j\ge2$ put
\[
 G_j=\left(1+\frac j{3.2}\right)\frac{2j^2}{2.56},\qquad
 K_j=(1+j+j^2)(\log j+jG_j).
\]
The estimate used for the two reference volumes is
\begin{equation}\label{eq:transfer-photon-logdet}
 \frac12\bigl(\mathcal L(jn)-j^3\mathcal L(n)\bigr)
              \le-(j^3-1)\log n+K_j
                   \qquad(n,j\in\Z,\ n,j\ge2).
\end{equation}
Its logarithmic term cancels the corresponding scale dependence
of the Gaussian normalization. We retain the finite-volume
calculation since a bulk determinant bound would lose this term.

Sum first in the first coordinate. For a nonzero transverse
momentum $k'=(k_2,k_3)$ define $\omega(k')>0$ by
\[
 2\cosh\omega(k')
     =2+\sum_{\mu=2,3}
                 \left(2-2\cos\frac{2\pi k_\mu}{n_\mu}\right).
\]
Factoring over the $n_1$-th roots of unity gives
\[
 \prod_{k_1\in\Z/n_1\Z}
       \left(2\cosh\omega-2\cos\frac{2\pi k_1}{n_1}\right)
                       =4\sinh^2(n_1\omega/2).
\]
For $k'=0$, omit $k_1=0$; the corresponding product is $n_1^2$.
Therefore
\begin{equation}\label{eq:transfer-photon-one-direction}
 \mathcal L(n_1,n_2,n_3)
   =2\log n_1+
       2\sum_{k'\ne0}\log\bigl(2\sinh(n_1\omega(k')/2)\bigr).
\end{equation}
For $z>0$ let $g(z)=-\log(1-e^{-2z})$.
The identity $2\sinh z=e^z(1-e^{-2z})$ shows that
\[
 \log(2\sinh(jz))-j\log(2\sinh z)\le jg(z).
\]
Consequently
\begin{multline}\label{eq:transfer-photon-side-difference}
 \mathcal L(jn_1,n_2,n_3)-j\mathcal L(n_1,n_2,n_3)\\
 \le2\log(jn_1)-2j\log n_1
                       +2j\sum_{k'\ne0}g(n_1\omega(k')/2).
\end{multline}

In the three applications below, $n_1=n$ and
$n_2,n_3\in\{n,jn\}$. Represent each transverse momentum by
$|k_\mu|\le n_\mu/2$ and put $p_\mu=2\pi k_\mu/n_\mu$.
For
$M(p)=\sum_{\mu=2,3}(2-2\cos p_\mu)$,
\[
 \frac4{\pi^2}|p|^2\le M(p)\le8,\qquad
 \omega=2\operatorname{arcsinh}(\sqrt M/2)\ge0.51|p|.
\]
The last inequality follows from concavity of
$\operatorname{arcsinh}$ on $[0,\sqrt2]$ and
$\operatorname{arcsinh}\sqrt2>0.81\sqrt2$.
Thus $n\omega(k')/2\ge c_j|k'|$, where $c_j=1.6/j$.
Since
$g(z)\le e^{-2z}(1+(2z)^{-1})$, and
$|k'|\ge(|k_2|+|k_3|)/\sqrt2$, one obtains
\begin{align*}
 \sum_{k'\ne0}g(n\omega(k')/2)
 &\le\left(1+\frac1{2c_j}\right)
   \left[
     \left(\sum_{i\in\Z}e^{-\sqrt2c_j|i|}\right)^2-1
   \right]\\
 &=\left(1+\frac1{2c_j}\right)
                  \sinh(c_j/\sqrt2)^{-2}
 \le\left(1+\frac1{2c_j}\right)\frac2{c_j^2}
 =G_j.
\end{align*}
Enlargement to all of $\Z^2\setminus\{0\}$ is valid because
the chosen representatives are distinct.

Finally telescope the enlargement of the three sides:
\begin{align*}
 \mathcal L(jn,jn,jn)-j^3\mathcal L(n,n,n)
 ={}&\bigl[\mathcal L(jn,jn,jn)-j\mathcal L(n,jn,jn)\bigr]\\
 &+j\bigl[\mathcal L(jn,n,jn)-j\mathcal L(n,n,jn)\bigr]\\
 &+j^2\bigl[\mathcal L(jn,n,n)-j\mathcal L(n,n,n)\bigr].
\end{align*}
Symmetry of $\mathcal L$ supplies this identity.
Each bracket is bounded by
$2\log j+2(1-j)\log n+2jG_j$ by
\eqref{eq:transfer-photon-side-difference}.
Multiplication by $1,j,j^2$ and summation prove
\eqref{eq:transfer-photon-logdet}.

\subsection{The photon, matter and flux contributions}

Fix $q>0$ and a real lattice mass. On a cubic torus with
$n\ge4$ sites per side, let $Z_n^{(w)}[1]$ be the
partition function in flux sector $w\in\Z^3$, and put
$P_n(w)=Z_n^{(w)}[1]/Z_n[1]$. The flux order is
$(23,13,12)$. Define
\begin{equation}\label{eq:transfer-flux-tail}
 \varsigma_n=\sum_{w\ne0}
              \exp\left(-\frac{2\pi^2|w|_2^2}{q^2n}\right).
\end{equation}
The notation $z_{r,n}$ below is the positive matter factor
of \eqref{eq:external-matter-factor}.

\begin{proposition}\label{prop:transfer-flux-defect}
For $n\ge4$ and an integer $j\ge2$,
\begin{align}
 Z_n[1]&=\frac{Z_n^{\mathrm{ph}}z_{r,n}}{P_n(0)},
                  \qquad 0<P_n(0)\le1,\label{eq:transfer-zero-factor}\\
 P_n(w)&\le\exp\left(-\frac{2\pi^2|w|_2^2}{q^2n}\right),
 \qquad P_n(0)\ge1-\varsigma_n. \label{eq:transfer-flux-weight}
\end{align}
If $q^2n\le1$, then
\begin{equation}\label{eq:transfer-flux-small}
 \varsigma_n\le7\exp\left(-\frac{2\pi^2}{q^2n}\right)<\frac18.
\end{equation}
The defect $D_j$ of \eqref{eq:transfer-reference-data}
decomposes as
\begin{align}
 D_j&=D_j^{\mathrm{ph}}+D_j^{\mathrm{mat}}+D_j^{\mathrm{fl}},
                                      \label{eq:transfer-defect-split}\\
 D_j^{\mathrm{mat}}&=j^3\log z_{r,n}-\log z_{r,jn},
 \qquad
 D_j^{\mathrm{fl}}=-j^3\log P_n(0)+\log P_{jn}(0),\notag\\
 D_j^{\mathrm{ph}}
   &=(j^3-1)\log\frac{2\pi}{q^2}
                +\tfrac12[\mathcal L(jn)-j^3\mathcal L(n)]\notag\\
   &\le(j^3-1)\log\frac{2\pi}{q^2n}+K_j.
                                      \label{eq:transfer-photon-defect}
\end{align}
If $\varsigma_n<1$, then
\begin{equation}\label{eq:transfer-flux-defect}
 D_j^{\mathrm{fl}}\le-j^3\log(1-\varsigma_n).
\end{equation}
All identities use the same charge and lattice mass at
the two sides; the mass can have either sign.
\end{proposition}

\begin{proof}
The zero-flux factorization \eqref{eq:external-zero-factorization}
gives $Z_n^{(0)}[1]=Z_n^{\mathrm{ph}}z_{r,n}>0$.
All sectors have nonnegative weights, and their absolutely
convergent sum is $Z_n[1]$. Hence
$0<P_n(0)\le1$ and \eqref{eq:transfer-zero-factor}.

For a configuration of flux $w$, let $h_w$ be the
constant two-cochain with value
$r_qw_{\mu\nu}/n^2$ on plaquettes of orientation $\mu\nu$.
The Fourier Hodge decomposition writes
$r_qz=dB+h$ with $h$ constant.
Summing over a coordinate two-torus annihilates $dB$
and gives $n^2h_{\mu\nu}=r_qw_{\mu\nu}$.
Thus $F^{\mathrm u}-h_w$ is exact. A constant cochain is orthogonal
to every exact cochain, by periodic summation by parts.
In particular $\langle h,F^{\mathrm u}\rangle=\langle h,h_w\rangle$
on sector $w$ whenever $h$ is constant.

Apply Gaussian domination \eqref{eq:slab-gaussian-domination}
with $h=h_{w'}$ and split the expectation into its sectors:
\[
 1\ge\sum_w
 \exp\bigl(\langle h_{w'},h_w\rangle-\|h_{w'}\|^2/2\bigr)P_n(w)
 \ge \exp(\|h_{w'}\|^2/2)P_n(w').
\]
Since $\|h_w\|^2=r_q^2|w|_2^2/n$, this proves the
first bound in \eqref{eq:transfer-flux-weight}.
The identity $\sum_wP_n(w)=1$ proves the second.

Set $x=\exp(-2\pi^2/(q^2n))$. Positivity permits the
factorization
\[
 \varsigma_n=\left(\sum_{i\in\Z}x^{i^2}\right)^3-1
 \le\left(1+\frac{2x}{1-x}\right)^3-1.
\]
For $q^2n\le1$, one has $x<2^{-18}$.
The bound $(1-x)^{-1}\le1+2x$ gives
$2x/(1-x)\le2.001x<10^{-5}$. Therefore
$\varsigma_n\le2.001x(3+3\cdot10^{-5}+10^{-10})
\le7x<1/8$. This also proves convergence of
\eqref{eq:transfer-flux-tail}.

Taking logarithms in \eqref{eq:transfer-zero-factor}
at $n$ and $jn$ gives \eqref{eq:transfer-defect-split}.
By \eqref{eq:external-photon-factor},
\[
 \log Z_n^{\mathrm{ph}}
   =(n^3+2)\log r_q+(n^3-1)\log(2\pi)-\mathcal L(n)/2.
\]
Its bulk terms cancel in the defect. The remaining
constant is $(j^3-1)\log(2\pi/q^2)$.
Equation \eqref{eq:transfer-photon-logdet} proves
\eqref{eq:transfer-photon-defect}. Finally
$\log P_{jn}(0)\le0$, and
$P_n(0)\ge1-\varsigma_n>0$ gives
\eqref{eq:transfer-flux-defect}.
\end{proof}

\subsection{Wick products on arbitrary cutoff tori}

Use the cell sources of Section~\ref{sec:reference-cells}
and write
\[
 N_*=8L^3,\qquad
 a_L=\frac{L^3-1}{(L-1)^3},\qquad
 D^{\mathrm{mat}}=L^3\log z_{r,n_*}-\log z_{r,Ln_*},
 \qquad
 \delta=\frac{D_L}{(L-1)^3N_*}.
\]
Fix $0<e\le1/3$, with $q=e\sqrt\eta$ and
$q^2n_*=e^2R^*$, and assume
\begin{equation}\label{eq:transfer-W3}
 6N_*e^2R^*\log(1/e)\le\pi^2.
\end{equation}
The choice of scales ensuring this condition is given
below. Let $\mathcal C_*=\{0,\ldots,2L-1\}^3$
and use the probe variance $v_\eta$ of
\eqref{eq:cells-orthogonality}. For real $c_2,c_4$ put
\begin{align}
 H_2(x;v)&=x^2-v,\qquad
 H_4(x;v)=x^4-6vx^2+3v^2,\notag\\
 O_m^{[c]}&=e^{-m}H_m(X_c;v_\eta)-c_m,\qquad m=2,4,
                                      \label{eq:transfer-Wick-observable}\\
 \Theta_m^{(0)}
  &=\left|
      \left[\prod_{c\in\mathcal C_*}
            (\partial_{t_c}^m-c_m)\right]\mathring\Phi(0)
    \right|^{1/N_*}.                  \label{eq:transfer-Theta}\\
 M_4&=600N_*^2v_\eta^2e^{-4}+|c_4|,\qquad
 M_2=12N_*v_\eta e^{-2}+|c_2|,\notag\\
 T_m&=\Theta_m^{(0)}+2e^6M_m,\qquad
 u_m=\exp(\delta)T_m.                 \label{eq:transfer-Wick-data}
\end{align}
Here $\mathring\Phi$ is the normalized external partition
function \eqref{eq:external-reference-partition}.
The notation $e^{-m}$ denotes an inverse power of the
charge.

\begin{proposition}\label{prop:transfer-Wick}
Assume the probe placement condition \eqref{eq:cells-fit}.
For $n=n_*$, or $n\in p\Z$ with $n\ge Ln_*$,
the all-flux law with the fixed real lattice mass obeys
\begin{align}
 0\le\delta
 &\le\frac{a_L}{N_*}\log\frac{2\pi}{e^2R^*}
       +\frac{K_L+D^{\mathrm{mat}}+1}{(L-1)^3N_*},
                                           \label{eq:transfer-delta}\\
 \left|\E_n^{\mathrm{af}}\prod_{c\in\mathcal C}
                        O_m^{[c]}\right|
 &\le u_m^{|\mathcal C|}              \label{eq:transfer-Wick-product}
\end{align}
for every set $\mathcal C$ of distinct cells and $m=2,4$.
In particular, if $\kappa_j$ denotes the cumulant of
$X_0$ on that torus, then
\begin{align}
 \left|e^{-2}(\kappa_2-v_\eta)-c_2\right|&\le u_2,
                                      \label{eq:transfer-second-cumulant}\\
 \left|e^{-4}\kappa_4-(c_4-3c_2^2)\right|
   &\le u_4+3u_2(2|c_2|+u_2).        \label{eq:transfer-fourth-cumulant}
\end{align}
The bounds are uniform in $n$.
Their right sides are formed on the two reference tori;
$\Theta_m^{(0)}$, and therefore $u_m$, may depend on
the placement.
\end{proposition}

\begin{proof}
Condition \eqref{eq:transfer-W3}, $N_*\ge1000$ and
$\log(1/e)\ge\log3>1$ imply $e^2R^*<1$.
For $x=\exp(-2\pi^2/(e^2R^*))$ they also give
\[
 x\le e^{12N_*},\qquad x^{1/(2N_*)}\le e^6.
\]
Apply Proposition~\ref{prop:transfer-flux-defect} at
$n_*$ and $j=L$. Since $\varsigma_{n_*}\le7x<1/8$,
\[
 D_L^{\mathrm{fl}}\le-L^3\log(1-\varsigma_{n_*})
              \le8L^3x=N_*x
              \le N_*3^{-12N_*}\le1.
\]
The photon estimate and the defect decomposition
now prove the upper bound in \eqref{eq:transfer-delta};
its lower bound is Proposition~\ref{prop:transfer-general}.

All the polynomials in \eqref{eq:transfer-Wick-observable}
are the same local integrands on every carrying torus,
by Section~\ref{sec:reference-cells}. They are real
elements of the cell observable class. Write
$O_m^\sharp=\prod_{c\in\mathcal C_*}O_m^{[c]}$ on
$C(n_*)$. Its all-flux expectation is nonnegative,
by the fold trace \eqref{eq:transfer-fold-trace}.
We estimate it by splitting flux zero from the
other sectors.

First, Gaussian domination and the moment bound
give, for every positive integer $i$,
\[
 \E_{n_*}^{\mathrm{af}}X_c^{2i}
                              \le(2i\exp(1)v_\eta)^i.
\]
Minkowski's inequality at exponent $2N_*$ yields
\begin{align*}
 \|H_4(X_c;v_\eta)\|_{2N_*}
   &\le[64\exp(2)N_*^2+24\exp(1)N_*+3]v_\eta^2
       \le600N_*^2v_\eta^2,\\
 \|H_2(X_c;v_\eta)\|_{2N_*}
   &\le[4\exp(1)N_*+1]v_\eta
       \le12N_*v_\eta.
\end{align*}
Thus $\|O_m^{[c]}\|_{2N_*}\le M_m$.
The finite moments justify all uses of H\"older
and Cauchy--Schwarz below, also when $v_\eta=0$.

On the zero-flux sector,
\eqref{eq:external-Wick-product} gives
\[
 \left|\E_{n_*}^{\mathrm{af}}
               [O_m^\sharp\mathbf1_{\{\mathrm{flux}=0\}}]\right|
   =P_{n_*}(0)\left|\E_{n_*}^{\mathrm{av}}O_m^\sharp\right|
   \le(\Theta_m^{(0)})^{N_*}.
\]
On the remaining sectors, Cauchy--Schwarz followed
by H\"older with $N_*$ factors gives
\begin{align*}
 \left|\E_{n_*}^{\mathrm{af}}
               [O_m^\sharp\mathbf1_{\{\mathrm{flux}\ne0\}}]\right|
 &\le\varsigma_{n_*}^{1/2}
                \bigl(\E_{n_*}^{\mathrm{af}}|O_m^\sharp|^2\bigr)^{1/2}\\
 &\le\varsigma_{n_*}^{1/2}
                \prod_{c\in\mathcal C_*}\|O_m^{[c]}\|_{2N_*}\\
 &\le\bigl(\varsigma_{n_*}^{1/(2N_*)}M_m\bigr)^{N_*}.
\end{align*}
Since $\varsigma_{n_*}^{1/(2N_*)}\le2e^6$,
the nonnegative $N_*$-th root satisfies
\[
 \|O_m\|_*=
       \bigl(\E_{n_*}^{\mathrm{af}}O_m^\sharp\bigr)^{1/N_*}
             \le\Theta_m^{(0)}+2e^6M_m=T_m.
\]
For nonnegative $a,b$ we used
$(a^{N_*}+b^{N_*})^{1/N_*}\le a+b$.
Proposition~\ref{prop:transfer-general}, with
$s_0=2L$ and $j=L$, proves
\eqref{eq:transfer-Wick-product}.

For one cell the same estimate gives
$|e^{-m}\E H_m(X_0;v_\eta)-c_m|\le u_m$.
Charge conjugation makes the odd moments zero.
Writing $y=e^{-2}\E H_2(X_0;v_\eta)$, one has
\[
 e^{-4}\kappa_4=e^{-4}\E H_4(X_0;v_\eta)-3y^2.
\]
The second-order conclusion follows at once.
For the fourth, use $|y-c_2|\le u_2$ and
$|y+c_2|\le2|c_2|+u_2$ to obtain
\eqref{eq:transfer-fourth-cumulant}.
\end{proof}

\subsection{Reference scales and the loss of charge exponent}

The reference side is chosen from the logarithmic charge
scale. Fix a power $M$ of $L$ with $M\ge L^2$ and an
integer $r\ge2$. With $h$ as in \eqref{eq:cutoff-scales},
set
\begin{equation}\label{eq:transfer-scale-choice}
 e_K=e\sqrt h,\qquad v_K=-\log e_K,\qquad
 r_K=\min\{L^i:L^i\ge v_K^r,\ i\in\Z_{\ge0}\},
 \qquad R^*=\frac{hMr_K}{L}.
\end{equation}
Assume $v_K\ge8$, $e\le1/3$, $e_K\log h\le1$, and
\begin{equation}\label{eq:transfer-tail-condition}
 6N_*Mv_K^r \exp(-v_K)
                 \left(v_K\exp(-v_K)+\tfrac12\right)\le\pi^2.
\end{equation}
These are finite charge conditions after $h,L,M,r$ are fixed.
They hold on a sufficiently small charge interval:
$v_K\to\infty$ as $e\to0$, and the polynomial times
exponential in \eqref{eq:transfer-tail-condition}
tends to zero.

Minimality gives $v_K^r\le r_K<Lv_K^r$.
Consequently $T^*=\log_L(Mr_K)-1$ is an integer at
least one, $R^*=hL^{T^*}$, and
$\ell_c=R^*/(2L)\to\infty$ as $e\to0$.
Moreover $1/L<\bar mh\le1$, so
\begin{equation}\label{eq:transfer-volume-error}
 \bar mR^*>\frac{Mr_K}{L^2}\ge r_K\ge v_K,
 \qquad \exp(-\bar mR^*)\le e\sqrt h.
\end{equation}
The same upper bound on $r_K$ gives
\begin{align*}
 e^2R^*\log(1/e)
 &<Mv_K^r \exp(-2v_K)
                          (v_K+\tfrac12\log h)\\
 &\le Mv_K^r \exp(-v_K)
                          (v_K\exp(-v_K)+\tfrac12).
\end{align*}
Thus \eqref{eq:transfer-tail-condition} implies
\eqref{eq:transfer-W3}.
The probe fits every sufficiently small charge and
sufficiently fine mesh, since $D_\beta^{\mathrm{buf}}$
is fixed and $\ell_c\to\infty$.

We now isolate the two reference estimates used in the
passage to arbitrary volumes. Their dependence on the
charge matters: a bounded partition defect still loses
a small power of $e$ through its photon contribution.
Let $C_{\mathrm{mat}}$ and $C_{\mathrm{ref}}$ be finite constants
independent of the charge, mesh and placement. Assume
that for $N\ge N_{\mathrm{ref}}(e)$ the mass assignment and reference
response satisfy
\begin{align}
 D^{\mathrm{mat}}&\le C_{\mathrm{mat}},\notag\\
 \Theta_4^{(0)}&\le C_{\mathrm{ref}} e^{\mathfrak e},\qquad
 (1+|c_2|)\Theta_2^{(0)}
                          \le C_{\mathrm{ref}} e^{\mathfrak e},
                                   \label{eq:transfer-terminal-input}
\end{align}
where
\begin{equation}\label{eq:transfer-charge-exponents}
 \mathfrak e=s(1/4-9\varepsilon)-12\varepsilon,\qquad
 \sigma=\mathfrak e-\frac{2a_L}{N_*}>0,\qquad
 0<s<1,\quad 0<\varepsilon\le1/200.
\end{equation}
Here $N_{\mathrm{ref}}(e)$ is a fixed finite mesh threshold.
For the chamber construction, \eqref{eq:term-matter-comparison}
and \eqref{eq:interacting-final-terminal-error} give
\eqref{eq:transfer-terminal-input} for
$0<e\le\min(e_G,e_\sharp(\beta))$.
The constants and cutoff thresholds are those fixed in
Section~\ref{sec:ordered-constants}.

On the reference torus take the free coefficients
$c_m=N_0^{(m)}(0)$ from
\eqref{eq:response-coefficient-relations}, using the
source associated with the same placed probe.
Proposition~\ref{prop:response-uniform-coefficients}
and the support fit give a constant, enlarged into
$C_{\mathrm{ref}}$ if necessary, such that
\begin{equation}\label{eq:transfer-free-input}
 |c_2|+|c_4|^{1/2}\le C_{\mathrm{ref}},\qquad
 |c_4-3c_2^2-b(f)|
              \le C_{\mathrm{ref}}(e+\eta^{1/4}).
\end{equation}
Indeed $R^*\ge D_\beta+1$, and
\eqref{eq:transfer-volume-error} bounds the finite-volume
error in Proposition~\ref{prop:response-fourth}.
The coefficients are real and independent of placement.

\begin{proposition}\label{prop:transfer-UV-inputs}
Fix the probe and the scale data above, and assume
\eqref{eq:cells-fit} and \eqref{eq:transfer-terminal-input}.
There is a finite
constant $C_U$, independent of $e,N,T$ and placement,
such that for all the stated reference charges and
$N\ge N_{\mathrm{ref}}(e)$,
\begin{equation}\label{eq:transfer-small-Wick}
 u_4\le C_Ue^\sigma,\qquad
 (1+|c_2|)u_2\le C_Ue^\sigma,\qquad
 |c_4-3c_2^2-b(f)|\le C_U(e^\sigma+\eta^{1/4}).
\end{equation}
On every cutoff torus with $T\ge T^*+1$, every finite
set $Y\subset(R^*/L)\Z^3$ of points incongruent
modulo $R\Z^3$ satisfies
\begin{equation}\label{eq:transfer-translated-Wick}
 \left|\E_{N,T}\prod_{y\in Y}
       W_m(F^\sharp(\tau_yf^\eta);v_\eta,c_m)\right|
                          \le u_m^{|Y|},\qquad m=2,4,
\end{equation}
where $W_m(x;v,c)=e^{-m}H_m(x;v)-c$.
The fourth cumulant of the fixed smooth test obeys
\begin{equation}\label{eq:transfer-smooth-cumulant}
 |\kappa_4(F^\sharp(f))-e^4b(f)|
 \le(8C_U+3C_U^2)e^{4+\sigma}
                  +C_U(\eta^2+e^4\eta^{1/4}).
\end{equation}
The tests $f^\eta$ have one fixed compact support,
$f^\eta\to f$ in $L^2$, and $v_\eta\to\|f\|_2^2$.
Thus \eqref{eq:transfer-small-Wick} and
\eqref{eq:transfer-translated-Wick} give exactly the
cutoff hypotheses of Proposition~\ref{prop:UV-component-cumulant}.
\end{proposition}

\begin{proof}
Write
\[
 s_\beta=\|f\|_2+\|\beta\|_{H^4},\qquad
 C_\Gamma^{\mathrm{tr}}=(2\pi)^{a_L/N_*}
         \exp\left(\frac{K_L+C_{\mathrm{mat}}+1}
                                  {(L-1)^3N_*}\right).
\]
Since $R^*\ge1$, \eqref{eq:transfer-delta} gives
\begin{equation}\label{eq:transfer-Gamma-loss}
 \exp(\delta)\le C_\Gamma^{\mathrm{tr}} e^{-2a_L/N_*}.
\end{equation}
The inverse power here is a power of the charge.
The source-curvature identity and Jensen's inequality
for the disjoint averaging cubes give
$v_\eta^{1/2}\le\|f^\eta\|_2\le s_\beta$.
Use this and \eqref{eq:transfer-free-input} in
\eqref{eq:transfer-Wick-data}:
\begin{align*}
 T_4
 &\le C_{\mathrm{ref}} e^{\mathfrak e}
           +1200N_*^2s_\beta^4e^2+2C_{\mathrm{ref}}^2e^6\\
 &\le(C_{\mathrm{ref}}+1200N_*^2s_\beta^4+2C_{\mathrm{ref}}^2)e^{\mathfrak e},\\
 (1+|c_2|)T_2
 &\le C_{\mathrm{ref}} e^{\mathfrak e}
      +2e^6(12N_*s_\beta^2e^{-2}+C_{\mathrm{ref}})(1+C_{\mathrm{ref}})\\
 &\le[C_{\mathrm{ref}}+2(12N_*s_\beta^2+C_{\mathrm{ref}})(1+C_{\mathrm{ref}})]
                                                    e^{\mathfrak e}.
\end{align*}
Here $0<\mathfrak e<1/4$, from
\eqref{eq:transfer-charge-exponents}.
One possible constant is
\begin{equation}\label{eq:transfer-CU}
 \begin{split}
 C_U=\max\{&
 C_\Gamma^{\mathrm{tr}}(C_{\mathrm{ref}}+1200N_*^2s_\beta^4+2C_{\mathrm{ref}}^2),\\
 &C_\Gamma^{\mathrm{tr}}[C_{\mathrm{ref}}+2(12N_*s_\beta^2+C_{\mathrm{ref}})(1+C_{\mathrm{ref}})],
              \ C_{\mathrm{ref}},\ 60s_\beta^4\}.
 \end{split}
\end{equation}
Multiplication by \eqref{eq:transfer-Gamma-loss}
proves the first two bounds in
\eqref{eq:transfer-small-Wick}. The last follows from
\eqref{eq:transfer-free-input}, since $e\le e^\sigma$.

For a cutoff side $n=8hL^{T+N}$ with $T\ge T^*+1$,
one has $n=8L^{T-T^*}n_*$, a multiple of $p$ with
$n\ge Ln_*$. For $y\in(R^*/L)\Z^3$ use the even
cell index $c(y)$ of
\eqref{eq:cells-translation-observable}.
Incongruent points give distinct cells.
Proposition~\ref{prop:transfer-Wick} applies to their
product, and lattice translation invariance removes
the common placement. This proves
\eqref{eq:transfer-translated-Wick}.
For a fixed finite set of distinct points it applies
as soon as $R>2\max_{y\in Y}|y|_\infty$.

It remains to replace the discrete smooth test
$f^\eta$ by $f$ in the cumulant. We give the
quantitative multilinear estimate used for this purpose.
For a real two-cochain $g$, charge conjugation and
Gaussian domination give
\[
 \E\cosh(t\langle g,F^{\mathrm u}\rangle)
                          \le\exp(t^2\|g\|^2/2).
\]
Extract its nonnegative term of degree $2i$ and choose
$t^2=2i/\|g\|^2$. This gives
\[
 \E\langle g,F^{\mathrm u}\rangle^{2i}
    \le(2i)!\left(\frac{\|g\|^2}{2i}\right)^i \exp(i).
\]
For $i=2$ the bound is $(3/2)\exp(2)\|g\|^4$.
H\"older bounds the four-field moment by
\[
 \tfrac32\exp(2)\prod_a\|g_a\|.
\]
The variance bound bounds each of its three pair subtractions
by $\prod_a\|g_a\|$. Since $(3/2)\exp(2)+3<15$,
the symmetric multilinear forms $K_2,K_4$ obey
\begin{equation}\label{eq:transfer-multilinear-cumulants}
 |K_2(g_1,g_2)|\le\|g_1\|\|g_2\|,\qquad
 |K_4(g_1,g_2,g_3,g_4)|\le15\prod_a\|g_a\|.
\end{equation}
For $g=0$ the corresponding moment is zero.

Let $g=\eta^{3/2}f_\eta$ and
$g'=\eta^{3/2}(f^\eta)_\eta$ on the cutoff torus.
The common support cube has side less than $R$.
The averaging bound and \eqref{eq:cells-test-error} give
\[
 \|g\|,\|g'\|\le s_\beta,\qquad
 \|g-g'\|\le\eta^2\|\beta\|_{H^4}.
\]
The observables paired with these cochains are
$F^\sharp(f)$ and $F^\sharp(f^\eta)$.
Telescope the four slots of $K_4$ and the two slots
of $K_2$ in \eqref{eq:transfer-multilinear-cumulants}.
This proves
\begin{align}
 |\kappa_4(F^\sharp(f))-\kappa_4(F^\sharp(f^\eta))|
                              &\le60s_\beta^4\eta^2,\notag\\
 |\kappa_2(F^\sharp(f))-\kappa_2(F^\sharp(f^\eta))|
                              &\le2s_\beta^2\eta^2.
                                   \label{eq:transfer-test-replacement}
\end{align}
One-cell transfer and \eqref{eq:transfer-small-Wick}
bound the fourth cumulant of $F^\sharp(f^\eta)$ by
\[
 |\kappa_4(F^\sharp(f^\eta))-e^4(c_4-3c_2^2)|
                     \le(7C_U+3C_U^2)e^{4+\sigma}.
\]
Adding the free-coefficient error and the first bound
in \eqref{eq:transfer-test-replacement} proves
\eqref{eq:transfer-smooth-cumulant}.
Finally the compact support, $L^2$ convergence and
variance convergence are respectively
\eqref{eq:cells-smooth-test}, \eqref{eq:cells-test-error}
and \eqref{eq:cells-variance-limit}.
\end{proof}


\section{Fourth cumulants of the limiting components}
\label{sec:ultraviolet-limits}

The fourth-cumulant estimate for a mixed limiting law does not by
itself control each ergodic component. We use the stronger estimate
on products of translated Wick polynomials. It bounds their
translation-invariant conditional expectations, and hence the
second and fourth moments of almost every component.

\subsection{From distinct-point products to the invariant part}

Let $U(y)$,
$y\in\Z^d$, be the operators of a measure-preserving action on
a probability space, where $d\ge1$, with $U(y)G=G\circ T_{-y}$. Let
$P_{\mathrm{lat}}$ be the orthogonal projection in $L^2$ onto
the vectors invariant under this action.

\begin{lemma}\label{lem:UV-invariant-product-bound}
Suppose that $G$ is real and belongs to every $L^p$, $p<\infty$.
If $u\ge0$ and
\begin{equation}\label{eq:UV-distinct-product-bound}
 \left|\E\prod_{a=1}^J U(y_a)G\right|\le u^J
\end{equation}
for every finite family of pairwise distinct points of $\Z^d$,
then $|P_{\mathrm{lat}}G|\le u$ almost surely.
\end{lemma}

\begin{proof}
For $B_n=\{0,\ldots,n-1\}^d$ put
$A_n=n^{-d}\sum_{y\in B_n}U(y)$. These contractions converge
strongly in $L^2$ to $P_{\mathrm{lat}}$.
Indeed, the orthogonal complement of the invariant space is the
closed span of $(U(y)-1)H$. On such a vector, cancellation of
the common terms in $B_n$ and $y+B_n$ gives
\[
 \|A_n(U(y)-1)H\|_2\le2|y|_1n^{-1}\|H\|_2.
\]
On invariant vectors the averages are the identity, proving
the asserted convergence by the orthogonal decomposition.

For an even integer $J\ge2$, expand $\E(A_nG)^J$.
The distinct-index terms obey \eqref{eq:UV-distinct-product-bound}.
There are at most $\binom J2 n^{d(J-1)}$ remaining tuples:
each has a repeated pair of indices. H\"older's inequality
bounds their individual expectations in modulus by
$\|G\|_J^J$. Therefore
\[
 \E|A_nG|^J
       \le u^J+\binom J2 n^{-d}\|G\|_J^J.
\]
Along a subsequence the $L^2$ convergence is almost sure.
Fatou's lemma gives $\E|P_{\mathrm{lat}}G|^J\le u^J$
for every even $J$. For $t>u$, Markov's inequality now gives
$\Pr(|P_{\mathrm{lat}}G|>t)\le(u/t)^J$.
Let $J\to\infty$ and then $t\downarrow u$ through a countable
sequence.
\end{proof}

Let a translation-invariant law $\mu$ on tempered fields have
the decomposition $\mu=\int\mu_\xi\,\dd\nu(\xi)$.
Write $I$ for its translation-invariant sigma algebra and
$P=\E_\mu[\,\cdot\mid I]$. Apply the lemma to translations
on $\lambda\Z^3$, with $\lambda>0$ fixed. The range of $P$
is contained in that of $P_{\mathrm{lat}}$, so
$P=PP_{\mathrm{lat}}$. Positivity of conditional expectation
therefore gives
\begin{equation}\label{eq:UV-component-product-bound}
 |PG|\le P|P_{\mathrm{lat}}G|\le u,\qquad
 |\mu_\xi(G)|\le u\quad\text{for }\nu\text{-almost every }\xi.
\end{equation}
The second assertion uses the conditional-law representation of
the decomposition. Integrability under the conditional laws and
the representation for this integrable $G$ follow by truncating
its positive and negative parts.

\subsection{Passing the Wick-product estimate to a limit law}

Fix the charge $e>0$ and a real compactly supported test $f$.
For $v\ge0$ define
\[
 H_2(x;v)=x^2-v,\qquad
 H_4(x;v)=x^4-6vx^2+3v^2,\qquad
 W_m(x;v,c)=e^{-m}H_m(x;v)-c\quad(m=2,4).
\]
Here $e^{-m}$ is the inverse power of the charge.
In the application $v$ is a variance, while the exponent in
the small-charge error is denoted by $\sigma$.

Take a joint cutoff sequence defining a limit law $\mu$.
Suppose that the lattice estimates give compactly
supported real tests $f_i$, all supported in one fixed cube,
and real numbers
\[
 v_i,\ c_{2,i},\ c_{4,i},\ u_{2,i},\ u_{4,i},
\]
with
$v_i\ge0$, such that
\begin{equation}\label{eq:UV-test-convergence}
 \|f_i-f\|_2\to0,\qquad
 v_i\to\|f\|_2^2,\qquad
 \sup_i(|c_{2,i}|+|c_{4,i}|)<\infty.
\end{equation}
For a fixed $\lambda>0$, assume that
\begin{equation}\label{eq:UV-cutoff-product-bound}
 \left|\E_i\prod_{y\in Y}
       W_m(F^\sharp(\tau_y f_i);v_i,c_{m,i})\right|
                     \le u_{m,i}^{|Y|},\qquad m=2,4,
\end{equation}
for every finite $Y\subset\lambda\Z^3$ and all sufficiently
large $i$ depending on $Y$. This is the translated-cell form
of the transfer estimate; the volume condition makes the
points of $Y$ incongruent modulo the torus.
The reference estimates also give
\begin{equation}\label{eq:UV-reference-bounds}
 \begin{split}
 0\le u_{4,i}&\le Ce^\sigma,\qquad
 0\le(1+|c_{2,i}|)u_{2,i}\le Ce^\sigma,\\
 |c_{4,i}-3c_{2,i}^2-b(f)|
                  &\le C(e^\sigma+\eta_i^{1/4}),
 \end{split}
\end{equation}
with $C$ independent of the cutoff, $\sigma>0$ and $0<e\le1$.
For the chamber laws, Proposition~\ref{prop:transfer-UV-inputs}
supplies these hypotheses throughout the charge interval
\eqref{eq:constants-probe-charge}.

Pass to a subsequence on which the four bounded sequences
$c_{2,i},c_{4,i},u_{2,i},u_{4,i}$ converge, with limits
$c_2,c_4,u_2,u_4$. The weak limit remains $\mu$.
Equation \eqref{eq:UV-reference-bounds} gives
\begin{equation}\label{eq:UV-reference-limit}
 u_4\le Ce^\sigma,\quad
 (1+|c_2|)u_2\le Ce^\sigma,\quad
 |c_4-3c_2^2-b(f)|\le Ce^\sigma.
\end{equation}
Set $G_m=W_m(F(f);\|f\|_2^2,c_m)$.
Then \eqref{eq:UV-cutoff-product-bound} passes to
\begin{equation}\label{eq:UV-law-product-bound}
 \left|\E_\mu\prod_{y\in Y}U(y)G_m\right|\le u_m^{|Y|}
                   \qquad(Y\subset\lambda\Z^3\text{ finite}).
\end{equation}

To justify this passage for the varying tests, expand the finite
product of Wick polynomials into monomials. For a monomial
of positive total degree $J$, telescope the replacement of each
$F^\sharp(\tau_y f_i)$ by $F^\sharp(\tau_y f)$.
All these tests have support in one open $R_i$-cube once $i$
is large. The test-map variance and product-moment bounds
then give an error at most
\[
 Jc_J B^{J-1}\|f_i-f\|_2,
 \qquad B=\sup_i\max(\|f_i\|_2,\|f\|_2).
\]
This tends to zero. The moments of the fixed tests converge
by Section~\ref{sec:continuum-limits}, and the polynomial
coefficients converge by \eqref{eq:UV-test-convergence}.
There are only finitely many monomials for each fixed $Y$.
This proves \eqref{eq:UV-law-product-bound}.

\subsection{The cumulant in an ergodic component}

\begin{proposition}\label{prop:UV-component-cumulant}
Under \eqref{eq:UV-test-convergence}--\eqref{eq:UV-reference-bounds},
suppose that almost every translation-ergodic component of $\mu$
satisfies Gaussian domination. Then
\begin{equation}\label{eq:UV-component-cumulant}
 \left|\kappa_4^{\mu_\xi}(f)-e^4b(f)\right|
                 \le(8C+3C^2)e^{4+\sigma}
\end{equation}
for almost every component. In particular, if $b(f)>0$ and
$(8C+3C^2)e^\sigma\le b(f)/2$, its fourth cumulant is positive.
\end{proposition}

\begin{proof}
Gaussian domination of $\mu$ makes the real polynomials $G_2,G_4$
belong to every $L^p(\mu)$. Apply
Lemma~\ref{lem:UV-invariant-product-bound} to
\eqref{eq:UV-law-product-bound}, and then use
\eqref{eq:UV-component-product-bound}. On one set of full
component measure,
\[
 |d_m|\le u_m,\qquad d_m=\E_{\mu_\xi}G_m,\qquad m=2,4.
\]
Intersect with the full-measure set of components having Gaussian
domination. For a component in this intersection, put $X=F(f)$
and $v=\|f\|_2^2$. Gaussian domination gives $\E_{\mu_\xi}X=0$
and all moments finite. Thus
\[
 \kappa_4^{\mu_\xi}(f)
  =\E_{\mu_\xi}H_4(X;v)
          -3\bigl(\E_{\mu_\xi}H_2(X;v)\bigr)^2.
\]
The definition of $G_m$ gives the exact identity
\begin{equation}\label{eq:UV-cumulant-identity}
 e^{-4}\kappa_4^{\mu_\xi}(f)
                =c_4+d_4-3(c_2+d_2)^2.
\end{equation}
Consequently \eqref{eq:UV-reference-limit} implies
\begin{align*}
 |e^{-4}\kappa_4^{\mu_\xi}(f)-b(f)|
 &\le |c_4-3c_2^2-b(f)|+u_4+6|c_2|u_2+3u_2^2\\
 &\le8Ce^\sigma+3C^2e^{2\sigma}
 \le(8C+3C^2)e^\sigma.
\end{align*}
This proves the assertion uniformly over the full-measure set
of components. A centered normal variable has zero fourth
cumulant, so positivity also proves that these laws are
non-Gaussian.
\end{proof}



\section{Tensor Schwinger functions}\label{sec:tensor-axioms}

The Gaussian bound has already defined a continuous map
$\Psi\mapsto Q_\Psi$ from $\Sspace_n$ to $L^2(\lambda)$.
To use law-level reflection positivity for arbitrary tensor tests,
we must show that support in a half-space makes this random polynomial
measurable with respect to the corresponding half-space sigma algebra. The Hermite
expansion used for the growth bound does not preserve supports.
We therefore use a compact Fourier expansion for this step.

\subsection{Products and supported tensor tests}

For $\Psi\in\Sspace_k$ and $\Phi\in\Sspace_l$, the tensor evaluation
obeys
\begin{equation}\label{eq:tensor-product-rule}
 Q_\Psi Q_\Phi=Q_{\Psi\otimes\Phi}\quad\text{in }L^1(\lambda).
\end{equation}
Both sides are separately continuous bilinear maps to $L^1$:
for the left side use the $L^2$ bounds and Cauchy--Schwarz, and for
the right side use continuity of the tensor product in Schwartz
seminorms. They agree for elementary tensors. Their linear span
is dense, so approximation in each variable proves
\eqref{eq:tensor-product-rule}. The same argument gives
\begin{equation}\label{eq:tensor-equivariance}
 Q_{\overline\Psi}=\overline{Q_\Psi},\qquad
 Q_{g\cdot\Psi}=\widehat g Q_\Psi
 \quad\text{if }g_*\lambda=\lambda.
\end{equation}
Here the second identity uses the exterior-square action, including
the tensor indices.

We recall the needed supported expansion. Let $O\subset\R^3$ be open
and $\Psi\in C_c^\infty(O^n)$. Its support projects into compact
sets $K_a\subset O$. Choose $\chi_a\in C_c^\infty(O)$ equal to one
on $K_a$, and a cube of side $s$ containing all these supports in
its interior. Extend $\Psi$ periodically from this cube in all its
variables. Its Fourier coefficients $c_{\mathbf k}$ obey, for
every integer $p$,
\[
 |c_{\mathbf k}|\le C_{p,s,\Psi}(1+|\mathbf k|^2)^{-p},
 \qquad \mathbf k=(k_1,\ldots,k_n)\in(\Z^3)^n.
\]
This follows by applying $(1-\Delta)^p$ and integrating by parts;
there are no boundary terms. Multiplication of the Fourier series by
$\prod_a\chi_a(y_a)$ gives the identity on all space
\begin{equation}\label{eq:supported-tensor-expansion}
 \Psi=\sum_{\mathbf k}c_{\mathbf k}
       \bigotimes_{a=1}^n
       \left(\chi_a(y_a)e^{2\pi i k_a\cdot y_a/s}\right).
\end{equation}
The seminorm $q_m$ of each tensor factor product grows at most as
$(1+|\mathbf k|^2)^{m/2}$. Taking $p=m+3n$ makes the series
absolutely convergent in that seminorm. The same argument applies
componentwise to two-form tensor tests. It also proves the density
of elementary tensors used above; arbitrary Schwartz tests are
approximated by growing compact cutoffs.

Each elementary tensor in \eqref{eq:supported-tensor-expansion}
has factors supported in $O$, so its random evaluation is
$\Sigma_O$-measurable and belongs to $L^2(\lambda)$ by the moment
bound. Continuity of $Q$ and closedness of $L^2(\Sigma_O,\lambda)$
give
\begin{equation}\label{eq:tensor-local-measurability}
 Q_\Psi\in L^2(\Sigma_O,\lambda)
       \qquad(\Psi\in C_c^\infty(O^n)).
\end{equation}

This extends to tests supported in the closed half-space
$H=\{x:v\cdot x\ge c\}$, with $|v|=1$. Choose a fixed compact
cutoff $\chi$ equal to one near zero and put
\[
 \Psi_j(y)=\prod_{a=1}^n\chi(y_a/j)\,
           \Psi(y_1-v/j,\ldots,y_n-v/j).
\]
Each $\Psi_j$ has compact support in
$O^n$, $O=\{x:v\cdot x>c\}$. Translation is continuous in Schwartz
space, and the growing cutoff tends to one in its action on each
Schwartz function. More explicitly, Leibniz' rule gives
\[
 q_m(\Psi_j-\Psi)\le C_{m,n,\chi}j^{-1}q_{m+1}(\Psi).
\]
Consequently $Q_{\Psi_j}\to Q_\Psi$ in $L^2(\lambda)$, and
\begin{equation}\label{eq:tensor-closed-halfspace}
 Q_\Psi\in L^2(\Sigma_O,\lambda)\qquad
                   (\supp\Psi\subset H^n).
\end{equation}
Thus tests meeting the boundary plane belong to the same
positive-side $L^2$ space.

\subsection{The five conditions}

\begin{proposition}\label{prop:tensor-OS}
Let $\lambda$ be a Gaussian-dominated law of real tempered two-forms.
Suppose it is invariant under Euclidean motions and reflection
positive at every plane through the origin. Then its tensor
Schwinger functions satisfy \textup{(E0$'$)}--\textup{(E3)} of
Theorem~\ref{thm:main}. If, for every $A,B\in L^2(\lambda)$,
\[
 \E_\lambda[\overline A\,U(a)B]\longrightarrow
       \overline{\E_\lambda A}\,\E_\lambda B
       \quad\text{as }|a|\longrightarrow\infty,
\]
they also satisfy \textup{(E4)}.
\end{proposition}

\begin{proof}
The Hermite estimate proves \textup{(E0$'$)}. Covariance follows from
\eqref{eq:tensor-equivariance} and invariance of expectation.
Permutation symmetry holds on elementary tensors because the random
field evaluations commute; continuity gives \textup{(E3)} on all
tensor tests.

For \textup{(E2)}, let $r$ be the stated reflection and $H$ the
chosen closed half-space. By \eqref{eq:tensor-closed-halfspace},
$A=\sum_{k=0}^MQ_{\Psi_k}$ lies in its positive-side $L^2$ space,
where $Q_{\Psi_0}=\Psi_0$. Reflection positivity extends from
bounded characters to this whole space by their density and the
boundedness of the reflection form. Equations
\eqref{eq:tensor-product-rule}--\eqref{eq:tensor-equivariance} give
\[
 \overline{\mathcal R_rQ_{\Psi_k}}
       =Q_{\Theta_k\Psi_k},\qquad
 Q_{\Theta_k\Psi_k}Q_{\Psi_l}
       =Q_{\Theta_k\Psi_k\otimes\Psi_l}\quad\text{in }L^1.
\]
Therefore
\[
 0\le\E_\lambda[(\overline{\mathcal R_rA})A]
   =\sum_{k,l=0}^MS_{k+l}(\Theta_k\Psi_k\otimes\Psi_l).
\]
This is the tensor reflection form with the conjugation and signs
in the first-page statement.

Finally, \eqref{eq:tensor-product-rule} and translation covariance
identify the expression in \textup{(E4)} with
$\E_\lambda[Q_\Psi U(a)Q_\Phi]$. Both random polynomials are in
$L^2(\lambda)$. Apply the stated clustering with
$A=\overline{Q_\Psi}$ and $B=Q_\Phi$ to obtain
$\E_\lambda Q_\Psi\,\E_\lambda Q_\Phi$.
\end{proof}

For the component laws constructed above, qualitative clustering
has already been obtained from the cone and the invariant projections.
Once rotation invariance is proved for each component, the nine-mirror
reflection form is transported to every plane, and the proposition
applies.


\section{Completion of the field-strength construction}
\label{sec:main-theorem-closure}

\begin{proof}[Proof of Theorem~\ref{thm:main}]
Fix the construction constants of
Section~\ref{sec:ordered-constants}, in their stated order.
In particular $\sigma>0$ in
\eqref{eq:transfer-charge-exponents}.
Take the positive physical ceiling $e_G$ in
\eqref{eq:constants-physical-charge}. Theorem
\ref{thm:rgfl-common-pair} supplies real bare mass shifts
$m_{0,N}(e)$, independent of the volume and twist,
for $0<e\le e_G$; below its defining mesh threshold they
are set to zero. Their bound
\eqref{eq:rgfl-bare-mass-bound} gives
$|m_{0,N}(e)|<1/8$ for every $N$.
The all-flux laws at \eqref{eq:cutoff-scales} are therefore
defined at every cutoff, even when the total lattice mass
has either sign.
The physical ceiling in the first-page statement is
$e_{\mathrm{phys}}=e_G$.
Proposition~\ref{prop:positive-probe} proves the existence
of the probes in \eqref{eq:positive-coefficient}.
Fix any such probe and $0<e\le e_\beta$, with
$e_\beta$ as in \eqref{eq:constants-probe-charge}.

For this probe,
Proposition~\ref{prop:interacting-positive-threshold}
provides $e_\sharp(\beta)>0$.
The terminal partition bound
\eqref{eq:term-matter-comparison} and the response estimate
\eqref{eq:interacting-final-terminal-error} supply
both parts of \eqref{eq:transfer-terminal-input}.
The free-coefficient bound
\eqref{eq:transfer-free-input} supplies its remaining
reference input. Thus Proposition~\ref{prop:transfer-UV-inputs}
applies throughout this interval, with one $C_U$
independent of the cutoff, volume, placement and limit law.
Its estimates hold whenever
$N\ge N_{\rm ref}(e)$ and $T\ge T^*(e)+1$.
These are finite thresholds for each fixed $e$.

\emph{Volume-first laws and their Euclidean components.}
The tightness estimate \eqref{eq:limit-uniform-tightness}
shows that every fixed-spacing volume-limit set
$\mathcal V_N$ is nonempty and that the collection of
these limits is tight on $H_{-4}$.
Consequently volume-first limit points exist.
Let $\mu$ be any one of them, and write
$\mu^{(N_i)}\to\mu$, where $N_i\to\infty$ and
$\mu^{(N_i)}\in\mathcal V_{N_i}$.
For each $i$, choose a volume index $T_i$ along a
defining subsequence for $\mu^{(N_i)}$ such that
\[
 T_i\ge i,\qquad
 d_{\rm w}(P_{N_i,T_i},\mu^{(N_i)})\le1/i,
\]
where $d_{\rm w}$ metrizes weak convergence on $H_{-4}$.
Then
$d_{\rm w}(P_{N_i,T_i},\mu)
\le1/i+d_{\rm w}(\mu^{(N_i)},\mu)\to0$.
This is the joint sequence in
Proposition~\ref{prop:limit-nine-mirrors};
both indices diverge, and all fixed finite-family
moments converge along it.

Proposition~\ref{prop:limit-nine-mirrors} gives closedness,
translation invariance, Gaussian domination with constant
one, and invariance and reflection positivity at all
planes in the nine mirror families.
Section~\ref{subsec:component-hypotheses} transfers
these properties to a measurable set of full measure
in the translation-ergodic decomposition.
The Gaussian bound holds there for all real Schwartz
tests at once, by the countable-density and Fatou
argument in \eqref{eq:component-G}.
The same section transfers each additional mirror
by transporting the conditional laws and using
uniqueness of the decomposition. Intersecting the
nine full-measure sets gives all mirror forms in
each component; these reflections also generate
the cubic group.

Theorem~\ref{thm:rotation-covariance} now applies
separately to $\mu$ and to every component in this
full-measure set. It proves their invariance under
the full Euclidean group and reflection positivity
at every plane.
For an ergodic component, the sharp cone and
\eqref{eq:invariant-projections-equal} identify the
translation-invariant projection with the constant
projection. Proposition~\ref{prop:all-direction-mixing}
therefore gives clustering for all $L^2$ observables
as $|a|\to\infty$.
Proposition~\ref{prop:tensor-OS} gives
\textup{(E0$'$)}--\textup{(E3)} for the parent law and
\textup{(E0$'$)}--\textup{(E4)} for these components,
with exactly the tensor action and growth bound of
the first-page statement.
In particular, supported tensor evaluations belong
to the positive-side $L^2$ space even when their
supports meet the reflection plane, by
\eqref{eq:tensor-closed-halfspace}.
For clustering the relevant identity is
\[
 S_{k+l}^\xi(\Psi\otimes\tau_a^{\otimes l}\Phi)
   =\E_{\mu_\xi}[Q_\Psi\,U(a)Q_\Phi],
                 \qquad Q_\Psi,Q_\Phi\in L^2(\mu_\xi).
\]
It gives \eqref{eq:main-cluster} on all tensor Schwartz
tests, using \(A=\overline{Q_\Psi}\) in the
$L^2$ clustering formula.

\emph{Fourth cumulants in the same components.}
Fix $e$ in the interval above and put
$\lambda=R^*(e)/L$.
This spacing of translated cells is independent of $N$.
Along the joint sequence just chosen, the reference
numbers $c_{2,i},c_{4,i},u_{2,i},u_{4,i}$ are bounded
and depend on $N_i$ and the chosen probe placement,
rather than on the outer volume or on a set of
translated cells.
Pass to one subsequence on which these four numbers
converge. Its weak limit is still $\mu$.

For any finite $Y\subset\lambda\Z^3$, all large $i$
satisfy the two cutoff thresholds above.
They also satisfy
$R_i>2(\max_{y\in Y}|y|_\infty+r_f)$,
where a fixed cube $[-r_f,r_f]^3$ contains
the supports of $f$ and $f^{\eta_i}$.
Thus the points of $Y$ are incongruent modulo
$R_i\Z^3$, and all the translated tests lie in
the interior of one period cube.
Equation~\eqref{eq:transfer-translated-Wick} applies:
\[
 \left|\E_{N_i,T_i}
       \prod_{y\in Y}
       W_m(F^\sharp(\tau_yf^{\eta_i});v_i,c_{m,i})\right|
                         \le u_{m,i}^{|Y|},\qquad m=2,4.
\]
For each fixed $Y$, expansion into monomials and the
test replacement estimate in
Section~\ref{sec:ultraviolet-limits} pass this inequality
to \eqref{eq:UV-law-product-bound}.
More explicitly, the replacement error for a monomial
of degree $J>0$ is at most
$Jc_JB^{J-1}\|f^{\eta_i}-f\|_2\to0$; the remaining
fixed-test moments converge along the joint sequence.
The variances and polynomial coefficients converge
as well. This works on the same coefficient subsequence
for every finite $Y$ in the countable lattice:
the index beyond which support containment holds
may depend on $Y$, while the limiting coefficients
and constants are common.

The limits obey \eqref{eq:UV-reference-limit}.
Lemma~\ref{lem:UV-invariant-product-bound}, followed by
\eqref{eq:UV-component-product-bound}, gives
$|d_m|\le u_m$ almost surely in the component parameter,
where
$d_m=\E_{\mu_\xi}W_m(F(f);\|f\|_2^2,c_m)$.
Intersect this set with the full-measure set used
for the OS conditions.
Proposition~\ref{prop:UV-component-cumulant} applies
there. Its exact identity is
\[
 e^{-4}\kappa_4^{\mu_\xi}(f)
                        =c_4+d_4-3(c_2+d_2)^2,
\]
and gives
$|\kappa_4^{\mu_\xi}(f)-e^4b(f)|
\le C_f e^{4+\sigma}$ with
$C_f=8C_U+3C_U^2$.
The coefficient limits may depend on the defining
subsequence; this constant does not.
Finally \eqref{eq:constants-probe-charge} gives
$C_fe^\sigma\le b(f)/2$ throughout the interval.
Then
$\kappa_4^{\mu_\xi}(f)\ge e^4b(f)/2>0$.
This proves \eqref{eq:main-cumulant} uniformly in
the chosen limit law and almost every component.
A centered Gaussian variable has zero fourth
cumulant, so every component in this full-measure
set is non-Gaussian.
\end{proof}


\appendix

\section{Polymer geometry}\label{app:polymer-geometry}

Let $L$ be odd, $M$ a power of $L$, and $\ell$ a multiple of $LM$
with $\ell\ge13LM$. The $M$-cubes have centres in
$M\Z^3/\ell\Z^3$. They partition the lattice sites; their closed
continuum versions may meet along faces and edges. A polymer is a
nonempty union of $M$-cubes connected through faces. Denote its number
of cubes by $|X|_M$. A net in the continuum torus is a connected
finite union of images of specified closed straight segments
in $\R^3$, a point being allowed as a segment.
Its length is the sum of the Euclidean lengths of those segments.
Set
\[
 Md_M(X)=\inf\{\operatorname{length}(N):N
             \text{ is a net meeting every closed cube of }X\},
\]
where the infimum also extends over the finite segment representations.
The same definition is used for finite face-connected unions of
$M$-cubes in $\R^3$, with nets in $\R^3$. Distances are sup distances.

\subsection{Size, diameter, and reblocking}

Put
\[
 c_X=54,\qquad \kappa_0=54\log72,\qquad
 K_0=e^{\kappa_0}/36,\qquad c_{\mathrm{rb}}=c_XK_0.
\]
The following estimates are uniform in $M$ and $\ell$:
\begin{align}
 d_M(X)&\le |X|_M-1,&
 |X|_M&\le c_X(d_M(X)+1),\label{eq:polymer-size-tree}\\
 d(p,q)&\le M(d_M(X)+2)
       &&(p,q\text{ in the closed cubes of }X),
                                  \label{eq:polymer-diameter}\\
 \sum_{X\supset\square}e^{-\kappa_0d_M(X)}&\le K_0,&
 \sum_{X\subset Y}e^{-\kappa_0d_M(X)}&\le K_0|Y|_M.
                                  \label{eq:polymer-counting}
\end{align}
Here $Y$ is any union of $M$-cubes. If $\ell/L$ satisfies the same
side conditions stated above, let $\bar X^L$ be the union of $LM$-cubes containing cubes of
$X$. There is a unique polymer $Y$ on the rescaled torus with
$\bar X^L=LY$, and
\begin{align}
 Ld_M(Y)&\le d_M(X),\label{eq:polymer-reblock-tree}\\
 \sum_{X:\bar X^L=LY}e^{-\kappa'd_M(X)}
      &\le c_{\mathrm{rb}}L^3
           e^{-L(\kappa'-\kappa_0-1)d_M(Y)}
                     \quad(\kappa'\ge\kappa_0).
                                  \label{eq:polymer-reblock-sum}
\end{align}

To prove the first inequality of \eqref{eq:polymer-size-tree}, order
the cubes along a spanning tree of their adjacency graph. A net
meeting one cube can be extended to an adjacent cube by a segment of
length at most $M$ perpendicular to their common face. Starting with
one point and repeating this construction gives a net of length at
most $M(|X|_M-1)$.

For the reverse inequality, subdivide the segments of any net $N$ at
their intersections and discard overlapping pieces. Its total length
does not increase. Doubling the edges of the resulting connected
graph gives a closed traversal of length at most
$2\operatorname{length}(N)$. Cut it into at most
$2\operatorname{length}(N)/M+1$ arcs of length at most $M$.
Each arc lifts to a curve in $\R^3$ whose range in each coordinate
has length at most $M$. Such an interval meets at most three closed
cube intervals, so an arc meets at most $27$ closed $M$-cubes.
Every cube of $X$ meets one of these arcs. Taking the infimum over
$N$ gives
\[
 |X|_M\le27(2d_M(X)+1)\le c_X(d_M(X)+1).
\]
For \eqref{eq:polymer-diameter}, take a net of length less than
$Md_M(X)+\delta$ and points $p',q'$ of the net in the respective
cubes of $p,q$. The path in the net from $p'$ to $q'$ has length
at most that of the net. Hence
$d(p,q)\le M+(Md_M(X)+\delta)+M$; let $\delta\downarrow0$.
The same argument holds in $\R^3$.

A polymer with $n$ cubes containing a fixed cube determines a closed
walk of length $2(n-1)$ starting there: double a spanning tree of its
adjacency graph and traverse it. Choose one such walk for each
polymer. This assignment is injective, since the visited cubes recover
the polymer, and there are at most $6^{2(n-1)}=36^{n-1}$ possible
walks. By the size estimate, $d_M(X)\ge n/c_X-1$. Summing over $n$
gives
\[
 \sum_{X\supset\square}e^{-\kappa_0d_M(X)}
       \le e^{\kappa_0}\sum_{n\ge1}36^{n-1}72^{-n}=K_0.
\]
Sum this bound over the cubes of $Y$ to obtain the second inequality
of \eqref{eq:polymer-counting}.

Every $M$-cube belongs to one $LM$-cube, since $L$ is odd, and adjacent
$M$-cubes belong to the same or to adjacent $LM$-cubes. Thus $\bar X^L$
is connected through faces. Projection by $c\mapsto c/L$ gives its
unique rescaled polymer $Y$. A net for $X$ meets every $LM$-cube of
$\bar X^L$; its rescaling has length divided by $L$. This proves
\eqref{eq:polymer-reblock-tree}. For \eqref{eq:polymer-reblock-sum},
split the decay into rates $\kappa_0$ and $\kappa'-\kappa_0$ and use
\eqref{eq:polymer-reblock-tree}. All terms satisfy $X\subset LY$,
which has $L^3|Y|_M$ fine cubes, so
\[
 \sum_{X:\bar X^L=LY}e^{-\kappa_0d_M(X)}
      \le K_0L^3|Y|_M\le c_XK_0L^3e^{d_M(Y)}.
\]
Combining the two factors proves the stated bound.

\subsection{Nonwrapping lifts}\label{subsec:polymer-lifts}

Let $\operatorname{pr}:\R^3\to\R^3/\ell\Z^3$ be the projection.
\begin{proposition}\label{prop:nonwrapping-polymer-lift}
If $X$ is a polymer and
\begin{equation}\label{eq:nonwrapping-polymer-condition}
 2M(d_M(X)+2)<\ell,
\end{equation}
then $X=\operatorname{pr}(\widetilde X)$ for a finite face-connected
union of $M$-cubes in $\R^3$. The projection is injective on
$\widetilde X$, preserves adjacency of its cubes in both directions,
and $d_M(\widetilde X)=d_M(X)$. The lift is unique up to a translation
in $\ell\Z^3$.
\end{proposition}

\begin{proof}
Put $D=M(d_M(X)+2)$, fix $p_0$ in a closed cube of $X$, and fix its
lift $\widetilde p_0$. The projection is injective on the closed sup
ball $\widetilde B=\{|p-\widetilde p_0|_\infty\le D\}$, since two
points there differ in every coordinate by at most $2D<\ell$.
It maps this ball onto the torus ball $B=\{p:d(p,p_0)\le D\}$.
Its inverse $\varphi:B\to\widetilde B$ is continuous, and is locally
isometric. The diameter estimate places every closed cube of $X$ in
$B$. A segment or a closed $M$-cube lying in $B$ lifts under
$\varphi$ to a segment of the same length or a closed $M$-cube:
the translates $\widetilde B+\ell n$ are pairwise disjoint, and a
connected lift lies in one of them. The same inverse maps shared
faces to shared faces. Thus $\widetilde X=\varphi(X)$ has the
asserted injectivity, adjacency, and connectedness.

Any other face-connected injective lift can be translated so that it
contains the chosen lift of the cube of $p_0$. There are no two
different neighbours of one lifted cube with the same projection:
their centres differ by at most $2M<\ell$. Propagate this observation
along a chain of adjacent cubes. Every cube of the other lift equals
the corresponding cube of $\widetilde X$, proving uniqueness.

Projecting a net for $\widetilde X$ gives a net for $X$ with no greater
length, so $d_M(X)\le d_M(\widetilde X)$. Conversely, let $N$ be a
net for $X$ of length at most $Md_M(X)+M$, and choose its point $q_0$
in the cube of $p_0$. Every point of $N$ is joined to $q_0$ by a
curve of length at most that of $N$, and $d(q_0,p_0)\le M$.
Therefore $N\subset B$. Its lift $\varphi(N)$ has the same length
and meets every cube of $\widetilde X$. Taking an infimizing sequence
of these nets proves the reverse inequality.
\end{proof}

For the determinant estimate, suppose in addition that
$M(d_M(X)+4)<\ell$ and $M\ge2$. The closed cubes of $\widetilde X$
have coordinate ranges of length at most $M(d_M(X)+2)$ by
\eqref{eq:polymer-diameter}. Enlarge these intervals by one mesh
$\eta_k\le1$ at each end and intersect with $\eta_k\Z^3$. The
resulting box $W$ contains every bond meeting the lifted polymer and
has each side of length at most $M(d_M(X)+3)$. This length, increased
by one mesh, is less than $\ell$. Thus the projection is injective
on $W$, and every torus bond meeting $X$ has exactly one lift in $W$.

\begin{thebibliography}{10}
\bibitem{OS1973}
K.~Osterwalder and R.~Schrader,
\emph{Axioms for Euclidean Green's functions},
Commun.\ Math.\ Phys.\ \textbf{31} (1973), 83--112.
\url{https://doi.org/10.1007/BF01645738}.

\bibitem{OS1975}
K.~Osterwalder and R.~Schrader,
\emph{Axioms for Euclidean Green's functions II},
Commun.\ Math.\ Phys.\ \textbf{42} (1975), 281--305.
\url{https://doi.org/10.1007/BF01608978}.

\bibitem{OsterwalderSeiler1978}
K.~Osterwalder and E.~Seiler,
\emph{Gauge field theories on a lattice},
Ann.\ Phys.\ \textbf{110} (1978), 440--471.
\url{https://doi.org/10.1016/0003-4916(78)90039-8}.

\bibitem{FILS1978}
J.~Fr\"ohlich, R.~Israel, E.~H.~Lieb, and B.~Simon,
\emph{Phase transitions and reflection positivity.
I. General theory and long range lattice models},
Commun.\ Math.\ Phys.\ \textbf{62} (1978), 1--34.
\url{https://doi.org/10.1007/BF01940327}.

\bibitem{DimockUVS2022}
J.~Dimock, \emph{Ultraviolet stability for QED in $d=3$},
Ann.\ Henri Poincar\'e \textbf{23} (2022), 2113--2205.
\url{https://doi.org/10.1007/s00023-021-01127-z}.

\bibitem{EtingofLieNotes} P. Etingof, \emph{Lie Groups and Lie Algebras},
MIT course notes, \S22.3, Theorem~22.14.
\url{https://ocw.mit.edu/courses/18-755-lie-groups-and-lie-algebras-ii-spring-2024/mit18_755_s24_lec_full.pdf}.
\bibitem{Hatcher2002} A. Hatcher, \emph{Algebraic Topology},
Cambridge University Press, 2002.
\url{https://math.uchicago.edu/~farb/algtop2023/Hatcher.pdf}.
\bibitem{PostnikovLin2019} A. Postnikov, \emph{Algebraic Combinatorics},
MIT 18.212, Lecture 23, notes by A. Lin, April 5, 2019,
Proposition 4.
\url{https://ocw.mit.edu/courses/18-212-algebraic-combinatorics-spring-2019/f6113618a3e085ba86b799f65934886d_MIT18_212S19_lec23.pdf}.
\bibitem{BrightMetric2023} P. Bright, \emph{Introduction to Metric Spaces},
MIT 18.S190, Lecture 5, January 24, 2023, Theorem 8.
\url{https://ocw.mit.edu/courses/18-s190-introduction-to-metric-spaces-january-iap-2023/mit18_s190iap23_lec5.pdf}.
\bibitem{PongeSchatten2025} R. Ponge, \emph{Singular Values and Schatten Classes},
lecture notes, 2025, slides 3--9 and 53--56.
\url{https://raphaelponge.org/wp-content/uploads/2025/07/summer-chap4-1.pdf}.
\end{thebibliography}
\end{document}
